\documentclass[11pt]{article}

\usepackage[margin=1in]{geometry}
\usepackage{amsmath,amssymb,amsthm,mathtools}
\usepackage{bm}
\usepackage{microtype}
\usepackage{enumitem}
\usepackage{booktabs}
\usepackage{array}
\usepackage{xcolor}
\usepackage{xr-hyper}
\usepackage{hyperref}
\usepackage[nameinlink,noabbrev]{cleveref}

\providecommand{\Fsloc}{\mathcal F_{\rm s.loc}}
\providecommand{\Qt}{Q_t}
\externaldocument[main-]{UOM/main}
\externaldocument[uomdesc-]{UOM/uom_descendant_note}

\hypersetup{
  colorlinks=true,
  linkcolor=blue!60!black,
  citecolor=blue!60!black,
  urlcolor=blue!60!black
}

% ---------- Theorem environments ----------
\theoremstyle{definition}
\newtheorem{definition}{Definition}[section]
\newtheorem{assumption}{Assumption}[section]
\theoremstyle{plain}
\newtheorem{theorem}[definition]{Theorem}
\newtheorem{lemma}[definition]{Lemma}
\newtheorem{proposition}[definition]{Proposition}
\newtheorem{corollary}[definition]{Corollary}
\theoremstyle{remark}
\newtheorem{remark}[definition]{Remark}
\crefname{definition}{definition}{definitions}
\Crefname{definition}{Definition}{Definitions}
\crefname{assumption}{assumption}{assumptions}
\Crefname{assumption}{Assumption}{Assumptions}
\crefname{theorem}{theorem}{theorems}
\Crefname{theorem}{Theorem}{Theorems}
\crefname{lemma}{lemma}{lemmas}
\Crefname{lemma}{Lemma}{Lemmas}
\crefname{proposition}{proposition}{propositions}
\Crefname{proposition}{Proposition}{Propositions}
\crefname{corollary}{corollary}{corollaries}
\Crefname{corollary}{Corollary}{Corollaries}
\crefname{remark}{remark}{remarks}
\Crefname{remark}{Remark}{Remarks}

% ---------- Notation ----------
\newcommand{\cH}{\mathcal H}
\newcommand{\cA}{\mathcal A}
\newcommand{\cD}{\mathcal D}
\newcommand{\cB}{\mathcal B}
\newcommand{\cC}{\mathcal C}
\newcommand{\cL}{\mathcal L}
\newcommand{\cK}{\mathcal K}
\newcommand{\cR}{\mathcal R}
\newcommand{\cG}{\mathcal G}
\newcommand{\Ztwo}{\mathbb Z_2}
\newcommand{\Sys}{\mathrm{Sys}}
\newcommand{\Time}{\mathrm{Time}}
\newcommand{\RG}{\mathrm{RG}}
\newcommand{\Tr}{\mathrm{Tr}}
\newcommand{\diag}{\mathrm{diag}}
\newcommand{\ad}{\mathrm{ad}}
\newcommand{\id}{\mathrm{id}}
\newcommand{\norm}[1]{\left\lVert #1\right\rVert}
\newcommand{\abs}[1]{\left\lvert #1\right\rvert}
\newcommand{\ip}[2]{\left\langle #1,#2\right\rangle}

\title{Nontrivial Summability and the Split in UOM and 2T-LS:\\
Completion-Sector Leak, Classicality, and Theoretical Derivation}
\author{Internal technical note}
\date{\today}

\begin{document}
\maketitle

\begin{abstract}
This note studies two coupled structures in Universe-Originating Measurement (UOM):
\emph{(i)} the nontrivial summability of the Brown--York (BY) increments $\{\Theta_k\}$ and
\emph{(ii)} the split as exposure of a completion sector into the next aeon's source.
The leak term in the SK/HJ wall balance is formulated as a completion-sector export flux
controlled by a signed odd split amplitude and carried by hidden Stinespring/commutant degrees of
freedom.
The note derives the continued-sector RP coefficient \(\alpha(\Lambda)\), a parity-based rank-one
wall-match normal form, a projected Gaussian affine realization of the odd split mode, a
centered/squeeze-dominated specialization of that affine dynamics, a positive decomposition of the
split-driving leak into fresh-injection and transport pieces, a low-rank projected stock law, and
an exported characterization interface for the canonical odd representative used downstream in
\(\mathbf{2T}_{\mathrm{Main}}\).
Where UOM does not determine a unique completion Hamiltonian, export projector, or one-wall/two-wall
matching, the corresponding steps are left as explicit \emph{hypotheses}.
Classicality is assigned to the coarse-grained accessible ledger, whereas RP/OS oddness is assigned
to the analytic-continued completion sector.
\end{abstract}

\tableofcontents

\section{Scope and inputs}
\label{sec:scope}

\paragraph{Primary sources (internal).}
We take as authoritative the UOM main manuscript and the 2T-LS rigorous translation note.
The simulators are used only as \emph{phenomenological inputs}: they probe reduced scalar closures of the wall balance and
help identify which qualitative leak behaviors do or do not produce split.
They are not treated as the primary source of theoretical definitions.

\paragraph{What ``nontrivial summability'' means here.}
There are two levels:
\begin{enumerate}[leftmargin=2em]
\item \textbf{Within one aeon:} beyond the trivial telescoping identity that follows from the bookkeeping update,
the observed phenomenon is that the \emph{implicit one-wall $\Lambda$-solve} can dynamically adjust the cutoff trajectory
so as to avoid overshoot \emph{at all finite ticks} while still exhausting the curvature budget:
\[
0<\Theta_k<R^{\rm abs}_k\ \text{for all finite }k,
\qquad
R^{\rm abs}_k\downarrow 0,
\qquad
\sum_{k\ge 0}\Theta_k=\abs{R_0}.
\]
This is what we call \emph{nontrivial (exact) summability} in this note; the ``flip'' is an asymptotic boundary point of
the solvable one-wall channel rather than a finite threshold crossing.
\item \textbf{Across aeons:} in the exhaustion/terminal branch the total available wall-work is summable (UOM Theorem
\texttt{thm:ads5-static}), yielding a terminal cutoff $\Lambda_\infty$ and decoupled tail behavior.
\end{enumerate}

\paragraph{What ``split'' means here.}
Operationally, a split is the end-of-aeon event in which the construction produces a \emph{new} dS source branch and a
\emph{new} AdS receiver branch. In UOM this is treated as compatible with BY charge conservation across a pillbox around
the split region (UOM App. \texttt{app:energy-split}). In 2T-LS, the split/continuation has a mathematically clean
signature: a parity cocycle $\omega$ whose nontrivial holonomy obstructs a global ``splitting'' (sectioning) of a
$\Ztwo$-cover (2T-LS Theorem \texttt{thm:non-splitting}).

\section{Mathematical objects: states, channels, and per-tick scalars}
\label{sec:objects}

\subsection{State spaces and admissible coarse views}
\label{subsec:state-spaces}

Let $\cH_{\rm tot}$ denote the total microscopic Hilbert space of the composite (source patches, commutant/bath bookkeeping,
receiver degrees of freedom, and any bulk/shadow auxiliaries retained).
We write $\cD(\cH)$ for density matrices on $\cH$.

UOM frequently avoids assuming a literal tensor factorization for ``accessible'' degrees of freedom and instead uses
operator-algebra restrictions and conditional expectations. A key UOM definition (label
\texttt{def:eta-bridge-leak}) introduces:
\begin{itemize}[leftmargin=1.6em]
\item a conditional expectation $\mathbb E_{\rm ptr}$ onto a pointer algebra, and
\item a commutant record map $\cR_k$ with commutative range (classical record algebra),
\end{itemize}
together with a bridge/leak instrument decomposition $\Phi_k=\cB_k+\cL_k$ of the tick CPTP map.

\subsection{Source-side QND dephaser and exported payload $\Delta I_k$}
\label{subsec:qnd-payload}

The source-side mechanism is a pointer-aligned QND absorption channel with generator of the form (UOM label
\texttt{lem:exp-tail}):
\begin{equation}
\cL(t)\rho \;=\; -\frac12 \sum_a \gamma_a(t)\,[L_a,[L_a,\rho]],
\qquad L_a=f_a(H_{{\rm sp},\Delta}),\qquad [L_a,H_{{\rm sp},\Delta}]=0.
\label{eq:qnd-generator}
\end{equation}
Across a tick window $I_k=[t_-,t_+]$ in a steady-actuator regime one obtains exponential contraction of the modular
free energy $F(t)=S(\rho_t\|\sigma_\Delta)$ relative to a band-limited BD/KMS anchor $\sigma_\Delta$:
\begin{equation}
F(t_+) \le e^{-\beta_k}F(t_-),
\qquad
\Delta I_k := F(t_-)-F(t_+)\ge (1-e^{-\beta_k})F(t_-).
\label{eq:payload-lower-bound}
\end{equation}
In reduced scalar closures this appears as:
\begin{equation}
\Delta I_k \approx (1-e^{-\beta_k})\,F_k,
\qquad
F_{k+1}=F_k-\Delta I_k,
\label{eq:sim-payload-update}
\end{equation}
where $F_k$ is represented by the scalar source stock \texttt{F\_dS} in \texttt{aeon\_driver.py}, and where
$\beta_k$ is computed from a schedule factor times the source Hubble scale and (in v1) a band-fraction modifier.

\paragraph{Matrix-element form (for later classicality bounds).}
The closed form of the dephasing channel on the band is (UOM label \texttt{eq:app-phi-closed})
\begin{equation}
\Phi_\Delta^{(k)}=\exp\!\Big\{-\tfrac12\sum_a \Gamma_a^{(k)}\,\ad_{L_a}^{\,2}\Big\},
\qquad
\Gamma_a^{(k)}=\int_{\mathbb R}\gamma_a^{(k)}(t)\,dt.
\label{eq:qnd-channel-closed}
\end{equation}
In a common eigenbasis of $\{L_a\}$, $\ad_{L_a}^{\,2}$ acts diagonally on matrix elements, yielding suppression of
off-diagonals by factors $\exp(-\Gamma^{(k)}_{ij})$ with
$\Gamma^{(k)}_{ij}=\tfrac12\sum_a\Gamma_a^{(k)}(\ell_a(i)-\ell_a(j))^2$.

\subsection{Receiver-side BY rectifier increment $\Theta_k$}
\label{subsec:theta}

UOM defines the BY anomaly/rectifier increment (UOM label \texttt{eq:theta-def}) as
\begin{equation}
\Theta_k
:=\gamma\,\Delta\!\langle R^2\rangle_k
\;=\;\underbrace{\frac{aG_4}{12\pi}}_{=\gamma}\,(8\pi G_4)^2\,
\big\langle|\cK|^2\big\rangle_k\,
C_{\rm LoG}\;\frac{A_k^2}{\sigma_{t,k}\,\sigma_{x,k}^3}.
\label{eq:theta-def-note}
\end{equation}
Equivalently, in the welded/throughput form (UOM label \texttt{eq:R2}) one writes a per-tick BY penalty
\begin{equation}
\Phi_{{\rm BY},k}
=
\gamma(8\pi G_4)^2\,C_{\rm LoG}\;
\frac{\Delta I_k}{s_{t,k}\,s_{x,k}^{3}\,G(s_{t,k},s_{x,k};\Lambda_k)}.
\label{eq:by-throughput-note}
\end{equation}
In reduced phenomenological closures, $\Theta_k$ is evaluated by a discretized Euclidean BY kernel integrating over a
frequency grid and an $S^3$ harmonic grid, with:
\begin{itemize}[leftmargin=1.6em]
\item welded widths $s_t(\Lambda)=\alpha_t/\Lambda$ and $s_x(\Lambda)=\alpha_x/(v_{\rm geom}\Lambda)$,
\item an evanescent suppression across the brane gap $d_{\rm sep}$ as $\exp(-d_{\rm sep}q)$ with
$q=\sqrt{\omega^2+k^2+m_g^2}$,
\item screening $K=q/(q+r(\Lambda)k^2)$ with $r(\Lambda)\simeq \ell_{\rm AdS5}/2 + C_{m3}\Lambda^2$,
\item degeneracy weights $D_\ell=(\ell+1)^2$.
\end{itemize}
\begin{remark}[Linearity in $\Delta I_k$ and the effective kernel $K(\Lambda)$]
Both the manuscript and the code implement $\Theta_k$ as \emph{linear} in the exported payload $\Delta I_k$ when the welded
shape family is fixed: the reduced closure chooses an amplitude $A_{\rm eff}$ so that (heuristically) ``signal power'' equals
$\Delta I_k$, and then integrates a quadratic response which is proportional to $A_{\rm eff}^2\propto \Delta I_k$.
Thus one can write
\begin{equation}
\Theta_k = K(\bar\Lambda_k;d_k,L_{4,k},\text{locks})\,\Delta I_k,
\label{eq:theta-factorization}
\end{equation}
for a (state-dependent) kernel factor $K$ extracted numerically as $\Theta_k/\Delta I_k$.
This factorization is explicit in UOM's receiver evolution summary (label \texttt{par:uom-main-138}).
\end{remark}

\section{Nontrivial summability within an aeon}
\label{sec:summability}

\subsection{Curvature bookkeeping and exact summability}
\label{subsec:curvature-bookkeeping}

UOM's receiver-side coarse evolution tracks a signed curvature variable (AdS negative, dS positive) or equivalently
its magnitude $\abs{R}$.
The simulators encode the update in one of two equivalent ways:
\begin{itemize}[leftmargin=1.6em]
\item \textbf{v1 (magnitude bookkeeping):} $R^{\rm abs}_{k+1}=R^{\rm abs}_k-\Theta_k$ with $R^{\rm abs}_0=\abs{R_0}$.
\item \textbf{v2 (signed bookkeeping):} $R_{k+1}=R_k+\Theta_k$ with $R_0=-\abs{R_0}<0$.
\end{itemize}
In either case, $\Theta_k\ge 0$ drives the receiver toward the flip.

\begin{remark}[Legacy ``first-crossing'' convention]
For comparison with UOM's manuscript language, one sometimes declares a flip at the first $N$ such that
$S_N:=\sum_{k=0}^N\Theta_k\ge \abs{R_0}$.
In the exact-summability regimes observed in the implicit $\Lambda$-solver experiments (no-leak and fixed-work-percentage
(i.e.\ proportional) leak models), one instead finds $S_N<\abs{R_0}$ for all finite $N$ but
$\lim_{N\to\infty}S_N=\abs{R_0}$.
Thus there is no finite ``flip tick'': the flip is realized as an asymptotic boundary point of the solvable one-wall
channel.
\end{remark}

\begin{proposition}[Telescoping identity (exact curvature-budget summation)]
Assume the within-aeon update is $R^{\rm abs}_{k+1}=R^{\rm abs}_k-\Theta_k$ without additional clamps and that
$\Theta_k\ge 0$.
Then for any $N$,
\begin{equation}
\sum_{k=0}^N \Theta_k = R^{\rm abs}_0 - R^{\rm abs}_{N+1}.
\label{eq:telescoping}
\end{equation}
In particular, if the update runs indefinitely with $R^{\rm abs}_{N+1}\downarrow 0$ as $N\to\infty$ (no overshoot /
no sign crossing), then the series is \emph{exactly} summable:
\begin{equation}
\sum_{k=0}^{\infty}\Theta_k \;=\; \abs{R_0}.
\label{eq:exact-summability}
\end{equation}
\end{proposition}

\begin{proof}[Derivation and mechanism]
\textbf{(i) Exact telescoping.}
Summing the within-aeon update $R^{\rm abs}_{k+1}=R^{\rm abs}_k-\Theta_k$ from $k=0$ to $N$ gives
\[
R^{\rm abs}_{N+1}=R^{\rm abs}_0-\sum_{k=0}^N\Theta_k,
\]
which is \eqref{eq:telescoping}.
This is not an approximation: it is the exact discrete bookkeeping of the receiver curvature budget used by the reduced
one-wall closures.

\textbf{(ii) Why the \emph{no-leak solvable branch} produces $R^{\rm abs}_k\downarrow 0$ (hence exact summability).}
The nontrivial experimental observation is not \eqref{eq:telescoping} but that, in the no-leak and
fixed-work-percentage-leak regimes where the implicit move equation is solvable, one gets a trajectory with
\emph{no finite flip/crossing tick}:
\[
0<\Theta_k<R^{\rm abs}_k\quad\text{for all finite }k,
\qquad\text{and}\qquad
R^{\rm abs}_k\downarrow 0.
\]
This is a feedback consequence of the \emph{implicit} $\Lambda$-solve.
At tick $k$ the solver does \emph{not} choose $\Theta_k$ directly; it chooses $\Lambda_{k+1}$ as a real root of the
energy-balance residual $f_k(\Lambda)=0$ (cf.\ \eqref{eq:residual}).
The resulting $\Theta_k$ is then the BY-kernel output at that $\Lambda_{k+1}$, i.e.\ a function
$\Theta_k=\Theta(\Lambda_{k+1};\Delta I_k,d_k,L_{4,k},\text{locks})$ which is linear in $\Delta I_k$ but is
strongly suppressed at large $\Lambda$ and/or large separation $d_k$ by the evanescent factor $\exp(-d_k q)$ in the kernel
(\cref{subsec:theta}).
Thus, when $R^{\rm abs}_k$ becomes small, the solver can keep the trajectory feasible (no overshoot) by moving to a
slightly larger cutoff so that $\Theta_k$ becomes correspondingly smaller.
In the solvable no-leak branch this yields an infinite ``micro-step'' regime in which
$\Lambda_{k+1}-\Lambda_k\to 0$ and $\Theta_k\to 0$ while still maintaining $\Theta_k<R^{\rm abs}_k$
\emph{until the budget itself vanishes}.

Formally, the exact-summability regime is the one-wall continuation in which (a) the move equation remains solvable for all
ticks (no split) and (b) exported payload does not terminate before the feasibility boundary is reached (``enough
coherence''). In that regime the only terminal point is $R^{\rm abs}=0$, reached asymptotically with $\Theta_k\to 0$, and
taking $N\to\infty$ in \eqref{eq:telescoping} yields \eqref{eq:exact-summability}.

\textbf{(iii) Why constant/additive leaks spoil exact summability (and create splits).}
Additive leaks enter as a positive shift $Q_{\rm leak}$ in the residual \eqref{eq:residual}.
Once $R^{\rm abs}_k$ is small enough that the residual minimum becomes strictly positive, the real root disappears and the
local quadratic continuation produces a complex-conjugate pair (``two complex solutions''); this is the split mechanism
discussed in \cref{subsec:root-vs-split}.
The trajectory then terminates before $R^{\rm abs}_k$ reaches $0$, and \eqref{eq:exact-summability} fails exactly by the
nonzero remainder $R^{\rm abs}_\infty$ in \eqref{eq:telescoping}.
\end{proof}

\begin{remark}[Why this is still nontrivial]
The identity \eqref{eq:telescoping} is algebraic once the update law is fixed. The \emph{nontrivial} content found in the
simulator experiments is that the \emph{implicit one-wall $\Lambda$-solve} dynamically adjusts the cutoff trajectory so as
to \emph{avoid overshoot} when a real solution exists: the solver selects $\Lambda_{k+1}$ satisfying the discrete
energy-balance residual \eqref{eq:residual}, and in the no-leak / fixed-work-percentage (proportional) leak regimes this
leads to
$0<\Theta_k<R^{\rm abs}_k$ for all ticks and hence $R^{\rm abs}_k\downarrow 0$ while $S_N\uparrow \abs{R_0}$.
By contrast, additive/rigid leak models can destroy the existence of real solutions for the $\Lambda$-move equation,
producing a residual minimum strictly above tolerance (numerically: loss of bracketing/root and appearance of an
``imaginary root'' diagnostic), which is interpreted as the onset of a split (see \cref{subsec:root-vs-split} and
\cref{sec:phenom}).
\end{remark}

\subsection{A sufficient summability bound from QND exponential tails}
\label{subsec:summability-bound}

In the steady-actuator regime, UOM provides a geometric/exponential tail for $\Delta I_k$ (label \texttt{eq:tail-geom} in
\texttt{lem:exp-tail}):
\begin{equation}
\Delta I_k \simeq F_0\,(1-e^{-\beta})\,e^{-\beta k},
\qquad \sum_{k\ge 0}\Delta I_k = F_0.
\label{eq:deltaI-summable}
\end{equation}
If $K(\bar\Lambda_k)$ is bounded above by $K_{\max}$ over the relevant tick range, then
\begin{equation}
\sum_{k\ge 0}\Theta_k
\;=\;\sum_{k\ge 0}K(\bar\Lambda_k)\Delta I_k
\;\le\;K_{\max}\sum_{k\ge 0}\Delta I_k
\;=\;K_{\max}F_0.
\label{eq:theta-summable-bound}
\end{equation}
This gives a clean \emph{sufficient} summability condition and makes explicit the two control levers: (i) the available
coherence stock $F_0$ and (ii) the magnitude of $K$ along the receiver trajectory.

\section{Wall balance and the implicit $\Lambda$-solver}
\label{sec:lambda-solver}

\subsection{Continuous wall RG balance}
\label{subsec:wall-balance}

UOM's SK/HJ wall balance is written as (UOM labels \texttt{eq:wall-ode} and \texttt{eq:ads5-wall-ode})
\begin{equation}
\mu_{\rm RG}(\Lambda)\,\dot\Lambda(t)
= (\mathcal P_{\rm in}-\mathcal P_{\rm out})(t) - \dot{\mathcal F}_{\rm rec}(t),
\label{eq:wall-ode-note}
\end{equation}
and in integrated/budget form
\begin{equation}
\Delta\lambda
\le \int W^{\rm avail}(t)\,dt,
\qquad \lambda'(\Lambda)=\mu_{\rm RG}(\Lambda).
\label{eq:wall-budget-note}
\end{equation}
Reduced one-wall closures treat the tick as an implicit step that chooses $\Lambda_{k+1}$ such that a discrete analogue of
this balance holds.

\subsection{Geometry-side free energy surrogate and discrete balance}
\label{subsec:geom-energy}

The phenomenological one-wall closures use a geometry-side surrogate
\begin{equation}
E_{\rm geom}(\Lambda,L_4;\kappa)
\;=\;
\zeta\,\frac{A+\kappa\,a_{\rm cft}\,\Lambda^2}{L_4},
\qquad
A=\frac{M_5^3}{\ell_{\rm AdS5}}.
\label{eq:Egeom}
\end{equation}
The discrete step compares $E_{\rm geom}$ between tick endpoints, with $L_4$ updated from the curvature bookkeeping
$L_4=\sqrt{12/\abs{R}}$.

\subsection{Wall-work kernel}
\label{subsec:wall-kernel}

The one-wall closure uses a wall penalty (work) kernel $J_{\rm wall}(\Lambda_k\to\Lambda)$ obtained by numerically
evaluating an integral along the log-$\Lambda$ path. Heuristically, it implements the tick budget
\begin{equation}
J_{\rm wall}(\Lambda_k\to\Lambda_{k+1}) \approx E_{\rm geom}(\Lambda_k,L_{4,k})-E_{\rm geom}(\Lambda_{k+1},L_{4,k+1}),
\label{eq:discrete-balance-no-leak}
\end{equation}
and the open-system closure generalizes this to include a nonnegative leak term $Q_{\rm leak}$ (UOM text around
\texttt{eq:ads5-wall-ode}):
\begin{equation}
E_{\rm geom}(\Lambda_k,L_{4,k})-E_{\rm geom}(\Lambda_{k+1},L_{4,k+1})
\approx
J_{\rm wall}(\Lambda_k\to\Lambda_{k+1}) + Q_{\rm leak,k}.
\label{eq:discrete-balance-leak}
\end{equation}

\subsection{Implicit solve and ``root vs split''}
\label{subsec:root-vs-split}

Define the residual at a candidate $\Lambda\ge \Lambda_k$:
\begin{equation}
f(\Lambda)
:=
J_{\rm wall}(\Lambda_k\to\Lambda)+Q_{\rm leak}(\Lambda)
-\Big(E_{\rm geom}(\Lambda_k,L_{4,k})-E_{\rm geom}(\Lambda,L_4(\Lambda))\Big).
\label{eq:residual}
\end{equation}
\begin{itemize}[leftmargin=1.6em]
\item \textbf{Move step:} if a root exists, choose $\Lambda_{k+1}$ so that $f(\Lambda_{k+1})=0$.
\item \textbf{Fallback + split (v2):} if no root is bracketed, minimize $f(\Lambda)$ over a grid (optionally refine).
If the \emph{minimized residual is strictly positive} (exceeds tolerance), trigger a split event at $\Lambda_\star$; if the
minimized residual is nonpositive, accept an approximate move with slack.
\end{itemize}
This is the operational meaning of the split diagnostic used in phenomenological one-wall solvers.

\begin{remark}[Bracketing failure $\neq$ loss of a real move solution]
\label{rem:bracketing-failure}
The implementations typically seed the sign-change search at a slightly increased cutoff
$\Lambda_{\rm lo}=\Lambda_k(1+\varepsilon_\Lambda)$ (e.g.\ $\varepsilon_\Lambda=10^{-9}$) to avoid the trivial fixed point.
In late-time regimes where the true solution satisfies $\Lambda_{k+1}-\Lambda_k=o(\varepsilon_\Lambda\Lambda_k)$, the
root can lie inside the \emph{micro-step} interval $[\Lambda_k,\Lambda_{\rm lo}]$ and is therefore missed by the bracketing
loop (which only evaluates $f$ for $\Lambda\ge\Lambda_{\rm lo}$).
The subsequent grid-min fallback can then return $\Lambda_{\rm lo}$ with a very small residual, even though a real root
exists.

In phenomenological implementations this issue appears whenever
$f(\Lambda_k)<0$ but $f(\Lambda_k(1+\varepsilon_\Lambda))>0$: then a real solution exists inside the excluded micro-step
interval.
Therefore a split should be diagnosed by the \emph{positivity of the minimized residual} (equivalently, under local
convexity, $f(\Lambda_\star)>0$ at the minimizer), not by ``no bracket found'' alone.
\end{remark}

\begin{remark}[``Two complex solutions'' as the analytic continuation of the move equation]
In reduced numerical closures, loss of a real move solution is naturally viewed as the emergence of two
complex-conjugate solutions of an (approximately) quadratic move equation.
One may therefore use a local diagnostic based on the second derivative \(f''\) in the variable \(\ell=\log\Lambda\).

Indeed, if locally
\[
f(\ell)\approx r+\tfrac12 f''(\ell_\star)\,(\ell-\ell_\star)^2,\qquad r:=f(\ell_\star)>0,
\]
then the equation $f(\ell)=0$ has no real solutions but has complex solutions
\[
\ell=\ell_\star\pm i\sqrt{\frac{2r}{f''(\ell_\star)}}.
\]
Thus the split trigger ``$\min f>0$'' (no feasible one-wall move) is naturally equivalent to ``the real root pair has moved
off the real axis'', and the proxy estimates the imaginary displacement scale. This matches the stated interpretation of
the split as the point where the one-wall continuation ceases to be physically realizable by a real cutoff move.
\end{remark}

\subsection{Receiver-wall leakage $Q_{\rm leak}$ and the RG-complement (split ancilla) sector}
\label{subsec:leak-theory}

\paragraph{Overview.}
UOM has an information-theoretic ``leakage'' notion (bridge/leak instrument, Definition \texttt{def:eta-bridge-leak}), but
the leak term relevant for the one-wall \(\Lambda\)-solver is different: it is a \emph{receiver-wall energy complement}
appearing in the SK/HJ balance.
In the open-system version of that balance one may include a nonnegative term \(Q_{\rm leak}\ge 0\) so that
\begin{equation}
\Delta E_{{\rm geom},k}\ =\ J_{{\rm wall},k}\ +\ Q_{{\rm leak},k},
\qquad Q_{{\rm leak},k}\ge 0,
\label{eq:wall-leak-balance}
\end{equation}
equivalently \(W^{\rm avail}=W_{\rm RG}+Q_{\rm leak}\) in integrated form.

\paragraph{Interpretation used below.}
Throughout the rest of the note we interpret \(Q_{{\rm leak},k}\) as the tickwise energy flux into the hidden completion
sector \(B\), namely \(Q_{{\rm leak},k}=\Delta E_{B,k}\), with \(B\) prepared by the welded history through the quadratic
SK data.
The detailed derivation of this statement, the Gaussian recursion for the \(B\)-covariance, and the resulting scalar split
criterion are given in \cref{sec:resolution}.

\paragraph{Placement of RP/OS diagnostics.}
UOM's active KR receiver is OS-positive and has vanishing RP block on the welded family, so any nontrivial RP witness must
be assigned to the analytic-continued completion sector rather than to the healthy receiver itself.
This placement is the central reason the current receiver-tangent proxy cannot be the fundamental leak law.

\section{Split in UOM: conservation, diagnostics, and dilation picture}
\label{sec:uom-split}

\subsection{Charge conservation across a split}
\label{subsec:by-charge-conservation}

UOM app. \texttt{app:energy-split} states: enclose the split region in a small Gaussian pillbox and integrate the
Hamiltonian constraint. With a static UV boundary, the Brown--York charge change is $\Delta Q_{\rm BY}=0$ and the pillbox
balance reproduces the wall ODE projected onto $\Sigma_\Lambda$. In short: the split is a redistribution of quasi-local
energy compatible with global conservation.

\subsection{Split identification diagnostic}
\label{subsec:split-ident}

UOM treats $\sigma_{\rm eff}\to\sigma_\star$ (equivalently $\Lambda_{4,{\rm R}}\to 0^+$) as a \emph{diagnostic}
of approach in the one-wall channel; in ``Model A'' the split onset is identified instead by the first crossing of a
two-wall vs one-wall free-energy difference $\Delta\mathcal F(\Lambda)$ (UOM assumption \texttt{ass:split-ident} and app.
\texttt{app:split-check}).

Reduced one-wall closures implement an \emph{analogous} idea: if the one-wall balance residual $f(\Lambda)$ cannot be made
small within tolerance, they interpret the failure as a split and use $\Lambda_\star$ minimizing the residual as the split
point.
This is a computational proxy for ``two-channel (two-wall) becomes favorable''.

\subsection{No-cloning resolution via Stinespring dilation}
\label{subsec:no-cloning}

UOM makes the split mechanism compatible with no-cloning by interpreting it as exposing part of a Stinespring dilation
ancilla as dynamical brane sector (UOM remark \texttt{rem:uom-main-007}): the two outputs (new dS source and new AdS
receiver) are correlated and do not carry independent copies of the UV state.
Operationally, UOM uses an algebraic notion of sub-cutoff equivalence (UOM definition \texttt{def:uom-main-002}) which does
not require a microscopic factorization.

\section{Split in 2T-LS: parity cocycle and non-splitting obstruction}
\label{sec:2t-split}

We recall only the pieces needed for the UOM split interpretation.

\subsection{Square-loop and parity cocycle $\omega$}
\label{subsec:omega}

In 2T-LS, one defines a typed syntax $\mathrm{Proc}_5$ of record-level morphisms with generators
(\texttt{time}, \texttt{rg}) and a ``square'' constructor $\mathrm{sq}(f;t,u,r)$ with $f\in\{\mathrm{diag},\mathrm{bridge}\}$
(2T-LS definition \texttt{def:Proc5}). One then defines:
\begin{itemize}[leftmargin=1.6em]
\item the canonical ``to-corner'' paths $\mathrm{toCorner}_f(t,u,r)$ and
\item the \emph{square-loop} $\mathrm{sqLoop}(t,u,r)$ comparing diag vs bridge procedures
(2T-LS definition \texttt{def:sqloop}).
\end{itemize}
The parity cocycle $\omega$ assigns $0$ to diag squares and $1$ to bridge squares, extends additively mod $2$ under
composition, and is invariant under inverses and casts (2T-LS definition
\texttt{def:limited-scope-theorem-translation-clean-003}).

The key lemma is that $\omega$ is well-defined and detects the square-loop:
\begin{equation}
\omega(\mathrm{sqLoop}(t,u,r))=1\qquad (t,u,r),
\label{eq:omega-sqloop}
\end{equation}
which is 2T-LS lemma \texttt{lem:omega}.

\subsection{$\Ztwo$-cover and non-splitting}
\label{subsec:nonsplitting}

Define the $\Ztwo$-cover $\pi:\widetilde{\mathcal G}\to\mathcal G$ whose objects are pairs $(r,b)$ with $b\in\Ztwo$
and whose morphisms $(r,b)\to(s,c)$ are those base morphisms $p:r\to s$ with $c=b+\omega(p)$ mod $2$
(2T-LS definition \texttt{def:limited-scope-theorem-translation-clean-004}).

\paragraph{Splitting $\Leftrightarrow$ coboundary (2T-LS).}
2T-LS makes the relationship between splittings and coboundaries explicit (lemma \texttt{lem:split-cob}):
there exists a splitting if and only if there exists a function $\sigma:\Sys\to\Ztwo$ such that
$\omega(p)=\sigma(r)+\sigma(s)\pmod 2$ for every morphism $p:r\to s$ in $\mathcal G$.
Equivalently, an \emph{odd loop} ($\omega(p)=1$ for some loop $p:r\to r$) obstructs splitting (corollary
\texttt{cor:odd-loop}).

Then:
\begin{theorem}[Non-splitting of the $\Ztwo$-cover (2T-LS)]
Under assumptions \textbf{(A1)}--\textbf{(A6)} of the 2T-LS paper, if $\Sys$ is nonempty then the cover does not admit a
splitting; equivalently, $\omega$ is not a coboundary on $\mathcal G$ (2T-LS theorem \texttt{thm:non-splitting}).
\end{theorem}
The obstruction can already be witnessed by the two parallel ``to-corner'' morphisms $p_{\rm diag}$ and $p_{\rm bridge}$
whose $\omega$-values differ (2T-LS proposition \texttt{prop:two-way-obstruction}).

\paragraph{Interpretation for UOM.}
At the level of \emph{record-accessible} descriptions, the split is precisely the place where one must choose between two
continuations that are not homotopic within a purely diagonal (classical) sector; the $\Ztwo$ holonomy is the formal
encoding of that ``odd'' choice. This matches the proposed identification of the split with the odd cocycle generator.

\section{Coarse-graining classicality reframing (analysis of the proposal)}
\label{sec:coarse-graining}

This section formalizes (and slightly sharpens) the proposed split reframing given in the prompt.

\subsection{One global accessibility/coarse-graining map}
\label{subsec:accessibility-map}

Fix a tick/aeon scale index $k$ and define a CPTP \emph{accessibility} (or RG coarse-graining) map
\begin{equation}
\cC_k:\cD(\cH_{\rm tot})\to \cD(\cH_{{\rm acc},k}),
\qquad
\rho_k^{\rm RG}:=\cC_k(\varrho_k^{\rm tot}).
\label{eq:accessibility}
\end{equation}
The point is that ``accessibility'' is defined \emph{after} applying $\cC_k$; no claim is made that the microscopic
5D state ever becomes diagonal.

\paragraph{UOM-compatible construction.}
UOM already provides the natural building blocks for $\cC_k$:
\begin{enumerate}[leftmargin=2em]
\item patchwise pointer-aligned QND dephasing,
\item conditional expectation $\mathbb E_{\rm ptr}$ onto the pointer algebra,
\item a commutant record map $\cR_k$ with commutative range,
\item an admissible bridge $\cB_k$ into the receiver code subspace, followed by any ``export'' selection of RG-visible
observables (e.g.\ DC anomaly coordinates).
\end{enumerate}
Composing these yields a single global CPTP map of the form \eqref{eq:accessibility}.

\subsection{Classicality defect at RG scale}
\label{subsec:classicality-defect}

Fix on $\cH_{{\rm acc},k}$ a preferred ``ledger'' basis (or commutative subalgebra) that defines what ``diagonal'' means.
Define the \emph{classicality defect} by trace distance to the diagonal set:
\begin{equation}
\mathrm{cl}_k
:=\inf_{\sigma\in\mathrm{Diag}(\cH_{{\rm acc},k})}\norm{\rho_k^{\rm RG}-\sigma}_1.
\label{eq:cl-def}
\end{equation}
This is a legitimate operational requirement: it quantifies how well an observer restricted to the RG-accessible degrees
of freedom can approximate the state by a classical probability distribution.

\paragraph{A computable surrogate.}
While \eqref{eq:cl-def} is the clean definition, in practice one often uses
\begin{equation}
\mathrm{Off}_k:=\norm{\rho_k^{\rm RG}-\diag(\rho_k^{\rm RG})}_1,
\label{eq:off-def}
\end{equation}
which satisfies $\mathrm{cl}_k\le \mathrm{Off}_k$ by feasibility of $\diag(\rho_k^{\rm RG})$.
(Note: for the trace norm, $\diag(\rho)$ need not be the unique optimizer in all dimensions; we will only use
$\mathrm{Off}_k$ as a controlled upper bound and as a diagnostic aligned with dephasing.)

\subsection{End-of-aeon classicality: where the bound is tightest}
\label{subsec:end-of-aeon}

Under pointer-basis dephasing, off-diagonals are suppressed multiplicatively in the pointer eigenbasis. Combined with
data-processing (CPTP contractivity of distinguishability), one expects the smallest classicality defect to occur after:
\begin{enumerate}[leftmargin=2em]
\item enough QND time has accumulated so that coherence is exponentially small,
\item patch reconciliation has produced a consistent global record basis, and
\item the accessibility map $\cC_k$ has ``forgotten'' phases (projects/averages onto commutative output).
\end{enumerate}
This motivates defining \emph{diagonal ledger points} at the \emph{end of an aeon}:
\begin{equation}
\inf_{\sigma\in\mathrm{Diag}}\norm{\rho_{n,{\rm end}}^{\rm RG}-\sigma}_1\le \varepsilon_n,
\qquad \varepsilon_n\downarrow 0.
\label{eq:end-classicality}
\end{equation}
It is consistent with UOM that the diagonal approximation is tightest late in the aeon and not necessarily after each tick.

\begin{lemma}[Off-diagonal suppression under pointer dephasing (diagnostic bound)]
Let $\Phi$ be a pointer-basis QND dephasing channel of the form \eqref{eq:qnd-channel-closed} on a finite-dimensional
space $\cH_{{\rm acc},k}$ with pointer eigenbasis $\{\lvert i\rangle\}$.
Let $\Pi_{\rm diag}$ denote the pinching (conditional expectation) onto the diagonal in this basis:
$\Pi_{\rm diag}(\rho)=\diag(\rho)$.
Then for $A:=\rho-\Pi_{\rm diag}(\rho)$ (zero diagonal) the matrix elements obey
\[
(\Phi(A))_{ij}=e^{-\Gamma_{ij}}A_{ij}\qquad (i\neq j)
\]
with $\Gamma_{ij}\ge 0$ determined by the $\Gamma_a^{(k)}$ and pointer gaps.
Consequently, if $\Gamma_{\min}:=\min_{i\neq j}\Gamma_{ij}$ then
\begin{equation}
\norm{\Phi(A)}_F \le e^{-\Gamma_{\min}}\norm{A}_F
\qquad\Longrightarrow\qquad
\mathrm{Off}(\Phi(\rho)) \le \sqrt{d}\,e^{-\Gamma_{\min}}\,\mathrm{Off}(\rho),
\label{eq:off-bound}
\end{equation}
where $d=\dim\cH_{{\rm acc},k}$, $\norm{\cdot}_F$ is the Frobenius norm, and $\mathrm{Off}(\rho)=\norm{\rho-\diag(\rho)}_1$.
\end{lemma}

\begin{remark}[Backflow, leakage, and the ``tail + leak'' corrections]
If the effective reduced dynamics is not exactly Markovian on the accessible space (e.g.\ due to memory kernels) or if the
tick map includes a leakage branch $\cL_k$ outside the commutative record/code sector, then
\eqref{eq:off-bound} is modified schematically to
\[
\mathrm{Off}_{k+1}
\;\lesssim\;
e^{-\gamma_k}\,\mathrm{Off}_k
\;+\;(\text{backflow tail})
\;+\;(\text{leakage}),
\]
matching the qualitative structure suggested in the prompt. UOM explicitly controls backflow tails in its source-side
analysis (e.g.\ App.\ \texttt{app:dS-collapse}, ``Exponentially decaying backflow'') and treats leakage via the instrument
decomposition in \texttt{def:eta-bridge-leak}.
\end{remark}

\subsection{Split as an odd lift / Stinespring refactorization}
\label{subsec:split-as-lift}

Let $\rho_{n,{\rm end}}^{\rm RG}$ denote the end-of-aeon accessible state.
Model the split by a (generally non-unique) \emph{section} (odd lift)
\begin{equation}
\mathrm{sec}_\Pi:\rho_{n,{\rm end}}^{\rm RG}\mapsto x_{n+1,{\rm start}}\in \mathcal S_{\rm SK},
\label{eq:secpi}
\end{equation}
where $\mathcal S_{\rm SK}$ is the space of Schwinger--Keldysh (two-time) states for the next aeon, including a new
receiver initialized in a pure ``vacuum-like'' code state together with a new source state.
Operationally, $\mathrm{sec}_\Pi$ is best understood as a Stinespring refactorization/isometry on an enlarged space
rather than as ``recohering the same subsystem''.

\paragraph{Connection to 2T-LS.}
The non-uniqueness / parity of the lift is exactly what the 2T-LS cocycle $\omega$ formalizes:
the choice between diag vs bridge corner procedures has intrinsic one-bit holonomy and cannot be absorbed into a
single global state label on records. Thus the split is the natural locus of the odd cocycle.

\subsection{Split ancilla $B$ and coherent birth: what drives it in UOM}
\label{subsec:split-ancilla}

This subsection develops the split proposal from the prompt into a UOM-faithful identification of
\emph{(i)} what the ``newborn'' sector $B$ is and \emph{(ii)} why it is generically coherent even if the inherited ledger
sector $A$ is (approximately) classical/QND at the end of the aeon.

\subsubsection*{Subsystem identification: inherited ledger $A$ vs completion sector $B$}
\label{subsec:AB-identification}

UOM's microscopic $5$D evolution is unitary; CPTP structure only appears after discarding degrees of freedom.
Therefore the split is \emph{not} a CPTP map acting on the old receiver alone.
The clean reading of the Stinespring/refactorization picture is:
\begin{itemize}[leftmargin=1.6em]
\item $A$ is the \emph{inherited receiver ledger sector}: the commuting family of record variables that constitute the
end-of-aeon accessible ``classical'' description (e.g.\ a DC anomaly/trace coordinate together with any other commuting
registers retained by $\cC_{n,{\rm end}}$).
\item $B$ is the \emph{completion / environment sector} that makes the receiver export isometric before reduction and that,
at the split, is reclassified as (part of) the \emph{new dS brane Hilbert space} (new source).
\end{itemize}
UOM offers two canonical realizations of such a completion sector:
\begin{enumerate}[leftmargin=2em]
\item \textbf{Receiver Stinespring completion.}
UOM explicitly enlarges the receiver code by a Weyl compensator $\sigma$ so that the welded export is an isometry, and
the trace anomaly appears only after tracing out $\sigma$ (UOM \texttt{par:uom-main-040}--\texttt{par:uom-main-044} and
\texttt{par:uom-main-077}).
In this reading, $B$ may include $\sigma$ (and allied completion modes such as the brane-local boundary fluctuation
$\mathcal S$ used in the portal sector).
\item \textbf{Commutant/purifying partner.}
On the source side, UOM's intrinsic bath is the commutant/purifying sector: patch marginals can look KMS/thermal because
the global state is purified across a complementary algebra (UOM \texttt{subsec:source-qnd-commutant} and
\texttt{subsec:commutant-stinespring}).
In this reading, $B$ is the canonical partner sector that becomes dynamical as the next aeon's source.
\end{enumerate}
Both choices implement the same structural point: $B$ is \emph{not a new knob} but ``whatever sector was previously
integrated out'' in the reduced description and is fixed by the same $5$D SK saddle.

\subsubsection*{Classical-control split unitary (ledger inheritance) and the coherence target}
\label{subsec:classical-control-split}

To keep the inherited ledger classical while creating a coherent newborn source, the split must be QND on the ledger:
\begin{equation}
U_{\rm split}
\;=\;
\sum_i \Pi_i^{(A)}\otimes U_i^{(B)},
\qquad
[U_{\rm split},\Pi_i^{(A)}\otimes \mathbf 1_B]=0\quad\forall i,
\label{eq:split-controlled}
\end{equation}
where $\{\Pi_i^{(A)}\}$ are orthogonal projectors onto the end-of-aeon ledger basis (or, equivalently, minimal
projections of the commutative record algebra).
If $\rho_A$ is diagonal in this basis (or $\mathrm{cl}_{n,{\rm end}}\ll 1$), then $A$ remains (approximately) diagonal
after the split and can be viewed as \emph{classical control}.

Recoherence, in the operational sense relevant for UOM's ``recohere $\to$ QND dephase'' picture, is a condition on $B$:
let $\Pi_\Delta$ denote the pointer/record projection used in the next aeon (band/pointer basis on the newborn source).
Then the newborn source must satisfy
\begin{equation}
\norm{\rho^{(B)}_{\rm start}-\Pi_\Delta(\rho^{(B)}_{\rm start})}_1 \;>\; 0,
\label{eq:B-coherence-target}
\end{equation}
otherwise the next aeon's QND absorber produces no new record increment (and no positive ``entropy production'' in the
Spohn/QND sense).

\subsubsection*{What UOM fixes about the welded source $\int J\cdot\mathcal O$}
\label{subsec:welded-source-fixed}

UOM's microscopic receiver action contains an explicit welded source term (UOM \texttt{par:uom-main-077}):
\begin{equation}
S_{\rm weld}[J]
\;=\;
\int_{\Sigma_{\rm rec}}\! d^4x\,\sqrt{-\gamma}\,\big(J\cdot \mathcal O_{\rm LoG}\big),
\qquad
J\equiv J_k\ \text{(tick-localized LoG envelope).}
\label{eq:Sweld-note}
\end{equation}
Two points are essential for the split proposal:
\begin{enumerate}[leftmargin=2em]
\item \textbf{$J$ is not arbitrary.}
In UOM it is a \emph{compiled/welded} image of the source-side payload into the receiver trace channel, regulated by the
receiver cutoff (UOM \texttt{par:uom-main-080}, and ``welding/regulation'' step in the proof of
\texttt{lem:ImW-noise}).
In the portal sector, $J$ is explicitly identified as a LoG projection of the Brown--York trace response
(UOM \texttt{eq:JLoG-def}): $J_{\rm LoG}=\kappa_{\rm HS}(\mathrm{LoG}\star T^\mu{}_\mu{}^{\rm (BY)})$.
\item \textbf{$J$ is a classical control parameter at the coarse-grained level.}
Once the ledger $A$ is diagonal, $J$ is effectively a function of the ledger history (the accepted payload/trace record),
so the controlled form \eqref{eq:split-controlled} is the natural way to package ``history-dependent preparation'' of the
newborn sector $B$.
\end{enumerate}

\subsubsection*{Quadratic generators and the squeeze picture: what UOM provides (and what it implies)}
\label{subsec:quadratic-generators}

UOM evaluates the \emph{Schwinger--Keldysh influence functional} on the welded source and integrates out bulk+brane fields
in the Gaussian sector at \emph{quadratic order in $J$} (UOM \texttt{par:uom-main-080}):
\begin{equation}
W[J]
=\frac12\!\int\!\frac{d^4k}{(2\pi)^4}\,
J(-k)\,\big(G_R(k)+i\,\mathcal N(k)\big)\,J(k)
+O(J^3),
\label{eq:W-quadratic-note}
\end{equation}
with $G_R$ retarded and $\mathcal N\ge 0$ a KMS/noise kernel (hence $\Im W\ge 0$; UOM \texttt{lem:ImW-noise}).
Equivalently, in Keldysh $(\mathrm{cl},\mathrm{q})$ variables,
\[
W[\tau]\;=\;-\langle \tau_{\rm q},G^{\rm R}\tau_{\rm cl}\rangle\;+\;\frac{i}{2}\langle\tau_{\rm q},N_\Lambda\tau_{\rm q}\rangle\;+\;O(\tau^3),
\]
as in the proof of UOM \texttt{lem:ImW-noise}.

\paragraph{Implication for split-state preparation.}
To quadratic order, the completion sector excited by $J$ is \emph{Gaussian}. For bosonic degrees of freedom, any Gaussian
state can be prepared from a vacuum-like reference by a product of a displacement (linear drive) and a symplectic
transformation (squeezing/Bogoliubov):
\begin{equation}
U_i^{(B)}\;\sim\;D(\alpha_i)\,S(\zeta_i),
\qquad
S(\zeta)=\exp\!\Big(\tfrac12\sum_{\alpha,\beta}\zeta_{\alpha\beta}\,b_\alpha^\dagger b_\beta^\dagger-\mathrm{h.c.}\Big),
\label{eq:gaussian-prep}
\end{equation}
with multi-mode creation/annihilation operators $b_\alpha$ on $B$.
In this representation:
\begin{itemize}[leftmargin=1.6em]
\item the \textbf{linear} welded source term \eqref{eq:Sweld-note} supplies the displacement/drive (coherent component),
and
\item the \textbf{quadratic} SK kernels in \eqref{eq:W-quadratic-note} fix the Gaussian covariance, hence the effective
squeeze/Bogoliubov parameters.
\end{itemize}
UOM does not presently write the squeeze generator as an operator $H_{B,{\rm sq}}$; rather, it provides the quadratic
influence data $(G_R,\mathcal N)$ and the regulated on-shell quadratic functionals (Fisher and noise forms) from which
one can reconstruct an equivalent Gaussian unitary description if one chooses an explicit mode basis on $B$.

\paragraph{Contact with $H^{(i)}_{B,{\rm drive}}$ and $H^{(i)}_{B,{\rm sq}}$.}
If one insists on a Hamiltonian/stroboscopic representation of the conditional preparation unitaries $U_i^{(B)}$, a
minimal parametrization consistent with \eqref{eq:gaussian-prep} is
\begin{equation}
U_i^{(B)}
\;\simeq\;
\mathcal T\exp\!\Big(-i\int\!dt\;\big(H^{(0)}_B(t)+H^{(i)}_{B,{\rm drive}}(t)+H^{(i)}_{B,{\rm sq}}(t)\big)\Big),
\label{eq:HB-drive-sq}
\end{equation}
where the drive is linear in fields/operators and the squeeze part is quadratic in canonical modes:
\[
H^{(i)}_{B,{\rm drive}}(t)=\int\!d^3x\;\big(J_i(t,x)\,\mathcal O_B(t,x)+\overline{J_i(t,x)}\,\mathcal O_B^\dagger(t,x)\big),
\]
\[
H^{(i)}_{B,{\rm sq}}(t)
=\frac{i}{2}\sum_{\alpha,\beta}\Big(\kappa^{(i)}_{\alpha\beta}(t)\,b_\alpha^\dagger b_\beta^\dagger
-\overline{\kappa^{(i)}_{\alpha\beta}(t)}\,b_\alpha b_\beta\Big)
\;+\;\sum_{\alpha,\beta}h^{(i)}_{\alpha\beta}(t)\,b_\alpha^\dagger b_\beta.
\]
In UOM terms, $J_i$ is supplied by the welded/compiled source history, and the quadratic kernels determining
$\kappa^{(i)}$ and $h^{(i)}$ are encoded by the same Gaussian SK reduction data that define $W[J]$.

\paragraph{Why analytic continuation naturally produces squeezing.}
UOM's AdS$\to$dS analytic continuation maps a Euclidean on-shell quadratic functional to an oscillatory Hartle--Hawking
wavefunctional (UOM \texttt{subsec:analytic-continuation}):
\[
\Psi_{\rm HH}[\phi_{(0)}]\propto e^{-S_{\rm os}^{\rm (EAdS)}[\phi_{(0)}]}
\ \longrightarrow\ 
e^{\,i S_{\rm eff}^{\rm (dS)}[\phi_{(0)}]}\,.
\]
For free/linearized modes this continuation is precisely the mechanism by which an ``in'' vacuum becomes a squeezed
state with respect to an ``out'' (static-patch/pointer) basis: the quadratic phase encodes the Bogoliubov mixing.
Thus UOM's existing quadratic (Gaussian) infrastructure already supplies the natural home for the
``quadratic generators'' in a split-as-recoherence picture.

\section{Phenomenological inputs from existing simulations}
\label{sec:phenom}

The simulator family is useful here only as a reduced phenomenology of the wall balance. Three lessons are robust and
theory-relevant.

\subsection{The exhaustion branch is physically meaningful}

The one-wall solver can remain real for all finite ticks and realize exact summability:
\[
0<\Theta_k<R_k^{\rm abs}\quad\text{for all finite }k,\qquad
R_k^{\rm abs}\downarrow 0,\qquad
\sum_{k\ge 0}\Theta_k=\abs{R_0}.
\]
This is the phenomenological evidence for the \emph{exhaustion branch}: split is not guaranteed, and proportional
or no-leak closures can stay on the solvable one-wall channel indefinitely.

\subsection{Residual positivity is the meaningful split diagnostic}

The theory-level split criterion is the positivity of the minimized discrete residual,
\[
\min_{\Lambda\ge \Lambda_k} f_k(\Lambda)>0,
\]
not mere failure of numerical bracketing and not the magnitude \(\min |f|\).
This is exactly the criterion compatible with the wall-balance interpretation of split as loss of a real one-wall
continuation.

\subsection{The current RP proxy is only a phenomenological probe}

The presently implemented quantity \(Q_{\rm rg}=\chi J_{\rm RP}/\beta_{\rm dS}\) is a useful reduced probe of when a
one-wall scalar closure becomes hard to realize, but it is not a derived completion-sector leak.
The key reason is structural: it is built from a receiver-side welded tangent, whereas the healthy KR receiver is
OS-positive and has vanishing RP block on the welded family.
The simulations therefore motivate two qualitative conclusions, but do not yet provide the microscopic leak law:
\begin{enumerate}[leftmargin=2em]
\item a leak capable of forcing split must be history-conditioned and non-proportional to the shrinking payload;
\item any RP-based witness must be placed on the analytic-continued completion sector rather than on the healthy
receiver tangent.
\end{enumerate}

\IfFileExists{completion_sector_split_model_body.tex}{%
  \section{Completion-sector leak model and split criterion}
\label{sec:resolution}

This section separates what is directly derived from UOM from the additional hypotheses required to turn those ingredients
into a concrete split-driving leak model. The new organizing principle is:
the primary hidden variable is a \emph{signed odd split amplitude} \(a_k\), while the positive stock
\(S_k:=a_k^2\) is the quantity that can enter scalar observables and leak laws.
We reserve \(R_k^{\rm abs}\) for the remaining curvature budget and
\[
\mathfrak r_k:=\min_{\Lambda\ge \Lambda_k} f_k(\Lambda)
\]
for the minimized one-wall residual.

\subsection{Odd split mode, output parity, and the rank-one normal form}
\label{subsec:leak-solution}

The only ``leak'' that directly competes with wall work in the \(\Lambda\)-solver is the open-system modification of the
wall energy balance allowed explicitly in UOM's SK/HJ wall equation:
\begin{equation}
\Delta E_{{\rm geom},k}\ =\ J_{{\rm wall},k}\ +\ Q_{{\rm leak},k},
\qquad Q_{{\rm leak},k}\ge 0,
\label{eq:resolution-wall-leak}
\end{equation}
This is the welded wall-balance form used in the earlier summability analysis.
This is a \emph{receiver-side energy complement} term. It must not be conflated with UOM's information-theoretic
bridge/leak instrument from the UOM main manuscript.

\paragraph{What is already derived in UOM.}
Four ingredients are fixed by the manuscript:
\begin{enumerate}[leftmargin=2em]
\item \textbf{Wall balance.} The reduced one-wall description may contain a nonnegative complement \(Q_{\rm leak}\) as in
\eqref{eq:resolution-wall-leak}.
\item \textbf{Quadratic welded SK data.} In the Gaussian sector,
\[
W_\Lambda[J]=\frac12\int J\,(G_{R,\Lambda}+i\,\mathcal N_\Lambda)\,J+O(J^3),
\qquad \mathcal N_\Lambda\ge 0.
\]
\item \textbf{Healthy receiver.} The active KR receiver has vanishing welded RP block, so a receiver-tangent RP penalty is
not the fundamental home of split-driving oddness.
\item \textbf{Continued RP defect.} The analytic-continued covariance has a negative OS mode with
\[
\lambda_{\min}(\mu_{ps},\Lambda)
=
-\alpha(\Lambda)\mu_{ps}^2+O(\mu_{ps}^4),
\qquad
\alpha(\Lambda)\sim \kappa_d\,\Lambda^{d+2}H^{-d}.
\]
\end{enumerate}

\begin{assumption}[Split-mode observability]
\label{ass:wall-match}
Near split onset there exists a one-dimensional odd line \(L_k\) with local coordinate \(a_k\in\mathbb R\) such that the
split involution acts by
\[
a_k\longmapsto -a_k.
\]
The true two-wall advantage
\[
\Delta F_k:=\Delta\mathcal F_{2{\rm w}-1{\rm w},k}
\]
and the minimized one-wall residual
\[
\mathfrak r_k:=\min_{\Lambda\ge \Lambda_k} f_k(\Lambda)
\]
are smooth even scalar outputs of this same local mode.
Equivalently, the 2T \(\Gamma\)-odd descendant direction and the UOM continued negative mode \(u_-(\Lambda_k)\) are
assumed to represent the same active split mode up to a smooth local identification.
\end{assumption}

\begin{lemma}[Evenness of the minimized residual]
\label{lem:residual-even}
Suppose the pre-minimized one-wall residual is a smooth function \(f_k(\Lambda,a)\) satisfying
\[
f_k(\Lambda,-a)=f_k(\Lambda,a),
\]
and suppose there is a unique local minimizer \(\Lambda_\ast(a)\) near \(a=0\) such that
\[
\partial_\Lambda f_k(\Lambda_\ast(a),a)=0,
\qquad
\partial_\Lambda^2 f_k(\Lambda_\ast(0),0)>0.
\]
Then
\[
\mathfrak r_k(a):=f_k(\Lambda_\ast(a),a)
\]
is a smooth even function of \(a\).
\end{lemma}

\begin{proof}
By the implicit function theorem, \(\Lambda_\ast(a)\) is smooth near \(a=0\).
Since \(f_k(\Lambda,a)\) is even in \(a\), the stationarity equation for \(a\) and \(-a\) is identical; uniqueness of the
local minimizer gives
\[
\Lambda_\ast(-a)=\Lambda_\ast(a).
\]
Therefore
\[
\mathfrak r_k(-a)=f_k(\Lambda_\ast(-a),-a)=f_k(\Lambda_\ast(a),a)=\mathfrak r_k(a),
\]
so \(\mathfrak r_k\) is smooth and even.
\end{proof}

\begin{proposition}[Rank-one normal form near split onset]
\label{prop:rank-one-wall-match}
Under \cref{ass:wall-match}, the two scalar outputs admit local even expansions
\begin{equation}
\Delta F_k(a)
=
\mu_k^{(2{\rm w})}+h_k a^2+O(a^4),
\label{eq:DeltaF-detuned}
\end{equation}
\begin{equation}
\mathfrak r_k(a)
=
\mu_k^{(1{\rm w})}+m_k a^2+O(a^4),
\label{eq:rres-detuned}
\end{equation}
with \(h_k,m_k>0\) on the active split line.
\end{proposition}

\begin{proof}
Smooth even functions of one real variable have Taylor series containing only even powers.
The constants \(\mu_k^{(2{\rm w})}\) and \(\mu_k^{(1{\rm w})}\) are the baseline detunings of the two-wall and one-wall
channels. Positivity of \(h_k\) and \(m_k\) expresses that the selected mode is active in both outputs.
\end{proof}

\begin{corollary}[Onset wall-match in rank one]
\label{cor:rank-one-wall-match}
On the onset manifold,
\[
\mu_k^{(1{\rm w})}=\mu_k^{(2{\rm w})}=0,
\]
so
\[
\Delta F_k=h_k a_k^2+O(a_k^4),
\qquad
\mathfrak r_k=m_k a_k^2+O(a_k^4).
\]
Eliminating \(a_k^2\) gives
\begin{equation}
\Delta F_k
=
\frac{h_k}{m_k}\,\mathfrak r_k+O(\mathfrak r_k^2).
\label{eq:rank-one-wall-match}
\end{equation}
\end{corollary}

\subsection{Low-rank generalization}
\label{subsec:alpha-derived}

If the active split sector is \(m\)-dimensional, with odd amplitude vector \(a_k\in\mathbb R^m\), parity gives quadratic
normal forms
\begin{equation}
\Delta F_k(a)
=
\mu_k^{(2{\rm w})}+a^\top H_k a+O(\|a\|^4),
\label{eq:DeltaF-lowrank}
\end{equation}
\begin{equation}
\mathfrak r_k(a)
=
\mu_k^{(1{\rm w})}+a^\top M_k a+O(\|a\|^4),
\label{eq:rres-lowrank}
\end{equation}
with \(H_k,M_k\succeq 0\) on the active split subspace.

\begin{assumption}[Low-rank comparability]
\label{ass:low-rank-comp}
On the active split subspace there exist \(c_k^-,c_k^+>0\) such that
\[
c_k^- M_k \preceq H_k \preceq c_k^+ M_k.
\]
\end{assumption}

\begin{proposition}[Low-rank wall-match comparability]
\label{prop:low-rank-wall-match}
Assume \cref{ass:wall-match,ass:low-rank-comp} and restrict to the onset manifold
\(\mu_k^{(1{\rm w})}=\mu_k^{(2{\rm w})}=0\).
Then for sufficiently small \(\mathfrak r_k\) there exists \(C_k>0\) such that
\begin{equation}
c_k^-\,\mathfrak r_k-C_k\mathfrak r_k^2
\le
\Delta F_k
\le
c_k^+\,\mathfrak r_k+C_k\mathfrak r_k^2.
\label{eq:low-rank-wall-match}
\end{equation}
\end{proposition}

\begin{proof}
On onset,
\[
\Delta F_k=a_k^\top H_k a_k+O(\|a_k\|^4),
\qquad
\mathfrak r_k=a_k^\top M_k a_k+O(\|a_k\|^4).
\]
The quadratic-form comparison gives
\[
c_k^- a_k^\top M_k a_k\le a_k^\top H_k a_k\le c_k^+ a_k^\top M_k a_k.
\]
Since \(\mathfrak r_k\) is itself quadratic in \(a_k\) to leading order, the quartic remainders are \(O(\mathfrak r_k^2)\),
which yields \eqref{eq:low-rank-wall-match}.
\end{proof}

\subsection{Gaussian realization of the odd mode and the leak stock law}
\label{subsec:B-gaussian}

The theorem path above does \emph{not} rely on a KMS-subtracted matrix.
In particular, we do not use
\[
\Sigma_{B,k}:=\bigl(C_{B,k}-C^{\rm KMS}_{\Lambda_k}\bigr)+m_{B,k}m_{B,k}^\top
\]
as a positive object, because it need not be positive semidefinite.
When a safe matrix observable is needed in the realization layer, the natural PSD quantity is
\begin{equation}
M_{B,k}:=C_{B,k}+m_{B,k}m_{B,k}^\top\succeq 0.
\label{eq:MB-def}
\end{equation}

\begin{assumption}[Gaussian completion realization]
\label{ass:gaussian-completion}
For each admissible cutoff \(\Lambda\) in the welded regime, the quadratic influence functional \(W_\Lambda[J]\) admits a
Gaussian dilation, unique up to symplectic equivalence, by a bosonic completion sector \(B_\Lambda\) with canonical
quadratures \(X_\Lambda\) and positive quadratic Hamiltonian
\[
H_{B,\Lambda}=\frac12 X_\Lambda^\top h_\Lambda X_\Lambda,
\qquad h_\Lambda\succeq 0.
\]
\end{assumption}

\begin{assumption}[Dominant odd-mode projection]
\label{ass:export-selection}
Inside the Gaussian completion of \cref{ass:gaussian-completion}, there exists a smooth unit vector
\[
u_k:=u_-(\Lambda_k)
\]
spanning the active odd line and separated from the orthogonal sector by a local spectral gap.
We write
\[
P_k^\perp:=I-u_k u_k^\top
\]
for the orthogonal projector.
\end{assumption}

\begin{assumption}[Classical-control gate]
\label{ass:classical-gate}
There exists a scalar gate \(g_{\rm cl}(k)\in[0,1]\), determined by the accessible classicality defect
\(\mathrm{cl}_k\) or \(\mathrm{Off}_k\), such that \(g_{\rm cl}(k)\ll 1\) at a fresh-aeon start and
\(g_{\rm cl}(k)\to 1\) late in the aeon.
\end{assumption}

\begin{proposition}[Projected Gaussian affine dynamics]
\label{prop:projected-affine}
Under \cref{ass:gaussian-completion,ass:export-selection}, one may choose a finite effective completion-mode basis in which
the completion quadratures obey a linear-Gaussian recursion
\begin{equation}
X_{k+1}=A_k X_k+B_k j_k^{\rm ctl}+\nu_k,
\label{eq:X-recursion}
\end{equation}
with
\[
\mathbb E[\nu_k\mid\mathcal F_k]=0,
\qquad
\mathbb E[\nu_k\nu_k^\top\mid\mathcal F_k]=Q_k\succeq 0.
\]
Consequently the PSD second moment
\[
M_{B,k}=C_{B,k}+m_{B,k}m_{B,k}^\top
\]
obeys
\begin{equation}
M_{B,k+1}=A_k M_{B,k}A_k^\top+Q_k+\text{drift terms}.
\label{eq:MB-recursion}
\end{equation}
Projecting onto the odd line via
\[
a_k:=u_k^\top X_k
\]
gives the scalar affine law
\begin{equation}
a_{k+1}
=
\rho_k a_k+\beta_k+\xi_k+\varepsilon_k^{\rm mix},
\label{eq:B-recursion-general}
\end{equation}
where
\[
\rho_k:=u_{k+1}^\top A_k u_k,\qquad
\beta_k:=u_{k+1}^\top B_k j_k^{\rm ctl},
\]
\[
\xi_k:=u_{k+1}^\top \nu_k,\qquad
\varepsilon_k^{\rm mix}:=u_{k+1}^\top A_k P_k^\perp X_k.
\]
Moreover,
\[
\mathbb E[\xi_k\mid\mathcal F_k]=0,
\qquad
\mathbb E[\xi_k^2\mid\mathcal F_k]=u_{k+1}^\top Q_k u_{k+1}\ge 0.
\]
If \cref{ass:classical-gate} also holds and the welded shape family is fixed, then
\begin{equation}
\mathbb E[\xi_k^2\mid\mathcal F_k]
=
\eta_k\,g_{\rm cl}(k)\,\Delta I_k+O(\Delta I_k^2,\delta\Lambda_k^2),
\label{eq:Zk-linear}
\end{equation}
with \(\eta_k\ge 0\).
\end{proposition}

\begin{proof}
The quadratic SK influence functional determines a Gaussian CPTP map on the finite active completion-mode sector, hence a
linear-Gaussian state-space realization of the form \eqref{eq:X-recursion}; \eqref{eq:MB-recursion} is the standard update
for the PSD second moment.
Decompose \(X_k=u_k a_k+P_k^\perp X_k\), apply \(\ip{u_{k+1}}{\cdot}\) to \eqref{eq:X-recursion}, and collect the projected
transport, deterministic drive, noise, and orthogonal-mode contamination to obtain \eqref{eq:B-recursion-general}.
The mean-zero and positive-variance statements for \(\xi_k\) follow directly from the projected Gaussian noise covariance
\(
u_{k+1}^\top Q_k u_{k+1}\ge 0
\).
Finally, at fixed welded shape family one has \(A_{{\rm eff},k}^2\propto \Delta I_k\), so the fresh covariance injection
into the odd mode is linear in the current payload at leading order; the gate \(g_{\rm cl}(k)\) suppresses this injection
near a fresh-aeon start and yields \eqref{eq:Zk-linear}.
\end{proof}

\begin{assumption}[Centered odd preparation]
\label{ass:centered-odd}
Near split onset, the active odd mode is prepared either branch-symmetrically or in a squeeze-dominated regime, so the
projected deterministic displacement vanishes and orthogonal-mode contamination is higher order:
\[
\beta_k=0,
\qquad
\varepsilon_k^{\rm mix}=O_2.
\]
\end{assumption}

\begin{corollary}[Centered odd-mode law]
\label{cor:centered-odd-mode}
Under
\cref{ass:gaussian-completion,ass:export-selection,ass:classical-gate,ass:centered-odd},
the projected affine law \eqref{eq:B-recursion-general} reduces to
\begin{equation}
a_{k+1}=\rho_k a_k+\xi_k+O_2,
\label{eq:B-recursion}
\end{equation}
with \(\mathbb E[\xi_k\mid\mathcal F_k]=0\) and variance scaling \eqref{eq:Zk-linear}.
\end{corollary}

\paragraph{Stochastic/Gaussian realization layer.}
At theorem level the positive stock is simply
\[
S_k:=a_k^2.
\]
The general projected Gaussian dynamics need not be centered.
Under \cref{ass:centered-odd}, however, the physically relevant late-time regime is the one in which the odd mode carries no
leading deterministic displacement and the fresh effect enters through its variance.
In that centered specialization it is natural to use the conditional second moment
\[
S_k:=\mathbb E[a_k^2\mid\mathcal F_k].
\]

\begin{proposition}[Odd-mode stock law]
\label{prop:odd-mode-stock}
Under
\cref{ass:gaussian-completion,ass:export-selection,ass:classical-gate,ass:centered-odd},
the conditional stock
\[
S_k:=\mathbb E[a_k^2\mid\mathcal F_k]
\]
satisfies
\begin{equation}
S_{k+1}
=
t_k\,S_k+\eta_k\,g_{\rm cl}(k)\,\Delta I_k
+O(S_k\Delta I_k,\Delta I_k^2,\delta\Lambda_k^2),
\qquad t_k=\rho_k^2.
\label{eq:Sk-update}
\end{equation}
\end{proposition}

\begin{proof}
Square \eqref{eq:B-recursion}:
\[
a_{k+1}^2=\rho_k^2 a_k^2+2\rho_k a_k\xi_k+\xi_k^2+O_2.
\]
Taking conditional expectation and using \(\mathbb E[\xi_k\mid\mathcal F_k]=0\) removes the mixed term.
Substituting \eqref{eq:Zk-linear} gives \eqref{eq:Sk-update}.
\end{proof}

\begin{assumption}[UV comparability of energetic stiffness]
\label{ass:uv-match}
There exist constants \(c_-,c_+>0\) and a monotone increasing energetic stiffness \(\widetilde\alpha(\Lambda)\) of the
active odd mode such that, for all sufficiently large \(\Lambda\),
\begin{equation}
c_-\,\alpha(\Lambda)\le \widetilde\alpha(\Lambda)\le c_+\,\alpha(\Lambda).
\label{eq:alpha-comparable}
\end{equation}
\end{assumption}

This is the precise place where the note passes from what is derived in UOM to a physical hypothesis.
The coefficient \(\alpha(\Lambda)\) is derived from the continued OS problem; the statement that the energetic stiffness of
the split mode is comparable to it is not yet a theorem.

With the stock \(S_k\) as the actual state variable, the positive split-driving export flux is modeled by
\begin{equation}
Q_{{\rm leak},k}^{\rm model}
=
\widetilde\alpha'(\Lambda_k)\,\delta\Lambda_k\,S_k
\;+\;
\widetilde\alpha(\Lambda_k)\,\eta_k\,g_{\rm cl}(k)\,\Delta I_k
\;+\; O_2,
\label{eq:Qleak-effective}
\end{equation}
where \(\delta\Lambda_k:=\Lambda_{k+1}-\Lambda_k\).
The first term is the transport cost of carrying already-prepared odd stock to larger cutoff; the second is the fresh
injection due to the current payload.
Under \cref{ass:uv-match}, this gives the scaling comparison
\begin{equation}
Q_{{\rm leak},k}^{\rm model}
\asymp
\alpha'(\Lambda_k)\,\delta\Lambda_k\,S_k
\;+\;
\alpha(\Lambda_k)\,\eta_k\,g_{\rm cl}(k)\,\Delta I_k.
\label{eq:Qleak-alpha}
\end{equation}

\paragraph{Current RP proxy reinterpreted.}
The quantity
\[
Q_{\rm rg}=\chi J_{\rm RP}/\beta_{\rm dS}
\]
therefore remains useful only as a phenomenological probe. It may correlate with \(a_k^2\) or with the stock \(S_k\), but
it is not the fundamental leak law because the healthy receiver itself has no quadratic RP block on the welded family.

\subsection{Approximate split condition from the odd-stock leak law}
\label{subsec:split-threshold}

Let \(R_k^{\rm abs}:=|R_k|\). In the micro-step regime the geometry-side surrogate gives
\[
\Delta E_{{\rm geom},k}
\approx
\frac{\zeta\,(A+c_\Lambda \Lambda_k^2)}{2\sqrt{12\,R_k^{\rm abs}}}\;\Theta_k
\;-\;
\frac{2\zeta\,c_\Lambda \Lambda_k\sqrt{R_k^{\rm abs}}}{\sqrt{12}}\;\delta\Lambda,
\]
while
\[
J_{{\rm wall},k}\approx \mu_{\rm RG}(\Lambda_k)\,\delta\Lambda,
\qquad
\mu_{\rm RG}(\Lambda)\simeq \mu_0+\mu_2\Lambda^2>0.
\]
Hence the local one-wall threshold remains
\begin{equation}
Q_{{\rm leak},k}
\gtrsim
\frac{\zeta\,(A+c_\Lambda \Lambda_k^2)}{2\sqrt{12\,R_k^{\rm abs}}}\;\Theta_k.
\label{eq:split-threshold-base}
\end{equation}
Since
\[
\Theta_k
=
K(\Lambda_k)\,\Delta I_k
\sim
K_0\,\Lambda_k^{-p}\,e^{-2d_{\rm sep}(\Lambda_k)\Lambda_k}\,\Delta I_k,
\qquad p\in\{10,14\},
\]
substituting \eqref{eq:Qleak-effective} gives
\begin{equation}
\widetilde\alpha'(\Lambda_k)\,\delta\Lambda_k\,S_k
\;+\;
\widetilde\alpha(\Lambda_k)\,\eta_k\,g_{\rm cl}(k)\,\Delta I_k
\gtrsim
\frac{\zeta\,(A+c_\Lambda \Lambda_k^2)}{2\sqrt{12\,R_k^{\rm abs}}}\,
K(\Lambda_k)\,\Delta I_k.
\label{eq:split-threshold-effective}
\end{equation}
This is the corrected split threshold. It shows the two late-time branches transparently:
\begin{enumerate}[leftmargin=2em]
\item if \(S_k\) remains negligible, only the fresh term remains, and the exhaustion branch can survive;
\item if \(S_k\) has accumulated, the transport term can dominate even when \(\Delta I_k\) is already tiny.
\end{enumerate}
Under \cref{ass:uv-match} and \(\alpha(\Lambda)\sim \kappa_d\Lambda^{d+2}H^{-d}\), the transport term scales like
\[
\widetilde\alpha'(\Lambda_k)\,\delta\Lambda_k\,S_k
\asymp
\Lambda_k^{d+1}\,\delta\Lambda_k\,S_k,
\]
which is exactly the mechanism by which late split can beat the shrinking fresh payload.

\subsection{Nondegeneracy clause for the normalized seam coefficient}
\label{subsec:separate-c-nondegeneracy}

The split threshold \eqref{eq:split-threshold-effective} is a statement about when the late-time
odd stock can overcome the shrinking fresh payload.
It is \emph{not} by itself a theorem that the normalized seam coefficient \(c_{\Lambda_k}\) on the
canonical compressed welded LoG line is nonzero.
Any theorem proving \(c_{\Lambda_k}\neq 0\) must therefore come from a separate overlap or transport
hypothesis, not from the threshold inequality alone.
Downstream, this clause is absorbed into the delayed selected-line package rather than carried as a
separate package summand.

\begin{definition}[Unit-source seam overlap scalar and no-destructive-cancellation margin]
\label{def:unit-source-overlap-margin}
Fix an accepted tick \(k\).
Assume the canonical compressed welded LoG line is defined and let
\[
\tau^{\mathrm{unit}}_{\Lambda_k}
\]
be a detector-factorized unit welded odd source.
Define the raw unit-source seam overlap scalar by
\begin{equation}
\widetilde c_{\Lambda_k}
:=
\big\langle\!\big\langle
\Pi_{\mathrm{phys},\Lambda_k}\mathsf T_{\Lambda_k}\!\big[\tau^{\mathrm{unit}}_{\Lambda_k}\big],\,
\delta C^{\mathrm{LoG}}_{\Lambda_k}
\big\rangle\!\big\rangle_{C_{\ast,\Lambda_k}}.
\label{eq:unit-source-overlap-scalar}
\end{equation}
A \emph{no-destructive-cancellation margin} at \(\Lambda_k\) is any real number
\[
\mu_{\Lambda_k}>0
\]
such that
\begin{equation}
\abs{\widetilde c_{\Lambda_k}}\ge \mu_{\Lambda_k}.
\label{eq:no-destr-cancel-margin}
\end{equation}
On a connected accepted-branch segment \(I\), a \emph{branchwise no-destructive-cancellation
margin} is a constant \(\mu_I>0\) such that
\begin{equation}
\abs{\widetilde c_{\Lambda}}\ge \mu_I
\qquad
\forall\,\Lambda\in I.
\label{eq:no-destr-cancel-margin-branch}
\end{equation}
\end{definition}

\begin{theorem}[Nondegeneracy criterion for \texorpdfstring{\(c_{\Lambda_k}\)}{c_Lambda_k}]
\label{thm:separate-c-nondegeneracy}
Fix an accepted tick \(k\).
Assume:
\begin{enumerate}[leftmargin=2em]
\item a detector-factorized unit welded odd source at \(\Lambda_k\);
\item the unit-source seam-faithfulness package, namely:
  \begin{enumerate}[label=(\alph*),leftmargin=2em]
  \item the raw seam pairing factorizes as
    \begin{equation}
    \widetilde s_{\Lambda_k}(f,y)
    =
    \sigma_{\mathrm{fill}}(f)\,
    \widetilde c_{\Lambda_k}\,
    \overline{\mathrm{obs}}(y)
    \qquad
    \forall f,\forall y;
    \label{eq:unit-source-seam-factorization}
    \end{equation}
  \item whenever
    \(
    \widetilde c_{\Lambda_k}\neq 0
    \)
    and
    \(
    \overline{\mathrm{obs}}(y)\neq 0
    \),
    the compressed welded LoG line is transversely physical, dominant-line realization holds, and
    the normalized seam coefficient
    \begin{equation}
    c_{\Lambda_k}
    :=
    \kappa_{\Lambda_k}^{-1/2}\,\widetilde c_{\Lambda_k}
    \label{eq:unit-source-normalized-c}
    \end{equation}
    is defined;
  \end{enumerate}
\item a no-destructive-cancellation margin at \(\Lambda_k\) in the sense of
  \cref{def:unit-source-overlap-margin};
\item a witness \(y\) with
  \[
  \overline{\mathrm{obs}}(y)\neq 0.
  \]
\end{enumerate}
Then:
\begin{enumerate}[leftmargin=2em]
\item the raw seam pairing is nonzero for every filler tag \(f\):
  \[
  \widetilde s_{\Lambda_k}(f,y)\neq 0;
  \]
\item the canonical compressed welded LoG line is transversely physical at \(\Lambda_k\);
\item dominant-line realization holds at \(\Lambda_k\);
\item the normalized seam coefficient is defined and nonzero:
  \[
  c_{\Lambda_k}\neq 0.
  \]
\end{enumerate}
If, in addition, for some filler tag \(f\) the source \(\tau^{f,y}_{\Lambda_k}\) is representable on
the characterization layer and admits a source-compatible descendant realization, then the
corresponding descendant representative satisfies
\[
\rho^{\mathrm{dom}}_{\Lambda_k}\!\big(J^{f,y}_{\Lambda_k}\big)\neq 0.
\]
\end{theorem}

\begin{proof}
The margin \eqref{eq:no-destr-cancel-margin} implies
\(
\widetilde c_{\Lambda_k}\neq 0
\).
Together with
\(
\overline{\mathrm{obs}}(y)\neq 0
\),
the factorization \eqref{eq:unit-source-seam-factorization} gives
\(
\widetilde s_{\Lambda_k}(f,y)\neq 0
\)
for every filler tag \(f\), proving item (1).
Item (2), item (3), and the definition and nonvanishing of
\(
c_{\Lambda_k}
\)
are exactly the remaining conclusions of the assumed unit-source seam-faithfulness package under the
same nondegeneracy hypothesis.
The final sentence is the corresponding witness-to-descendant comparison step under representability
and source-compatible descendant realization.
\end{proof}

\begin{corollary}[Branchwise sign-stable nondegeneracy]
\label{cor:branchwise-c-nondegeneracy}
Let \(I\) be a connected accepted-branch segment.
Assume:
\begin{enumerate}[leftmargin=2em]
\item the branchwise normalized seam coefficient
  \[
  \Lambda\mapsto c_{\Lambda}
  \]
  is continuous on \(I\);
\item a branchwise no-destructive-cancellation margin on \(I\) in the sense of
  \cref{def:unit-source-overlap-margin}.
\end{enumerate}
Then
\[
c_{\Lambda}\neq 0
\qquad
\forall\,\Lambda\in I.
\]
In particular, the sign of \(c_{\Lambda}\) is constant on \(I\).
If there also exist a fixed LS witness \(y\) with
\(
\overline{\mathrm{obs}}(y)\neq 0
\)
and a branchwise source-compatible descendant realization satisfying
\[
\rho^{\mathrm{dom}}_{\Lambda}\!\big(J^{f,y}_{\Lambda}\big)
=
\sigma_{\mathrm{fill}}(f)\,
c_{\Lambda}\,
\overline{\mathrm{obs}}(y),
\]
then the sign of
\(
\rho^{\mathrm{dom}}_{\Lambda}(J^{f,y}_{\Lambda})
\)
is constant on \(I\) as well.
\end{corollary}

\begin{proof}
By \eqref{eq:no-destr-cancel-margin-branch}, the scalar
\(
\widetilde c_{\Lambda}
\)
never vanishes on \(I\).
Hence neither does
\(
c_{\Lambda}
\),
because the normalization only rescales the raw overlap scalar by a positive branchwise factor.
A continuous nonzero real-valued function on a connected segment cannot change sign, proving the
constancy of \(\operatorname{sgn}(c_{\Lambda})\).
The final statement follows immediately from the displayed branchwise comparison identity and the
fixed nonzero scalar \(\overline{\mathrm{obs}}(y)\).
\end{proof}

\begin{remark}[Logical separation from the split threshold]
\label{rem:separate-c-nondegeneracy}
The hypotheses of \cref{thm:separate-c-nondegeneracy,cor:branchwise-c-nondegeneracy} are
intentionally separate from \eqref{eq:split-threshold-effective}.
The threshold governs the emergence of split-side stock; the no-destructive-cancellation margin
governs overlap of the unit welded odd source with the canonical compressed dominant line.
Accordingly, the present note treats \(c_{\Lambda}\neq 0\) as the transport/overlap clause later
imported into the delayed selected-line package rather than as a consequence of the split
threshold.
\end{remark}
%
}{%
  \IfFileExists{UOM/completion_sector_split_model_body.tex}{%
    \section{Completion-sector leak model and split criterion}
\label{sec:resolution}

This section separates what is directly derived from UOM from the additional hypotheses required to turn those ingredients
into a concrete split-driving leak model. The new organizing principle is:
the primary hidden variable is a \emph{signed odd split amplitude} \(a_k\), while the positive stock
\(S_k:=a_k^2\) is the quantity that can enter scalar observables and leak laws.
We reserve \(R_k^{\rm abs}\) for the remaining curvature budget and
\[
\mathfrak r_k:=\min_{\Lambda\ge \Lambda_k} f_k(\Lambda)
\]
for the minimized one-wall residual.

\subsection{Odd split mode, output parity, and the rank-one normal form}
\label{subsec:leak-solution}

The only ``leak'' that directly competes with wall work in the \(\Lambda\)-solver is the open-system modification of the
wall energy balance allowed explicitly in UOM's SK/HJ wall equation:
\begin{equation}
\Delta E_{{\rm geom},k}\ =\ J_{{\rm wall},k}\ +\ Q_{{\rm leak},k},
\qquad Q_{{\rm leak},k}\ge 0,
\label{eq:resolution-wall-leak}
\end{equation}
This is the welded wall-balance form used in the earlier summability analysis.
This is a \emph{receiver-side energy complement} term. It must not be conflated with UOM's information-theoretic
bridge/leak instrument from the UOM main manuscript.

\paragraph{What is already derived in UOM.}
Four ingredients are fixed by the manuscript:
\begin{enumerate}[leftmargin=2em]
\item \textbf{Wall balance.} The reduced one-wall description may contain a nonnegative complement \(Q_{\rm leak}\) as in
\eqref{eq:resolution-wall-leak}.
\item \textbf{Quadratic welded SK data.} In the Gaussian sector,
\[
W_\Lambda[J]=\frac12\int J\,(G_{R,\Lambda}+i\,\mathcal N_\Lambda)\,J+O(J^3),
\qquad \mathcal N_\Lambda\ge 0.
\]
\item \textbf{Healthy receiver.} The active KR receiver has vanishing welded RP block, so a receiver-tangent RP penalty is
not the fundamental home of split-driving oddness.
\item \textbf{Continued RP defect.} The analytic-continued covariance has a negative OS mode with
\[
\lambda_{\min}(\mu_{ps},\Lambda)
=
-\alpha(\Lambda)\mu_{ps}^2+O(\mu_{ps}^4),
\qquad
\alpha(\Lambda)\sim \kappa_d\,\Lambda^{d+2}H^{-d}.
\]
\end{enumerate}

\begin{assumption}[Split-mode observability]
\label{ass:wall-match}
Near split onset there exists a one-dimensional odd line \(L_k\) with local coordinate \(a_k\in\mathbb R\) such that the
split involution acts by
\[
a_k\longmapsto -a_k.
\]
The true two-wall advantage
\[
\Delta F_k:=\Delta\mathcal F_{2{\rm w}-1{\rm w},k}
\]
and the minimized one-wall residual
\[
\mathfrak r_k:=\min_{\Lambda\ge \Lambda_k} f_k(\Lambda)
\]
are smooth even scalar outputs of this same local mode.
Equivalently, the 2T \(\Gamma\)-odd descendant direction and the UOM continued negative mode \(u_-(\Lambda_k)\) are
assumed to represent the same active split mode up to a smooth local identification.
\end{assumption}

\begin{lemma}[Evenness of the minimized residual]
\label{lem:residual-even}
Suppose the pre-minimized one-wall residual is a smooth function \(f_k(\Lambda,a)\) satisfying
\[
f_k(\Lambda,-a)=f_k(\Lambda,a),
\]
and suppose there is a unique local minimizer \(\Lambda_\ast(a)\) near \(a=0\) such that
\[
\partial_\Lambda f_k(\Lambda_\ast(a),a)=0,
\qquad
\partial_\Lambda^2 f_k(\Lambda_\ast(0),0)>0.
\]
Then
\[
\mathfrak r_k(a):=f_k(\Lambda_\ast(a),a)
\]
is a smooth even function of \(a\).
\end{lemma}

\begin{proof}
By the implicit function theorem, \(\Lambda_\ast(a)\) is smooth near \(a=0\).
Since \(f_k(\Lambda,a)\) is even in \(a\), the stationarity equation for \(a\) and \(-a\) is identical; uniqueness of the
local minimizer gives
\[
\Lambda_\ast(-a)=\Lambda_\ast(a).
\]
Therefore
\[
\mathfrak r_k(-a)=f_k(\Lambda_\ast(-a),-a)=f_k(\Lambda_\ast(a),a)=\mathfrak r_k(a),
\]
so \(\mathfrak r_k\) is smooth and even.
\end{proof}

\begin{proposition}[Rank-one normal form near split onset]
\label{prop:rank-one-wall-match}
Under \cref{ass:wall-match}, the two scalar outputs admit local even expansions
\begin{equation}
\Delta F_k(a)
=
\mu_k^{(2{\rm w})}+h_k a^2+O(a^4),
\label{eq:DeltaF-detuned}
\end{equation}
\begin{equation}
\mathfrak r_k(a)
=
\mu_k^{(1{\rm w})}+m_k a^2+O(a^4),
\label{eq:rres-detuned}
\end{equation}
with \(h_k,m_k>0\) on the active split line.
\end{proposition}

\begin{proof}
Smooth even functions of one real variable have Taylor series containing only even powers.
The constants \(\mu_k^{(2{\rm w})}\) and \(\mu_k^{(1{\rm w})}\) are the baseline detunings of the two-wall and one-wall
channels. Positivity of \(h_k\) and \(m_k\) expresses that the selected mode is active in both outputs.
\end{proof}

\begin{corollary}[Onset wall-match in rank one]
\label{cor:rank-one-wall-match}
On the onset manifold,
\[
\mu_k^{(1{\rm w})}=\mu_k^{(2{\rm w})}=0,
\]
so
\[
\Delta F_k=h_k a_k^2+O(a_k^4),
\qquad
\mathfrak r_k=m_k a_k^2+O(a_k^4).
\]
Eliminating \(a_k^2\) gives
\begin{equation}
\Delta F_k
=
\frac{h_k}{m_k}\,\mathfrak r_k+O(\mathfrak r_k^2).
\label{eq:rank-one-wall-match}
\end{equation}
\end{corollary}

\subsection{Low-rank generalization}
\label{subsec:alpha-derived}

If the active split sector is \(m\)-dimensional, with odd amplitude vector \(a_k\in\mathbb R^m\), parity gives quadratic
normal forms
\begin{equation}
\Delta F_k(a)
=
\mu_k^{(2{\rm w})}+a^\top H_k a+O(\|a\|^4),
\label{eq:DeltaF-lowrank}
\end{equation}
\begin{equation}
\mathfrak r_k(a)
=
\mu_k^{(1{\rm w})}+a^\top M_k a+O(\|a\|^4),
\label{eq:rres-lowrank}
\end{equation}
with \(H_k,M_k\succeq 0\) on the active split subspace.

\begin{assumption}[Low-rank comparability]
\label{ass:low-rank-comp}
On the active split subspace there exist \(c_k^-,c_k^+>0\) such that
\[
c_k^- M_k \preceq H_k \preceq c_k^+ M_k.
\]
\end{assumption}

\begin{proposition}[Low-rank wall-match comparability]
\label{prop:low-rank-wall-match}
Assume \cref{ass:wall-match,ass:low-rank-comp} and restrict to the onset manifold
\(\mu_k^{(1{\rm w})}=\mu_k^{(2{\rm w})}=0\).
Then for sufficiently small \(\mathfrak r_k\) there exists \(C_k>0\) such that
\begin{equation}
c_k^-\,\mathfrak r_k-C_k\mathfrak r_k^2
\le
\Delta F_k
\le
c_k^+\,\mathfrak r_k+C_k\mathfrak r_k^2.
\label{eq:low-rank-wall-match}
\end{equation}
\end{proposition}

\begin{proof}
On onset,
\[
\Delta F_k=a_k^\top H_k a_k+O(\|a_k\|^4),
\qquad
\mathfrak r_k=a_k^\top M_k a_k+O(\|a_k\|^4).
\]
The quadratic-form comparison gives
\[
c_k^- a_k^\top M_k a_k\le a_k^\top H_k a_k\le c_k^+ a_k^\top M_k a_k.
\]
Since \(\mathfrak r_k\) is itself quadratic in \(a_k\) to leading order, the quartic remainders are \(O(\mathfrak r_k^2)\),
which yields \eqref{eq:low-rank-wall-match}.
\end{proof}

\subsection{Gaussian realization of the odd mode and the leak stock law}
\label{subsec:B-gaussian}

The theorem path above does \emph{not} rely on a KMS-subtracted matrix.
In particular, we do not use
\[
\Sigma_{B,k}:=\bigl(C_{B,k}-C^{\rm KMS}_{\Lambda_k}\bigr)+m_{B,k}m_{B,k}^\top
\]
as a positive object, because it need not be positive semidefinite.
When a safe matrix observable is needed in the realization layer, the natural PSD quantity is
\begin{equation}
M_{B,k}:=C_{B,k}+m_{B,k}m_{B,k}^\top\succeq 0.
\label{eq:MB-def}
\end{equation}

\begin{assumption}[Gaussian completion realization]
\label{ass:gaussian-completion}
For each admissible cutoff \(\Lambda\) in the welded regime, the quadratic influence functional \(W_\Lambda[J]\) admits a
Gaussian dilation, unique up to symplectic equivalence, by a bosonic completion sector \(B_\Lambda\) with canonical
quadratures \(X_\Lambda\) and positive quadratic Hamiltonian
\[
H_{B,\Lambda}=\frac12 X_\Lambda^\top h_\Lambda X_\Lambda,
\qquad h_\Lambda\succeq 0.
\]
\end{assumption}

\begin{assumption}[Dominant odd-mode projection]
\label{ass:export-selection}
Inside the Gaussian completion of \cref{ass:gaussian-completion}, there exists a smooth unit vector
\[
u_k:=u_-(\Lambda_k)
\]
spanning the active odd line and separated from the orthogonal sector by a local spectral gap.
We write
\[
P_k^\perp:=I-u_k u_k^\top
\]
for the orthogonal projector.
\end{assumption}

\begin{assumption}[Classical-control gate]
\label{ass:classical-gate}
There exists a scalar gate \(g_{\rm cl}(k)\in[0,1]\), determined by the accessible classicality defect
\(\mathrm{cl}_k\) or \(\mathrm{Off}_k\), such that \(g_{\rm cl}(k)\ll 1\) at a fresh-aeon start and
\(g_{\rm cl}(k)\to 1\) late in the aeon.
\end{assumption}

\begin{proposition}[Projected Gaussian affine dynamics]
\label{prop:projected-affine}
Under \cref{ass:gaussian-completion,ass:export-selection}, one may choose a finite effective completion-mode basis in which
the completion quadratures obey a linear-Gaussian recursion
\begin{equation}
X_{k+1}=A_k X_k+B_k j_k^{\rm ctl}+\nu_k,
\label{eq:X-recursion}
\end{equation}
with
\[
\mathbb E[\nu_k\mid\mathcal F_k]=0,
\qquad
\mathbb E[\nu_k\nu_k^\top\mid\mathcal F_k]=Q_k\succeq 0.
\]
Consequently the PSD second moment
\[
M_{B,k}=C_{B,k}+m_{B,k}m_{B,k}^\top
\]
obeys
\begin{equation}
M_{B,k+1}=A_k M_{B,k}A_k^\top+Q_k+\text{drift terms}.
\label{eq:MB-recursion}
\end{equation}
Projecting onto the odd line via
\[
a_k:=u_k^\top X_k
\]
gives the scalar affine law
\begin{equation}
a_{k+1}
=
\rho_k a_k+\beta_k+\xi_k+\varepsilon_k^{\rm mix},
\label{eq:B-recursion-general}
\end{equation}
where
\[
\rho_k:=u_{k+1}^\top A_k u_k,\qquad
\beta_k:=u_{k+1}^\top B_k j_k^{\rm ctl},
\]
\[
\xi_k:=u_{k+1}^\top \nu_k,\qquad
\varepsilon_k^{\rm mix}:=u_{k+1}^\top A_k P_k^\perp X_k.
\]
Moreover,
\[
\mathbb E[\xi_k\mid\mathcal F_k]=0,
\qquad
\mathbb E[\xi_k^2\mid\mathcal F_k]=u_{k+1}^\top Q_k u_{k+1}\ge 0.
\]
If \cref{ass:classical-gate} also holds and the welded shape family is fixed, then
\begin{equation}
\mathbb E[\xi_k^2\mid\mathcal F_k]
=
\eta_k\,g_{\rm cl}(k)\,\Delta I_k+O(\Delta I_k^2,\delta\Lambda_k^2),
\label{eq:Zk-linear}
\end{equation}
with \(\eta_k\ge 0\).
\end{proposition}

\begin{proof}
The quadratic SK influence functional determines a Gaussian CPTP map on the finite active completion-mode sector, hence a
linear-Gaussian state-space realization of the form \eqref{eq:X-recursion}; \eqref{eq:MB-recursion} is the standard update
for the PSD second moment.
Decompose \(X_k=u_k a_k+P_k^\perp X_k\), apply \(\ip{u_{k+1}}{\cdot}\) to \eqref{eq:X-recursion}, and collect the projected
transport, deterministic drive, noise, and orthogonal-mode contamination to obtain \eqref{eq:B-recursion-general}.
The mean-zero and positive-variance statements for \(\xi_k\) follow directly from the projected Gaussian noise covariance
\(
u_{k+1}^\top Q_k u_{k+1}\ge 0
\).
Finally, at fixed welded shape family one has \(A_{{\rm eff},k}^2\propto \Delta I_k\), so the fresh covariance injection
into the odd mode is linear in the current payload at leading order; the gate \(g_{\rm cl}(k)\) suppresses this injection
near a fresh-aeon start and yields \eqref{eq:Zk-linear}.
\end{proof}

\begin{assumption}[Centered odd preparation]
\label{ass:centered-odd}
Near split onset, the active odd mode is prepared either branch-symmetrically or in a squeeze-dominated regime, so the
projected deterministic displacement vanishes and orthogonal-mode contamination is higher order:
\[
\beta_k=0,
\qquad
\varepsilon_k^{\rm mix}=O_2.
\]
\end{assumption}

\begin{corollary}[Centered odd-mode law]
\label{cor:centered-odd-mode}
Under
\cref{ass:gaussian-completion,ass:export-selection,ass:classical-gate,ass:centered-odd},
the projected affine law \eqref{eq:B-recursion-general} reduces to
\begin{equation}
a_{k+1}=\rho_k a_k+\xi_k+O_2,
\label{eq:B-recursion}
\end{equation}
with \(\mathbb E[\xi_k\mid\mathcal F_k]=0\) and variance scaling \eqref{eq:Zk-linear}.
\end{corollary}

\paragraph{Stochastic/Gaussian realization layer.}
At theorem level the positive stock is simply
\[
S_k:=a_k^2.
\]
The general projected Gaussian dynamics need not be centered.
Under \cref{ass:centered-odd}, however, the physically relevant late-time regime is the one in which the odd mode carries no
leading deterministic displacement and the fresh effect enters through its variance.
In that centered specialization it is natural to use the conditional second moment
\[
S_k:=\mathbb E[a_k^2\mid\mathcal F_k].
\]

\begin{proposition}[Odd-mode stock law]
\label{prop:odd-mode-stock}
Under
\cref{ass:gaussian-completion,ass:export-selection,ass:classical-gate,ass:centered-odd},
the conditional stock
\[
S_k:=\mathbb E[a_k^2\mid\mathcal F_k]
\]
satisfies
\begin{equation}
S_{k+1}
=
t_k\,S_k+\eta_k\,g_{\rm cl}(k)\,\Delta I_k
+O(S_k\Delta I_k,\Delta I_k^2,\delta\Lambda_k^2),
\qquad t_k=\rho_k^2.
\label{eq:Sk-update}
\end{equation}
\end{proposition}

\begin{proof}
Square \eqref{eq:B-recursion}:
\[
a_{k+1}^2=\rho_k^2 a_k^2+2\rho_k a_k\xi_k+\xi_k^2+O_2.
\]
Taking conditional expectation and using \(\mathbb E[\xi_k\mid\mathcal F_k]=0\) removes the mixed term.
Substituting \eqref{eq:Zk-linear} gives \eqref{eq:Sk-update}.
\end{proof}

\begin{assumption}[UV comparability of energetic stiffness]
\label{ass:uv-match}
There exist constants \(c_-,c_+>0\) and a monotone increasing energetic stiffness \(\widetilde\alpha(\Lambda)\) of the
active odd mode such that, for all sufficiently large \(\Lambda\),
\begin{equation}
c_-\,\alpha(\Lambda)\le \widetilde\alpha(\Lambda)\le c_+\,\alpha(\Lambda).
\label{eq:alpha-comparable}
\end{equation}
\end{assumption}

This is the precise place where the note passes from what is derived in UOM to a physical hypothesis.
The coefficient \(\alpha(\Lambda)\) is derived from the continued OS problem; the statement that the energetic stiffness of
the split mode is comparable to it is not yet a theorem.

With the stock \(S_k\) as the actual state variable, the positive split-driving export flux is modeled by
\begin{equation}
Q_{{\rm leak},k}^{\rm model}
=
\widetilde\alpha'(\Lambda_k)\,\delta\Lambda_k\,S_k
\;+\;
\widetilde\alpha(\Lambda_k)\,\eta_k\,g_{\rm cl}(k)\,\Delta I_k
\;+\; O_2,
\label{eq:Qleak-effective}
\end{equation}
where \(\delta\Lambda_k:=\Lambda_{k+1}-\Lambda_k\).
The first term is the transport cost of carrying already-prepared odd stock to larger cutoff; the second is the fresh
injection due to the current payload.
Under \cref{ass:uv-match}, this gives the scaling comparison
\begin{equation}
Q_{{\rm leak},k}^{\rm model}
\asymp
\alpha'(\Lambda_k)\,\delta\Lambda_k\,S_k
\;+\;
\alpha(\Lambda_k)\,\eta_k\,g_{\rm cl}(k)\,\Delta I_k.
\label{eq:Qleak-alpha}
\end{equation}

\paragraph{Current RP proxy reinterpreted.}
The quantity
\[
Q_{\rm rg}=\chi J_{\rm RP}/\beta_{\rm dS}
\]
therefore remains useful only as a phenomenological probe. It may correlate with \(a_k^2\) or with the stock \(S_k\), but
it is not the fundamental leak law because the healthy receiver itself has no quadratic RP block on the welded family.

\subsection{Approximate split condition from the odd-stock leak law}
\label{subsec:split-threshold}

Let \(R_k^{\rm abs}:=|R_k|\). In the micro-step regime the geometry-side surrogate gives
\[
\Delta E_{{\rm geom},k}
\approx
\frac{\zeta\,(A+c_\Lambda \Lambda_k^2)}{2\sqrt{12\,R_k^{\rm abs}}}\;\Theta_k
\;-\;
\frac{2\zeta\,c_\Lambda \Lambda_k\sqrt{R_k^{\rm abs}}}{\sqrt{12}}\;\delta\Lambda,
\]
while
\[
J_{{\rm wall},k}\approx \mu_{\rm RG}(\Lambda_k)\,\delta\Lambda,
\qquad
\mu_{\rm RG}(\Lambda)\simeq \mu_0+\mu_2\Lambda^2>0.
\]
Hence the local one-wall threshold remains
\begin{equation}
Q_{{\rm leak},k}
\gtrsim
\frac{\zeta\,(A+c_\Lambda \Lambda_k^2)}{2\sqrt{12\,R_k^{\rm abs}}}\;\Theta_k.
\label{eq:split-threshold-base}
\end{equation}
Since
\[
\Theta_k
=
K(\Lambda_k)\,\Delta I_k
\sim
K_0\,\Lambda_k^{-p}\,e^{-2d_{\rm sep}(\Lambda_k)\Lambda_k}\,\Delta I_k,
\qquad p\in\{10,14\},
\]
substituting \eqref{eq:Qleak-effective} gives
\begin{equation}
\widetilde\alpha'(\Lambda_k)\,\delta\Lambda_k\,S_k
\;+\;
\widetilde\alpha(\Lambda_k)\,\eta_k\,g_{\rm cl}(k)\,\Delta I_k
\gtrsim
\frac{\zeta\,(A+c_\Lambda \Lambda_k^2)}{2\sqrt{12\,R_k^{\rm abs}}}\,
K(\Lambda_k)\,\Delta I_k.
\label{eq:split-threshold-effective}
\end{equation}
This is the corrected split threshold. It shows the two late-time branches transparently:
\begin{enumerate}[leftmargin=2em]
\item if \(S_k\) remains negligible, only the fresh term remains, and the exhaustion branch can survive;
\item if \(S_k\) has accumulated, the transport term can dominate even when \(\Delta I_k\) is already tiny.
\end{enumerate}
Under \cref{ass:uv-match} and \(\alpha(\Lambda)\sim \kappa_d\Lambda^{d+2}H^{-d}\), the transport term scales like
\[
\widetilde\alpha'(\Lambda_k)\,\delta\Lambda_k\,S_k
\asymp
\Lambda_k^{d+1}\,\delta\Lambda_k\,S_k,
\]
which is exactly the mechanism by which late split can beat the shrinking fresh payload.

\subsection{Nondegeneracy clause for the normalized seam coefficient}
\label{subsec:separate-c-nondegeneracy}

The split threshold \eqref{eq:split-threshold-effective} is a statement about when the late-time
odd stock can overcome the shrinking fresh payload.
It is \emph{not} by itself a theorem that the normalized seam coefficient \(c_{\Lambda_k}\) on the
canonical compressed welded LoG line is nonzero.
Any theorem proving \(c_{\Lambda_k}\neq 0\) must therefore come from a separate overlap or transport
hypothesis, not from the threshold inequality alone.
Downstream, this clause is absorbed into the delayed selected-line package rather than carried as a
separate package summand.

\begin{definition}[Unit-source seam overlap scalar and no-destructive-cancellation margin]
\label{def:unit-source-overlap-margin}
Fix an accepted tick \(k\).
Assume the canonical compressed welded LoG line is defined and let
\[
\tau^{\mathrm{unit}}_{\Lambda_k}
\]
be a detector-factorized unit welded odd source.
Define the raw unit-source seam overlap scalar by
\begin{equation}
\widetilde c_{\Lambda_k}
:=
\big\langle\!\big\langle
\Pi_{\mathrm{phys},\Lambda_k}\mathsf T_{\Lambda_k}\!\big[\tau^{\mathrm{unit}}_{\Lambda_k}\big],\,
\delta C^{\mathrm{LoG}}_{\Lambda_k}
\big\rangle\!\big\rangle_{C_{\ast,\Lambda_k}}.
\label{eq:unit-source-overlap-scalar}
\end{equation}
A \emph{no-destructive-cancellation margin} at \(\Lambda_k\) is any real number
\[
\mu_{\Lambda_k}>0
\]
such that
\begin{equation}
\abs{\widetilde c_{\Lambda_k}}\ge \mu_{\Lambda_k}.
\label{eq:no-destr-cancel-margin}
\end{equation}
On a connected accepted-branch segment \(I\), a \emph{branchwise no-destructive-cancellation
margin} is a constant \(\mu_I>0\) such that
\begin{equation}
\abs{\widetilde c_{\Lambda}}\ge \mu_I
\qquad
\forall\,\Lambda\in I.
\label{eq:no-destr-cancel-margin-branch}
\end{equation}
\end{definition}

\begin{theorem}[Nondegeneracy criterion for \texorpdfstring{\(c_{\Lambda_k}\)}{c_Lambda_k}]
\label{thm:separate-c-nondegeneracy}
Fix an accepted tick \(k\).
Assume:
\begin{enumerate}[leftmargin=2em]
\item a detector-factorized unit welded odd source at \(\Lambda_k\);
\item the unit-source seam-faithfulness package, namely:
  \begin{enumerate}[label=(\alph*),leftmargin=2em]
  \item the raw seam pairing factorizes as
    \begin{equation}
    \widetilde s_{\Lambda_k}(f,y)
    =
    \sigma_{\mathrm{fill}}(f)\,
    \widetilde c_{\Lambda_k}\,
    \overline{\mathrm{obs}}(y)
    \qquad
    \forall f,\forall y;
    \label{eq:unit-source-seam-factorization}
    \end{equation}
  \item whenever
    \(
    \widetilde c_{\Lambda_k}\neq 0
    \)
    and
    \(
    \overline{\mathrm{obs}}(y)\neq 0
    \),
    the compressed welded LoG line is transversely physical, dominant-line realization holds, and
    the normalized seam coefficient
    \begin{equation}
    c_{\Lambda_k}
    :=
    \kappa_{\Lambda_k}^{-1/2}\,\widetilde c_{\Lambda_k}
    \label{eq:unit-source-normalized-c}
    \end{equation}
    is defined;
  \end{enumerate}
\item a no-destructive-cancellation margin at \(\Lambda_k\) in the sense of
  \cref{def:unit-source-overlap-margin};
\item a witness \(y\) with
  \[
  \overline{\mathrm{obs}}(y)\neq 0.
  \]
\end{enumerate}
Then:
\begin{enumerate}[leftmargin=2em]
\item the raw seam pairing is nonzero for every filler tag \(f\):
  \[
  \widetilde s_{\Lambda_k}(f,y)\neq 0;
  \]
\item the canonical compressed welded LoG line is transversely physical at \(\Lambda_k\);
\item dominant-line realization holds at \(\Lambda_k\);
\item the normalized seam coefficient is defined and nonzero:
  \[
  c_{\Lambda_k}\neq 0.
  \]
\end{enumerate}
If, in addition, for some filler tag \(f\) the source \(\tau^{f,y}_{\Lambda_k}\) is representable on
the characterization layer and admits a source-compatible descendant realization, then the
corresponding descendant representative satisfies
\[
\rho^{\mathrm{dom}}_{\Lambda_k}\!\big(J^{f,y}_{\Lambda_k}\big)\neq 0.
\]
\end{theorem}

\begin{proof}
The margin \eqref{eq:no-destr-cancel-margin} implies
\(
\widetilde c_{\Lambda_k}\neq 0
\).
Together with
\(
\overline{\mathrm{obs}}(y)\neq 0
\),
the factorization \eqref{eq:unit-source-seam-factorization} gives
\(
\widetilde s_{\Lambda_k}(f,y)\neq 0
\)
for every filler tag \(f\), proving item (1).
Item (2), item (3), and the definition and nonvanishing of
\(
c_{\Lambda_k}
\)
are exactly the remaining conclusions of the assumed unit-source seam-faithfulness package under the
same nondegeneracy hypothesis.
The final sentence is the corresponding witness-to-descendant comparison step under representability
and source-compatible descendant realization.
\end{proof}

\begin{corollary}[Branchwise sign-stable nondegeneracy]
\label{cor:branchwise-c-nondegeneracy}
Let \(I\) be a connected accepted-branch segment.
Assume:
\begin{enumerate}[leftmargin=2em]
\item the branchwise normalized seam coefficient
  \[
  \Lambda\mapsto c_{\Lambda}
  \]
  is continuous on \(I\);
\item a branchwise no-destructive-cancellation margin on \(I\) in the sense of
  \cref{def:unit-source-overlap-margin}.
\end{enumerate}
Then
\[
c_{\Lambda}\neq 0
\qquad
\forall\,\Lambda\in I.
\]
In particular, the sign of \(c_{\Lambda}\) is constant on \(I\).
If there also exist a fixed LS witness \(y\) with
\(
\overline{\mathrm{obs}}(y)\neq 0
\)
and a branchwise source-compatible descendant realization satisfying
\[
\rho^{\mathrm{dom}}_{\Lambda}\!\big(J^{f,y}_{\Lambda}\big)
=
\sigma_{\mathrm{fill}}(f)\,
c_{\Lambda}\,
\overline{\mathrm{obs}}(y),
\]
then the sign of
\(
\rho^{\mathrm{dom}}_{\Lambda}(J^{f,y}_{\Lambda})
\)
is constant on \(I\) as well.
\end{corollary}

\begin{proof}
By \eqref{eq:no-destr-cancel-margin-branch}, the scalar
\(
\widetilde c_{\Lambda}
\)
never vanishes on \(I\).
Hence neither does
\(
c_{\Lambda}
\),
because the normalization only rescales the raw overlap scalar by a positive branchwise factor.
A continuous nonzero real-valued function on a connected segment cannot change sign, proving the
constancy of \(\operatorname{sgn}(c_{\Lambda})\).
The final statement follows immediately from the displayed branchwise comparison identity and the
fixed nonzero scalar \(\overline{\mathrm{obs}}(y)\).
\end{proof}

\begin{remark}[Logical separation from the split threshold]
\label{rem:separate-c-nondegeneracy}
The hypotheses of \cref{thm:separate-c-nondegeneracy,cor:branchwise-c-nondegeneracy} are
intentionally separate from \eqref{eq:split-threshold-effective}.
The threshold governs the emergence of split-side stock; the no-destructive-cancellation margin
governs overlap of the unit welded odd source with the canonical compressed dominant line.
Accordingly, the present note treats \(c_{\Lambda}\neq 0\) as the transport/overlap clause later
imported into the delayed selected-line package rather than as a consequence of the split
threshold.
\end{remark}
%
  }{%
    \section{Completion-sector leak model and split criterion}
\label{sec:resolution}

This section separates what is directly derived from UOM from the additional hypotheses required to turn those ingredients
into a concrete split-driving leak model. The new organizing principle is:
the primary hidden variable is a \emph{signed odd split amplitude} \(a_k\), while the positive stock
\(S_k:=a_k^2\) is the quantity that can enter scalar observables and leak laws.
We reserve \(R_k^{\rm abs}\) for the remaining curvature budget and
\[
\mathfrak r_k:=\min_{\Lambda\ge \Lambda_k} f_k(\Lambda)
\]
for the minimized one-wall residual.

\subsection{Odd split mode, output parity, and the rank-one normal form}
\label{subsec:leak-solution}

The only ``leak'' that directly competes with wall work in the \(\Lambda\)-solver is the open-system modification of the
wall energy balance allowed explicitly in UOM's SK/HJ wall equation:
\begin{equation}
\Delta E_{{\rm geom},k}\ =\ J_{{\rm wall},k}\ +\ Q_{{\rm leak},k},
\qquad Q_{{\rm leak},k}\ge 0,
\label{eq:resolution-wall-leak}
\end{equation}
This is the welded wall-balance form used in the earlier summability analysis.
This is a \emph{receiver-side energy complement} term. It must not be conflated with UOM's information-theoretic
bridge/leak instrument from the UOM main manuscript.

\paragraph{What is already derived in UOM.}
Four ingredients are fixed by the manuscript:
\begin{enumerate}[leftmargin=2em]
\item \textbf{Wall balance.} The reduced one-wall description may contain a nonnegative complement \(Q_{\rm leak}\) as in
\eqref{eq:resolution-wall-leak}.
\item \textbf{Quadratic welded SK data.} In the Gaussian sector,
\[
W_\Lambda[J]=\frac12\int J\,(G_{R,\Lambda}+i\,\mathcal N_\Lambda)\,J+O(J^3),
\qquad \mathcal N_\Lambda\ge 0.
\]
\item \textbf{Healthy receiver.} The active KR receiver has vanishing welded RP block, so a receiver-tangent RP penalty is
not the fundamental home of split-driving oddness.
\item \textbf{Continued RP defect.} The analytic-continued covariance has a negative OS mode with
\[
\lambda_{\min}(\mu_{ps},\Lambda)
=
-\alpha(\Lambda)\mu_{ps}^2+O(\mu_{ps}^4),
\qquad
\alpha(\Lambda)\sim \kappa_d\,\Lambda^{d+2}H^{-d}.
\]
\end{enumerate}

\begin{assumption}[Split-mode observability]
\label{ass:wall-match}
Near split onset there exists a one-dimensional odd line \(L_k\) with local coordinate \(a_k\in\mathbb R\) such that the
split involution acts by
\[
a_k\longmapsto -a_k.
\]
The true two-wall advantage
\[
\Delta F_k:=\Delta\mathcal F_{2{\rm w}-1{\rm w},k}
\]
and the minimized one-wall residual
\[
\mathfrak r_k:=\min_{\Lambda\ge \Lambda_k} f_k(\Lambda)
\]
are smooth even scalar outputs of this same local mode.
Equivalently, the 2T \(\Gamma\)-odd descendant direction and the UOM continued negative mode \(u_-(\Lambda_k)\) are
assumed to represent the same active split mode up to a smooth local identification.
\end{assumption}

\begin{lemma}[Evenness of the minimized residual]
\label{lem:residual-even}
Suppose the pre-minimized one-wall residual is a smooth function \(f_k(\Lambda,a)\) satisfying
\[
f_k(\Lambda,-a)=f_k(\Lambda,a),
\]
and suppose there is a unique local minimizer \(\Lambda_\ast(a)\) near \(a=0\) such that
\[
\partial_\Lambda f_k(\Lambda_\ast(a),a)=0,
\qquad
\partial_\Lambda^2 f_k(\Lambda_\ast(0),0)>0.
\]
Then
\[
\mathfrak r_k(a):=f_k(\Lambda_\ast(a),a)
\]
is a smooth even function of \(a\).
\end{lemma}

\begin{proof}
By the implicit function theorem, \(\Lambda_\ast(a)\) is smooth near \(a=0\).
Since \(f_k(\Lambda,a)\) is even in \(a\), the stationarity equation for \(a\) and \(-a\) is identical; uniqueness of the
local minimizer gives
\[
\Lambda_\ast(-a)=\Lambda_\ast(a).
\]
Therefore
\[
\mathfrak r_k(-a)=f_k(\Lambda_\ast(-a),-a)=f_k(\Lambda_\ast(a),a)=\mathfrak r_k(a),
\]
so \(\mathfrak r_k\) is smooth and even.
\end{proof}

\begin{proposition}[Rank-one normal form near split onset]
\label{prop:rank-one-wall-match}
Under \cref{ass:wall-match}, the two scalar outputs admit local even expansions
\begin{equation}
\Delta F_k(a)
=
\mu_k^{(2{\rm w})}+h_k a^2+O(a^4),
\label{eq:DeltaF-detuned}
\end{equation}
\begin{equation}
\mathfrak r_k(a)
=
\mu_k^{(1{\rm w})}+m_k a^2+O(a^4),
\label{eq:rres-detuned}
\end{equation}
with \(h_k,m_k>0\) on the active split line.
\end{proposition}

\begin{proof}
Smooth even functions of one real variable have Taylor series containing only even powers.
The constants \(\mu_k^{(2{\rm w})}\) and \(\mu_k^{(1{\rm w})}\) are the baseline detunings of the two-wall and one-wall
channels. Positivity of \(h_k\) and \(m_k\) expresses that the selected mode is active in both outputs.
\end{proof}

\begin{corollary}[Onset wall-match in rank one]
\label{cor:rank-one-wall-match}
On the onset manifold,
\[
\mu_k^{(1{\rm w})}=\mu_k^{(2{\rm w})}=0,
\]
so
\[
\Delta F_k=h_k a_k^2+O(a_k^4),
\qquad
\mathfrak r_k=m_k a_k^2+O(a_k^4).
\]
Eliminating \(a_k^2\) gives
\begin{equation}
\Delta F_k
=
\frac{h_k}{m_k}\,\mathfrak r_k+O(\mathfrak r_k^2).
\label{eq:rank-one-wall-match}
\end{equation}
\end{corollary}

\subsection{Low-rank generalization}
\label{subsec:alpha-derived}

If the active split sector is \(m\)-dimensional, with odd amplitude vector \(a_k\in\mathbb R^m\), parity gives quadratic
normal forms
\begin{equation}
\Delta F_k(a)
=
\mu_k^{(2{\rm w})}+a^\top H_k a+O(\|a\|^4),
\label{eq:DeltaF-lowrank}
\end{equation}
\begin{equation}
\mathfrak r_k(a)
=
\mu_k^{(1{\rm w})}+a^\top M_k a+O(\|a\|^4),
\label{eq:rres-lowrank}
\end{equation}
with \(H_k,M_k\succeq 0\) on the active split subspace.

\begin{assumption}[Low-rank comparability]
\label{ass:low-rank-comp}
On the active split subspace there exist \(c_k^-,c_k^+>0\) such that
\[
c_k^- M_k \preceq H_k \preceq c_k^+ M_k.
\]
\end{assumption}

\begin{proposition}[Low-rank wall-match comparability]
\label{prop:low-rank-wall-match}
Assume \cref{ass:wall-match,ass:low-rank-comp} and restrict to the onset manifold
\(\mu_k^{(1{\rm w})}=\mu_k^{(2{\rm w})}=0\).
Then for sufficiently small \(\mathfrak r_k\) there exists \(C_k>0\) such that
\begin{equation}
c_k^-\,\mathfrak r_k-C_k\mathfrak r_k^2
\le
\Delta F_k
\le
c_k^+\,\mathfrak r_k+C_k\mathfrak r_k^2.
\label{eq:low-rank-wall-match}
\end{equation}
\end{proposition}

\begin{proof}
On onset,
\[
\Delta F_k=a_k^\top H_k a_k+O(\|a_k\|^4),
\qquad
\mathfrak r_k=a_k^\top M_k a_k+O(\|a_k\|^4).
\]
The quadratic-form comparison gives
\[
c_k^- a_k^\top M_k a_k\le a_k^\top H_k a_k\le c_k^+ a_k^\top M_k a_k.
\]
Since \(\mathfrak r_k\) is itself quadratic in \(a_k\) to leading order, the quartic remainders are \(O(\mathfrak r_k^2)\),
which yields \eqref{eq:low-rank-wall-match}.
\end{proof}

\subsection{Gaussian realization of the odd mode and the leak stock law}
\label{subsec:B-gaussian}

The theorem path above does \emph{not} rely on a KMS-subtracted matrix.
In particular, we do not use
\[
\Sigma_{B,k}:=\bigl(C_{B,k}-C^{\rm KMS}_{\Lambda_k}\bigr)+m_{B,k}m_{B,k}^\top
\]
as a positive object, because it need not be positive semidefinite.
When a safe matrix observable is needed in the realization layer, the natural PSD quantity is
\begin{equation}
M_{B,k}:=C_{B,k}+m_{B,k}m_{B,k}^\top\succeq 0.
\label{eq:MB-def}
\end{equation}

\begin{assumption}[Gaussian completion realization]
\label{ass:gaussian-completion}
For each admissible cutoff \(\Lambda\) in the welded regime, the quadratic influence functional \(W_\Lambda[J]\) admits a
Gaussian dilation, unique up to symplectic equivalence, by a bosonic completion sector \(B_\Lambda\) with canonical
quadratures \(X_\Lambda\) and positive quadratic Hamiltonian
\[
H_{B,\Lambda}=\frac12 X_\Lambda^\top h_\Lambda X_\Lambda,
\qquad h_\Lambda\succeq 0.
\]
\end{assumption}

\begin{assumption}[Dominant odd-mode projection]
\label{ass:export-selection}
Inside the Gaussian completion of \cref{ass:gaussian-completion}, there exists a smooth unit vector
\[
u_k:=u_-(\Lambda_k)
\]
spanning the active odd line and separated from the orthogonal sector by a local spectral gap.
We write
\[
P_k^\perp:=I-u_k u_k^\top
\]
for the orthogonal projector.
\end{assumption}

\begin{assumption}[Classical-control gate]
\label{ass:classical-gate}
There exists a scalar gate \(g_{\rm cl}(k)\in[0,1]\), determined by the accessible classicality defect
\(\mathrm{cl}_k\) or \(\mathrm{Off}_k\), such that \(g_{\rm cl}(k)\ll 1\) at a fresh-aeon start and
\(g_{\rm cl}(k)\to 1\) late in the aeon.
\end{assumption}

\begin{proposition}[Projected Gaussian affine dynamics]
\label{prop:projected-affine}
Under \cref{ass:gaussian-completion,ass:export-selection}, one may choose a finite effective completion-mode basis in which
the completion quadratures obey a linear-Gaussian recursion
\begin{equation}
X_{k+1}=A_k X_k+B_k j_k^{\rm ctl}+\nu_k,
\label{eq:X-recursion}
\end{equation}
with
\[
\mathbb E[\nu_k\mid\mathcal F_k]=0,
\qquad
\mathbb E[\nu_k\nu_k^\top\mid\mathcal F_k]=Q_k\succeq 0.
\]
Consequently the PSD second moment
\[
M_{B,k}=C_{B,k}+m_{B,k}m_{B,k}^\top
\]
obeys
\begin{equation}
M_{B,k+1}=A_k M_{B,k}A_k^\top+Q_k+\text{drift terms}.
\label{eq:MB-recursion}
\end{equation}
Projecting onto the odd line via
\[
a_k:=u_k^\top X_k
\]
gives the scalar affine law
\begin{equation}
a_{k+1}
=
\rho_k a_k+\beta_k+\xi_k+\varepsilon_k^{\rm mix},
\label{eq:B-recursion-general}
\end{equation}
where
\[
\rho_k:=u_{k+1}^\top A_k u_k,\qquad
\beta_k:=u_{k+1}^\top B_k j_k^{\rm ctl},
\]
\[
\xi_k:=u_{k+1}^\top \nu_k,\qquad
\varepsilon_k^{\rm mix}:=u_{k+1}^\top A_k P_k^\perp X_k.
\]
Moreover,
\[
\mathbb E[\xi_k\mid\mathcal F_k]=0,
\qquad
\mathbb E[\xi_k^2\mid\mathcal F_k]=u_{k+1}^\top Q_k u_{k+1}\ge 0.
\]
If \cref{ass:classical-gate} also holds and the welded shape family is fixed, then
\begin{equation}
\mathbb E[\xi_k^2\mid\mathcal F_k]
=
\eta_k\,g_{\rm cl}(k)\,\Delta I_k+O(\Delta I_k^2,\delta\Lambda_k^2),
\label{eq:Zk-linear}
\end{equation}
with \(\eta_k\ge 0\).
\end{proposition}

\begin{proof}
The quadratic SK influence functional determines a Gaussian CPTP map on the finite active completion-mode sector, hence a
linear-Gaussian state-space realization of the form \eqref{eq:X-recursion}; \eqref{eq:MB-recursion} is the standard update
for the PSD second moment.
Decompose \(X_k=u_k a_k+P_k^\perp X_k\), apply \(\ip{u_{k+1}}{\cdot}\) to \eqref{eq:X-recursion}, and collect the projected
transport, deterministic drive, noise, and orthogonal-mode contamination to obtain \eqref{eq:B-recursion-general}.
The mean-zero and positive-variance statements for \(\xi_k\) follow directly from the projected Gaussian noise covariance
\(
u_{k+1}^\top Q_k u_{k+1}\ge 0
\).
Finally, at fixed welded shape family one has \(A_{{\rm eff},k}^2\propto \Delta I_k\), so the fresh covariance injection
into the odd mode is linear in the current payload at leading order; the gate \(g_{\rm cl}(k)\) suppresses this injection
near a fresh-aeon start and yields \eqref{eq:Zk-linear}.
\end{proof}

\begin{assumption}[Centered odd preparation]
\label{ass:centered-odd}
Near split onset, the active odd mode is prepared either branch-symmetrically or in a squeeze-dominated regime, so the
projected deterministic displacement vanishes and orthogonal-mode contamination is higher order:
\[
\beta_k=0,
\qquad
\varepsilon_k^{\rm mix}=O_2.
\]
\end{assumption}

\begin{corollary}[Centered odd-mode law]
\label{cor:centered-odd-mode}
Under
\cref{ass:gaussian-completion,ass:export-selection,ass:classical-gate,ass:centered-odd},
the projected affine law \eqref{eq:B-recursion-general} reduces to
\begin{equation}
a_{k+1}=\rho_k a_k+\xi_k+O_2,
\label{eq:B-recursion}
\end{equation}
with \(\mathbb E[\xi_k\mid\mathcal F_k]=0\) and variance scaling \eqref{eq:Zk-linear}.
\end{corollary}

\paragraph{Stochastic/Gaussian realization layer.}
At theorem level the positive stock is simply
\[
S_k:=a_k^2.
\]
The general projected Gaussian dynamics need not be centered.
Under \cref{ass:centered-odd}, however, the physically relevant late-time regime is the one in which the odd mode carries no
leading deterministic displacement and the fresh effect enters through its variance.
In that centered specialization it is natural to use the conditional second moment
\[
S_k:=\mathbb E[a_k^2\mid\mathcal F_k].
\]

\begin{proposition}[Odd-mode stock law]
\label{prop:odd-mode-stock}
Under
\cref{ass:gaussian-completion,ass:export-selection,ass:classical-gate,ass:centered-odd},
the conditional stock
\[
S_k:=\mathbb E[a_k^2\mid\mathcal F_k]
\]
satisfies
\begin{equation}
S_{k+1}
=
t_k\,S_k+\eta_k\,g_{\rm cl}(k)\,\Delta I_k
+O(S_k\Delta I_k,\Delta I_k^2,\delta\Lambda_k^2),
\qquad t_k=\rho_k^2.
\label{eq:Sk-update}
\end{equation}
\end{proposition}

\begin{proof}
Square \eqref{eq:B-recursion}:
\[
a_{k+1}^2=\rho_k^2 a_k^2+2\rho_k a_k\xi_k+\xi_k^2+O_2.
\]
Taking conditional expectation and using \(\mathbb E[\xi_k\mid\mathcal F_k]=0\) removes the mixed term.
Substituting \eqref{eq:Zk-linear} gives \eqref{eq:Sk-update}.
\end{proof}

\begin{assumption}[UV comparability of energetic stiffness]
\label{ass:uv-match}
There exist constants \(c_-,c_+>0\) and a monotone increasing energetic stiffness \(\widetilde\alpha(\Lambda)\) of the
active odd mode such that, for all sufficiently large \(\Lambda\),
\begin{equation}
c_-\,\alpha(\Lambda)\le \widetilde\alpha(\Lambda)\le c_+\,\alpha(\Lambda).
\label{eq:alpha-comparable}
\end{equation}
\end{assumption}

This is the precise place where the note passes from what is derived in UOM to a physical hypothesis.
The coefficient \(\alpha(\Lambda)\) is derived from the continued OS problem; the statement that the energetic stiffness of
the split mode is comparable to it is not yet a theorem.

With the stock \(S_k\) as the actual state variable, the positive split-driving export flux is modeled by
\begin{equation}
Q_{{\rm leak},k}^{\rm model}
=
\widetilde\alpha'(\Lambda_k)\,\delta\Lambda_k\,S_k
\;+\;
\widetilde\alpha(\Lambda_k)\,\eta_k\,g_{\rm cl}(k)\,\Delta I_k
\;+\; O_2,
\label{eq:Qleak-effective}
\end{equation}
where \(\delta\Lambda_k:=\Lambda_{k+1}-\Lambda_k\).
The first term is the transport cost of carrying already-prepared odd stock to larger cutoff; the second is the fresh
injection due to the current payload.
Under \cref{ass:uv-match}, this gives the scaling comparison
\begin{equation}
Q_{{\rm leak},k}^{\rm model}
\asymp
\alpha'(\Lambda_k)\,\delta\Lambda_k\,S_k
\;+\;
\alpha(\Lambda_k)\,\eta_k\,g_{\rm cl}(k)\,\Delta I_k.
\label{eq:Qleak-alpha}
\end{equation}

\paragraph{Current RP proxy reinterpreted.}
The quantity
\[
Q_{\rm rg}=\chi J_{\rm RP}/\beta_{\rm dS}
\]
therefore remains useful only as a phenomenological probe. It may correlate with \(a_k^2\) or with the stock \(S_k\), but
it is not the fundamental leak law because the healthy receiver itself has no quadratic RP block on the welded family.

\subsection{Approximate split condition from the odd-stock leak law}
\label{subsec:split-threshold}

Let \(R_k^{\rm abs}:=|R_k|\). In the micro-step regime the geometry-side surrogate gives
\[
\Delta E_{{\rm geom},k}
\approx
\frac{\zeta\,(A+c_\Lambda \Lambda_k^2)}{2\sqrt{12\,R_k^{\rm abs}}}\;\Theta_k
\;-\;
\frac{2\zeta\,c_\Lambda \Lambda_k\sqrt{R_k^{\rm abs}}}{\sqrt{12}}\;\delta\Lambda,
\]
while
\[
J_{{\rm wall},k}\approx \mu_{\rm RG}(\Lambda_k)\,\delta\Lambda,
\qquad
\mu_{\rm RG}(\Lambda)\simeq \mu_0+\mu_2\Lambda^2>0.
\]
Hence the local one-wall threshold remains
\begin{equation}
Q_{{\rm leak},k}
\gtrsim
\frac{\zeta\,(A+c_\Lambda \Lambda_k^2)}{2\sqrt{12\,R_k^{\rm abs}}}\;\Theta_k.
\label{eq:split-threshold-base}
\end{equation}
Since
\[
\Theta_k
=
K(\Lambda_k)\,\Delta I_k
\sim
K_0\,\Lambda_k^{-p}\,e^{-2d_{\rm sep}(\Lambda_k)\Lambda_k}\,\Delta I_k,
\qquad p\in\{10,14\},
\]
substituting \eqref{eq:Qleak-effective} gives
\begin{equation}
\widetilde\alpha'(\Lambda_k)\,\delta\Lambda_k\,S_k
\;+\;
\widetilde\alpha(\Lambda_k)\,\eta_k\,g_{\rm cl}(k)\,\Delta I_k
\gtrsim
\frac{\zeta\,(A+c_\Lambda \Lambda_k^2)}{2\sqrt{12\,R_k^{\rm abs}}}\,
K(\Lambda_k)\,\Delta I_k.
\label{eq:split-threshold-effective}
\end{equation}
This is the corrected split threshold. It shows the two late-time branches transparently:
\begin{enumerate}[leftmargin=2em]
\item if \(S_k\) remains negligible, only the fresh term remains, and the exhaustion branch can survive;
\item if \(S_k\) has accumulated, the transport term can dominate even when \(\Delta I_k\) is already tiny.
\end{enumerate}
Under \cref{ass:uv-match} and \(\alpha(\Lambda)\sim \kappa_d\Lambda^{d+2}H^{-d}\), the transport term scales like
\[
\widetilde\alpha'(\Lambda_k)\,\delta\Lambda_k\,S_k
\asymp
\Lambda_k^{d+1}\,\delta\Lambda_k\,S_k,
\]
which is exactly the mechanism by which late split can beat the shrinking fresh payload.

\subsection{Nondegeneracy clause for the normalized seam coefficient}
\label{subsec:separate-c-nondegeneracy}

The split threshold \eqref{eq:split-threshold-effective} is a statement about when the late-time
odd stock can overcome the shrinking fresh payload.
It is \emph{not} by itself a theorem that the normalized seam coefficient \(c_{\Lambda_k}\) on the
canonical compressed welded LoG line is nonzero.
Any theorem proving \(c_{\Lambda_k}\neq 0\) must therefore come from a separate overlap or transport
hypothesis, not from the threshold inequality alone.
Downstream, this clause is absorbed into the delayed selected-line package rather than carried as a
separate package summand.

\begin{definition}[Unit-source seam overlap scalar and no-destructive-cancellation margin]
\label{def:unit-source-overlap-margin}
Fix an accepted tick \(k\).
Assume the canonical compressed welded LoG line is defined and let
\[
\tau^{\mathrm{unit}}_{\Lambda_k}
\]
be a detector-factorized unit welded odd source.
Define the raw unit-source seam overlap scalar by
\begin{equation}
\widetilde c_{\Lambda_k}
:=
\big\langle\!\big\langle
\Pi_{\mathrm{phys},\Lambda_k}\mathsf T_{\Lambda_k}\!\big[\tau^{\mathrm{unit}}_{\Lambda_k}\big],\,
\delta C^{\mathrm{LoG}}_{\Lambda_k}
\big\rangle\!\big\rangle_{C_{\ast,\Lambda_k}}.
\label{eq:unit-source-overlap-scalar}
\end{equation}
A \emph{no-destructive-cancellation margin} at \(\Lambda_k\) is any real number
\[
\mu_{\Lambda_k}>0
\]
such that
\begin{equation}
\abs{\widetilde c_{\Lambda_k}}\ge \mu_{\Lambda_k}.
\label{eq:no-destr-cancel-margin}
\end{equation}
On a connected accepted-branch segment \(I\), a \emph{branchwise no-destructive-cancellation
margin} is a constant \(\mu_I>0\) such that
\begin{equation}
\abs{\widetilde c_{\Lambda}}\ge \mu_I
\qquad
\forall\,\Lambda\in I.
\label{eq:no-destr-cancel-margin-branch}
\end{equation}
\end{definition}

\begin{theorem}[Nondegeneracy criterion for \texorpdfstring{\(c_{\Lambda_k}\)}{c_Lambda_k}]
\label{thm:separate-c-nondegeneracy}
Fix an accepted tick \(k\).
Assume:
\begin{enumerate}[leftmargin=2em]
\item a detector-factorized unit welded odd source at \(\Lambda_k\);
\item the unit-source seam-faithfulness package, namely:
  \begin{enumerate}[label=(\alph*),leftmargin=2em]
  \item the raw seam pairing factorizes as
    \begin{equation}
    \widetilde s_{\Lambda_k}(f,y)
    =
    \sigma_{\mathrm{fill}}(f)\,
    \widetilde c_{\Lambda_k}\,
    \overline{\mathrm{obs}}(y)
    \qquad
    \forall f,\forall y;
    \label{eq:unit-source-seam-factorization}
    \end{equation}
  \item whenever
    \(
    \widetilde c_{\Lambda_k}\neq 0
    \)
    and
    \(
    \overline{\mathrm{obs}}(y)\neq 0
    \),
    the compressed welded LoG line is transversely physical, dominant-line realization holds, and
    the normalized seam coefficient
    \begin{equation}
    c_{\Lambda_k}
    :=
    \kappa_{\Lambda_k}^{-1/2}\,\widetilde c_{\Lambda_k}
    \label{eq:unit-source-normalized-c}
    \end{equation}
    is defined;
  \end{enumerate}
\item a no-destructive-cancellation margin at \(\Lambda_k\) in the sense of
  \cref{def:unit-source-overlap-margin};
\item a witness \(y\) with
  \[
  \overline{\mathrm{obs}}(y)\neq 0.
  \]
\end{enumerate}
Then:
\begin{enumerate}[leftmargin=2em]
\item the raw seam pairing is nonzero for every filler tag \(f\):
  \[
  \widetilde s_{\Lambda_k}(f,y)\neq 0;
  \]
\item the canonical compressed welded LoG line is transversely physical at \(\Lambda_k\);
\item dominant-line realization holds at \(\Lambda_k\);
\item the normalized seam coefficient is defined and nonzero:
  \[
  c_{\Lambda_k}\neq 0.
  \]
\end{enumerate}
If, in addition, for some filler tag \(f\) the source \(\tau^{f,y}_{\Lambda_k}\) is representable on
the characterization layer and admits a source-compatible descendant realization, then the
corresponding descendant representative satisfies
\[
\rho^{\mathrm{dom}}_{\Lambda_k}\!\big(J^{f,y}_{\Lambda_k}\big)\neq 0.
\]
\end{theorem}

\begin{proof}
The margin \eqref{eq:no-destr-cancel-margin} implies
\(
\widetilde c_{\Lambda_k}\neq 0
\).
Together with
\(
\overline{\mathrm{obs}}(y)\neq 0
\),
the factorization \eqref{eq:unit-source-seam-factorization} gives
\(
\widetilde s_{\Lambda_k}(f,y)\neq 0
\)
for every filler tag \(f\), proving item (1).
Item (2), item (3), and the definition and nonvanishing of
\(
c_{\Lambda_k}
\)
are exactly the remaining conclusions of the assumed unit-source seam-faithfulness package under the
same nondegeneracy hypothesis.
The final sentence is the corresponding witness-to-descendant comparison step under representability
and source-compatible descendant realization.
\end{proof}

\begin{corollary}[Branchwise sign-stable nondegeneracy]
\label{cor:branchwise-c-nondegeneracy}
Let \(I\) be a connected accepted-branch segment.
Assume:
\begin{enumerate}[leftmargin=2em]
\item the branchwise normalized seam coefficient
  \[
  \Lambda\mapsto c_{\Lambda}
  \]
  is continuous on \(I\);
\item a branchwise no-destructive-cancellation margin on \(I\) in the sense of
  \cref{def:unit-source-overlap-margin}.
\end{enumerate}
Then
\[
c_{\Lambda}\neq 0
\qquad
\forall\,\Lambda\in I.
\]
In particular, the sign of \(c_{\Lambda}\) is constant on \(I\).
If there also exist a fixed LS witness \(y\) with
\(
\overline{\mathrm{obs}}(y)\neq 0
\)
and a branchwise source-compatible descendant realization satisfying
\[
\rho^{\mathrm{dom}}_{\Lambda}\!\big(J^{f,y}_{\Lambda}\big)
=
\sigma_{\mathrm{fill}}(f)\,
c_{\Lambda}\,
\overline{\mathrm{obs}}(y),
\]
then the sign of
\(
\rho^{\mathrm{dom}}_{\Lambda}(J^{f,y}_{\Lambda})
\)
is constant on \(I\) as well.
\end{corollary}

\begin{proof}
By \eqref{eq:no-destr-cancel-margin-branch}, the scalar
\(
\widetilde c_{\Lambda}
\)
never vanishes on \(I\).
Hence neither does
\(
c_{\Lambda}
\),
because the normalization only rescales the raw overlap scalar by a positive branchwise factor.
A continuous nonzero real-valued function on a connected segment cannot change sign, proving the
constancy of \(\operatorname{sgn}(c_{\Lambda})\).
The final statement follows immediately from the displayed branchwise comparison identity and the
fixed nonzero scalar \(\overline{\mathrm{obs}}(y)\).
\end{proof}

\begin{remark}[Logical separation from the split threshold]
\label{rem:separate-c-nondegeneracy}
The hypotheses of \cref{thm:separate-c-nondegeneracy,cor:branchwise-c-nondegeneracy} are
intentionally separate from \eqref{eq:split-threshold-effective}.
The threshold governs the emergence of split-side stock; the no-destructive-cancellation margin
governs overlap of the unit welded odd source with the canonical compressed dominant line.
Accordingly, the present note treats \(c_{\Lambda}\neq 0\) as the transport/overlap clause later
imported into the delayed selected-line package rather than as a consequence of the split
threshold.
\end{remark}
%
  }%
}

\def\UOMTransportFragmentLoaded{}

\IfFileExists{compressed_log_spectral_selection_note.tex}{%
  \ifdefined\UOMTransportFragmentLoaded\else
\documentclass[11pt]{article}
\usepackage[margin=1in]{geometry}
\usepackage{amsmath,amssymb,amsthm,mathtools}
\usepackage{bm}
\usepackage{microtype}
\usepackage{enumitem}
\usepackage{booktabs}
\usepackage{array}
\usepackage{xcolor}
\usepackage{xr-hyper}
\usepackage{hyperref}
\usepackage[nameinlink,noabbrev]{cleveref}

\providecommand{\Fsloc}{\mathcal F_{\rm s.loc}}
\providecommand{\Qt}{Q_t}
\externaldocument[main-]{UOM/main}
\externaldocument[uomdesc-]{UOM/uom_descendant_note}
\externaldocument{UOM/completion_sector_split_model_note}

\hypersetup{
  colorlinks=true,
  linkcolor=blue!60!black,
  citecolor=blue!60!black,
  urlcolor=blue!60!black
}

\theoremstyle{definition}
\newtheorem{definition}{Definition}[section]
\newtheorem{assumption}{Assumption}[section]
\theoremstyle{plain}
\newtheorem{theorem}[definition]{Theorem}
\newtheorem{lemma}[definition]{Lemma}
\newtheorem{proposition}[definition]{Proposition}
\newtheorem{corollary}[definition]{Corollary}
\theoremstyle{remark}
\newtheorem{remark}[definition]{Remark}
\crefname{definition}{definition}{definitions}
\Crefname{definition}{Definition}{Definitions}
\crefname{assumption}{assumption}{assumptions}
\Crefname{assumption}{Assumption}{Assumptions}
\crefname{theorem}{theorem}{theorems}
\Crefname{theorem}{Theorem}{Theorems}
\crefname{lemma}{lemma}{lemmas}
\Crefname{lemma}{Lemma}{Lemmas}
\crefname{proposition}{proposition}{propositions}
\Crefname{proposition}{Proposition}{Propositions}
\crefname{corollary}{corollary}{corollaries}
\Crefname{corollary}{Corollary}{Corollaries}
\crefname{remark}{remark}{remarks}
\Crefname{remark}{Remark}{Remarks}

\newcommand{\cH}{\mathcal H}
\newcommand{\cA}{\mathcal A}
\newcommand{\cD}{\mathcal D}
\newcommand{\cB}{\mathcal B}
\newcommand{\cC}{\mathcal C}
\newcommand{\cL}{\mathcal L}
\newcommand{\cK}{\mathcal K}
\newcommand{\cR}{\mathcal R}
\newcommand{\cG}{\mathcal G}
\newcommand{\Ztwo}{\mathbb Z_2}
\newcommand{\Sys}{\mathrm{Sys}}
\newcommand{\Time}{\mathrm{Time}}
\newcommand{\RG}{\mathrm{RG}}
\newcommand{\Tr}{\mathrm{Tr}}
\newcommand{\diag}{\mathrm{diag}}
\newcommand{\ad}{\mathrm{ad}}
\newcommand{\id}{\mathrm{id}}
\newcommand{\norm}[1]{\left\lVert #1\right\rVert}
\newcommand{\abs}[1]{\left\lvert #1\right\rvert}
\newcommand{\ip}[2]{\left\langle #1,#2\right\rangle}

\externaldocument{UOM/odd_selector_note}
\externaldocument{UOM/canonical_odd_interface_note}
\externaldocument{UOM/active_response_transport_program_note}
\title{Canonical Odd Spectral Selection on the Compressed Welded LoG Channel}
\author{}
\date{}
\begin{document}
\maketitle
\fi

\section{Canonical Odd Spectral Selection on the Compressed Welded LoG Channel}
\label{sec:compressed-log-spectral-selection}

Fix an accepted tick \(k\) on a welded admissible branch. This section isolates the compressed
welded LoG channel and proves that the continued odd quadratic form canonically selects a negative
rank-one odd line on that channel. The transport of the descendant class into that line is recorded
separately in \cref{thm:char-weak-interface}; here it is used only in
\cref{cor:compressed-transport-hit}.

The macro source-side temporal envelope is DoG, the macro source-side spatial envelope is LoG, the
micro receiver temporal envelope is DoG, and the micro receiver spatial envelope is Gaussian. By
Corollary~\ref{main-cor:welded-log}, the welded envelope is LoG in both time and space, with scales
adding in quadrature.

\begin{definition}[Canonical welded LoG profile]
\label{def:compressed-log-profile}
Let \(t_k\) be the accepted tick time and let \(\Omega_\ast\in S^3\) denote the common spatial center
selected by the co-centering/history-lock mechanism. Let
\[
s_{t,k}:=\sqrt{\Sigma_k^2+\sigma_{t,k}^2},
\qquad
s_{x,k}:=\sqrt{\rho_k^2+\sigma_{x,k}^2}
\]
be the effective welded scales. Define the normalized welded LoG pulse by
\begin{equation}
u^{\mathrm{LoG}}_{\Lambda_k}(t,\Omega)
:=
\mathcal N_k\,
\frac{d^2}{dt^2}g_{s_{t,k}}(t-t_k)\,
\Delta_{S^3}g_{s_{x,k}}(\Omega,\Omega_\ast),
\label{eq:ulog-def}
\end{equation}
where \(\mathcal N_k>0\) is chosen so that
\(
\|u^{\mathrm{LoG}}_{\Lambda_k}\|_{L^2(\mathbb R\times S^3)}=1
\).
We write
\begin{equation}
\mathcal H^{\mathrm{LoG}}_{\Lambda_k}
:=
\mathbb R\,u^{\mathrm{LoG}}_{\Lambda_k}
\label{eq:Hlog-def}
\end{equation}
for the compressed welded LoG channel.
\end{definition}

\begin{proposition}[Exact welded LoG form]
\label{prop:compressed-log-exact}
The normalized profile \(u^{\mathrm{LoG}}_{\Lambda_k}\) of
\cref{def:compressed-log-profile} is exactly the welded macro/micro profile produced by:
\begin{enumerate}[leftmargin=2em]
\item macro temporal DoG convolved with micro temporal DoG;
\item macro spatial LoG convolved with micro spatial Gaussian.
\end{enumerate}
In particular the welded profile is LoG in time and LoG in space, with scales adding in quadrature.
\end{proposition}

\begin{proof}
The temporal and spatial welded identities are exactly the two statements of
Corollary~\ref{main-cor:welded-log}. After multiplying the factors and normalizing in \(L^2\), one obtains
\eqref{eq:ulog-def}.
\end{proof}

\begin{proposition}[Band-edge locking and unique center]
\label{prop:compressed-log-lock}
Assume the accepted tick lies on an admissible band-locked welded branch in the sense of
Definition~\ref{main-def:eta-bridge-leak}. Then:
\begin{enumerate}[leftmargin=2em]
\item the effective welded widths satisfy
\[
s_{t,k}=\frac{\alpha_t^\ast}{\Lambda_k},
\qquad
s_{x,k}=\frac{\alpha_x^\ast}{v\,\Lambda_k},
\]
with \(\alpha_t^\ast,\alpha_x^\ast>0\) independent of \(k\) along a fixed admissible branch;
\item the macro and micro temporal centers coincide at \(t_k\), and the macro and micro spatial
  centers coincide at \(\Omega_\ast\);
\item once the receiver trace sector carries a nonzero inherited LoG background, the absolute center
\(\Omega_\ast\) is uniquely locked for all subsequent accepted ticks.
\end{enumerate}
Consequently, after normalization, the admissible welded pulse space at fixed accepted tick is the
one-dimensional line \(\mathcal H^{\mathrm{LoG}}_{\Lambda_k}\).
\end{proposition}

\begin{proof}
The width lock is the band-locking clause of Definition~\ref{main-def:eta-bridge-leak}. The
temporal and spatial co-centering statements are \eqref{main-eq:Ct} and \eqref{main-eq:Cx}, and the
persistence of the spatial center is exactly Proposition~\ref{main-prop:history-center-lock}.

Therefore, at fixed accepted tick, the shape, widths, and center are fixed. The only remaining freedom
is an overall real amplitude. After \(L^2\)-normalization, that freedom reduces to the sign
\(\pm u^{\mathrm{LoG}}_{\Lambda_k}\). Hence the admissible welded channel is precisely the
one-dimensional real line \(\mathcal H^{\mathrm{LoG}}_{\Lambda_k}\).
\end{proof}

\begin{proposition}[Smooth dependence along a branch]
\label{prop:compressed-log-smooth}
Along any admissible branch on which \(\Lambda\mapsto \alpha_t^\ast,\alpha_x^\ast\) are constant and
the locked center \(\Omega_\ast\) is fixed, the map
\[
\Lambda\longmapsto u^{\mathrm{LoG}}_\Lambda
\]
is \(C^1\) in \(L^2(\mathbb R\times S^3)\). Consequently the rank-one orthogonal projector
\begin{equation}
P^{\mathrm{LoG}}_\Lambda
:=
u^{\mathrm{LoG}}_\Lambda\otimes u^{\mathrm{LoG}}_\Lambda
\label{eq:Plog-def}
\end{equation}
is \(C^1\) along the branch.
\end{proposition}

\begin{proof}
Along the branch one has
\(
s_t(\Lambda)=\alpha_t^\ast/\Lambda
\)
and
\(
s_x(\Lambda)=\alpha_x^\ast/(v\Lambda)
\),
so both scales depend \(C^\infty\)-smoothly on \(\Lambda\). The Gaussian heat kernel and its time and
spatial derivatives depend smoothly on the scale parameter, hence the unnormalized welded LoG profile
depends \(C^1\)-smoothly on \(\Lambda\) in \(L^2\). Its \(L^2\)-norm is strictly positive, so
normalization preserves \(C^1\)-dependence. Formula \eqref{eq:Plog-def} then gives a \(C^1\)
projector-valued map.
\end{proof}

\begin{proposition}[Well-defined restriction of the continued quadratic form]
\label{prop:compressed-form-well-defined}
Assume the welded linear-response/Gaussian radiative-sector package of
Paragraph~\ref{main-para:uplift-decomp} and Theorem~\ref{main-thm:tick-fidelity-RP}. Then the
analytically continued odd quadratic response form is well-defined on the compressed welded channel
\(\mathcal H^{\mathrm{LoG}}_{\Lambda_k}\).
\end{proposition}

\begin{proof}
By Paragraph~\ref{main-para:uplift-decomp}, the welded source enters the Gaussian radiative sector
through the linear response map \(\mathsf T_\Lambda\), and Theorem~\ref{main-thm:tick-fidelity-RP}
supplies the finite-cutoff bounded Gaussian response package on the welded envelope class. By
\cref{prop:compressed-log-lock}, the compressed welded channel is a one-dimensional subspace of that
admissible class. Restricting any continuous quadratic form to a one-dimensional subspace is
automatic. Hence the analytically continued odd quadratic response form restricts canonically to
\(\mathcal H^{\mathrm{LoG}}_{\Lambda_k}\).
\end{proof}

\begin{definition}[Compressed continued odd quadratic coefficient]
\label{def:compressed-odd-coeff}
Fix a witness window \(\mu_{ps}>0\) in the accepted regime. Let
\[
\mathfrak q^{\mathrm{cont}}_{\Lambda_k,\mu_{ps}}
:
\mathcal H^{\mathrm{LoG}}_{\Lambda_k}\times \mathcal H^{\mathrm{LoG}}_{\Lambda_k}
\to
\mathbb R
\]
denote the analytically continued odd quadratic response form restricted to the compressed welded LoG
channel. Since \(\mathcal H^{\mathrm{LoG}}_{\Lambda_k}\) is one-dimensional, there is a unique scalar
\(\beta_k\in\mathbb R\) such that
\begin{equation}
\mathfrak q^{\mathrm{cont}}_{\Lambda_k,\mu_{ps}}(X,Y)
=
\beta_k\,
\langle X,u^{\mathrm{LoG}}_{\Lambda_k}\rangle_{L^2}
\langle Y,u^{\mathrm{LoG}}_{\Lambda_k}\rangle_{L^2}
\qquad
\forall X,Y\in\mathcal H^{\mathrm{LoG}}_{\Lambda_k}.
\label{eq:qcont-beta}
\end{equation}
Equivalently,
\[
\beta_k
=
\mathfrak q^{\mathrm{cont}}_{\Lambda_k,\mu_{ps}}
\big(
u^{\mathrm{LoG}}_{\Lambda_k},
u^{\mathrm{LoG}}_{\Lambda_k}
\big).
\]
\end{definition}

\begin{proposition}[Rigorous control of the compressed coefficient]
\label{prop:compressed-beta-estimate}
Assume:
\begin{enumerate}[leftmargin=2em]
\item the smooth RP degradation law of Theorem~\ref{main-thm:smooth-RP};
\item the approximate intertwiner bound of Lemma~\ref{main-lem:approx-intertwiner-new};
\item the temporal DoG commutator bound of Lemma~\ref{main-lem:DoG-comm-new};
\item the exponentially decaying memory-tail package of
  Lemma~\ref{main-lem:app-exp-backflow} and Proposition~\ref{main-prop:app-tick-progress}.
\end{enumerate}
Then there exist constants \(C_{\rm int},C_{\rm DoG},C_{\rm mem},C_4>0\) and \(\eta>0\), depending only on the
fixed welded family and not on \(k\), such that
\begin{equation}
\beta_k
=
-\,\alpha(\Lambda_k)\,\mu_{ps}^2
\;+\;
r_k,
\label{eq:beta-splitting}
\end{equation}
with remainder bounded by
\begin{equation}
|r_k|
\le
C_{\rm int}\,\varepsilon_k^{\eta}
\;+\;
C_{\rm DoG}\,\Phi(\sigma_{t,k}\Omega_{\rm gap})
\;+\;
C_{\rm mem}\,e^{-\Delta_k/\tau_b}
\;+\;
C_4\,\mu_{ps}^4.
\label{eq:beta-remainder}
\end{equation}
\end{proposition}

\begin{proof}
By Theorem~\ref{main-thm:smooth-RP}, the ideal analytically continued Gaussian OS problem gives
\(
\lambda_{\min}(\mu_{ps},\Lambda_k)
=
-\alpha(\Lambda_k)\mu_{ps}^2+O(\mu_{ps}^4)
\)
on the welded Gaussian reference kernel. Restricting that quadratic form to the canonical normalized
welded LoG profile \(u^{\mathrm{LoG}}_{\Lambda_k}\) produces the leading term
\(-\alpha(\Lambda_k)\mu_{ps}^2\), with the \(O(\mu_{ps}^4)\) correction bounded by
\(C_4\mu_{ps}^4\).

The regulated intertwiner contribution is bounded by
Lemma~\ref{main-lem:approx-intertwiner-new}, the temporal DoG contribution is bounded by
Lemma~\ref{main-lem:DoG-comm-new}, and the intrinsic KMS bath contributes only the
Born-Nakajima-Zwanzig memory tail controlled by Lemma~\ref{main-lem:app-exp-backflow} and
Proposition~\ref{main-prop:app-tick-progress}. Since \(u^{\mathrm{LoG}}_{\Lambda_k}\) is
normalized, each of these operator bounds transfers directly to the scalar quadratic coefficient.
Summing the four contributions gives \eqref{eq:beta-splitting}--\eqref{eq:beta-remainder}.
\end{proof}

\begin{definition}[Certified compressed spectral regime]
\label{def:compressed-spectral-regime}
We say that an accepted tick \(k\) lies in the \emph{certified compressed spectral regime} if
\begin{equation}
\alpha(\Lambda_k)\,\mu_{ps}^2
>
2\Big(
C_{\rm int}\,\varepsilon_k^{\eta}
\;+\;
C_{\rm DoG}\,\Phi(\sigma_{t,k}\Omega_{\rm gap})
\;+\;
C_{\rm mem}\,e^{-\Delta_k/\tau_b}
\;+\;
C_4\,\mu_{ps}^4
\Big).
\label{eq:compressed-certified}
\end{equation}
\end{definition}

\begin{theorem}[Finite-cutoff canonical odd spectral selection on the compressed welded LoG channel]
\label{thm:finite-cutoff-canonical-odd-spectral-selection}
Assume
\cref{prop:compressed-log-exact,prop:compressed-log-lock,prop:compressed-log-smooth,prop:compressed-form-well-defined,prop:compressed-beta-estimate}.
Let \(k\) be an accepted tick in the certified compressed spectral regime of
\cref{def:compressed-spectral-regime}. Then:
\begin{enumerate}[leftmargin=2em]
\item the compressed welded channel \(\mathcal H^{\mathrm{LoG}}_{\Lambda_k}\) is a canonical
one-dimensional odd line;
\item the compressed continued odd quadratic form is represented by the rank-one self-adjoint operator
\begin{equation}
K^{\mathrm{LoG}}_{\Lambda_k}
:=
\beta_k\,P^{\mathrm{LoG}}_{\Lambda_k},
\label{eq:Klog-def}
\end{equation}
with \(\beta_k<0\);
\item the unique negative spectral projector of \(K^{\mathrm{LoG}}_{\Lambda_k}\) is
\begin{equation}
P_{-,\Lambda_k}^{\mathrm{LoG}}
=
P^{\mathrm{LoG}}_{\Lambda_k}
=
u^{\mathrm{LoG}}_{\Lambda_k}\otimes u^{\mathrm{LoG}}_{\Lambda_k};
\label{eq:Pminus-log}
\end{equation}
\item along any admissible branch, the projector \(P_{-,\Lambda}^{\mathrm{LoG}}\) is \(C^1\), hence
the compressed active odd line is canonically and smoothly transported.
\end{enumerate}
\end{theorem}

\begin{proof}
By \cref{prop:compressed-log-lock}, the admissible welded channel at fixed accepted tick is exactly the
one-dimensional line \(\mathcal H^{\mathrm{LoG}}_{\Lambda_k}\). By
\cref{def:compressed-odd-coeff}, the compressed continued odd quadratic form on that line is represented
by the scalar \(\beta_k\), equivalently by the self-adjoint rank-one operator
\eqref{eq:Klog-def}.

By \cref{prop:compressed-beta-estimate},
\[
\beta_k
\le
-\,\alpha(\Lambda_k)\mu_{ps}^2
\;+\;
\Big(
C_{\rm int}\,\varepsilon_k^{\eta}
\;+\;
C_{\rm DoG}\,\Phi(\sigma_{t,k}\Omega_{\rm gap})
\;+\;
C_{\rm mem}\,e^{-\Delta_k/\tau_b}
\;+\;
C_4\,\mu_{ps}^4
\Big).
\]
Under the certified regime inequality \eqref{eq:compressed-certified}, the bracketed term is strictly
smaller than \(\tfrac12\alpha(\Lambda_k)\mu_{ps}^2\). Therefore
\[
\beta_k
<
-\frac12\,\alpha(\Lambda_k)\mu_{ps}^2
<
0.
\]
Thus \(K^{\mathrm{LoG}}_{\Lambda_k}\) has exactly one negative eigenvalue, namely \(\beta_k\), and
its unique negative spectral subspace is the entire line
\(\mathcal H^{\mathrm{LoG}}_{\Lambda_k}\). Since
\(
P^{\mathrm{LoG}}_{\Lambda_k}
=
u^{\mathrm{LoG}}_{\Lambda_k}\otimes u^{\mathrm{LoG}}_{\Lambda_k}
\),
the negative spectral projector is exactly \eqref{eq:Pminus-log}.

Finally, \cref{prop:compressed-log-smooth} gives \(C^1\)-dependence of
\(\Lambda\mapsto u^{\mathrm{LoG}}_\Lambda\), hence also of
\(\Lambda\mapsto P_{-,\Lambda}^{\mathrm{LoG}}\). This proves the canonical branchwise transport of the
compressed active odd line.
\end{proof}

\begin{corollary}[Canonical compressed odd projector]
\label{cor:compressed-export-selection}
For each accepted tick in the certified compressed spectral regime, the compressed, co-centered,
band-locked welded LoG channel carries the canonical rank-one projector
\[
P_{-,\Lambda_k}^{\mathrm{LoG}}
\]
onto its negative odd line. Equivalently,
\[
\mathcal H_{\Lambda_k}^{\mathrm{LoG}}
=
\operatorname{ran} P_{-,\Lambda_k}^{\mathrm{LoG}}.
\]
\end{corollary}

\begin{proof}
This is exactly \cref{thm:finite-cutoff-canonical-odd-spectral-selection}: the theorem identifies the
negative spectral projector with the rank-one projector
\(
u^{\mathrm{LoG}}_{\Lambda_k}\otimes u^{\mathrm{LoG}}_{\Lambda_k}
\),
whose range is \(\mathcal H^{\mathrm{LoG}}_{\Lambda_k}\).
\end{proof}

\begin{corollary}[Completion-sector realization of the compressed projector]
\label{cor:compressed-export-selection-completion}
Assume \cref{ass:gaussian-completion}. Then, for each accepted tick in the certified compressed
spectral regime, any Gaussian completion realization identifies the compressed projector
\(
P_{-,\Lambda_k}^{\mathrm{LoG}}
\)
with a unique rank-one active odd line in the completion sector. Equivalently, there exists a unit
vector
\[
u_-(\Lambda_k)
\]
spanning that realized line such that
\[
P_{-,\Lambda_k}^{\mathrm{LoG}}
=
u_-(\Lambda_k)\otimes u_-(\Lambda_k)
\]
in the chosen realization.
\end{corollary}

\begin{proof}
Under \cref{ass:gaussian-completion}, the finite active completion-mode sector is defined only up to
symplectic equivalence. A one-dimensional real line is invariant under such equivalences as a subspace,
even though a basis vector on that line is not canonical before normalization. By
\cref{thm:finite-cutoff-canonical-odd-spectral-selection}, the compressed odd line is already
canonically fixed abstractly by the projector \(P_{-,\Lambda_k}^{\mathrm{LoG}}\). Choosing the unique
unit vector on that line in a given realization yields the stated formula, and changing the completion
realization changes only the coordinate representative, not the underlying rank-one line.
\end{proof}

\begin{remark}[Scope of the compressed spectral theorem]
\label{rem:compressed-spectral-scope}
Theorem \cref{thm:finite-cutoff-canonical-odd-spectral-selection} determines the odd line only on
the compressed welded LoG channel. It does not classify the full completion-sector odd subspace, and
the completion-side statement of \cref{cor:compressed-export-selection-completion} is only a
realization statement for this rank-one projector.
\end{remark}

\begin{corollary}[Nontrivial descendant hit of the compressed dominant line]
\label{cor:compressed-transport-hit}
Assume \cref{thm:char-weak-interface}. Then for each accepted tick \(k\), there exists a
representative
\[
J^{\mathrm{desc}}_{\Lambda_k}\in[\Xi^{(1)}_{\Lambda_k}]
\]
whose response has nonzero component on the canonically selected compressed dominant odd line.
Equivalently, the transported descendant sector maps nontrivially into the compressed line fixed by
\(
P_{-,\Lambda_k}^{\mathrm{LoG}}
\).
\end{corollary}

\begin{proof}
This is exactly item (1) of \cref{thm:char-weak-interface}.
\end{proof}

\ifdefined\UOMTransportFragmentLoaded\else
\end{document}
\fi
%
}{%
  \IfFileExists{UOM/compressed_log_spectral_selection_note.tex}{%
    \ifdefined\UOMTransportFragmentLoaded\else
\documentclass[11pt]{article}
\usepackage[margin=1in]{geometry}
\usepackage{amsmath,amssymb,amsthm,mathtools}
\usepackage{bm}
\usepackage{microtype}
\usepackage{enumitem}
\usepackage{booktabs}
\usepackage{array}
\usepackage{xcolor}
\usepackage{xr-hyper}
\usepackage{hyperref}
\usepackage[nameinlink,noabbrev]{cleveref}

\providecommand{\Fsloc}{\mathcal F_{\rm s.loc}}
\providecommand{\Qt}{Q_t}
\externaldocument[main-]{UOM/main}
\externaldocument[uomdesc-]{UOM/uom_descendant_note}
\externaldocument{UOM/completion_sector_split_model_note}

\hypersetup{
  colorlinks=true,
  linkcolor=blue!60!black,
  citecolor=blue!60!black,
  urlcolor=blue!60!black
}

\theoremstyle{definition}
\newtheorem{definition}{Definition}[section]
\newtheorem{assumption}{Assumption}[section]
\theoremstyle{plain}
\newtheorem{theorem}[definition]{Theorem}
\newtheorem{lemma}[definition]{Lemma}
\newtheorem{proposition}[definition]{Proposition}
\newtheorem{corollary}[definition]{Corollary}
\theoremstyle{remark}
\newtheorem{remark}[definition]{Remark}
\crefname{definition}{definition}{definitions}
\Crefname{definition}{Definition}{Definitions}
\crefname{assumption}{assumption}{assumptions}
\Crefname{assumption}{Assumption}{Assumptions}
\crefname{theorem}{theorem}{theorems}
\Crefname{theorem}{Theorem}{Theorems}
\crefname{lemma}{lemma}{lemmas}
\Crefname{lemma}{Lemma}{Lemmas}
\crefname{proposition}{proposition}{propositions}
\Crefname{proposition}{Proposition}{Propositions}
\crefname{corollary}{corollary}{corollaries}
\Crefname{corollary}{Corollary}{Corollaries}
\crefname{remark}{remark}{remarks}
\Crefname{remark}{Remark}{Remarks}

\newcommand{\cH}{\mathcal H}
\newcommand{\cA}{\mathcal A}
\newcommand{\cD}{\mathcal D}
\newcommand{\cB}{\mathcal B}
\newcommand{\cC}{\mathcal C}
\newcommand{\cL}{\mathcal L}
\newcommand{\cK}{\mathcal K}
\newcommand{\cR}{\mathcal R}
\newcommand{\cG}{\mathcal G}
\newcommand{\Ztwo}{\mathbb Z_2}
\newcommand{\Sys}{\mathrm{Sys}}
\newcommand{\Time}{\mathrm{Time}}
\newcommand{\RG}{\mathrm{RG}}
\newcommand{\Tr}{\mathrm{Tr}}
\newcommand{\diag}{\mathrm{diag}}
\newcommand{\ad}{\mathrm{ad}}
\newcommand{\id}{\mathrm{id}}
\newcommand{\norm}[1]{\left\lVert #1\right\rVert}
\newcommand{\abs}[1]{\left\lvert #1\right\rvert}
\newcommand{\ip}[2]{\left\langle #1,#2\right\rangle}

\externaldocument{UOM/odd_selector_note}
\externaldocument{UOM/canonical_odd_interface_note}
\externaldocument{UOM/active_response_transport_program_note}
\title{Canonical Odd Spectral Selection on the Compressed Welded LoG Channel}
\author{}
\date{}
\begin{document}
\maketitle
\fi

\section{Canonical Odd Spectral Selection on the Compressed Welded LoG Channel}
\label{sec:compressed-log-spectral-selection}

Fix an accepted tick \(k\) on a welded admissible branch. This section isolates the compressed
welded LoG channel and proves that the continued odd quadratic form canonically selects a negative
rank-one odd line on that channel. The transport of the descendant class into that line is recorded
separately in \cref{thm:char-weak-interface}; here it is used only in
\cref{cor:compressed-transport-hit}.

The macro source-side temporal envelope is DoG, the macro source-side spatial envelope is LoG, the
micro receiver temporal envelope is DoG, and the micro receiver spatial envelope is Gaussian. By
Corollary~\ref{main-cor:welded-log}, the welded envelope is LoG in both time and space, with scales
adding in quadrature.

\begin{definition}[Canonical welded LoG profile]
\label{def:compressed-log-profile}
Let \(t_k\) be the accepted tick time and let \(\Omega_\ast\in S^3\) denote the common spatial center
selected by the co-centering/history-lock mechanism. Let
\[
s_{t,k}:=\sqrt{\Sigma_k^2+\sigma_{t,k}^2},
\qquad
s_{x,k}:=\sqrt{\rho_k^2+\sigma_{x,k}^2}
\]
be the effective welded scales. Define the normalized welded LoG pulse by
\begin{equation}
u^{\mathrm{LoG}}_{\Lambda_k}(t,\Omega)
:=
\mathcal N_k\,
\frac{d^2}{dt^2}g_{s_{t,k}}(t-t_k)\,
\Delta_{S^3}g_{s_{x,k}}(\Omega,\Omega_\ast),
\label{eq:ulog-def}
\end{equation}
where \(\mathcal N_k>0\) is chosen so that
\(
\|u^{\mathrm{LoG}}_{\Lambda_k}\|_{L^2(\mathbb R\times S^3)}=1
\).
We write
\begin{equation}
\mathcal H^{\mathrm{LoG}}_{\Lambda_k}
:=
\mathbb R\,u^{\mathrm{LoG}}_{\Lambda_k}
\label{eq:Hlog-def}
\end{equation}
for the compressed welded LoG channel.
\end{definition}

\begin{proposition}[Exact welded LoG form]
\label{prop:compressed-log-exact}
The normalized profile \(u^{\mathrm{LoG}}_{\Lambda_k}\) of
\cref{def:compressed-log-profile} is exactly the welded macro/micro profile produced by:
\begin{enumerate}[leftmargin=2em]
\item macro temporal DoG convolved with micro temporal DoG;
\item macro spatial LoG convolved with micro spatial Gaussian.
\end{enumerate}
In particular the welded profile is LoG in time and LoG in space, with scales adding in quadrature.
\end{proposition}

\begin{proof}
The temporal and spatial welded identities are exactly the two statements of
Corollary~\ref{main-cor:welded-log}. After multiplying the factors and normalizing in \(L^2\), one obtains
\eqref{eq:ulog-def}.
\end{proof}

\begin{proposition}[Band-edge locking and unique center]
\label{prop:compressed-log-lock}
Assume the accepted tick lies on an admissible band-locked welded branch in the sense of
Definition~\ref{main-def:eta-bridge-leak}. Then:
\begin{enumerate}[leftmargin=2em]
\item the effective welded widths satisfy
\[
s_{t,k}=\frac{\alpha_t^\ast}{\Lambda_k},
\qquad
s_{x,k}=\frac{\alpha_x^\ast}{v\,\Lambda_k},
\]
with \(\alpha_t^\ast,\alpha_x^\ast>0\) independent of \(k\) along a fixed admissible branch;
\item the macro and micro temporal centers coincide at \(t_k\), and the macro and micro spatial
  centers coincide at \(\Omega_\ast\);
\item once the receiver trace sector carries a nonzero inherited LoG background, the absolute center
\(\Omega_\ast\) is uniquely locked for all subsequent accepted ticks.
\end{enumerate}
Consequently, after normalization, the admissible welded pulse space at fixed accepted tick is the
one-dimensional line \(\mathcal H^{\mathrm{LoG}}_{\Lambda_k}\).
\end{proposition}

\begin{proof}
The width lock is the band-locking clause of Definition~\ref{main-def:eta-bridge-leak}. The
temporal and spatial co-centering statements are \eqref{main-eq:Ct} and \eqref{main-eq:Cx}, and the
persistence of the spatial center is exactly Proposition~\ref{main-prop:history-center-lock}.

Therefore, at fixed accepted tick, the shape, widths, and center are fixed. The only remaining freedom
is an overall real amplitude. After \(L^2\)-normalization, that freedom reduces to the sign
\(\pm u^{\mathrm{LoG}}_{\Lambda_k}\). Hence the admissible welded channel is precisely the
one-dimensional real line \(\mathcal H^{\mathrm{LoG}}_{\Lambda_k}\).
\end{proof}

\begin{proposition}[Smooth dependence along a branch]
\label{prop:compressed-log-smooth}
Along any admissible branch on which \(\Lambda\mapsto \alpha_t^\ast,\alpha_x^\ast\) are constant and
the locked center \(\Omega_\ast\) is fixed, the map
\[
\Lambda\longmapsto u^{\mathrm{LoG}}_\Lambda
\]
is \(C^1\) in \(L^2(\mathbb R\times S^3)\). Consequently the rank-one orthogonal projector
\begin{equation}
P^{\mathrm{LoG}}_\Lambda
:=
u^{\mathrm{LoG}}_\Lambda\otimes u^{\mathrm{LoG}}_\Lambda
\label{eq:Plog-def}
\end{equation}
is \(C^1\) along the branch.
\end{proposition}

\begin{proof}
Along the branch one has
\(
s_t(\Lambda)=\alpha_t^\ast/\Lambda
\)
and
\(
s_x(\Lambda)=\alpha_x^\ast/(v\Lambda)
\),
so both scales depend \(C^\infty\)-smoothly on \(\Lambda\). The Gaussian heat kernel and its time and
spatial derivatives depend smoothly on the scale parameter, hence the unnormalized welded LoG profile
depends \(C^1\)-smoothly on \(\Lambda\) in \(L^2\). Its \(L^2\)-norm is strictly positive, so
normalization preserves \(C^1\)-dependence. Formula \eqref{eq:Plog-def} then gives a \(C^1\)
projector-valued map.
\end{proof}

\begin{proposition}[Well-defined restriction of the continued quadratic form]
\label{prop:compressed-form-well-defined}
Assume the welded linear-response/Gaussian radiative-sector package of
Paragraph~\ref{main-para:uplift-decomp} and Theorem~\ref{main-thm:tick-fidelity-RP}. Then the
analytically continued odd quadratic response form is well-defined on the compressed welded channel
\(\mathcal H^{\mathrm{LoG}}_{\Lambda_k}\).
\end{proposition}

\begin{proof}
By Paragraph~\ref{main-para:uplift-decomp}, the welded source enters the Gaussian radiative sector
through the linear response map \(\mathsf T_\Lambda\), and Theorem~\ref{main-thm:tick-fidelity-RP}
supplies the finite-cutoff bounded Gaussian response package on the welded envelope class. By
\cref{prop:compressed-log-lock}, the compressed welded channel is a one-dimensional subspace of that
admissible class. Restricting any continuous quadratic form to a one-dimensional subspace is
automatic. Hence the analytically continued odd quadratic response form restricts canonically to
\(\mathcal H^{\mathrm{LoG}}_{\Lambda_k}\).
\end{proof}

\begin{definition}[Compressed continued odd quadratic coefficient]
\label{def:compressed-odd-coeff}
Fix a witness window \(\mu_{ps}>0\) in the accepted regime. Let
\[
\mathfrak q^{\mathrm{cont}}_{\Lambda_k,\mu_{ps}}
:
\mathcal H^{\mathrm{LoG}}_{\Lambda_k}\times \mathcal H^{\mathrm{LoG}}_{\Lambda_k}
\to
\mathbb R
\]
denote the analytically continued odd quadratic response form restricted to the compressed welded LoG
channel. Since \(\mathcal H^{\mathrm{LoG}}_{\Lambda_k}\) is one-dimensional, there is a unique scalar
\(\beta_k\in\mathbb R\) such that
\begin{equation}
\mathfrak q^{\mathrm{cont}}_{\Lambda_k,\mu_{ps}}(X,Y)
=
\beta_k\,
\langle X,u^{\mathrm{LoG}}_{\Lambda_k}\rangle_{L^2}
\langle Y,u^{\mathrm{LoG}}_{\Lambda_k}\rangle_{L^2}
\qquad
\forall X,Y\in\mathcal H^{\mathrm{LoG}}_{\Lambda_k}.
\label{eq:qcont-beta}
\end{equation}
Equivalently,
\[
\beta_k
=
\mathfrak q^{\mathrm{cont}}_{\Lambda_k,\mu_{ps}}
\big(
u^{\mathrm{LoG}}_{\Lambda_k},
u^{\mathrm{LoG}}_{\Lambda_k}
\big).
\]
\end{definition}

\begin{proposition}[Rigorous control of the compressed coefficient]
\label{prop:compressed-beta-estimate}
Assume:
\begin{enumerate}[leftmargin=2em]
\item the smooth RP degradation law of Theorem~\ref{main-thm:smooth-RP};
\item the approximate intertwiner bound of Lemma~\ref{main-lem:approx-intertwiner-new};
\item the temporal DoG commutator bound of Lemma~\ref{main-lem:DoG-comm-new};
\item the exponentially decaying memory-tail package of
  Lemma~\ref{main-lem:app-exp-backflow} and Proposition~\ref{main-prop:app-tick-progress}.
\end{enumerate}
Then there exist constants \(C_{\rm int},C_{\rm DoG},C_{\rm mem},C_4>0\) and \(\eta>0\), depending only on the
fixed welded family and not on \(k\), such that
\begin{equation}
\beta_k
=
-\,\alpha(\Lambda_k)\,\mu_{ps}^2
\;+\;
r_k,
\label{eq:beta-splitting}
\end{equation}
with remainder bounded by
\begin{equation}
|r_k|
\le
C_{\rm int}\,\varepsilon_k^{\eta}
\;+\;
C_{\rm DoG}\,\Phi(\sigma_{t,k}\Omega_{\rm gap})
\;+\;
C_{\rm mem}\,e^{-\Delta_k/\tau_b}
\;+\;
C_4\,\mu_{ps}^4.
\label{eq:beta-remainder}
\end{equation}
\end{proposition}

\begin{proof}
By Theorem~\ref{main-thm:smooth-RP}, the ideal analytically continued Gaussian OS problem gives
\(
\lambda_{\min}(\mu_{ps},\Lambda_k)
=
-\alpha(\Lambda_k)\mu_{ps}^2+O(\mu_{ps}^4)
\)
on the welded Gaussian reference kernel. Restricting that quadratic form to the canonical normalized
welded LoG profile \(u^{\mathrm{LoG}}_{\Lambda_k}\) produces the leading term
\(-\alpha(\Lambda_k)\mu_{ps}^2\), with the \(O(\mu_{ps}^4)\) correction bounded by
\(C_4\mu_{ps}^4\).

The regulated intertwiner contribution is bounded by
Lemma~\ref{main-lem:approx-intertwiner-new}, the temporal DoG contribution is bounded by
Lemma~\ref{main-lem:DoG-comm-new}, and the intrinsic KMS bath contributes only the
Born-Nakajima-Zwanzig memory tail controlled by Lemma~\ref{main-lem:app-exp-backflow} and
Proposition~\ref{main-prop:app-tick-progress}. Since \(u^{\mathrm{LoG}}_{\Lambda_k}\) is
normalized, each of these operator bounds transfers directly to the scalar quadratic coefficient.
Summing the four contributions gives \eqref{eq:beta-splitting}--\eqref{eq:beta-remainder}.
\end{proof}

\begin{definition}[Certified compressed spectral regime]
\label{def:compressed-spectral-regime}
We say that an accepted tick \(k\) lies in the \emph{certified compressed spectral regime} if
\begin{equation}
\alpha(\Lambda_k)\,\mu_{ps}^2
>
2\Big(
C_{\rm int}\,\varepsilon_k^{\eta}
\;+\;
C_{\rm DoG}\,\Phi(\sigma_{t,k}\Omega_{\rm gap})
\;+\;
C_{\rm mem}\,e^{-\Delta_k/\tau_b}
\;+\;
C_4\,\mu_{ps}^4
\Big).
\label{eq:compressed-certified}
\end{equation}
\end{definition}

\begin{theorem}[Finite-cutoff canonical odd spectral selection on the compressed welded LoG channel]
\label{thm:finite-cutoff-canonical-odd-spectral-selection}
Assume
\cref{prop:compressed-log-exact,prop:compressed-log-lock,prop:compressed-log-smooth,prop:compressed-form-well-defined,prop:compressed-beta-estimate}.
Let \(k\) be an accepted tick in the certified compressed spectral regime of
\cref{def:compressed-spectral-regime}. Then:
\begin{enumerate}[leftmargin=2em]
\item the compressed welded channel \(\mathcal H^{\mathrm{LoG}}_{\Lambda_k}\) is a canonical
one-dimensional odd line;
\item the compressed continued odd quadratic form is represented by the rank-one self-adjoint operator
\begin{equation}
K^{\mathrm{LoG}}_{\Lambda_k}
:=
\beta_k\,P^{\mathrm{LoG}}_{\Lambda_k},
\label{eq:Klog-def}
\end{equation}
with \(\beta_k<0\);
\item the unique negative spectral projector of \(K^{\mathrm{LoG}}_{\Lambda_k}\) is
\begin{equation}
P_{-,\Lambda_k}^{\mathrm{LoG}}
=
P^{\mathrm{LoG}}_{\Lambda_k}
=
u^{\mathrm{LoG}}_{\Lambda_k}\otimes u^{\mathrm{LoG}}_{\Lambda_k};
\label{eq:Pminus-log}
\end{equation}
\item along any admissible branch, the projector \(P_{-,\Lambda}^{\mathrm{LoG}}\) is \(C^1\), hence
the compressed active odd line is canonically and smoothly transported.
\end{enumerate}
\end{theorem}

\begin{proof}
By \cref{prop:compressed-log-lock}, the admissible welded channel at fixed accepted tick is exactly the
one-dimensional line \(\mathcal H^{\mathrm{LoG}}_{\Lambda_k}\). By
\cref{def:compressed-odd-coeff}, the compressed continued odd quadratic form on that line is represented
by the scalar \(\beta_k\), equivalently by the self-adjoint rank-one operator
\eqref{eq:Klog-def}.

By \cref{prop:compressed-beta-estimate},
\[
\beta_k
\le
-\,\alpha(\Lambda_k)\mu_{ps}^2
\;+\;
\Big(
C_{\rm int}\,\varepsilon_k^{\eta}
\;+\;
C_{\rm DoG}\,\Phi(\sigma_{t,k}\Omega_{\rm gap})
\;+\;
C_{\rm mem}\,e^{-\Delta_k/\tau_b}
\;+\;
C_4\,\mu_{ps}^4
\Big).
\]
Under the certified regime inequality \eqref{eq:compressed-certified}, the bracketed term is strictly
smaller than \(\tfrac12\alpha(\Lambda_k)\mu_{ps}^2\). Therefore
\[
\beta_k
<
-\frac12\,\alpha(\Lambda_k)\mu_{ps}^2
<
0.
\]
Thus \(K^{\mathrm{LoG}}_{\Lambda_k}\) has exactly one negative eigenvalue, namely \(\beta_k\), and
its unique negative spectral subspace is the entire line
\(\mathcal H^{\mathrm{LoG}}_{\Lambda_k}\). Since
\(
P^{\mathrm{LoG}}_{\Lambda_k}
=
u^{\mathrm{LoG}}_{\Lambda_k}\otimes u^{\mathrm{LoG}}_{\Lambda_k}
\),
the negative spectral projector is exactly \eqref{eq:Pminus-log}.

Finally, \cref{prop:compressed-log-smooth} gives \(C^1\)-dependence of
\(\Lambda\mapsto u^{\mathrm{LoG}}_\Lambda\), hence also of
\(\Lambda\mapsto P_{-,\Lambda}^{\mathrm{LoG}}\). This proves the canonical branchwise transport of the
compressed active odd line.
\end{proof}

\begin{corollary}[Canonical compressed odd projector]
\label{cor:compressed-export-selection}
For each accepted tick in the certified compressed spectral regime, the compressed, co-centered,
band-locked welded LoG channel carries the canonical rank-one projector
\[
P_{-,\Lambda_k}^{\mathrm{LoG}}
\]
onto its negative odd line. Equivalently,
\[
\mathcal H_{\Lambda_k}^{\mathrm{LoG}}
=
\operatorname{ran} P_{-,\Lambda_k}^{\mathrm{LoG}}.
\]
\end{corollary}

\begin{proof}
This is exactly \cref{thm:finite-cutoff-canonical-odd-spectral-selection}: the theorem identifies the
negative spectral projector with the rank-one projector
\(
u^{\mathrm{LoG}}_{\Lambda_k}\otimes u^{\mathrm{LoG}}_{\Lambda_k}
\),
whose range is \(\mathcal H^{\mathrm{LoG}}_{\Lambda_k}\).
\end{proof}

\begin{corollary}[Completion-sector realization of the compressed projector]
\label{cor:compressed-export-selection-completion}
Assume \cref{ass:gaussian-completion}. Then, for each accepted tick in the certified compressed
spectral regime, any Gaussian completion realization identifies the compressed projector
\(
P_{-,\Lambda_k}^{\mathrm{LoG}}
\)
with a unique rank-one active odd line in the completion sector. Equivalently, there exists a unit
vector
\[
u_-(\Lambda_k)
\]
spanning that realized line such that
\[
P_{-,\Lambda_k}^{\mathrm{LoG}}
=
u_-(\Lambda_k)\otimes u_-(\Lambda_k)
\]
in the chosen realization.
\end{corollary}

\begin{proof}
Under \cref{ass:gaussian-completion}, the finite active completion-mode sector is defined only up to
symplectic equivalence. A one-dimensional real line is invariant under such equivalences as a subspace,
even though a basis vector on that line is not canonical before normalization. By
\cref{thm:finite-cutoff-canonical-odd-spectral-selection}, the compressed odd line is already
canonically fixed abstractly by the projector \(P_{-,\Lambda_k}^{\mathrm{LoG}}\). Choosing the unique
unit vector on that line in a given realization yields the stated formula, and changing the completion
realization changes only the coordinate representative, not the underlying rank-one line.
\end{proof}

\begin{remark}[Scope of the compressed spectral theorem]
\label{rem:compressed-spectral-scope}
Theorem \cref{thm:finite-cutoff-canonical-odd-spectral-selection} determines the odd line only on
the compressed welded LoG channel. It does not classify the full completion-sector odd subspace, and
the completion-side statement of \cref{cor:compressed-export-selection-completion} is only a
realization statement for this rank-one projector.
\end{remark}

\begin{corollary}[Nontrivial descendant hit of the compressed dominant line]
\label{cor:compressed-transport-hit}
Assume \cref{thm:char-weak-interface}. Then for each accepted tick \(k\), there exists a
representative
\[
J^{\mathrm{desc}}_{\Lambda_k}\in[\Xi^{(1)}_{\Lambda_k}]
\]
whose response has nonzero component on the canonically selected compressed dominant odd line.
Equivalently, the transported descendant sector maps nontrivially into the compressed line fixed by
\(
P_{-,\Lambda_k}^{\mathrm{LoG}}
\).
\end{corollary}

\begin{proof}
This is exactly item (1) of \cref{thm:char-weak-interface}.
\end{proof}

\ifdefined\UOMTransportFragmentLoaded\else
\end{document}
\fi
%
  }{%
    \ifdefined\UOMTransportFragmentLoaded\else
\documentclass[11pt]{article}
\usepackage[margin=1in]{geometry}
\usepackage{amsmath,amssymb,amsthm,mathtools}
\usepackage{bm}
\usepackage{microtype}
\usepackage{enumitem}
\usepackage{booktabs}
\usepackage{array}
\usepackage{xcolor}
\usepackage{xr-hyper}
\usepackage{hyperref}
\usepackage[nameinlink,noabbrev]{cleveref}

\providecommand{\Fsloc}{\mathcal F_{\rm s.loc}}
\providecommand{\Qt}{Q_t}
\externaldocument[main-]{UOM/main}
\externaldocument[uomdesc-]{UOM/uom_descendant_note}
\externaldocument{UOM/completion_sector_split_model_note}

\hypersetup{
  colorlinks=true,
  linkcolor=blue!60!black,
  citecolor=blue!60!black,
  urlcolor=blue!60!black
}

\theoremstyle{definition}
\newtheorem{definition}{Definition}[section]
\newtheorem{assumption}{Assumption}[section]
\theoremstyle{plain}
\newtheorem{theorem}[definition]{Theorem}
\newtheorem{lemma}[definition]{Lemma}
\newtheorem{proposition}[definition]{Proposition}
\newtheorem{corollary}[definition]{Corollary}
\theoremstyle{remark}
\newtheorem{remark}[definition]{Remark}
\crefname{definition}{definition}{definitions}
\Crefname{definition}{Definition}{Definitions}
\crefname{assumption}{assumption}{assumptions}
\Crefname{assumption}{Assumption}{Assumptions}
\crefname{theorem}{theorem}{theorems}
\Crefname{theorem}{Theorem}{Theorems}
\crefname{lemma}{lemma}{lemmas}
\Crefname{lemma}{Lemma}{Lemmas}
\crefname{proposition}{proposition}{propositions}
\Crefname{proposition}{Proposition}{Propositions}
\crefname{corollary}{corollary}{corollaries}
\Crefname{corollary}{Corollary}{Corollaries}
\crefname{remark}{remark}{remarks}
\Crefname{remark}{Remark}{Remarks}

\newcommand{\cH}{\mathcal H}
\newcommand{\cA}{\mathcal A}
\newcommand{\cD}{\mathcal D}
\newcommand{\cB}{\mathcal B}
\newcommand{\cC}{\mathcal C}
\newcommand{\cL}{\mathcal L}
\newcommand{\cK}{\mathcal K}
\newcommand{\cR}{\mathcal R}
\newcommand{\cG}{\mathcal G}
\newcommand{\Ztwo}{\mathbb Z_2}
\newcommand{\Sys}{\mathrm{Sys}}
\newcommand{\Time}{\mathrm{Time}}
\newcommand{\RG}{\mathrm{RG}}
\newcommand{\Tr}{\mathrm{Tr}}
\newcommand{\diag}{\mathrm{diag}}
\newcommand{\ad}{\mathrm{ad}}
\newcommand{\id}{\mathrm{id}}
\newcommand{\norm}[1]{\left\lVert #1\right\rVert}
\newcommand{\abs}[1]{\left\lvert #1\right\rvert}
\newcommand{\ip}[2]{\left\langle #1,#2\right\rangle}

\externaldocument{UOM/odd_selector_note}
\externaldocument{UOM/canonical_odd_interface_note}
\externaldocument{UOM/active_response_transport_program_note}
\title{Canonical Odd Spectral Selection on the Compressed Welded LoG Channel}
\author{}
\date{}
\begin{document}
\maketitle
\fi

\section{Canonical Odd Spectral Selection on the Compressed Welded LoG Channel}
\label{sec:compressed-log-spectral-selection}

Fix an accepted tick \(k\) on a welded admissible branch. This section isolates the compressed
welded LoG channel and proves that the continued odd quadratic form canonically selects a negative
rank-one odd line on that channel. The transport of the descendant class into that line is recorded
separately in \cref{thm:char-weak-interface}; here it is used only in
\cref{cor:compressed-transport-hit}.

The macro source-side temporal envelope is DoG, the macro source-side spatial envelope is LoG, the
micro receiver temporal envelope is DoG, and the micro receiver spatial envelope is Gaussian. By
Corollary~\ref{main-cor:welded-log}, the welded envelope is LoG in both time and space, with scales
adding in quadrature.

\begin{definition}[Canonical welded LoG profile]
\label{def:compressed-log-profile}
Let \(t_k\) be the accepted tick time and let \(\Omega_\ast\in S^3\) denote the common spatial center
selected by the co-centering/history-lock mechanism. Let
\[
s_{t,k}:=\sqrt{\Sigma_k^2+\sigma_{t,k}^2},
\qquad
s_{x,k}:=\sqrt{\rho_k^2+\sigma_{x,k}^2}
\]
be the effective welded scales. Define the normalized welded LoG pulse by
\begin{equation}
u^{\mathrm{LoG}}_{\Lambda_k}(t,\Omega)
:=
\mathcal N_k\,
\frac{d^2}{dt^2}g_{s_{t,k}}(t-t_k)\,
\Delta_{S^3}g_{s_{x,k}}(\Omega,\Omega_\ast),
\label{eq:ulog-def}
\end{equation}
where \(\mathcal N_k>0\) is chosen so that
\(
\|u^{\mathrm{LoG}}_{\Lambda_k}\|_{L^2(\mathbb R\times S^3)}=1
\).
We write
\begin{equation}
\mathcal H^{\mathrm{LoG}}_{\Lambda_k}
:=
\mathbb R\,u^{\mathrm{LoG}}_{\Lambda_k}
\label{eq:Hlog-def}
\end{equation}
for the compressed welded LoG channel.
\end{definition}

\begin{proposition}[Exact welded LoG form]
\label{prop:compressed-log-exact}
The normalized profile \(u^{\mathrm{LoG}}_{\Lambda_k}\) of
\cref{def:compressed-log-profile} is exactly the welded macro/micro profile produced by:
\begin{enumerate}[leftmargin=2em]
\item macro temporal DoG convolved with micro temporal DoG;
\item macro spatial LoG convolved with micro spatial Gaussian.
\end{enumerate}
In particular the welded profile is LoG in time and LoG in space, with scales adding in quadrature.
\end{proposition}

\begin{proof}
The temporal and spatial welded identities are exactly the two statements of
Corollary~\ref{main-cor:welded-log}. After multiplying the factors and normalizing in \(L^2\), one obtains
\eqref{eq:ulog-def}.
\end{proof}

\begin{proposition}[Band-edge locking and unique center]
\label{prop:compressed-log-lock}
Assume the accepted tick lies on an admissible band-locked welded branch in the sense of
Definition~\ref{main-def:eta-bridge-leak}. Then:
\begin{enumerate}[leftmargin=2em]
\item the effective welded widths satisfy
\[
s_{t,k}=\frac{\alpha_t^\ast}{\Lambda_k},
\qquad
s_{x,k}=\frac{\alpha_x^\ast}{v\,\Lambda_k},
\]
with \(\alpha_t^\ast,\alpha_x^\ast>0\) independent of \(k\) along a fixed admissible branch;
\item the macro and micro temporal centers coincide at \(t_k\), and the macro and micro spatial
  centers coincide at \(\Omega_\ast\);
\item once the receiver trace sector carries a nonzero inherited LoG background, the absolute center
\(\Omega_\ast\) is uniquely locked for all subsequent accepted ticks.
\end{enumerate}
Consequently, after normalization, the admissible welded pulse space at fixed accepted tick is the
one-dimensional line \(\mathcal H^{\mathrm{LoG}}_{\Lambda_k}\).
\end{proposition}

\begin{proof}
The width lock is the band-locking clause of Definition~\ref{main-def:eta-bridge-leak}. The
temporal and spatial co-centering statements are \eqref{main-eq:Ct} and \eqref{main-eq:Cx}, and the
persistence of the spatial center is exactly Proposition~\ref{main-prop:history-center-lock}.

Therefore, at fixed accepted tick, the shape, widths, and center are fixed. The only remaining freedom
is an overall real amplitude. After \(L^2\)-normalization, that freedom reduces to the sign
\(\pm u^{\mathrm{LoG}}_{\Lambda_k}\). Hence the admissible welded channel is precisely the
one-dimensional real line \(\mathcal H^{\mathrm{LoG}}_{\Lambda_k}\).
\end{proof}

\begin{proposition}[Smooth dependence along a branch]
\label{prop:compressed-log-smooth}
Along any admissible branch on which \(\Lambda\mapsto \alpha_t^\ast,\alpha_x^\ast\) are constant and
the locked center \(\Omega_\ast\) is fixed, the map
\[
\Lambda\longmapsto u^{\mathrm{LoG}}_\Lambda
\]
is \(C^1\) in \(L^2(\mathbb R\times S^3)\). Consequently the rank-one orthogonal projector
\begin{equation}
P^{\mathrm{LoG}}_\Lambda
:=
u^{\mathrm{LoG}}_\Lambda\otimes u^{\mathrm{LoG}}_\Lambda
\label{eq:Plog-def}
\end{equation}
is \(C^1\) along the branch.
\end{proposition}

\begin{proof}
Along the branch one has
\(
s_t(\Lambda)=\alpha_t^\ast/\Lambda
\)
and
\(
s_x(\Lambda)=\alpha_x^\ast/(v\Lambda)
\),
so both scales depend \(C^\infty\)-smoothly on \(\Lambda\). The Gaussian heat kernel and its time and
spatial derivatives depend smoothly on the scale parameter, hence the unnormalized welded LoG profile
depends \(C^1\)-smoothly on \(\Lambda\) in \(L^2\). Its \(L^2\)-norm is strictly positive, so
normalization preserves \(C^1\)-dependence. Formula \eqref{eq:Plog-def} then gives a \(C^1\)
projector-valued map.
\end{proof}

\begin{proposition}[Well-defined restriction of the continued quadratic form]
\label{prop:compressed-form-well-defined}
Assume the welded linear-response/Gaussian radiative-sector package of
Paragraph~\ref{main-para:uplift-decomp} and Theorem~\ref{main-thm:tick-fidelity-RP}. Then the
analytically continued odd quadratic response form is well-defined on the compressed welded channel
\(\mathcal H^{\mathrm{LoG}}_{\Lambda_k}\).
\end{proposition}

\begin{proof}
By Paragraph~\ref{main-para:uplift-decomp}, the welded source enters the Gaussian radiative sector
through the linear response map \(\mathsf T_\Lambda\), and Theorem~\ref{main-thm:tick-fidelity-RP}
supplies the finite-cutoff bounded Gaussian response package on the welded envelope class. By
\cref{prop:compressed-log-lock}, the compressed welded channel is a one-dimensional subspace of that
admissible class. Restricting any continuous quadratic form to a one-dimensional subspace is
automatic. Hence the analytically continued odd quadratic response form restricts canonically to
\(\mathcal H^{\mathrm{LoG}}_{\Lambda_k}\).
\end{proof}

\begin{definition}[Compressed continued odd quadratic coefficient]
\label{def:compressed-odd-coeff}
Fix a witness window \(\mu_{ps}>0\) in the accepted regime. Let
\[
\mathfrak q^{\mathrm{cont}}_{\Lambda_k,\mu_{ps}}
:
\mathcal H^{\mathrm{LoG}}_{\Lambda_k}\times \mathcal H^{\mathrm{LoG}}_{\Lambda_k}
\to
\mathbb R
\]
denote the analytically continued odd quadratic response form restricted to the compressed welded LoG
channel. Since \(\mathcal H^{\mathrm{LoG}}_{\Lambda_k}\) is one-dimensional, there is a unique scalar
\(\beta_k\in\mathbb R\) such that
\begin{equation}
\mathfrak q^{\mathrm{cont}}_{\Lambda_k,\mu_{ps}}(X,Y)
=
\beta_k\,
\langle X,u^{\mathrm{LoG}}_{\Lambda_k}\rangle_{L^2}
\langle Y,u^{\mathrm{LoG}}_{\Lambda_k}\rangle_{L^2}
\qquad
\forall X,Y\in\mathcal H^{\mathrm{LoG}}_{\Lambda_k}.
\label{eq:qcont-beta}
\end{equation}
Equivalently,
\[
\beta_k
=
\mathfrak q^{\mathrm{cont}}_{\Lambda_k,\mu_{ps}}
\big(
u^{\mathrm{LoG}}_{\Lambda_k},
u^{\mathrm{LoG}}_{\Lambda_k}
\big).
\]
\end{definition}

\begin{proposition}[Rigorous control of the compressed coefficient]
\label{prop:compressed-beta-estimate}
Assume:
\begin{enumerate}[leftmargin=2em]
\item the smooth RP degradation law of Theorem~\ref{main-thm:smooth-RP};
\item the approximate intertwiner bound of Lemma~\ref{main-lem:approx-intertwiner-new};
\item the temporal DoG commutator bound of Lemma~\ref{main-lem:DoG-comm-new};
\item the exponentially decaying memory-tail package of
  Lemma~\ref{main-lem:app-exp-backflow} and Proposition~\ref{main-prop:app-tick-progress}.
\end{enumerate}
Then there exist constants \(C_{\rm int},C_{\rm DoG},C_{\rm mem},C_4>0\) and \(\eta>0\), depending only on the
fixed welded family and not on \(k\), such that
\begin{equation}
\beta_k
=
-\,\alpha(\Lambda_k)\,\mu_{ps}^2
\;+\;
r_k,
\label{eq:beta-splitting}
\end{equation}
with remainder bounded by
\begin{equation}
|r_k|
\le
C_{\rm int}\,\varepsilon_k^{\eta}
\;+\;
C_{\rm DoG}\,\Phi(\sigma_{t,k}\Omega_{\rm gap})
\;+\;
C_{\rm mem}\,e^{-\Delta_k/\tau_b}
\;+\;
C_4\,\mu_{ps}^4.
\label{eq:beta-remainder}
\end{equation}
\end{proposition}

\begin{proof}
By Theorem~\ref{main-thm:smooth-RP}, the ideal analytically continued Gaussian OS problem gives
\(
\lambda_{\min}(\mu_{ps},\Lambda_k)
=
-\alpha(\Lambda_k)\mu_{ps}^2+O(\mu_{ps}^4)
\)
on the welded Gaussian reference kernel. Restricting that quadratic form to the canonical normalized
welded LoG profile \(u^{\mathrm{LoG}}_{\Lambda_k}\) produces the leading term
\(-\alpha(\Lambda_k)\mu_{ps}^2\), with the \(O(\mu_{ps}^4)\) correction bounded by
\(C_4\mu_{ps}^4\).

The regulated intertwiner contribution is bounded by
Lemma~\ref{main-lem:approx-intertwiner-new}, the temporal DoG contribution is bounded by
Lemma~\ref{main-lem:DoG-comm-new}, and the intrinsic KMS bath contributes only the
Born-Nakajima-Zwanzig memory tail controlled by Lemma~\ref{main-lem:app-exp-backflow} and
Proposition~\ref{main-prop:app-tick-progress}. Since \(u^{\mathrm{LoG}}_{\Lambda_k}\) is
normalized, each of these operator bounds transfers directly to the scalar quadratic coefficient.
Summing the four contributions gives \eqref{eq:beta-splitting}--\eqref{eq:beta-remainder}.
\end{proof}

\begin{definition}[Certified compressed spectral regime]
\label{def:compressed-spectral-regime}
We say that an accepted tick \(k\) lies in the \emph{certified compressed spectral regime} if
\begin{equation}
\alpha(\Lambda_k)\,\mu_{ps}^2
>
2\Big(
C_{\rm int}\,\varepsilon_k^{\eta}
\;+\;
C_{\rm DoG}\,\Phi(\sigma_{t,k}\Omega_{\rm gap})
\;+\;
C_{\rm mem}\,e^{-\Delta_k/\tau_b}
\;+\;
C_4\,\mu_{ps}^4
\Big).
\label{eq:compressed-certified}
\end{equation}
\end{definition}

\begin{theorem}[Finite-cutoff canonical odd spectral selection on the compressed welded LoG channel]
\label{thm:finite-cutoff-canonical-odd-spectral-selection}
Assume
\cref{prop:compressed-log-exact,prop:compressed-log-lock,prop:compressed-log-smooth,prop:compressed-form-well-defined,prop:compressed-beta-estimate}.
Let \(k\) be an accepted tick in the certified compressed spectral regime of
\cref{def:compressed-spectral-regime}. Then:
\begin{enumerate}[leftmargin=2em]
\item the compressed welded channel \(\mathcal H^{\mathrm{LoG}}_{\Lambda_k}\) is a canonical
one-dimensional odd line;
\item the compressed continued odd quadratic form is represented by the rank-one self-adjoint operator
\begin{equation}
K^{\mathrm{LoG}}_{\Lambda_k}
:=
\beta_k\,P^{\mathrm{LoG}}_{\Lambda_k},
\label{eq:Klog-def}
\end{equation}
with \(\beta_k<0\);
\item the unique negative spectral projector of \(K^{\mathrm{LoG}}_{\Lambda_k}\) is
\begin{equation}
P_{-,\Lambda_k}^{\mathrm{LoG}}
=
P^{\mathrm{LoG}}_{\Lambda_k}
=
u^{\mathrm{LoG}}_{\Lambda_k}\otimes u^{\mathrm{LoG}}_{\Lambda_k};
\label{eq:Pminus-log}
\end{equation}
\item along any admissible branch, the projector \(P_{-,\Lambda}^{\mathrm{LoG}}\) is \(C^1\), hence
the compressed active odd line is canonically and smoothly transported.
\end{enumerate}
\end{theorem}

\begin{proof}
By \cref{prop:compressed-log-lock}, the admissible welded channel at fixed accepted tick is exactly the
one-dimensional line \(\mathcal H^{\mathrm{LoG}}_{\Lambda_k}\). By
\cref{def:compressed-odd-coeff}, the compressed continued odd quadratic form on that line is represented
by the scalar \(\beta_k\), equivalently by the self-adjoint rank-one operator
\eqref{eq:Klog-def}.

By \cref{prop:compressed-beta-estimate},
\[
\beta_k
\le
-\,\alpha(\Lambda_k)\mu_{ps}^2
\;+\;
\Big(
C_{\rm int}\,\varepsilon_k^{\eta}
\;+\;
C_{\rm DoG}\,\Phi(\sigma_{t,k}\Omega_{\rm gap})
\;+\;
C_{\rm mem}\,e^{-\Delta_k/\tau_b}
\;+\;
C_4\,\mu_{ps}^4
\Big).
\]
Under the certified regime inequality \eqref{eq:compressed-certified}, the bracketed term is strictly
smaller than \(\tfrac12\alpha(\Lambda_k)\mu_{ps}^2\). Therefore
\[
\beta_k
<
-\frac12\,\alpha(\Lambda_k)\mu_{ps}^2
<
0.
\]
Thus \(K^{\mathrm{LoG}}_{\Lambda_k}\) has exactly one negative eigenvalue, namely \(\beta_k\), and
its unique negative spectral subspace is the entire line
\(\mathcal H^{\mathrm{LoG}}_{\Lambda_k}\). Since
\(
P^{\mathrm{LoG}}_{\Lambda_k}
=
u^{\mathrm{LoG}}_{\Lambda_k}\otimes u^{\mathrm{LoG}}_{\Lambda_k}
\),
the negative spectral projector is exactly \eqref{eq:Pminus-log}.

Finally, \cref{prop:compressed-log-smooth} gives \(C^1\)-dependence of
\(\Lambda\mapsto u^{\mathrm{LoG}}_\Lambda\), hence also of
\(\Lambda\mapsto P_{-,\Lambda}^{\mathrm{LoG}}\). This proves the canonical branchwise transport of the
compressed active odd line.
\end{proof}

\begin{corollary}[Canonical compressed odd projector]
\label{cor:compressed-export-selection}
For each accepted tick in the certified compressed spectral regime, the compressed, co-centered,
band-locked welded LoG channel carries the canonical rank-one projector
\[
P_{-,\Lambda_k}^{\mathrm{LoG}}
\]
onto its negative odd line. Equivalently,
\[
\mathcal H_{\Lambda_k}^{\mathrm{LoG}}
=
\operatorname{ran} P_{-,\Lambda_k}^{\mathrm{LoG}}.
\]
\end{corollary}

\begin{proof}
This is exactly \cref{thm:finite-cutoff-canonical-odd-spectral-selection}: the theorem identifies the
negative spectral projector with the rank-one projector
\(
u^{\mathrm{LoG}}_{\Lambda_k}\otimes u^{\mathrm{LoG}}_{\Lambda_k}
\),
whose range is \(\mathcal H^{\mathrm{LoG}}_{\Lambda_k}\).
\end{proof}

\begin{corollary}[Completion-sector realization of the compressed projector]
\label{cor:compressed-export-selection-completion}
Assume \cref{ass:gaussian-completion}. Then, for each accepted tick in the certified compressed
spectral regime, any Gaussian completion realization identifies the compressed projector
\(
P_{-,\Lambda_k}^{\mathrm{LoG}}
\)
with a unique rank-one active odd line in the completion sector. Equivalently, there exists a unit
vector
\[
u_-(\Lambda_k)
\]
spanning that realized line such that
\[
P_{-,\Lambda_k}^{\mathrm{LoG}}
=
u_-(\Lambda_k)\otimes u_-(\Lambda_k)
\]
in the chosen realization.
\end{corollary}

\begin{proof}
Under \cref{ass:gaussian-completion}, the finite active completion-mode sector is defined only up to
symplectic equivalence. A one-dimensional real line is invariant under such equivalences as a subspace,
even though a basis vector on that line is not canonical before normalization. By
\cref{thm:finite-cutoff-canonical-odd-spectral-selection}, the compressed odd line is already
canonically fixed abstractly by the projector \(P_{-,\Lambda_k}^{\mathrm{LoG}}\). Choosing the unique
unit vector on that line in a given realization yields the stated formula, and changing the completion
realization changes only the coordinate representative, not the underlying rank-one line.
\end{proof}

\begin{remark}[Scope of the compressed spectral theorem]
\label{rem:compressed-spectral-scope}
Theorem \cref{thm:finite-cutoff-canonical-odd-spectral-selection} determines the odd line only on
the compressed welded LoG channel. It does not classify the full completion-sector odd subspace, and
the completion-side statement of \cref{cor:compressed-export-selection-completion} is only a
realization statement for this rank-one projector.
\end{remark}

\begin{corollary}[Nontrivial descendant hit of the compressed dominant line]
\label{cor:compressed-transport-hit}
Assume \cref{thm:char-weak-interface}. Then for each accepted tick \(k\), there exists a
representative
\[
J^{\mathrm{desc}}_{\Lambda_k}\in[\Xi^{(1)}_{\Lambda_k}]
\]
whose response has nonzero component on the canonically selected compressed dominant odd line.
Equivalently, the transported descendant sector maps nontrivially into the compressed line fixed by
\(
P_{-,\Lambda_k}^{\mathrm{LoG}}
\).
\end{corollary}

\begin{proof}
This is exactly item (1) of \cref{thm:char-weak-interface}.
\end{proof}

\ifdefined\UOMTransportFragmentLoaded\else
\end{document}
\fi
%
  }%
}

\IfFileExists{odd_selector_note.tex}{%
  \ifdefined\UOMTransportFragmentLoaded\else
\documentclass[11pt]{article}
\usepackage[margin=1in]{geometry}
\usepackage{amsmath,amssymb,amsthm,mathtools}
\usepackage{bm}
\usepackage{microtype}
\usepackage{enumitem}
\usepackage{booktabs}
\usepackage{array}
\usepackage{xcolor}
\usepackage{xr-hyper}
\usepackage{hyperref}
\usepackage[nameinlink,noabbrev]{cleveref}

\providecommand{\Fsloc}{\mathcal F_{\rm s.loc}}
\providecommand{\Qt}{Q_t}
\externaldocument[main-]{UOM/main}
\externaldocument[uomdesc-]{UOM/uom_descendant_note}
\externaldocument{UOM/completion_sector_split_model_note}

\hypersetup{
  colorlinks=true,
  linkcolor=blue!60!black,
  citecolor=blue!60!black,
  urlcolor=blue!60!black
}

\theoremstyle{definition}
\newtheorem{definition}{Definition}[section]
\newtheorem{assumption}{Assumption}[section]
\theoremstyle{plain}
\newtheorem{theorem}[definition]{Theorem}
\newtheorem{lemma}[definition]{Lemma}
\newtheorem{proposition}[definition]{Proposition}
\newtheorem{corollary}[definition]{Corollary}
\theoremstyle{remark}
\newtheorem{remark}[definition]{Remark}
\crefname{definition}{definition}{definitions}
\Crefname{definition}{Definition}{Definitions}
\crefname{assumption}{assumption}{assumptions}
\Crefname{assumption}{Assumption}{Assumptions}
\crefname{theorem}{theorem}{theorems}
\Crefname{theorem}{Theorem}{Theorems}
\crefname{lemma}{lemma}{lemmas}
\Crefname{lemma}{Lemma}{Lemmas}
\crefname{proposition}{proposition}{propositions}
\Crefname{proposition}{Proposition}{Propositions}
\crefname{corollary}{corollary}{corollaries}
\Crefname{corollary}{Corollary}{Corollaries}
\crefname{remark}{remark}{remarks}
\Crefname{remark}{Remark}{Remarks}

\newcommand{\cH}{\mathcal H}
\newcommand{\cA}{\mathcal A}
\newcommand{\cD}{\mathcal D}
\newcommand{\cB}{\mathcal B}
\newcommand{\cC}{\mathcal C}
\newcommand{\cL}{\mathcal L}
\newcommand{\cK}{\mathcal K}
\newcommand{\cR}{\mathcal R}
\newcommand{\cG}{\mathcal G}
\newcommand{\Ztwo}{\mathbb Z_2}
\newcommand{\Sys}{\mathrm{Sys}}
\newcommand{\Time}{\mathrm{Time}}
\newcommand{\RG}{\mathrm{RG}}
\newcommand{\Tr}{\mathrm{Tr}}
\newcommand{\diag}{\mathrm{diag}}
\newcommand{\ad}{\mathrm{ad}}
\newcommand{\id}{\mathrm{id}}
\newcommand{\norm}[1]{\left\lVert #1\right\rVert}
\newcommand{\abs}[1]{\left\lvert #1\right\rvert}
\newcommand{\ip}[2]{\left\langle #1,#2\right\rangle}

\externaldocument{UOM/compressed_log_spectral_selection_note}
\externaldocument{UOM/canonical_odd_interface_note}
\externaldocument{UOM/active_response_transport_program_note}
\title{Raw Odd Selector and Scheme-Stable Affine Output}
\author{}
\date{}
\begin{document}
\maketitle
\fi

\section{Raw Odd Selector and Scheme-Stable Affine Output}
\label{sec:odd-selector}

This section isolates the raw selector step on the odd descendant side of the UOM-to-Main
interface used in Definition~\ref{main-def:uom-main-theorem-assumptions}.
At fixed finite cutoff, it records the \(\Gamma\)-odd affine output determined by a chosen
sufficiently local anomaly representative and its same-scheme and inter-scheme translation laws.

\paragraph{Convention on the two \texorpdfstring{\(\mathbb Z_2\)}{Z2} gradings.}
The cohomological degree for the \(Q_t\)-complex is written with bars:
\[
(\cdot)_{\bar 0},
\qquad
(\cdot)_{\bar 1}.
\]
Physical parity refers to the swap involution \(\Gamma\), with eigenspaces
\[
\Gamma F=\pm F.
\]
Thus the selector output constructed below is cohomologically degree \(0\) but physically odd.
Unless bar-degree notation is used explicitly, ``odd'' below means \(\Gamma\)-odd rather than
\(Q_t\)-degree \(1\).

Work at one accepted tick \(k\) and one branch value \(u\) at cutoff \(\Lambda_k\).
Use the raw UOM descendant package of
\cref{uomdesc-thm:uomdesc-raw-package,uomdesc-thm:uomdesc-conditional-odd-package}.
In particular:
\begin{enumerate}[leftmargin=2em]
\item a sufficiently local good monotone representative
  \[
  \mathcal O^{(0)}_{\Lambda_k}(u);
  \]
\item its canonical anomaly representative
  \[
  \mathcal A^{(1)}_{\Lambda_k}(u)
  :=
  Q^{\rm phys}_{u,\Lambda_k}\mathcal O^{(0)}_{\Lambda_k}(u);
  \]
\item a swap involution
  \[
  \Gamma:
  \big(\mathcal F_{\rm s.loc}^{\Lambda_k}(\mathcal X_k)\big)_{\bar\imath}
  \to
  \big(\mathcal F_{\rm s.loc}^{\Lambda_k}(\mathcal X_k)\big)_{\bar\imath}
  \qquad
  (\bar\imath\in\{\bar 0,\bar 1\})
  \]
  preserving sufficient locality and cohomological degree, and commuting with \(Q_t\):
  \[
  \Gamma Q_t=Q_t\Gamma;
  \]
\item swap-evenness of the canonical core,
  \[
  \Gamma\mathcal O^{(0)}_{\Lambda_k}(u)=\mathcal O^{(0)}_{\Lambda_k}(u),
  \qquad
  \Gamma\mathcal A^{(1)}_{\Lambda_k}(u)=\mathcal A^{(1)}_{\Lambda_k}(u);
  \]
\item swap covariance of the fixed-scheme physical radial differential on the chosen monotone
  sector,
  \[
  \Gamma Q^{\rm phys}_{u,\Lambda_k}
  =
  Q^{\rm phys}_{u,\Lambda_k}\Gamma;
  \]
\end{enumerate}
Items (1) and (2) are contained in \cref{uomdesc-thm:uomdesc-raw-package};
item (3) uses \cref{uomdesc-def:uomdesc-gamma,uomdesc-lem:uomdesc-gamma-comm};
item (4) uses \cref{uomdesc-prop:uomdesc-gamma-even-core}; and item (5) uses
\cref{uomdesc-prop:uomdesc-swap-split}.

Fix a chosen \(\Gamma\)-stable selector class of sufficiently local odd representatives of the same
\(Q_t\)-cohomology class as \(\mathcal A^{(1)}_{\Lambda_k}(u)\), and restrict to representatives
for which the exact descent equation
\begin{equation}
Q^{\rm phys}_{u,\Lambda_k}\mathcal O^{(0)}_{\Lambda_k}(u)
+
Q_t\widetilde{\Xi}^{(1)}_{\Lambda_k}(u)
=
\widetilde{\mathcal A}^{(1)}_{\Lambda_k}(u)
\label{eq:selector-full-descent}
\end{equation}
admits at least one sufficiently local even solution
\(
\widetilde{\Xi}^{(1)}_{\Lambda_k}(u)
\in
\big(\mathcal F_{\rm s.loc}^{\Lambda_k}(\mathcal X_k)\big)_{\bar 0}
\).

\begin{definition}[Selector class and canonical odd kernel channel]
\label{def:selector-class}
Let
\[
\mathcal C_{{\rm sel},\Lambda_k}(u)
\subset
\big(\mathcal F_{\rm s.loc}^{\Lambda_k}(\mathcal X_k)\big)_{\bar 1}
\]
be the chosen \(\Gamma\)-stable class of sufficiently local odd representatives of the same
degree-one \(Q_t\)-cohomology class as \(\mathcal A^{(1)}_{\Lambda_k}(u)\) for which
\eqref{eq:selector-full-descent} admits a sufficiently local even solution.
For
\(
\widetilde{\mathcal A}^{(1)}_{\Lambda_k}(u)\in\mathcal C_{{\rm sel},\Lambda_k}(u)
\),
define its swap-even and swap-odd parts by
\[
\widetilde{\mathcal A}^{(1)}_{\Lambda_k,\pm}(u)
:=
\frac12\Big(
\widetilde{\mathcal A}^{(1)}_{\Lambda_k}(u)
\pm
\Gamma\widetilde{\mathcal A}^{(1)}_{\Lambda_k}(u)
\Big).
\]
Define the canonical odd kernel channel by
\[
\mathcal K^-_{\Lambda_k}
:=
\ker(Q_t)
\cap
\big(\mathcal F_{\rm s.loc}^{\Lambda_k}(\mathcal X_k)\big)_{\bar 0}
\cap
\ker(\Gamma+\mathrm{id}).
\]
\end{definition}

\begin{proposition}[Odd defect of a selector representative is always \(Q_t\)-exact]
\label{prop:selector-odd-defect}
Let
\(
\widetilde{\mathcal A}^{(1)}_{\Lambda_k}(u)\in\mathcal C_{{\rm sel},\Lambda_k}(u)
\).
Then:
\begin{enumerate}[leftmargin=2em]
\item
  \(
  \widetilde{\mathcal A}^{(1)}_{\Lambda_k,+}(u)
  \)
  and
  \(
  \widetilde{\mathcal A}^{(1)}_{\Lambda_k,-}(u)
  \)
  are sufficiently local;
\item
  \(
  \widetilde{\mathcal A}^{(1)}_{\Lambda_k,+}(u)
  \)
  is a swap-even representative of the same \(Q_t\)-cohomology class as
  \(\mathcal A^{(1)}_{\Lambda_k}(u)\);
\item
  \[
  \widetilde{\mathcal A}^{(1)}_{\Lambda_k,-}(u)
  \in
  \operatorname{im}(Q_t)
  \cap
  \big(\mathcal F_{\rm s.loc}^{\Lambda_k}(\mathcal X_k)\big)_{\bar 1}
  \cap
  \ker(\Gamma+\mathrm{id}).
  \]
\end{enumerate}
\end{proposition}

\begin{proof}
Sufficient locality is preserved by linear combinations, proving item (1).

Because
\(
\widetilde{\mathcal A}^{(1)}_{\Lambda_k}(u)
\)
and
\(
\mathcal A^{(1)}_{\Lambda_k}(u)
\)
represent the same \(Q_t\)-cohomology class, there exists
\(
Y_{\Lambda_k}(u)\in
\big(\mathcal F_{\rm s.loc}^{\Lambda_k}(\mathcal X_k)\big)_{\bar 0}
\)
such that
\[
\widetilde{\mathcal A}^{(1)}_{\Lambda_k}(u)
-
\mathcal A^{(1)}_{\Lambda_k}(u)
=
Q_tY_{\Lambda_k}(u).
\]
Taking the swap-even part and using
\(
\Gamma Q_t=Q_t\Gamma
\)
and
\(
\Gamma\mathcal A^{(1)}_{\Lambda_k}(u)=\mathcal A^{(1)}_{\Lambda_k}(u)
\),
we obtain
\[
\widetilde{\mathcal A}^{(1)}_{\Lambda_k,+}(u)
-
\mathcal A^{(1)}_{\Lambda_k}(u)
=
Q_tY_{\Lambda_k,+}(u),
\]
so item (2) follows.

Taking the swap-odd part of the same identity gives
\[
\widetilde{\mathcal A}^{(1)}_{\Lambda_k,-}(u)
=
Q_tY_{\Lambda_k,-}(u).
\]
By construction,
\(
\Gamma\widetilde{\mathcal A}^{(1)}_{\Lambda_k,-}(u)
=
-\widetilde{\mathcal A}^{(1)}_{\Lambda_k,-}(u)
\),
which proves item (3).
\end{proof}

\begin{definition}[Selected odd descendant affine class]
\label{def:selector-affine-class}
For
\(
\widetilde{\mathcal A}^{(1)}_{\Lambda_k}(u)\in\mathcal C_{{\rm sel},\Lambda_k}(u)
\),
define the associated raw odd selector output by
\[
\mathfrak X_{\Lambda_k}^{{\rm sel},-}
\big(u;\widetilde{\mathcal A}^{(1)}_{\Lambda_k}\big)
:=
\left\{
\Xi^-_{\Lambda_k}(u)
\in
\big(\mathcal F_{\rm s.loc}^{\Lambda_k}(\mathcal X_k)\big)_{\bar 0}
\;\middle|\;
\Gamma\Xi^-_{\Lambda_k}(u)=-\Xi^-_{\Lambda_k}(u),
\quad
Q_t\Xi^-_{\Lambda_k}(u)
=
\widetilde{\mathcal A}^{(1)}_{\Lambda_k,-}(u)
\right\}.
\]
\end{definition}

\begin{remark}[Dependence only on the swap-odd defect]
\label{rem:selector-odd-defect-only}
By \cref{thm:selector-decomposition}, the affine class
\(
\mathfrak X_{\Lambda_k}^{{\rm sel},-}
\big(u;\widetilde{\mathcal A}^{(1)}_{\Lambda_k}\big)
\)
depends only on the swap-odd defect
\(
\widetilde{\mathcal A}^{(1)}_{\Lambda_k,-}(u)
\).
\end{remark}

\begin{theorem}[Selector decomposition into even representative channel and odd output]
\label{thm:selector-decomposition}
Let
\(
\widetilde{\mathcal A}^{(1)}_{\Lambda_k}(u)\in\mathcal C_{{\rm sel},\Lambda_k}(u)
\),
and let
\(
\widetilde{\Xi}^{(1)}_{\Lambda_k}(u)
\)
be any sufficiently local even solution of \eqref{eq:selector-full-descent}.
Define
\[
\widetilde{\Xi}^{(1)}_{\Lambda_k,\pm}(u)
:=
\frac12\Big(
\widetilde{\Xi}^{(1)}_{\Lambda_k}(u)
\pm
\Gamma\widetilde{\Xi}^{(1)}_{\Lambda_k}(u)
\Big).
\]
Then:
\begin{enumerate}[leftmargin=2em]
\item the swap-even and swap-odd channels decouple as
  \begin{align}
  Q^{\rm phys}_{u,\Lambda_k}\mathcal O^{(0)}_{\Lambda_k}(u)
  +
  Q_t\widetilde{\Xi}^{(1)}_{\Lambda_k,+}(u)
  &=
  \widetilde{\mathcal A}^{(1)}_{\Lambda_k,+}(u),
  \label{eq:selector-even-channel}
  \\
  Q_t\widetilde{\Xi}^{(1)}_{\Lambda_k,-}(u)
  &=
  \widetilde{\mathcal A}^{(1)}_{\Lambda_k,-}(u);
  \label{eq:selector-odd-channel}
  \end{align}
\item the raw odd selector output
  \(
  \mathfrak X_{\Lambda_k}^{{\rm sel},-}
  \big(u;\widetilde{\mathcal A}^{(1)}_{\Lambda_k}\big)
  \)
  contains the odd projection of any full solution; in particular it is nonempty;
\item it is an affine space over the canonical odd kernel channel:
  \[
  \mathfrak X_{\Lambda_k}^{{\rm sel},-}
  \big(u;\widetilde{\mathcal A}^{(1)}_{\Lambda_k}\big)
  =
  \widetilde{\Xi}^{(1)}_{\Lambda_k,-}(u)
  +
  \mathcal K^-_{\Lambda_k}.
  \]
\end{enumerate}
\end{theorem}

\begin{proof}
Apply \(\Gamma\) to \eqref{eq:selector-full-descent}.
Using
\(
\Gamma Q^{\rm phys}_{u,\Lambda_k}
=
Q^{\rm phys}_{u,\Lambda_k}\Gamma
\)
on the chosen monotone sector and
\(
\Gamma\mathcal O^{(0)}_{\Lambda_k}(u)=\mathcal O^{(0)}_{\Lambda_k}(u)
\),
we obtain
\[
Q^{\rm phys}_{u,\Lambda_k}\mathcal O^{(0)}_{\Lambda_k}(u)
+
Q_t\Gamma\widetilde{\Xi}^{(1)}_{\Lambda_k}(u)
=
\Gamma\widetilde{\mathcal A}^{(1)}_{\Lambda_k}(u).
\]
Adding and subtracting this equation from \eqref{eq:selector-full-descent} yields
\eqref{eq:selector-even-channel} and \eqref{eq:selector-odd-channel}, proving item (1).

Item (2) is immediate from item (1): the odd projection
\(
\widetilde{\Xi}^{(1)}_{\Lambda_k,-}(u)
\)
of any full solution belongs to the set in
\cref{def:selector-affine-class}.

Now let
\(
\Xi^-_{\Lambda_k,1}(u)
\)
and
\(
\Xi^-_{\Lambda_k,2}(u)
\)
lie in
\(
\mathfrak X_{\Lambda_k}^{{\rm sel},-}
\big(u;\widetilde{\mathcal A}^{(1)}_{\Lambda_k}\big)
\).
Subtracting their defining equations gives
\[
Q_t\big(
\Xi^-_{\Lambda_k,1}(u)-\Xi^-_{\Lambda_k,2}(u)
\big)
=
0.
\]
Because both are \(\Gamma\)-odd, their difference lies in
\(
\mathcal K^-_{\Lambda_k}
\).
Conversely, if
\(
K^-_{\Lambda_k}\in\mathcal K^-_{\Lambda_k}
\),
then
\[
Q_t\big(
\widetilde{\Xi}^{(1)}_{\Lambda_k,-}(u)+K^-_{\Lambda_k}
\big)
=
\widetilde{\mathcal A}^{(1)}_{\Lambda_k,-}(u),
\]
and the sum remains \(\Gamma\)-odd.
This proves item (3).
\end{proof}

\begin{corollary}[The two selector channels]
\label{cor:selector-two-channels}
At fixed cutoff, the selector problem splits into exactly two channels.
\begin{enumerate}[leftmargin=2em]
\item \textbf{Kernel channel.}
  If
  \(
  \widetilde{\mathcal A}^{(1)}_{\Lambda_k,-}(u)=0
  \),
  then
  \[
  \mathfrak X_{\Lambda_k}^{{\rm sel},-}
  \big(u;\widetilde{\mathcal A}^{(1)}_{\Lambda_k}\big)
  =
  \mathcal K^-_{\Lambda_k}.
  \]
  In particular, the canonical anomaly representative
  \(
  \widetilde{\mathcal A}^{(1)}_{\Lambda_k}(u)=\mathcal A^{(1)}_{\Lambda_k}(u)
  \)
  produces only the canonical odd kernel sector.
\item \textbf{Representative channel.}
  If
  \(
  \widetilde{\mathcal A}^{(1)}_{\Lambda_k,-}(u)\neq 0
  \),
  then the raw odd selector output is the translated affine space
  \[
  \widetilde{\Xi}^{(1)}_{\Lambda_k,-}(u)+\mathcal K^-_{\Lambda_k}
  \]
  solving
  \[
  Q_t\Xi^-_{\Lambda_k}(u)
  =
  \widetilde{\mathcal A}^{(1)}_{\Lambda_k,-}(u).
  \]
\end{enumerate}
Thus the selector output is determined by the canonical odd kernel channel together with the
swap-odd part of the chosen anomaly representative.
\end{corollary}

\begin{proof}
This is an immediate specialization of \cref{thm:selector-decomposition}.
\end{proof}

\begin{remark}[Kernel activity versus representative-driven activity]
\label{rem:selector-two-sources}
By \cref{cor:selector-two-channels}, the selector output is either the canonical odd kernel channel
\(\mathcal K^-_{\Lambda_k}\) or a translate of that channel by a particular solution of
\(
Q_t\Xi^-_{\Lambda_k}(u)=\widetilde{\mathcal A}^{(1)}_{\Lambda_k,-}(u)
\).
\end{remark}

\begin{proposition}[Same-scheme affine translation of selector outputs]
\label{prop:selector-same-scheme}
Let
\[
\widetilde{\mathcal A}^{(1)}_{\Lambda_k}(u),
\qquad
\widehat{\mathcal A}^{(1)}_{\Lambda_k}(u)
\in
\mathcal C_{{\rm sel},\Lambda_k}(u)
\]
be selector representatives in the same finite-cutoff scheme, and suppose there exists an even
sufficiently local cochain \(Y_{\Lambda_k}(u)\) such that
\begin{equation}
\widehat{\mathcal A}^{(1)}_{\Lambda_k}(u)
-
\widetilde{\mathcal A}^{(1)}_{\Lambda_k}(u)
=
Q_tY_{\Lambda_k}(u).
\label{eq:selector-same-scheme}
\end{equation}
Define
\[
Y^-_{\Lambda_k}(u)
:=
\frac12\Big(
Y_{\Lambda_k}(u)-\Gamma Y_{\Lambda_k}(u)
\Big).
\]
Then
\[
\mathfrak X_{\Lambda_k}^{{\rm sel},-}
\big(u;\widehat{\mathcal A}^{(1)}_{\Lambda_k}\big)
=
\mathfrak X_{\Lambda_k}^{{\rm sel},-}
\big(u;\widetilde{\mathcal A}^{(1)}_{\Lambda_k}\big)
+
Y^-_{\Lambda_k}(u).
\]
If \(Y'_{\Lambda_k}(u)\) is another primitive satisfying
\eqref{eq:selector-same-scheme}, then
\[
Y_{\Lambda_k}^{\prime -}(u)-Y^-_{\Lambda_k}(u)\in \mathcal K^-_{\Lambda_k}.
\]
In particular, if
\(
\widehat{\mathcal A}^{(1)}_{\Lambda_k,-}(u)
=
\widetilde{\mathcal A}^{(1)}_{\Lambda_k,-}(u)
\),
then
\[
\mathfrak X_{\Lambda_k}^{{\rm sel},-}
\big(u;\widehat{\mathcal A}^{(1)}_{\Lambda_k}\big)
=
\mathfrak X_{\Lambda_k}^{{\rm sel},-}
\big(u;\widetilde{\mathcal A}^{(1)}_{\Lambda_k}\big).
\]
\end{proposition}

\begin{proof}
Taking swap-odd parts of \eqref{eq:selector-same-scheme} gives
\[
\widehat{\mathcal A}^{(1)}_{\Lambda_k,-}(u)
-
\widetilde{\mathcal A}^{(1)}_{\Lambda_k,-}(u)
=
Q_tY^-_{\Lambda_k}(u).
\]
Let
\(
\Xi^-_{\Lambda_k}(u)
\in
\mathfrak X_{\Lambda_k}^{{\rm sel},-}
\big(u;\widetilde{\mathcal A}^{(1)}_{\Lambda_k}\big)
\).
Then
\[
Q_t\Xi^-_{\Lambda_k}(u)
=
\widetilde{\mathcal A}^{(1)}_{\Lambda_k,-}(u),
\]
so
\[
Q_t\big(
\Xi^-_{\Lambda_k}(u)+Y^-_{\Lambda_k}(u)
\big)
=
\widehat{\mathcal A}^{(1)}_{\Lambda_k,-}(u).
\]
Because both terms are \(\Gamma\)-odd, the sum lies in
\(
\mathfrak X_{\Lambda_k}^{{\rm sel},-}
\big(u;\widehat{\mathcal A}^{(1)}_{\Lambda_k}\big)
\).
This proves one inclusion, and the reverse inclusion follows by symmetry with \(-Y^-_{\Lambda_k}(u)\).

If \(Y'_{\Lambda_k}(u)\) is another primitive, then
\(
Q_t(Y'_{\Lambda_k}(u)-Y_{\Lambda_k}(u))=0
\).
Taking swap-odd parts yields
\[
Q_t\big(
Y_{\Lambda_k}^{\prime -}(u)-Y^-_{\Lambda_k}(u)
\big)=0.
\]
Since the difference is also \(\Gamma\)-odd, it lies in \(\mathcal K^-_{\Lambda_k}\).
If
\(
\widehat{\mathcal A}^{(1)}_{\Lambda_k,-}(u)
=
\widetilde{\mathcal A}^{(1)}_{\Lambda_k,-}(u)
\),
then one may choose \(Y^-_{\Lambda_k}(u)=0\), and the last claim follows.
\end{proof}

\begin{theorem}[Scheme-stable affine translation of the odd selector output]
\label{thm:selector-scheme-stability}
Let \(S\) and \(S'\) be cohomologically admissibly comparable finite-cutoff schemes on the same
chosen monotone sector.
Assume in addition that:
\begin{enumerate}[leftmargin=2em]
\item the canonical anomaly representatives
  \(
  \mathcal A^{(1)}_{k,S}(u)
  \)
  and
  \(
  \mathcal A^{(1)}_{k,S'}(u)
  \)
  are swap-even;
\item there exist chosen \(\Gamma\)-stable selector classes
  \[
  \mathcal C_{{\rm sel},\Lambda_k}^{S}(u),
  \qquad
  \mathcal C_{{\rm sel},\Lambda_k}^{S'}(u),
  \]
  and chosen selector representatives
  \[
  \widehat{\mathcal A}^{(1)}_{k,S}(u)\in\mathcal C_{{\rm sel},\Lambda_k}^{S}(u),
  \qquad
  \widehat{\mathcal A}^{(1)}_{k,S'}(u)\in\mathcal C_{{\rm sel},\Lambda_k}^{S'}(u),
  \]
  and an even sufficiently local cochain \(Y_k(u)\) such that
  \begin{equation}
  \widehat{\mathcal A}^{(1)}_{k,S'}(u)
  -
  \widehat{\mathcal A}^{(1)}_{k,S}(u)
  =
  \mathcal A^{(1)}_{k,S'}(u)
  -
  \mathcal A^{(1)}_{k,S}(u)
  +
  Q_tY_k(u).
  \label{eq:selector-exact-comparison}
  \end{equation}
\end{enumerate}
Define the odd comparison primitive
\[
Y^-_k(u)
:=
\frac12\Big(
Y_k(u)-\Gamma Y_k(u)
\Big).
\]
Then
\[
\mathfrak X_{{\Lambda_k},S'}^{{\rm sel},-}
\big(u;\widehat{\mathcal A}^{(1)}_{k,S'}\big)
=
\mathfrak X_{{\Lambda_k},S}^{{\rm sel},-}
\big(u;\widehat{\mathcal A}^{(1)}_{k,S}\big)
+
Y^-_k(u).
\]
If \(Y'_k(u)\) is another comparison primitive satisfying
\eqref{eq:selector-exact-comparison} for the same chosen pair of representatives, then
\[
Y_k^{\prime -}(u)-Y^-_k(u)\in \mathcal K^-_{\Lambda_k}.
\]
\end{theorem}

\begin{proof}
Taking swap-odd parts of \eqref{eq:selector-exact-comparison} and using the swap-evenness of the
canonical anomaly representatives gives
\[
\widehat{\mathcal A}^{(1)}_{k,S',-}(u)
- 
\widehat{\mathcal A}^{(1)}_{k,S,-}(u)
=
Q_tY^-_k(u).
\]
The same affine-fiber argument as in \cref{prop:selector-same-scheme} therefore gives the
translation statement by \(Y^-_k(u)\), and the same odd-kernel argument gives the ambiguity claim
for alternative primitives.
\end{proof}

\begin{remark}[Selector output handed to the downstream notes]
\label{rem:selector-handoff}
The raw selector datum is the affine class
\[
\mathfrak X_{\Lambda_k}^{{\rm sel},-}
\big(u;\widetilde{\mathcal A}^{(1)}_{\Lambda_k}\big)
\]
attached to the chosen selector representative.
Its ambiguity descent, transported dominant-response value, and weak-interface image are treated in
\cref{thm:dominant-quotient-descent,thm:active-response-transport-minimal,thm:char-weak-interface}.
\end{remark}

\begin{remark}[Relation to the Hodge-normalized layer]
\label{rem:selector-vs-hodge}
The Hodge-normalized comparison layer for admissible alternative representatives is treated in
\cref{uomdesc-thm:uomdesc-conditional-hodge-package}.
The present section uses only raw sufficiently local selector classes.
\end{remark}

\ifdefined\UOMTransportFragmentLoaded\else
\end{document}
\fi
%
}{%
  \IfFileExists{UOM/odd_selector_note.tex}{%
    \ifdefined\UOMTransportFragmentLoaded\else
\documentclass[11pt]{article}
\usepackage[margin=1in]{geometry}
\usepackage{amsmath,amssymb,amsthm,mathtools}
\usepackage{bm}
\usepackage{microtype}
\usepackage{enumitem}
\usepackage{booktabs}
\usepackage{array}
\usepackage{xcolor}
\usepackage{xr-hyper}
\usepackage{hyperref}
\usepackage[nameinlink,noabbrev]{cleveref}

\providecommand{\Fsloc}{\mathcal F_{\rm s.loc}}
\providecommand{\Qt}{Q_t}
\externaldocument[main-]{UOM/main}
\externaldocument[uomdesc-]{UOM/uom_descendant_note}
\externaldocument{UOM/completion_sector_split_model_note}

\hypersetup{
  colorlinks=true,
  linkcolor=blue!60!black,
  citecolor=blue!60!black,
  urlcolor=blue!60!black
}

\theoremstyle{definition}
\newtheorem{definition}{Definition}[section]
\newtheorem{assumption}{Assumption}[section]
\theoremstyle{plain}
\newtheorem{theorem}[definition]{Theorem}
\newtheorem{lemma}[definition]{Lemma}
\newtheorem{proposition}[definition]{Proposition}
\newtheorem{corollary}[definition]{Corollary}
\theoremstyle{remark}
\newtheorem{remark}[definition]{Remark}
\crefname{definition}{definition}{definitions}
\Crefname{definition}{Definition}{Definitions}
\crefname{assumption}{assumption}{assumptions}
\Crefname{assumption}{Assumption}{Assumptions}
\crefname{theorem}{theorem}{theorems}
\Crefname{theorem}{Theorem}{Theorems}
\crefname{lemma}{lemma}{lemmas}
\Crefname{lemma}{Lemma}{Lemmas}
\crefname{proposition}{proposition}{propositions}
\Crefname{proposition}{Proposition}{Propositions}
\crefname{corollary}{corollary}{corollaries}
\Crefname{corollary}{Corollary}{Corollaries}
\crefname{remark}{remark}{remarks}
\Crefname{remark}{Remark}{Remarks}

\newcommand{\cH}{\mathcal H}
\newcommand{\cA}{\mathcal A}
\newcommand{\cD}{\mathcal D}
\newcommand{\cB}{\mathcal B}
\newcommand{\cC}{\mathcal C}
\newcommand{\cL}{\mathcal L}
\newcommand{\cK}{\mathcal K}
\newcommand{\cR}{\mathcal R}
\newcommand{\cG}{\mathcal G}
\newcommand{\Ztwo}{\mathbb Z_2}
\newcommand{\Sys}{\mathrm{Sys}}
\newcommand{\Time}{\mathrm{Time}}
\newcommand{\RG}{\mathrm{RG}}
\newcommand{\Tr}{\mathrm{Tr}}
\newcommand{\diag}{\mathrm{diag}}
\newcommand{\ad}{\mathrm{ad}}
\newcommand{\id}{\mathrm{id}}
\newcommand{\norm}[1]{\left\lVert #1\right\rVert}
\newcommand{\abs}[1]{\left\lvert #1\right\rvert}
\newcommand{\ip}[2]{\left\langle #1,#2\right\rangle}

\externaldocument{UOM/compressed_log_spectral_selection_note}
\externaldocument{UOM/canonical_odd_interface_note}
\externaldocument{UOM/active_response_transport_program_note}
\title{Raw Odd Selector and Scheme-Stable Affine Output}
\author{}
\date{}
\begin{document}
\maketitle
\fi

\section{Raw Odd Selector and Scheme-Stable Affine Output}
\label{sec:odd-selector}

This section isolates the raw selector step on the odd descendant side of the UOM-to-Main
interface used in Definition~\ref{main-def:uom-main-theorem-assumptions}.
At fixed finite cutoff, it records the \(\Gamma\)-odd affine output determined by a chosen
sufficiently local anomaly representative and its same-scheme and inter-scheme translation laws.

\paragraph{Convention on the two \texorpdfstring{\(\mathbb Z_2\)}{Z2} gradings.}
The cohomological degree for the \(Q_t\)-complex is written with bars:
\[
(\cdot)_{\bar 0},
\qquad
(\cdot)_{\bar 1}.
\]
Physical parity refers to the swap involution \(\Gamma\), with eigenspaces
\[
\Gamma F=\pm F.
\]
Thus the selector output constructed below is cohomologically degree \(0\) but physically odd.
Unless bar-degree notation is used explicitly, ``odd'' below means \(\Gamma\)-odd rather than
\(Q_t\)-degree \(1\).

Work at one accepted tick \(k\) and one branch value \(u\) at cutoff \(\Lambda_k\).
Use the raw UOM descendant package of
\cref{uomdesc-thm:uomdesc-raw-package,uomdesc-thm:uomdesc-conditional-odd-package}.
In particular:
\begin{enumerate}[leftmargin=2em]
\item a sufficiently local good monotone representative
  \[
  \mathcal O^{(0)}_{\Lambda_k}(u);
  \]
\item its canonical anomaly representative
  \[
  \mathcal A^{(1)}_{\Lambda_k}(u)
  :=
  Q^{\rm phys}_{u,\Lambda_k}\mathcal O^{(0)}_{\Lambda_k}(u);
  \]
\item a swap involution
  \[
  \Gamma:
  \big(\mathcal F_{\rm s.loc}^{\Lambda_k}(\mathcal X_k)\big)_{\bar\imath}
  \to
  \big(\mathcal F_{\rm s.loc}^{\Lambda_k}(\mathcal X_k)\big)_{\bar\imath}
  \qquad
  (\bar\imath\in\{\bar 0,\bar 1\})
  \]
  preserving sufficient locality and cohomological degree, and commuting with \(Q_t\):
  \[
  \Gamma Q_t=Q_t\Gamma;
  \]
\item swap-evenness of the canonical core,
  \[
  \Gamma\mathcal O^{(0)}_{\Lambda_k}(u)=\mathcal O^{(0)}_{\Lambda_k}(u),
  \qquad
  \Gamma\mathcal A^{(1)}_{\Lambda_k}(u)=\mathcal A^{(1)}_{\Lambda_k}(u);
  \]
\item swap covariance of the fixed-scheme physical radial differential on the chosen monotone
  sector,
  \[
  \Gamma Q^{\rm phys}_{u,\Lambda_k}
  =
  Q^{\rm phys}_{u,\Lambda_k}\Gamma;
  \]
\end{enumerate}
Items (1) and (2) are contained in \cref{uomdesc-thm:uomdesc-raw-package};
item (3) uses \cref{uomdesc-def:uomdesc-gamma,uomdesc-lem:uomdesc-gamma-comm};
item (4) uses \cref{uomdesc-prop:uomdesc-gamma-even-core}; and item (5) uses
\cref{uomdesc-prop:uomdesc-swap-split}.

Fix a chosen \(\Gamma\)-stable selector class of sufficiently local odd representatives of the same
\(Q_t\)-cohomology class as \(\mathcal A^{(1)}_{\Lambda_k}(u)\), and restrict to representatives
for which the exact descent equation
\begin{equation}
Q^{\rm phys}_{u,\Lambda_k}\mathcal O^{(0)}_{\Lambda_k}(u)
+
Q_t\widetilde{\Xi}^{(1)}_{\Lambda_k}(u)
=
\widetilde{\mathcal A}^{(1)}_{\Lambda_k}(u)
\label{eq:selector-full-descent}
\end{equation}
admits at least one sufficiently local even solution
\(
\widetilde{\Xi}^{(1)}_{\Lambda_k}(u)
\in
\big(\mathcal F_{\rm s.loc}^{\Lambda_k}(\mathcal X_k)\big)_{\bar 0}
\).

\begin{definition}[Selector class and canonical odd kernel channel]
\label{def:selector-class}
Let
\[
\mathcal C_{{\rm sel},\Lambda_k}(u)
\subset
\big(\mathcal F_{\rm s.loc}^{\Lambda_k}(\mathcal X_k)\big)_{\bar 1}
\]
be the chosen \(\Gamma\)-stable class of sufficiently local odd representatives of the same
degree-one \(Q_t\)-cohomology class as \(\mathcal A^{(1)}_{\Lambda_k}(u)\) for which
\eqref{eq:selector-full-descent} admits a sufficiently local even solution.
For
\(
\widetilde{\mathcal A}^{(1)}_{\Lambda_k}(u)\in\mathcal C_{{\rm sel},\Lambda_k}(u)
\),
define its swap-even and swap-odd parts by
\[
\widetilde{\mathcal A}^{(1)}_{\Lambda_k,\pm}(u)
:=
\frac12\Big(
\widetilde{\mathcal A}^{(1)}_{\Lambda_k}(u)
\pm
\Gamma\widetilde{\mathcal A}^{(1)}_{\Lambda_k}(u)
\Big).
\]
Define the canonical odd kernel channel by
\[
\mathcal K^-_{\Lambda_k}
:=
\ker(Q_t)
\cap
\big(\mathcal F_{\rm s.loc}^{\Lambda_k}(\mathcal X_k)\big)_{\bar 0}
\cap
\ker(\Gamma+\mathrm{id}).
\]
\end{definition}

\begin{proposition}[Odd defect of a selector representative is always \(Q_t\)-exact]
\label{prop:selector-odd-defect}
Let
\(
\widetilde{\mathcal A}^{(1)}_{\Lambda_k}(u)\in\mathcal C_{{\rm sel},\Lambda_k}(u)
\).
Then:
\begin{enumerate}[leftmargin=2em]
\item
  \(
  \widetilde{\mathcal A}^{(1)}_{\Lambda_k,+}(u)
  \)
  and
  \(
  \widetilde{\mathcal A}^{(1)}_{\Lambda_k,-}(u)
  \)
  are sufficiently local;
\item
  \(
  \widetilde{\mathcal A}^{(1)}_{\Lambda_k,+}(u)
  \)
  is a swap-even representative of the same \(Q_t\)-cohomology class as
  \(\mathcal A^{(1)}_{\Lambda_k}(u)\);
\item
  \[
  \widetilde{\mathcal A}^{(1)}_{\Lambda_k,-}(u)
  \in
  \operatorname{im}(Q_t)
  \cap
  \big(\mathcal F_{\rm s.loc}^{\Lambda_k}(\mathcal X_k)\big)_{\bar 1}
  \cap
  \ker(\Gamma+\mathrm{id}).
  \]
\end{enumerate}
\end{proposition}

\begin{proof}
Sufficient locality is preserved by linear combinations, proving item (1).

Because
\(
\widetilde{\mathcal A}^{(1)}_{\Lambda_k}(u)
\)
and
\(
\mathcal A^{(1)}_{\Lambda_k}(u)
\)
represent the same \(Q_t\)-cohomology class, there exists
\(
Y_{\Lambda_k}(u)\in
\big(\mathcal F_{\rm s.loc}^{\Lambda_k}(\mathcal X_k)\big)_{\bar 0}
\)
such that
\[
\widetilde{\mathcal A}^{(1)}_{\Lambda_k}(u)
-
\mathcal A^{(1)}_{\Lambda_k}(u)
=
Q_tY_{\Lambda_k}(u).
\]
Taking the swap-even part and using
\(
\Gamma Q_t=Q_t\Gamma
\)
and
\(
\Gamma\mathcal A^{(1)}_{\Lambda_k}(u)=\mathcal A^{(1)}_{\Lambda_k}(u)
\),
we obtain
\[
\widetilde{\mathcal A}^{(1)}_{\Lambda_k,+}(u)
-
\mathcal A^{(1)}_{\Lambda_k}(u)
=
Q_tY_{\Lambda_k,+}(u),
\]
so item (2) follows.

Taking the swap-odd part of the same identity gives
\[
\widetilde{\mathcal A}^{(1)}_{\Lambda_k,-}(u)
=
Q_tY_{\Lambda_k,-}(u).
\]
By construction,
\(
\Gamma\widetilde{\mathcal A}^{(1)}_{\Lambda_k,-}(u)
=
-\widetilde{\mathcal A}^{(1)}_{\Lambda_k,-}(u)
\),
which proves item (3).
\end{proof}

\begin{definition}[Selected odd descendant affine class]
\label{def:selector-affine-class}
For
\(
\widetilde{\mathcal A}^{(1)}_{\Lambda_k}(u)\in\mathcal C_{{\rm sel},\Lambda_k}(u)
\),
define the associated raw odd selector output by
\[
\mathfrak X_{\Lambda_k}^{{\rm sel},-}
\big(u;\widetilde{\mathcal A}^{(1)}_{\Lambda_k}\big)
:=
\left\{
\Xi^-_{\Lambda_k}(u)
\in
\big(\mathcal F_{\rm s.loc}^{\Lambda_k}(\mathcal X_k)\big)_{\bar 0}
\;\middle|\;
\Gamma\Xi^-_{\Lambda_k}(u)=-\Xi^-_{\Lambda_k}(u),
\quad
Q_t\Xi^-_{\Lambda_k}(u)
=
\widetilde{\mathcal A}^{(1)}_{\Lambda_k,-}(u)
\right\}.
\]
\end{definition}

\begin{remark}[Dependence only on the swap-odd defect]
\label{rem:selector-odd-defect-only}
By \cref{thm:selector-decomposition}, the affine class
\(
\mathfrak X_{\Lambda_k}^{{\rm sel},-}
\big(u;\widetilde{\mathcal A}^{(1)}_{\Lambda_k}\big)
\)
depends only on the swap-odd defect
\(
\widetilde{\mathcal A}^{(1)}_{\Lambda_k,-}(u)
\).
\end{remark}

\begin{theorem}[Selector decomposition into even representative channel and odd output]
\label{thm:selector-decomposition}
Let
\(
\widetilde{\mathcal A}^{(1)}_{\Lambda_k}(u)\in\mathcal C_{{\rm sel},\Lambda_k}(u)
\),
and let
\(
\widetilde{\Xi}^{(1)}_{\Lambda_k}(u)
\)
be any sufficiently local even solution of \eqref{eq:selector-full-descent}.
Define
\[
\widetilde{\Xi}^{(1)}_{\Lambda_k,\pm}(u)
:=
\frac12\Big(
\widetilde{\Xi}^{(1)}_{\Lambda_k}(u)
\pm
\Gamma\widetilde{\Xi}^{(1)}_{\Lambda_k}(u)
\Big).
\]
Then:
\begin{enumerate}[leftmargin=2em]
\item the swap-even and swap-odd channels decouple as
  \begin{align}
  Q^{\rm phys}_{u,\Lambda_k}\mathcal O^{(0)}_{\Lambda_k}(u)
  +
  Q_t\widetilde{\Xi}^{(1)}_{\Lambda_k,+}(u)
  &=
  \widetilde{\mathcal A}^{(1)}_{\Lambda_k,+}(u),
  \label{eq:selector-even-channel}
  \\
  Q_t\widetilde{\Xi}^{(1)}_{\Lambda_k,-}(u)
  &=
  \widetilde{\mathcal A}^{(1)}_{\Lambda_k,-}(u);
  \label{eq:selector-odd-channel}
  \end{align}
\item the raw odd selector output
  \(
  \mathfrak X_{\Lambda_k}^{{\rm sel},-}
  \big(u;\widetilde{\mathcal A}^{(1)}_{\Lambda_k}\big)
  \)
  contains the odd projection of any full solution; in particular it is nonempty;
\item it is an affine space over the canonical odd kernel channel:
  \[
  \mathfrak X_{\Lambda_k}^{{\rm sel},-}
  \big(u;\widetilde{\mathcal A}^{(1)}_{\Lambda_k}\big)
  =
  \widetilde{\Xi}^{(1)}_{\Lambda_k,-}(u)
  +
  \mathcal K^-_{\Lambda_k}.
  \]
\end{enumerate}
\end{theorem}

\begin{proof}
Apply \(\Gamma\) to \eqref{eq:selector-full-descent}.
Using
\(
\Gamma Q^{\rm phys}_{u,\Lambda_k}
=
Q^{\rm phys}_{u,\Lambda_k}\Gamma
\)
on the chosen monotone sector and
\(
\Gamma\mathcal O^{(0)}_{\Lambda_k}(u)=\mathcal O^{(0)}_{\Lambda_k}(u)
\),
we obtain
\[
Q^{\rm phys}_{u,\Lambda_k}\mathcal O^{(0)}_{\Lambda_k}(u)
+
Q_t\Gamma\widetilde{\Xi}^{(1)}_{\Lambda_k}(u)
=
\Gamma\widetilde{\mathcal A}^{(1)}_{\Lambda_k}(u).
\]
Adding and subtracting this equation from \eqref{eq:selector-full-descent} yields
\eqref{eq:selector-even-channel} and \eqref{eq:selector-odd-channel}, proving item (1).

Item (2) is immediate from item (1): the odd projection
\(
\widetilde{\Xi}^{(1)}_{\Lambda_k,-}(u)
\)
of any full solution belongs to the set in
\cref{def:selector-affine-class}.

Now let
\(
\Xi^-_{\Lambda_k,1}(u)
\)
and
\(
\Xi^-_{\Lambda_k,2}(u)
\)
lie in
\(
\mathfrak X_{\Lambda_k}^{{\rm sel},-}
\big(u;\widetilde{\mathcal A}^{(1)}_{\Lambda_k}\big)
\).
Subtracting their defining equations gives
\[
Q_t\big(
\Xi^-_{\Lambda_k,1}(u)-\Xi^-_{\Lambda_k,2}(u)
\big)
=
0.
\]
Because both are \(\Gamma\)-odd, their difference lies in
\(
\mathcal K^-_{\Lambda_k}
\).
Conversely, if
\(
K^-_{\Lambda_k}\in\mathcal K^-_{\Lambda_k}
\),
then
\[
Q_t\big(
\widetilde{\Xi}^{(1)}_{\Lambda_k,-}(u)+K^-_{\Lambda_k}
\big)
=
\widetilde{\mathcal A}^{(1)}_{\Lambda_k,-}(u),
\]
and the sum remains \(\Gamma\)-odd.
This proves item (3).
\end{proof}

\begin{corollary}[The two selector channels]
\label{cor:selector-two-channels}
At fixed cutoff, the selector problem splits into exactly two channels.
\begin{enumerate}[leftmargin=2em]
\item \textbf{Kernel channel.}
  If
  \(
  \widetilde{\mathcal A}^{(1)}_{\Lambda_k,-}(u)=0
  \),
  then
  \[
  \mathfrak X_{\Lambda_k}^{{\rm sel},-}
  \big(u;\widetilde{\mathcal A}^{(1)}_{\Lambda_k}\big)
  =
  \mathcal K^-_{\Lambda_k}.
  \]
  In particular, the canonical anomaly representative
  \(
  \widetilde{\mathcal A}^{(1)}_{\Lambda_k}(u)=\mathcal A^{(1)}_{\Lambda_k}(u)
  \)
  produces only the canonical odd kernel sector.
\item \textbf{Representative channel.}
  If
  \(
  \widetilde{\mathcal A}^{(1)}_{\Lambda_k,-}(u)\neq 0
  \),
  then the raw odd selector output is the translated affine space
  \[
  \widetilde{\Xi}^{(1)}_{\Lambda_k,-}(u)+\mathcal K^-_{\Lambda_k}
  \]
  solving
  \[
  Q_t\Xi^-_{\Lambda_k}(u)
  =
  \widetilde{\mathcal A}^{(1)}_{\Lambda_k,-}(u).
  \]
\end{enumerate}
Thus the selector output is determined by the canonical odd kernel channel together with the
swap-odd part of the chosen anomaly representative.
\end{corollary}

\begin{proof}
This is an immediate specialization of \cref{thm:selector-decomposition}.
\end{proof}

\begin{remark}[Kernel activity versus representative-driven activity]
\label{rem:selector-two-sources}
By \cref{cor:selector-two-channels}, the selector output is either the canonical odd kernel channel
\(\mathcal K^-_{\Lambda_k}\) or a translate of that channel by a particular solution of
\(
Q_t\Xi^-_{\Lambda_k}(u)=\widetilde{\mathcal A}^{(1)}_{\Lambda_k,-}(u)
\).
\end{remark}

\begin{proposition}[Same-scheme affine translation of selector outputs]
\label{prop:selector-same-scheme}
Let
\[
\widetilde{\mathcal A}^{(1)}_{\Lambda_k}(u),
\qquad
\widehat{\mathcal A}^{(1)}_{\Lambda_k}(u)
\in
\mathcal C_{{\rm sel},\Lambda_k}(u)
\]
be selector representatives in the same finite-cutoff scheme, and suppose there exists an even
sufficiently local cochain \(Y_{\Lambda_k}(u)\) such that
\begin{equation}
\widehat{\mathcal A}^{(1)}_{\Lambda_k}(u)
-
\widetilde{\mathcal A}^{(1)}_{\Lambda_k}(u)
=
Q_tY_{\Lambda_k}(u).
\label{eq:selector-same-scheme}
\end{equation}
Define
\[
Y^-_{\Lambda_k}(u)
:=
\frac12\Big(
Y_{\Lambda_k}(u)-\Gamma Y_{\Lambda_k}(u)
\Big).
\]
Then
\[
\mathfrak X_{\Lambda_k}^{{\rm sel},-}
\big(u;\widehat{\mathcal A}^{(1)}_{\Lambda_k}\big)
=
\mathfrak X_{\Lambda_k}^{{\rm sel},-}
\big(u;\widetilde{\mathcal A}^{(1)}_{\Lambda_k}\big)
+
Y^-_{\Lambda_k}(u).
\]
If \(Y'_{\Lambda_k}(u)\) is another primitive satisfying
\eqref{eq:selector-same-scheme}, then
\[
Y_{\Lambda_k}^{\prime -}(u)-Y^-_{\Lambda_k}(u)\in \mathcal K^-_{\Lambda_k}.
\]
In particular, if
\(
\widehat{\mathcal A}^{(1)}_{\Lambda_k,-}(u)
=
\widetilde{\mathcal A}^{(1)}_{\Lambda_k,-}(u)
\),
then
\[
\mathfrak X_{\Lambda_k}^{{\rm sel},-}
\big(u;\widehat{\mathcal A}^{(1)}_{\Lambda_k}\big)
=
\mathfrak X_{\Lambda_k}^{{\rm sel},-}
\big(u;\widetilde{\mathcal A}^{(1)}_{\Lambda_k}\big).
\]
\end{proposition}

\begin{proof}
Taking swap-odd parts of \eqref{eq:selector-same-scheme} gives
\[
\widehat{\mathcal A}^{(1)}_{\Lambda_k,-}(u)
-
\widetilde{\mathcal A}^{(1)}_{\Lambda_k,-}(u)
=
Q_tY^-_{\Lambda_k}(u).
\]
Let
\(
\Xi^-_{\Lambda_k}(u)
\in
\mathfrak X_{\Lambda_k}^{{\rm sel},-}
\big(u;\widetilde{\mathcal A}^{(1)}_{\Lambda_k}\big)
\).
Then
\[
Q_t\Xi^-_{\Lambda_k}(u)
=
\widetilde{\mathcal A}^{(1)}_{\Lambda_k,-}(u),
\]
so
\[
Q_t\big(
\Xi^-_{\Lambda_k}(u)+Y^-_{\Lambda_k}(u)
\big)
=
\widehat{\mathcal A}^{(1)}_{\Lambda_k,-}(u).
\]
Because both terms are \(\Gamma\)-odd, the sum lies in
\(
\mathfrak X_{\Lambda_k}^{{\rm sel},-}
\big(u;\widehat{\mathcal A}^{(1)}_{\Lambda_k}\big)
\).
This proves one inclusion, and the reverse inclusion follows by symmetry with \(-Y^-_{\Lambda_k}(u)\).

If \(Y'_{\Lambda_k}(u)\) is another primitive, then
\(
Q_t(Y'_{\Lambda_k}(u)-Y_{\Lambda_k}(u))=0
\).
Taking swap-odd parts yields
\[
Q_t\big(
Y_{\Lambda_k}^{\prime -}(u)-Y^-_{\Lambda_k}(u)
\big)=0.
\]
Since the difference is also \(\Gamma\)-odd, it lies in \(\mathcal K^-_{\Lambda_k}\).
If
\(
\widehat{\mathcal A}^{(1)}_{\Lambda_k,-}(u)
=
\widetilde{\mathcal A}^{(1)}_{\Lambda_k,-}(u)
\),
then one may choose \(Y^-_{\Lambda_k}(u)=0\), and the last claim follows.
\end{proof}

\begin{theorem}[Scheme-stable affine translation of the odd selector output]
\label{thm:selector-scheme-stability}
Let \(S\) and \(S'\) be cohomologically admissibly comparable finite-cutoff schemes on the same
chosen monotone sector.
Assume in addition that:
\begin{enumerate}[leftmargin=2em]
\item the canonical anomaly representatives
  \(
  \mathcal A^{(1)}_{k,S}(u)
  \)
  and
  \(
  \mathcal A^{(1)}_{k,S'}(u)
  \)
  are swap-even;
\item there exist chosen \(\Gamma\)-stable selector classes
  \[
  \mathcal C_{{\rm sel},\Lambda_k}^{S}(u),
  \qquad
  \mathcal C_{{\rm sel},\Lambda_k}^{S'}(u),
  \]
  and chosen selector representatives
  \[
  \widehat{\mathcal A}^{(1)}_{k,S}(u)\in\mathcal C_{{\rm sel},\Lambda_k}^{S}(u),
  \qquad
  \widehat{\mathcal A}^{(1)}_{k,S'}(u)\in\mathcal C_{{\rm sel},\Lambda_k}^{S'}(u),
  \]
  and an even sufficiently local cochain \(Y_k(u)\) such that
  \begin{equation}
  \widehat{\mathcal A}^{(1)}_{k,S'}(u)
  -
  \widehat{\mathcal A}^{(1)}_{k,S}(u)
  =
  \mathcal A^{(1)}_{k,S'}(u)
  -
  \mathcal A^{(1)}_{k,S}(u)
  +
  Q_tY_k(u).
  \label{eq:selector-exact-comparison}
  \end{equation}
\end{enumerate}
Define the odd comparison primitive
\[
Y^-_k(u)
:=
\frac12\Big(
Y_k(u)-\Gamma Y_k(u)
\Big).
\]
Then
\[
\mathfrak X_{{\Lambda_k},S'}^{{\rm sel},-}
\big(u;\widehat{\mathcal A}^{(1)}_{k,S'}\big)
=
\mathfrak X_{{\Lambda_k},S}^{{\rm sel},-}
\big(u;\widehat{\mathcal A}^{(1)}_{k,S}\big)
+
Y^-_k(u).
\]
If \(Y'_k(u)\) is another comparison primitive satisfying
\eqref{eq:selector-exact-comparison} for the same chosen pair of representatives, then
\[
Y_k^{\prime -}(u)-Y^-_k(u)\in \mathcal K^-_{\Lambda_k}.
\]
\end{theorem}

\begin{proof}
Taking swap-odd parts of \eqref{eq:selector-exact-comparison} and using the swap-evenness of the
canonical anomaly representatives gives
\[
\widehat{\mathcal A}^{(1)}_{k,S',-}(u)
- 
\widehat{\mathcal A}^{(1)}_{k,S,-}(u)
=
Q_tY^-_k(u).
\]
The same affine-fiber argument as in \cref{prop:selector-same-scheme} therefore gives the
translation statement by \(Y^-_k(u)\), and the same odd-kernel argument gives the ambiguity claim
for alternative primitives.
\end{proof}

\begin{remark}[Selector output handed to the downstream notes]
\label{rem:selector-handoff}
The raw selector datum is the affine class
\[
\mathfrak X_{\Lambda_k}^{{\rm sel},-}
\big(u;\widetilde{\mathcal A}^{(1)}_{\Lambda_k}\big)
\]
attached to the chosen selector representative.
Its ambiguity descent, transported dominant-response value, and weak-interface image are treated in
\cref{thm:dominant-quotient-descent,thm:active-response-transport-minimal,thm:char-weak-interface}.
\end{remark}

\begin{remark}[Relation to the Hodge-normalized layer]
\label{rem:selector-vs-hodge}
The Hodge-normalized comparison layer for admissible alternative representatives is treated in
\cref{uomdesc-thm:uomdesc-conditional-hodge-package}.
The present section uses only raw sufficiently local selector classes.
\end{remark}

\ifdefined\UOMTransportFragmentLoaded\else
\end{document}
\fi
%
  }{%
    \ifdefined\UOMTransportFragmentLoaded\else
\documentclass[11pt]{article}
\usepackage[margin=1in]{geometry}
\usepackage{amsmath,amssymb,amsthm,mathtools}
\usepackage{bm}
\usepackage{microtype}
\usepackage{enumitem}
\usepackage{booktabs}
\usepackage{array}
\usepackage{xcolor}
\usepackage{xr-hyper}
\usepackage{hyperref}
\usepackage[nameinlink,noabbrev]{cleveref}

\providecommand{\Fsloc}{\mathcal F_{\rm s.loc}}
\providecommand{\Qt}{Q_t}
\externaldocument[main-]{UOM/main}
\externaldocument[uomdesc-]{UOM/uom_descendant_note}
\externaldocument{UOM/completion_sector_split_model_note}

\hypersetup{
  colorlinks=true,
  linkcolor=blue!60!black,
  citecolor=blue!60!black,
  urlcolor=blue!60!black
}

\theoremstyle{definition}
\newtheorem{definition}{Definition}[section]
\newtheorem{assumption}{Assumption}[section]
\theoremstyle{plain}
\newtheorem{theorem}[definition]{Theorem}
\newtheorem{lemma}[definition]{Lemma}
\newtheorem{proposition}[definition]{Proposition}
\newtheorem{corollary}[definition]{Corollary}
\theoremstyle{remark}
\newtheorem{remark}[definition]{Remark}
\crefname{definition}{definition}{definitions}
\Crefname{definition}{Definition}{Definitions}
\crefname{assumption}{assumption}{assumptions}
\Crefname{assumption}{Assumption}{Assumptions}
\crefname{theorem}{theorem}{theorems}
\Crefname{theorem}{Theorem}{Theorems}
\crefname{lemma}{lemma}{lemmas}
\Crefname{lemma}{Lemma}{Lemmas}
\crefname{proposition}{proposition}{propositions}
\Crefname{proposition}{Proposition}{Propositions}
\crefname{corollary}{corollary}{corollaries}
\Crefname{corollary}{Corollary}{Corollaries}
\crefname{remark}{remark}{remarks}
\Crefname{remark}{Remark}{Remarks}

\newcommand{\cH}{\mathcal H}
\newcommand{\cA}{\mathcal A}
\newcommand{\cD}{\mathcal D}
\newcommand{\cB}{\mathcal B}
\newcommand{\cC}{\mathcal C}
\newcommand{\cL}{\mathcal L}
\newcommand{\cK}{\mathcal K}
\newcommand{\cR}{\mathcal R}
\newcommand{\cG}{\mathcal G}
\newcommand{\Ztwo}{\mathbb Z_2}
\newcommand{\Sys}{\mathrm{Sys}}
\newcommand{\Time}{\mathrm{Time}}
\newcommand{\RG}{\mathrm{RG}}
\newcommand{\Tr}{\mathrm{Tr}}
\newcommand{\diag}{\mathrm{diag}}
\newcommand{\ad}{\mathrm{ad}}
\newcommand{\id}{\mathrm{id}}
\newcommand{\norm}[1]{\left\lVert #1\right\rVert}
\newcommand{\abs}[1]{\left\lvert #1\right\rvert}
\newcommand{\ip}[2]{\left\langle #1,#2\right\rangle}

\externaldocument{UOM/compressed_log_spectral_selection_note}
\externaldocument{UOM/canonical_odd_interface_note}
\externaldocument{UOM/active_response_transport_program_note}
\title{Raw Odd Selector and Scheme-Stable Affine Output}
\author{}
\date{}
\begin{document}
\maketitle
\fi

\section{Raw Odd Selector and Scheme-Stable Affine Output}
\label{sec:odd-selector}

This section isolates the raw selector step on the odd descendant side of the UOM-to-Main
interface used in Definition~\ref{main-def:uom-main-theorem-assumptions}.
At fixed finite cutoff, it records the \(\Gamma\)-odd affine output determined by a chosen
sufficiently local anomaly representative and its same-scheme and inter-scheme translation laws.

\paragraph{Convention on the two \texorpdfstring{\(\mathbb Z_2\)}{Z2} gradings.}
The cohomological degree for the \(Q_t\)-complex is written with bars:
\[
(\cdot)_{\bar 0},
\qquad
(\cdot)_{\bar 1}.
\]
Physical parity refers to the swap involution \(\Gamma\), with eigenspaces
\[
\Gamma F=\pm F.
\]
Thus the selector output constructed below is cohomologically degree \(0\) but physically odd.
Unless bar-degree notation is used explicitly, ``odd'' below means \(\Gamma\)-odd rather than
\(Q_t\)-degree \(1\).

Work at one accepted tick \(k\) and one branch value \(u\) at cutoff \(\Lambda_k\).
Use the raw UOM descendant package of
\cref{uomdesc-thm:uomdesc-raw-package,uomdesc-thm:uomdesc-conditional-odd-package}.
In particular:
\begin{enumerate}[leftmargin=2em]
\item a sufficiently local good monotone representative
  \[
  \mathcal O^{(0)}_{\Lambda_k}(u);
  \]
\item its canonical anomaly representative
  \[
  \mathcal A^{(1)}_{\Lambda_k}(u)
  :=
  Q^{\rm phys}_{u,\Lambda_k}\mathcal O^{(0)}_{\Lambda_k}(u);
  \]
\item a swap involution
  \[
  \Gamma:
  \big(\mathcal F_{\rm s.loc}^{\Lambda_k}(\mathcal X_k)\big)_{\bar\imath}
  \to
  \big(\mathcal F_{\rm s.loc}^{\Lambda_k}(\mathcal X_k)\big)_{\bar\imath}
  \qquad
  (\bar\imath\in\{\bar 0,\bar 1\})
  \]
  preserving sufficient locality and cohomological degree, and commuting with \(Q_t\):
  \[
  \Gamma Q_t=Q_t\Gamma;
  \]
\item swap-evenness of the canonical core,
  \[
  \Gamma\mathcal O^{(0)}_{\Lambda_k}(u)=\mathcal O^{(0)}_{\Lambda_k}(u),
  \qquad
  \Gamma\mathcal A^{(1)}_{\Lambda_k}(u)=\mathcal A^{(1)}_{\Lambda_k}(u);
  \]
\item swap covariance of the fixed-scheme physical radial differential on the chosen monotone
  sector,
  \[
  \Gamma Q^{\rm phys}_{u,\Lambda_k}
  =
  Q^{\rm phys}_{u,\Lambda_k}\Gamma;
  \]
\end{enumerate}
Items (1) and (2) are contained in \cref{uomdesc-thm:uomdesc-raw-package};
item (3) uses \cref{uomdesc-def:uomdesc-gamma,uomdesc-lem:uomdesc-gamma-comm};
item (4) uses \cref{uomdesc-prop:uomdesc-gamma-even-core}; and item (5) uses
\cref{uomdesc-prop:uomdesc-swap-split}.

Fix a chosen \(\Gamma\)-stable selector class of sufficiently local odd representatives of the same
\(Q_t\)-cohomology class as \(\mathcal A^{(1)}_{\Lambda_k}(u)\), and restrict to representatives
for which the exact descent equation
\begin{equation}
Q^{\rm phys}_{u,\Lambda_k}\mathcal O^{(0)}_{\Lambda_k}(u)
+
Q_t\widetilde{\Xi}^{(1)}_{\Lambda_k}(u)
=
\widetilde{\mathcal A}^{(1)}_{\Lambda_k}(u)
\label{eq:selector-full-descent}
\end{equation}
admits at least one sufficiently local even solution
\(
\widetilde{\Xi}^{(1)}_{\Lambda_k}(u)
\in
\big(\mathcal F_{\rm s.loc}^{\Lambda_k}(\mathcal X_k)\big)_{\bar 0}
\).

\begin{definition}[Selector class and canonical odd kernel channel]
\label{def:selector-class}
Let
\[
\mathcal C_{{\rm sel},\Lambda_k}(u)
\subset
\big(\mathcal F_{\rm s.loc}^{\Lambda_k}(\mathcal X_k)\big)_{\bar 1}
\]
be the chosen \(\Gamma\)-stable class of sufficiently local odd representatives of the same
degree-one \(Q_t\)-cohomology class as \(\mathcal A^{(1)}_{\Lambda_k}(u)\) for which
\eqref{eq:selector-full-descent} admits a sufficiently local even solution.
For
\(
\widetilde{\mathcal A}^{(1)}_{\Lambda_k}(u)\in\mathcal C_{{\rm sel},\Lambda_k}(u)
\),
define its swap-even and swap-odd parts by
\[
\widetilde{\mathcal A}^{(1)}_{\Lambda_k,\pm}(u)
:=
\frac12\Big(
\widetilde{\mathcal A}^{(1)}_{\Lambda_k}(u)
\pm
\Gamma\widetilde{\mathcal A}^{(1)}_{\Lambda_k}(u)
\Big).
\]
Define the canonical odd kernel channel by
\[
\mathcal K^-_{\Lambda_k}
:=
\ker(Q_t)
\cap
\big(\mathcal F_{\rm s.loc}^{\Lambda_k}(\mathcal X_k)\big)_{\bar 0}
\cap
\ker(\Gamma+\mathrm{id}).
\]
\end{definition}

\begin{proposition}[Odd defect of a selector representative is always \(Q_t\)-exact]
\label{prop:selector-odd-defect}
Let
\(
\widetilde{\mathcal A}^{(1)}_{\Lambda_k}(u)\in\mathcal C_{{\rm sel},\Lambda_k}(u)
\).
Then:
\begin{enumerate}[leftmargin=2em]
\item
  \(
  \widetilde{\mathcal A}^{(1)}_{\Lambda_k,+}(u)
  \)
  and
  \(
  \widetilde{\mathcal A}^{(1)}_{\Lambda_k,-}(u)
  \)
  are sufficiently local;
\item
  \(
  \widetilde{\mathcal A}^{(1)}_{\Lambda_k,+}(u)
  \)
  is a swap-even representative of the same \(Q_t\)-cohomology class as
  \(\mathcal A^{(1)}_{\Lambda_k}(u)\);
\item
  \[
  \widetilde{\mathcal A}^{(1)}_{\Lambda_k,-}(u)
  \in
  \operatorname{im}(Q_t)
  \cap
  \big(\mathcal F_{\rm s.loc}^{\Lambda_k}(\mathcal X_k)\big)_{\bar 1}
  \cap
  \ker(\Gamma+\mathrm{id}).
  \]
\end{enumerate}
\end{proposition}

\begin{proof}
Sufficient locality is preserved by linear combinations, proving item (1).

Because
\(
\widetilde{\mathcal A}^{(1)}_{\Lambda_k}(u)
\)
and
\(
\mathcal A^{(1)}_{\Lambda_k}(u)
\)
represent the same \(Q_t\)-cohomology class, there exists
\(
Y_{\Lambda_k}(u)\in
\big(\mathcal F_{\rm s.loc}^{\Lambda_k}(\mathcal X_k)\big)_{\bar 0}
\)
such that
\[
\widetilde{\mathcal A}^{(1)}_{\Lambda_k}(u)
-
\mathcal A^{(1)}_{\Lambda_k}(u)
=
Q_tY_{\Lambda_k}(u).
\]
Taking the swap-even part and using
\(
\Gamma Q_t=Q_t\Gamma
\)
and
\(
\Gamma\mathcal A^{(1)}_{\Lambda_k}(u)=\mathcal A^{(1)}_{\Lambda_k}(u)
\),
we obtain
\[
\widetilde{\mathcal A}^{(1)}_{\Lambda_k,+}(u)
-
\mathcal A^{(1)}_{\Lambda_k}(u)
=
Q_tY_{\Lambda_k,+}(u),
\]
so item (2) follows.

Taking the swap-odd part of the same identity gives
\[
\widetilde{\mathcal A}^{(1)}_{\Lambda_k,-}(u)
=
Q_tY_{\Lambda_k,-}(u).
\]
By construction,
\(
\Gamma\widetilde{\mathcal A}^{(1)}_{\Lambda_k,-}(u)
=
-\widetilde{\mathcal A}^{(1)}_{\Lambda_k,-}(u)
\),
which proves item (3).
\end{proof}

\begin{definition}[Selected odd descendant affine class]
\label{def:selector-affine-class}
For
\(
\widetilde{\mathcal A}^{(1)}_{\Lambda_k}(u)\in\mathcal C_{{\rm sel},\Lambda_k}(u)
\),
define the associated raw odd selector output by
\[
\mathfrak X_{\Lambda_k}^{{\rm sel},-}
\big(u;\widetilde{\mathcal A}^{(1)}_{\Lambda_k}\big)
:=
\left\{
\Xi^-_{\Lambda_k}(u)
\in
\big(\mathcal F_{\rm s.loc}^{\Lambda_k}(\mathcal X_k)\big)_{\bar 0}
\;\middle|\;
\Gamma\Xi^-_{\Lambda_k}(u)=-\Xi^-_{\Lambda_k}(u),
\quad
Q_t\Xi^-_{\Lambda_k}(u)
=
\widetilde{\mathcal A}^{(1)}_{\Lambda_k,-}(u)
\right\}.
\]
\end{definition}

\begin{remark}[Dependence only on the swap-odd defect]
\label{rem:selector-odd-defect-only}
By \cref{thm:selector-decomposition}, the affine class
\(
\mathfrak X_{\Lambda_k}^{{\rm sel},-}
\big(u;\widetilde{\mathcal A}^{(1)}_{\Lambda_k}\big)
\)
depends only on the swap-odd defect
\(
\widetilde{\mathcal A}^{(1)}_{\Lambda_k,-}(u)
\).
\end{remark}

\begin{theorem}[Selector decomposition into even representative channel and odd output]
\label{thm:selector-decomposition}
Let
\(
\widetilde{\mathcal A}^{(1)}_{\Lambda_k}(u)\in\mathcal C_{{\rm sel},\Lambda_k}(u)
\),
and let
\(
\widetilde{\Xi}^{(1)}_{\Lambda_k}(u)
\)
be any sufficiently local even solution of \eqref{eq:selector-full-descent}.
Define
\[
\widetilde{\Xi}^{(1)}_{\Lambda_k,\pm}(u)
:=
\frac12\Big(
\widetilde{\Xi}^{(1)}_{\Lambda_k}(u)
\pm
\Gamma\widetilde{\Xi}^{(1)}_{\Lambda_k}(u)
\Big).
\]
Then:
\begin{enumerate}[leftmargin=2em]
\item the swap-even and swap-odd channels decouple as
  \begin{align}
  Q^{\rm phys}_{u,\Lambda_k}\mathcal O^{(0)}_{\Lambda_k}(u)
  +
  Q_t\widetilde{\Xi}^{(1)}_{\Lambda_k,+}(u)
  &=
  \widetilde{\mathcal A}^{(1)}_{\Lambda_k,+}(u),
  \label{eq:selector-even-channel}
  \\
  Q_t\widetilde{\Xi}^{(1)}_{\Lambda_k,-}(u)
  &=
  \widetilde{\mathcal A}^{(1)}_{\Lambda_k,-}(u);
  \label{eq:selector-odd-channel}
  \end{align}
\item the raw odd selector output
  \(
  \mathfrak X_{\Lambda_k}^{{\rm sel},-}
  \big(u;\widetilde{\mathcal A}^{(1)}_{\Lambda_k}\big)
  \)
  contains the odd projection of any full solution; in particular it is nonempty;
\item it is an affine space over the canonical odd kernel channel:
  \[
  \mathfrak X_{\Lambda_k}^{{\rm sel},-}
  \big(u;\widetilde{\mathcal A}^{(1)}_{\Lambda_k}\big)
  =
  \widetilde{\Xi}^{(1)}_{\Lambda_k,-}(u)
  +
  \mathcal K^-_{\Lambda_k}.
  \]
\end{enumerate}
\end{theorem}

\begin{proof}
Apply \(\Gamma\) to \eqref{eq:selector-full-descent}.
Using
\(
\Gamma Q^{\rm phys}_{u,\Lambda_k}
=
Q^{\rm phys}_{u,\Lambda_k}\Gamma
\)
on the chosen monotone sector and
\(
\Gamma\mathcal O^{(0)}_{\Lambda_k}(u)=\mathcal O^{(0)}_{\Lambda_k}(u)
\),
we obtain
\[
Q^{\rm phys}_{u,\Lambda_k}\mathcal O^{(0)}_{\Lambda_k}(u)
+
Q_t\Gamma\widetilde{\Xi}^{(1)}_{\Lambda_k}(u)
=
\Gamma\widetilde{\mathcal A}^{(1)}_{\Lambda_k}(u).
\]
Adding and subtracting this equation from \eqref{eq:selector-full-descent} yields
\eqref{eq:selector-even-channel} and \eqref{eq:selector-odd-channel}, proving item (1).

Item (2) is immediate from item (1): the odd projection
\(
\widetilde{\Xi}^{(1)}_{\Lambda_k,-}(u)
\)
of any full solution belongs to the set in
\cref{def:selector-affine-class}.

Now let
\(
\Xi^-_{\Lambda_k,1}(u)
\)
and
\(
\Xi^-_{\Lambda_k,2}(u)
\)
lie in
\(
\mathfrak X_{\Lambda_k}^{{\rm sel},-}
\big(u;\widetilde{\mathcal A}^{(1)}_{\Lambda_k}\big)
\).
Subtracting their defining equations gives
\[
Q_t\big(
\Xi^-_{\Lambda_k,1}(u)-\Xi^-_{\Lambda_k,2}(u)
\big)
=
0.
\]
Because both are \(\Gamma\)-odd, their difference lies in
\(
\mathcal K^-_{\Lambda_k}
\).
Conversely, if
\(
K^-_{\Lambda_k}\in\mathcal K^-_{\Lambda_k}
\),
then
\[
Q_t\big(
\widetilde{\Xi}^{(1)}_{\Lambda_k,-}(u)+K^-_{\Lambda_k}
\big)
=
\widetilde{\mathcal A}^{(1)}_{\Lambda_k,-}(u),
\]
and the sum remains \(\Gamma\)-odd.
This proves item (3).
\end{proof}

\begin{corollary}[The two selector channels]
\label{cor:selector-two-channels}
At fixed cutoff, the selector problem splits into exactly two channels.
\begin{enumerate}[leftmargin=2em]
\item \textbf{Kernel channel.}
  If
  \(
  \widetilde{\mathcal A}^{(1)}_{\Lambda_k,-}(u)=0
  \),
  then
  \[
  \mathfrak X_{\Lambda_k}^{{\rm sel},-}
  \big(u;\widetilde{\mathcal A}^{(1)}_{\Lambda_k}\big)
  =
  \mathcal K^-_{\Lambda_k}.
  \]
  In particular, the canonical anomaly representative
  \(
  \widetilde{\mathcal A}^{(1)}_{\Lambda_k}(u)=\mathcal A^{(1)}_{\Lambda_k}(u)
  \)
  produces only the canonical odd kernel sector.
\item \textbf{Representative channel.}
  If
  \(
  \widetilde{\mathcal A}^{(1)}_{\Lambda_k,-}(u)\neq 0
  \),
  then the raw odd selector output is the translated affine space
  \[
  \widetilde{\Xi}^{(1)}_{\Lambda_k,-}(u)+\mathcal K^-_{\Lambda_k}
  \]
  solving
  \[
  Q_t\Xi^-_{\Lambda_k}(u)
  =
  \widetilde{\mathcal A}^{(1)}_{\Lambda_k,-}(u).
  \]
\end{enumerate}
Thus the selector output is determined by the canonical odd kernel channel together with the
swap-odd part of the chosen anomaly representative.
\end{corollary}

\begin{proof}
This is an immediate specialization of \cref{thm:selector-decomposition}.
\end{proof}

\begin{remark}[Kernel activity versus representative-driven activity]
\label{rem:selector-two-sources}
By \cref{cor:selector-two-channels}, the selector output is either the canonical odd kernel channel
\(\mathcal K^-_{\Lambda_k}\) or a translate of that channel by a particular solution of
\(
Q_t\Xi^-_{\Lambda_k}(u)=\widetilde{\mathcal A}^{(1)}_{\Lambda_k,-}(u)
\).
\end{remark}

\begin{proposition}[Same-scheme affine translation of selector outputs]
\label{prop:selector-same-scheme}
Let
\[
\widetilde{\mathcal A}^{(1)}_{\Lambda_k}(u),
\qquad
\widehat{\mathcal A}^{(1)}_{\Lambda_k}(u)
\in
\mathcal C_{{\rm sel},\Lambda_k}(u)
\]
be selector representatives in the same finite-cutoff scheme, and suppose there exists an even
sufficiently local cochain \(Y_{\Lambda_k}(u)\) such that
\begin{equation}
\widehat{\mathcal A}^{(1)}_{\Lambda_k}(u)
-
\widetilde{\mathcal A}^{(1)}_{\Lambda_k}(u)
=
Q_tY_{\Lambda_k}(u).
\label{eq:selector-same-scheme}
\end{equation}
Define
\[
Y^-_{\Lambda_k}(u)
:=
\frac12\Big(
Y_{\Lambda_k}(u)-\Gamma Y_{\Lambda_k}(u)
\Big).
\]
Then
\[
\mathfrak X_{\Lambda_k}^{{\rm sel},-}
\big(u;\widehat{\mathcal A}^{(1)}_{\Lambda_k}\big)
=
\mathfrak X_{\Lambda_k}^{{\rm sel},-}
\big(u;\widetilde{\mathcal A}^{(1)}_{\Lambda_k}\big)
+
Y^-_{\Lambda_k}(u).
\]
If \(Y'_{\Lambda_k}(u)\) is another primitive satisfying
\eqref{eq:selector-same-scheme}, then
\[
Y_{\Lambda_k}^{\prime -}(u)-Y^-_{\Lambda_k}(u)\in \mathcal K^-_{\Lambda_k}.
\]
In particular, if
\(
\widehat{\mathcal A}^{(1)}_{\Lambda_k,-}(u)
=
\widetilde{\mathcal A}^{(1)}_{\Lambda_k,-}(u)
\),
then
\[
\mathfrak X_{\Lambda_k}^{{\rm sel},-}
\big(u;\widehat{\mathcal A}^{(1)}_{\Lambda_k}\big)
=
\mathfrak X_{\Lambda_k}^{{\rm sel},-}
\big(u;\widetilde{\mathcal A}^{(1)}_{\Lambda_k}\big).
\]
\end{proposition}

\begin{proof}
Taking swap-odd parts of \eqref{eq:selector-same-scheme} gives
\[
\widehat{\mathcal A}^{(1)}_{\Lambda_k,-}(u)
-
\widetilde{\mathcal A}^{(1)}_{\Lambda_k,-}(u)
=
Q_tY^-_{\Lambda_k}(u).
\]
Let
\(
\Xi^-_{\Lambda_k}(u)
\in
\mathfrak X_{\Lambda_k}^{{\rm sel},-}
\big(u;\widetilde{\mathcal A}^{(1)}_{\Lambda_k}\big)
\).
Then
\[
Q_t\Xi^-_{\Lambda_k}(u)
=
\widetilde{\mathcal A}^{(1)}_{\Lambda_k,-}(u),
\]
so
\[
Q_t\big(
\Xi^-_{\Lambda_k}(u)+Y^-_{\Lambda_k}(u)
\big)
=
\widehat{\mathcal A}^{(1)}_{\Lambda_k,-}(u).
\]
Because both terms are \(\Gamma\)-odd, the sum lies in
\(
\mathfrak X_{\Lambda_k}^{{\rm sel},-}
\big(u;\widehat{\mathcal A}^{(1)}_{\Lambda_k}\big)
\).
This proves one inclusion, and the reverse inclusion follows by symmetry with \(-Y^-_{\Lambda_k}(u)\).

If \(Y'_{\Lambda_k}(u)\) is another primitive, then
\(
Q_t(Y'_{\Lambda_k}(u)-Y_{\Lambda_k}(u))=0
\).
Taking swap-odd parts yields
\[
Q_t\big(
Y_{\Lambda_k}^{\prime -}(u)-Y^-_{\Lambda_k}(u)
\big)=0.
\]
Since the difference is also \(\Gamma\)-odd, it lies in \(\mathcal K^-_{\Lambda_k}\).
If
\(
\widehat{\mathcal A}^{(1)}_{\Lambda_k,-}(u)
=
\widetilde{\mathcal A}^{(1)}_{\Lambda_k,-}(u)
\),
then one may choose \(Y^-_{\Lambda_k}(u)=0\), and the last claim follows.
\end{proof}

\begin{theorem}[Scheme-stable affine translation of the odd selector output]
\label{thm:selector-scheme-stability}
Let \(S\) and \(S'\) be cohomologically admissibly comparable finite-cutoff schemes on the same
chosen monotone sector.
Assume in addition that:
\begin{enumerate}[leftmargin=2em]
\item the canonical anomaly representatives
  \(
  \mathcal A^{(1)}_{k,S}(u)
  \)
  and
  \(
  \mathcal A^{(1)}_{k,S'}(u)
  \)
  are swap-even;
\item there exist chosen \(\Gamma\)-stable selector classes
  \[
  \mathcal C_{{\rm sel},\Lambda_k}^{S}(u),
  \qquad
  \mathcal C_{{\rm sel},\Lambda_k}^{S'}(u),
  \]
  and chosen selector representatives
  \[
  \widehat{\mathcal A}^{(1)}_{k,S}(u)\in\mathcal C_{{\rm sel},\Lambda_k}^{S}(u),
  \qquad
  \widehat{\mathcal A}^{(1)}_{k,S'}(u)\in\mathcal C_{{\rm sel},\Lambda_k}^{S'}(u),
  \]
  and an even sufficiently local cochain \(Y_k(u)\) such that
  \begin{equation}
  \widehat{\mathcal A}^{(1)}_{k,S'}(u)
  -
  \widehat{\mathcal A}^{(1)}_{k,S}(u)
  =
  \mathcal A^{(1)}_{k,S'}(u)
  -
  \mathcal A^{(1)}_{k,S}(u)
  +
  Q_tY_k(u).
  \label{eq:selector-exact-comparison}
  \end{equation}
\end{enumerate}
Define the odd comparison primitive
\[
Y^-_k(u)
:=
\frac12\Big(
Y_k(u)-\Gamma Y_k(u)
\Big).
\]
Then
\[
\mathfrak X_{{\Lambda_k},S'}^{{\rm sel},-}
\big(u;\widehat{\mathcal A}^{(1)}_{k,S'}\big)
=
\mathfrak X_{{\Lambda_k},S}^{{\rm sel},-}
\big(u;\widehat{\mathcal A}^{(1)}_{k,S}\big)
+
Y^-_k(u).
\]
If \(Y'_k(u)\) is another comparison primitive satisfying
\eqref{eq:selector-exact-comparison} for the same chosen pair of representatives, then
\[
Y_k^{\prime -}(u)-Y^-_k(u)\in \mathcal K^-_{\Lambda_k}.
\]
\end{theorem}

\begin{proof}
Taking swap-odd parts of \eqref{eq:selector-exact-comparison} and using the swap-evenness of the
canonical anomaly representatives gives
\[
\widehat{\mathcal A}^{(1)}_{k,S',-}(u)
- 
\widehat{\mathcal A}^{(1)}_{k,S,-}(u)
=
Q_tY^-_k(u).
\]
The same affine-fiber argument as in \cref{prop:selector-same-scheme} therefore gives the
translation statement by \(Y^-_k(u)\), and the same odd-kernel argument gives the ambiguity claim
for alternative primitives.
\end{proof}

\begin{remark}[Selector output handed to the downstream notes]
\label{rem:selector-handoff}
The raw selector datum is the affine class
\[
\mathfrak X_{\Lambda_k}^{{\rm sel},-}
\big(u;\widetilde{\mathcal A}^{(1)}_{\Lambda_k}\big)
\]
attached to the chosen selector representative.
Its ambiguity descent, transported dominant-response value, and weak-interface image are treated in
\cref{thm:dominant-quotient-descent,thm:active-response-transport-minimal,thm:char-weak-interface}.
\end{remark}

\begin{remark}[Relation to the Hodge-normalized layer]
\label{rem:selector-vs-hodge}
The Hodge-normalized comparison layer for admissible alternative representatives is treated in
\cref{uomdesc-thm:uomdesc-conditional-hodge-package}.
The present section uses only raw sufficiently local selector classes.
\end{remark}

\ifdefined\UOMTransportFragmentLoaded\else
\end{document}
\fi
%
  }%
}

\IfFileExists{canonical_odd_interface_note.tex}{%
  \ifdefined\UOMTransportFragmentLoaded\else
\documentclass[11pt]{article}
\usepackage[margin=1in]{geometry}
\usepackage{amsmath,amssymb,amsthm,mathtools}
\usepackage{bm}
\usepackage{microtype}
\usepackage{enumitem}
\usepackage{booktabs}
\usepackage{array}
\usepackage{xcolor}
\usepackage{xr-hyper}
\usepackage{hyperref}
\usepackage[nameinlink,noabbrev]{cleveref}

\providecommand{\Fsloc}{\mathcal F_{\rm s.loc}}
\providecommand{\Qt}{Q_t}
\externaldocument[main-]{UOM/main}
\externaldocument[uomdesc-]{UOM/uom_descendant_note}
\externaldocument{UOM/completion_sector_split_model_note}

\hypersetup{
  colorlinks=true,
  linkcolor=blue!60!black,
  citecolor=blue!60!black,
  urlcolor=blue!60!black
}

\theoremstyle{definition}
\newtheorem{definition}{Definition}[section]
\newtheorem{assumption}{Assumption}[section]
\theoremstyle{plain}
\newtheorem{theorem}[definition]{Theorem}
\newtheorem{lemma}[definition]{Lemma}
\newtheorem{proposition}[definition]{Proposition}
\newtheorem{corollary}[definition]{Corollary}
\theoremstyle{remark}
\newtheorem{remark}[definition]{Remark}
\crefname{definition}{definition}{definitions}
\Crefname{definition}{Definition}{Definitions}
\crefname{assumption}{assumption}{assumptions}
\Crefname{assumption}{Assumption}{Assumptions}
\crefname{theorem}{theorem}{theorems}
\Crefname{theorem}{Theorem}{Theorems}
\crefname{lemma}{lemma}{lemmas}
\Crefname{lemma}{Lemma}{Lemmas}
\crefname{proposition}{proposition}{propositions}
\Crefname{proposition}{Proposition}{Propositions}
\crefname{corollary}{corollary}{corollaries}
\Crefname{corollary}{Corollary}{Corollaries}
\crefname{remark}{remark}{remarks}
\Crefname{remark}{Remark}{Remarks}

\newcommand{\cH}{\mathcal H}
\newcommand{\cA}{\mathcal A}
\newcommand{\cD}{\mathcal D}
\newcommand{\cB}{\mathcal B}
\newcommand{\cC}{\mathcal C}
\newcommand{\cL}{\mathcal L}
\newcommand{\cK}{\mathcal K}
\newcommand{\cR}{\mathcal R}
\newcommand{\cG}{\mathcal G}
\newcommand{\Ztwo}{\mathbb Z_2}
\newcommand{\Sys}{\mathrm{Sys}}
\newcommand{\Time}{\mathrm{Time}}
\newcommand{\RG}{\mathrm{RG}}
\newcommand{\Tr}{\mathrm{Tr}}
\newcommand{\diag}{\mathrm{diag}}
\newcommand{\ad}{\mathrm{ad}}
\newcommand{\id}{\mathrm{id}}
\newcommand{\norm}[1]{\left\lVert #1\right\rVert}
\newcommand{\abs}[1]{\left\lvert #1\right\rvert}
\newcommand{\ip}[2]{\left\langle #1,#2\right\rangle}

\externaldocument{UOM/compressed_log_spectral_selection_note}
\externaldocument{UOM/odd_selector_note}
\externaldocument{UOM/active_response_transport_program_note}
\title{Polarized Characterization Interface}
\author{}
\date{}
\begin{document}
\maketitle
\fi

\section{Polarized Characterization Interface}
\label{sec:char-export}

This section records the finite-cutoff characterization data extracted from the split/leak
construction and used in the 2T Main assumption package
Definition~\ref{main-def:uom-main-theorem-assumptions}. In this note, an ``accepted tick'' denotes the local
accepted-branch slice at cutoff \(\Lambda_k\).

For each accepted tick \(k\), write
\[
a_{\Lambda_k}:=a_k,
\qquad
S_{\Lambda_k}:=a_{\Lambda_k}^2.
\]
The variable \(a_{\Lambda_k}\) is the distinguished odd split coordinate identified in
\cref{subsec:leak-solution,subsec:B-gaussian}; the even quantity \(S_{\Lambda_k}\) is the
positive stock that enters the leak law and the split threshold.

\begin{definition}[Characterization algebra, operational triviality, and even null ideal]
\label{def:char-data}
Fix an accepted tick \(k\) and let \(E^-_{\Lambda_k}\) denote the physically odd sufficiently local
generator space at cutoff \(\Lambda_k\), i.e.\ the odd sector for the transported physical involution
\(\Gamma\). Let \(E^+_{\Lambda_k}\) denote the corresponding physically even sufficiently local
observable space.
Define the core characterization algebra by
\begin{equation}
\mathcal O^{\mathrm{char}}_{\Lambda_k}
:=
\mathrm{alg}\{a_{\Lambda_k},\,S_{\Lambda_k}\}.
\label{eq:Ochar-def}
\end{equation}
Define the operationally interface-trivial odd subspace by
\begin{equation}
\mathcal T^{\mathrm{op}}_{\Lambda_k}
:=
\overline{\left\{
T\in E^-_{\Lambda_k}:
\{T,O\}_{\mathrm P,\Lambda_k}=0
\ \forall O\in\mathcal O^{\mathrm{char}}_{\Lambda_k}
\right\}}.
\label{eq:Tchar-def}
\end{equation}
Define the even interface-null ideal by
\begin{equation}
\mathcal N^{\mathrm{even}}_{\Lambda_k}
:=
\overline{\left\{
E\in E^+_{\Lambda_k}:
\{E,O\}_{\mathrm P,\Lambda_k}=0
\ \forall O\in\mathcal O^{\mathrm{char}}_{\Lambda_k}
\right\}}.
\label{eq:Neven-def}
\end{equation}
\end{definition}

\begin{remark}[Do not identify the characterization observable with the generator]
\label{rem:a-vs-Jact}
The odd observable \(a_{\Lambda_k}\) characterizes the active completion-sector channel, but it is
not itself the odd Hamiltonian generator class appearing in item \textbf{(M4)} of
Definition~\ref{main-def:uom-main-theorem-assumptions}.
That generator is denoted \(J^{\mathrm{act}}_{\Lambda_k}\in E^-_{\Lambda_k}\) below.
The quotient statement belongs to \([J^{\mathrm{act}}_{\Lambda_k}]\), not to
\([a_{\Lambda_k}]\).
\end{remark}

\begin{definition}[Intrinsic covariance-tangent response map]
\label{def:intrinsic-cov-response}
Fix an accepted tick \(k\). Let \(\mathcal T^{\mathrm{cov}}_{\Lambda_k}\) denote the finite-cutoff
receiver covariance-tangent space at the UOM reference Gaussian state
\(\rho_{\ast,\Lambda_k}\), equipped with the Fisher pairing
\[
\big\langle\!\big\langle \delta C,\delta C'\big\rangle\!\big\rangle_{C_{\ast,\Lambda_k}}
:=
\tfrac12\,\mathrm{tr}\!\big(
C_{\ast,\Lambda_k}^{-1}\delta C\,C_{\ast,\Lambda_k}^{-1}\delta C'
\big).
\]
For admissible band-limited receiver probes \(f,g\), let
\[
\mathcal Q^{\Lambda_k}_{f,g}\in E^+_{\Lambda_k}
\]
denote the corresponding sufficiently local quadratic observable, normalized so that on centered
Gaussian states with covariance \(C\),
\[
\big\langle \mathcal Q^{\Lambda_k}_{f,g}\big\rangle_{\rho(C)}
=
\langle f,Cg\rangle.
\]
For \(F\in E^-_{\Lambda_k}\), define \(\mathcal R_{\Lambda_k}(F)\in\mathcal T^{\mathrm{cov}}_{\Lambda_k}\)
by the matrix elements
\[
\big\langle f,\mathcal R_{\Lambda_k}(F)g\big\rangle
:=
\Big\langle
\{F,\mathcal Q^{\Lambda_k}_{f,g}\}_{\mathrm P,\Lambda_k}
\Big\rangle_{\rho_{\ast,\Lambda_k}}.
\]
We call
\[
\mathcal R_{\Lambda_k}:E^-_{\Lambda_k}\to\mathcal T^{\mathrm{cov}}_{\Lambda_k}
\]
the \emph{intrinsic covariance-tangent response map}.
\end{definition}

\begin{proposition}[Intrinsic definition and completion independence]
\label{prop:intrinsic-cov-response}
The map \(\mathcal R_{\Lambda_k}\) of \cref{def:intrinsic-cov-response} is continuous and linear.
It depends only on the finite-cutoff UOM Peierls package and the UOM covariance/Fisher package.
If a Gaussian completion realization is chosen, then \(\mathcal R_{\Lambda_k}(F)\) is represented by
the same completion-sector covariance tangent that one obtains by linearizing the completion
quadratures. Thus Gaussian completion is a realization of the response map, not extra structure
required to define it.
\end{proposition}

\begin{proof}
For fixed \(F\in E^-_{\Lambda_k}\), the finite-cutoff UOM Peierls package makes
\[
\delta_F:=\{F,-\}_{\mathrm P,\Lambda_k}
\]
a continuous derivation on the sufficiently local algebra. Applying \(\delta_F\) to the quadratic
probes \(\mathcal Q^{\Lambda_k}_{f,g}\) and evaluating in the reference Gaussian state produces a
continuous symmetric bilinear form in \(f,g\). In the UOM Gaussian radiative sector, such bilinear
forms are exactly covariance tangents at \(C_{\ast,\Lambda_k}\), equipped with the Fisher metric.
Linearity and continuity in \(F\) are inherited directly from the Peierls bracket. No completion
coordinates are used in this construction.

If a Gaussian completion is later chosen, the same second moments are encoded by the covariance
matrix of the completion quadratures. The bilinear form defined above and the completion-sector
covariance tangent therefore represent the same infinitesimal response object in two different
coordinate systems.
\end{proof}

\begin{definition}[Physical intrinsic covariance-response map]
\label{def:phys-cov-response}
Fix an accepted tick \(k\). Let
\[
\mathcal T^{\mathrm{phys}}_{\Lambda_k}
:=
\Pi_{\mathrm{phys},\Lambda_k}\mathcal T^{\mathrm{cov}}_{\Lambda_k}
\subset \mathcal T^{\mathrm{cov}}_{\Lambda_k}.
\]
Define the physical intrinsic covariance-response map by
\[
R^{\mathrm{phys}}_{\Lambda_k}
:=
\Pi_{\mathrm{phys},\Lambda_k}\mathcal R_{\Lambda_k}
:
E^-_{\Lambda_k}\to \mathcal T^{\mathrm{phys}}_{\Lambda_k}.
\]
\end{definition}

\begin{proposition}[Physical intrinsic response map]
\label{prop:phys-cov-response}
The map \(R^{\mathrm{phys}}_{\Lambda_k}\) of \cref{def:phys-cov-response} is continuous and linear.
\end{proposition}

\begin{proof}
This is immediate from \cref{prop:intrinsic-cov-response} and the continuity and linearity of the
finite-cutoff physical projector \(\Pi_{\mathrm{phys},\Lambda_k}\).
\end{proof}

\begin{proposition}[Reduction of operational triviality to the active odd coordinate]
\label{prop:char-op-from-a}
Fix an accepted tick \(k\). If \(F\in E^-_{\Lambda_k}\) satisfies
\[
\{F,a_{\Lambda_k}\}_{\mathrm P,\Lambda_k}=0,
\]
then \(F\in\mathcal T^{\mathrm{op}}_{\Lambda_k}\).
\end{proposition}

\begin{proof}
Since \(S_{\Lambda_k}=a_{\Lambda_k}^2\) and the Peierls bracket is a derivation,
\[
\{F,S_{\Lambda_k}\}_{\mathrm P,\Lambda_k}
=
\{F,a_{\Lambda_k}^2\}_{\mathrm P,\Lambda_k}
=
2\,a_{\Lambda_k}\,\{F,a_{\Lambda_k}\}_{\mathrm P,\Lambda_k}
=
0.
\]
The algebra \(\mathcal O^{\mathrm{char}}_{\Lambda_k}=\mathrm{alg}\{a_{\Lambda_k},S_{\Lambda_k}\}\)
is generated by \(a_{\Lambda_k}\) and \(S_{\Lambda_k}\), and the Peierls bracket acts on it by the
Leibniz rule. Hence \(\{F,O\}_{\mathrm P,\Lambda_k}=0\) for every
\(O\in\mathcal O^{\mathrm{char}}_{\Lambda_k}\), proving
\(F\in\mathcal T^{\mathrm{op}}_{\Lambda_k}\).
\end{proof}

\begin{assumption}[Active odd generator and intrinsic-response factorization]
\label{ass:response-factorization}
For each accepted tick \(k\), there exist a distinguished active odd generator
\[
J^{\mathrm{act}}_{\Lambda_k}\in E^-_{\Lambda_k}
\]
and a nonzero physical covariance tangent
\[
C^{\mathrm{act}}_{\Lambda_k}\in \mathcal T^{\mathrm{phys}}_{\Lambda_k}
\]
and a continuous nonzero linear functional
\[
\rho_{\Lambda_k}:E^-_{\Lambda_k}\to\mathbb R
\]
normalized by
\[
\rho_{\Lambda_k}(J^{\mathrm{act}}_{\Lambda_k})=1,
\]
such that
\[
R^{\mathrm{phys}}_{\Lambda_k}(J^{\mathrm{act}}_{\Lambda_k})
=
C^{\mathrm{act}}_{\Lambda_k}
\neq 0
\]
and
\begin{equation}
R^{\mathrm{phys}}_{\Lambda_k}(F)
=
\rho_{\Lambda_k}(F)\,C^{\mathrm{act}}_{\Lambda_k}
\qquad
\forall F\in E^-_{\Lambda_k}
\label{eq:rho-factorization}
\end{equation}
holds on the strong sector under consideration.
\end{assumption}

\begin{remark}[Strong normalization]
\label{rem:rho-strong}
If one further chooses the Fisher-unit normalization
\(
\widehat C^{\mathrm{act}}_{\Lambda_k}
:=
\|C^{\mathrm{act}}_{\Lambda_k}\|^{-1}_{C_{\ast,\Lambda_k}}\,C^{\mathrm{act}}_{\Lambda_k},
\)
then
\eqref{eq:rho-factorization} sharpens to
\[
R^{\mathrm{phys}}_{\Lambda_k}(F)
=
\widehat{\rho}_{\Lambda_k}(F)\,\widehat C^{\mathrm{act}}_{\Lambda_k}
\]
with
\(
\widehat{\rho}_{\Lambda_k}
:=
\|C^{\mathrm{act}}_{\Lambda_k}\|_{C_{\ast,\Lambda_k}}\,\rho_{\Lambda_k}.
\)
The export below does not require this extra normalization; it only needs the scalar response
coefficient \(\rho_{\Lambda_k}\) relative to a fixed nonzero active tangent
\(C^{\mathrm{act}}_{\Lambda_k}\).
\end{remark}

\begin{assumption}[Physical-response sufficiency]
\label{ass:response-sufficiency}
For each accepted tick \(k\), the operationally trivial odd directions are exactly the kernel of the
physical intrinsic covariance-response map:
\[
 \ker R^{\mathrm{phys}}_{\Lambda_k}
 =
 \mathcal T^{\mathrm{op}}_{\Lambda_k}.
\]
Equivalently, an odd generator has zero physical intrinsic covariance response if and only if it is
operationally trivial on the full core characterization algebra \(\mathcal O^{\mathrm{char}}_{\Lambda_k}\).
\end{assumption}

\begin{remark}[Why the strong package must be stated intrinsically]
\label{rem:response-quotient-obstruction}
The earlier quotient-valued formulation of the strong package in
\(
E^+_{\Lambda_k}/\mathcal N^{\mathrm{even}}_{\Lambda_k}
\)
is not stable.
Because the Peierls bracket here is the ordinary derivation bracket and \(S_{\Lambda_k}=a_{\Lambda_k}^2\),
the even observable \(\{F,a_{\Lambda_k}\}_{\mathrm P,\Lambda_k}\) already Poisson-commutes with
both \(a_{\Lambda_k}\) and \(S_{\Lambda_k}\), hence lies automatically in
\(\mathcal N^{\mathrm{even}}_{\Lambda_k}\).
The correct codomain for factorization and kernel identification is therefore the intrinsic
physical covariance-response space \(\mathcal T^{\mathrm{phys}}_{\Lambda_k}\), not the even
observable quotient.
\end{remark}

\begin{proposition}[Operational triviality is the kernel of the active response]
\label{prop:T-kerrho}
Assume \cref{ass:response-factorization,ass:response-sufficiency}. Then
\[
\mathcal T^{\mathrm{op}}_{\Lambda_k}=\ker\rho_{\Lambda_k}.
\]
\end{proposition}

\begin{proof}
Because \(C^{\mathrm{act}}_{\Lambda_k}\neq 0\), the factorization identity
\eqref{eq:rho-factorization} implies
\[
R^{\mathrm{phys}}_{\Lambda_k}(F)=0
\qquad\Longleftrightarrow\qquad
\rho_{\Lambda_k}(F)=0
\]
for every \(F\in E^-_{\Lambda_k}\).
Thus
\(
\ker R^{\mathrm{phys}}_{\Lambda_k}
=
\ker\rho_{\Lambda_k}.
\)
Apply \cref{ass:response-sufficiency}.
\end{proof}

\begin{corollary}[Rank-one active odd quotient]
\label{cor:rank-one-char}
Assume \cref{ass:response-factorization,ass:response-sufficiency}. Then
\[
E^-_{\Lambda_k}/\mathcal T^{\mathrm{op}}_{\Lambda_k}\cong \mathbb R,
\]
and the quotient is generated by the class of the active odd generator
\([J^{\mathrm{act}}_{\Lambda_k}]\).
\end{corollary}

\begin{proof}
By \cref{prop:T-kerrho}, the quotient is canonically isomorphic to the image of the nonzero
functional \(\rho_{\Lambda_k}\). Since \(\rho_{\Lambda_k}\) is real-valued and
\(\rho_{\Lambda_k}(J^{\mathrm{act}}_{\Lambda_k})=1\), its image is \(\mathbb R\) and the quotient
is generated by \([J^{\mathrm{act}}_{\Lambda_k}]\).
\end{proof}

\begin{proposition}[BRST/gauge-exact compatibility on the characterization layer]
\label{prop:ambiguity-brst-char}
For each accepted tick \(k\), the odd parts of the BRST/gauge-exact directions
\[
Q_tX+Q_uY
\]
lie in \(\mathcal T^{\mathrm{op}}_{\Lambda_k}\).
\end{proposition}

\begin{proof}
The exported characterization observables \(a_{\Lambda_k}\) and \(S_{\Lambda_k}\) are physical
finite-cutoff observables on the characterization layer. The BRST/gauge-exact odd directions
\(Q_tX+Q_uY\) represent redundancies of that physical layer, so their Peierls action vanishes on
\(\mathcal O^{\mathrm{char}}_{\Lambda_k}\). By \eqref{eq:Tchar-def}, their odd parts therefore lie
in \(\mathcal T^{\mathrm{op}}_{\Lambda_k}\).
\end{proof}

\begin{proposition}[Boundary-admissible divergence compatibility]
\label{prop:ambiguity-div-char}
Let \(K^\mu\) be a sufficiently local current whose renormalized boundary flux vanishes on the core
characterization algebra \(\mathcal O^{\mathrm{char}}_{\Lambda_k}\). Then the odd part of
\[
\partial_\mu K^\mu
\]
lies in \(\mathcal T^{\mathrm{op}}_{\Lambda_k}\).
\end{proposition}

\begin{proof}
For the admissible currents in the statement, the renormalized Stokes reduction of
\(\partial_\mu K^\mu\) produces no boundary term on the characterization layer. Hence its Peierls
action on every \(O\in\mathcal O^{\mathrm{char}}_{\Lambda_k}\) vanishes, and the odd part belongs
to \(\mathcal T^{\mathrm{op}}_{\Lambda_k}\) by \eqref{eq:Tchar-def}.
\end{proof}

\begin{definition}[Admissible odd counterterms on the characterization layer]
\label{def:ambiguity-ct-char}
For each accepted tick \(k\), an odd sufficiently local counterterm ambiguity
\[
\Delta_{\rm ct}^{-}\in E^-_{\Lambda_k}
\]
is called \emph{admissible on the characterization layer} if its intrinsic covariance response has
no physical projection:
\[
\Pi_{\mathrm{phys},\Lambda_k}\,\mathcal R_{\Lambda_k}(\Delta_{\rm ct}^{-})=0.
\]
Equivalently, admissible odd counterterms are precisely those local scheme shifts whose response is
confined to the RP, Weyl, conformal, or other even-null monitor sectors and does not enter the
compressed active odd line.
\end{definition}

\begin{proposition}[Counterterm rigidity on the characterization layer]
\label{prop:ambiguity-ct-char}
Assume that the characterization observable \(a_{\Lambda_k}\) is extracted from the physical
covariance response so that
\[
\Pi_{\mathrm{phys},\Lambda_k}\mathcal R_{\Lambda_k}(F)=0
\quad\Longrightarrow\quad
\{F,a_{\Lambda_k}\}_{\mathrm P,\Lambda_k}=0
\qquad
\forall F\in E^-_{\Lambda_k}.
\]
Then every admissible odd counterterm ambiguity in the sense of \cref{def:ambiguity-ct-char} acts
trivially on the core characterization algebra:
\[
\Delta_{\rm ct}^{-}\in \mathcal T^{\mathrm{op}}_{\Lambda_k}.
\]
\end{proposition}

\begin{proof}
If \(F=\Delta_{\rm ct}^{-}\) is admissible, then this physical projection vanishes by definition.
The displayed hypothesis therefore gives
\[
\{\Delta_{\rm ct}^{-},a_{\Lambda_k}\}_{\mathrm P,\Lambda_k}=0,
\]
and since \(S_{\Lambda_k}=a_{\Lambda_k}^2\), the derivation property yields
\[
\{\Delta_{\rm ct}^{-},S_{\Lambda_k}\}_{\mathrm P,\Lambda_k}=0.
\]
By \cref{prop:char-op-from-a}, the first identity already implies
\(\Delta_{\rm ct}^{-}\in\mathcal T^{\mathrm{op}}_{\Lambda_k}\).
\end{proof}

\begin{remark}[Active stiffness versus full calibration]
\label{rem:char-calibration-split}
The split/leak analysis determines the active stiffness \(\widetilde\alpha(\Lambda_k)\) on the
distinguished line generated by \(J^{\mathrm{act}}_{\Lambda_k}\), with
\[
\widetilde\alpha(\Lambda_k)\asymp \alpha(\Lambda_k)\sim \Lambda_k^{d+2}
\]
by \cref{ass:uv-match}. What UOM does \emph{not} determine canonically is a positive bilinear form
on the whole odd generator space \(E^-_{\Lambda_k}\). That extension is auxiliary interface data,
not a UOM-dynamical output.
\end{remark}

\begin{assumption}[Auxiliary coercive core on the operationally trivial sector]
\label{ass:aux-char-core}
For each accepted tick \(k\), choose a continuous symmetric coercive bilinear form
\[
\langle\cdot,\cdot\rangle^{\mathrm{aux}}_{\Lambda_k}
:
\mathcal T^{\mathrm{op}}_{\Lambda_k}\times \mathcal T^{\mathrm{op}}_{\Lambda_k}\to\mathbb R
\]
whose induced norm Hilbertizes \(\mathcal T^{\mathrm{op}}_{\Lambda_k}\).
\end{assumption}

\begin{proposition}[Strong-package calibration extension from active stiffness]
\label{prop:char-calibration-extension}
Assume \cref{ass:response-factorization,ass:response-sufficiency,ass:aux-char-core}. Set
\[
\chi_{\Lambda_k}:=\rho_{\Lambda_k},
\qquad
\mathcal T_{\Lambda_k}:=\mathcal T^{\mathrm{op}}_{\Lambda_k}.
\]
Then the formula
\begin{equation}
\mathsf B_{\Lambda_k}(F,G)
:=
\widetilde\alpha(\Lambda_k)\,\chi_{\Lambda_k}(F)\chi_{\Lambda_k}(G)
+
\left\langle
F-\chi_{\Lambda_k}(F)J^{\mathrm{act}}_{\Lambda_k},
\,G-\chi_{\Lambda_k}(G)J^{\mathrm{act}}_{\Lambda_k}
\right\rangle^{\mathrm{aux}}_{\Lambda_k}
\label{eq:B-extension}
\end{equation}
defines a continuous symmetric positive definite bilinear form on \(E^-_{\Lambda_k}\) satisfying
\[
\mathsf B_{\Lambda_k}(J^{\mathrm{act}}_{\Lambda_k},J^{\mathrm{act}}_{\Lambda_k})
=
\widetilde\alpha(\Lambda_k),
\qquad
\mathsf B_{\Lambda_k}(J^{\mathrm{act}}_{\Lambda_k},T)=0
\quad
\forall T\in\mathcal T_{\Lambda_k}.
\]
\end{proposition}

\begin{proof}
By \cref{prop:T-kerrho}, every \(F\in E^-_{\Lambda_k}\) has a unique decomposition
\[
F=\chi_{\Lambda_k}(F)\,J^{\mathrm{act}}_{\Lambda_k}+T_F,
\qquad
T_F\in \mathcal T_{\Lambda_k}.
\]
Equation \eqref{eq:B-extension} is therefore well-defined, continuous, and symmetric.
Positivity is immediate from \(\widetilde\alpha(\Lambda_k)>0\) together with coercivity of
\(\langle\cdot,\cdot\rangle^{\mathrm{aux}}_{\Lambda_k}\) on \(\mathcal T_{\Lambda_k}\).
If \(\mathsf B_{\Lambda_k}(F,F)=0\), then both \(\chi_{\Lambda_k}(F)=0\) and
\(\langle T_F,T_F\rangle^{\mathrm{aux}}_{\Lambda_k}=0\), hence \(T_F=0\) and \(F=0\).
The displayed normalization and orthogonality properties follow by substitution into
\eqref{eq:B-extension}.
\end{proof}

\begin{corollary}[Descendant ambiguity compatibility on the characterization layer]
\label{cor:ambiguity-char}
For each accepted tick \(k\), the odd-projected ambiguity directions coming from
\[
Q_tX+Q_uY,
\qquad
\partial_\mu K^\mu,
\qquad
\Delta_{\rm ct}^{-}
\]
belong to \(\mathcal T^{\mathrm{op}}_{\Lambda_k}\) whenever \(K^\mu\) is taken from the
boundary-admissible class of \cref{prop:ambiguity-div-char} and \(\Delta_{\rm ct}^{-}\) is
admissible in the sense of \cref{def:ambiguity-ct-char}.
\end{corollary}

\begin{proof}
Combine \cref{prop:ambiguity-brst-char,prop:ambiguity-div-char,prop:ambiguity-ct-char}.
\end{proof}

\begin{theorem}[Strong rank-one realization of the characterization layer]
\label{thm:uom-char-export}
Assume
\cref{ass:export-selection,ass:centered-odd,ass:uv-match,ass:response-factorization,ass:response-sufficiency,ass:aux-char-core}.
Then for each accepted tick \(k\), the split/leak construction exports the UOM-side data
\begin{equation}
\big(
a_{\Lambda_k},\,
S_{\Lambda_k},\,
J^{\mathrm{act}}_{\Lambda_k},\,
C^{\mathrm{act}}_{\Lambda_k},\,
\widetilde\alpha(\Lambda_k),\,
\rho_{\Lambda_k}
\big),
\label{eq:uom-char-realization}
\end{equation}
where:
\begin{enumerate}[leftmargin=2em]
\item \(\mathcal O^{\mathrm{char}}_{\Lambda_k}=\mathrm{alg}\{a_{\Lambda_k},S_{\Lambda_k}\}\) is the
  minimal completion-sector characterization algebra;
\item the operationally trivial subspace satisfies
  \[
  \mathcal T_{\Lambda_k}:=\mathcal T^{\mathrm{op}}_{\Lambda_k}
  =
  \ker R^{\mathrm{phys}}_{\Lambda_k}
  =
  \ker\rho_{\Lambda_k};
  \]
\item the intrinsic physical response is one-dimensional on the strong sector:
  \[
  R^{\mathrm{phys}}_{\Lambda_k}(F)=\rho_{\Lambda_k}(F)\,C^{\mathrm{act}}_{\Lambda_k},
  \qquad
  R^{\mathrm{phys}}_{\Lambda_k}(J^{\mathrm{act}}_{\Lambda_k})=C^{\mathrm{act}}_{\Lambda_k}\neq 0;
  \]
\item the active odd quotient is rank one and is generated by
  \([J^{\mathrm{act}}_{\Lambda_k}]\);
\item the abstract orientation functional of the transported interface is realized by
  \[
  \chi_{\Lambda_k}:=\rho_{\Lambda_k};
  \]
\item any admissible calibration form \(\mathsf B_{\Lambda_k}\) in this strong package may be obtained by
  extending the active stiffness \(\widetilde\alpha(\Lambda_k)\) with
  \cref{prop:char-calibration-extension}.
\end{enumerate}
\end{theorem}

\begin{proof}
\Cref{def:char-data} packages the active completion-sector channel extracted from the split/leak
construction: \(a_{\Lambda_k}\) is the distinguished odd split coordinate singled out by
\cref{subsec:leak-solution,subsec:B-gaussian}, while \(S_{\Lambda_k}=a_{\Lambda_k}^2\) is the even
stock entering \eqref{eq:Qleak-effective} and \eqref{eq:split-threshold-effective}. The active odd
generator \(J^{\mathrm{act}}_{\Lambda_k}\), active physical covariance tangent
\(C^{\mathrm{act}}_{\Lambda_k}\), and scalar response coefficient \(\rho_{\Lambda_k}\) are
exported by \cref{ass:response-factorization}; \cref{prop:T-kerrho,cor:rank-one-char} then identify
the abstract interface-trivial subspace and the corresponding rank-one quotient.
The split/leak law fixes the active stiffness \(\widetilde\alpha(\Lambda_k)\), and
\cref{prop:char-calibration-extension} upgrades that one-line datum to an admissible full
calibration form once the auxiliary coercive core is chosen.
Finally, \cref{cor:ambiguity-char} places the odd-projected descendant ambiguities in
\(\mathcal T_{\Lambda_k}\), yielding the stated ambiguity compatibility on the characterization
layer.
\end{proof}

\begin{theorem}[Intrinsic weak transported descendant-class interface]
\label{thm:char-weak-interface}
Assume the hypotheses of \cref{thm:active-response-transport-minimal} on a connected accepted-branch
segment \(I\).
Then for every \(\Lambda\in I\), the sufficiently local odd generator layer carries the weak
transported interface data
\[
\big(
E^-_{\Lambda},\,
\mathcal A^-_{\Lambda},\,
[\Xi^{(1)}_{\Lambda}],\,
\chi_{\Lambda},\,
\varepsilon_{\Lambda}
\big)
\]
with
\[
\chi_{\Lambda}
:=
\nu_{\Lambda}^{-1}\rho^{\mathrm{dom}}_{\Lambda},
\]
and the following conclusions hold.
\begin{enumerate}[leftmargin=2em]
\item there exists a representative
  \[
  J^{\mathrm{desc}}_{\Lambda}\in[\Xi^{(1)}_{\Lambda}]
  \]
  such that
  \[
  \rho^{\mathrm{dom}}_{\Lambda}\!\big(J^{\mathrm{desc}}_{\Lambda}\big)\neq 0;
  \]
  equivalently, the transported descendant sector has nonzero overlap with the canonically selected
  compressed dominant odd line;
\item the orientation functional is realized intrinsically on the sufficiently local generator layer:
  \[
  \mathcal A^-_{\Lambda}\subset\ker\chi_{\Lambda},
  \qquad
  \chi_{\Lambda}(J)=\varepsilon_{\Lambda}
  \quad
  \forall J\in[\Xi^{(1)}_{\Lambda}];
  \]
\item along \(I\), the sign \(\varepsilon_{\Lambda}\) is locally constant.
\end{enumerate}
In particular, the weak transported interface is realized intrinsically, without appeal to completion
coordinates, under the stated hypotheses.
\end{theorem}

\begin{proof}
The functional \(\rho^{\mathrm{dom}}_{\Lambda}\) is defined intrinsically from the sufficiently local
generator layer by \cref{def:intrinsic-cov-response} together with the dominant-line prescription of
\cref{def:dominant-response-functional}.
Item (1) is exactly the nonzero-overlap conclusion of item (3) in
\cref{thm:active-response-transport-minimal}, item (2) is exactly item (4) there, and item (3) is
item (5).
\end{proof}

\begin{theorem}[Quadratic-probe realization of the weak activity supplement]
\label{thm:char-weak-activity}
Assume the hypotheses of \cref{thm:active-response-transport-minimal} on a connected accepted-branch
segment \(I\).
For each \(\Lambda\in I\), define the even quadratic probe sector by
\[
\mathcal P^+_{\Lambda}
:=
\operatorname{span}
\bigl\{
\mathcal Q^{\Lambda}_{f,g}
:
\text{\(f,g\) admissible band-limited receiver probes}
\bigr\}
\subset E^+_{\Lambda}.
\]
Then, for every \(F\in E^-_{\Lambda}\),
\[
\{F,A\}_{\mathrm P,\Lambda}=0
\qquad
\forall A\in\mathcal P^+_{\Lambda}
\]
implies
\[
\chi_{\Lambda}(F)=0.
\]
Equivalently, if \(\chi_{\Lambda}(F)\neq 0\), then there exists a sufficiently local even quadratic
probe \(A\in\mathcal P^+_{\Lambda}\) such that
\[
\{F,A\}_{\mathrm P,\Lambda}\neq 0.
\]
In particular, the weak interface data of \cref{thm:char-weak-interface} carry a canonical
activity witness sector built from the intrinsic quadratic receiver probes.
\end{theorem}

\begin{proof}
Assume
\[
\{F,A\}_{\mathrm P,\Lambda}=0
\qquad
\forall A\in\mathcal P^+_{\Lambda}.
\]
In particular,
\[
\{F,\mathcal Q^{\Lambda}_{f,g}\}_{\mathrm P,\Lambda}=0
\]
for every admissible pair \(f,g\).
By \cref{def:intrinsic-cov-response}, this means
\[
\big\langle f,\mathcal R_{\Lambda}(F)g\big\rangle
=
\Big\langle
\{F,\mathcal Q^{\Lambda}_{f,g}\}_{\mathrm P,\Lambda}
\Big\rangle_{\rho_{\ast,\Lambda}}
=
0
\]
for all admissible \(f,g\), hence
\(
\mathcal R_{\Lambda}(F)=0
\).
Therefore
\[
\Pi_{\mathrm{phys},\Lambda}\mathcal R_{\Lambda}(F)=0,
\]
and the definition of the raw dominant-response coefficient in
\cref{def:dominant-response-functional} gives
\[
\widetilde\rho^{\mathrm{dom}}_{\Lambda}(F)=0.
\]
By item (1) of \cref{thm:active-response-transport-minimal},
\(
\delta C^{\mathrm{LoG}}_{\Lambda}\neq 0
\),
so dominant-line realization holds and \(\rho^{\mathrm{dom}}_{\Lambda}\) is well defined.
Hence
\[
\rho^{\mathrm{dom}}_{\Lambda}(F)=0.
\]
Now apply \cref{thm:char-weak-interface}, which identifies
\(
\chi_{\Lambda}
=
\nu_{\Lambda}^{-1}\rho^{\mathrm{dom}}_{\Lambda}
\),
to obtain
\[
\chi_{\Lambda}(F)=0.
\]
The contrapositive gives the equivalent detection statement.
\end{proof}

\begin{remark}[Correct scalar form of representative comparisons]
\label{rem:char-weak-interface-rho-first}
The weak export theorem should be read \(\rho^{\mathrm{dom}}\)-first on the ambiguity quotient.
Because
\[
\chi_{\Lambda}
:=
\nu_{\Lambda}^{-1}\rho^{\mathrm{dom}}_{\Lambda},
\]
any comparison statement that depends only on dominant pairing is naturally a statement about the
descended scalar
\[
\overline{\rho}^{\mathrm{dom}}_{\Lambda}
:
E^-_{\Lambda}/\mathcal A^-_{\Lambda}\to\mathbb R
\]
or, equivalently, about its value on the transported descendant affine class.
Such a statement may be rewritten in \(\chi_{\Lambda}\) form only when the normalization is made
explicit; for example, if \(\nu_{\Lambda}=1\), then the same scalar value is represented by a
descended \(\bar\chi_{\Lambda}\).
Without that normalization the stronger unscaled \(\chi\)-compatibility form is false in general.
\end{remark}

\begin{corollary}[UOM-stiffness specialization of the weak transported calibration]
\label{cor:char-weak-calibration}
Assume \cref{ass:uv-match} and fix an accepted tick \(k\) for which the conclusions of
\cref{thm:char-weak-interface} hold.
Suppose \(\ker\chi_{\Lambda_k}\) admits a continuous symmetric coercive bilinear form whose induced
norm Hilbertizes it, and choose any
\[
J^\sharp_{\Lambda_k}\in E^-_{\Lambda_k}
\qquad
\text{with}
\qquad
\chi_{\Lambda_k}(J^\sharp_{\Lambda_k})=1.
\]
Then \cref{prop:minimal-transport-calibration} may be applied with
\[
\beta_{\Lambda_k}:=\widetilde\alpha(\Lambda_k)
\]
to produce a continuous symmetric positive definite bilinear form
\[
\mathsf B^{\widetilde\alpha}_{\Lambda_k}
:
E^-_{\Lambda_k}\times E^-_{\Lambda_k}\to\mathbb R
\]
such that
\[
\mathsf B^{\widetilde\alpha}_{\Lambda_k}(J^\sharp_{\Lambda_k},J^\sharp_{\Lambda_k})
=
\widetilde\alpha(\Lambda_k),
\qquad
\mathsf B^{\widetilde\alpha}_{\Lambda_k}(J^\sharp_{\Lambda_k},K)=0
\quad
\forall K\in\ker\chi_{\Lambda_k}.
\]
In particular,
\[
\mathcal A^-_{\Lambda_k}\subset\ker\chi_{\Lambda_k}
\]
implies
\[
\mathsf B^{\widetilde\alpha}_{\Lambda_k}(J^\sharp_{\Lambda_k},A)=0
\qquad
\forall A\in\mathcal A^-_{\Lambda_k}.
\]
\end{corollary}

\begin{proof}
By \cref{rem:char-calibration-split}, the split/leak package fixes the positive stiffness scale
\(\widetilde\alpha(\Lambda_k)\), and \cref{ass:uv-match} records its positivity and asymptotic
size.
Apply \cref{prop:minimal-transport-calibration} with
\(\beta_{\Lambda_k}=\widetilde\alpha(\Lambda_k)\).
The final orthogonality statement follows from item (2) of \cref{thm:char-weak-interface}.
\end{proof}

\begin{remark}[Inherited ledger compatibility]
\label{rem:ledger-compat}
The inherited commutative ledger algebra is intentionally not part of the core characterization
algebra \(\mathcal O^{\mathrm{char}}_{\Lambda_k}\).
Its role is separate: the split model requires the active odd representative to be completion-sector
nontrivial while acting trivially, to leading order, on the inherited classical ledger.
This is the UOM-side compatibility statement behind the slogan
``oddness in the completion sector, classicality in the accessible ledger.''
\end{remark}

\ifdefined\UOMTransportFragmentLoaded\else
\end{document}
\fi
%
}{%
  \IfFileExists{UOM/canonical_odd_interface_note.tex}{%
    \ifdefined\UOMTransportFragmentLoaded\else
\documentclass[11pt]{article}
\usepackage[margin=1in]{geometry}
\usepackage{amsmath,amssymb,amsthm,mathtools}
\usepackage{bm}
\usepackage{microtype}
\usepackage{enumitem}
\usepackage{booktabs}
\usepackage{array}
\usepackage{xcolor}
\usepackage{xr-hyper}
\usepackage{hyperref}
\usepackage[nameinlink,noabbrev]{cleveref}

\providecommand{\Fsloc}{\mathcal F_{\rm s.loc}}
\providecommand{\Qt}{Q_t}
\externaldocument[main-]{UOM/main}
\externaldocument[uomdesc-]{UOM/uom_descendant_note}
\externaldocument{UOM/completion_sector_split_model_note}

\hypersetup{
  colorlinks=true,
  linkcolor=blue!60!black,
  citecolor=blue!60!black,
  urlcolor=blue!60!black
}

\theoremstyle{definition}
\newtheorem{definition}{Definition}[section]
\newtheorem{assumption}{Assumption}[section]
\theoremstyle{plain}
\newtheorem{theorem}[definition]{Theorem}
\newtheorem{lemma}[definition]{Lemma}
\newtheorem{proposition}[definition]{Proposition}
\newtheorem{corollary}[definition]{Corollary}
\theoremstyle{remark}
\newtheorem{remark}[definition]{Remark}
\crefname{definition}{definition}{definitions}
\Crefname{definition}{Definition}{Definitions}
\crefname{assumption}{assumption}{assumptions}
\Crefname{assumption}{Assumption}{Assumptions}
\crefname{theorem}{theorem}{theorems}
\Crefname{theorem}{Theorem}{Theorems}
\crefname{lemma}{lemma}{lemmas}
\Crefname{lemma}{Lemma}{Lemmas}
\crefname{proposition}{proposition}{propositions}
\Crefname{proposition}{Proposition}{Propositions}
\crefname{corollary}{corollary}{corollaries}
\Crefname{corollary}{Corollary}{Corollaries}
\crefname{remark}{remark}{remarks}
\Crefname{remark}{Remark}{Remarks}

\newcommand{\cH}{\mathcal H}
\newcommand{\cA}{\mathcal A}
\newcommand{\cD}{\mathcal D}
\newcommand{\cB}{\mathcal B}
\newcommand{\cC}{\mathcal C}
\newcommand{\cL}{\mathcal L}
\newcommand{\cK}{\mathcal K}
\newcommand{\cR}{\mathcal R}
\newcommand{\cG}{\mathcal G}
\newcommand{\Ztwo}{\mathbb Z_2}
\newcommand{\Sys}{\mathrm{Sys}}
\newcommand{\Time}{\mathrm{Time}}
\newcommand{\RG}{\mathrm{RG}}
\newcommand{\Tr}{\mathrm{Tr}}
\newcommand{\diag}{\mathrm{diag}}
\newcommand{\ad}{\mathrm{ad}}
\newcommand{\id}{\mathrm{id}}
\newcommand{\norm}[1]{\left\lVert #1\right\rVert}
\newcommand{\abs}[1]{\left\lvert #1\right\rvert}
\newcommand{\ip}[2]{\left\langle #1,#2\right\rangle}

\externaldocument{UOM/compressed_log_spectral_selection_note}
\externaldocument{UOM/odd_selector_note}
\externaldocument{UOM/active_response_transport_program_note}
\title{Polarized Characterization Interface}
\author{}
\date{}
\begin{document}
\maketitle
\fi

\section{Polarized Characterization Interface}
\label{sec:char-export}

This section records the finite-cutoff characterization data extracted from the split/leak
construction and used in the 2T Main assumption package
Definition~\ref{main-def:uom-main-theorem-assumptions}. In this note, an ``accepted tick'' denotes the local
accepted-branch slice at cutoff \(\Lambda_k\).

For each accepted tick \(k\), write
\[
a_{\Lambda_k}:=a_k,
\qquad
S_{\Lambda_k}:=a_{\Lambda_k}^2.
\]
The variable \(a_{\Lambda_k}\) is the distinguished odd split coordinate identified in
\cref{subsec:leak-solution,subsec:B-gaussian}; the even quantity \(S_{\Lambda_k}\) is the
positive stock that enters the leak law and the split threshold.

\begin{definition}[Characterization algebra, operational triviality, and even null ideal]
\label{def:char-data}
Fix an accepted tick \(k\) and let \(E^-_{\Lambda_k}\) denote the physically odd sufficiently local
generator space at cutoff \(\Lambda_k\), i.e.\ the odd sector for the transported physical involution
\(\Gamma\). Let \(E^+_{\Lambda_k}\) denote the corresponding physically even sufficiently local
observable space.
Define the core characterization algebra by
\begin{equation}
\mathcal O^{\mathrm{char}}_{\Lambda_k}
:=
\mathrm{alg}\{a_{\Lambda_k},\,S_{\Lambda_k}\}.
\label{eq:Ochar-def}
\end{equation}
Define the operationally interface-trivial odd subspace by
\begin{equation}
\mathcal T^{\mathrm{op}}_{\Lambda_k}
:=
\overline{\left\{
T\in E^-_{\Lambda_k}:
\{T,O\}_{\mathrm P,\Lambda_k}=0
\ \forall O\in\mathcal O^{\mathrm{char}}_{\Lambda_k}
\right\}}.
\label{eq:Tchar-def}
\end{equation}
Define the even interface-null ideal by
\begin{equation}
\mathcal N^{\mathrm{even}}_{\Lambda_k}
:=
\overline{\left\{
E\in E^+_{\Lambda_k}:
\{E,O\}_{\mathrm P,\Lambda_k}=0
\ \forall O\in\mathcal O^{\mathrm{char}}_{\Lambda_k}
\right\}}.
\label{eq:Neven-def}
\end{equation}
\end{definition}

\begin{remark}[Do not identify the characterization observable with the generator]
\label{rem:a-vs-Jact}
The odd observable \(a_{\Lambda_k}\) characterizes the active completion-sector channel, but it is
not itself the odd Hamiltonian generator class appearing in item \textbf{(M4)} of
Definition~\ref{main-def:uom-main-theorem-assumptions}.
That generator is denoted \(J^{\mathrm{act}}_{\Lambda_k}\in E^-_{\Lambda_k}\) below.
The quotient statement belongs to \([J^{\mathrm{act}}_{\Lambda_k}]\), not to
\([a_{\Lambda_k}]\).
\end{remark}

\begin{definition}[Intrinsic covariance-tangent response map]
\label{def:intrinsic-cov-response}
Fix an accepted tick \(k\). Let \(\mathcal T^{\mathrm{cov}}_{\Lambda_k}\) denote the finite-cutoff
receiver covariance-tangent space at the UOM reference Gaussian state
\(\rho_{\ast,\Lambda_k}\), equipped with the Fisher pairing
\[
\big\langle\!\big\langle \delta C,\delta C'\big\rangle\!\big\rangle_{C_{\ast,\Lambda_k}}
:=
\tfrac12\,\mathrm{tr}\!\big(
C_{\ast,\Lambda_k}^{-1}\delta C\,C_{\ast,\Lambda_k}^{-1}\delta C'
\big).
\]
For admissible band-limited receiver probes \(f,g\), let
\[
\mathcal Q^{\Lambda_k}_{f,g}\in E^+_{\Lambda_k}
\]
denote the corresponding sufficiently local quadratic observable, normalized so that on centered
Gaussian states with covariance \(C\),
\[
\big\langle \mathcal Q^{\Lambda_k}_{f,g}\big\rangle_{\rho(C)}
=
\langle f,Cg\rangle.
\]
For \(F\in E^-_{\Lambda_k}\), define \(\mathcal R_{\Lambda_k}(F)\in\mathcal T^{\mathrm{cov}}_{\Lambda_k}\)
by the matrix elements
\[
\big\langle f,\mathcal R_{\Lambda_k}(F)g\big\rangle
:=
\Big\langle
\{F,\mathcal Q^{\Lambda_k}_{f,g}\}_{\mathrm P,\Lambda_k}
\Big\rangle_{\rho_{\ast,\Lambda_k}}.
\]
We call
\[
\mathcal R_{\Lambda_k}:E^-_{\Lambda_k}\to\mathcal T^{\mathrm{cov}}_{\Lambda_k}
\]
the \emph{intrinsic covariance-tangent response map}.
\end{definition}

\begin{proposition}[Intrinsic definition and completion independence]
\label{prop:intrinsic-cov-response}
The map \(\mathcal R_{\Lambda_k}\) of \cref{def:intrinsic-cov-response} is continuous and linear.
It depends only on the finite-cutoff UOM Peierls package and the UOM covariance/Fisher package.
If a Gaussian completion realization is chosen, then \(\mathcal R_{\Lambda_k}(F)\) is represented by
the same completion-sector covariance tangent that one obtains by linearizing the completion
quadratures. Thus Gaussian completion is a realization of the response map, not extra structure
required to define it.
\end{proposition}

\begin{proof}
For fixed \(F\in E^-_{\Lambda_k}\), the finite-cutoff UOM Peierls package makes
\[
\delta_F:=\{F,-\}_{\mathrm P,\Lambda_k}
\]
a continuous derivation on the sufficiently local algebra. Applying \(\delta_F\) to the quadratic
probes \(\mathcal Q^{\Lambda_k}_{f,g}\) and evaluating in the reference Gaussian state produces a
continuous symmetric bilinear form in \(f,g\). In the UOM Gaussian radiative sector, such bilinear
forms are exactly covariance tangents at \(C_{\ast,\Lambda_k}\), equipped with the Fisher metric.
Linearity and continuity in \(F\) are inherited directly from the Peierls bracket. No completion
coordinates are used in this construction.

If a Gaussian completion is later chosen, the same second moments are encoded by the covariance
matrix of the completion quadratures. The bilinear form defined above and the completion-sector
covariance tangent therefore represent the same infinitesimal response object in two different
coordinate systems.
\end{proof}

\begin{definition}[Physical intrinsic covariance-response map]
\label{def:phys-cov-response}
Fix an accepted tick \(k\). Let
\[
\mathcal T^{\mathrm{phys}}_{\Lambda_k}
:=
\Pi_{\mathrm{phys},\Lambda_k}\mathcal T^{\mathrm{cov}}_{\Lambda_k}
\subset \mathcal T^{\mathrm{cov}}_{\Lambda_k}.
\]
Define the physical intrinsic covariance-response map by
\[
R^{\mathrm{phys}}_{\Lambda_k}
:=
\Pi_{\mathrm{phys},\Lambda_k}\mathcal R_{\Lambda_k}
:
E^-_{\Lambda_k}\to \mathcal T^{\mathrm{phys}}_{\Lambda_k}.
\]
\end{definition}

\begin{proposition}[Physical intrinsic response map]
\label{prop:phys-cov-response}
The map \(R^{\mathrm{phys}}_{\Lambda_k}\) of \cref{def:phys-cov-response} is continuous and linear.
\end{proposition}

\begin{proof}
This is immediate from \cref{prop:intrinsic-cov-response} and the continuity and linearity of the
finite-cutoff physical projector \(\Pi_{\mathrm{phys},\Lambda_k}\).
\end{proof}

\begin{proposition}[Reduction of operational triviality to the active odd coordinate]
\label{prop:char-op-from-a}
Fix an accepted tick \(k\). If \(F\in E^-_{\Lambda_k}\) satisfies
\[
\{F,a_{\Lambda_k}\}_{\mathrm P,\Lambda_k}=0,
\]
then \(F\in\mathcal T^{\mathrm{op}}_{\Lambda_k}\).
\end{proposition}

\begin{proof}
Since \(S_{\Lambda_k}=a_{\Lambda_k}^2\) and the Peierls bracket is a derivation,
\[
\{F,S_{\Lambda_k}\}_{\mathrm P,\Lambda_k}
=
\{F,a_{\Lambda_k}^2\}_{\mathrm P,\Lambda_k}
=
2\,a_{\Lambda_k}\,\{F,a_{\Lambda_k}\}_{\mathrm P,\Lambda_k}
=
0.
\]
The algebra \(\mathcal O^{\mathrm{char}}_{\Lambda_k}=\mathrm{alg}\{a_{\Lambda_k},S_{\Lambda_k}\}\)
is generated by \(a_{\Lambda_k}\) and \(S_{\Lambda_k}\), and the Peierls bracket acts on it by the
Leibniz rule. Hence \(\{F,O\}_{\mathrm P,\Lambda_k}=0\) for every
\(O\in\mathcal O^{\mathrm{char}}_{\Lambda_k}\), proving
\(F\in\mathcal T^{\mathrm{op}}_{\Lambda_k}\).
\end{proof}

\begin{assumption}[Active odd generator and intrinsic-response factorization]
\label{ass:response-factorization}
For each accepted tick \(k\), there exist a distinguished active odd generator
\[
J^{\mathrm{act}}_{\Lambda_k}\in E^-_{\Lambda_k}
\]
and a nonzero physical covariance tangent
\[
C^{\mathrm{act}}_{\Lambda_k}\in \mathcal T^{\mathrm{phys}}_{\Lambda_k}
\]
and a continuous nonzero linear functional
\[
\rho_{\Lambda_k}:E^-_{\Lambda_k}\to\mathbb R
\]
normalized by
\[
\rho_{\Lambda_k}(J^{\mathrm{act}}_{\Lambda_k})=1,
\]
such that
\[
R^{\mathrm{phys}}_{\Lambda_k}(J^{\mathrm{act}}_{\Lambda_k})
=
C^{\mathrm{act}}_{\Lambda_k}
\neq 0
\]
and
\begin{equation}
R^{\mathrm{phys}}_{\Lambda_k}(F)
=
\rho_{\Lambda_k}(F)\,C^{\mathrm{act}}_{\Lambda_k}
\qquad
\forall F\in E^-_{\Lambda_k}
\label{eq:rho-factorization}
\end{equation}
holds on the strong sector under consideration.
\end{assumption}

\begin{remark}[Strong normalization]
\label{rem:rho-strong}
If one further chooses the Fisher-unit normalization
\(
\widehat C^{\mathrm{act}}_{\Lambda_k}
:=
\|C^{\mathrm{act}}_{\Lambda_k}\|^{-1}_{C_{\ast,\Lambda_k}}\,C^{\mathrm{act}}_{\Lambda_k},
\)
then
\eqref{eq:rho-factorization} sharpens to
\[
R^{\mathrm{phys}}_{\Lambda_k}(F)
=
\widehat{\rho}_{\Lambda_k}(F)\,\widehat C^{\mathrm{act}}_{\Lambda_k}
\]
with
\(
\widehat{\rho}_{\Lambda_k}
:=
\|C^{\mathrm{act}}_{\Lambda_k}\|_{C_{\ast,\Lambda_k}}\,\rho_{\Lambda_k}.
\)
The export below does not require this extra normalization; it only needs the scalar response
coefficient \(\rho_{\Lambda_k}\) relative to a fixed nonzero active tangent
\(C^{\mathrm{act}}_{\Lambda_k}\).
\end{remark}

\begin{assumption}[Physical-response sufficiency]
\label{ass:response-sufficiency}
For each accepted tick \(k\), the operationally trivial odd directions are exactly the kernel of the
physical intrinsic covariance-response map:
\[
 \ker R^{\mathrm{phys}}_{\Lambda_k}
 =
 \mathcal T^{\mathrm{op}}_{\Lambda_k}.
\]
Equivalently, an odd generator has zero physical intrinsic covariance response if and only if it is
operationally trivial on the full core characterization algebra \(\mathcal O^{\mathrm{char}}_{\Lambda_k}\).
\end{assumption}

\begin{remark}[Why the strong package must be stated intrinsically]
\label{rem:response-quotient-obstruction}
The earlier quotient-valued formulation of the strong package in
\(
E^+_{\Lambda_k}/\mathcal N^{\mathrm{even}}_{\Lambda_k}
\)
is not stable.
Because the Peierls bracket here is the ordinary derivation bracket and \(S_{\Lambda_k}=a_{\Lambda_k}^2\),
the even observable \(\{F,a_{\Lambda_k}\}_{\mathrm P,\Lambda_k}\) already Poisson-commutes with
both \(a_{\Lambda_k}\) and \(S_{\Lambda_k}\), hence lies automatically in
\(\mathcal N^{\mathrm{even}}_{\Lambda_k}\).
The correct codomain for factorization and kernel identification is therefore the intrinsic
physical covariance-response space \(\mathcal T^{\mathrm{phys}}_{\Lambda_k}\), not the even
observable quotient.
\end{remark}

\begin{proposition}[Operational triviality is the kernel of the active response]
\label{prop:T-kerrho}
Assume \cref{ass:response-factorization,ass:response-sufficiency}. Then
\[
\mathcal T^{\mathrm{op}}_{\Lambda_k}=\ker\rho_{\Lambda_k}.
\]
\end{proposition}

\begin{proof}
Because \(C^{\mathrm{act}}_{\Lambda_k}\neq 0\), the factorization identity
\eqref{eq:rho-factorization} implies
\[
R^{\mathrm{phys}}_{\Lambda_k}(F)=0
\qquad\Longleftrightarrow\qquad
\rho_{\Lambda_k}(F)=0
\]
for every \(F\in E^-_{\Lambda_k}\).
Thus
\(
\ker R^{\mathrm{phys}}_{\Lambda_k}
=
\ker\rho_{\Lambda_k}.
\)
Apply \cref{ass:response-sufficiency}.
\end{proof}

\begin{corollary}[Rank-one active odd quotient]
\label{cor:rank-one-char}
Assume \cref{ass:response-factorization,ass:response-sufficiency}. Then
\[
E^-_{\Lambda_k}/\mathcal T^{\mathrm{op}}_{\Lambda_k}\cong \mathbb R,
\]
and the quotient is generated by the class of the active odd generator
\([J^{\mathrm{act}}_{\Lambda_k}]\).
\end{corollary}

\begin{proof}
By \cref{prop:T-kerrho}, the quotient is canonically isomorphic to the image of the nonzero
functional \(\rho_{\Lambda_k}\). Since \(\rho_{\Lambda_k}\) is real-valued and
\(\rho_{\Lambda_k}(J^{\mathrm{act}}_{\Lambda_k})=1\), its image is \(\mathbb R\) and the quotient
is generated by \([J^{\mathrm{act}}_{\Lambda_k}]\).
\end{proof}

\begin{proposition}[BRST/gauge-exact compatibility on the characterization layer]
\label{prop:ambiguity-brst-char}
For each accepted tick \(k\), the odd parts of the BRST/gauge-exact directions
\[
Q_tX+Q_uY
\]
lie in \(\mathcal T^{\mathrm{op}}_{\Lambda_k}\).
\end{proposition}

\begin{proof}
The exported characterization observables \(a_{\Lambda_k}\) and \(S_{\Lambda_k}\) are physical
finite-cutoff observables on the characterization layer. The BRST/gauge-exact odd directions
\(Q_tX+Q_uY\) represent redundancies of that physical layer, so their Peierls action vanishes on
\(\mathcal O^{\mathrm{char}}_{\Lambda_k}\). By \eqref{eq:Tchar-def}, their odd parts therefore lie
in \(\mathcal T^{\mathrm{op}}_{\Lambda_k}\).
\end{proof}

\begin{proposition}[Boundary-admissible divergence compatibility]
\label{prop:ambiguity-div-char}
Let \(K^\mu\) be a sufficiently local current whose renormalized boundary flux vanishes on the core
characterization algebra \(\mathcal O^{\mathrm{char}}_{\Lambda_k}\). Then the odd part of
\[
\partial_\mu K^\mu
\]
lies in \(\mathcal T^{\mathrm{op}}_{\Lambda_k}\).
\end{proposition}

\begin{proof}
For the admissible currents in the statement, the renormalized Stokes reduction of
\(\partial_\mu K^\mu\) produces no boundary term on the characterization layer. Hence its Peierls
action on every \(O\in\mathcal O^{\mathrm{char}}_{\Lambda_k}\) vanishes, and the odd part belongs
to \(\mathcal T^{\mathrm{op}}_{\Lambda_k}\) by \eqref{eq:Tchar-def}.
\end{proof}

\begin{definition}[Admissible odd counterterms on the characterization layer]
\label{def:ambiguity-ct-char}
For each accepted tick \(k\), an odd sufficiently local counterterm ambiguity
\[
\Delta_{\rm ct}^{-}\in E^-_{\Lambda_k}
\]
is called \emph{admissible on the characterization layer} if its intrinsic covariance response has
no physical projection:
\[
\Pi_{\mathrm{phys},\Lambda_k}\,\mathcal R_{\Lambda_k}(\Delta_{\rm ct}^{-})=0.
\]
Equivalently, admissible odd counterterms are precisely those local scheme shifts whose response is
confined to the RP, Weyl, conformal, or other even-null monitor sectors and does not enter the
compressed active odd line.
\end{definition}

\begin{proposition}[Counterterm rigidity on the characterization layer]
\label{prop:ambiguity-ct-char}
Assume that the characterization observable \(a_{\Lambda_k}\) is extracted from the physical
covariance response so that
\[
\Pi_{\mathrm{phys},\Lambda_k}\mathcal R_{\Lambda_k}(F)=0
\quad\Longrightarrow\quad
\{F,a_{\Lambda_k}\}_{\mathrm P,\Lambda_k}=0
\qquad
\forall F\in E^-_{\Lambda_k}.
\]
Then every admissible odd counterterm ambiguity in the sense of \cref{def:ambiguity-ct-char} acts
trivially on the core characterization algebra:
\[
\Delta_{\rm ct}^{-}\in \mathcal T^{\mathrm{op}}_{\Lambda_k}.
\]
\end{proposition}

\begin{proof}
If \(F=\Delta_{\rm ct}^{-}\) is admissible, then this physical projection vanishes by definition.
The displayed hypothesis therefore gives
\[
\{\Delta_{\rm ct}^{-},a_{\Lambda_k}\}_{\mathrm P,\Lambda_k}=0,
\]
and since \(S_{\Lambda_k}=a_{\Lambda_k}^2\), the derivation property yields
\[
\{\Delta_{\rm ct}^{-},S_{\Lambda_k}\}_{\mathrm P,\Lambda_k}=0.
\]
By \cref{prop:char-op-from-a}, the first identity already implies
\(\Delta_{\rm ct}^{-}\in\mathcal T^{\mathrm{op}}_{\Lambda_k}\).
\end{proof}

\begin{remark}[Active stiffness versus full calibration]
\label{rem:char-calibration-split}
The split/leak analysis determines the active stiffness \(\widetilde\alpha(\Lambda_k)\) on the
distinguished line generated by \(J^{\mathrm{act}}_{\Lambda_k}\), with
\[
\widetilde\alpha(\Lambda_k)\asymp \alpha(\Lambda_k)\sim \Lambda_k^{d+2}
\]
by \cref{ass:uv-match}. What UOM does \emph{not} determine canonically is a positive bilinear form
on the whole odd generator space \(E^-_{\Lambda_k}\). That extension is auxiliary interface data,
not a UOM-dynamical output.
\end{remark}

\begin{assumption}[Auxiliary coercive core on the operationally trivial sector]
\label{ass:aux-char-core}
For each accepted tick \(k\), choose a continuous symmetric coercive bilinear form
\[
\langle\cdot,\cdot\rangle^{\mathrm{aux}}_{\Lambda_k}
:
\mathcal T^{\mathrm{op}}_{\Lambda_k}\times \mathcal T^{\mathrm{op}}_{\Lambda_k}\to\mathbb R
\]
whose induced norm Hilbertizes \(\mathcal T^{\mathrm{op}}_{\Lambda_k}\).
\end{assumption}

\begin{proposition}[Strong-package calibration extension from active stiffness]
\label{prop:char-calibration-extension}
Assume \cref{ass:response-factorization,ass:response-sufficiency,ass:aux-char-core}. Set
\[
\chi_{\Lambda_k}:=\rho_{\Lambda_k},
\qquad
\mathcal T_{\Lambda_k}:=\mathcal T^{\mathrm{op}}_{\Lambda_k}.
\]
Then the formula
\begin{equation}
\mathsf B_{\Lambda_k}(F,G)
:=
\widetilde\alpha(\Lambda_k)\,\chi_{\Lambda_k}(F)\chi_{\Lambda_k}(G)
+
\left\langle
F-\chi_{\Lambda_k}(F)J^{\mathrm{act}}_{\Lambda_k},
\,G-\chi_{\Lambda_k}(G)J^{\mathrm{act}}_{\Lambda_k}
\right\rangle^{\mathrm{aux}}_{\Lambda_k}
\label{eq:B-extension}
\end{equation}
defines a continuous symmetric positive definite bilinear form on \(E^-_{\Lambda_k}\) satisfying
\[
\mathsf B_{\Lambda_k}(J^{\mathrm{act}}_{\Lambda_k},J^{\mathrm{act}}_{\Lambda_k})
=
\widetilde\alpha(\Lambda_k),
\qquad
\mathsf B_{\Lambda_k}(J^{\mathrm{act}}_{\Lambda_k},T)=0
\quad
\forall T\in\mathcal T_{\Lambda_k}.
\]
\end{proposition}

\begin{proof}
By \cref{prop:T-kerrho}, every \(F\in E^-_{\Lambda_k}\) has a unique decomposition
\[
F=\chi_{\Lambda_k}(F)\,J^{\mathrm{act}}_{\Lambda_k}+T_F,
\qquad
T_F\in \mathcal T_{\Lambda_k}.
\]
Equation \eqref{eq:B-extension} is therefore well-defined, continuous, and symmetric.
Positivity is immediate from \(\widetilde\alpha(\Lambda_k)>0\) together with coercivity of
\(\langle\cdot,\cdot\rangle^{\mathrm{aux}}_{\Lambda_k}\) on \(\mathcal T_{\Lambda_k}\).
If \(\mathsf B_{\Lambda_k}(F,F)=0\), then both \(\chi_{\Lambda_k}(F)=0\) and
\(\langle T_F,T_F\rangle^{\mathrm{aux}}_{\Lambda_k}=0\), hence \(T_F=0\) and \(F=0\).
The displayed normalization and orthogonality properties follow by substitution into
\eqref{eq:B-extension}.
\end{proof}

\begin{corollary}[Descendant ambiguity compatibility on the characterization layer]
\label{cor:ambiguity-char}
For each accepted tick \(k\), the odd-projected ambiguity directions coming from
\[
Q_tX+Q_uY,
\qquad
\partial_\mu K^\mu,
\qquad
\Delta_{\rm ct}^{-}
\]
belong to \(\mathcal T^{\mathrm{op}}_{\Lambda_k}\) whenever \(K^\mu\) is taken from the
boundary-admissible class of \cref{prop:ambiguity-div-char} and \(\Delta_{\rm ct}^{-}\) is
admissible in the sense of \cref{def:ambiguity-ct-char}.
\end{corollary}

\begin{proof}
Combine \cref{prop:ambiguity-brst-char,prop:ambiguity-div-char,prop:ambiguity-ct-char}.
\end{proof}

\begin{theorem}[Strong rank-one realization of the characterization layer]
\label{thm:uom-char-export}
Assume
\cref{ass:export-selection,ass:centered-odd,ass:uv-match,ass:response-factorization,ass:response-sufficiency,ass:aux-char-core}.
Then for each accepted tick \(k\), the split/leak construction exports the UOM-side data
\begin{equation}
\big(
a_{\Lambda_k},\,
S_{\Lambda_k},\,
J^{\mathrm{act}}_{\Lambda_k},\,
C^{\mathrm{act}}_{\Lambda_k},\,
\widetilde\alpha(\Lambda_k),\,
\rho_{\Lambda_k}
\big),
\label{eq:uom-char-realization}
\end{equation}
where:
\begin{enumerate}[leftmargin=2em]
\item \(\mathcal O^{\mathrm{char}}_{\Lambda_k}=\mathrm{alg}\{a_{\Lambda_k},S_{\Lambda_k}\}\) is the
  minimal completion-sector characterization algebra;
\item the operationally trivial subspace satisfies
  \[
  \mathcal T_{\Lambda_k}:=\mathcal T^{\mathrm{op}}_{\Lambda_k}
  =
  \ker R^{\mathrm{phys}}_{\Lambda_k}
  =
  \ker\rho_{\Lambda_k};
  \]
\item the intrinsic physical response is one-dimensional on the strong sector:
  \[
  R^{\mathrm{phys}}_{\Lambda_k}(F)=\rho_{\Lambda_k}(F)\,C^{\mathrm{act}}_{\Lambda_k},
  \qquad
  R^{\mathrm{phys}}_{\Lambda_k}(J^{\mathrm{act}}_{\Lambda_k})=C^{\mathrm{act}}_{\Lambda_k}\neq 0;
  \]
\item the active odd quotient is rank one and is generated by
  \([J^{\mathrm{act}}_{\Lambda_k}]\);
\item the abstract orientation functional of the transported interface is realized by
  \[
  \chi_{\Lambda_k}:=\rho_{\Lambda_k};
  \]
\item any admissible calibration form \(\mathsf B_{\Lambda_k}\) in this strong package may be obtained by
  extending the active stiffness \(\widetilde\alpha(\Lambda_k)\) with
  \cref{prop:char-calibration-extension}.
\end{enumerate}
\end{theorem}

\begin{proof}
\Cref{def:char-data} packages the active completion-sector channel extracted from the split/leak
construction: \(a_{\Lambda_k}\) is the distinguished odd split coordinate singled out by
\cref{subsec:leak-solution,subsec:B-gaussian}, while \(S_{\Lambda_k}=a_{\Lambda_k}^2\) is the even
stock entering \eqref{eq:Qleak-effective} and \eqref{eq:split-threshold-effective}. The active odd
generator \(J^{\mathrm{act}}_{\Lambda_k}\), active physical covariance tangent
\(C^{\mathrm{act}}_{\Lambda_k}\), and scalar response coefficient \(\rho_{\Lambda_k}\) are
exported by \cref{ass:response-factorization}; \cref{prop:T-kerrho,cor:rank-one-char} then identify
the abstract interface-trivial subspace and the corresponding rank-one quotient.
The split/leak law fixes the active stiffness \(\widetilde\alpha(\Lambda_k)\), and
\cref{prop:char-calibration-extension} upgrades that one-line datum to an admissible full
calibration form once the auxiliary coercive core is chosen.
Finally, \cref{cor:ambiguity-char} places the odd-projected descendant ambiguities in
\(\mathcal T_{\Lambda_k}\), yielding the stated ambiguity compatibility on the characterization
layer.
\end{proof}

\begin{theorem}[Intrinsic weak transported descendant-class interface]
\label{thm:char-weak-interface}
Assume the hypotheses of \cref{thm:active-response-transport-minimal} on a connected accepted-branch
segment \(I\).
Then for every \(\Lambda\in I\), the sufficiently local odd generator layer carries the weak
transported interface data
\[
\big(
E^-_{\Lambda},\,
\mathcal A^-_{\Lambda},\,
[\Xi^{(1)}_{\Lambda}],\,
\chi_{\Lambda},\,
\varepsilon_{\Lambda}
\big)
\]
with
\[
\chi_{\Lambda}
:=
\nu_{\Lambda}^{-1}\rho^{\mathrm{dom}}_{\Lambda},
\]
and the following conclusions hold.
\begin{enumerate}[leftmargin=2em]
\item there exists a representative
  \[
  J^{\mathrm{desc}}_{\Lambda}\in[\Xi^{(1)}_{\Lambda}]
  \]
  such that
  \[
  \rho^{\mathrm{dom}}_{\Lambda}\!\big(J^{\mathrm{desc}}_{\Lambda}\big)\neq 0;
  \]
  equivalently, the transported descendant sector has nonzero overlap with the canonically selected
  compressed dominant odd line;
\item the orientation functional is realized intrinsically on the sufficiently local generator layer:
  \[
  \mathcal A^-_{\Lambda}\subset\ker\chi_{\Lambda},
  \qquad
  \chi_{\Lambda}(J)=\varepsilon_{\Lambda}
  \quad
  \forall J\in[\Xi^{(1)}_{\Lambda}];
  \]
\item along \(I\), the sign \(\varepsilon_{\Lambda}\) is locally constant.
\end{enumerate}
In particular, the weak transported interface is realized intrinsically, without appeal to completion
coordinates, under the stated hypotheses.
\end{theorem}

\begin{proof}
The functional \(\rho^{\mathrm{dom}}_{\Lambda}\) is defined intrinsically from the sufficiently local
generator layer by \cref{def:intrinsic-cov-response} together with the dominant-line prescription of
\cref{def:dominant-response-functional}.
Item (1) is exactly the nonzero-overlap conclusion of item (3) in
\cref{thm:active-response-transport-minimal}, item (2) is exactly item (4) there, and item (3) is
item (5).
\end{proof}

\begin{theorem}[Quadratic-probe realization of the weak activity supplement]
\label{thm:char-weak-activity}
Assume the hypotheses of \cref{thm:active-response-transport-minimal} on a connected accepted-branch
segment \(I\).
For each \(\Lambda\in I\), define the even quadratic probe sector by
\[
\mathcal P^+_{\Lambda}
:=
\operatorname{span}
\bigl\{
\mathcal Q^{\Lambda}_{f,g}
:
\text{\(f,g\) admissible band-limited receiver probes}
\bigr\}
\subset E^+_{\Lambda}.
\]
Then, for every \(F\in E^-_{\Lambda}\),
\[
\{F,A\}_{\mathrm P,\Lambda}=0
\qquad
\forall A\in\mathcal P^+_{\Lambda}
\]
implies
\[
\chi_{\Lambda}(F)=0.
\]
Equivalently, if \(\chi_{\Lambda}(F)\neq 0\), then there exists a sufficiently local even quadratic
probe \(A\in\mathcal P^+_{\Lambda}\) such that
\[
\{F,A\}_{\mathrm P,\Lambda}\neq 0.
\]
In particular, the weak interface data of \cref{thm:char-weak-interface} carry a canonical
activity witness sector built from the intrinsic quadratic receiver probes.
\end{theorem}

\begin{proof}
Assume
\[
\{F,A\}_{\mathrm P,\Lambda}=0
\qquad
\forall A\in\mathcal P^+_{\Lambda}.
\]
In particular,
\[
\{F,\mathcal Q^{\Lambda}_{f,g}\}_{\mathrm P,\Lambda}=0
\]
for every admissible pair \(f,g\).
By \cref{def:intrinsic-cov-response}, this means
\[
\big\langle f,\mathcal R_{\Lambda}(F)g\big\rangle
=
\Big\langle
\{F,\mathcal Q^{\Lambda}_{f,g}\}_{\mathrm P,\Lambda}
\Big\rangle_{\rho_{\ast,\Lambda}}
=
0
\]
for all admissible \(f,g\), hence
\(
\mathcal R_{\Lambda}(F)=0
\).
Therefore
\[
\Pi_{\mathrm{phys},\Lambda}\mathcal R_{\Lambda}(F)=0,
\]
and the definition of the raw dominant-response coefficient in
\cref{def:dominant-response-functional} gives
\[
\widetilde\rho^{\mathrm{dom}}_{\Lambda}(F)=0.
\]
By item (1) of \cref{thm:active-response-transport-minimal},
\(
\delta C^{\mathrm{LoG}}_{\Lambda}\neq 0
\),
so dominant-line realization holds and \(\rho^{\mathrm{dom}}_{\Lambda}\) is well defined.
Hence
\[
\rho^{\mathrm{dom}}_{\Lambda}(F)=0.
\]
Now apply \cref{thm:char-weak-interface}, which identifies
\(
\chi_{\Lambda}
=
\nu_{\Lambda}^{-1}\rho^{\mathrm{dom}}_{\Lambda}
\),
to obtain
\[
\chi_{\Lambda}(F)=0.
\]
The contrapositive gives the equivalent detection statement.
\end{proof}

\begin{remark}[Correct scalar form of representative comparisons]
\label{rem:char-weak-interface-rho-first}
The weak export theorem should be read \(\rho^{\mathrm{dom}}\)-first on the ambiguity quotient.
Because
\[
\chi_{\Lambda}
:=
\nu_{\Lambda}^{-1}\rho^{\mathrm{dom}}_{\Lambda},
\]
any comparison statement that depends only on dominant pairing is naturally a statement about the
descended scalar
\[
\overline{\rho}^{\mathrm{dom}}_{\Lambda}
:
E^-_{\Lambda}/\mathcal A^-_{\Lambda}\to\mathbb R
\]
or, equivalently, about its value on the transported descendant affine class.
Such a statement may be rewritten in \(\chi_{\Lambda}\) form only when the normalization is made
explicit; for example, if \(\nu_{\Lambda}=1\), then the same scalar value is represented by a
descended \(\bar\chi_{\Lambda}\).
Without that normalization the stronger unscaled \(\chi\)-compatibility form is false in general.
\end{remark}

\begin{corollary}[UOM-stiffness specialization of the weak transported calibration]
\label{cor:char-weak-calibration}
Assume \cref{ass:uv-match} and fix an accepted tick \(k\) for which the conclusions of
\cref{thm:char-weak-interface} hold.
Suppose \(\ker\chi_{\Lambda_k}\) admits a continuous symmetric coercive bilinear form whose induced
norm Hilbertizes it, and choose any
\[
J^\sharp_{\Lambda_k}\in E^-_{\Lambda_k}
\qquad
\text{with}
\qquad
\chi_{\Lambda_k}(J^\sharp_{\Lambda_k})=1.
\]
Then \cref{prop:minimal-transport-calibration} may be applied with
\[
\beta_{\Lambda_k}:=\widetilde\alpha(\Lambda_k)
\]
to produce a continuous symmetric positive definite bilinear form
\[
\mathsf B^{\widetilde\alpha}_{\Lambda_k}
:
E^-_{\Lambda_k}\times E^-_{\Lambda_k}\to\mathbb R
\]
such that
\[
\mathsf B^{\widetilde\alpha}_{\Lambda_k}(J^\sharp_{\Lambda_k},J^\sharp_{\Lambda_k})
=
\widetilde\alpha(\Lambda_k),
\qquad
\mathsf B^{\widetilde\alpha}_{\Lambda_k}(J^\sharp_{\Lambda_k},K)=0
\quad
\forall K\in\ker\chi_{\Lambda_k}.
\]
In particular,
\[
\mathcal A^-_{\Lambda_k}\subset\ker\chi_{\Lambda_k}
\]
implies
\[
\mathsf B^{\widetilde\alpha}_{\Lambda_k}(J^\sharp_{\Lambda_k},A)=0
\qquad
\forall A\in\mathcal A^-_{\Lambda_k}.
\]
\end{corollary}

\begin{proof}
By \cref{rem:char-calibration-split}, the split/leak package fixes the positive stiffness scale
\(\widetilde\alpha(\Lambda_k)\), and \cref{ass:uv-match} records its positivity and asymptotic
size.
Apply \cref{prop:minimal-transport-calibration} with
\(\beta_{\Lambda_k}=\widetilde\alpha(\Lambda_k)\).
The final orthogonality statement follows from item (2) of \cref{thm:char-weak-interface}.
\end{proof}

\begin{remark}[Inherited ledger compatibility]
\label{rem:ledger-compat}
The inherited commutative ledger algebra is intentionally not part of the core characterization
algebra \(\mathcal O^{\mathrm{char}}_{\Lambda_k}\).
Its role is separate: the split model requires the active odd representative to be completion-sector
nontrivial while acting trivially, to leading order, on the inherited classical ledger.
This is the UOM-side compatibility statement behind the slogan
``oddness in the completion sector, classicality in the accessible ledger.''
\end{remark}

\ifdefined\UOMTransportFragmentLoaded\else
\end{document}
\fi
%
  }{%
    \ifdefined\UOMTransportFragmentLoaded\else
\documentclass[11pt]{article}
\usepackage[margin=1in]{geometry}
\usepackage{amsmath,amssymb,amsthm,mathtools}
\usepackage{bm}
\usepackage{microtype}
\usepackage{enumitem}
\usepackage{booktabs}
\usepackage{array}
\usepackage{xcolor}
\usepackage{xr-hyper}
\usepackage{hyperref}
\usepackage[nameinlink,noabbrev]{cleveref}

\providecommand{\Fsloc}{\mathcal F_{\rm s.loc}}
\providecommand{\Qt}{Q_t}
\externaldocument[main-]{UOM/main}
\externaldocument[uomdesc-]{UOM/uom_descendant_note}
\externaldocument{UOM/completion_sector_split_model_note}

\hypersetup{
  colorlinks=true,
  linkcolor=blue!60!black,
  citecolor=blue!60!black,
  urlcolor=blue!60!black
}

\theoremstyle{definition}
\newtheorem{definition}{Definition}[section]
\newtheorem{assumption}{Assumption}[section]
\theoremstyle{plain}
\newtheorem{theorem}[definition]{Theorem}
\newtheorem{lemma}[definition]{Lemma}
\newtheorem{proposition}[definition]{Proposition}
\newtheorem{corollary}[definition]{Corollary}
\theoremstyle{remark}
\newtheorem{remark}[definition]{Remark}
\crefname{definition}{definition}{definitions}
\Crefname{definition}{Definition}{Definitions}
\crefname{assumption}{assumption}{assumptions}
\Crefname{assumption}{Assumption}{Assumptions}
\crefname{theorem}{theorem}{theorems}
\Crefname{theorem}{Theorem}{Theorems}
\crefname{lemma}{lemma}{lemmas}
\Crefname{lemma}{Lemma}{Lemmas}
\crefname{proposition}{proposition}{propositions}
\Crefname{proposition}{Proposition}{Propositions}
\crefname{corollary}{corollary}{corollaries}
\Crefname{corollary}{Corollary}{Corollaries}
\crefname{remark}{remark}{remarks}
\Crefname{remark}{Remark}{Remarks}

\newcommand{\cH}{\mathcal H}
\newcommand{\cA}{\mathcal A}
\newcommand{\cD}{\mathcal D}
\newcommand{\cB}{\mathcal B}
\newcommand{\cC}{\mathcal C}
\newcommand{\cL}{\mathcal L}
\newcommand{\cK}{\mathcal K}
\newcommand{\cR}{\mathcal R}
\newcommand{\cG}{\mathcal G}
\newcommand{\Ztwo}{\mathbb Z_2}
\newcommand{\Sys}{\mathrm{Sys}}
\newcommand{\Time}{\mathrm{Time}}
\newcommand{\RG}{\mathrm{RG}}
\newcommand{\Tr}{\mathrm{Tr}}
\newcommand{\diag}{\mathrm{diag}}
\newcommand{\ad}{\mathrm{ad}}
\newcommand{\id}{\mathrm{id}}
\newcommand{\norm}[1]{\left\lVert #1\right\rVert}
\newcommand{\abs}[1]{\left\lvert #1\right\rvert}
\newcommand{\ip}[2]{\left\langle #1,#2\right\rangle}

\externaldocument{UOM/compressed_log_spectral_selection_note}
\externaldocument{UOM/odd_selector_note}
\externaldocument{UOM/active_response_transport_program_note}
\title{Polarized Characterization Interface}
\author{}
\date{}
\begin{document}
\maketitle
\fi

\section{Polarized Characterization Interface}
\label{sec:char-export}

This section records the finite-cutoff characterization data extracted from the split/leak
construction and used in the 2T Main assumption package
Definition~\ref{main-def:uom-main-theorem-assumptions}. In this note, an ``accepted tick'' denotes the local
accepted-branch slice at cutoff \(\Lambda_k\).

For each accepted tick \(k\), write
\[
a_{\Lambda_k}:=a_k,
\qquad
S_{\Lambda_k}:=a_{\Lambda_k}^2.
\]
The variable \(a_{\Lambda_k}\) is the distinguished odd split coordinate identified in
\cref{subsec:leak-solution,subsec:B-gaussian}; the even quantity \(S_{\Lambda_k}\) is the
positive stock that enters the leak law and the split threshold.

\begin{definition}[Characterization algebra, operational triviality, and even null ideal]
\label{def:char-data}
Fix an accepted tick \(k\) and let \(E^-_{\Lambda_k}\) denote the physically odd sufficiently local
generator space at cutoff \(\Lambda_k\), i.e.\ the odd sector for the transported physical involution
\(\Gamma\). Let \(E^+_{\Lambda_k}\) denote the corresponding physically even sufficiently local
observable space.
Define the core characterization algebra by
\begin{equation}
\mathcal O^{\mathrm{char}}_{\Lambda_k}
:=
\mathrm{alg}\{a_{\Lambda_k},\,S_{\Lambda_k}\}.
\label{eq:Ochar-def}
\end{equation}
Define the operationally interface-trivial odd subspace by
\begin{equation}
\mathcal T^{\mathrm{op}}_{\Lambda_k}
:=
\overline{\left\{
T\in E^-_{\Lambda_k}:
\{T,O\}_{\mathrm P,\Lambda_k}=0
\ \forall O\in\mathcal O^{\mathrm{char}}_{\Lambda_k}
\right\}}.
\label{eq:Tchar-def}
\end{equation}
Define the even interface-null ideal by
\begin{equation}
\mathcal N^{\mathrm{even}}_{\Lambda_k}
:=
\overline{\left\{
E\in E^+_{\Lambda_k}:
\{E,O\}_{\mathrm P,\Lambda_k}=0
\ \forall O\in\mathcal O^{\mathrm{char}}_{\Lambda_k}
\right\}}.
\label{eq:Neven-def}
\end{equation}
\end{definition}

\begin{remark}[Do not identify the characterization observable with the generator]
\label{rem:a-vs-Jact}
The odd observable \(a_{\Lambda_k}\) characterizes the active completion-sector channel, but it is
not itself the odd Hamiltonian generator class appearing in item \textbf{(M4)} of
Definition~\ref{main-def:uom-main-theorem-assumptions}.
That generator is denoted \(J^{\mathrm{act}}_{\Lambda_k}\in E^-_{\Lambda_k}\) below.
The quotient statement belongs to \([J^{\mathrm{act}}_{\Lambda_k}]\), not to
\([a_{\Lambda_k}]\).
\end{remark}

\begin{definition}[Intrinsic covariance-tangent response map]
\label{def:intrinsic-cov-response}
Fix an accepted tick \(k\). Let \(\mathcal T^{\mathrm{cov}}_{\Lambda_k}\) denote the finite-cutoff
receiver covariance-tangent space at the UOM reference Gaussian state
\(\rho_{\ast,\Lambda_k}\), equipped with the Fisher pairing
\[
\big\langle\!\big\langle \delta C,\delta C'\big\rangle\!\big\rangle_{C_{\ast,\Lambda_k}}
:=
\tfrac12\,\mathrm{tr}\!\big(
C_{\ast,\Lambda_k}^{-1}\delta C\,C_{\ast,\Lambda_k}^{-1}\delta C'
\big).
\]
For admissible band-limited receiver probes \(f,g\), let
\[
\mathcal Q^{\Lambda_k}_{f,g}\in E^+_{\Lambda_k}
\]
denote the corresponding sufficiently local quadratic observable, normalized so that on centered
Gaussian states with covariance \(C\),
\[
\big\langle \mathcal Q^{\Lambda_k}_{f,g}\big\rangle_{\rho(C)}
=
\langle f,Cg\rangle.
\]
For \(F\in E^-_{\Lambda_k}\), define \(\mathcal R_{\Lambda_k}(F)\in\mathcal T^{\mathrm{cov}}_{\Lambda_k}\)
by the matrix elements
\[
\big\langle f,\mathcal R_{\Lambda_k}(F)g\big\rangle
:=
\Big\langle
\{F,\mathcal Q^{\Lambda_k}_{f,g}\}_{\mathrm P,\Lambda_k}
\Big\rangle_{\rho_{\ast,\Lambda_k}}.
\]
We call
\[
\mathcal R_{\Lambda_k}:E^-_{\Lambda_k}\to\mathcal T^{\mathrm{cov}}_{\Lambda_k}
\]
the \emph{intrinsic covariance-tangent response map}.
\end{definition}

\begin{proposition}[Intrinsic definition and completion independence]
\label{prop:intrinsic-cov-response}
The map \(\mathcal R_{\Lambda_k}\) of \cref{def:intrinsic-cov-response} is continuous and linear.
It depends only on the finite-cutoff UOM Peierls package and the UOM covariance/Fisher package.
If a Gaussian completion realization is chosen, then \(\mathcal R_{\Lambda_k}(F)\) is represented by
the same completion-sector covariance tangent that one obtains by linearizing the completion
quadratures. Thus Gaussian completion is a realization of the response map, not extra structure
required to define it.
\end{proposition}

\begin{proof}
For fixed \(F\in E^-_{\Lambda_k}\), the finite-cutoff UOM Peierls package makes
\[
\delta_F:=\{F,-\}_{\mathrm P,\Lambda_k}
\]
a continuous derivation on the sufficiently local algebra. Applying \(\delta_F\) to the quadratic
probes \(\mathcal Q^{\Lambda_k}_{f,g}\) and evaluating in the reference Gaussian state produces a
continuous symmetric bilinear form in \(f,g\). In the UOM Gaussian radiative sector, such bilinear
forms are exactly covariance tangents at \(C_{\ast,\Lambda_k}\), equipped with the Fisher metric.
Linearity and continuity in \(F\) are inherited directly from the Peierls bracket. No completion
coordinates are used in this construction.

If a Gaussian completion is later chosen, the same second moments are encoded by the covariance
matrix of the completion quadratures. The bilinear form defined above and the completion-sector
covariance tangent therefore represent the same infinitesimal response object in two different
coordinate systems.
\end{proof}

\begin{definition}[Physical intrinsic covariance-response map]
\label{def:phys-cov-response}
Fix an accepted tick \(k\). Let
\[
\mathcal T^{\mathrm{phys}}_{\Lambda_k}
:=
\Pi_{\mathrm{phys},\Lambda_k}\mathcal T^{\mathrm{cov}}_{\Lambda_k}
\subset \mathcal T^{\mathrm{cov}}_{\Lambda_k}.
\]
Define the physical intrinsic covariance-response map by
\[
R^{\mathrm{phys}}_{\Lambda_k}
:=
\Pi_{\mathrm{phys},\Lambda_k}\mathcal R_{\Lambda_k}
:
E^-_{\Lambda_k}\to \mathcal T^{\mathrm{phys}}_{\Lambda_k}.
\]
\end{definition}

\begin{proposition}[Physical intrinsic response map]
\label{prop:phys-cov-response}
The map \(R^{\mathrm{phys}}_{\Lambda_k}\) of \cref{def:phys-cov-response} is continuous and linear.
\end{proposition}

\begin{proof}
This is immediate from \cref{prop:intrinsic-cov-response} and the continuity and linearity of the
finite-cutoff physical projector \(\Pi_{\mathrm{phys},\Lambda_k}\).
\end{proof}

\begin{proposition}[Reduction of operational triviality to the active odd coordinate]
\label{prop:char-op-from-a}
Fix an accepted tick \(k\). If \(F\in E^-_{\Lambda_k}\) satisfies
\[
\{F,a_{\Lambda_k}\}_{\mathrm P,\Lambda_k}=0,
\]
then \(F\in\mathcal T^{\mathrm{op}}_{\Lambda_k}\).
\end{proposition}

\begin{proof}
Since \(S_{\Lambda_k}=a_{\Lambda_k}^2\) and the Peierls bracket is a derivation,
\[
\{F,S_{\Lambda_k}\}_{\mathrm P,\Lambda_k}
=
\{F,a_{\Lambda_k}^2\}_{\mathrm P,\Lambda_k}
=
2\,a_{\Lambda_k}\,\{F,a_{\Lambda_k}\}_{\mathrm P,\Lambda_k}
=
0.
\]
The algebra \(\mathcal O^{\mathrm{char}}_{\Lambda_k}=\mathrm{alg}\{a_{\Lambda_k},S_{\Lambda_k}\}\)
is generated by \(a_{\Lambda_k}\) and \(S_{\Lambda_k}\), and the Peierls bracket acts on it by the
Leibniz rule. Hence \(\{F,O\}_{\mathrm P,\Lambda_k}=0\) for every
\(O\in\mathcal O^{\mathrm{char}}_{\Lambda_k}\), proving
\(F\in\mathcal T^{\mathrm{op}}_{\Lambda_k}\).
\end{proof}

\begin{assumption}[Active odd generator and intrinsic-response factorization]
\label{ass:response-factorization}
For each accepted tick \(k\), there exist a distinguished active odd generator
\[
J^{\mathrm{act}}_{\Lambda_k}\in E^-_{\Lambda_k}
\]
and a nonzero physical covariance tangent
\[
C^{\mathrm{act}}_{\Lambda_k}\in \mathcal T^{\mathrm{phys}}_{\Lambda_k}
\]
and a continuous nonzero linear functional
\[
\rho_{\Lambda_k}:E^-_{\Lambda_k}\to\mathbb R
\]
normalized by
\[
\rho_{\Lambda_k}(J^{\mathrm{act}}_{\Lambda_k})=1,
\]
such that
\[
R^{\mathrm{phys}}_{\Lambda_k}(J^{\mathrm{act}}_{\Lambda_k})
=
C^{\mathrm{act}}_{\Lambda_k}
\neq 0
\]
and
\begin{equation}
R^{\mathrm{phys}}_{\Lambda_k}(F)
=
\rho_{\Lambda_k}(F)\,C^{\mathrm{act}}_{\Lambda_k}
\qquad
\forall F\in E^-_{\Lambda_k}
\label{eq:rho-factorization}
\end{equation}
holds on the strong sector under consideration.
\end{assumption}

\begin{remark}[Strong normalization]
\label{rem:rho-strong}
If one further chooses the Fisher-unit normalization
\(
\widehat C^{\mathrm{act}}_{\Lambda_k}
:=
\|C^{\mathrm{act}}_{\Lambda_k}\|^{-1}_{C_{\ast,\Lambda_k}}\,C^{\mathrm{act}}_{\Lambda_k},
\)
then
\eqref{eq:rho-factorization} sharpens to
\[
R^{\mathrm{phys}}_{\Lambda_k}(F)
=
\widehat{\rho}_{\Lambda_k}(F)\,\widehat C^{\mathrm{act}}_{\Lambda_k}
\]
with
\(
\widehat{\rho}_{\Lambda_k}
:=
\|C^{\mathrm{act}}_{\Lambda_k}\|_{C_{\ast,\Lambda_k}}\,\rho_{\Lambda_k}.
\)
The export below does not require this extra normalization; it only needs the scalar response
coefficient \(\rho_{\Lambda_k}\) relative to a fixed nonzero active tangent
\(C^{\mathrm{act}}_{\Lambda_k}\).
\end{remark}

\begin{assumption}[Physical-response sufficiency]
\label{ass:response-sufficiency}
For each accepted tick \(k\), the operationally trivial odd directions are exactly the kernel of the
physical intrinsic covariance-response map:
\[
 \ker R^{\mathrm{phys}}_{\Lambda_k}
 =
 \mathcal T^{\mathrm{op}}_{\Lambda_k}.
\]
Equivalently, an odd generator has zero physical intrinsic covariance response if and only if it is
operationally trivial on the full core characterization algebra \(\mathcal O^{\mathrm{char}}_{\Lambda_k}\).
\end{assumption}

\begin{remark}[Why the strong package must be stated intrinsically]
\label{rem:response-quotient-obstruction}
The earlier quotient-valued formulation of the strong package in
\(
E^+_{\Lambda_k}/\mathcal N^{\mathrm{even}}_{\Lambda_k}
\)
is not stable.
Because the Peierls bracket here is the ordinary derivation bracket and \(S_{\Lambda_k}=a_{\Lambda_k}^2\),
the even observable \(\{F,a_{\Lambda_k}\}_{\mathrm P,\Lambda_k}\) already Poisson-commutes with
both \(a_{\Lambda_k}\) and \(S_{\Lambda_k}\), hence lies automatically in
\(\mathcal N^{\mathrm{even}}_{\Lambda_k}\).
The correct codomain for factorization and kernel identification is therefore the intrinsic
physical covariance-response space \(\mathcal T^{\mathrm{phys}}_{\Lambda_k}\), not the even
observable quotient.
\end{remark}

\begin{proposition}[Operational triviality is the kernel of the active response]
\label{prop:T-kerrho}
Assume \cref{ass:response-factorization,ass:response-sufficiency}. Then
\[
\mathcal T^{\mathrm{op}}_{\Lambda_k}=\ker\rho_{\Lambda_k}.
\]
\end{proposition}

\begin{proof}
Because \(C^{\mathrm{act}}_{\Lambda_k}\neq 0\), the factorization identity
\eqref{eq:rho-factorization} implies
\[
R^{\mathrm{phys}}_{\Lambda_k}(F)=0
\qquad\Longleftrightarrow\qquad
\rho_{\Lambda_k}(F)=0
\]
for every \(F\in E^-_{\Lambda_k}\).
Thus
\(
\ker R^{\mathrm{phys}}_{\Lambda_k}
=
\ker\rho_{\Lambda_k}.
\)
Apply \cref{ass:response-sufficiency}.
\end{proof}

\begin{corollary}[Rank-one active odd quotient]
\label{cor:rank-one-char}
Assume \cref{ass:response-factorization,ass:response-sufficiency}. Then
\[
E^-_{\Lambda_k}/\mathcal T^{\mathrm{op}}_{\Lambda_k}\cong \mathbb R,
\]
and the quotient is generated by the class of the active odd generator
\([J^{\mathrm{act}}_{\Lambda_k}]\).
\end{corollary}

\begin{proof}
By \cref{prop:T-kerrho}, the quotient is canonically isomorphic to the image of the nonzero
functional \(\rho_{\Lambda_k}\). Since \(\rho_{\Lambda_k}\) is real-valued and
\(\rho_{\Lambda_k}(J^{\mathrm{act}}_{\Lambda_k})=1\), its image is \(\mathbb R\) and the quotient
is generated by \([J^{\mathrm{act}}_{\Lambda_k}]\).
\end{proof}

\begin{proposition}[BRST/gauge-exact compatibility on the characterization layer]
\label{prop:ambiguity-brst-char}
For each accepted tick \(k\), the odd parts of the BRST/gauge-exact directions
\[
Q_tX+Q_uY
\]
lie in \(\mathcal T^{\mathrm{op}}_{\Lambda_k}\).
\end{proposition}

\begin{proof}
The exported characterization observables \(a_{\Lambda_k}\) and \(S_{\Lambda_k}\) are physical
finite-cutoff observables on the characterization layer. The BRST/gauge-exact odd directions
\(Q_tX+Q_uY\) represent redundancies of that physical layer, so their Peierls action vanishes on
\(\mathcal O^{\mathrm{char}}_{\Lambda_k}\). By \eqref{eq:Tchar-def}, their odd parts therefore lie
in \(\mathcal T^{\mathrm{op}}_{\Lambda_k}\).
\end{proof}

\begin{proposition}[Boundary-admissible divergence compatibility]
\label{prop:ambiguity-div-char}
Let \(K^\mu\) be a sufficiently local current whose renormalized boundary flux vanishes on the core
characterization algebra \(\mathcal O^{\mathrm{char}}_{\Lambda_k}\). Then the odd part of
\[
\partial_\mu K^\mu
\]
lies in \(\mathcal T^{\mathrm{op}}_{\Lambda_k}\).
\end{proposition}

\begin{proof}
For the admissible currents in the statement, the renormalized Stokes reduction of
\(\partial_\mu K^\mu\) produces no boundary term on the characterization layer. Hence its Peierls
action on every \(O\in\mathcal O^{\mathrm{char}}_{\Lambda_k}\) vanishes, and the odd part belongs
to \(\mathcal T^{\mathrm{op}}_{\Lambda_k}\) by \eqref{eq:Tchar-def}.
\end{proof}

\begin{definition}[Admissible odd counterterms on the characterization layer]
\label{def:ambiguity-ct-char}
For each accepted tick \(k\), an odd sufficiently local counterterm ambiguity
\[
\Delta_{\rm ct}^{-}\in E^-_{\Lambda_k}
\]
is called \emph{admissible on the characterization layer} if its intrinsic covariance response has
no physical projection:
\[
\Pi_{\mathrm{phys},\Lambda_k}\,\mathcal R_{\Lambda_k}(\Delta_{\rm ct}^{-})=0.
\]
Equivalently, admissible odd counterterms are precisely those local scheme shifts whose response is
confined to the RP, Weyl, conformal, or other even-null monitor sectors and does not enter the
compressed active odd line.
\end{definition}

\begin{proposition}[Counterterm rigidity on the characterization layer]
\label{prop:ambiguity-ct-char}
Assume that the characterization observable \(a_{\Lambda_k}\) is extracted from the physical
covariance response so that
\[
\Pi_{\mathrm{phys},\Lambda_k}\mathcal R_{\Lambda_k}(F)=0
\quad\Longrightarrow\quad
\{F,a_{\Lambda_k}\}_{\mathrm P,\Lambda_k}=0
\qquad
\forall F\in E^-_{\Lambda_k}.
\]
Then every admissible odd counterterm ambiguity in the sense of \cref{def:ambiguity-ct-char} acts
trivially on the core characterization algebra:
\[
\Delta_{\rm ct}^{-}\in \mathcal T^{\mathrm{op}}_{\Lambda_k}.
\]
\end{proposition}

\begin{proof}
If \(F=\Delta_{\rm ct}^{-}\) is admissible, then this physical projection vanishes by definition.
The displayed hypothesis therefore gives
\[
\{\Delta_{\rm ct}^{-},a_{\Lambda_k}\}_{\mathrm P,\Lambda_k}=0,
\]
and since \(S_{\Lambda_k}=a_{\Lambda_k}^2\), the derivation property yields
\[
\{\Delta_{\rm ct}^{-},S_{\Lambda_k}\}_{\mathrm P,\Lambda_k}=0.
\]
By \cref{prop:char-op-from-a}, the first identity already implies
\(\Delta_{\rm ct}^{-}\in\mathcal T^{\mathrm{op}}_{\Lambda_k}\).
\end{proof}

\begin{remark}[Active stiffness versus full calibration]
\label{rem:char-calibration-split}
The split/leak analysis determines the active stiffness \(\widetilde\alpha(\Lambda_k)\) on the
distinguished line generated by \(J^{\mathrm{act}}_{\Lambda_k}\), with
\[
\widetilde\alpha(\Lambda_k)\asymp \alpha(\Lambda_k)\sim \Lambda_k^{d+2}
\]
by \cref{ass:uv-match}. What UOM does \emph{not} determine canonically is a positive bilinear form
on the whole odd generator space \(E^-_{\Lambda_k}\). That extension is auxiliary interface data,
not a UOM-dynamical output.
\end{remark}

\begin{assumption}[Auxiliary coercive core on the operationally trivial sector]
\label{ass:aux-char-core}
For each accepted tick \(k\), choose a continuous symmetric coercive bilinear form
\[
\langle\cdot,\cdot\rangle^{\mathrm{aux}}_{\Lambda_k}
:
\mathcal T^{\mathrm{op}}_{\Lambda_k}\times \mathcal T^{\mathrm{op}}_{\Lambda_k}\to\mathbb R
\]
whose induced norm Hilbertizes \(\mathcal T^{\mathrm{op}}_{\Lambda_k}\).
\end{assumption}

\begin{proposition}[Strong-package calibration extension from active stiffness]
\label{prop:char-calibration-extension}
Assume \cref{ass:response-factorization,ass:response-sufficiency,ass:aux-char-core}. Set
\[
\chi_{\Lambda_k}:=\rho_{\Lambda_k},
\qquad
\mathcal T_{\Lambda_k}:=\mathcal T^{\mathrm{op}}_{\Lambda_k}.
\]
Then the formula
\begin{equation}
\mathsf B_{\Lambda_k}(F,G)
:=
\widetilde\alpha(\Lambda_k)\,\chi_{\Lambda_k}(F)\chi_{\Lambda_k}(G)
+
\left\langle
F-\chi_{\Lambda_k}(F)J^{\mathrm{act}}_{\Lambda_k},
\,G-\chi_{\Lambda_k}(G)J^{\mathrm{act}}_{\Lambda_k}
\right\rangle^{\mathrm{aux}}_{\Lambda_k}
\label{eq:B-extension}
\end{equation}
defines a continuous symmetric positive definite bilinear form on \(E^-_{\Lambda_k}\) satisfying
\[
\mathsf B_{\Lambda_k}(J^{\mathrm{act}}_{\Lambda_k},J^{\mathrm{act}}_{\Lambda_k})
=
\widetilde\alpha(\Lambda_k),
\qquad
\mathsf B_{\Lambda_k}(J^{\mathrm{act}}_{\Lambda_k},T)=0
\quad
\forall T\in\mathcal T_{\Lambda_k}.
\]
\end{proposition}

\begin{proof}
By \cref{prop:T-kerrho}, every \(F\in E^-_{\Lambda_k}\) has a unique decomposition
\[
F=\chi_{\Lambda_k}(F)\,J^{\mathrm{act}}_{\Lambda_k}+T_F,
\qquad
T_F\in \mathcal T_{\Lambda_k}.
\]
Equation \eqref{eq:B-extension} is therefore well-defined, continuous, and symmetric.
Positivity is immediate from \(\widetilde\alpha(\Lambda_k)>0\) together with coercivity of
\(\langle\cdot,\cdot\rangle^{\mathrm{aux}}_{\Lambda_k}\) on \(\mathcal T_{\Lambda_k}\).
If \(\mathsf B_{\Lambda_k}(F,F)=0\), then both \(\chi_{\Lambda_k}(F)=0\) and
\(\langle T_F,T_F\rangle^{\mathrm{aux}}_{\Lambda_k}=0\), hence \(T_F=0\) and \(F=0\).
The displayed normalization and orthogonality properties follow by substitution into
\eqref{eq:B-extension}.
\end{proof}

\begin{corollary}[Descendant ambiguity compatibility on the characterization layer]
\label{cor:ambiguity-char}
For each accepted tick \(k\), the odd-projected ambiguity directions coming from
\[
Q_tX+Q_uY,
\qquad
\partial_\mu K^\mu,
\qquad
\Delta_{\rm ct}^{-}
\]
belong to \(\mathcal T^{\mathrm{op}}_{\Lambda_k}\) whenever \(K^\mu\) is taken from the
boundary-admissible class of \cref{prop:ambiguity-div-char} and \(\Delta_{\rm ct}^{-}\) is
admissible in the sense of \cref{def:ambiguity-ct-char}.
\end{corollary}

\begin{proof}
Combine \cref{prop:ambiguity-brst-char,prop:ambiguity-div-char,prop:ambiguity-ct-char}.
\end{proof}

\begin{theorem}[Strong rank-one realization of the characterization layer]
\label{thm:uom-char-export}
Assume
\cref{ass:export-selection,ass:centered-odd,ass:uv-match,ass:response-factorization,ass:response-sufficiency,ass:aux-char-core}.
Then for each accepted tick \(k\), the split/leak construction exports the UOM-side data
\begin{equation}
\big(
a_{\Lambda_k},\,
S_{\Lambda_k},\,
J^{\mathrm{act}}_{\Lambda_k},\,
C^{\mathrm{act}}_{\Lambda_k},\,
\widetilde\alpha(\Lambda_k),\,
\rho_{\Lambda_k}
\big),
\label{eq:uom-char-realization}
\end{equation}
where:
\begin{enumerate}[leftmargin=2em]
\item \(\mathcal O^{\mathrm{char}}_{\Lambda_k}=\mathrm{alg}\{a_{\Lambda_k},S_{\Lambda_k}\}\) is the
  minimal completion-sector characterization algebra;
\item the operationally trivial subspace satisfies
  \[
  \mathcal T_{\Lambda_k}:=\mathcal T^{\mathrm{op}}_{\Lambda_k}
  =
  \ker R^{\mathrm{phys}}_{\Lambda_k}
  =
  \ker\rho_{\Lambda_k};
  \]
\item the intrinsic physical response is one-dimensional on the strong sector:
  \[
  R^{\mathrm{phys}}_{\Lambda_k}(F)=\rho_{\Lambda_k}(F)\,C^{\mathrm{act}}_{\Lambda_k},
  \qquad
  R^{\mathrm{phys}}_{\Lambda_k}(J^{\mathrm{act}}_{\Lambda_k})=C^{\mathrm{act}}_{\Lambda_k}\neq 0;
  \]
\item the active odd quotient is rank one and is generated by
  \([J^{\mathrm{act}}_{\Lambda_k}]\);
\item the abstract orientation functional of the transported interface is realized by
  \[
  \chi_{\Lambda_k}:=\rho_{\Lambda_k};
  \]
\item any admissible calibration form \(\mathsf B_{\Lambda_k}\) in this strong package may be obtained by
  extending the active stiffness \(\widetilde\alpha(\Lambda_k)\) with
  \cref{prop:char-calibration-extension}.
\end{enumerate}
\end{theorem}

\begin{proof}
\Cref{def:char-data} packages the active completion-sector channel extracted from the split/leak
construction: \(a_{\Lambda_k}\) is the distinguished odd split coordinate singled out by
\cref{subsec:leak-solution,subsec:B-gaussian}, while \(S_{\Lambda_k}=a_{\Lambda_k}^2\) is the even
stock entering \eqref{eq:Qleak-effective} and \eqref{eq:split-threshold-effective}. The active odd
generator \(J^{\mathrm{act}}_{\Lambda_k}\), active physical covariance tangent
\(C^{\mathrm{act}}_{\Lambda_k}\), and scalar response coefficient \(\rho_{\Lambda_k}\) are
exported by \cref{ass:response-factorization}; \cref{prop:T-kerrho,cor:rank-one-char} then identify
the abstract interface-trivial subspace and the corresponding rank-one quotient.
The split/leak law fixes the active stiffness \(\widetilde\alpha(\Lambda_k)\), and
\cref{prop:char-calibration-extension} upgrades that one-line datum to an admissible full
calibration form once the auxiliary coercive core is chosen.
Finally, \cref{cor:ambiguity-char} places the odd-projected descendant ambiguities in
\(\mathcal T_{\Lambda_k}\), yielding the stated ambiguity compatibility on the characterization
layer.
\end{proof}

\begin{theorem}[Intrinsic weak transported descendant-class interface]
\label{thm:char-weak-interface}
Assume the hypotheses of \cref{thm:active-response-transport-minimal} on a connected accepted-branch
segment \(I\).
Then for every \(\Lambda\in I\), the sufficiently local odd generator layer carries the weak
transported interface data
\[
\big(
E^-_{\Lambda},\,
\mathcal A^-_{\Lambda},\,
[\Xi^{(1)}_{\Lambda}],\,
\chi_{\Lambda},\,
\varepsilon_{\Lambda}
\big)
\]
with
\[
\chi_{\Lambda}
:=
\nu_{\Lambda}^{-1}\rho^{\mathrm{dom}}_{\Lambda},
\]
and the following conclusions hold.
\begin{enumerate}[leftmargin=2em]
\item there exists a representative
  \[
  J^{\mathrm{desc}}_{\Lambda}\in[\Xi^{(1)}_{\Lambda}]
  \]
  such that
  \[
  \rho^{\mathrm{dom}}_{\Lambda}\!\big(J^{\mathrm{desc}}_{\Lambda}\big)\neq 0;
  \]
  equivalently, the transported descendant sector has nonzero overlap with the canonically selected
  compressed dominant odd line;
\item the orientation functional is realized intrinsically on the sufficiently local generator layer:
  \[
  \mathcal A^-_{\Lambda}\subset\ker\chi_{\Lambda},
  \qquad
  \chi_{\Lambda}(J)=\varepsilon_{\Lambda}
  \quad
  \forall J\in[\Xi^{(1)}_{\Lambda}];
  \]
\item along \(I\), the sign \(\varepsilon_{\Lambda}\) is locally constant.
\end{enumerate}
In particular, the weak transported interface is realized intrinsically, without appeal to completion
coordinates, under the stated hypotheses.
\end{theorem}

\begin{proof}
The functional \(\rho^{\mathrm{dom}}_{\Lambda}\) is defined intrinsically from the sufficiently local
generator layer by \cref{def:intrinsic-cov-response} together with the dominant-line prescription of
\cref{def:dominant-response-functional}.
Item (1) is exactly the nonzero-overlap conclusion of item (3) in
\cref{thm:active-response-transport-minimal}, item (2) is exactly item (4) there, and item (3) is
item (5).
\end{proof}

\begin{theorem}[Quadratic-probe realization of the weak activity supplement]
\label{thm:char-weak-activity}
Assume the hypotheses of \cref{thm:active-response-transport-minimal} on a connected accepted-branch
segment \(I\).
For each \(\Lambda\in I\), define the even quadratic probe sector by
\[
\mathcal P^+_{\Lambda}
:=
\operatorname{span}
\bigl\{
\mathcal Q^{\Lambda}_{f,g}
:
\text{\(f,g\) admissible band-limited receiver probes}
\bigr\}
\subset E^+_{\Lambda}.
\]
Then, for every \(F\in E^-_{\Lambda}\),
\[
\{F,A\}_{\mathrm P,\Lambda}=0
\qquad
\forall A\in\mathcal P^+_{\Lambda}
\]
implies
\[
\chi_{\Lambda}(F)=0.
\]
Equivalently, if \(\chi_{\Lambda}(F)\neq 0\), then there exists a sufficiently local even quadratic
probe \(A\in\mathcal P^+_{\Lambda}\) such that
\[
\{F,A\}_{\mathrm P,\Lambda}\neq 0.
\]
In particular, the weak interface data of \cref{thm:char-weak-interface} carry a canonical
activity witness sector built from the intrinsic quadratic receiver probes.
\end{theorem}

\begin{proof}
Assume
\[
\{F,A\}_{\mathrm P,\Lambda}=0
\qquad
\forall A\in\mathcal P^+_{\Lambda}.
\]
In particular,
\[
\{F,\mathcal Q^{\Lambda}_{f,g}\}_{\mathrm P,\Lambda}=0
\]
for every admissible pair \(f,g\).
By \cref{def:intrinsic-cov-response}, this means
\[
\big\langle f,\mathcal R_{\Lambda}(F)g\big\rangle
=
\Big\langle
\{F,\mathcal Q^{\Lambda}_{f,g}\}_{\mathrm P,\Lambda}
\Big\rangle_{\rho_{\ast,\Lambda}}
=
0
\]
for all admissible \(f,g\), hence
\(
\mathcal R_{\Lambda}(F)=0
\).
Therefore
\[
\Pi_{\mathrm{phys},\Lambda}\mathcal R_{\Lambda}(F)=0,
\]
and the definition of the raw dominant-response coefficient in
\cref{def:dominant-response-functional} gives
\[
\widetilde\rho^{\mathrm{dom}}_{\Lambda}(F)=0.
\]
By item (1) of \cref{thm:active-response-transport-minimal},
\(
\delta C^{\mathrm{LoG}}_{\Lambda}\neq 0
\),
so dominant-line realization holds and \(\rho^{\mathrm{dom}}_{\Lambda}\) is well defined.
Hence
\[
\rho^{\mathrm{dom}}_{\Lambda}(F)=0.
\]
Now apply \cref{thm:char-weak-interface}, which identifies
\(
\chi_{\Lambda}
=
\nu_{\Lambda}^{-1}\rho^{\mathrm{dom}}_{\Lambda}
\),
to obtain
\[
\chi_{\Lambda}(F)=0.
\]
The contrapositive gives the equivalent detection statement.
\end{proof}

\begin{remark}[Correct scalar form of representative comparisons]
\label{rem:char-weak-interface-rho-first}
The weak export theorem should be read \(\rho^{\mathrm{dom}}\)-first on the ambiguity quotient.
Because
\[
\chi_{\Lambda}
:=
\nu_{\Lambda}^{-1}\rho^{\mathrm{dom}}_{\Lambda},
\]
any comparison statement that depends only on dominant pairing is naturally a statement about the
descended scalar
\[
\overline{\rho}^{\mathrm{dom}}_{\Lambda}
:
E^-_{\Lambda}/\mathcal A^-_{\Lambda}\to\mathbb R
\]
or, equivalently, about its value on the transported descendant affine class.
Such a statement may be rewritten in \(\chi_{\Lambda}\) form only when the normalization is made
explicit; for example, if \(\nu_{\Lambda}=1\), then the same scalar value is represented by a
descended \(\bar\chi_{\Lambda}\).
Without that normalization the stronger unscaled \(\chi\)-compatibility form is false in general.
\end{remark}

\begin{corollary}[UOM-stiffness specialization of the weak transported calibration]
\label{cor:char-weak-calibration}
Assume \cref{ass:uv-match} and fix an accepted tick \(k\) for which the conclusions of
\cref{thm:char-weak-interface} hold.
Suppose \(\ker\chi_{\Lambda_k}\) admits a continuous symmetric coercive bilinear form whose induced
norm Hilbertizes it, and choose any
\[
J^\sharp_{\Lambda_k}\in E^-_{\Lambda_k}
\qquad
\text{with}
\qquad
\chi_{\Lambda_k}(J^\sharp_{\Lambda_k})=1.
\]
Then \cref{prop:minimal-transport-calibration} may be applied with
\[
\beta_{\Lambda_k}:=\widetilde\alpha(\Lambda_k)
\]
to produce a continuous symmetric positive definite bilinear form
\[
\mathsf B^{\widetilde\alpha}_{\Lambda_k}
:
E^-_{\Lambda_k}\times E^-_{\Lambda_k}\to\mathbb R
\]
such that
\[
\mathsf B^{\widetilde\alpha}_{\Lambda_k}(J^\sharp_{\Lambda_k},J^\sharp_{\Lambda_k})
=
\widetilde\alpha(\Lambda_k),
\qquad
\mathsf B^{\widetilde\alpha}_{\Lambda_k}(J^\sharp_{\Lambda_k},K)=0
\quad
\forall K\in\ker\chi_{\Lambda_k}.
\]
In particular,
\[
\mathcal A^-_{\Lambda_k}\subset\ker\chi_{\Lambda_k}
\]
implies
\[
\mathsf B^{\widetilde\alpha}_{\Lambda_k}(J^\sharp_{\Lambda_k},A)=0
\qquad
\forall A\in\mathcal A^-_{\Lambda_k}.
\]
\end{corollary}

\begin{proof}
By \cref{rem:char-calibration-split}, the split/leak package fixes the positive stiffness scale
\(\widetilde\alpha(\Lambda_k)\), and \cref{ass:uv-match} records its positivity and asymptotic
size.
Apply \cref{prop:minimal-transport-calibration} with
\(\beta_{\Lambda_k}=\widetilde\alpha(\Lambda_k)\).
The final orthogonality statement follows from item (2) of \cref{thm:char-weak-interface}.
\end{proof}

\begin{remark}[Inherited ledger compatibility]
\label{rem:ledger-compat}
The inherited commutative ledger algebra is intentionally not part of the core characterization
algebra \(\mathcal O^{\mathrm{char}}_{\Lambda_k}\).
Its role is separate: the split model requires the active odd representative to be completion-sector
nontrivial while acting trivially, to leading order, on the inherited classical ledger.
This is the UOM-side compatibility statement behind the slogan
``oddness in the completion sector, classicality in the accessible ledger.''
\end{remark}

\ifdefined\UOMTransportFragmentLoaded\else
\end{document}
\fi
%
  }%
}

\IfFileExists{active_response_transport_program_note.tex}{%
  \ifdefined\UOMTransportFragmentLoaded\else
\documentclass[11pt]{article}
\usepackage[margin=1in]{geometry}
\usepackage{amsmath,amssymb,amsthm,mathtools}
\usepackage{bm}
\usepackage{microtype}
\usepackage{enumitem}
\usepackage{booktabs}
\usepackage{array}
\usepackage{xcolor}
\usepackage{xr-hyper}
\usepackage{hyperref}
\usepackage[nameinlink,noabbrev]{cleveref}

\providecommand{\Fsloc}{\mathcal F_{\rm s.loc}}
\providecommand{\Qt}{Q_t}
\externaldocument[main-]{UOM/main}
\externaldocument[uomdesc-]{UOM/uom_descendant_note}
\externaldocument{UOM/completion_sector_split_model_note}

\hypersetup{
  colorlinks=true,
  linkcolor=blue!60!black,
  citecolor=blue!60!black,
  urlcolor=blue!60!black
}

\theoremstyle{definition}
\newtheorem{definition}{Definition}[section]
\newtheorem{assumption}{Assumption}[section]
\theoremstyle{plain}
\newtheorem{theorem}[definition]{Theorem}
\newtheorem{lemma}[definition]{Lemma}
\newtheorem{proposition}[definition]{Proposition}
\newtheorem{corollary}[definition]{Corollary}
\theoremstyle{remark}
\newtheorem{remark}[definition]{Remark}
\crefname{definition}{definition}{definitions}
\Crefname{definition}{Definition}{Definitions}
\crefname{assumption}{assumption}{assumptions}
\Crefname{assumption}{Assumption}{Assumptions}
\crefname{theorem}{theorem}{theorems}
\Crefname{theorem}{Theorem}{Theorems}
\crefname{lemma}{lemma}{lemmas}
\Crefname{lemma}{Lemma}{Lemmas}
\crefname{proposition}{proposition}{propositions}
\Crefname{proposition}{Proposition}{Propositions}
\crefname{corollary}{corollary}{corollaries}
\Crefname{corollary}{Corollary}{Corollaries}
\crefname{remark}{remark}{remarks}
\Crefname{remark}{Remark}{Remarks}

\newcommand{\cH}{\mathcal H}
\newcommand{\cA}{\mathcal A}
\newcommand{\cD}{\mathcal D}
\newcommand{\cB}{\mathcal B}
\newcommand{\cC}{\mathcal C}
\newcommand{\cL}{\mathcal L}
\newcommand{\cK}{\mathcal K}
\newcommand{\cR}{\mathcal R}
\newcommand{\cG}{\mathcal G}
\newcommand{\Ztwo}{\mathbb Z_2}
\newcommand{\Sys}{\mathrm{Sys}}
\newcommand{\Time}{\mathrm{Time}}
\newcommand{\RG}{\mathrm{RG}}
\newcommand{\Tr}{\mathrm{Tr}}
\newcommand{\diag}{\mathrm{diag}}
\newcommand{\ad}{\mathrm{ad}}
\newcommand{\id}{\mathrm{id}}
\newcommand{\norm}[1]{\left\lVert #1\right\rVert}
\newcommand{\abs}[1]{\left\lvert #1\right\rvert}
\newcommand{\ip}[2]{\left\langle #1,#2\right\rangle}

\externaldocument{UOM/compressed_log_spectral_selection_note}
\externaldocument{UOM/odd_selector_note}
\externaldocument{UOM/canonical_odd_interface_note}
\externaldocument[splitlift-]{UOM/split_controlled_completion_lift_note}
\title{Finite-Cutoff Active Response Transport Theorem}
\author{}
\date{}
\begin{document}
\maketitle
\fi

\section{Finite-Cutoff Active Response Transport Theorem}
\label{sec:active-response-transport-program}

This section assembles a finite-cutoff transport theorem for the odd descendant / Hamiltonization
lane of the 2T Main assumption package Definition~\ref{main-def:uom-main-theorem-assumptions}.
Accordingly, the phrase ``accepted tick'' remains literal UOM cutoff language throughout this note:
the transport theorem is pointwise on the accepted branch at fixed \(\Lambda_k\).

\begin{remark}[Physical odd sector and branchwise ambient-space conventions]
\label{rem:transport-conventions}
Throughout this note, the superscript \({}^-\) refers to physical \(\Gamma\)-odd parity on the
characterization layer, not to cohomological oddness.
Every ambiguity direction appearing below is therefore understood after passage to the physically
odd characterization sector \(E^-_{\Lambda_k}\).

All branchwise continuity statements are made after the standard accepted-branch identifications of
the relevant finite-cutoff source spaces, covariance-tangent spaces, and physical subspaces with
fixed ambient finite-cutoff spaces.
Operator-norm and Fisher-norm continuity are always understood with respect to those identifications.
\end{remark}

\begin{definition}[Odd ambiguity subspace on the characterization layer]
\label{def:transport-ambiguity-subspace}
For each accepted tick \(k\), let
\[
\mathcal A^-_{\Lambda_k}\subset E^-_{\Lambda_k}
\]
denote the closed linear span of the odd characterization-layer projections of the descendant
ambiguity directions
\[
\big(Q_tX+Q_uY\big)^-,\qquad \big(\partial_\mu K^\mu\big)^-,\qquad \Delta_{\rm ct}^{-},
\]
where \(K^\mu\) is boundary-admissible on the characterization layer and
\(\Delta_{\rm ct}^{-}\) is admissible in the sense of \cref{def:ambiguity-ct-char}.
Equivalently, \(\mathcal A^-_{\Lambda_k}\) is the odd ambiguity space already placed inside
\(\mathcal T^{\mathrm{op}}_{\Lambda_k}\) by \cref{cor:ambiguity-char}.
\end{definition}

\begin{definition}[Dominant compressed response line and scalar coefficient]
\label{def:dominant-response-functional}
Fix an accepted tick \(k\). Let
\[
u^{\mathrm{LoG}}_{\Lambda_k}
\]
and
\[
P_{-,\Lambda_k}^{\mathrm{LoG}}
\]
be the canonically selected compressed odd line and projector furnished by
\cref{cor:compressed-export-selection}. Define the physical covariance tangent induced by that line by
\[
\delta C^{\mathrm{LoG}}_{\Lambda_k}
:=
\Pi_{\mathrm{phys},\Lambda_k}\,\mathsf T_{\Lambda_k}[u^{\mathrm{LoG}}_{\Lambda_k}],
\]
and, whenever this tangent is nonzero, write
\[
L^{\mathrm{dom}}_{\Lambda_k}
:=
\mathbb R\,\widehat{\delta C}^{\mathrm{LoG}}_{\Lambda_k}
\]
for its Fisher-normalized line.
The raw dominant-response coefficient is
\begin{equation}
\widetilde{\rho}^{\mathrm{dom}}_{\Lambda_k}(F)
:=
\big\langle\!\big\langle
\Pi_{\mathrm{phys},\Lambda_k}\mathcal R_{\Lambda_k}(F),\,
\delta C^{\mathrm{LoG}}_{\Lambda_k}
\big\rangle\!\big\rangle_{C_{\ast,\Lambda_k}},
\qquad
F\in E^-_{\Lambda_k}.
\label{eq:rho-dom-raw-def}
\end{equation}
Whenever \(\delta C^{\mathrm{LoG}}_{\Lambda_k}\neq 0\), equivalently whenever dominant-line
realization holds at \(\Lambda_k\), define the normalized dominant-response coefficient by
\begin{equation}
\rho^{\mathrm{dom}}_{\Lambda_k}(F)
:=
\kappa_{\Lambda_k}^{-1/2}\,
\widetilde{\rho}^{\mathrm{dom}}_{\Lambda_k}(F),
\qquad
F\in E^-_{\Lambda_k}.
\label{eq:rho-dom-def}
\end{equation}
\end{definition}

\begin{proposition}[Well-posedness and continuity of the dominant-response coefficients]
\label{prop:rho-dom-continuity}
Fix an accepted tick \(k\).
Then the raw dominant-response coefficient
\[
\widetilde{\rho}^{\mathrm{dom}}_{\Lambda_k}:E^-_{\Lambda_k}\to\mathbb R
\]
is continuous and linear.
If dominant-line realization holds at \(\Lambda_k\), then the normalized dominant-response
coefficient
\[
\rho^{\mathrm{dom}}_{\Lambda_k}:E^-_{\Lambda_k}\to\mathbb R
\]
is also continuous and linear.
\end{proposition}

\begin{proof}
By \cref{prop:intrinsic-cov-response}, the intrinsic response map
\(
\mathcal R_{\Lambda_k}:E^-_{\Lambda_k}\to\mathcal T^{\mathrm{cov}}_{\Lambda_k}
\)
is continuous and linear.
Pairing its physical projection against the fixed covariance tangent
\(
\delta C^{\mathrm{LoG}}_{\Lambda_k}
\)
in the Fisher metric therefore gives a continuous linear functional, which is exactly
\(\widetilde{\rho}^{\mathrm{dom}}_{\Lambda_k}\).
If dominant-line realization holds, then \(\kappa_{\Lambda_k}>0\) by
\cref{prop:dominant-line-equivalent}, so \(\rho^{\mathrm{dom}}_{\Lambda_k}\) is obtained from
\(\widetilde{\rho}^{\mathrm{dom}}_{\Lambda_k}\) by multiplication by the fixed scalar
\(\kappa_{\Lambda_k}^{-1/2}\).
\end{proof}

\subsection*{Dominant-line realization}
\label{subsec:dominant-line-realization}

The compressed spectral-selection note already supplies the one-dimensional source line
\(\mathcal H^{\mathrm{LoG}}_{\Lambda_k}\).
The present transport note needs only the nonvanishing and branchwise continuity of its
\emph{physical} covariance push-forward.

\begin{definition}[Compressed-line transversality]
\label{def:compressed-line-transversality}
For an accepted tick \(k\), say that the welded response is \emph{transversely physical on the
compressed line} if
\begin{equation}
\Pi_{\mathrm{phys},\Lambda_k}\,\mathsf T_{\Lambda_k}\!\restriction_{\mathcal H^{\mathrm{LoG}}_{\Lambda_k}}
\neq 0.
\label{eq:compressed-transversality}
\end{equation}
Equivalently, the unique generator \(u^{\mathrm{LoG}}_{\Lambda_k}\) of the compressed odd line has
nonzero physical covariance response.
\end{definition}

\begin{definition}[Dominant-response norm]
\label{def:dominant-response-norm}
Whenever \(\delta C^{\mathrm{LoG}}_{\Lambda_k}\) is defined, set
\begin{equation}
\kappa_{\Lambda_k}
:=
\big\langle\!\big\langle
\delta C^{\mathrm{LoG}}_{\Lambda_k},\,
\delta C^{\mathrm{LoG}}_{\Lambda_k}
\big\rangle\!\big\rangle_{C_{\ast,\Lambda_k}}.
\label{eq:kappa-dom-def}
\end{equation}
This is the Fisher squared norm of the physical covariance tangent induced by the compressed welded
LoG source line.
\end{definition}

\begin{proposition}[Equivalent formulations of dominant-line realization]
\label{prop:dominant-line-equivalent}
For a fixed accepted tick \(k\), the following are equivalent.
\begin{enumerate}[leftmargin=2em]
\item the compressed welded LoG line is transversely physical in the sense of
  \cref{def:compressed-line-transversality};
\item the distinguished physical covariance tangent is nonzero,
  \[
  \delta C^{\mathrm{LoG}}_{\Lambda_k}\neq 0;
  \]
\item the dominant-response norm is strictly positive,
  \[
  \kappa_{\Lambda_k}>0;
  \]
\item there exists a physical covariance tangent
  \(
  B_{\Lambda_k}\in \mathcal T^{\mathrm{phys}}_{\Lambda_k}
  \)
  such that
  \[
  \big\langle\!\big\langle
  \delta C^{\mathrm{LoG}}_{\Lambda_k},\,B_{\Lambda_k}
  \big\rangle\!\big\rangle_{C_{\ast,\Lambda_k}}
  \neq 0;
  \]
\item in any Gaussian completion realization of
  \cref{cor:compressed-export-selection-completion}, the completion-sector covariance tangent
  representing \(\delta C^{\mathrm{LoG}}_{\Lambda_k}\) is nonzero.
\end{enumerate}
\end{proposition}

\begin{proof}
The equivalence of (1) and (2) is immediate from
\eqref{eq:compressed-transversality} and the definition
\[
\delta C^{\mathrm{LoG}}_{\Lambda_k}
=
\Pi_{\mathrm{phys},\Lambda_k}\mathsf T_{\Lambda_k}[u^{\mathrm{LoG}}_{\Lambda_k}],
\]
because \(\mathcal H^{\mathrm{LoG}}_{\Lambda_k}\) is the one-dimensional line generated by
\(u^{\mathrm{LoG}}_{\Lambda_k}\).

The equivalence of (2) and (3) follows from positive definiteness of the Fisher metric on the
covariance tangent space: a vector has positive Fisher norm if and only if it is nonzero.

Statement (3) implies (4) by taking
\(
B_{\Lambda_k}:=\delta C^{\mathrm{LoG}}_{\Lambda_k}
\),
which belongs to the physical tangent space by construction.
Conversely, if (4) holds then \(\delta C^{\mathrm{LoG}}_{\Lambda_k}\) cannot vanish, so (2) holds.

Finally, \cref{prop:intrinsic-cov-response,cor:compressed-export-selection-completion} identify the
abstract covariance tangent with any Gaussian completion representative of the same tangent object.
Therefore the abstract tangent is nonzero exactly when its completion realization is nonzero, giving
the equivalence of (2) and (5).
\end{proof}

\begin{theorem}[Branchwise realization and persistence of the dominant line]
\label{thm:dominant-line-persistence}
Let \(I\) be a connected accepted-branch segment in cutoff space.
Assume that:
\begin{enumerate}[leftmargin=2em]
\item the compressed projector \(P_{-,\Lambda}^{\mathrm{LoG}}\) varies continuously on \(I\);
\item the physical welded response operator
  \[
  \Lambda\longmapsto \Pi_{\mathrm{phys},\Lambda}\mathsf T_{\Lambda}
  \]
  is continuous on \(I\) in the operator topology induced by the Fisher norm on covariance
  tangents;
\item there exists a constant \(c_I>0\) such that
  \[
  \kappa_{\Lambda}\ge c_I
  \qquad
  \forall \Lambda\in I.
  \]
\end{enumerate}
Then:
\begin{enumerate}[leftmargin=2em]
\item \(\delta C^{\mathrm{LoG}}_{\Lambda}\neq 0\) for every \(\Lambda\in I\);
\item the normalized dominant vector
  \[
  \widehat{\delta C}^{\mathrm{LoG}}_{\Lambda}
  :=
  \kappa_{\Lambda}^{-1/2}\,\delta C^{\mathrm{LoG}}_{\Lambda}
  \]
  is continuous on \(I\);
\item the dominant physical line
  \[
  L^{\mathrm{dom}}_{\Lambda}
  =
  \mathbb R\,\widehat{\delta C}^{\mathrm{LoG}}_{\Lambda}
  \]
  is canonically transported along \(I\);
\item the Fisher-orthogonal rank-one projector
  \[
  \Pi^{\mathrm{dom}}_{\Lambda}
  :=
  \widehat{\delta C}^{\mathrm{LoG}}_{\Lambda}
  \otimes
  \widehat{\delta C}^{\mathrm{LoG}}_{\Lambda}
  \]
  depends continuously on \(\Lambda\in I\).
\end{enumerate}
\end{theorem}

\begin{proof}
By \cref{cor:compressed-export-selection}, the compressed projector
\(
P_{-,\Lambda}^{\mathrm{LoG}}
\)
and hence the normalized source vector \(u^{\mathrm{LoG}}_{\Lambda}\) vary continuously on the branch.
By the assumed continuity of \(\Pi_{\mathrm{phys},\Lambda}\mathsf T_{\Lambda}\), the map
\[
\Lambda\longmapsto
\delta C^{\mathrm{LoG}}_{\Lambda}
=
\Pi_{\mathrm{phys},\Lambda}\mathsf T_{\Lambda}[u^{\mathrm{LoG}}_{\Lambda}]
\]
is continuous in the Fisher norm.

The lower bound \(\kappa_{\Lambda}\ge c_I>0\) implies
\(
\delta C^{\mathrm{LoG}}_{\Lambda}\neq 0
\)
for every \(\Lambda\in I\), proving item (1).
Since the normalization map \(v\mapsto v/\|v\|_{C_\ast}\) is continuous on the complement of the
zero section, continuity of \(\delta C^{\mathrm{LoG}}_{\Lambda}\) and the positive lower bound on its
norm imply continuity of \(\widehat{\delta C}^{\mathrm{LoG}}_{\Lambda}\), proving item (2).
Items (3) and (4) are immediate consequences of item (2).
\end{proof}

\begin{corollary}[Local persistence from one-tick transversality]
\label{cor:dominant-line-local-persistence}
Assume the continuity hypotheses of \cref{thm:dominant-line-persistence}.
If the compressed line is transversely physical at one accepted tick \(\Lambda_{k_0}\), then there
exists a neighborhood \(I_{k_0}\) of \(\Lambda_{k_0}\) on the same accepted branch such that
\[
\delta C^{\mathrm{LoG}}_{\Lambda}\neq 0
\qquad
\forall \Lambda\in I_{k_0}.
\]
Hence the dominant physical line is locally well defined and continuously transported near
\(\Lambda_{k_0}\).
\end{corollary}

\begin{proof}
By \cref{prop:dominant-line-equivalent}, transversality at \(\Lambda_{k_0}\) is equivalent to
\(
\kappa_{\Lambda_{k_0}}>0
\).
Continuity of \(\kappa_{\Lambda}\) then gives a neighborhood \(I_{k_0}\) and a constant
\(c_{k_0}>0\) such that
\(
\kappa_{\Lambda}\ge c_{k_0}
\)
for all \(\Lambda\in I_{k_0}\).
Apply \cref{thm:dominant-line-persistence}.
\end{proof}

\begin{remark}[Practical witness criteria]
\label{rem:dominant-line-witnesses}
For applications, any one of the following suffices to discharge dominant-line realization at a fixed
accepted tick:
\begin{enumerate}[leftmargin=2em]
\item a direct lower bound \(\kappa_{\Lambda_k}>0\);
\item a physical witness tangent \(B_{\Lambda_k}\in\mathcal T^{\mathrm{phys}}_{\Lambda_k}\) with
  nonzero Fisher pairing against \(\delta C^{\mathrm{LoG}}_{\Lambda_k}\);
\item a Gaussian completion calculation showing that the completion covariance of the realized active
  odd mode changes nontrivially under the compressed welded LoG source.
\end{enumerate}
Each witness proves the same theorem-level fact by \cref{prop:dominant-line-equivalent}.
\end{remark}

\subsection*{Descent to the ambiguity quotient}
\label{subsec:ambiguity-quotient-descent}

Once the dominant line has been realized, the next task is purely structural:
show that the dominant-response coefficient is insensitive to ambiguity moves.
This is the exact point at which the transported descendant class becomes an affine class in the
sense required by Main rather than a collection of representatives.

\begin{definition}[Dominant-response quotient and descended descendant point]
\label{def:dominant-response-quotient}
For each accepted tick \(k\), define the dominant-response quotient by
\[
\mathfrak D^-_{\Lambda_k}
:=
E^-_{\Lambda_k}/\mathcal A^-_{\Lambda_k},
\]
with quotient map
\[
q_{\Lambda_k}:E^-_{\Lambda_k}\to \mathfrak D^-_{\Lambda_k}.
\]
If
\[
[\Xi^{(1)}_{\Lambda_k}]
=
J_{0,\Lambda_k}+\mathcal A^-_{\Lambda_k},
\]
then its image in the quotient is a single point,
\[
[\![\Xi^{(1)}_{\Lambda_k}]\!]
:=
q_{\Lambda_k}\!\big(J_{0,\Lambda_k}\big)
\in \mathfrak D^-_{\Lambda_k},
\]
independent of the chosen representative \(J_{0,\Lambda_k}\).
\end{definition}

\begin{proposition}[Universal descent criterion for the dominant-response coefficient]
\label{prop:dominant-quotient-universal}
Assume dominant-line realization at an accepted tick \(k\), so that
\(\rho^{\mathrm{dom}}_{\Lambda_k}\) is defined by \eqref{eq:rho-dom-def}.
Then the following are equivalent.
\begin{enumerate}[leftmargin=2em]
\item there exists a unique continuous linear functional
  \[
  \overline{\rho}^{\mathrm{dom}}_{\Lambda_k}:\mathfrak D^-_{\Lambda_k}\to\mathbb R
  \]
  such that
  \[
  \rho^{\mathrm{dom}}_{\Lambda_k}
  =
  \overline{\rho}^{\mathrm{dom}}_{\Lambda_k}\circ q_{\Lambda_k};
  \]
\item the ambiguity subspace is contained in the kernel:
  \[
  \mathcal A^-_{\Lambda_k}\subset \ker\rho^{\mathrm{dom}}_{\Lambda_k};
  \]
\item the dominant-response coefficient is constant on each ambiguity coset:
  \[
  \rho^{\mathrm{dom}}_{\Lambda_k}(F+A)
  =
  \rho^{\mathrm{dom}}_{\Lambda_k}(F)
  \qquad
  \forall F\in E^-_{\Lambda_k},\ \forall A\in\mathcal A^-_{\Lambda_k}.
  \]
\end{enumerate}
\end{proposition}

\begin{proof}
The equivalence of (2) and (3) is immediate from linearity:
\[
\rho^{\mathrm{dom}}_{\Lambda_k}(F+A)-\rho^{\mathrm{dom}}_{\Lambda_k}(F)
=
\rho^{\mathrm{dom}}_{\Lambda_k}(A).
\]

If (2) holds, define
\[
\overline{\rho}^{\mathrm{dom}}_{\Lambda_k}\big(q_{\Lambda_k}(F)\big)
:=
\rho^{\mathrm{dom}}_{\Lambda_k}(F).
\]
This is well defined precisely because changing representatives by an element of
\(\mathcal A^-_{\Lambda_k}\) does not change the value of \(\rho^{\mathrm{dom}}_{\Lambda_k}\).
Linearity and continuity are inherited from \(\rho^{\mathrm{dom}}_{\Lambda_k}\), and uniqueness
follows because \(q_{\Lambda_k}\) is surjective.

Conversely, if (1) holds and \(A\in\mathcal A^-_{\Lambda_k}\), then \(q_{\Lambda_k}(A)=0\), hence
\[
\rho^{\mathrm{dom}}_{\Lambda_k}(A)
=
\overline{\rho}^{\mathrm{dom}}_{\Lambda_k}\!\big(q_{\Lambda_k}(A)\big)
=
\overline{\rho}^{\mathrm{dom}}_{\Lambda_k}(0)
=
0.
\]
Thus (1) implies (2).
\end{proof}

\begin{corollary}[Constancy on the transported descendant affine class]
\label{cor:dominant-quotient-constancy}
Assume \cref{prop:dominant-quotient-universal}. Then
\[
\rho^{\mathrm{dom}}_{\Lambda_k}(J)
=
\overline{\rho}^{\mathrm{dom}}_{\Lambda_k}\big([\![\Xi^{(1)}_{\Lambda_k}]\!]\big)
\qquad
\forall J\in[\Xi^{(1)}_{\Lambda_k}].
\]
In particular, once the descended quotient functional exists, the descendant affine class carries a
single unambiguous dominant-response value.
\end{corollary}

\begin{proof}
Every \(J\in[\Xi^{(1)}_{\Lambda_k}]\) has the form
\(
J=J_{0,\Lambda_k}+A
\)
with \(A\in\mathcal A^-_{\Lambda_k}\), so
\[
q_{\Lambda_k}(J)=q_{\Lambda_k}(J_{0,\Lambda_k})=[\![\Xi^{(1)}_{\Lambda_k}]\!].
\]
Apply the factorization identity of \cref{prop:dominant-quotient-universal}.
\end{proof}

\begin{remark}[False stronger form and normalized special case]
\label{rem:dominant-quotient-false-stronger-form}
The quotient statement above is intrinsically a statement about the descended dominant scalar
\[
\overline{\rho}^{\mathrm{dom}}_{\Lambda_k}\big([\![\Xi^{(1)}_{\Lambda_k}]\!]\big),
\]
not about a bare \(\chi\)-compatibility condition.
At the weak transported-interface layer one has
\[
\chi_{\Lambda_k}
=
\nu_{\Lambda_k}^{-1}\rho^{\mathrm{dom}}_{\Lambda_k},
\]
so a representative-search condition of the form
\[
\rho^{\mathrm{dom}}_{\Lambda_k}(J)=s
\]
is equivalent to
\[
\chi_{\Lambda_k}(J)=\nu_{\Lambda_k}^{-1}s,
\]
not in general to \(\chi_{\Lambda_k}(J)=s\).
Only after an explicit normalization, for example \(\nu_{\Lambda_k}=1\), may the same scalar
comparison be rewritten in descended \(\bar\chi_{\Lambda_k}\) form without changing the value.
The finite-cutoff Wolfram audit contains exact countermodels to the stronger unscaled
\(\chi\)-level claim.
\end{remark}

\begin{proposition}[Generator-by-generator reduction]
\label{prop:dominant-quotient-generators}
Assume dominant-line realization at an accepted tick \(k\).
Suppose the dominant-response coefficient vanishes on the three generating ambiguity families:
\begin{enumerate}[leftmargin=2em]
\item for every odd-projected BRST/gauge-exact direction \(\big(Q_tX+Q_uY\big)^-\),
  \[
  \rho^{\mathrm{dom}}_{\Lambda_k}\!\big(\big(Q_tX+Q_uY\big)^-\big)=0;
  \]
\item for every odd-projected boundary-admissible divergence \(\big(\partial_\mu K^\mu\big)^-\),
  \[
  \rho^{\mathrm{dom}}_{\Lambda_k}\!\big(\big(\partial_\mu K^\mu\big)^-\big)=0;
  \]
\item for every admissible odd counterterm \(\Delta^-_{\rm ct}\),
  \[
  \rho^{\mathrm{dom}}_{\Lambda_k}(\Delta^-_{\rm ct})=0.
  \]
\end{enumerate}
Then
\[
\mathcal A^-_{\Lambda_k}\subset\ker\rho^{\mathrm{dom}}_{\Lambda_k},
\]
hence \(\rho^{\mathrm{dom}}_{\Lambda_k}\) descends uniquely to the quotient
\(\mathfrak D^-_{\Lambda_k}\).
\end{proposition}

\begin{proof}
By \cref{def:transport-ambiguity-subspace}, \(\mathcal A^-_{\Lambda_k}\) is the closed linear span of
the three families listed in the statement. Since \(\rho^{\mathrm{dom}}_{\Lambda_k}\) is linear, it
vanishes on the algebraic span of those generators. Since it is also continuous, it vanishes on the
closure of that span. Therefore every element of \(\mathcal A^-_{\Lambda_k}\) lies in the kernel.
Now apply \cref{prop:dominant-quotient-universal}.
\end{proof}

\begin{theorem}[Descent to the ambiguity quotient]
\label{thm:dominant-quotient-descent}
Fix an accepted tick \(k\) and assume dominant-line realization.
If the three vanishing statements of \cref{prop:dominant-quotient-generators} hold, then there
exists a unique descended dominant
functional
\[
\overline{\rho}^{\mathrm{dom}}_{\Lambda_k}:\mathfrak D^-_{\Lambda_k}\to\mathbb R
\]
such that
\[
\rho^{\mathrm{dom}}_{\Lambda_k}
=
\overline{\rho}^{\mathrm{dom}}_{\Lambda_k}\circ q_{\Lambda_k},
\]
and the transported descendant affine class has a single well-defined quotient value
\[
\overline{\rho}^{\mathrm{dom}}_{\Lambda_k}\big([\![\Xi^{(1)}_{\Lambda_k}]\!]\big).
\]
\end{theorem}

\begin{proof}
Apply \cref{prop:dominant-quotient-generators,cor:dominant-quotient-constancy}.
\end{proof}

\subsection*{Nonzero overlap on the descendant class}
\label{subsec:descendant-overlap}

After dominant-line realization and quotient descent, the descendant class defines a single point in
the ambiguity quotient, and the dominant-response functional descends to that quotient. This
subsection studies whether that descended descendant point is \emph{active}, i.e.\ not annihilated
by the descended dominant functional.

\begin{definition}[Dominant covariance projector]
\label{def:dominant-cov-projector}
Assume dominant-line realization at an accepted tick \(k\).
Define the Fisher-orthogonal rank-one projector onto the dominant physical line by
\begin{equation}
\Pi^{\mathrm{dom}}_{\Lambda_k}(X)
:=
\big\langle\!\big\langle
X,\,
\widehat{\delta C}^{\mathrm{LoG}}_{\Lambda_k}
\big\rangle\!\big\rangle_{C_{\ast,\Lambda_k}}\,
\widehat{\delta C}^{\mathrm{LoG}}_{\Lambda_k},
\qquad
X\in\mathcal T^{\mathrm{cov}}_{\Lambda_k}.
\label{eq:Pi-dom-def}
\end{equation}
\end{definition}

\begin{proposition}[Representative-level witness criteria]
\label{prop:descendant-overlap-witness}
Assume dominant-line realization at an accepted tick \(k\).
For any odd generator \(F\in E^-_{\Lambda_k}\), the following are equivalent.
\begin{enumerate}[leftmargin=2em]
\item the dominant-response coefficient is nonzero:
  \[
  \rho^{\mathrm{dom}}_{\Lambda_k}(F)\neq 0;
  \]
\item the dominant projection of the physical covariance response is nonzero:
  \[
  \Pi^{\mathrm{dom}}_{\Lambda_k}\Pi_{\mathrm{phys},\Lambda_k}\mathcal R_{\Lambda_k}(F)\neq 0;
  \]
\item the physical covariance response has nonzero Fisher pairing with the dominant tangent:
  \[
  \big\langle\!\big\langle
  \Pi_{\mathrm{phys},\Lambda_k}\mathcal R_{\Lambda_k}(F),\,
  \delta C^{\mathrm{LoG}}_{\Lambda_k}
  \big\rangle\!\big\rangle_{C_{\ast,\Lambda_k}}
  \neq 0;
  \]
\item in any Gaussian completion realization of
  \cref{cor:compressed-export-selection-completion}, the completion representative of
  \(\Pi_{\mathrm{phys},\Lambda_k}\mathcal R_{\Lambda_k}(F)\) has nonzero component on the realized
  active odd line.
\end{enumerate}
\end{proposition}

\begin{proof}
By \cref{def:dominant-cov-projector},
\[
\Pi^{\mathrm{dom}}_{\Lambda_k}\Pi_{\mathrm{phys},\Lambda_k}\mathcal R_{\Lambda_k}(F)
=
\rho^{\mathrm{dom}}_{\Lambda_k}(F)\,
\widehat{\delta C}^{\mathrm{LoG}}_{\Lambda_k}.
\]
Since \(\widehat{\delta C}^{\mathrm{LoG}}_{\Lambda_k}\neq 0\), this proves the equivalence of (1)
and (2).

Statement (3) is equivalent to (1) because
\[
\big\langle\!\big\langle
\Pi_{\mathrm{phys},\Lambda_k}\mathcal R_{\Lambda_k}(F),\,
\delta C^{\mathrm{LoG}}_{\Lambda_k}
\big\rangle\!\big\rangle_{C_{\ast,\Lambda_k}}
=
\|\delta C^{\mathrm{LoG}}_{\Lambda_k}\|_{C_{\ast,\Lambda_k}}\,
\rho^{\mathrm{dom}}_{\Lambda_k}(F),
\]
and the prefactor is strictly positive by dominant-line realization.

Finally, by \cref{prop:intrinsic-cov-response,cor:compressed-export-selection-completion}, the
abstract covariance tangent and its Gaussian completion representative describe the same physical
response object in different coordinates. Therefore the abstract dominant projection is nonzero if
and only if the completion representative has nonzero component on the realized active line,
proving the equivalence of (2) and (4).
\end{proof}

\begin{definition}[Descendant overlap margin]
\label{def:descendant-overlap-margin}
Assume dominant-line realization and quotient descent at an accepted tick \(k\).
Define the descendant overlap margin by
\begin{equation}
\nu_{\Lambda_k}
:=
\left|
\overline{\rho}^{\mathrm{dom}}_{\Lambda_k}\big([\![\Xi^{(1)}_{\Lambda_k}]\!]\big)
\right|.
\label{eq:descendant-overlap-margin}
\end{equation}
We say that the transported descendant class is \emph{dominant-active} if
\(
\nu_{\Lambda_k}>0
\).
\end{definition}

\begin{proposition}[Equivalent formulations of nonzero overlap on the descendant class]
\label{prop:descendant-overlap-equivalent}
Assume dominant-line realization and quotient descent at an accepted tick \(k\).
Then the following are equivalent.
\begin{enumerate}[leftmargin=2em]
\item the transported descendant sector has nonzero overlap with the canonically selected compressed
  dominant odd line;
\item there exists one representative \(J^{\mathrm{desc}}_{\Lambda_k}\in[\Xi^{(1)}_{\Lambda_k}]\)
  such that
  \[
  \Pi^{\mathrm{dom}}_{\Lambda_k}\Pi_{\mathrm{phys},\Lambda_k}
  \mathcal R_{\Lambda_k}(J^{\mathrm{desc}}_{\Lambda_k})
  \neq 0;
  \]
\item every representative \(J\in[\Xi^{(1)}_{\Lambda_k}]\) has the same nonzero dominant-response
  value:
  \[
  \rho^{\mathrm{dom}}_{\Lambda_k}(J)
  =
  \overline{\rho}^{\mathrm{dom}}_{\Lambda_k}\big([\![\Xi^{(1)}_{\Lambda_k}]\!]\big)
  \neq 0;
  \]
\item the descended descendant point is not annihilated by the descended dominant functional:
  \[
  [\![\Xi^{(1)}_{\Lambda_k}]\!]\notin \ker\overline{\rho}^{\mathrm{dom}}_{\Lambda_k};
  \]
\item the descendant overlap margin is strictly positive:
  \[
  \nu_{\Lambda_k}>0.
  \]
\end{enumerate}
\end{proposition}

\begin{proof}
The equivalence of (1) and (2) follows immediately from
\cref{prop:descendant-overlap-witness}.

If (2) holds, then
\(
\rho^{\mathrm{dom}}_{\Lambda_k}(J^{\mathrm{desc}}_{\Lambda_k})\neq 0
\)
by \cref{prop:descendant-overlap-witness}. By
\cref{cor:dominant-quotient-constancy}, every other representative of
\([\Xi^{(1)}_{\Lambda_k}]\) has the same value, proving (3).
Conversely, (3) obviously implies (1).

The equivalence of (3) and (4) is immediate from the identity
\[
\rho^{\mathrm{dom}}_{\Lambda_k}(J)
=
\overline{\rho}^{\mathrm{dom}}_{\Lambda_k}\big([\![\Xi^{(1)}_{\Lambda_k}]\!]\big)
\qquad
\forall J\in[\Xi^{(1)}_{\Lambda_k}]
\]
of \cref{cor:dominant-quotient-constancy}.

Finally, (4) and (5) are equivalent by the definition of \(\nu_{\Lambda_k}\) in
\cref{def:descendant-overlap-margin}.
\end{proof}

\begin{theorem}[Nonzero overlap on the descendant class]
\label{thm:descendant-overlap}
Fix an accepted tick \(k\), and assume dominant-line realization together with quotient descent.
If any one of the equivalent conditions in \cref{prop:descendant-overlap-equivalent} holds, then the
following conclusions obtain.
In particular, it is enough to construct a single representative
\[
J^{\mathrm{desc}}_{\Lambda_k}\in[\Xi^{(1)}_{\Lambda_k}]
\]
whose dominant projected covariance response is nonzero.

When this holds, the quotient descendant point is dominant-active, the scalar value
\[
\overline{\rho}^{\mathrm{dom}}_{\Lambda_k}\big([\![\Xi^{(1)}_{\Lambda_k}]\!]\big)
\]
is nonzero, and the sign
\begin{equation}
\varepsilon_{\Lambda_k}
:=
\operatorname{sgn}\!\Big(
\overline{\rho}^{\mathrm{dom}}_{\Lambda_k}\big([\![\Xi^{(1)}_{\Lambda_k}]\!]\big)
\Big)
\in\{\pm1\}
\label{eq:epsilon-from-overlap}
\end{equation}
is well defined.
Thus the transported descendant sector has nonzero overlap with the canonically selected compressed
dominant odd line.
\end{theorem}

\begin{proof}
The equivalence package is exactly \cref{prop:descendant-overlap-equivalent}.
Once one representative has nonzero dominant projected response, the same proposition gives the
nonzero scalar value on the descended descendant point. The sign in \eqref{eq:epsilon-from-overlap}
is therefore well defined.
\end{proof}

\subsection*{Conditional mod-null reentry on the selected scalar line}
\label{subsec:conditional-mod-null-reentry}

\begin{definition}[Parity-stable mod-null subspace transverse to the selected scalar line]
\label{def:mod-null-transverse-datum}
Fix an accepted tick \(k\), assume dominant-line realization and quotient descent, and let
\[
\overline L^{\mathrm{Main}}_{\Lambda_k}
=
\mathbb R\,\overline\ell^{\mathrm{Main}}_{\Lambda_k}
\subset
\mathfrak D^-_{\Lambda_k}
\]
be a one-dimensional selected scalar line in the ambiguity quotient.
A \emph{parity-stable mod-null datum transverse to}
\(
\overline L^{\mathrm{Main}}_{\Lambda_k}
\)
consists of a linear subspace
\[
N^{\mathrm{mn}}_{\Lambda_k}\subset \mathfrak D^-_{\Lambda_k}
\]
such that:
\begin{enumerate}[leftmargin=2em,label=(\alph*)]
\item \(N^{\mathrm{mn}}_{\Lambda_k}\) is stable under the descended physical parity on
  \(\mathfrak D^-_{\Lambda_k}\);
\item the selected scalar line is transverse to \(N^{\mathrm{mn}}_{\Lambda_k}\):
  \[
  \overline L^{\mathrm{Main}}_{\Lambda_k}\cap N^{\mathrm{mn}}_{\Lambda_k}=\{0\}.
  \]
\end{enumerate}
Define
\[
\mathcal V^{1,\mathrm{mn}}_{\Lambda_k}
:=
\mathfrak D^-_{\Lambda_k}/N^{\mathrm{mn}}_{\Lambda_k}
\]
with quotient map
\[
\pi^{\mathrm{mn}}_{\Lambda_k}:
\mathfrak D^-_{\Lambda_k}\to \mathcal V^{1,\mathrm{mn}}_{\Lambda_k}.
\]
Because \(\mathfrak D^-_{\Lambda_k}\) is itself a quotient of the physical \(\Gamma\)-odd sector,
the descended parity acts on \(\mathfrak D^-_{\Lambda_k}\) by \(-\mathrm{id}\).
Thus the parity-stability clause yields a descended involution
\[
\Gamma^{\mathrm{mn}}_{\Lambda_k}
\bigl(
\pi^{\mathrm{mn}}_{\Lambda_k}(x)
\bigr)
:=
\pi^{\mathrm{mn}}_{\Lambda_k}(-x),
\qquad
x\in \mathfrak D^-_{\Lambda_k}.
\]
\end{definition}

\begin{definition}[Dominant-null parity-stable mod-null datum]
\label{def:dominant-null-mod-null-datum}
Fix an accepted tick \(k\), assume dominant-line realization and quotient descent, and let
\[
\overline L^{\mathrm{Main}}_{\Lambda_k}
=
\mathbb R\,\overline\ell^{\mathrm{Main}}_{\Lambda_k}
\subset
\mathfrak D^-_{\Lambda_k}
\]
be a one-dimensional selected scalar line in the ambiguity quotient.
A parity-stable mod-null datum
\[
N^{\mathrm{mn}}_{\Lambda_k}\subset \mathfrak D^-_{\Lambda_k}
\]
is called \emph{dominant-null} if it is contained in the kernel of the descended dominant
functional:
\[
N^{\mathrm{mn}}_{\Lambda_k}\subset \ker \overline{\rho}^{\mathrm{dom}}_{\Lambda_k}.
\]
It is understood that
\(
N^{\mathrm{mn}}_{\Lambda_k}
\)
is stable under the descended physical parity on
\(
\mathfrak D^-_{\Lambda_k}
\),
and that the associated quotient
\(
\mathfrak D^-_{\Lambda_k}/N^{\mathrm{mn}}_{\Lambda_k}
\)
carries the descended involution from \cref{def:mod-null-transverse-datum}.
The canonical dominant-kernel choice is isolated separately in
\cref{def:canonical-dominant-null-datum}.
The transversality lemma below uses only this kernel-subordination property.
\end{definition}

\begin{definition}[Canonical dominant-null datum]
\label{def:canonical-dominant-null-datum}
Fix an accepted tick \(k\) and assume dominant-line realization together with quotient descent.
Define the \emph{canonical dominant-null datum} by
\[
N^{\mathrm{dom-null}}_{\Lambda_k}
:=
\ker \overline{\rho}^{\mathrm{dom}}_{\Lambda_k}
\subset
\mathfrak D^-_{\Lambda_k}.
\]
This datum depends only on the descended dominant functional on the ambiguity quotient.
No additional reduced-welded or completion-side choice is built into the definition.
\end{definition}

\begin{lemma}[Dominant-kernel criterion for selected-line transversality]
\label{lem:dominant-kernel-transversality}
Fix an accepted tick \(k\), assume dominant-line realization together with quotient descent, and let
\[
\overline L^{\mathrm{Main}}_{\Lambda_k}
=
\mathbb R\,\overline\ell^{\mathrm{Main}}_{\Lambda_k}
\subset
\mathfrak D^-_{\Lambda_k}
\]
be a one-dimensional selected scalar line with
\[
\overline{\rho}^{\mathrm{dom}}_{\Lambda_k}
\!\bigl(
\overline\ell^{\mathrm{Main}}_{\Lambda_k}
\bigr)
\neq 0.
\]
If
\[
N^{\mathrm{mn}}_{\Lambda_k}\subset \mathfrak D^-_{\Lambda_k}
\]
is a dominant-null parity-stable mod-null datum in the sense of
\cref{def:dominant-null-mod-null-datum},
then
\[
\overline L^{\mathrm{Main}}_{\Lambda_k}\cap N^{\mathrm{mn}}_{\Lambda_k}=\{0\}.
\]
So every dominant-null parity-stable mod-null datum is automatically transverse to the selected
scalar line.
\end{lemma}

\begin{proof}
Let
\(
x\in \overline L^{\mathrm{Main}}_{\Lambda_k}\cap N^{\mathrm{mn}}_{\Lambda_k}
\).
Because
\(
\overline L^{\mathrm{Main}}_{\Lambda_k}
=
\mathbb R\,\overline\ell^{\mathrm{Main}}_{\Lambda_k}
\),
there exists \(a\in\mathbb R\) with
\(
x=a\,\overline\ell^{\mathrm{Main}}_{\Lambda_k}
\).
Since
\(
N^{\mathrm{mn}}_{\Lambda_k}\subset
\ker \overline{\rho}^{\mathrm{dom}}_{\Lambda_k}
\),
we have
\[
0
=
\overline{\rho}^{\mathrm{dom}}_{\Lambda_k}(x)
=
a\,
\overline{\rho}^{\mathrm{dom}}_{\Lambda_k}
\!\bigl(
\overline\ell^{\mathrm{Main}}_{\Lambda_k}
\bigr).
\]
The coefficient
\(
\overline{\rho}^{\mathrm{dom}}_{\Lambda_k}
(\overline\ell^{\mathrm{Main}}_{\Lambda_k})
\)
is nonzero by hypothesis, so \(a=0\).
Therefore \(x=0\), proving
\(
\overline L^{\mathrm{Main}}_{\Lambda_k}\cap N^{\mathrm{mn}}_{\Lambda_k}=\{0\}
\).
\end{proof}

\begin{theorem}[The dominant kernel is the canonical dominant-null datum]
\label{thm:canonical-dominant-null-datum}
Fix an accepted tick \(k\), assume dominant-line realization together with quotient descent, and let
\[
\overline L^{\mathrm{Main}}_{\Lambda_k}
=
\mathbb R\,\overline\ell^{\mathrm{Main}}_{\Lambda_k}
\subset
\mathfrak D^-_{\Lambda_k}
\]
be a one-dimensional selected scalar line with
\[
\overline{\rho}^{\mathrm{dom}}_{\Lambda_k}
\!\bigl(
\overline\ell^{\mathrm{Main}}_{\Lambda_k}
\bigr)
\neq 0.
\]
Define
\[
N^{\mathrm{dom-null}}_{\Lambda_k}
:=
\ker \overline{\rho}^{\mathrm{dom}}_{\Lambda_k}
\subset
\mathfrak D^-_{\Lambda_k}.
\]
Then:
\begin{enumerate}[leftmargin=2em]
\item \(N^{\mathrm{dom-null}}_{\Lambda_k}\) is stable under the descended physical parity on
  \(\mathfrak D^-_{\Lambda_k}\);
\item \(N^{\mathrm{dom-null}}_{\Lambda_k}\) is a dominant-null parity-stable mod-null datum in the
  sense of \cref{def:dominant-null-mod-null-datum};
\item the selected scalar line is automatically transverse to the canonical dominant-null datum:
  \[
  \overline L^{\mathrm{Main}}_{\Lambda_k}
  \cap
  N^{\mathrm{dom-null}}_{\Lambda_k}
  =
  \{0\}.
  \]
\end{enumerate}
Hence the canonical dominant kernel is admissible as the mod-null datum for the conditional
reentry theorem.
\end{theorem}

\begin{proof}
Let \(x\in N^{\mathrm{dom-null}}_{\Lambda_k}\).
Because the descended physical parity acts on \(\mathfrak D^-_{\Lambda_k}\) by \(-\mathrm{id}\),
we have
\[
\overline{\rho}^{\mathrm{dom}}_{\Lambda_k}(-x)
=
-\overline{\rho}^{\mathrm{dom}}_{\Lambda_k}(x)
=
0.
\]
So \(-x\in N^{\mathrm{dom-null}}_{\Lambda_k}\), proving parity stability.
Item \textnormal{(2)} is then exactly \cref{def:dominant-null-mod-null-datum} with
\(
N^{\mathrm{mn}}_{\Lambda_k}=N^{\mathrm{dom-null}}_{\Lambda_k}
\).
Item \textnormal{(3)} is \cref{lem:dominant-kernel-transversality} applied to this canonical
choice.
\end{proof}

\begin{lemma}[Transversality reduces mod-null reentry to quotient construction]
\label{lem:mod-null-transversality-reduction}
Fix an accepted tick \(k\), assume dominant-line realization together with quotient descent, and let
\[
\overline L^{\mathrm{Main}}_{\Lambda_k}
=
\mathbb R\,\overline\ell^{\mathrm{Main}}_{\Lambda_k}
\subset
\mathfrak D^-_{\Lambda_k}
\]
be a one-dimensional selected scalar line with
\[
\overline{\rho}^{\mathrm{dom}}_{\Lambda_k}
\!\bigl(
\overline\ell^{\mathrm{Main}}_{\Lambda_k}
\bigr)
\neq 0.
\]
Assume moreover a parity-stable mod-null datum transverse to
\(
\overline L^{\mathrm{Main}}_{\Lambda_k}
\)
in the sense of \cref{def:mod-null-transverse-datum}.
Define
\[
\ell^{\mathrm{sel,mn}}_{\Lambda_k}
:=
\pi^{\mathrm{mn}}_{\Lambda_k}
\!\bigl(
\overline\ell^{\mathrm{Main}}_{\Lambda_k}
\bigr)
\]
and
\[
L^{\mathrm{sel,mn}}_{\Lambda_k}
:=
\mathbb R\,\ell^{\mathrm{sel,mn}}_{\Lambda_k}
\subset
\mathcal V^{1,\mathrm{mn}}_{\Lambda_k}.
\]
Then:
\begin{enumerate}[leftmargin=2em]
\item the restriction of
  \(
  \pi^{\mathrm{mn}}_{\Lambda_k}
  \)
  to
  \(
  \overline L^{\mathrm{Main}}_{\Lambda_k}
  \)
  is injective;
\item \(L^{\mathrm{sel,mn}}_{\Lambda_k}\) is a well-defined one-dimensional line in the mod-null
  quotient, generated by the nonzero class
  \(
  \ell^{\mathrm{sel,mn}}_{\Lambda_k}
  \);
\item the mod-null selected realization map
  \[
  \mathfrak R^{\mathrm{mn}}_{\Lambda_k}:
  L^{\mathrm{sel,mn}}_{\Lambda_k}
  \longrightarrow
  \mathbb R\,u^{\mathrm{LoG}}_{\Lambda_k},
  \qquad
  \mathfrak R^{\mathrm{mn}}_{\Lambda_k}
  \!\bigl(
  a\,\ell^{\mathrm{sel,mn}}_{\Lambda_k}
  \bigr)
  :=
  a\,
  \overline{\rho}^{\mathrm{dom}}_{\Lambda_k}
  \!\bigl(
  \overline\ell^{\mathrm{Main}}_{\Lambda_k}
  \bigr)\,
  u^{\mathrm{LoG}}_{\Lambda_k}
  \]
  is well defined and faithful;
\item the descended involution acts by sign:
  \[
  \Gamma^{\mathrm{mn}}_{\Lambda_k}
  \!\bigl(
  \ell^{\mathrm{sel,mn}}_{\Lambda_k}
  \bigr)
  =
  -
  \ell^{\mathrm{sel,mn}}_{\Lambda_k}.
  \]
\end{enumerate}
In particular, once the transversality condition
\(
\overline L^{\mathrm{Main}}_{\Lambda_k}\cap N^{\mathrm{mn}}_{\Lambda_k}=\{0\}
\)
is known, the remaining mod-null reentry step reduces to quotient construction plus parity descent.
\end{lemma}

\begin{proof}
If
\(
x\in \overline L^{\mathrm{Main}}_{\Lambda_k}
\)
and
\(
\pi^{\mathrm{mn}}_{\Lambda_k}(x)=0
\),
then
\(
x\in N^{\mathrm{mn}}_{\Lambda_k}
\)
by definition of the quotient kernel.
Hence
\(
x\in
\overline L^{\mathrm{Main}}_{\Lambda_k}\cap N^{\mathrm{mn}}_{\Lambda_k}
=
\{0\}
\),
which proves injectivity.
Since
\(
\overline\ell^{\mathrm{Main}}_{\Lambda_k}\neq 0
\),
its image
\(
\ell^{\mathrm{sel,mn}}_{\Lambda_k}
\)
is therefore nonzero, and the image of a one-dimensional line under an injective linear map is
again one-dimensional.
This proves items \textnormal{(1)} and \textnormal{(2)}.

Injectivity on
\(
\overline L^{\mathrm{Main}}_{\Lambda_k}
\)
implies that every element of
\(
L^{\mathrm{sel,mn}}_{\Lambda_k}
\)
has a unique expression
\(
a\,\ell^{\mathrm{sel,mn}}_{\Lambda_k}
\),
so the displayed formula defines
\(
\mathfrak R^{\mathrm{mn}}_{\Lambda_k}
\)
unambiguously.
Its coefficient
\(
\overline{\rho}^{\mathrm{dom}}_{\Lambda_k}
(\overline\ell^{\mathrm{Main}}_{\Lambda_k})
\)
is nonzero by hypothesis, and
\(
u^{\mathrm{LoG}}_{\Lambda_k}\neq 0
\)
by \cref{cor:compressed-export-selection}.
Therefore
\(
\mathfrak R^{\mathrm{mn}}_{\Lambda_k}
(a\,\ell^{\mathrm{sel,mn}}_{\Lambda_k})
=
0
\)
implies
\(
a=0
\),
so the map is faithful.
This proves item \textnormal{(3)}.

By construction of the descended involution,
\[
\Gamma^{\mathrm{mn}}_{\Lambda_k}
\!\bigl(
\ell^{\mathrm{sel,mn}}_{\Lambda_k}
\bigr)
=
\Gamma^{\mathrm{mn}}_{\Lambda_k}
\!\Bigl(
\pi^{\mathrm{mn}}_{\Lambda_k}
\!\bigl(
\overline\ell^{\mathrm{Main}}_{\Lambda_k}
\bigr)
\Bigr)
=
\pi^{\mathrm{mn}}_{\Lambda_k}
\!\bigl(
-\overline\ell^{\mathrm{Main}}_{\Lambda_k}
\bigr)
=
-\ell^{\mathrm{sel,mn}}_{\Lambda_k},
\]
which proves item \textnormal{(4)}.
\end{proof}

\begin{theorem}[Conditional mod-null reentry on the selected scalar line]
\label{thm:conditional-mod-null-reentry}
Fix an accepted tick \(k\) and the next accepted receiver step \(k+1\).
Assume the hypotheses and notation of
\cref{splitlift-thm:splitlift-scalar-descendant-lift}.
In particular, let
\[
\overline L^{\mathrm{Main}}_{\Lambda_{k+1}}
=
\mathbb R\,\overline\ell^{\mathrm{Main}}_{\Lambda_{k+1}}
\subset
\mathfrak D^-_{\Lambda_{k+1}}
\]
be the nonzero selected scalar descendant line supplied there.
Assume moreover a parity-stable mod-null datum transverse to
\(
\overline L^{\mathrm{Main}}_{\Lambda_{k+1}}
\)
in the sense of \cref{def:mod-null-transverse-datum}.
Define
\[
\ell^{\mathrm{sel,mn}}_{\Lambda_{k+1}}
:=
\pi^{\mathrm{mn}}_{\Lambda_{k+1}}
\!\bigl(
\overline\ell^{\mathrm{Main}}_{\Lambda_{k+1}}
\bigr)
\]
and
\[
L^{\mathrm{sel,mn}}_{\Lambda_{k+1}}
:=
\mathbb R\,\ell^{\mathrm{sel,mn}}_{\Lambda_{k+1}}
\subset
\mathcal V^{1,\mathrm{mn}}_{\Lambda_{k+1}}.
\]
Then:
\begin{enumerate}[leftmargin=2em]
\item \(L^{\mathrm{sel,mn}}_{\Lambda_{k+1}}\) is a well-defined one-dimensional selected descendant
  line in the mod-null quotient;
\item the induced mod-null selected realization map
  \[
  \mathfrak R^{\mathrm{mn}}_{\Lambda_{k+1}}:
  L^{\mathrm{sel,mn}}_{\Lambda_{k+1}}
  \longrightarrow
  \mathbb R\,u^{\mathrm{LoG}}_{\Lambda_{k+1}},
  \qquad
  \mathfrak R^{\mathrm{mn}}_{\Lambda_{k+1}}
  \!\bigl(
  a\,\ell^{\mathrm{sel,mn}}_{\Lambda_{k+1}}
  \bigr)
  :=
  a\,
  \overline{\rho}^{\mathrm{dom}}_{\Lambda_{k+1}}
  \!\bigl(
  \overline\ell^{\mathrm{Main}}_{\Lambda_{k+1}}
  \bigr)\,
  u^{\mathrm{LoG}}_{\Lambda_{k+1}}
  \]
  is well defined and faithful;
\item the descended involution acts by sign:
  \[
  \Gamma^{\mathrm{mn}}_{\Lambda_{k+1}}
  \!\bigl(
  \ell^{\mathrm{sel,mn}}_{\Lambda_{k+1}}
  \bigr)
  =
  -
  \ell^{\mathrm{sel,mn}}_{\Lambda_{k+1}}.
  \]
\end{enumerate}
\end{theorem}

\begin{proof}
By \cref{splitlift-thm:splitlift-scalar-descendant-lift},
\[
\overline\ell^{\mathrm{Main}}_{\Lambda_{k+1}}\neq 0
\qquad\text{and}\qquad
\overline{\rho}^{\mathrm{dom}}_{\Lambda_{k+1}}
\!\bigl(
\overline\ell^{\mathrm{Main}}_{\Lambda_{k+1}}
\bigr)
\neq 0.
\]
So every hypothesis of
\cref{lem:mod-null-transversality-reduction}
is in force at \(\Lambda_{k+1}\).
Applying that lemma gives items \textnormal{(1)}--\textnormal{(3)} immediately.
\end{proof}

\begin{corollary}[Canonical dominant-kernel discharge of the mod-null reentry step]
\label{cor:dominant-kernel-mod-null-reentry}
Fix an accepted tick \(k\) and the next accepted receiver step \(k+1\).
Assume the hypotheses and notation of
\cref{splitlift-thm:splitlift-scalar-descendant-lift}.
Define the canonical dominant-null datum
\[
N^{\mathrm{dom-null}}_{\Lambda_{k+1}}
:=
\ker \overline{\rho}^{\mathrm{dom}}_{\Lambda_{k+1}}
\subset
\mathfrak D^-_{\Lambda_{k+1}}.
\]
Then \(N^{\mathrm{dom-null}}_{\Lambda_{k+1}}\) is admissible as the mod-null datum of
\cref{thm:conditional-mod-null-reentry}, the transversality hypothesis of that theorem is
automatic, and hence the conclusions of that theorem hold with this canonical datum.
\end{corollary}

\begin{proof}
By \cref{splitlift-thm:splitlift-scalar-descendant-lift},
\[
\overline{\rho}^{\mathrm{dom}}_{\Lambda_{k+1}}
\!\bigl(
\overline\ell^{\mathrm{Main}}_{\Lambda_{k+1}}
\bigr)
\neq 0.
\]
So \cref{thm:canonical-dominant-null-datum} applies at \(\Lambda_{k+1}\).
It shows that
\(
N^{\mathrm{dom-null}}_{\Lambda_{k+1}}
\)
is a dominant-null parity-stable mod-null datum and that
\[
\overline L^{\mathrm{Main}}_{\Lambda_{k+1}}
\cap
N^{\mathrm{dom-null}}_{\Lambda_{k+1}}
=
\{0\}.
\]
This is exactly the admissibility and transversality clause required by
\cref{thm:conditional-mod-null-reentry}, so that theorem now applies directly with the canonical
datum.
\end{proof}

\begin{remark}[Conditional character of the reentry step]
\label{rem:conditional-mod-null-reentry}
\cref{thm:conditional-mod-null-reentry} is a quotient-level reentry statement.
After
\cref{lem:dominant-kernel-transversality,thm:canonical-dominant-null-datum,cor:dominant-kernel-mod-null-reentry},
the canonical dominant-null datum is fixed as
\(
\ker\overline{\rho}^{\mathrm{dom}}_{\Lambda_{k+1}}
\)
on \(\mathfrak D^-_{\Lambda_{k+1}}\).
What remains is passage to the mod-null quotient together with parity descent on the odd line.
Accordingly, the theorem should not be read as a statement on the raw descendant ambiguity quotient
or on any stronger physical quotient.
\end{remark}

\subsection*{Branchwise orientation transport}
\label{subsec:branchwise-orientation}

This subsection studies transport of the descended descendant sign along a connected accepted
branch. The sign can be transported once the descended descendant value varies continuously with the
cutoff. The compressed spectral-selection theorem already gives continuity of the dominant
projector, and \cref{thm:dominant-line-persistence} gives continuity of the dominant physical line.
We combine those inputs with continuity of the projected descendant response.

\begin{definition}[Branchwise descendant overlap function]
\label{def:branchwise-overlap-function}
Let \(I\) be a connected accepted-branch segment.
Assume dominant-line realization and quotient descent for every \(\Lambda\in I\).
Suppose there is a branchwise choice of representatives
\[
J^{\mathrm{desc}}_{\Lambda}\in[\Xi^{(1)}_{\Lambda}]
\qquad
(\Lambda\in I)
\]
such that the projected covariance response
\[
\Lambda\longmapsto
\Pi_{\mathrm{phys},\Lambda}\mathcal R_{\Lambda}(J^{\mathrm{desc}}_{\Lambda})
\]
is continuous in the Fisher topology.
Define the associated branchwise descendant overlap function by
\begin{equation}
m_I(\Lambda)
:=
\rho^{\mathrm{dom}}_{\Lambda}(J^{\mathrm{desc}}_{\Lambda})
=
\overline{\rho}^{\mathrm{dom}}_{\Lambda}\big([\![\Xi^{(1)}_{\Lambda}]\!]\big).
\label{eq:branchwise-overlap-function}
\end{equation}
\end{definition}

\begin{proposition}[Continuity of the descended overlap scalar]
\label{prop:branchwise-overlap-continuity}
Let \(I\) be a connected accepted-branch segment.
Assume:
\begin{enumerate}[leftmargin=2em]
\item the hypotheses of \cref{thm:dominant-line-persistence} on \(I\);
\item quotient descent holds for every \(\Lambda\in I\);
\item the branchwise choice of representatives in
  \cref{def:branchwise-overlap-function} exists.
\end{enumerate}
Then the scalar function \(m_I:I\to\mathbb R\) is well defined, independent of the chosen
branchwise representatives, and continuous.
Moreover,
\begin{equation}
|m_I(\Lambda)|
=
\nu_{\Lambda}
\qquad
\forall \Lambda\in I.
\label{eq:branchwise-overlap-margin}
\end{equation}
\end{proposition}

\begin{proof}
Independence of the chosen branchwise representatives is exactly the coset constancy proved in
\cref{cor:dominant-quotient-constancy}: if
\(
J'^{\mathrm{desc}}_{\Lambda}-J^{\mathrm{desc}}_{\Lambda}\in\mathcal A^-_{\Lambda}
\),
then
\[
\rho^{\mathrm{dom}}_{\Lambda}(J'^{\mathrm{desc}}_{\Lambda})
=
\rho^{\mathrm{dom}}_{\Lambda}(J^{\mathrm{desc}}_{\Lambda})
=
\overline{\rho}^{\mathrm{dom}}_{\Lambda}\big([\![\Xi^{(1)}_{\Lambda}]\!]\big).
\]
Thus \(m_I\) is well defined.

By \cref{thm:dominant-line-persistence}, the normalized dominant vector
\(
\widehat{\delta C}^{\mathrm{LoG}}_{\Lambda}
\)
and the dominant projector
\(
\Pi^{\mathrm{dom}}_{\Lambda}
\)
depend continuously on \(\Lambda\in I\).
By assumption, the response vector
\[
X_{\Lambda}
:=
\Pi_{\mathrm{phys},\Lambda}\mathcal R_{\Lambda}(J^{\mathrm{desc}}_{\Lambda})
\]
is also continuous on \(I\).
Using \cref{def:dominant-cov-projector}, we have
\[
\Pi^{\mathrm{dom}}_{\Lambda}X_{\Lambda}
=
m_I(\Lambda)\,\widehat{\delta C}^{\mathrm{LoG}}_{\Lambda}.
\]
The left-hand side is continuous because both factors are continuous, while
\(
\widehat{\delta C}^{\mathrm{LoG}}_{\Lambda}\neq 0
\)
for all \(\Lambda\in I\) by \cref{thm:dominant-line-persistence}.
Therefore the scalar coefficient \(m_I(\Lambda)\) is continuous.

Finally, \eqref{eq:branchwise-overlap-margin} is just the definition of \(\nu_{\Lambda}\) from
\cref{def:descendant-overlap-margin} together with \eqref{eq:branchwise-overlap-function}.
\end{proof}

\begin{theorem}[Branchwise orientation transport on the descendant class]
\label{thm:branchwise-orientation-transport}
Let \(I\) be a connected accepted-branch segment and assume the hypotheses of
\cref{prop:branchwise-overlap-continuity}.
Then:
\begin{enumerate}[leftmargin=2em]
\item on every connected subsegment \(I'\subset I\) such that
  \[
  \nu_{\Lambda}>0
  \qquad
  \forall \Lambda\in I',
  \]
  the sign
  \begin{equation}
  \varepsilon_{\Lambda}
  :=
  \operatorname{sgn} m_I(\Lambda)
  =
  \operatorname{sgn}\!\Big(
  \overline{\rho}^{\mathrm{dom}}_{\Lambda}\big([\![\Xi^{(1)}_{\Lambda}]\!]\big)
  \Big)
  \in\{\pm1\}
  \label{eq:branchwise-epsilon}
  \end{equation}
  is constant on \(I'\);
\item if \(\nu_{\Lambda_0}>0\) at one point \(\Lambda_0\in I\), then there exists a neighborhood
  \(I_0\subset I\) of \(\Lambda_0\) on which \(\varepsilon_{\Lambda}\) is constant and equal to
  \(\varepsilon_{\Lambda_0}\);
\item under the present hypotheses, a change of sign can occur only at a cutoff where the
  descendant overlap closes
  \(
  \nu_{\Lambda}=0
  \).
\end{enumerate}
In particular, once dominant-line realization, quotient descent, and nonzero overlap are
established on an accepted branch segment and the projected descendant response varies continuously
there, the branchwise sign statement of \cref{thm:active-response-transport-minimal} follows.
\end{theorem}

\begin{proof}
By \cref{prop:branchwise-overlap-continuity}, \(m_I(\Lambda)\) is a continuous real-valued function
on \(I\).
On any connected subsegment \(I'\) where \(\nu_{\Lambda}=|m_I(\Lambda)|>0\), the function \(m_I\)
never vanishes.
A continuous nonvanishing real-valued function on a connected set has constant sign, proving item
(1).

Item (2) is the local version of the same argument: if \(m_I(\Lambda_0)\neq 0\), continuity gives a
neighborhood \(I_0\) on which \(m_I\) keeps the sign it has at \(\Lambda_0\).

For item (3), any sign change would force \(m_I\) to pass through zero by the intermediate value
theorem, hence \(\nu_{\Lambda}=|m_I(\Lambda)|=0\) at some intermediate cutoff.
\end{proof}

\begin{corollary}[Local orientation persistence from one-tick activity]
\label{cor:branchwise-orientation-local}
Let \(I\) be a connected accepted-branch segment and assume the hypotheses of
\cref{prop:branchwise-overlap-continuity}.
If the transported descendant class is dominant-active at one cutoff \(\Lambda_0\in I\), i.e.
\[
\nu_{\Lambda_0}>0,
\]
then there exists a neighborhood \(I_0\subset I\) of \(\Lambda_0\) such that
\[
\varepsilon_{\Lambda}=\varepsilon_{\Lambda_0}
\qquad
\forall \Lambda\in I_0.
\]
\end{corollary}

\begin{proof}
This is exactly item (2) of \cref{thm:branchwise-orientation-transport}.
\end{proof}

\subsection*{Witness-indexed seam realization and scalar bridge}
\label{subsec:witness-seam-bridge}

The bridge used below factors through a common LS witness datum rather than a direct identification
of a UOM seam source with a descendant representative.
On the UOM side one realizes a filler tag together with a detector witness by an admissible odd
welded source, measures the resulting dominant physical pairing, and then compares that scalar with
the odd representative attached to the same witness.

\begin{definition}[Accepted record payload of a welded odd source]
\label{def:accepted-record-payload}
Fix an accepted tick \(k\).
For an admissible welded odd trace-channel source \(\tau\), let
\[
\rho_{\Lambda_k}[\tau]
\]
denote the receiver state at cutoff \(\Lambda_k\) produced by that source in the welded UOM channel.
Define its accepted record payload by
\begin{equation}
\mathcal P^{\mathrm{acc}}_{\Lambda_k}[\tau]
:=
\Pi_{\Delta(\Lambda_k)}\!\big(\rho_{\Lambda_k}[\tau]\big).
\label{eq:accepted-record-payload}
\end{equation}
This is the diagonal/record part retained by the finite-cutoff acceptance map.
\end{definition}

\begin{definition}[Witness-indexed seam realization at an accepted tick]
\label{def:witness-seam-realization}
Fix an accepted tick \(k\).
Let \(y\) be a witness on the LS difference object, and write \(\mathrm{flip}(y)\) for the induced
LS involution.
A \emph{witness-indexed seam realization} at \(\Lambda_k\) is a choice of admissible welded odd
trace-channel sources
\[
\tau^{f,y}_{\Lambda_k},
\qquad
f\in\{\mathrm{diag},\mathrm{bridge}\},
\]
such that:
\begin{enumerate}[leftmargin=2em]
\item \textbf{filler toggle changes the odd seam sign:}
  \[
  \tau^{\mathrm{bridge},y}_{\Lambda_k}
  =
  -\,\tau^{\mathrm{diag},y}_{\Lambda_k};
  \]
\item \textbf{witness flip changes the odd seam sign:}
  \[
  \tau^{f,\mathrm{flip}(y)}_{\Lambda_k}
  =
  -\,\tau^{f,y}_{\Lambda_k}
  \qquad
  \forall f\in\{\mathrm{diag},\mathrm{bridge}\};
  \]
\item \textbf{accepted record payload is unchanged under filler toggle:}
  \[
  \mathcal P^{\mathrm{acc}}_{\Lambda_k}\!\big[\tau^{\mathrm{bridge},y}_{\Lambda_k}\big]
  =
  \mathcal P^{\mathrm{acc}}_{\Lambda_k}\!\big[\tau^{\mathrm{diag},y}_{\Lambda_k}\big].
  \]
\end{enumerate}
\end{definition}

\begin{remark}[Accepted-record invariance versus seam-scalar use]
\label{rem:accepted-record-vs-seam}
The accepted-record invariance clause in \cref{def:witness-seam-realization} encodes the physical
indistinguishability of filler toggle at the accepted-record layer.
The seam-scalar theorems below use only the induced odd sign relations in the source channel.
\end{remark}

\begin{definition}[Raw and normalized seam pairings]
\label{def:raw-seam-pairing}
Fix an accepted tick \(k\) and a witness-indexed seam realization as in
\cref{def:witness-seam-realization}.
For \(f\in\{\mathrm{diag},\mathrm{bridge}\}\) and witness \(y\), define the \emph{raw seam
pairing} by
\begin{equation}
\widetilde s_{\Lambda_k}(f,y)
:=
\big\langle\!\big\langle
\Pi_{\mathrm{phys},\Lambda_k}\mathsf T_{\Lambda_k}\!\big[\tau^{f,y}_{\Lambda_k}\big],\,
\delta C^{\mathrm{LoG}}_{\Lambda_k}
\big\rangle\!\big\rangle_{C_{\ast,\Lambda_k}}.
\label{eq:raw-seam-pairing}
\end{equation}
Whenever dominant-line realization holds so that
\(
\widehat{\delta C}^{\mathrm{LoG}}_{\Lambda_k}
\)
is defined, also set
\begin{equation}
s_{\Lambda_k}(f,y)
:=
\big\langle\!\big\langle
\Pi_{\mathrm{phys},\Lambda_k}\mathsf T_{\Lambda_k}\!\big[\tau^{f,y}_{\Lambda_k}\big],\,
\widehat{\delta C}^{\mathrm{LoG}}_{\Lambda_k}
\big\rangle\!\big\rangle_{C_{\ast,\Lambda_k}}.
\label{eq:normalized-seam-pairing}
\end{equation}
\end{definition}

\begin{proposition}[Filler and witness parity of the seam scalar]
\label{prop:seam-scalar-signs}
Fix an accepted tick \(k\) and assume a witness-indexed seam realization at \(\Lambda_k\).
Then the raw seam pairing satisfies
\begin{equation}
\widetilde s_{\Lambda_k}(\mathrm{bridge},y)
=
-\,\widetilde s_{\Lambda_k}(\mathrm{diag},y),
\qquad
\widetilde s_{\Lambda_k}(f,\mathrm{flip}(y))
=
-\,\widetilde s_{\Lambda_k}(f,y).
\label{eq:raw-seam-signs}
\end{equation}
If dominant-line realization also holds, then the same identities hold for the normalized seam
pairing:
\begin{equation}
s_{\Lambda_k}(\mathrm{bridge},y)
=
-\,s_{\Lambda_k}(\mathrm{diag},y),
\qquad
s_{\Lambda_k}(f,\mathrm{flip}(y))
=
-\,s_{\Lambda_k}(f,y).
\label{eq:normalized-seam-signs}
\end{equation}
\end{proposition}

\begin{proof}
The map
\(
\tau\mapsto \Pi_{\mathrm{phys},\Lambda_k}\mathsf T_{\Lambda_k}[\tau]
\)
is linear because \(\mathsf T_{\Lambda_k}\) and \(\Pi_{\mathrm{phys},\Lambda_k}\) are linear.
The Fisher pairing is bilinear.
Applying the sign relations of \cref{def:witness-seam-realization} therefore gives
\[
\Pi_{\mathrm{phys},\Lambda_k}\mathsf T_{\Lambda_k}\!\big[\tau^{\mathrm{bridge},y}_{\Lambda_k}\big]
=
-\,\Pi_{\mathrm{phys},\Lambda_k}\mathsf T_{\Lambda_k}\!\big[\tau^{\mathrm{diag},y}_{\Lambda_k}\big]
\]
and
\[
\Pi_{\mathrm{phys},\Lambda_k}\mathsf T_{\Lambda_k}\!\big[\tau^{f,\mathrm{flip}(y)}_{\Lambda_k}\big]
=
-\,\Pi_{\mathrm{phys},\Lambda_k}\mathsf T_{\Lambda_k}\!\big[\tau^{f,y}_{\Lambda_k}\big].
\]
Pairing with \(\delta C^{\mathrm{LoG}}_{\Lambda_k}\) proves \eqref{eq:raw-seam-signs}.
If dominant-line realization holds, then
\(
\widehat{\delta C}^{\mathrm{LoG}}_{\Lambda_k}
\)
is a fixed nonzero normalization of the same vector, so the same argument proves
\eqref{eq:normalized-seam-signs}.
\end{proof}

\begin{definition}[Detector-factorized seam coupling]
\label{def:detector-factorized-seam}
Fix an accepted tick \(k\).
A witness-indexed seam realization is called \emph{detector-factorized} if there exists a single
admissible welded odd trace-channel source
\[
\tau^{\mathrm{unit}}_{\Lambda_k}
\]
and the LS detector factor
\[
\overline{\mathrm{obs}}
:
\text{LS difference object}\to\mathbb R
\]
such that, for every filler tag \(f\) and witness \(y\),
\begin{equation}
\tau^{f,y}_{\Lambda_k}
=
\sigma_{\mathrm{fill}}(f)\,
\overline{\mathrm{obs}}(y)\,
\tau^{\mathrm{unit}}_{\Lambda_k},
\qquad
\sigma_{\mathrm{fill}}(\mathrm{diag})=+1,\quad
\sigma_{\mathrm{fill}}(\mathrm{bridge})=-1.
\label{eq:detector-factorized-seam}
\end{equation}
\end{definition}

\begin{theorem}[Seam faithfulness scalar]
\label{thm:seam-faithfulness-scalar}
Fix an accepted tick \(k\) and assume a detector-factorized witness-indexed seam realization at
\(\Lambda_k\).
Define the fixed-cutoff seam coefficient by
\begin{equation}
\widetilde c_{\Lambda_k}
:=
\big\langle\!\big\langle
\Pi_{\mathrm{phys},\Lambda_k}\mathsf T_{\Lambda_k}\!\big[\tau^{\mathrm{unit}}_{\Lambda_k}\big],\,
\delta C^{\mathrm{LoG}}_{\Lambda_k}
\big\rangle\!\big\rangle_{C_{\ast,\Lambda_k}}.
\label{eq:seam-faithfulness-coeff}
\end{equation}
Then, for every filler tag \(f\) and witness \(y\),
\begin{equation}
\widetilde s_{\Lambda_k}(f,y)
=
\sigma_{\mathrm{fill}}(f)\,
\widetilde c_{\Lambda_k}\,
\overline{\mathrm{obs}}(y).
\label{eq:raw-seam-factorization}
\end{equation}
In particular, if
\[
\widetilde c_{\Lambda_k}\neq 0
\qquad\text{and}\qquad
\overline{\mathrm{obs}}(y)\neq 0,
\]
then the raw seam pairing is nonzero:
\[
\widetilde s_{\Lambda_k}(f,y)\neq 0.
\]
Consequently the compressed welded LoG line is transversely physical at \(\Lambda_k\), hence
dominant-line realization holds at that cutoff by \cref{prop:dominant-line-equivalent}.
When dominant-line realization holds, the normalized seam coefficient
\begin{equation}
c_{\Lambda_k}
:=
\kappa_{\Lambda_k}^{-1/2}\,\widetilde c_{\Lambda_k}
\label{eq:normalized-seam-coeff}
\end{equation}
is defined and
\begin{equation}
s_{\Lambda_k}(f,y)
=
\sigma_{\mathrm{fill}}(f)\,
c_{\Lambda_k}\,
\overline{\mathrm{obs}}(y).
\label{eq:normalized-seam-factorization}
\end{equation}
\end{theorem}

\begin{proof}
Substitute the detector-factorized form \eqref{eq:detector-factorized-seam} into the definition
\eqref{eq:raw-seam-pairing}.
Linearity of \(\mathsf T_{\Lambda_k}\) and bilinearity of the Fisher pairing give
\[
\widetilde s_{\Lambda_k}(f,y)
=
\sigma_{\mathrm{fill}}(f)\,\overline{\mathrm{obs}}(y)\,
\big\langle\!\big\langle
\Pi_{\mathrm{phys},\Lambda_k}\mathsf T_{\Lambda_k}\!\big[\tau^{\mathrm{unit}}_{\Lambda_k}\big],\,
\delta C^{\mathrm{LoG}}_{\Lambda_k}
\big\rangle\!\big\rangle_{C_{\ast,\Lambda_k}},
\]
which is exactly \eqref{eq:raw-seam-factorization}.

If \(\widetilde c_{\Lambda_k}\neq 0\) and \(\overline{\mathrm{obs}}(y)\neq 0\), then
\(
\widetilde s_{\Lambda_k}(f,y)\neq 0
\).
Set
\[
B_{\Lambda_k}
:=
\Pi_{\mathrm{phys},\Lambda_k}\mathsf T_{\Lambda_k}\!\big[\tau^{f,y}_{\Lambda_k}\big].
\]
Then
\[
\big\langle\!\big\langle
\delta C^{\mathrm{LoG}}_{\Lambda_k},\,B_{\Lambda_k}
\big\rangle\!\big\rangle_{C_{\ast,\Lambda_k}}
=
\widetilde s_{\Lambda_k}(f,y)
\neq 0.
\]
By item (4) of \cref{prop:dominant-line-equivalent}, this proves transversality of the compressed
line and hence dominant-line realization at \(\Lambda_k\).

Once dominant-line realization holds, divide \eqref{eq:raw-seam-factorization} by
\(
\|\delta C^{\mathrm{LoG}}_{\Lambda_k}\|_{C_{\ast,\Lambda_k}}
=
\kappa_{\Lambda_k}^{1/2}
\)
to obtain \eqref{eq:normalized-seam-factorization}.
\end{proof}

\begin{remark}[Downstream role of the seam-coefficient hypothesis]
The fixed-cutoff nonvanishing hypothesis
\(
\widetilde c_{\Lambda_k}\neq 0
\)
is exactly the input that, once dominant-line realization defines
\(
c_{\Lambda_k}
\),
is imported downstream as the nondegeneracy clause
\(
c_{\Lambda_k}\neq 0
\)
inside the delayed selected-line package.
This note does not derive that clause from split-threshold analysis alone.
\end{remark}

\begin{definition}[Odd source representability on the characterization layer]
\label{def:odd-source-representability}
Fix an accepted tick \(k\).
Let \(\tau\) be an admissible welded odd trace-channel source.
We say that \(\tau\) is \emph{representable on the characterization layer} if there exists an odd
sufficiently local generator
\[
F_{\tau,\Lambda_k}\in E^-_{\Lambda_k}
\]
such that its intrinsic physical covariance response agrees with the UOM welded response:
\begin{equation}
\Pi_{\mathrm{phys},\Lambda_k}\mathcal R_{\Lambda_k}\!\big(F_{\tau,\Lambda_k}\big)
=
\Pi_{\mathrm{phys},\Lambda_k}\mathsf T_{\Lambda_k}[\tau].
\label{eq:odd-source-representability}
\end{equation}
This is the strong source-insertion realization statement.
\end{definition}

\begin{theorem}[UOM source-insertion representability theorem]
\label{thm:odd-source-representability}
Fix an accepted tick \(k\).
Let \(\tau\) be an admissible welded odd trace-channel source at cutoff \(\Lambda_k\), in the
finite-cutoff UOM sense: its receiver-head profile is band-limited at \(\Lambda_k\) and compactly
supported inside the accepted slab.
Then there exists a canonically associated odd sufficiently local generator
\[
F_{\tau,\Lambda_k}\in E^-_{\Lambda_k}
\]
such that
\begin{equation}
\Pi_{\mathrm{phys},\Lambda_k}\mathcal R_{\Lambda_k}\!\big(F_{\tau,\Lambda_k}\big)
=
\Pi_{\mathrm{phys},\Lambda_k}\mathsf T_{\Lambda_k}[\tau].
\label{eq:odd-source-representability-thm}
\end{equation}
In particular, every witness-indexed seam source used in the present note is representable on the
characterization layer.
\end{theorem}

\begin{proof}
By the finite-cutoff UOM Peierls package, admissible receiver-head sources at cutoff \(\Lambda_k\)
are band-limited boundary test data and the sufficiently local class
\(\mathcal F_{\rm s.loc}^{\Lambda_k}\) contains all compact smearings of local boundary densities.
Let
\[
F_{\tau,\Lambda_k}
\]
be the linear source functional obtained by pairing the odd trace-channel source profile \(\tau\)
with the corresponding odd trace-channel boundary field on the receiver head.
Because this is the linear case of the sufficiently local form and \(\tau\) is physically odd,
\[
F_{\tau,\Lambda_k}\in E^-_{\Lambda_k}.
\]

Now apply the UOM source insertion map from the finite-cutoff retarded-kernel package.
The source \(\tau\) is sent to the bulk forcing term in the linearized equation; the unique
retarded solution is obtained by the retarded Green operator; and boundary restriction gives the
retarded covariance variation. By the UOM identification of the induced retarded kernel with the SK
mixed second derivative, this covariance variation is exactly the welded source-to-covariance
response
\[
\mathsf T_{\Lambda_k}[\tau].
\]

On the other hand, by the definition of the intrinsic covariance-response map, the matrix elements
of
\(
\mathcal R_{\Lambda_k}(F_{\tau,\Lambda_k})
\)
are computed by the Peierls action of \(F_{\tau,\Lambda_k}\) on the sufficiently local quadratic
probes representing boundary covariance entries.
That Peierls variation is precisely the same retarded second-moment variation generated by the
source insertion \(\tau\).
Therefore
\[
\mathcal R_{\Lambda_k}(F_{\tau,\Lambda_k})
=
\mathsf T_{\Lambda_k}[\tau]
\]
as covariance tangents.
Applying the physical projector \(\Pi_{\mathrm{phys},\Lambda_k}\) gives
\eqref{eq:odd-source-representability-thm}.

The witness-indexed seam sources of \cref{def:witness-seam-realization} are, by construction, live
admissible welded odd trace-channel sources at the accepted tick, so the same argument applies to
each of them.
\end{proof}

\begin{proposition}[Dominant pairing for a representable odd source]
\label{prop:odd-source-pairing}
Fix an accepted tick \(k\) and let \(\tau\) be an admissible welded odd trace-channel source that is
representable in the sense of \cref{thm:odd-source-representability}.
Then
\begin{equation}
\big\langle\!\big\langle
\Pi_{\mathrm{phys},\Lambda_k}\mathcal R_{\Lambda_k}\!\big(F_{\tau,\Lambda_k}\big),\,
\delta C^{\mathrm{LoG}}_{\Lambda_k}
\big\rangle\!\big\rangle_{C_{\ast,\Lambda_k}}
=
\big\langle\!\big\langle
\Pi_{\mathrm{phys},\Lambda_k}\mathsf T_{\Lambda_k}[\tau],\,
\delta C^{\mathrm{LoG}}_{\Lambda_k}
\big\rangle\!\big\rangle_{C_{\ast,\Lambda_k}}.
\label{eq:representable-raw-pairing}
\end{equation}
If dominant-line realization also holds, then
\begin{equation}
\rho^{\mathrm{dom}}_{\Lambda_k}\!\big(F_{\tau,\Lambda_k}\big)
=
\big\langle\!\big\langle
\Pi_{\mathrm{phys},\Lambda_k}\mathsf T_{\Lambda_k}[\tau],\,
\widehat{\delta C}^{\mathrm{LoG}}_{\Lambda_k}
\big\rangle\!\big\rangle_{C_{\ast,\Lambda_k}}.
\label{eq:representable-normalized-pairing}
\end{equation}
\end{proposition}

\begin{proof}
Equation \eqref{eq:odd-source-representability} identifies the two physical covariance tangents.
Pairing both sides with \(\delta C^{\mathrm{LoG}}_{\Lambda_k}\) gives
\eqref{eq:representable-raw-pairing}.
If dominant-line realization holds, pair with the normalized vector
\(
\widehat{\delta C}^{\mathrm{LoG}}_{\Lambda_k}
\)
instead and use the definition \eqref{eq:rho-dom-def} of
\(
\rho^{\mathrm{dom}}_{\Lambda_k}
\)
to obtain \eqref{eq:representable-normalized-pairing}.
\end{proof}

\begin{corollary}[Unit welded odd source representability and scalar identification]
\label{cor:unit-source-representability}
Fix an accepted tick \(k\).
Assume a detector-factorized witness-indexed seam realization at \(\Lambda_k\) in the sense of
\cref{def:detector-factorized-seam}, and assume dominant-line realization at \(\Lambda_k\).
Then the unit welded odd source
\(
\tau^{\mathrm{unit}}_{\Lambda_k}
\)
is representable on the characterization layer, with canonical representable generator
\[
F_{\tau^{\mathrm{unit}}_{\Lambda_k},\Lambda_k}\in E^-_{\Lambda_k},
\]
and
\begin{equation}
\rho^{\mathrm{dom}}_{\Lambda_k}
\!\bigl(
F_{\tau^{\mathrm{unit}}_{\Lambda_k},\Lambda_k}
\bigr)
=
\big\langle\!\big\langle
\Pi_{\mathrm{phys},\Lambda_k}\mathsf T_{\Lambda_k}\!\big[\tau^{\mathrm{unit}}_{\Lambda_k}\big],\,
\widehat{\delta C}^{\mathrm{LoG}}_{\Lambda_k}
\big\rangle\!\big\rangle_{C_{\ast,\Lambda_k}}
=
c_{\Lambda_k}.
\label{eq:unit-source-representability-scalar}
\end{equation}
\end{corollary}

\begin{proof}
Representability of
\(
\tau^{\mathrm{unit}}_{\Lambda_k}
\)
is the specialization of \cref{thm:odd-source-representability} to the admissible welded odd trace
source singled out by \cref{def:detector-factorized-seam}.
Because dominant-line realization holds, \cref{prop:odd-source-pairing} applies to that
representable source and gives the first equality in \eqref{eq:unit-source-representability-scalar}.
The second equality is exactly the definition of the normalized seam coefficient
\(
c_{\Lambda_k}
\)
from \cref{thm:seam-faithfulness-scalar}.
\end{proof}

\begin{definition}[Scalar comparison shadow of the witness-indexed bridge]
\label{def:witness-descendant-comparison}
Fix an accepted tick \(k\).
Let \(y\) be an LS witness and let
\(
\tau^{f,y}_{\Lambda_k}
\)
be a witness-indexed seam realization source for a filler tag \(f\).
A \emph{scalar comparison shadow} for that data is a choice of representative
\[
J^{f,y}_{\Lambda_k}\in[\Xi^{(1)}_{\Lambda_k}]
\]
such that the dominant raw pairing of the 2T-side representative agrees with the dominant raw
pairing of the corresponding representable UOM source:
\begin{equation}
\big\langle\!\big\langle
\Pi_{\mathrm{phys},\Lambda_k}\mathcal R_{\Lambda_k}\!\big(J^{f,y}_{\Lambda_k}\big),\,
\delta C^{\mathrm{LoG}}_{\Lambda_k}
\big\rangle\!\big\rangle_{C_{\ast,\Lambda_k}}
=
\big\langle\!\big\langle
\Pi_{\mathrm{phys},\Lambda_k}\mathcal R_{\Lambda_k}\!\big(F_{\tau^{f,y}_{\Lambda_k},\Lambda_k}\big),\,
\delta C^{\mathrm{LoG}}_{\Lambda_k}
\big\rangle\!\big\rangle_{C_{\ast,\Lambda_k}}.
\label{eq:witness-descendant-comparison}
\end{equation}
When dominant-line realization holds, this is equivalently the equality of normalized dominant
pairings
\begin{equation}
\rho^{\mathrm{dom}}_{\Lambda_k}\!\big(J^{f,y}_{\Lambda_k}\big)
=
\rho^{\mathrm{dom}}_{\Lambda_k}\!\big(F_{\tau^{f,y}_{\Lambda_k},\Lambda_k}\big).
\label{eq:witness-descendant-comparison-normalized}
\end{equation}
\end{definition}

\begin{definition}[Witness-indexed source-compatible descendant realization]
\label{def:source-compatible-witness-realization}
Fix an accepted tick \(k\), a filler tag \(f\), a witness \(y\), and a witness-indexed seam source
\(
\tau^{f,y}_{\Lambda_k}
\).
A representative
\[
J^{f,y}_{\Lambda_k}\in[\Xi^{(1)}_{\Lambda_k}]
\]
is called a \emph{witness-indexed source-compatible descendant realization} if it realizes, on the
same LS witness datum, the Limited Scope output and odd Hamiltonization clauses from item
\textbf{(M1)} and item \textbf{(M4)} of Definition~\ref{main-def:uom-main-theorem-assumptions} and differs
from the source-insertion representative only by an ambiguity direction:
\begin{equation}
J^{f,y}_{\Lambda_k}
-
F_{\tau^{f,y}_{\Lambda_k},\Lambda_k}
\in
\mathcal A^-_{\Lambda_k},
\label{eq:source-compatible-witness-realization}
\end{equation}
where
\(
F_{\tau^{f,y}_{\Lambda_k},\Lambda_k}
\)
is the representable source generator furnished by
\cref{thm:odd-source-representability}.
Equivalently, the witness-indexed transported descendant class and the witness-indexed UOM
source-insertion realization define the same point of the ambiguity quotient
\(
E^-_{\Lambda_k}/\mathcal A^-_{\Lambda_k}
\).
\end{definition}
\begin{theorem}[Same-witness descended-scalar comparison from a source-compatible descendant realization]
\label{thm:witness-descended-scalar-comparison}
Fix an accepted tick \(k\), assume dominant-line realization at \(\Lambda_k\), and let
\(
\tau^{f,y}_{\Lambda_k}
\)
be a witness-indexed seam realization source.
Assume:
\begin{enumerate}[leftmargin=2em]
\item the source \(\tau^{f,y}_{\Lambda_k}\) is represented by the canonical generator
  \(
  F_{\tau^{f,y}_{\Lambda_k},\Lambda_k}
  \)
  from \cref{thm:odd-source-representability};
\item there exists a witness-indexed source-compatible descendant realization
  \(
  J^{f,y}_{\Lambda_k}\in[\Xi^{(1)}_{\Lambda_k}]
  \)
  in the sense of \cref{def:source-compatible-witness-realization};
\item the three vanishing statements of \cref{prop:dominant-quotient-generators} hold at
  \(\Lambda_k\).
\end{enumerate}
Then the raw and normalized comparison identities of
\cref{def:witness-descendant-comparison} hold:
\begin{align}
\big\langle\!\big\langle
\Pi_{\mathrm{phys},\Lambda_k}\mathcal R_{\Lambda_k}\!\big(J^{f,y}_{\Lambda_k}\big),\,
\delta C^{\mathrm{LoG}}_{\Lambda_k}
\big\rangle\!\big\rangle_{C_{\ast,\Lambda_k}}
&=
\big\langle\!\big\langle
\Pi_{\mathrm{phys},\Lambda_k}\mathcal R_{\Lambda_k}\!\big(F_{\tau^{f,y}_{\Lambda_k},\Lambda_k}\big),\,
\delta C^{\mathrm{LoG}}_{\Lambda_k}
\big\rangle\!\big\rangle_{C_{\ast,\Lambda_k}},
\label{eq:witness-descendant-comparison-thm-raw}
\\
\rho^{\mathrm{dom}}_{\Lambda_k}\!\big(J^{f,y}_{\Lambda_k}\big)
&=
\rho^{\mathrm{dom}}_{\Lambda_k}\!\big(F_{\tau^{f,y}_{\Lambda_k},\Lambda_k}\big).
\label{eq:witness-descendant-comparison-thm-normalized}
\end{align}
Moreover,
\begin{equation}
\rho^{\mathrm{dom}}_{\Lambda_k}\!\big(F_{\tau^{f,y}_{\Lambda_k},\Lambda_k}\big)
=
\big\langle\!\big\langle
\Pi_{\mathrm{phys},\Lambda_k}\mathsf T_{\Lambda_k}\!\big[\tau^{f,y}_{\Lambda_k}\big],\,
\widehat{\delta C}^{\mathrm{LoG}}_{\Lambda_k}
\big\rangle\!\big\rangle_{C_{\ast,\Lambda_k}}.
\label{eq:witness-descendant-comparison-thm-seam}
\end{equation}
If quotient descent also holds at \(\Lambda_k\), then the common value equals the descended
dominant scalar on the transported descendant affine class:
\begin{equation}
\overline{\rho}^{\mathrm{dom}}_{\Lambda_k}\big([\![\Xi^{(1)}_{\Lambda_k}]\!]\big)
=
\rho^{\mathrm{dom}}_{\Lambda_k}\!\big(J^{f,y}_{\Lambda_k}\big)
=
\rho^{\mathrm{dom}}_{\Lambda_k}\!\big(F_{\tau^{f,y}_{\Lambda_k},\Lambda_k}\big).
\label{eq:witness-descendant-comparison-thm-descended}
\end{equation}
\end{theorem}

\begin{proof}
Set
\[
A^{f,y}_{\Lambda_k}
:=
J^{f,y}_{\Lambda_k}
-
F_{\tau^{f,y}_{\Lambda_k},\Lambda_k}.
\]
By source compatibility,
\(
A^{f,y}_{\Lambda_k}\in\mathcal A^-_{\Lambda_k}
\).

Because dominant-line realization holds and the three vanishing statements of
\cref{prop:dominant-quotient-generators} are assumed, that proposition gives
\[
\mathcal A^-_{\Lambda_k}\subset\ker\rho^{\mathrm{dom}}_{\Lambda_k}.
\]
Hence
\[
\rho^{\mathrm{dom}}_{\Lambda_k}\!\big(J^{f,y}_{\Lambda_k}\big)
=
\rho^{\mathrm{dom}}_{\Lambda_k}\!\big(
F_{\tau^{f,y}_{\Lambda_k},\Lambda_k}
+
A^{f,y}_{\Lambda_k}
\big)
=
\rho^{\mathrm{dom}}_{\Lambda_k}\!\big(F_{\tau^{f,y}_{\Lambda_k},\Lambda_k}\big),
\]
proving \eqref{eq:witness-descendant-comparison-thm-normalized}.

Since dominant-line realization implies
\(
\delta C^{\mathrm{LoG}}_{\Lambda_k}\neq 0
\)
and
\(
\rho^{\mathrm{dom}}_{\Lambda_k}
=
\kappa_{\Lambda_k}^{-1/2}\widetilde{\rho}^{\mathrm{dom}}_{\Lambda_k},
\)
the raw pairing equality \eqref{eq:witness-descendant-comparison-thm-raw} is just the same identity
multiplied by the positive factor \(\kappa_{\Lambda_k}^{1/2}\).

Equation \eqref{eq:witness-descendant-comparison-thm-seam} is exactly the normalized pairing formula
of \cref{prop:odd-source-pairing} applied to the representable source
\(
\tau^{f,y}_{\Lambda_k}
\).

If quotient descent also holds, then \cref{cor:dominant-quotient-constancy} gives
\[
\rho^{\mathrm{dom}}_{\Lambda_k}\!\big(J^{f,y}_{\Lambda_k}\big)
=
\overline{\rho}^{\mathrm{dom}}_{\Lambda_k}\big([\![\Xi^{(1)}_{\Lambda_k}]\!]\big)
\]
for every representative
\(
J^{f,y}_{\Lambda_k}\in[\Xi^{(1)}_{\Lambda_k}],
\)
which yields \eqref{eq:witness-descendant-comparison-thm-descended} after combining with
\eqref{eq:witness-descendant-comparison-thm-normalized}.
\end{proof}

\begin{theorem}[Fixed-cutoff seam witness criterion for descendant overlap]
\label{thm:fixed-cutoff-seam-witness}
Fix an accepted tick \(k\).
Assume:
\begin{enumerate}[leftmargin=2em]
\item a detector-factorized witness-indexed seam realization at \(\Lambda_k\);
\item a witness \(y\) with
  \[
  \overline{\mathrm{obs}}(y)\neq 0;
  \]
\item the seam coefficient is nonzero,
  \[
  \widetilde c_{\Lambda_k}\neq 0;
  \]
\item for the chosen filler tag \(f\), the source \(\tau^{f,y}_{\Lambda_k}\) is representable on the
  characterization layer in the sense of \cref{def:odd-source-representability};
\item there exists a witness-indexed source-compatible descendant realization in the sense of
  \cref{def:source-compatible-witness-realization}.
\end{enumerate}
Then:
\begin{enumerate}[leftmargin=2em]
\item the compressed welded LoG line is transversely physical at \(\Lambda_k\), hence dominant-line
  realization holds at that cutoff;
\item the chosen 2T representative satisfies
  \[
  \rho^{\mathrm{dom}}_{\Lambda_k}\!\big(J^{f,y}_{\Lambda_k}\big)\neq 0;
  \]
\item therefore the transported descendant class has nonzero overlap with the canonically selected
  compressed dominant odd line at \(\Lambda_k\);
\item if quotient descent also holds at \(\Lambda_k\), then
  \cref{thm:descendant-overlap} applies and the transported descendant class is dominant-active with
  well-defined sign \(\varepsilon_{\Lambda_k}\).
\end{enumerate}
\end{theorem}

\begin{proof}
By \cref{thm:seam-faithfulness-scalar},
\[
\widetilde s_{\Lambda_k}(f,y)
=
\sigma_{\mathrm{fill}}(f)\,
\widetilde c_{\Lambda_k}\,
\overline{\mathrm{obs}}(y)
\neq 0.
\]
The same theorem therefore gives item (1): the compressed line is transversely physical, hence
dominant-line realization holds at \(\Lambda_k\).

By \cref{prop:odd-source-pairing},
\[
\big\langle\!\big\langle
\Pi_{\mathrm{phys},\Lambda_k}\mathcal R_{\Lambda_k}\!\big(F_{\tau^{f,y}_{\Lambda_k},\Lambda_k}\big),\,
\delta C^{\mathrm{LoG}}_{\Lambda_k}
\big\rangle\!\big\rangle_{C_{\ast,\Lambda_k}}
=
\widetilde s_{\Lambda_k}(f,y)
\neq 0.
\]
Applying the comparison identity \eqref{eq:witness-descendant-comparison} now gives
\[
\big\langle\!\big\langle
\Pi_{\mathrm{phys},\Lambda_k}\mathcal R_{\Lambda_k}\!\big(J^{f,y}_{\Lambda_k}\big),\,
\delta C^{\mathrm{LoG}}_{\Lambda_k}
\big\rangle\!\big\rangle_{C_{\ast,\Lambda_k}}
\neq 0.
\]
Because dominant-line realization has already been established, \(\delta C^{\mathrm{LoG}}_{\Lambda_k}\)
is nonzero and \(\widehat{\delta C}^{\mathrm{LoG}}_{\Lambda_k}\) is defined.
By \cref{prop:descendant-overlap-witness}, this is equivalent to
\(
\rho^{\mathrm{dom}}_{\Lambda_k}(J^{f,y}_{\Lambda_k})\neq 0
\),
proving item (2).

Item (3) is exactly the nonzero-overlap condition from
\cref{prop:descendant-overlap-equivalent}.
If quotient descent is also available, apply \cref{thm:descendant-overlap} to obtain item (4).
\end{proof}

\subsection*{Smooth-taper branch continuity and witness transport}
\label{subsec:smooth-taper-branch-continuity}

The previous fixed-cutoff theorem isolates the nonzero witness needed for activity.
To transport that witness along an accepted branch one needs a continuity statement for the welded
response.
At theorem level, the clean analytic choice is the smooth receiver taper.
The hard projector may still be used, but only on branch segments where its rank is constant.

\begin{lemma}[Smooth-taper continuity of the receiver band projector]
\label{lem:smooth-taper-projector}
Let \(I\) be a connected accepted-branch segment.
Suppose the receiver band is implemented by a smooth taper
\[
P^{\mathrm{(rec)}}_{\mathrm{band}}(\Lambda)
=
\sum_{\ell,m} w_\ell(\Lambda)\,|Y_{\ell m}\rangle\langle Y_{\ell m}|,
\qquad
0\le w_\ell(\Lambda)\le 1,
\]
on the receiver one-particle space, where
\[
\sup_{\ell\ge 0}|w_\ell(\Lambda')-w_\ell(\Lambda)|
\longrightarrow 0
\qquad
\text{as }\Lambda'\to\Lambda
\]
for every \(\Lambda\in I\).
Then
\[
\Lambda\longmapsto P^{\mathrm{(rec)}}_{\mathrm{band}}(\Lambda)
\]
is continuous on \(I\) in operator norm.
If instead the hard projector
\(
\sum_{\ell\le \ell_\star(\Lambda)}|Y_{\ell m}\rangle\langle Y_{\ell m}|
\)
is used, the same conclusion holds on every connected subsegment of \(I\) on which
\(\ell_\star(\Lambda)\) is constant.
\end{lemma}

\begin{proof}
In the spherical-harmonic basis \(\{Y_{\ell m}\}\), the difference
\[
P^{\mathrm{(rec)}}_{\mathrm{band}}(\Lambda')
-
P^{\mathrm{(rec)}}_{\mathrm{band}}(\Lambda)
\]
is diagonal with eigenvalues \(w_\ell(\Lambda')-w_\ell(\Lambda)\).
Therefore
\[
\big\|
P^{\mathrm{(rec)}}_{\mathrm{band}}(\Lambda')
-
P^{\mathrm{(rec)}}_{\mathrm{band}}(\Lambda)
\big\|_{\mathrm{op}}
=
\sup_{\ell\ge 0}|w_\ell(\Lambda')-w_\ell(\Lambda)|.
\]
The assumed uniform continuity of the taper coefficients is exactly the statement that this operator
norm tends to zero as \(\Lambda'\to\Lambda\).

For the hard projector, the coefficients take only the values \(0\) and \(1\), so discontinuities can
occur only when \(\ell_\star(\Lambda)\) jumps.
On a subsegment where \(\ell_\star(\Lambda)\) is constant, the projector is constant and therefore
continuous.
\end{proof}

\begin{theorem}[Smooth-taper branch continuity of the welded physical response]
\label{thm:smooth-taper-welded-response}
Let \(I\) be a connected accepted-branch segment.
Write the welded source-to-covariance map in the UOM factorized form
\begin{equation}
\mathsf T_\Lambda
=
\Pi_{L(\Lambda)}M_{\varepsilon,\Lambda}
\circ
\mathsf K_{{\rm comp},\Lambda}
\circ
\mathsf G_{{\rm ret},\Lambda}
\circ
\mathsf W_{{\rm src},\Lambda}.
\label{eq:factorized-welded-map}
\end{equation}
Assume:
\begin{enumerate}[leftmargin=2em]
\item the receiver band is implemented by a smooth taper satisfying
  \cref{lem:smooth-taper-projector};
\item the operator families
  \[
  \Lambda\longmapsto
  \Pi_{L(\Lambda)}M_{\varepsilon,\Lambda},\qquad
  \Lambda\longmapsto
  \mathsf K_{{\rm comp},\Lambda},\qquad
  \Lambda\longmapsto
  \mathsf G_{{\rm ret},\Lambda},\qquad
  \Lambda\longmapsto
  \mathsf W_{{\rm src},\Lambda},\qquad
  \Lambda\longmapsto
  \Pi_{\mathrm{phys},\Lambda}
  \]
  are continuous on \(I\) in operator norm on the corresponding band-limited source and covariance
  tangent spaces;
\item the source family
  \[
  \Lambda\longmapsto \tau_\Lambda
  \]
  is continuous on \(I\) in the welded source norm.
\end{enumerate}
Then the physical covariance response
\begin{equation}
X_\Lambda[\tau_\Lambda]
:=
\Pi_{\mathrm{phys},\Lambda}\mathsf T_\Lambda[\tau_\Lambda]
\label{eq:XLambda-tau-def}
\end{equation}
depends continuously on \(\Lambda\in I\) in the Fisher norm.
If the hard receiver projector is used instead, the same conclusion holds on each connected
subsegment where \(\ell_\star(\Lambda)\) is constant.
\end{theorem}

\begin{proof}
By \cref{lem:smooth-taper-projector}, the first factor
\(
\Pi_{L(\Lambda)}M_{\varepsilon,\Lambda}
\)
is operator-norm continuous along the branch.
By assumption, the remaining three UOM factors in \eqref{eq:factorized-welded-map} and the physical
projector \(\Pi_{\mathrm{phys},\Lambda}\) are also operator-norm continuous.
Hence the composed operator family
\[
\Lambda\longmapsto
\Pi_{\mathrm{phys},\Lambda}\mathsf T_\Lambda
\]
is continuous in operator norm on \(I\).
Applying this continuous family to the continuous source vector \(\tau_\Lambda\) gives continuity of
\(
X_\Lambda[\tau_\Lambda]
\)
in the Fisher norm.

For the hard projector, the same argument applies on every connected subsegment where
\(\ell_\star(\Lambda)\) is constant, because there the first factor is constant.
\end{proof}

\begin{corollary}[Branchwise continuity of the seam coefficient]
\label{cor:branchwise-seam-coeff}
Let \(I\) be a connected accepted-branch segment and assume:
\begin{enumerate}[leftmargin=2em]
\item the hypotheses of \cref{thm:dominant-line-persistence} on \(I\);
\item a detector-factorized branchwise seam realization on \(I\) with continuous unit source
  \[
  \Lambda\longmapsto \tau^{\mathrm{unit}}_{\Lambda};
  \]
\item the hypotheses of \cref{thm:smooth-taper-welded-response} for that unit source.
\end{enumerate}
Define the branchwise normalized seam coefficient by
\begin{equation}
c_\Lambda
:=
\big\langle\!\big\langle
\Pi_{\mathrm{phys},\Lambda}\mathsf T_\Lambda\!\big[\tau^{\mathrm{unit}}_{\Lambda}\big],\,
\widehat{\delta C}^{\mathrm{LoG}}_{\Lambda}
\big\rangle\!\big\rangle_{C_{\ast,\Lambda}}.
\label{eq:branchwise-seam-coeff}
\end{equation}
Then \(c_\Lambda\) is continuous on \(I\).
Consequently, for every fixed filler tag \(f\) and witness \(y\), the normalized seam scalar
\[
s_\Lambda(f,y)
=
\sigma_{\mathrm{fill}}(f)\,
c_\Lambda\,
\overline{\mathrm{obs}}(y)
\]
is continuous on \(I\).
\end{corollary}

\begin{proof}
By \cref{thm:smooth-taper-welded-response}, the response vector
\[
\Lambda\longmapsto
\Pi_{\mathrm{phys},\Lambda}\mathsf T_\Lambda\!\big[\tau^{\mathrm{unit}}_{\Lambda}\big]
\]
is continuous in the Fisher norm.
By \cref{thm:dominant-line-persistence}, the normalized dominant vector
\(
\widehat{\delta C}^{\mathrm{LoG}}_{\Lambda}
\)
is also continuous on \(I\).
Pairing these two continuous vectors in the Fisher metric gives continuity of \(c_\Lambda\).
The final statement is then immediate from
\eqref{eq:normalized-seam-factorization}.
\end{proof}

\begin{theorem}[Branchwise seam witness transport]
\label{thm:branchwise-seam-witness}
Let \(I\) be a connected accepted-branch segment.
Assume:
\begin{enumerate}[leftmargin=2em]
\item the hypotheses of \cref{cor:branchwise-seam-coeff};
\item quotient descent holds for every \(\Lambda\in I\);
\item there exists a fixed LS witness \(y\) with
  \[
  \overline{\mathrm{obs}}(y)\neq 0;
  \]
\item for some fixed filler tag \(f\), there is a branchwise family of representatives
  \[
  J^{f,y}_{\Lambda}\in[\Xi^{(1)}_{\Lambda}]
  \qquad
  (\Lambda\in I)
  \]
  such that, for every \(\Lambda\in I\), the corresponding seam source
  \(\tau^{f,y}_{\Lambda}\) is representable on the characterization layer and the 2T-side
  witness-to-descendant comparison identity
  \eqref{eq:witness-descendant-comparison-normalized} holds.
\end{enumerate}
Then:
\begin{enumerate}[leftmargin=2em]
\item the branchwise dominant-response scalar of the chosen descendant representatives is
  \begin{equation}
  \rho^{\mathrm{dom}}_{\Lambda}\!\big(J^{f,y}_{\Lambda}\big)
  =
  \sigma_{\mathrm{fill}}(f)\,
  c_\Lambda\,
  \overline{\mathrm{obs}}(y);
  \label{eq:branchwise-seam-descendant-scalar}
  \end{equation}
\item the scalar function
  \(
  \Lambda\mapsto \rho^{\mathrm{dom}}_{\Lambda}(J^{f,y}_{\Lambda})
  \)
  is continuous on \(I\);
\item on every connected subsegment \(I'\subset I\) on which \(c_\Lambda\neq 0\), the sign
  \[
  \operatorname{sgn}\!\big(
  \rho^{\mathrm{dom}}_{\Lambda}(J^{f,y}_{\Lambda})
  \big)
  \]
  is constant;
\item because quotient descent holds, the same constant sign is the transported orientation sign of
  the descendant affine class on \(I'\).
\end{enumerate}
\end{theorem}

\begin{proof}
For each \(\Lambda\in I\), representability of \(\tau^{f,y}_{\Lambda}\) and
\cref{prop:odd-source-pairing} give
\[
\rho^{\mathrm{dom}}_{\Lambda}\!\big(F_{\tau^{f,y}_{\Lambda},\Lambda}\big)
=
\big\langle\!\big\langle
\Pi_{\mathrm{phys},\Lambda}\mathsf T_\Lambda\!\big[\tau^{f,y}_{\Lambda}\big],\,
\widehat{\delta C}^{\mathrm{LoG}}_{\Lambda}
\big\rangle\!\big\rangle_{C_{\ast,\Lambda}}
=
s_\Lambda(f,y).
\]
Applying the branchwise comparison identity
\eqref{eq:witness-descendant-comparison-normalized} then yields
\[
\rho^{\mathrm{dom}}_{\Lambda}\!\big(J^{f,y}_{\Lambda}\big)
=
s_\Lambda(f,y).
\]
Now use \cref{cor:branchwise-seam-coeff} to substitute
\(
s_\Lambda(f,y)
=
\sigma_{\mathrm{fill}}(f)c_\Lambda\overline{\mathrm{obs}}(y)
\),
which proves item (1).
Item (2) follows because \(c_\Lambda\) is continuous and
\(\overline{\mathrm{obs}}(y)\) is a fixed nonzero scalar.
On any connected subsegment where \(c_\Lambda\neq 0\), the right-hand side of
\eqref{eq:branchwise-seam-descendant-scalar} is a continuous nonvanishing real-valued function, so
its sign is constant.
This proves item (3).

Finally, quotient descent implies that \(\rho^{\mathrm{dom}}_{\Lambda}\) is constant on
\([\Xi^{(1)}_{\Lambda}]\), so the sign of one representative is the sign of the whole descendant
class.
This gives item (4).
\end{proof}

\begin{corollary}[Seam package discharge of the minimal scalar transport input]
\label{cor:branchwise-seam-minimal-input}
Let \(I\) be a connected accepted-branch segment and assume the hypotheses of
\cref{thm:branchwise-seam-witness}.
Then the branchwise family
\[
J^{f,y}_{\Lambda}\in[\Xi^{(1)}_{\Lambda}]
\qquad
(\Lambda\in I)
\]
satisfies
\[
\Lambda\longmapsto
\rho^{\mathrm{dom}}_{\Lambda}\!\big(J^{f,y}_{\Lambda}\big)
=
\sigma_{\mathrm{fill}}(f)\,
c_\Lambda\,
\overline{\mathrm{obs}}(y)
\]
continuously on \(I\).
On every connected subsegment \(I'\subset I\) such that
\[
c_\Lambda\neq 0
\qquad
\forall \Lambda\in I',
\]
assumption \textup{(3)} of \cref{thm:active-response-transport-minimal} holds in its minimal
scalar form with
\[
J^{\mathrm{desc}}_{\Lambda}:=J^{f,y}_{\Lambda}
\qquad
(\Lambda\in I').
\]
If, in addition, the three generator-wise vanishing statements of
\cref{prop:dominant-quotient-generators} hold for every \(\Lambda\in I'\), then
\cref{thm:active-response-transport-minimal} applies on \(I'\).
\end{corollary}

\begin{proof}
Continuity of the displayed scalar function is exactly item (2) of
\cref{thm:branchwise-seam-witness}, together with the explicit formula of item (1).
On \(I'\), the hypotheses
\(
c_\Lambda\neq 0
\)
and
\(
\overline{\mathrm{obs}}(y)\neq 0
\)
imply
\[
\rho^{\mathrm{dom}}_{\Lambda}\!\big(J^{f,y}_{\Lambda}\big)\neq 0
\qquad
\forall \Lambda\in I'.
\]
Together with the continuity statement already proved above, this is exactly assumption
\textup{(3)} of \cref{thm:active-response-transport-minimal} in its minimal scalar form.
The hypotheses of \cref{thm:branchwise-seam-witness} already include those of
\cref{cor:branchwise-seam-coeff}, and hence the dominant-line persistence input of
\cref{thm:active-response-transport-minimal}.
If the three generator-wise vanishing statements also hold on \(I'\), then every hypothesis of
\cref{thm:active-response-transport-minimal} is in force there.
\end{proof}

\begin{theorem}[Finite-cutoff Active Response Transport Theorem]
\label{thm:active-response-transport-minimal}
Let \(I\) be a connected accepted-branch segment.
For each \(\Lambda\in I\), let
\[
[\Xi^{(1)}_{\Lambda}]
:=
J_{0,\Lambda}+\mathcal A^-_{\Lambda}
\subset E^-_{\Lambda}
\]
be the parity-transported odd descendant affine class modulo the ambiguity space of
\cref{def:transport-ambiguity-subspace}.
Assume:
\begin{enumerate}[leftmargin=2em]
\item the hypotheses of \cref{thm:dominant-line-persistence} hold on \(I\);
\item for every \(\Lambda\in I\), the three vanishing statements of
  \cref{prop:dominant-quotient-generators} hold;
\item there exists a branchwise family of representatives
  \[
  J^{\mathrm{desc}}_{\Lambda}\in[\Xi^{(1)}_{\Lambda}]
  \qquad
  (\Lambda\in I)
  \]
  such that
  \[
  \rho^{\mathrm{dom}}_{\Lambda}(J^{\mathrm{desc}}_{\Lambda})
  \neq 0
  \qquad
  \forall \Lambda\in I,
  \]
  and such that the dominant-response scalar
  \[
  \Lambda\longmapsto
  \rho^{\mathrm{dom}}_{\Lambda}(J^{\mathrm{desc}}_{\Lambda})
  \]
  is continuous on \(I\).
\end{enumerate}
Then, for every \(\Lambda\in I\):
\begin{enumerate}[leftmargin=2em]
\item the compressed dominant covariance tangent is nonzero,
  \[
  \delta C^{\mathrm{LoG}}_{\Lambda}\neq 0,
  \]
  so the dominant line \(L^{\mathrm{dom}}_{\Lambda}\) of
  \cref{def:dominant-response-functional} is canonically defined;
\item the dominant-response coefficient \(\rho^{\mathrm{dom}}_{\Lambda}\) annihilates the odd
  ambiguity space,
  \[
  \rho^{\mathrm{dom}}_{\Lambda}(A)=0
  \qquad
  \forall A\in\mathcal A^-_{\Lambda};
  \]
\item the transported descendant affine class carries a unique nonzero dominant-response value
  \[
  \rho^{\mathrm{dom}}_{\Lambda}(J)
  =
  \overline{\rho}^{\mathrm{dom}}_{\Lambda}\big([\![\Xi^{(1)}_{\Lambda}]\!]\big)
  \neq 0
  \qquad
  \forall J\in[\Xi^{(1)}_{\Lambda}],
  \]
  and hence a well-defined sign
  \[
  \varepsilon_{\Lambda}
  :=
  \operatorname{sgn}\!\Big(
  \overline{\rho}^{\mathrm{dom}}_{\Lambda}\big([\![\Xi^{(1)}_{\Lambda}]\!]\big)
  \Big)
  \in\{\pm1\};
  \]
\item the normalized orientation functional
  \begin{equation}
  \chi_{\Lambda}
  :=
  \nu_{\Lambda}^{-1}\rho^{\mathrm{dom}}_{\Lambda}
  \label{eq:chi-from-rho-dom}
  \end{equation}
  is well defined, annihilates \(\mathcal A^-_{\Lambda}\), and satisfies
  \[
  \chi_{\Lambda}(J)=\varepsilon_{\Lambda}
  \qquad
  \forall J\in[\Xi^{(1)}_{\Lambda}];
  \]
\item along \(I\), the sign \(\varepsilon_{\Lambda}\) is locally constant once an initial
  orientation is fixed.
\end{enumerate}
\end{theorem}

\begin{proof}
Item (1) is exactly the conclusion of \cref{thm:dominant-line-persistence}.

By assumption, the three generator-wise vanishing statements of
\cref{prop:dominant-quotient-generators} hold for every \(\Lambda\in I\).
Therefore \cref{prop:dominant-quotient-generators,thm:dominant-quotient-descent} give item (2)
and the existence of the descended functional
\(
\overline{\rho}^{\mathrm{dom}}_{\Lambda}
\)
at each cutoff.

Applying \cref{prop:descendant-overlap-equivalent,thm:descendant-overlap} pointwise to the
nonvanishing scalar assumption
\(
\rho^{\mathrm{dom}}_{\Lambda}(J^{\mathrm{desc}}_{\Lambda})\neq 0
\)
now yields
item (3): the entire descendant affine class has one common nonzero dominant-response value, and
its sign \(\varepsilon_{\Lambda}\) is well defined.

Since \(\nu_{\Lambda}>0\) by item (3), the normalized functional \(\chi_{\Lambda}\) of
\eqref{eq:chi-from-rho-dom} is well defined.
Item (2) implies that \(\chi_{\Lambda}\) annihilates \(\mathcal A^-_{\Lambda}\), and item (3)
gives
\[
\chi_{\Lambda}(J)
=
\frac{\rho^{\mathrm{dom}}_{\Lambda}(J)}{
\left|
\overline{\rho}^{\mathrm{dom}}_{\Lambda}\big([\![\Xi^{(1)}_{\Lambda}]\!]\big)
\right|
}
=
\varepsilon_{\Lambda}
\qquad
\forall J\in[\Xi^{(1)}_{\Lambda}],
\]
proving item (4).

Finally, assumption (3) provides a continuous nonvanishing real-valued branchwise descendant scalar,
while items (1)--(3) above identify that scalar with the transported orientation sign of the full
descendant affine class.
Since
\(
\Lambda\mapsto \rho^{\mathrm{dom}}_{\Lambda}(J^{\mathrm{desc}}_{\Lambda})
\)
is continuous and nonzero on the connected set \(I\), its sign is constant, hence in particular
locally constant.
Using item (3), this is exactly the transported orientation sign of the full descendant affine
class, proving item (5).
\end{proof}

\begin{proposition}[Ambiguity-triviality on the compressed welded image]
\label{prop:reduced-welded-ambiguity-trivial}
Assume the hypotheses of \cref{thm:active-response-transport-minimal} at an accepted tick \(k\).
Define the reduced welded image of an odd generator \(F\in E^-_{\Lambda_k}\) by
\[
\mathsf W^{\mathrm{red}}_{\Lambda_k}(F)
:=
\rho^{\mathrm{dom}}_{\Lambda_k}(F)\,P_{-,\Lambda_k}^{\mathrm{LoG}},
\]
where \(P_{-,\Lambda_k}^{\mathrm{LoG}}\) is the canonical compressed odd projector of
\cref{cor:compressed-export-selection}.
Then
\[
\mathsf W^{\mathrm{red}}_{\Lambda_k}(F+A)
=
\mathsf W^{\mathrm{red}}_{\Lambda_k}(F)
\qquad
\forall F\in E^-_{\Lambda_k},\ \forall A\in\mathcal A^-_{\Lambda_k}.
\]
Equivalently, \(\mathsf W^{\mathrm{red}}_{\Lambda_k}\) factors through the ambiguity quotient
\(E^-_{\Lambda_k}/\mathcal A^-_{\Lambda_k}\).
In particular, if \(J_1,J_2\in[\Xi^{(1)}_{\Lambda_k}]\), then
\[
\mathsf W^{\mathrm{red}}_{\Lambda_k}(J_1)
=
\mathsf W^{\mathrm{red}}_{\Lambda_k}(J_2).
\]
\end{proposition}

\begin{proof}
Item (2) of \cref{thm:active-response-transport-minimal} gives
\[
\rho^{\mathrm{dom}}_{\Lambda_k}(A)=0
\qquad
\forall A\in\mathcal A^-_{\Lambda_k}.
\]
Therefore, for every \(F\in E^-_{\Lambda_k}\) and \(A\in\mathcal A^-_{\Lambda_k}\),
\[
\mathsf W^{\mathrm{red}}_{\Lambda_k}(F+A)
=
\rho^{\mathrm{dom}}_{\Lambda_k}(F+A)\,P_{-,\Lambda_k}^{\mathrm{LoG}}
=
\big(\rho^{\mathrm{dom}}_{\Lambda_k}(F)+\rho^{\mathrm{dom}}_{\Lambda_k}(A)\big)
P_{-,\Lambda_k}^{\mathrm{LoG}}
=
\rho^{\mathrm{dom}}_{\Lambda_k}(F)\,P_{-,\Lambda_k}^{\mathrm{LoG}}.
\]
Since \cref{cor:compressed-export-selection} fixes \(P_{-,\Lambda_k}^{\mathrm{LoG}}\) canonically on
the compressed, co-centered, band-locked welded channel, the right-hand side depends only on the
ambiguity class of \(F\).
If \(J_2-J_1\in\mathcal A^-_{\Lambda_k}\), apply the same identity with \(F=J_1\) and
\(A=J_2-J_1\).
\end{proof}

\begin{proposition}[Construction of an admissible calibration for the minimal transported interface]
\label{prop:minimal-transport-calibration}
Fix an accepted tick \(k\). Suppose a weak transported-interface package is given:
\[
\mathcal A^-_{\Lambda_k}\subset E^-_{\Lambda_k},
\qquad
[\Xi^{(1)}_{\Lambda_k}]=J_{0,\Lambda_k}+\mathcal A^-_{\Lambda_k},
\]
together with a continuous nonzero linear functional
\[
\chi_{\Lambda_k}:E^-_{\Lambda_k}\to\mathbb R
\]
and a sign
\(
\varepsilon_{\Lambda_k}\in\{\pm1\}
\)
such that
\[
\mathcal A^-_{\Lambda_k}\subset\ker\chi_{\Lambda_k},
\qquad
\chi_{\Lambda_k}(J)=\varepsilon_{\Lambda_k}
\quad
\forall J\in[\Xi^{(1)}_{\Lambda_k}].
\]
Assume in addition that \(\ker\chi_{\Lambda_k}\) admits a continuous symmetric coercive bilinear form
whose induced norm Hilbertizes it.
Choose any vector \(J^\sharp_{\Lambda_k}\in E^-_{\Lambda_k}\) satisfying
\[
\chi_{\Lambda_k}(J^\sharp_{\Lambda_k})=1,
\]
any scale \(\beta_{\Lambda_k}>0\), and any continuous symmetric coercive bilinear form
\[
\langle\cdot,\cdot\rangle^{\mathrm{aux},\chi}_{\Lambda_k}
:
\ker\chi_{\Lambda_k}\times\ker\chi_{\Lambda_k}\to\mathbb R
\]
whose induced norm Hilbertizes \(\ker\chi_{\Lambda_k}\).
Define
\begin{equation}
\mathsf B_{\Lambda_k}(F,G)
:=
\beta_{\Lambda_k}\,\chi_{\Lambda_k}(F)\chi_{\Lambda_k}(G)
+
\left\langle
F-\chi_{\Lambda_k}(F)J^\sharp_{\Lambda_k},
\,
G-\chi_{\Lambda_k}(G)J^\sharp_{\Lambda_k}
\right\rangle^{\mathrm{aux},\chi}_{\Lambda_k}.
\label{eq:B-minimal-transport}
\end{equation}
Then \(\mathsf B_{\Lambda_k}\) is a continuous symmetric positive definite bilinear form on
\(E^-_{\Lambda_k}\) whose induced norm Hilbertizes \(E^-_{\Lambda_k}\), and
\[
\mathsf B_{\Lambda_k}(J^\sharp_{\Lambda_k},K)=0
\qquad
\forall K\in\ker\chi_{\Lambda_k}.
\]
Hence \(\mathsf B_{\Lambda_k}(J^\sharp_{\Lambda_k},K)=0\) in particular for every
\(K\in\mathcal A^-_{\Lambda_k}\), and
\[
\mathsf B_{\Lambda_k}(J^\sharp_{\Lambda_k},J^\sharp_{\Lambda_k})
=
\beta_{\Lambda_k}.
\]
Moreover, for any representative \(J^\Xi_{\Lambda_k}\in[\Xi^{(1)}_{\Lambda_k}]\), the choice
\[
J^\sharp_{\Lambda_k}:=\varepsilon_{\Lambda_k}J^\Xi_{\Lambda_k}
\]
is admissible.
\end{proposition}

\begin{proof}
Because \(\chi_{\Lambda_k}(J^\sharp_{\Lambda_k})=1\), the linear map
\[
P^\chi_{\Lambda_k}(F)
:=
F-\chi_{\Lambda_k}(F)J^\sharp_{\Lambda_k}
\]
takes values in \(\ker\chi_{\Lambda_k}\):
\[
\chi_{\Lambda_k}\!\big(P^\chi_{\Lambda_k}(F)\big)
=
\chi_{\Lambda_k}(F)-\chi_{\Lambda_k}(F)\chi_{\Lambda_k}(J^\sharp_{\Lambda_k})
=
0.
\]
Thus every \(F\in E^-_{\Lambda_k}\) has the decomposition
\[
F=\chi_{\Lambda_k}(F)J^\sharp_{\Lambda_k}+P^\chi_{\Lambda_k}(F),
\qquad
P^\chi_{\Lambda_k}(F)\in\ker\chi_{\Lambda_k}.
\]
This decomposition is unique: if \(F=aJ^\sharp_{\Lambda_k}+K\) with \(K\in\ker\chi_{\Lambda_k}\),
applying \(\chi_{\Lambda_k}\) gives \(a=\chi_{\Lambda_k}(F)\).

Equation \eqref{eq:B-minimal-transport} is therefore well defined, continuous, and symmetric.
If \(\mathsf B_{\Lambda_k}(F,F)=0\), then
\[
\beta_{\Lambda_k}\chi_{\Lambda_k}(F)^2=0
\]
and
\[
\left\langle
P^\chi_{\Lambda_k}(F),P^\chi_{\Lambda_k}(F)
\right\rangle^{\mathrm{aux},\chi}_{\Lambda_k}
=
0.
\]
Since \(\beta_{\Lambda_k}>0\), we get \(\chi_{\Lambda_k}(F)=0\), hence
\(P^\chi_{\Lambda_k}(F)=F\). Coercivity on \(\ker\chi_{\Lambda_k}\) then implies \(F=0\), proving
positive definiteness.

The map
\[
F\longmapsto
\big(
\chi_{\Lambda_k}(F),P^\chi_{\Lambda_k}(F)
\big)
\]
is a continuous linear isomorphism from \(E^-_{\Lambda_k}\) onto
\(
\mathbb R\oplus\ker\chi_{\Lambda_k}
\),
with inverse
\[
(a,K)\longmapsto aJ^\sharp_{\Lambda_k}+K.
\]
Under this isomorphism, \(\mathsf B_{\Lambda_k}\) becomes the product Hilbert form
\[
(a,K),(b,L)\longmapsto
\beta_{\Lambda_k}ab+\langle K,L\rangle^{\mathrm{aux},\chi}_{\Lambda_k},
\]
so its induced norm Hilbertizes \(E^-_{\Lambda_k}\).

Finally, if \(K\in\ker\chi_{\Lambda_k}\), then \(P^\chi_{\Lambda_k}(J^\sharp_{\Lambda_k})=0\) and
\(\chi_{\Lambda_k}(K)=0\), hence
\[
\mathsf B_{\Lambda_k}(J^\sharp_{\Lambda_k},K)=0.
\]
Setting \(K=J^\sharp_{\Lambda_k}\) gives
\(
\mathsf B_{\Lambda_k}(J^\sharp_{\Lambda_k},J^\sharp_{\Lambda_k})=\beta_{\Lambda_k}
\).
If \(J^\Xi_{\Lambda_k}\in[\Xi^{(1)}_{\Lambda_k}]\), then
\(
\chi_{\Lambda_k}(J^\Xi_{\Lambda_k})=\varepsilon_{\Lambda_k}
\),
so
\(
J^\sharp_{\Lambda_k}:=\varepsilon_{\Lambda_k}J^\Xi_{\Lambda_k}
\)
satisfies
\(
\chi_{\Lambda_k}(J^\sharp_{\Lambda_k})=1
\).
\end{proof}

\begin{remark}[Calibration existence versus stiffness matching]
\label{rem:minimal-transport-calibration-scope}
\Cref{prop:minimal-transport-calibration} is a construction theorem, not an unconditional existence
theorem: it extends a previously chosen Hilbertizing form on \(\ker\chi_{\Lambda_k}\) to an
admissible bilinear form on the whole weak transported interface.
It does \emph{not} identify the scale
\(\beta_{\Lambda_k}\) with the split/leak stiffness \(\widetilde\alpha(\Lambda_k)\).
That physical identification belongs to a later strong or physical layer.
\end{remark}

\begin{theorem}[Rank-one descendant-sector upgrade]
\label{thm:active-response-transport-strong}
Assume \cref{thm:active-response-transport-minimal}, and fix an accepted tick \(k\) on the branch.
Suppose in addition that on the transported
descendant sector the projected response map
\[
F\longmapsto
\Pi^{\mathrm{dom}}_{\Lambda_k}\Pi_{\mathrm{phys},\Lambda_k}\mathcal R_{\Lambda_k}(F),
\qquad
\Pi^{\mathrm{dom}}_{\Lambda_k}
:=
\widehat{\delta C}^{\mathrm{LoG}}_{\Lambda_k}\otimes
\widehat{\delta C}^{\mathrm{LoG}}_{\Lambda_k},
\]
has kernel exactly \(\mathcal A^-_{\Lambda_k}\).
Then the stronger conclusions follow:
\begin{enumerate}[leftmargin=2em]
\item the descendant quotient is rank one,
  \[
  E^-_{\Lambda_k}/\mathcal A^-_{\Lambda_k}\cong \mathbb R;
  \]
\item the orientation kernel is exactly the ambiguity space,
  \[
  \ker\chi_{\Lambda_k}
  =
  \ker\rho^{\mathrm{dom}}_{\Lambda_k}
  =
  \mathcal A^-_{\Lambda_k},
  \]
  so the descendant sector satisfies the rank-one kernel-identification package with quotient
  generated by any vector \(J^{\mathrm{act}}_{\Lambda_k}\) satisfying
  \(\chi_{\Lambda_k}(J^{\mathrm{act}}_{\Lambda_k})=1\);
  \item after choosing any such \(J^{\mathrm{act}}_{\Lambda_k}\), any positive scale
  \(\beta_{\Lambda_k}\), and any auxiliary coercive Hilbertizing form on
  \(\ker\chi_{\Lambda_k}\), \cref{prop:minimal-transport-calibration} yields an admissible
  calibration form.
  Recovering the stiffness-matched formula of \cref{prop:char-calibration-extension} requires
  additional physical identifications, namely
  \(\chi_{\Lambda_k}=\rho_{\Lambda_k}\), \(\ker\chi_{\Lambda_k}=\mathcal T^{\mathrm{op}}_{\Lambda_k}\),
  agreement of the chosen normalized generator with the split/leak active generator, and
  \(\beta_{\Lambda_k}=\widetilde\alpha(\Lambda_k)\).
\end{enumerate}
\end{theorem}

\begin{proof}
By \cref{thm:active-response-transport-minimal}, the dominant line is defined at \(\Lambda_k\), the
scalar \(\nu_{\Lambda_k}\) is positive, and
\[
\chi_{\Lambda_k}
=
\nu_{\Lambda_k}^{-1}\rho^{\mathrm{dom}}_{\Lambda_k}.
\]
Hence
\(
\ker\chi_{\Lambda_k}=\ker\rho^{\mathrm{dom}}_{\Lambda_k}
\).

For every \(F\in E^-_{\Lambda_k}\), the definition of \(\rho^{\mathrm{dom}}_{\Lambda_k}\) gives
\[
\Pi^{\mathrm{dom}}_{\Lambda_k}\Pi_{\mathrm{phys},\Lambda_k}\mathcal R_{\Lambda_k}(F)
=
\rho^{\mathrm{dom}}_{\Lambda_k}(F)\,
\widehat{\delta C}^{\mathrm{LoG}}_{\Lambda_k}.
\]
Therefore the kernel of the projected response map is exactly
\(
\ker\rho^{\mathrm{dom}}_{\Lambda_k}
\).
By the additional hypothesis, that kernel is \(\mathcal A^-_{\Lambda_k}\), proving item (2).

By \cref{thm:active-response-transport-minimal}, there exists
\(
J^{\mathrm{desc}}_{\Lambda_k}\in[\Xi^{(1)}_{\Lambda_k}]
\)
with
\(
\rho^{\mathrm{dom}}_{\Lambda_k}(J^{\mathrm{desc}}_{\Lambda_k})\neq 0
\),
hence the projected response map is nonzero.
Its image is automatically contained in the one-dimensional line
\(
L^{\mathrm{dom}}_{\Lambda_k}
\),
so the first isomorphism theorem gives
\[
E^-_{\Lambda_k}/\mathcal A^-_{\Lambda_k}
\cong
\operatorname{im}\!\big(
\Pi^{\mathrm{dom}}_{\Lambda_k}\Pi_{\mathrm{phys},\Lambda_k}\mathcal R_{\Lambda_k}
\big)
\cong
\mathbb R,
\]
which is item (1).
Any vector \(J^{\mathrm{act}}_{\Lambda_k}\) with \(\chi_{\Lambda_k}(J^{\mathrm{act}}_{\Lambda_k})=1\)
lies outside \(\ker\chi_{\Lambda_k}=\mathcal A^-_{\Lambda_k}\), so its quotient class generates this
one-dimensional quotient, completing item (2).

Item (3) is now an application of \cref{prop:minimal-transport-calibration} with the weak-interface
data supplied by \cref{thm:active-response-transport-minimal} and the chosen normalized vector
\(J^{\mathrm{act}}_{\Lambda_k}\).
The final sentence identifies when this weak calibration construction coincides with the
stiffness-matched formula of \cref{prop:char-calibration-extension}.
\end{proof}

\ifdefined\UOMTransportFragmentLoaded\else
\end{document}
\fi
%
}{%
  \IfFileExists{UOM/active_response_transport_program_note.tex}{%
    \ifdefined\UOMTransportFragmentLoaded\else
\documentclass[11pt]{article}
\usepackage[margin=1in]{geometry}
\usepackage{amsmath,amssymb,amsthm,mathtools}
\usepackage{bm}
\usepackage{microtype}
\usepackage{enumitem}
\usepackage{booktabs}
\usepackage{array}
\usepackage{xcolor}
\usepackage{xr-hyper}
\usepackage{hyperref}
\usepackage[nameinlink,noabbrev]{cleveref}

\providecommand{\Fsloc}{\mathcal F_{\rm s.loc}}
\providecommand{\Qt}{Q_t}
\externaldocument[main-]{UOM/main}
\externaldocument[uomdesc-]{UOM/uom_descendant_note}
\externaldocument{UOM/completion_sector_split_model_note}

\hypersetup{
  colorlinks=true,
  linkcolor=blue!60!black,
  citecolor=blue!60!black,
  urlcolor=blue!60!black
}

\theoremstyle{definition}
\newtheorem{definition}{Definition}[section]
\newtheorem{assumption}{Assumption}[section]
\theoremstyle{plain}
\newtheorem{theorem}[definition]{Theorem}
\newtheorem{lemma}[definition]{Lemma}
\newtheorem{proposition}[definition]{Proposition}
\newtheorem{corollary}[definition]{Corollary}
\theoremstyle{remark}
\newtheorem{remark}[definition]{Remark}
\crefname{definition}{definition}{definitions}
\Crefname{definition}{Definition}{Definitions}
\crefname{assumption}{assumption}{assumptions}
\Crefname{assumption}{Assumption}{Assumptions}
\crefname{theorem}{theorem}{theorems}
\Crefname{theorem}{Theorem}{Theorems}
\crefname{lemma}{lemma}{lemmas}
\Crefname{lemma}{Lemma}{Lemmas}
\crefname{proposition}{proposition}{propositions}
\Crefname{proposition}{Proposition}{Propositions}
\crefname{corollary}{corollary}{corollaries}
\Crefname{corollary}{Corollary}{Corollaries}
\crefname{remark}{remark}{remarks}
\Crefname{remark}{Remark}{Remarks}

\newcommand{\cH}{\mathcal H}
\newcommand{\cA}{\mathcal A}
\newcommand{\cD}{\mathcal D}
\newcommand{\cB}{\mathcal B}
\newcommand{\cC}{\mathcal C}
\newcommand{\cL}{\mathcal L}
\newcommand{\cK}{\mathcal K}
\newcommand{\cR}{\mathcal R}
\newcommand{\cG}{\mathcal G}
\newcommand{\Ztwo}{\mathbb Z_2}
\newcommand{\Sys}{\mathrm{Sys}}
\newcommand{\Time}{\mathrm{Time}}
\newcommand{\RG}{\mathrm{RG}}
\newcommand{\Tr}{\mathrm{Tr}}
\newcommand{\diag}{\mathrm{diag}}
\newcommand{\ad}{\mathrm{ad}}
\newcommand{\id}{\mathrm{id}}
\newcommand{\norm}[1]{\left\lVert #1\right\rVert}
\newcommand{\abs}[1]{\left\lvert #1\right\rvert}
\newcommand{\ip}[2]{\left\langle #1,#2\right\rangle}

\externaldocument{UOM/compressed_log_spectral_selection_note}
\externaldocument{UOM/odd_selector_note}
\externaldocument{UOM/canonical_odd_interface_note}
\externaldocument[splitlift-]{UOM/split_controlled_completion_lift_note}
\title{Finite-Cutoff Active Response Transport Theorem}
\author{}
\date{}
\begin{document}
\maketitle
\fi

\section{Finite-Cutoff Active Response Transport Theorem}
\label{sec:active-response-transport-program}

This section assembles a finite-cutoff transport theorem for the odd descendant / Hamiltonization
lane of the 2T Main assumption package Definition~\ref{main-def:uom-main-theorem-assumptions}.
Accordingly, the phrase ``accepted tick'' remains literal UOM cutoff language throughout this note:
the transport theorem is pointwise on the accepted branch at fixed \(\Lambda_k\).

\begin{remark}[Physical odd sector and branchwise ambient-space conventions]
\label{rem:transport-conventions}
Throughout this note, the superscript \({}^-\) refers to physical \(\Gamma\)-odd parity on the
characterization layer, not to cohomological oddness.
Every ambiguity direction appearing below is therefore understood after passage to the physically
odd characterization sector \(E^-_{\Lambda_k}\).

All branchwise continuity statements are made after the standard accepted-branch identifications of
the relevant finite-cutoff source spaces, covariance-tangent spaces, and physical subspaces with
fixed ambient finite-cutoff spaces.
Operator-norm and Fisher-norm continuity are always understood with respect to those identifications.
\end{remark}

\begin{definition}[Odd ambiguity subspace on the characterization layer]
\label{def:transport-ambiguity-subspace}
For each accepted tick \(k\), let
\[
\mathcal A^-_{\Lambda_k}\subset E^-_{\Lambda_k}
\]
denote the closed linear span of the odd characterization-layer projections of the descendant
ambiguity directions
\[
\big(Q_tX+Q_uY\big)^-,\qquad \big(\partial_\mu K^\mu\big)^-,\qquad \Delta_{\rm ct}^{-},
\]
where \(K^\mu\) is boundary-admissible on the characterization layer and
\(\Delta_{\rm ct}^{-}\) is admissible in the sense of \cref{def:ambiguity-ct-char}.
Equivalently, \(\mathcal A^-_{\Lambda_k}\) is the odd ambiguity space already placed inside
\(\mathcal T^{\mathrm{op}}_{\Lambda_k}\) by \cref{cor:ambiguity-char}.
\end{definition}

\begin{definition}[Dominant compressed response line and scalar coefficient]
\label{def:dominant-response-functional}
Fix an accepted tick \(k\). Let
\[
u^{\mathrm{LoG}}_{\Lambda_k}
\]
and
\[
P_{-,\Lambda_k}^{\mathrm{LoG}}
\]
be the canonically selected compressed odd line and projector furnished by
\cref{cor:compressed-export-selection}. Define the physical covariance tangent induced by that line by
\[
\delta C^{\mathrm{LoG}}_{\Lambda_k}
:=
\Pi_{\mathrm{phys},\Lambda_k}\,\mathsf T_{\Lambda_k}[u^{\mathrm{LoG}}_{\Lambda_k}],
\]
and, whenever this tangent is nonzero, write
\[
L^{\mathrm{dom}}_{\Lambda_k}
:=
\mathbb R\,\widehat{\delta C}^{\mathrm{LoG}}_{\Lambda_k}
\]
for its Fisher-normalized line.
The raw dominant-response coefficient is
\begin{equation}
\widetilde{\rho}^{\mathrm{dom}}_{\Lambda_k}(F)
:=
\big\langle\!\big\langle
\Pi_{\mathrm{phys},\Lambda_k}\mathcal R_{\Lambda_k}(F),\,
\delta C^{\mathrm{LoG}}_{\Lambda_k}
\big\rangle\!\big\rangle_{C_{\ast,\Lambda_k}},
\qquad
F\in E^-_{\Lambda_k}.
\label{eq:rho-dom-raw-def}
\end{equation}
Whenever \(\delta C^{\mathrm{LoG}}_{\Lambda_k}\neq 0\), equivalently whenever dominant-line
realization holds at \(\Lambda_k\), define the normalized dominant-response coefficient by
\begin{equation}
\rho^{\mathrm{dom}}_{\Lambda_k}(F)
:=
\kappa_{\Lambda_k}^{-1/2}\,
\widetilde{\rho}^{\mathrm{dom}}_{\Lambda_k}(F),
\qquad
F\in E^-_{\Lambda_k}.
\label{eq:rho-dom-def}
\end{equation}
\end{definition}

\begin{proposition}[Well-posedness and continuity of the dominant-response coefficients]
\label{prop:rho-dom-continuity}
Fix an accepted tick \(k\).
Then the raw dominant-response coefficient
\[
\widetilde{\rho}^{\mathrm{dom}}_{\Lambda_k}:E^-_{\Lambda_k}\to\mathbb R
\]
is continuous and linear.
If dominant-line realization holds at \(\Lambda_k\), then the normalized dominant-response
coefficient
\[
\rho^{\mathrm{dom}}_{\Lambda_k}:E^-_{\Lambda_k}\to\mathbb R
\]
is also continuous and linear.
\end{proposition}

\begin{proof}
By \cref{prop:intrinsic-cov-response}, the intrinsic response map
\(
\mathcal R_{\Lambda_k}:E^-_{\Lambda_k}\to\mathcal T^{\mathrm{cov}}_{\Lambda_k}
\)
is continuous and linear.
Pairing its physical projection against the fixed covariance tangent
\(
\delta C^{\mathrm{LoG}}_{\Lambda_k}
\)
in the Fisher metric therefore gives a continuous linear functional, which is exactly
\(\widetilde{\rho}^{\mathrm{dom}}_{\Lambda_k}\).
If dominant-line realization holds, then \(\kappa_{\Lambda_k}>0\) by
\cref{prop:dominant-line-equivalent}, so \(\rho^{\mathrm{dom}}_{\Lambda_k}\) is obtained from
\(\widetilde{\rho}^{\mathrm{dom}}_{\Lambda_k}\) by multiplication by the fixed scalar
\(\kappa_{\Lambda_k}^{-1/2}\).
\end{proof}

\subsection*{Dominant-line realization}
\label{subsec:dominant-line-realization}

The compressed spectral-selection note already supplies the one-dimensional source line
\(\mathcal H^{\mathrm{LoG}}_{\Lambda_k}\).
The present transport note needs only the nonvanishing and branchwise continuity of its
\emph{physical} covariance push-forward.

\begin{definition}[Compressed-line transversality]
\label{def:compressed-line-transversality}
For an accepted tick \(k\), say that the welded response is \emph{transversely physical on the
compressed line} if
\begin{equation}
\Pi_{\mathrm{phys},\Lambda_k}\,\mathsf T_{\Lambda_k}\!\restriction_{\mathcal H^{\mathrm{LoG}}_{\Lambda_k}}
\neq 0.
\label{eq:compressed-transversality}
\end{equation}
Equivalently, the unique generator \(u^{\mathrm{LoG}}_{\Lambda_k}\) of the compressed odd line has
nonzero physical covariance response.
\end{definition}

\begin{definition}[Dominant-response norm]
\label{def:dominant-response-norm}
Whenever \(\delta C^{\mathrm{LoG}}_{\Lambda_k}\) is defined, set
\begin{equation}
\kappa_{\Lambda_k}
:=
\big\langle\!\big\langle
\delta C^{\mathrm{LoG}}_{\Lambda_k},\,
\delta C^{\mathrm{LoG}}_{\Lambda_k}
\big\rangle\!\big\rangle_{C_{\ast,\Lambda_k}}.
\label{eq:kappa-dom-def}
\end{equation}
This is the Fisher squared norm of the physical covariance tangent induced by the compressed welded
LoG source line.
\end{definition}

\begin{proposition}[Equivalent formulations of dominant-line realization]
\label{prop:dominant-line-equivalent}
For a fixed accepted tick \(k\), the following are equivalent.
\begin{enumerate}[leftmargin=2em]
\item the compressed welded LoG line is transversely physical in the sense of
  \cref{def:compressed-line-transversality};
\item the distinguished physical covariance tangent is nonzero,
  \[
  \delta C^{\mathrm{LoG}}_{\Lambda_k}\neq 0;
  \]
\item the dominant-response norm is strictly positive,
  \[
  \kappa_{\Lambda_k}>0;
  \]
\item there exists a physical covariance tangent
  \(
  B_{\Lambda_k}\in \mathcal T^{\mathrm{phys}}_{\Lambda_k}
  \)
  such that
  \[
  \big\langle\!\big\langle
  \delta C^{\mathrm{LoG}}_{\Lambda_k},\,B_{\Lambda_k}
  \big\rangle\!\big\rangle_{C_{\ast,\Lambda_k}}
  \neq 0;
  \]
\item in any Gaussian completion realization of
  \cref{cor:compressed-export-selection-completion}, the completion-sector covariance tangent
  representing \(\delta C^{\mathrm{LoG}}_{\Lambda_k}\) is nonzero.
\end{enumerate}
\end{proposition}

\begin{proof}
The equivalence of (1) and (2) is immediate from
\eqref{eq:compressed-transversality} and the definition
\[
\delta C^{\mathrm{LoG}}_{\Lambda_k}
=
\Pi_{\mathrm{phys},\Lambda_k}\mathsf T_{\Lambda_k}[u^{\mathrm{LoG}}_{\Lambda_k}],
\]
because \(\mathcal H^{\mathrm{LoG}}_{\Lambda_k}\) is the one-dimensional line generated by
\(u^{\mathrm{LoG}}_{\Lambda_k}\).

The equivalence of (2) and (3) follows from positive definiteness of the Fisher metric on the
covariance tangent space: a vector has positive Fisher norm if and only if it is nonzero.

Statement (3) implies (4) by taking
\(
B_{\Lambda_k}:=\delta C^{\mathrm{LoG}}_{\Lambda_k}
\),
which belongs to the physical tangent space by construction.
Conversely, if (4) holds then \(\delta C^{\mathrm{LoG}}_{\Lambda_k}\) cannot vanish, so (2) holds.

Finally, \cref{prop:intrinsic-cov-response,cor:compressed-export-selection-completion} identify the
abstract covariance tangent with any Gaussian completion representative of the same tangent object.
Therefore the abstract tangent is nonzero exactly when its completion realization is nonzero, giving
the equivalence of (2) and (5).
\end{proof}

\begin{theorem}[Branchwise realization and persistence of the dominant line]
\label{thm:dominant-line-persistence}
Let \(I\) be a connected accepted-branch segment in cutoff space.
Assume that:
\begin{enumerate}[leftmargin=2em]
\item the compressed projector \(P_{-,\Lambda}^{\mathrm{LoG}}\) varies continuously on \(I\);
\item the physical welded response operator
  \[
  \Lambda\longmapsto \Pi_{\mathrm{phys},\Lambda}\mathsf T_{\Lambda}
  \]
  is continuous on \(I\) in the operator topology induced by the Fisher norm on covariance
  tangents;
\item there exists a constant \(c_I>0\) such that
  \[
  \kappa_{\Lambda}\ge c_I
  \qquad
  \forall \Lambda\in I.
  \]
\end{enumerate}
Then:
\begin{enumerate}[leftmargin=2em]
\item \(\delta C^{\mathrm{LoG}}_{\Lambda}\neq 0\) for every \(\Lambda\in I\);
\item the normalized dominant vector
  \[
  \widehat{\delta C}^{\mathrm{LoG}}_{\Lambda}
  :=
  \kappa_{\Lambda}^{-1/2}\,\delta C^{\mathrm{LoG}}_{\Lambda}
  \]
  is continuous on \(I\);
\item the dominant physical line
  \[
  L^{\mathrm{dom}}_{\Lambda}
  =
  \mathbb R\,\widehat{\delta C}^{\mathrm{LoG}}_{\Lambda}
  \]
  is canonically transported along \(I\);
\item the Fisher-orthogonal rank-one projector
  \[
  \Pi^{\mathrm{dom}}_{\Lambda}
  :=
  \widehat{\delta C}^{\mathrm{LoG}}_{\Lambda}
  \otimes
  \widehat{\delta C}^{\mathrm{LoG}}_{\Lambda}
  \]
  depends continuously on \(\Lambda\in I\).
\end{enumerate}
\end{theorem}

\begin{proof}
By \cref{cor:compressed-export-selection}, the compressed projector
\(
P_{-,\Lambda}^{\mathrm{LoG}}
\)
and hence the normalized source vector \(u^{\mathrm{LoG}}_{\Lambda}\) vary continuously on the branch.
By the assumed continuity of \(\Pi_{\mathrm{phys},\Lambda}\mathsf T_{\Lambda}\), the map
\[
\Lambda\longmapsto
\delta C^{\mathrm{LoG}}_{\Lambda}
=
\Pi_{\mathrm{phys},\Lambda}\mathsf T_{\Lambda}[u^{\mathrm{LoG}}_{\Lambda}]
\]
is continuous in the Fisher norm.

The lower bound \(\kappa_{\Lambda}\ge c_I>0\) implies
\(
\delta C^{\mathrm{LoG}}_{\Lambda}\neq 0
\)
for every \(\Lambda\in I\), proving item (1).
Since the normalization map \(v\mapsto v/\|v\|_{C_\ast}\) is continuous on the complement of the
zero section, continuity of \(\delta C^{\mathrm{LoG}}_{\Lambda}\) and the positive lower bound on its
norm imply continuity of \(\widehat{\delta C}^{\mathrm{LoG}}_{\Lambda}\), proving item (2).
Items (3) and (4) are immediate consequences of item (2).
\end{proof}

\begin{corollary}[Local persistence from one-tick transversality]
\label{cor:dominant-line-local-persistence}
Assume the continuity hypotheses of \cref{thm:dominant-line-persistence}.
If the compressed line is transversely physical at one accepted tick \(\Lambda_{k_0}\), then there
exists a neighborhood \(I_{k_0}\) of \(\Lambda_{k_0}\) on the same accepted branch such that
\[
\delta C^{\mathrm{LoG}}_{\Lambda}\neq 0
\qquad
\forall \Lambda\in I_{k_0}.
\]
Hence the dominant physical line is locally well defined and continuously transported near
\(\Lambda_{k_0}\).
\end{corollary}

\begin{proof}
By \cref{prop:dominant-line-equivalent}, transversality at \(\Lambda_{k_0}\) is equivalent to
\(
\kappa_{\Lambda_{k_0}}>0
\).
Continuity of \(\kappa_{\Lambda}\) then gives a neighborhood \(I_{k_0}\) and a constant
\(c_{k_0}>0\) such that
\(
\kappa_{\Lambda}\ge c_{k_0}
\)
for all \(\Lambda\in I_{k_0}\).
Apply \cref{thm:dominant-line-persistence}.
\end{proof}

\begin{remark}[Practical witness criteria]
\label{rem:dominant-line-witnesses}
For applications, any one of the following suffices to discharge dominant-line realization at a fixed
accepted tick:
\begin{enumerate}[leftmargin=2em]
\item a direct lower bound \(\kappa_{\Lambda_k}>0\);
\item a physical witness tangent \(B_{\Lambda_k}\in\mathcal T^{\mathrm{phys}}_{\Lambda_k}\) with
  nonzero Fisher pairing against \(\delta C^{\mathrm{LoG}}_{\Lambda_k}\);
\item a Gaussian completion calculation showing that the completion covariance of the realized active
  odd mode changes nontrivially under the compressed welded LoG source.
\end{enumerate}
Each witness proves the same theorem-level fact by \cref{prop:dominant-line-equivalent}.
\end{remark}

\subsection*{Descent to the ambiguity quotient}
\label{subsec:ambiguity-quotient-descent}

Once the dominant line has been realized, the next task is purely structural:
show that the dominant-response coefficient is insensitive to ambiguity moves.
This is the exact point at which the transported descendant class becomes an affine class in the
sense required by Main rather than a collection of representatives.

\begin{definition}[Dominant-response quotient and descended descendant point]
\label{def:dominant-response-quotient}
For each accepted tick \(k\), define the dominant-response quotient by
\[
\mathfrak D^-_{\Lambda_k}
:=
E^-_{\Lambda_k}/\mathcal A^-_{\Lambda_k},
\]
with quotient map
\[
q_{\Lambda_k}:E^-_{\Lambda_k}\to \mathfrak D^-_{\Lambda_k}.
\]
If
\[
[\Xi^{(1)}_{\Lambda_k}]
=
J_{0,\Lambda_k}+\mathcal A^-_{\Lambda_k},
\]
then its image in the quotient is a single point,
\[
[\![\Xi^{(1)}_{\Lambda_k}]\!]
:=
q_{\Lambda_k}\!\big(J_{0,\Lambda_k}\big)
\in \mathfrak D^-_{\Lambda_k},
\]
independent of the chosen representative \(J_{0,\Lambda_k}\).
\end{definition}

\begin{proposition}[Universal descent criterion for the dominant-response coefficient]
\label{prop:dominant-quotient-universal}
Assume dominant-line realization at an accepted tick \(k\), so that
\(\rho^{\mathrm{dom}}_{\Lambda_k}\) is defined by \eqref{eq:rho-dom-def}.
Then the following are equivalent.
\begin{enumerate}[leftmargin=2em]
\item there exists a unique continuous linear functional
  \[
  \overline{\rho}^{\mathrm{dom}}_{\Lambda_k}:\mathfrak D^-_{\Lambda_k}\to\mathbb R
  \]
  such that
  \[
  \rho^{\mathrm{dom}}_{\Lambda_k}
  =
  \overline{\rho}^{\mathrm{dom}}_{\Lambda_k}\circ q_{\Lambda_k};
  \]
\item the ambiguity subspace is contained in the kernel:
  \[
  \mathcal A^-_{\Lambda_k}\subset \ker\rho^{\mathrm{dom}}_{\Lambda_k};
  \]
\item the dominant-response coefficient is constant on each ambiguity coset:
  \[
  \rho^{\mathrm{dom}}_{\Lambda_k}(F+A)
  =
  \rho^{\mathrm{dom}}_{\Lambda_k}(F)
  \qquad
  \forall F\in E^-_{\Lambda_k},\ \forall A\in\mathcal A^-_{\Lambda_k}.
  \]
\end{enumerate}
\end{proposition}

\begin{proof}
The equivalence of (2) and (3) is immediate from linearity:
\[
\rho^{\mathrm{dom}}_{\Lambda_k}(F+A)-\rho^{\mathrm{dom}}_{\Lambda_k}(F)
=
\rho^{\mathrm{dom}}_{\Lambda_k}(A).
\]

If (2) holds, define
\[
\overline{\rho}^{\mathrm{dom}}_{\Lambda_k}\big(q_{\Lambda_k}(F)\big)
:=
\rho^{\mathrm{dom}}_{\Lambda_k}(F).
\]
This is well defined precisely because changing representatives by an element of
\(\mathcal A^-_{\Lambda_k}\) does not change the value of \(\rho^{\mathrm{dom}}_{\Lambda_k}\).
Linearity and continuity are inherited from \(\rho^{\mathrm{dom}}_{\Lambda_k}\), and uniqueness
follows because \(q_{\Lambda_k}\) is surjective.

Conversely, if (1) holds and \(A\in\mathcal A^-_{\Lambda_k}\), then \(q_{\Lambda_k}(A)=0\), hence
\[
\rho^{\mathrm{dom}}_{\Lambda_k}(A)
=
\overline{\rho}^{\mathrm{dom}}_{\Lambda_k}\!\big(q_{\Lambda_k}(A)\big)
=
\overline{\rho}^{\mathrm{dom}}_{\Lambda_k}(0)
=
0.
\]
Thus (1) implies (2).
\end{proof}

\begin{corollary}[Constancy on the transported descendant affine class]
\label{cor:dominant-quotient-constancy}
Assume \cref{prop:dominant-quotient-universal}. Then
\[
\rho^{\mathrm{dom}}_{\Lambda_k}(J)
=
\overline{\rho}^{\mathrm{dom}}_{\Lambda_k}\big([\![\Xi^{(1)}_{\Lambda_k}]\!]\big)
\qquad
\forall J\in[\Xi^{(1)}_{\Lambda_k}].
\]
In particular, once the descended quotient functional exists, the descendant affine class carries a
single unambiguous dominant-response value.
\end{corollary}

\begin{proof}
Every \(J\in[\Xi^{(1)}_{\Lambda_k}]\) has the form
\(
J=J_{0,\Lambda_k}+A
\)
with \(A\in\mathcal A^-_{\Lambda_k}\), so
\[
q_{\Lambda_k}(J)=q_{\Lambda_k}(J_{0,\Lambda_k})=[\![\Xi^{(1)}_{\Lambda_k}]\!].
\]
Apply the factorization identity of \cref{prop:dominant-quotient-universal}.
\end{proof}

\begin{remark}[False stronger form and normalized special case]
\label{rem:dominant-quotient-false-stronger-form}
The quotient statement above is intrinsically a statement about the descended dominant scalar
\[
\overline{\rho}^{\mathrm{dom}}_{\Lambda_k}\big([\![\Xi^{(1)}_{\Lambda_k}]\!]\big),
\]
not about a bare \(\chi\)-compatibility condition.
At the weak transported-interface layer one has
\[
\chi_{\Lambda_k}
=
\nu_{\Lambda_k}^{-1}\rho^{\mathrm{dom}}_{\Lambda_k},
\]
so a representative-search condition of the form
\[
\rho^{\mathrm{dom}}_{\Lambda_k}(J)=s
\]
is equivalent to
\[
\chi_{\Lambda_k}(J)=\nu_{\Lambda_k}^{-1}s,
\]
not in general to \(\chi_{\Lambda_k}(J)=s\).
Only after an explicit normalization, for example \(\nu_{\Lambda_k}=1\), may the same scalar
comparison be rewritten in descended \(\bar\chi_{\Lambda_k}\) form without changing the value.
The finite-cutoff Wolfram audit contains exact countermodels to the stronger unscaled
\(\chi\)-level claim.
\end{remark}

\begin{proposition}[Generator-by-generator reduction]
\label{prop:dominant-quotient-generators}
Assume dominant-line realization at an accepted tick \(k\).
Suppose the dominant-response coefficient vanishes on the three generating ambiguity families:
\begin{enumerate}[leftmargin=2em]
\item for every odd-projected BRST/gauge-exact direction \(\big(Q_tX+Q_uY\big)^-\),
  \[
  \rho^{\mathrm{dom}}_{\Lambda_k}\!\big(\big(Q_tX+Q_uY\big)^-\big)=0;
  \]
\item for every odd-projected boundary-admissible divergence \(\big(\partial_\mu K^\mu\big)^-\),
  \[
  \rho^{\mathrm{dom}}_{\Lambda_k}\!\big(\big(\partial_\mu K^\mu\big)^-\big)=0;
  \]
\item for every admissible odd counterterm \(\Delta^-_{\rm ct}\),
  \[
  \rho^{\mathrm{dom}}_{\Lambda_k}(\Delta^-_{\rm ct})=0.
  \]
\end{enumerate}
Then
\[
\mathcal A^-_{\Lambda_k}\subset\ker\rho^{\mathrm{dom}}_{\Lambda_k},
\]
hence \(\rho^{\mathrm{dom}}_{\Lambda_k}\) descends uniquely to the quotient
\(\mathfrak D^-_{\Lambda_k}\).
\end{proposition}

\begin{proof}
By \cref{def:transport-ambiguity-subspace}, \(\mathcal A^-_{\Lambda_k}\) is the closed linear span of
the three families listed in the statement. Since \(\rho^{\mathrm{dom}}_{\Lambda_k}\) is linear, it
vanishes on the algebraic span of those generators. Since it is also continuous, it vanishes on the
closure of that span. Therefore every element of \(\mathcal A^-_{\Lambda_k}\) lies in the kernel.
Now apply \cref{prop:dominant-quotient-universal}.
\end{proof}

\begin{theorem}[Descent to the ambiguity quotient]
\label{thm:dominant-quotient-descent}
Fix an accepted tick \(k\) and assume dominant-line realization.
If the three vanishing statements of \cref{prop:dominant-quotient-generators} hold, then there
exists a unique descended dominant
functional
\[
\overline{\rho}^{\mathrm{dom}}_{\Lambda_k}:\mathfrak D^-_{\Lambda_k}\to\mathbb R
\]
such that
\[
\rho^{\mathrm{dom}}_{\Lambda_k}
=
\overline{\rho}^{\mathrm{dom}}_{\Lambda_k}\circ q_{\Lambda_k},
\]
and the transported descendant affine class has a single well-defined quotient value
\[
\overline{\rho}^{\mathrm{dom}}_{\Lambda_k}\big([\![\Xi^{(1)}_{\Lambda_k}]\!]\big).
\]
\end{theorem}

\begin{proof}
Apply \cref{prop:dominant-quotient-generators,cor:dominant-quotient-constancy}.
\end{proof}

\subsection*{Nonzero overlap on the descendant class}
\label{subsec:descendant-overlap}

After dominant-line realization and quotient descent, the descendant class defines a single point in
the ambiguity quotient, and the dominant-response functional descends to that quotient. This
subsection studies whether that descended descendant point is \emph{active}, i.e.\ not annihilated
by the descended dominant functional.

\begin{definition}[Dominant covariance projector]
\label{def:dominant-cov-projector}
Assume dominant-line realization at an accepted tick \(k\).
Define the Fisher-orthogonal rank-one projector onto the dominant physical line by
\begin{equation}
\Pi^{\mathrm{dom}}_{\Lambda_k}(X)
:=
\big\langle\!\big\langle
X,\,
\widehat{\delta C}^{\mathrm{LoG}}_{\Lambda_k}
\big\rangle\!\big\rangle_{C_{\ast,\Lambda_k}}\,
\widehat{\delta C}^{\mathrm{LoG}}_{\Lambda_k},
\qquad
X\in\mathcal T^{\mathrm{cov}}_{\Lambda_k}.
\label{eq:Pi-dom-def}
\end{equation}
\end{definition}

\begin{proposition}[Representative-level witness criteria]
\label{prop:descendant-overlap-witness}
Assume dominant-line realization at an accepted tick \(k\).
For any odd generator \(F\in E^-_{\Lambda_k}\), the following are equivalent.
\begin{enumerate}[leftmargin=2em]
\item the dominant-response coefficient is nonzero:
  \[
  \rho^{\mathrm{dom}}_{\Lambda_k}(F)\neq 0;
  \]
\item the dominant projection of the physical covariance response is nonzero:
  \[
  \Pi^{\mathrm{dom}}_{\Lambda_k}\Pi_{\mathrm{phys},\Lambda_k}\mathcal R_{\Lambda_k}(F)\neq 0;
  \]
\item the physical covariance response has nonzero Fisher pairing with the dominant tangent:
  \[
  \big\langle\!\big\langle
  \Pi_{\mathrm{phys},\Lambda_k}\mathcal R_{\Lambda_k}(F),\,
  \delta C^{\mathrm{LoG}}_{\Lambda_k}
  \big\rangle\!\big\rangle_{C_{\ast,\Lambda_k}}
  \neq 0;
  \]
\item in any Gaussian completion realization of
  \cref{cor:compressed-export-selection-completion}, the completion representative of
  \(\Pi_{\mathrm{phys},\Lambda_k}\mathcal R_{\Lambda_k}(F)\) has nonzero component on the realized
  active odd line.
\end{enumerate}
\end{proposition}

\begin{proof}
By \cref{def:dominant-cov-projector},
\[
\Pi^{\mathrm{dom}}_{\Lambda_k}\Pi_{\mathrm{phys},\Lambda_k}\mathcal R_{\Lambda_k}(F)
=
\rho^{\mathrm{dom}}_{\Lambda_k}(F)\,
\widehat{\delta C}^{\mathrm{LoG}}_{\Lambda_k}.
\]
Since \(\widehat{\delta C}^{\mathrm{LoG}}_{\Lambda_k}\neq 0\), this proves the equivalence of (1)
and (2).

Statement (3) is equivalent to (1) because
\[
\big\langle\!\big\langle
\Pi_{\mathrm{phys},\Lambda_k}\mathcal R_{\Lambda_k}(F),\,
\delta C^{\mathrm{LoG}}_{\Lambda_k}
\big\rangle\!\big\rangle_{C_{\ast,\Lambda_k}}
=
\|\delta C^{\mathrm{LoG}}_{\Lambda_k}\|_{C_{\ast,\Lambda_k}}\,
\rho^{\mathrm{dom}}_{\Lambda_k}(F),
\]
and the prefactor is strictly positive by dominant-line realization.

Finally, by \cref{prop:intrinsic-cov-response,cor:compressed-export-selection-completion}, the
abstract covariance tangent and its Gaussian completion representative describe the same physical
response object in different coordinates. Therefore the abstract dominant projection is nonzero if
and only if the completion representative has nonzero component on the realized active line,
proving the equivalence of (2) and (4).
\end{proof}

\begin{definition}[Descendant overlap margin]
\label{def:descendant-overlap-margin}
Assume dominant-line realization and quotient descent at an accepted tick \(k\).
Define the descendant overlap margin by
\begin{equation}
\nu_{\Lambda_k}
:=
\left|
\overline{\rho}^{\mathrm{dom}}_{\Lambda_k}\big([\![\Xi^{(1)}_{\Lambda_k}]\!]\big)
\right|.
\label{eq:descendant-overlap-margin}
\end{equation}
We say that the transported descendant class is \emph{dominant-active} if
\(
\nu_{\Lambda_k}>0
\).
\end{definition}

\begin{proposition}[Equivalent formulations of nonzero overlap on the descendant class]
\label{prop:descendant-overlap-equivalent}
Assume dominant-line realization and quotient descent at an accepted tick \(k\).
Then the following are equivalent.
\begin{enumerate}[leftmargin=2em]
\item the transported descendant sector has nonzero overlap with the canonically selected compressed
  dominant odd line;
\item there exists one representative \(J^{\mathrm{desc}}_{\Lambda_k}\in[\Xi^{(1)}_{\Lambda_k}]\)
  such that
  \[
  \Pi^{\mathrm{dom}}_{\Lambda_k}\Pi_{\mathrm{phys},\Lambda_k}
  \mathcal R_{\Lambda_k}(J^{\mathrm{desc}}_{\Lambda_k})
  \neq 0;
  \]
\item every representative \(J\in[\Xi^{(1)}_{\Lambda_k}]\) has the same nonzero dominant-response
  value:
  \[
  \rho^{\mathrm{dom}}_{\Lambda_k}(J)
  =
  \overline{\rho}^{\mathrm{dom}}_{\Lambda_k}\big([\![\Xi^{(1)}_{\Lambda_k}]\!]\big)
  \neq 0;
  \]
\item the descended descendant point is not annihilated by the descended dominant functional:
  \[
  [\![\Xi^{(1)}_{\Lambda_k}]\!]\notin \ker\overline{\rho}^{\mathrm{dom}}_{\Lambda_k};
  \]
\item the descendant overlap margin is strictly positive:
  \[
  \nu_{\Lambda_k}>0.
  \]
\end{enumerate}
\end{proposition}

\begin{proof}
The equivalence of (1) and (2) follows immediately from
\cref{prop:descendant-overlap-witness}.

If (2) holds, then
\(
\rho^{\mathrm{dom}}_{\Lambda_k}(J^{\mathrm{desc}}_{\Lambda_k})\neq 0
\)
by \cref{prop:descendant-overlap-witness}. By
\cref{cor:dominant-quotient-constancy}, every other representative of
\([\Xi^{(1)}_{\Lambda_k}]\) has the same value, proving (3).
Conversely, (3) obviously implies (1).

The equivalence of (3) and (4) is immediate from the identity
\[
\rho^{\mathrm{dom}}_{\Lambda_k}(J)
=
\overline{\rho}^{\mathrm{dom}}_{\Lambda_k}\big([\![\Xi^{(1)}_{\Lambda_k}]\!]\big)
\qquad
\forall J\in[\Xi^{(1)}_{\Lambda_k}]
\]
of \cref{cor:dominant-quotient-constancy}.

Finally, (4) and (5) are equivalent by the definition of \(\nu_{\Lambda_k}\) in
\cref{def:descendant-overlap-margin}.
\end{proof}

\begin{theorem}[Nonzero overlap on the descendant class]
\label{thm:descendant-overlap}
Fix an accepted tick \(k\), and assume dominant-line realization together with quotient descent.
If any one of the equivalent conditions in \cref{prop:descendant-overlap-equivalent} holds, then the
following conclusions obtain.
In particular, it is enough to construct a single representative
\[
J^{\mathrm{desc}}_{\Lambda_k}\in[\Xi^{(1)}_{\Lambda_k}]
\]
whose dominant projected covariance response is nonzero.

When this holds, the quotient descendant point is dominant-active, the scalar value
\[
\overline{\rho}^{\mathrm{dom}}_{\Lambda_k}\big([\![\Xi^{(1)}_{\Lambda_k}]\!]\big)
\]
is nonzero, and the sign
\begin{equation}
\varepsilon_{\Lambda_k}
:=
\operatorname{sgn}\!\Big(
\overline{\rho}^{\mathrm{dom}}_{\Lambda_k}\big([\![\Xi^{(1)}_{\Lambda_k}]\!]\big)
\Big)
\in\{\pm1\}
\label{eq:epsilon-from-overlap}
\end{equation}
is well defined.
Thus the transported descendant sector has nonzero overlap with the canonically selected compressed
dominant odd line.
\end{theorem}

\begin{proof}
The equivalence package is exactly \cref{prop:descendant-overlap-equivalent}.
Once one representative has nonzero dominant projected response, the same proposition gives the
nonzero scalar value on the descended descendant point. The sign in \eqref{eq:epsilon-from-overlap}
is therefore well defined.
\end{proof}

\subsection*{Conditional mod-null reentry on the selected scalar line}
\label{subsec:conditional-mod-null-reentry}

\begin{definition}[Parity-stable mod-null subspace transverse to the selected scalar line]
\label{def:mod-null-transverse-datum}
Fix an accepted tick \(k\), assume dominant-line realization and quotient descent, and let
\[
\overline L^{\mathrm{Main}}_{\Lambda_k}
=
\mathbb R\,\overline\ell^{\mathrm{Main}}_{\Lambda_k}
\subset
\mathfrak D^-_{\Lambda_k}
\]
be a one-dimensional selected scalar line in the ambiguity quotient.
A \emph{parity-stable mod-null datum transverse to}
\(
\overline L^{\mathrm{Main}}_{\Lambda_k}
\)
consists of a linear subspace
\[
N^{\mathrm{mn}}_{\Lambda_k}\subset \mathfrak D^-_{\Lambda_k}
\]
such that:
\begin{enumerate}[leftmargin=2em,label=(\alph*)]
\item \(N^{\mathrm{mn}}_{\Lambda_k}\) is stable under the descended physical parity on
  \(\mathfrak D^-_{\Lambda_k}\);
\item the selected scalar line is transverse to \(N^{\mathrm{mn}}_{\Lambda_k}\):
  \[
  \overline L^{\mathrm{Main}}_{\Lambda_k}\cap N^{\mathrm{mn}}_{\Lambda_k}=\{0\}.
  \]
\end{enumerate}
Define
\[
\mathcal V^{1,\mathrm{mn}}_{\Lambda_k}
:=
\mathfrak D^-_{\Lambda_k}/N^{\mathrm{mn}}_{\Lambda_k}
\]
with quotient map
\[
\pi^{\mathrm{mn}}_{\Lambda_k}:
\mathfrak D^-_{\Lambda_k}\to \mathcal V^{1,\mathrm{mn}}_{\Lambda_k}.
\]
Because \(\mathfrak D^-_{\Lambda_k}\) is itself a quotient of the physical \(\Gamma\)-odd sector,
the descended parity acts on \(\mathfrak D^-_{\Lambda_k}\) by \(-\mathrm{id}\).
Thus the parity-stability clause yields a descended involution
\[
\Gamma^{\mathrm{mn}}_{\Lambda_k}
\bigl(
\pi^{\mathrm{mn}}_{\Lambda_k}(x)
\bigr)
:=
\pi^{\mathrm{mn}}_{\Lambda_k}(-x),
\qquad
x\in \mathfrak D^-_{\Lambda_k}.
\]
\end{definition}

\begin{definition}[Dominant-null parity-stable mod-null datum]
\label{def:dominant-null-mod-null-datum}
Fix an accepted tick \(k\), assume dominant-line realization and quotient descent, and let
\[
\overline L^{\mathrm{Main}}_{\Lambda_k}
=
\mathbb R\,\overline\ell^{\mathrm{Main}}_{\Lambda_k}
\subset
\mathfrak D^-_{\Lambda_k}
\]
be a one-dimensional selected scalar line in the ambiguity quotient.
A parity-stable mod-null datum
\[
N^{\mathrm{mn}}_{\Lambda_k}\subset \mathfrak D^-_{\Lambda_k}
\]
is called \emph{dominant-null} if it is contained in the kernel of the descended dominant
functional:
\[
N^{\mathrm{mn}}_{\Lambda_k}\subset \ker \overline{\rho}^{\mathrm{dom}}_{\Lambda_k}.
\]
It is understood that
\(
N^{\mathrm{mn}}_{\Lambda_k}
\)
is stable under the descended physical parity on
\(
\mathfrak D^-_{\Lambda_k}
\),
and that the associated quotient
\(
\mathfrak D^-_{\Lambda_k}/N^{\mathrm{mn}}_{\Lambda_k}
\)
carries the descended involution from \cref{def:mod-null-transverse-datum}.
The canonical dominant-kernel choice is isolated separately in
\cref{def:canonical-dominant-null-datum}.
The transversality lemma below uses only this kernel-subordination property.
\end{definition}

\begin{definition}[Canonical dominant-null datum]
\label{def:canonical-dominant-null-datum}
Fix an accepted tick \(k\) and assume dominant-line realization together with quotient descent.
Define the \emph{canonical dominant-null datum} by
\[
N^{\mathrm{dom-null}}_{\Lambda_k}
:=
\ker \overline{\rho}^{\mathrm{dom}}_{\Lambda_k}
\subset
\mathfrak D^-_{\Lambda_k}.
\]
This datum depends only on the descended dominant functional on the ambiguity quotient.
No additional reduced-welded or completion-side choice is built into the definition.
\end{definition}

\begin{lemma}[Dominant-kernel criterion for selected-line transversality]
\label{lem:dominant-kernel-transversality}
Fix an accepted tick \(k\), assume dominant-line realization together with quotient descent, and let
\[
\overline L^{\mathrm{Main}}_{\Lambda_k}
=
\mathbb R\,\overline\ell^{\mathrm{Main}}_{\Lambda_k}
\subset
\mathfrak D^-_{\Lambda_k}
\]
be a one-dimensional selected scalar line with
\[
\overline{\rho}^{\mathrm{dom}}_{\Lambda_k}
\!\bigl(
\overline\ell^{\mathrm{Main}}_{\Lambda_k}
\bigr)
\neq 0.
\]
If
\[
N^{\mathrm{mn}}_{\Lambda_k}\subset \mathfrak D^-_{\Lambda_k}
\]
is a dominant-null parity-stable mod-null datum in the sense of
\cref{def:dominant-null-mod-null-datum},
then
\[
\overline L^{\mathrm{Main}}_{\Lambda_k}\cap N^{\mathrm{mn}}_{\Lambda_k}=\{0\}.
\]
So every dominant-null parity-stable mod-null datum is automatically transverse to the selected
scalar line.
\end{lemma}

\begin{proof}
Let
\(
x\in \overline L^{\mathrm{Main}}_{\Lambda_k}\cap N^{\mathrm{mn}}_{\Lambda_k}
\).
Because
\(
\overline L^{\mathrm{Main}}_{\Lambda_k}
=
\mathbb R\,\overline\ell^{\mathrm{Main}}_{\Lambda_k}
\),
there exists \(a\in\mathbb R\) with
\(
x=a\,\overline\ell^{\mathrm{Main}}_{\Lambda_k}
\).
Since
\(
N^{\mathrm{mn}}_{\Lambda_k}\subset
\ker \overline{\rho}^{\mathrm{dom}}_{\Lambda_k}
\),
we have
\[
0
=
\overline{\rho}^{\mathrm{dom}}_{\Lambda_k}(x)
=
a\,
\overline{\rho}^{\mathrm{dom}}_{\Lambda_k}
\!\bigl(
\overline\ell^{\mathrm{Main}}_{\Lambda_k}
\bigr).
\]
The coefficient
\(
\overline{\rho}^{\mathrm{dom}}_{\Lambda_k}
(\overline\ell^{\mathrm{Main}}_{\Lambda_k})
\)
is nonzero by hypothesis, so \(a=0\).
Therefore \(x=0\), proving
\(
\overline L^{\mathrm{Main}}_{\Lambda_k}\cap N^{\mathrm{mn}}_{\Lambda_k}=\{0\}
\).
\end{proof}

\begin{theorem}[The dominant kernel is the canonical dominant-null datum]
\label{thm:canonical-dominant-null-datum}
Fix an accepted tick \(k\), assume dominant-line realization together with quotient descent, and let
\[
\overline L^{\mathrm{Main}}_{\Lambda_k}
=
\mathbb R\,\overline\ell^{\mathrm{Main}}_{\Lambda_k}
\subset
\mathfrak D^-_{\Lambda_k}
\]
be a one-dimensional selected scalar line with
\[
\overline{\rho}^{\mathrm{dom}}_{\Lambda_k}
\!\bigl(
\overline\ell^{\mathrm{Main}}_{\Lambda_k}
\bigr)
\neq 0.
\]
Define
\[
N^{\mathrm{dom-null}}_{\Lambda_k}
:=
\ker \overline{\rho}^{\mathrm{dom}}_{\Lambda_k}
\subset
\mathfrak D^-_{\Lambda_k}.
\]
Then:
\begin{enumerate}[leftmargin=2em]
\item \(N^{\mathrm{dom-null}}_{\Lambda_k}\) is stable under the descended physical parity on
  \(\mathfrak D^-_{\Lambda_k}\);
\item \(N^{\mathrm{dom-null}}_{\Lambda_k}\) is a dominant-null parity-stable mod-null datum in the
  sense of \cref{def:dominant-null-mod-null-datum};
\item the selected scalar line is automatically transverse to the canonical dominant-null datum:
  \[
  \overline L^{\mathrm{Main}}_{\Lambda_k}
  \cap
  N^{\mathrm{dom-null}}_{\Lambda_k}
  =
  \{0\}.
  \]
\end{enumerate}
Hence the canonical dominant kernel is admissible as the mod-null datum for the conditional
reentry theorem.
\end{theorem}

\begin{proof}
Let \(x\in N^{\mathrm{dom-null}}_{\Lambda_k}\).
Because the descended physical parity acts on \(\mathfrak D^-_{\Lambda_k}\) by \(-\mathrm{id}\),
we have
\[
\overline{\rho}^{\mathrm{dom}}_{\Lambda_k}(-x)
=
-\overline{\rho}^{\mathrm{dom}}_{\Lambda_k}(x)
=
0.
\]
So \(-x\in N^{\mathrm{dom-null}}_{\Lambda_k}\), proving parity stability.
Item \textnormal{(2)} is then exactly \cref{def:dominant-null-mod-null-datum} with
\(
N^{\mathrm{mn}}_{\Lambda_k}=N^{\mathrm{dom-null}}_{\Lambda_k}
\).
Item \textnormal{(3)} is \cref{lem:dominant-kernel-transversality} applied to this canonical
choice.
\end{proof}

\begin{lemma}[Transversality reduces mod-null reentry to quotient construction]
\label{lem:mod-null-transversality-reduction}
Fix an accepted tick \(k\), assume dominant-line realization together with quotient descent, and let
\[
\overline L^{\mathrm{Main}}_{\Lambda_k}
=
\mathbb R\,\overline\ell^{\mathrm{Main}}_{\Lambda_k}
\subset
\mathfrak D^-_{\Lambda_k}
\]
be a one-dimensional selected scalar line with
\[
\overline{\rho}^{\mathrm{dom}}_{\Lambda_k}
\!\bigl(
\overline\ell^{\mathrm{Main}}_{\Lambda_k}
\bigr)
\neq 0.
\]
Assume moreover a parity-stable mod-null datum transverse to
\(
\overline L^{\mathrm{Main}}_{\Lambda_k}
\)
in the sense of \cref{def:mod-null-transverse-datum}.
Define
\[
\ell^{\mathrm{sel,mn}}_{\Lambda_k}
:=
\pi^{\mathrm{mn}}_{\Lambda_k}
\!\bigl(
\overline\ell^{\mathrm{Main}}_{\Lambda_k}
\bigr)
\]
and
\[
L^{\mathrm{sel,mn}}_{\Lambda_k}
:=
\mathbb R\,\ell^{\mathrm{sel,mn}}_{\Lambda_k}
\subset
\mathcal V^{1,\mathrm{mn}}_{\Lambda_k}.
\]
Then:
\begin{enumerate}[leftmargin=2em]
\item the restriction of
  \(
  \pi^{\mathrm{mn}}_{\Lambda_k}
  \)
  to
  \(
  \overline L^{\mathrm{Main}}_{\Lambda_k}
  \)
  is injective;
\item \(L^{\mathrm{sel,mn}}_{\Lambda_k}\) is a well-defined one-dimensional line in the mod-null
  quotient, generated by the nonzero class
  \(
  \ell^{\mathrm{sel,mn}}_{\Lambda_k}
  \);
\item the mod-null selected realization map
  \[
  \mathfrak R^{\mathrm{mn}}_{\Lambda_k}:
  L^{\mathrm{sel,mn}}_{\Lambda_k}
  \longrightarrow
  \mathbb R\,u^{\mathrm{LoG}}_{\Lambda_k},
  \qquad
  \mathfrak R^{\mathrm{mn}}_{\Lambda_k}
  \!\bigl(
  a\,\ell^{\mathrm{sel,mn}}_{\Lambda_k}
  \bigr)
  :=
  a\,
  \overline{\rho}^{\mathrm{dom}}_{\Lambda_k}
  \!\bigl(
  \overline\ell^{\mathrm{Main}}_{\Lambda_k}
  \bigr)\,
  u^{\mathrm{LoG}}_{\Lambda_k}
  \]
  is well defined and faithful;
\item the descended involution acts by sign:
  \[
  \Gamma^{\mathrm{mn}}_{\Lambda_k}
  \!\bigl(
  \ell^{\mathrm{sel,mn}}_{\Lambda_k}
  \bigr)
  =
  -
  \ell^{\mathrm{sel,mn}}_{\Lambda_k}.
  \]
\end{enumerate}
In particular, once the transversality condition
\(
\overline L^{\mathrm{Main}}_{\Lambda_k}\cap N^{\mathrm{mn}}_{\Lambda_k}=\{0\}
\)
is known, the remaining mod-null reentry step reduces to quotient construction plus parity descent.
\end{lemma}

\begin{proof}
If
\(
x\in \overline L^{\mathrm{Main}}_{\Lambda_k}
\)
and
\(
\pi^{\mathrm{mn}}_{\Lambda_k}(x)=0
\),
then
\(
x\in N^{\mathrm{mn}}_{\Lambda_k}
\)
by definition of the quotient kernel.
Hence
\(
x\in
\overline L^{\mathrm{Main}}_{\Lambda_k}\cap N^{\mathrm{mn}}_{\Lambda_k}
=
\{0\}
\),
which proves injectivity.
Since
\(
\overline\ell^{\mathrm{Main}}_{\Lambda_k}\neq 0
\),
its image
\(
\ell^{\mathrm{sel,mn}}_{\Lambda_k}
\)
is therefore nonzero, and the image of a one-dimensional line under an injective linear map is
again one-dimensional.
This proves items \textnormal{(1)} and \textnormal{(2)}.

Injectivity on
\(
\overline L^{\mathrm{Main}}_{\Lambda_k}
\)
implies that every element of
\(
L^{\mathrm{sel,mn}}_{\Lambda_k}
\)
has a unique expression
\(
a\,\ell^{\mathrm{sel,mn}}_{\Lambda_k}
\),
so the displayed formula defines
\(
\mathfrak R^{\mathrm{mn}}_{\Lambda_k}
\)
unambiguously.
Its coefficient
\(
\overline{\rho}^{\mathrm{dom}}_{\Lambda_k}
(\overline\ell^{\mathrm{Main}}_{\Lambda_k})
\)
is nonzero by hypothesis, and
\(
u^{\mathrm{LoG}}_{\Lambda_k}\neq 0
\)
by \cref{cor:compressed-export-selection}.
Therefore
\(
\mathfrak R^{\mathrm{mn}}_{\Lambda_k}
(a\,\ell^{\mathrm{sel,mn}}_{\Lambda_k})
=
0
\)
implies
\(
a=0
\),
so the map is faithful.
This proves item \textnormal{(3)}.

By construction of the descended involution,
\[
\Gamma^{\mathrm{mn}}_{\Lambda_k}
\!\bigl(
\ell^{\mathrm{sel,mn}}_{\Lambda_k}
\bigr)
=
\Gamma^{\mathrm{mn}}_{\Lambda_k}
\!\Bigl(
\pi^{\mathrm{mn}}_{\Lambda_k}
\!\bigl(
\overline\ell^{\mathrm{Main}}_{\Lambda_k}
\bigr)
\Bigr)
=
\pi^{\mathrm{mn}}_{\Lambda_k}
\!\bigl(
-\overline\ell^{\mathrm{Main}}_{\Lambda_k}
\bigr)
=
-\ell^{\mathrm{sel,mn}}_{\Lambda_k},
\]
which proves item \textnormal{(4)}.
\end{proof}

\begin{theorem}[Conditional mod-null reentry on the selected scalar line]
\label{thm:conditional-mod-null-reentry}
Fix an accepted tick \(k\) and the next accepted receiver step \(k+1\).
Assume the hypotheses and notation of
\cref{splitlift-thm:splitlift-scalar-descendant-lift}.
In particular, let
\[
\overline L^{\mathrm{Main}}_{\Lambda_{k+1}}
=
\mathbb R\,\overline\ell^{\mathrm{Main}}_{\Lambda_{k+1}}
\subset
\mathfrak D^-_{\Lambda_{k+1}}
\]
be the nonzero selected scalar descendant line supplied there.
Assume moreover a parity-stable mod-null datum transverse to
\(
\overline L^{\mathrm{Main}}_{\Lambda_{k+1}}
\)
in the sense of \cref{def:mod-null-transverse-datum}.
Define
\[
\ell^{\mathrm{sel,mn}}_{\Lambda_{k+1}}
:=
\pi^{\mathrm{mn}}_{\Lambda_{k+1}}
\!\bigl(
\overline\ell^{\mathrm{Main}}_{\Lambda_{k+1}}
\bigr)
\]
and
\[
L^{\mathrm{sel,mn}}_{\Lambda_{k+1}}
:=
\mathbb R\,\ell^{\mathrm{sel,mn}}_{\Lambda_{k+1}}
\subset
\mathcal V^{1,\mathrm{mn}}_{\Lambda_{k+1}}.
\]
Then:
\begin{enumerate}[leftmargin=2em]
\item \(L^{\mathrm{sel,mn}}_{\Lambda_{k+1}}\) is a well-defined one-dimensional selected descendant
  line in the mod-null quotient;
\item the induced mod-null selected realization map
  \[
  \mathfrak R^{\mathrm{mn}}_{\Lambda_{k+1}}:
  L^{\mathrm{sel,mn}}_{\Lambda_{k+1}}
  \longrightarrow
  \mathbb R\,u^{\mathrm{LoG}}_{\Lambda_{k+1}},
  \qquad
  \mathfrak R^{\mathrm{mn}}_{\Lambda_{k+1}}
  \!\bigl(
  a\,\ell^{\mathrm{sel,mn}}_{\Lambda_{k+1}}
  \bigr)
  :=
  a\,
  \overline{\rho}^{\mathrm{dom}}_{\Lambda_{k+1}}
  \!\bigl(
  \overline\ell^{\mathrm{Main}}_{\Lambda_{k+1}}
  \bigr)\,
  u^{\mathrm{LoG}}_{\Lambda_{k+1}}
  \]
  is well defined and faithful;
\item the descended involution acts by sign:
  \[
  \Gamma^{\mathrm{mn}}_{\Lambda_{k+1}}
  \!\bigl(
  \ell^{\mathrm{sel,mn}}_{\Lambda_{k+1}}
  \bigr)
  =
  -
  \ell^{\mathrm{sel,mn}}_{\Lambda_{k+1}}.
  \]
\end{enumerate}
\end{theorem}

\begin{proof}
By \cref{splitlift-thm:splitlift-scalar-descendant-lift},
\[
\overline\ell^{\mathrm{Main}}_{\Lambda_{k+1}}\neq 0
\qquad\text{and}\qquad
\overline{\rho}^{\mathrm{dom}}_{\Lambda_{k+1}}
\!\bigl(
\overline\ell^{\mathrm{Main}}_{\Lambda_{k+1}}
\bigr)
\neq 0.
\]
So every hypothesis of
\cref{lem:mod-null-transversality-reduction}
is in force at \(\Lambda_{k+1}\).
Applying that lemma gives items \textnormal{(1)}--\textnormal{(3)} immediately.
\end{proof}

\begin{corollary}[Canonical dominant-kernel discharge of the mod-null reentry step]
\label{cor:dominant-kernel-mod-null-reentry}
Fix an accepted tick \(k\) and the next accepted receiver step \(k+1\).
Assume the hypotheses and notation of
\cref{splitlift-thm:splitlift-scalar-descendant-lift}.
Define the canonical dominant-null datum
\[
N^{\mathrm{dom-null}}_{\Lambda_{k+1}}
:=
\ker \overline{\rho}^{\mathrm{dom}}_{\Lambda_{k+1}}
\subset
\mathfrak D^-_{\Lambda_{k+1}}.
\]
Then \(N^{\mathrm{dom-null}}_{\Lambda_{k+1}}\) is admissible as the mod-null datum of
\cref{thm:conditional-mod-null-reentry}, the transversality hypothesis of that theorem is
automatic, and hence the conclusions of that theorem hold with this canonical datum.
\end{corollary}

\begin{proof}
By \cref{splitlift-thm:splitlift-scalar-descendant-lift},
\[
\overline{\rho}^{\mathrm{dom}}_{\Lambda_{k+1}}
\!\bigl(
\overline\ell^{\mathrm{Main}}_{\Lambda_{k+1}}
\bigr)
\neq 0.
\]
So \cref{thm:canonical-dominant-null-datum} applies at \(\Lambda_{k+1}\).
It shows that
\(
N^{\mathrm{dom-null}}_{\Lambda_{k+1}}
\)
is a dominant-null parity-stable mod-null datum and that
\[
\overline L^{\mathrm{Main}}_{\Lambda_{k+1}}
\cap
N^{\mathrm{dom-null}}_{\Lambda_{k+1}}
=
\{0\}.
\]
This is exactly the admissibility and transversality clause required by
\cref{thm:conditional-mod-null-reentry}, so that theorem now applies directly with the canonical
datum.
\end{proof}

\begin{remark}[Conditional character of the reentry step]
\label{rem:conditional-mod-null-reentry}
\cref{thm:conditional-mod-null-reentry} is a quotient-level reentry statement.
After
\cref{lem:dominant-kernel-transversality,thm:canonical-dominant-null-datum,cor:dominant-kernel-mod-null-reentry},
the canonical dominant-null datum is fixed as
\(
\ker\overline{\rho}^{\mathrm{dom}}_{\Lambda_{k+1}}
\)
on \(\mathfrak D^-_{\Lambda_{k+1}}\).
What remains is passage to the mod-null quotient together with parity descent on the odd line.
Accordingly, the theorem should not be read as a statement on the raw descendant ambiguity quotient
or on any stronger physical quotient.
\end{remark}

\subsection*{Branchwise orientation transport}
\label{subsec:branchwise-orientation}

This subsection studies transport of the descended descendant sign along a connected accepted
branch. The sign can be transported once the descended descendant value varies continuously with the
cutoff. The compressed spectral-selection theorem already gives continuity of the dominant
projector, and \cref{thm:dominant-line-persistence} gives continuity of the dominant physical line.
We combine those inputs with continuity of the projected descendant response.

\begin{definition}[Branchwise descendant overlap function]
\label{def:branchwise-overlap-function}
Let \(I\) be a connected accepted-branch segment.
Assume dominant-line realization and quotient descent for every \(\Lambda\in I\).
Suppose there is a branchwise choice of representatives
\[
J^{\mathrm{desc}}_{\Lambda}\in[\Xi^{(1)}_{\Lambda}]
\qquad
(\Lambda\in I)
\]
such that the projected covariance response
\[
\Lambda\longmapsto
\Pi_{\mathrm{phys},\Lambda}\mathcal R_{\Lambda}(J^{\mathrm{desc}}_{\Lambda})
\]
is continuous in the Fisher topology.
Define the associated branchwise descendant overlap function by
\begin{equation}
m_I(\Lambda)
:=
\rho^{\mathrm{dom}}_{\Lambda}(J^{\mathrm{desc}}_{\Lambda})
=
\overline{\rho}^{\mathrm{dom}}_{\Lambda}\big([\![\Xi^{(1)}_{\Lambda}]\!]\big).
\label{eq:branchwise-overlap-function}
\end{equation}
\end{definition}

\begin{proposition}[Continuity of the descended overlap scalar]
\label{prop:branchwise-overlap-continuity}
Let \(I\) be a connected accepted-branch segment.
Assume:
\begin{enumerate}[leftmargin=2em]
\item the hypotheses of \cref{thm:dominant-line-persistence} on \(I\);
\item quotient descent holds for every \(\Lambda\in I\);
\item the branchwise choice of representatives in
  \cref{def:branchwise-overlap-function} exists.
\end{enumerate}
Then the scalar function \(m_I:I\to\mathbb R\) is well defined, independent of the chosen
branchwise representatives, and continuous.
Moreover,
\begin{equation}
|m_I(\Lambda)|
=
\nu_{\Lambda}
\qquad
\forall \Lambda\in I.
\label{eq:branchwise-overlap-margin}
\end{equation}
\end{proposition}

\begin{proof}
Independence of the chosen branchwise representatives is exactly the coset constancy proved in
\cref{cor:dominant-quotient-constancy}: if
\(
J'^{\mathrm{desc}}_{\Lambda}-J^{\mathrm{desc}}_{\Lambda}\in\mathcal A^-_{\Lambda}
\),
then
\[
\rho^{\mathrm{dom}}_{\Lambda}(J'^{\mathrm{desc}}_{\Lambda})
=
\rho^{\mathrm{dom}}_{\Lambda}(J^{\mathrm{desc}}_{\Lambda})
=
\overline{\rho}^{\mathrm{dom}}_{\Lambda}\big([\![\Xi^{(1)}_{\Lambda}]\!]\big).
\]
Thus \(m_I\) is well defined.

By \cref{thm:dominant-line-persistence}, the normalized dominant vector
\(
\widehat{\delta C}^{\mathrm{LoG}}_{\Lambda}
\)
and the dominant projector
\(
\Pi^{\mathrm{dom}}_{\Lambda}
\)
depend continuously on \(\Lambda\in I\).
By assumption, the response vector
\[
X_{\Lambda}
:=
\Pi_{\mathrm{phys},\Lambda}\mathcal R_{\Lambda}(J^{\mathrm{desc}}_{\Lambda})
\]
is also continuous on \(I\).
Using \cref{def:dominant-cov-projector}, we have
\[
\Pi^{\mathrm{dom}}_{\Lambda}X_{\Lambda}
=
m_I(\Lambda)\,\widehat{\delta C}^{\mathrm{LoG}}_{\Lambda}.
\]
The left-hand side is continuous because both factors are continuous, while
\(
\widehat{\delta C}^{\mathrm{LoG}}_{\Lambda}\neq 0
\)
for all \(\Lambda\in I\) by \cref{thm:dominant-line-persistence}.
Therefore the scalar coefficient \(m_I(\Lambda)\) is continuous.

Finally, \eqref{eq:branchwise-overlap-margin} is just the definition of \(\nu_{\Lambda}\) from
\cref{def:descendant-overlap-margin} together with \eqref{eq:branchwise-overlap-function}.
\end{proof}

\begin{theorem}[Branchwise orientation transport on the descendant class]
\label{thm:branchwise-orientation-transport}
Let \(I\) be a connected accepted-branch segment and assume the hypotheses of
\cref{prop:branchwise-overlap-continuity}.
Then:
\begin{enumerate}[leftmargin=2em]
\item on every connected subsegment \(I'\subset I\) such that
  \[
  \nu_{\Lambda}>0
  \qquad
  \forall \Lambda\in I',
  \]
  the sign
  \begin{equation}
  \varepsilon_{\Lambda}
  :=
  \operatorname{sgn} m_I(\Lambda)
  =
  \operatorname{sgn}\!\Big(
  \overline{\rho}^{\mathrm{dom}}_{\Lambda}\big([\![\Xi^{(1)}_{\Lambda}]\!]\big)
  \Big)
  \in\{\pm1\}
  \label{eq:branchwise-epsilon}
  \end{equation}
  is constant on \(I'\);
\item if \(\nu_{\Lambda_0}>0\) at one point \(\Lambda_0\in I\), then there exists a neighborhood
  \(I_0\subset I\) of \(\Lambda_0\) on which \(\varepsilon_{\Lambda}\) is constant and equal to
  \(\varepsilon_{\Lambda_0}\);
\item under the present hypotheses, a change of sign can occur only at a cutoff where the
  descendant overlap closes
  \(
  \nu_{\Lambda}=0
  \).
\end{enumerate}
In particular, once dominant-line realization, quotient descent, and nonzero overlap are
established on an accepted branch segment and the projected descendant response varies continuously
there, the branchwise sign statement of \cref{thm:active-response-transport-minimal} follows.
\end{theorem}

\begin{proof}
By \cref{prop:branchwise-overlap-continuity}, \(m_I(\Lambda)\) is a continuous real-valued function
on \(I\).
On any connected subsegment \(I'\) where \(\nu_{\Lambda}=|m_I(\Lambda)|>0\), the function \(m_I\)
never vanishes.
A continuous nonvanishing real-valued function on a connected set has constant sign, proving item
(1).

Item (2) is the local version of the same argument: if \(m_I(\Lambda_0)\neq 0\), continuity gives a
neighborhood \(I_0\) on which \(m_I\) keeps the sign it has at \(\Lambda_0\).

For item (3), any sign change would force \(m_I\) to pass through zero by the intermediate value
theorem, hence \(\nu_{\Lambda}=|m_I(\Lambda)|=0\) at some intermediate cutoff.
\end{proof}

\begin{corollary}[Local orientation persistence from one-tick activity]
\label{cor:branchwise-orientation-local}
Let \(I\) be a connected accepted-branch segment and assume the hypotheses of
\cref{prop:branchwise-overlap-continuity}.
If the transported descendant class is dominant-active at one cutoff \(\Lambda_0\in I\), i.e.
\[
\nu_{\Lambda_0}>0,
\]
then there exists a neighborhood \(I_0\subset I\) of \(\Lambda_0\) such that
\[
\varepsilon_{\Lambda}=\varepsilon_{\Lambda_0}
\qquad
\forall \Lambda\in I_0.
\]
\end{corollary}

\begin{proof}
This is exactly item (2) of \cref{thm:branchwise-orientation-transport}.
\end{proof}

\subsection*{Witness-indexed seam realization and scalar bridge}
\label{subsec:witness-seam-bridge}

The bridge used below factors through a common LS witness datum rather than a direct identification
of a UOM seam source with a descendant representative.
On the UOM side one realizes a filler tag together with a detector witness by an admissible odd
welded source, measures the resulting dominant physical pairing, and then compares that scalar with
the odd representative attached to the same witness.

\begin{definition}[Accepted record payload of a welded odd source]
\label{def:accepted-record-payload}
Fix an accepted tick \(k\).
For an admissible welded odd trace-channel source \(\tau\), let
\[
\rho_{\Lambda_k}[\tau]
\]
denote the receiver state at cutoff \(\Lambda_k\) produced by that source in the welded UOM channel.
Define its accepted record payload by
\begin{equation}
\mathcal P^{\mathrm{acc}}_{\Lambda_k}[\tau]
:=
\Pi_{\Delta(\Lambda_k)}\!\big(\rho_{\Lambda_k}[\tau]\big).
\label{eq:accepted-record-payload}
\end{equation}
This is the diagonal/record part retained by the finite-cutoff acceptance map.
\end{definition}

\begin{definition}[Witness-indexed seam realization at an accepted tick]
\label{def:witness-seam-realization}
Fix an accepted tick \(k\).
Let \(y\) be a witness on the LS difference object, and write \(\mathrm{flip}(y)\) for the induced
LS involution.
A \emph{witness-indexed seam realization} at \(\Lambda_k\) is a choice of admissible welded odd
trace-channel sources
\[
\tau^{f,y}_{\Lambda_k},
\qquad
f\in\{\mathrm{diag},\mathrm{bridge}\},
\]
such that:
\begin{enumerate}[leftmargin=2em]
\item \textbf{filler toggle changes the odd seam sign:}
  \[
  \tau^{\mathrm{bridge},y}_{\Lambda_k}
  =
  -\,\tau^{\mathrm{diag},y}_{\Lambda_k};
  \]
\item \textbf{witness flip changes the odd seam sign:}
  \[
  \tau^{f,\mathrm{flip}(y)}_{\Lambda_k}
  =
  -\,\tau^{f,y}_{\Lambda_k}
  \qquad
  \forall f\in\{\mathrm{diag},\mathrm{bridge}\};
  \]
\item \textbf{accepted record payload is unchanged under filler toggle:}
  \[
  \mathcal P^{\mathrm{acc}}_{\Lambda_k}\!\big[\tau^{\mathrm{bridge},y}_{\Lambda_k}\big]
  =
  \mathcal P^{\mathrm{acc}}_{\Lambda_k}\!\big[\tau^{\mathrm{diag},y}_{\Lambda_k}\big].
  \]
\end{enumerate}
\end{definition}

\begin{remark}[Accepted-record invariance versus seam-scalar use]
\label{rem:accepted-record-vs-seam}
The accepted-record invariance clause in \cref{def:witness-seam-realization} encodes the physical
indistinguishability of filler toggle at the accepted-record layer.
The seam-scalar theorems below use only the induced odd sign relations in the source channel.
\end{remark}

\begin{definition}[Raw and normalized seam pairings]
\label{def:raw-seam-pairing}
Fix an accepted tick \(k\) and a witness-indexed seam realization as in
\cref{def:witness-seam-realization}.
For \(f\in\{\mathrm{diag},\mathrm{bridge}\}\) and witness \(y\), define the \emph{raw seam
pairing} by
\begin{equation}
\widetilde s_{\Lambda_k}(f,y)
:=
\big\langle\!\big\langle
\Pi_{\mathrm{phys},\Lambda_k}\mathsf T_{\Lambda_k}\!\big[\tau^{f,y}_{\Lambda_k}\big],\,
\delta C^{\mathrm{LoG}}_{\Lambda_k}
\big\rangle\!\big\rangle_{C_{\ast,\Lambda_k}}.
\label{eq:raw-seam-pairing}
\end{equation}
Whenever dominant-line realization holds so that
\(
\widehat{\delta C}^{\mathrm{LoG}}_{\Lambda_k}
\)
is defined, also set
\begin{equation}
s_{\Lambda_k}(f,y)
:=
\big\langle\!\big\langle
\Pi_{\mathrm{phys},\Lambda_k}\mathsf T_{\Lambda_k}\!\big[\tau^{f,y}_{\Lambda_k}\big],\,
\widehat{\delta C}^{\mathrm{LoG}}_{\Lambda_k}
\big\rangle\!\big\rangle_{C_{\ast,\Lambda_k}}.
\label{eq:normalized-seam-pairing}
\end{equation}
\end{definition}

\begin{proposition}[Filler and witness parity of the seam scalar]
\label{prop:seam-scalar-signs}
Fix an accepted tick \(k\) and assume a witness-indexed seam realization at \(\Lambda_k\).
Then the raw seam pairing satisfies
\begin{equation}
\widetilde s_{\Lambda_k}(\mathrm{bridge},y)
=
-\,\widetilde s_{\Lambda_k}(\mathrm{diag},y),
\qquad
\widetilde s_{\Lambda_k}(f,\mathrm{flip}(y))
=
-\,\widetilde s_{\Lambda_k}(f,y).
\label{eq:raw-seam-signs}
\end{equation}
If dominant-line realization also holds, then the same identities hold for the normalized seam
pairing:
\begin{equation}
s_{\Lambda_k}(\mathrm{bridge},y)
=
-\,s_{\Lambda_k}(\mathrm{diag},y),
\qquad
s_{\Lambda_k}(f,\mathrm{flip}(y))
=
-\,s_{\Lambda_k}(f,y).
\label{eq:normalized-seam-signs}
\end{equation}
\end{proposition}

\begin{proof}
The map
\(
\tau\mapsto \Pi_{\mathrm{phys},\Lambda_k}\mathsf T_{\Lambda_k}[\tau]
\)
is linear because \(\mathsf T_{\Lambda_k}\) and \(\Pi_{\mathrm{phys},\Lambda_k}\) are linear.
The Fisher pairing is bilinear.
Applying the sign relations of \cref{def:witness-seam-realization} therefore gives
\[
\Pi_{\mathrm{phys},\Lambda_k}\mathsf T_{\Lambda_k}\!\big[\tau^{\mathrm{bridge},y}_{\Lambda_k}\big]
=
-\,\Pi_{\mathrm{phys},\Lambda_k}\mathsf T_{\Lambda_k}\!\big[\tau^{\mathrm{diag},y}_{\Lambda_k}\big]
\]
and
\[
\Pi_{\mathrm{phys},\Lambda_k}\mathsf T_{\Lambda_k}\!\big[\tau^{f,\mathrm{flip}(y)}_{\Lambda_k}\big]
=
-\,\Pi_{\mathrm{phys},\Lambda_k}\mathsf T_{\Lambda_k}\!\big[\tau^{f,y}_{\Lambda_k}\big].
\]
Pairing with \(\delta C^{\mathrm{LoG}}_{\Lambda_k}\) proves \eqref{eq:raw-seam-signs}.
If dominant-line realization holds, then
\(
\widehat{\delta C}^{\mathrm{LoG}}_{\Lambda_k}
\)
is a fixed nonzero normalization of the same vector, so the same argument proves
\eqref{eq:normalized-seam-signs}.
\end{proof}

\begin{definition}[Detector-factorized seam coupling]
\label{def:detector-factorized-seam}
Fix an accepted tick \(k\).
A witness-indexed seam realization is called \emph{detector-factorized} if there exists a single
admissible welded odd trace-channel source
\[
\tau^{\mathrm{unit}}_{\Lambda_k}
\]
and the LS detector factor
\[
\overline{\mathrm{obs}}
:
\text{LS difference object}\to\mathbb R
\]
such that, for every filler tag \(f\) and witness \(y\),
\begin{equation}
\tau^{f,y}_{\Lambda_k}
=
\sigma_{\mathrm{fill}}(f)\,
\overline{\mathrm{obs}}(y)\,
\tau^{\mathrm{unit}}_{\Lambda_k},
\qquad
\sigma_{\mathrm{fill}}(\mathrm{diag})=+1,\quad
\sigma_{\mathrm{fill}}(\mathrm{bridge})=-1.
\label{eq:detector-factorized-seam}
\end{equation}
\end{definition}

\begin{theorem}[Seam faithfulness scalar]
\label{thm:seam-faithfulness-scalar}
Fix an accepted tick \(k\) and assume a detector-factorized witness-indexed seam realization at
\(\Lambda_k\).
Define the fixed-cutoff seam coefficient by
\begin{equation}
\widetilde c_{\Lambda_k}
:=
\big\langle\!\big\langle
\Pi_{\mathrm{phys},\Lambda_k}\mathsf T_{\Lambda_k}\!\big[\tau^{\mathrm{unit}}_{\Lambda_k}\big],\,
\delta C^{\mathrm{LoG}}_{\Lambda_k}
\big\rangle\!\big\rangle_{C_{\ast,\Lambda_k}}.
\label{eq:seam-faithfulness-coeff}
\end{equation}
Then, for every filler tag \(f\) and witness \(y\),
\begin{equation}
\widetilde s_{\Lambda_k}(f,y)
=
\sigma_{\mathrm{fill}}(f)\,
\widetilde c_{\Lambda_k}\,
\overline{\mathrm{obs}}(y).
\label{eq:raw-seam-factorization}
\end{equation}
In particular, if
\[
\widetilde c_{\Lambda_k}\neq 0
\qquad\text{and}\qquad
\overline{\mathrm{obs}}(y)\neq 0,
\]
then the raw seam pairing is nonzero:
\[
\widetilde s_{\Lambda_k}(f,y)\neq 0.
\]
Consequently the compressed welded LoG line is transversely physical at \(\Lambda_k\), hence
dominant-line realization holds at that cutoff by \cref{prop:dominant-line-equivalent}.
When dominant-line realization holds, the normalized seam coefficient
\begin{equation}
c_{\Lambda_k}
:=
\kappa_{\Lambda_k}^{-1/2}\,\widetilde c_{\Lambda_k}
\label{eq:normalized-seam-coeff}
\end{equation}
is defined and
\begin{equation}
s_{\Lambda_k}(f,y)
=
\sigma_{\mathrm{fill}}(f)\,
c_{\Lambda_k}\,
\overline{\mathrm{obs}}(y).
\label{eq:normalized-seam-factorization}
\end{equation}
\end{theorem}

\begin{proof}
Substitute the detector-factorized form \eqref{eq:detector-factorized-seam} into the definition
\eqref{eq:raw-seam-pairing}.
Linearity of \(\mathsf T_{\Lambda_k}\) and bilinearity of the Fisher pairing give
\[
\widetilde s_{\Lambda_k}(f,y)
=
\sigma_{\mathrm{fill}}(f)\,\overline{\mathrm{obs}}(y)\,
\big\langle\!\big\langle
\Pi_{\mathrm{phys},\Lambda_k}\mathsf T_{\Lambda_k}\!\big[\tau^{\mathrm{unit}}_{\Lambda_k}\big],\,
\delta C^{\mathrm{LoG}}_{\Lambda_k}
\big\rangle\!\big\rangle_{C_{\ast,\Lambda_k}},
\]
which is exactly \eqref{eq:raw-seam-factorization}.

If \(\widetilde c_{\Lambda_k}\neq 0\) and \(\overline{\mathrm{obs}}(y)\neq 0\), then
\(
\widetilde s_{\Lambda_k}(f,y)\neq 0
\).
Set
\[
B_{\Lambda_k}
:=
\Pi_{\mathrm{phys},\Lambda_k}\mathsf T_{\Lambda_k}\!\big[\tau^{f,y}_{\Lambda_k}\big].
\]
Then
\[
\big\langle\!\big\langle
\delta C^{\mathrm{LoG}}_{\Lambda_k},\,B_{\Lambda_k}
\big\rangle\!\big\rangle_{C_{\ast,\Lambda_k}}
=
\widetilde s_{\Lambda_k}(f,y)
\neq 0.
\]
By item (4) of \cref{prop:dominant-line-equivalent}, this proves transversality of the compressed
line and hence dominant-line realization at \(\Lambda_k\).

Once dominant-line realization holds, divide \eqref{eq:raw-seam-factorization} by
\(
\|\delta C^{\mathrm{LoG}}_{\Lambda_k}\|_{C_{\ast,\Lambda_k}}
=
\kappa_{\Lambda_k}^{1/2}
\)
to obtain \eqref{eq:normalized-seam-factorization}.
\end{proof}

\begin{remark}[Downstream role of the seam-coefficient hypothesis]
The fixed-cutoff nonvanishing hypothesis
\(
\widetilde c_{\Lambda_k}\neq 0
\)
is exactly the input that, once dominant-line realization defines
\(
c_{\Lambda_k}
\),
is imported downstream as the nondegeneracy clause
\(
c_{\Lambda_k}\neq 0
\)
inside the delayed selected-line package.
This note does not derive that clause from split-threshold analysis alone.
\end{remark}

\begin{definition}[Odd source representability on the characterization layer]
\label{def:odd-source-representability}
Fix an accepted tick \(k\).
Let \(\tau\) be an admissible welded odd trace-channel source.
We say that \(\tau\) is \emph{representable on the characterization layer} if there exists an odd
sufficiently local generator
\[
F_{\tau,\Lambda_k}\in E^-_{\Lambda_k}
\]
such that its intrinsic physical covariance response agrees with the UOM welded response:
\begin{equation}
\Pi_{\mathrm{phys},\Lambda_k}\mathcal R_{\Lambda_k}\!\big(F_{\tau,\Lambda_k}\big)
=
\Pi_{\mathrm{phys},\Lambda_k}\mathsf T_{\Lambda_k}[\tau].
\label{eq:odd-source-representability}
\end{equation}
This is the strong source-insertion realization statement.
\end{definition}

\begin{theorem}[UOM source-insertion representability theorem]
\label{thm:odd-source-representability}
Fix an accepted tick \(k\).
Let \(\tau\) be an admissible welded odd trace-channel source at cutoff \(\Lambda_k\), in the
finite-cutoff UOM sense: its receiver-head profile is band-limited at \(\Lambda_k\) and compactly
supported inside the accepted slab.
Then there exists a canonically associated odd sufficiently local generator
\[
F_{\tau,\Lambda_k}\in E^-_{\Lambda_k}
\]
such that
\begin{equation}
\Pi_{\mathrm{phys},\Lambda_k}\mathcal R_{\Lambda_k}\!\big(F_{\tau,\Lambda_k}\big)
=
\Pi_{\mathrm{phys},\Lambda_k}\mathsf T_{\Lambda_k}[\tau].
\label{eq:odd-source-representability-thm}
\end{equation}
In particular, every witness-indexed seam source used in the present note is representable on the
characterization layer.
\end{theorem}

\begin{proof}
By the finite-cutoff UOM Peierls package, admissible receiver-head sources at cutoff \(\Lambda_k\)
are band-limited boundary test data and the sufficiently local class
\(\mathcal F_{\rm s.loc}^{\Lambda_k}\) contains all compact smearings of local boundary densities.
Let
\[
F_{\tau,\Lambda_k}
\]
be the linear source functional obtained by pairing the odd trace-channel source profile \(\tau\)
with the corresponding odd trace-channel boundary field on the receiver head.
Because this is the linear case of the sufficiently local form and \(\tau\) is physically odd,
\[
F_{\tau,\Lambda_k}\in E^-_{\Lambda_k}.
\]

Now apply the UOM source insertion map from the finite-cutoff retarded-kernel package.
The source \(\tau\) is sent to the bulk forcing term in the linearized equation; the unique
retarded solution is obtained by the retarded Green operator; and boundary restriction gives the
retarded covariance variation. By the UOM identification of the induced retarded kernel with the SK
mixed second derivative, this covariance variation is exactly the welded source-to-covariance
response
\[
\mathsf T_{\Lambda_k}[\tau].
\]

On the other hand, by the definition of the intrinsic covariance-response map, the matrix elements
of
\(
\mathcal R_{\Lambda_k}(F_{\tau,\Lambda_k})
\)
are computed by the Peierls action of \(F_{\tau,\Lambda_k}\) on the sufficiently local quadratic
probes representing boundary covariance entries.
That Peierls variation is precisely the same retarded second-moment variation generated by the
source insertion \(\tau\).
Therefore
\[
\mathcal R_{\Lambda_k}(F_{\tau,\Lambda_k})
=
\mathsf T_{\Lambda_k}[\tau]
\]
as covariance tangents.
Applying the physical projector \(\Pi_{\mathrm{phys},\Lambda_k}\) gives
\eqref{eq:odd-source-representability-thm}.

The witness-indexed seam sources of \cref{def:witness-seam-realization} are, by construction, live
admissible welded odd trace-channel sources at the accepted tick, so the same argument applies to
each of them.
\end{proof}

\begin{proposition}[Dominant pairing for a representable odd source]
\label{prop:odd-source-pairing}
Fix an accepted tick \(k\) and let \(\tau\) be an admissible welded odd trace-channel source that is
representable in the sense of \cref{thm:odd-source-representability}.
Then
\begin{equation}
\big\langle\!\big\langle
\Pi_{\mathrm{phys},\Lambda_k}\mathcal R_{\Lambda_k}\!\big(F_{\tau,\Lambda_k}\big),\,
\delta C^{\mathrm{LoG}}_{\Lambda_k}
\big\rangle\!\big\rangle_{C_{\ast,\Lambda_k}}
=
\big\langle\!\big\langle
\Pi_{\mathrm{phys},\Lambda_k}\mathsf T_{\Lambda_k}[\tau],\,
\delta C^{\mathrm{LoG}}_{\Lambda_k}
\big\rangle\!\big\rangle_{C_{\ast,\Lambda_k}}.
\label{eq:representable-raw-pairing}
\end{equation}
If dominant-line realization also holds, then
\begin{equation}
\rho^{\mathrm{dom}}_{\Lambda_k}\!\big(F_{\tau,\Lambda_k}\big)
=
\big\langle\!\big\langle
\Pi_{\mathrm{phys},\Lambda_k}\mathsf T_{\Lambda_k}[\tau],\,
\widehat{\delta C}^{\mathrm{LoG}}_{\Lambda_k}
\big\rangle\!\big\rangle_{C_{\ast,\Lambda_k}}.
\label{eq:representable-normalized-pairing}
\end{equation}
\end{proposition}

\begin{proof}
Equation \eqref{eq:odd-source-representability} identifies the two physical covariance tangents.
Pairing both sides with \(\delta C^{\mathrm{LoG}}_{\Lambda_k}\) gives
\eqref{eq:representable-raw-pairing}.
If dominant-line realization holds, pair with the normalized vector
\(
\widehat{\delta C}^{\mathrm{LoG}}_{\Lambda_k}
\)
instead and use the definition \eqref{eq:rho-dom-def} of
\(
\rho^{\mathrm{dom}}_{\Lambda_k}
\)
to obtain \eqref{eq:representable-normalized-pairing}.
\end{proof}

\begin{corollary}[Unit welded odd source representability and scalar identification]
\label{cor:unit-source-representability}
Fix an accepted tick \(k\).
Assume a detector-factorized witness-indexed seam realization at \(\Lambda_k\) in the sense of
\cref{def:detector-factorized-seam}, and assume dominant-line realization at \(\Lambda_k\).
Then the unit welded odd source
\(
\tau^{\mathrm{unit}}_{\Lambda_k}
\)
is representable on the characterization layer, with canonical representable generator
\[
F_{\tau^{\mathrm{unit}}_{\Lambda_k},\Lambda_k}\in E^-_{\Lambda_k},
\]
and
\begin{equation}
\rho^{\mathrm{dom}}_{\Lambda_k}
\!\bigl(
F_{\tau^{\mathrm{unit}}_{\Lambda_k},\Lambda_k}
\bigr)
=
\big\langle\!\big\langle
\Pi_{\mathrm{phys},\Lambda_k}\mathsf T_{\Lambda_k}\!\big[\tau^{\mathrm{unit}}_{\Lambda_k}\big],\,
\widehat{\delta C}^{\mathrm{LoG}}_{\Lambda_k}
\big\rangle\!\big\rangle_{C_{\ast,\Lambda_k}}
=
c_{\Lambda_k}.
\label{eq:unit-source-representability-scalar}
\end{equation}
\end{corollary}

\begin{proof}
Representability of
\(
\tau^{\mathrm{unit}}_{\Lambda_k}
\)
is the specialization of \cref{thm:odd-source-representability} to the admissible welded odd trace
source singled out by \cref{def:detector-factorized-seam}.
Because dominant-line realization holds, \cref{prop:odd-source-pairing} applies to that
representable source and gives the first equality in \eqref{eq:unit-source-representability-scalar}.
The second equality is exactly the definition of the normalized seam coefficient
\(
c_{\Lambda_k}
\)
from \cref{thm:seam-faithfulness-scalar}.
\end{proof}

\begin{definition}[Scalar comparison shadow of the witness-indexed bridge]
\label{def:witness-descendant-comparison}
Fix an accepted tick \(k\).
Let \(y\) be an LS witness and let
\(
\tau^{f,y}_{\Lambda_k}
\)
be a witness-indexed seam realization source for a filler tag \(f\).
A \emph{scalar comparison shadow} for that data is a choice of representative
\[
J^{f,y}_{\Lambda_k}\in[\Xi^{(1)}_{\Lambda_k}]
\]
such that the dominant raw pairing of the 2T-side representative agrees with the dominant raw
pairing of the corresponding representable UOM source:
\begin{equation}
\big\langle\!\big\langle
\Pi_{\mathrm{phys},\Lambda_k}\mathcal R_{\Lambda_k}\!\big(J^{f,y}_{\Lambda_k}\big),\,
\delta C^{\mathrm{LoG}}_{\Lambda_k}
\big\rangle\!\big\rangle_{C_{\ast,\Lambda_k}}
=
\big\langle\!\big\langle
\Pi_{\mathrm{phys},\Lambda_k}\mathcal R_{\Lambda_k}\!\big(F_{\tau^{f,y}_{\Lambda_k},\Lambda_k}\big),\,
\delta C^{\mathrm{LoG}}_{\Lambda_k}
\big\rangle\!\big\rangle_{C_{\ast,\Lambda_k}}.
\label{eq:witness-descendant-comparison}
\end{equation}
When dominant-line realization holds, this is equivalently the equality of normalized dominant
pairings
\begin{equation}
\rho^{\mathrm{dom}}_{\Lambda_k}\!\big(J^{f,y}_{\Lambda_k}\big)
=
\rho^{\mathrm{dom}}_{\Lambda_k}\!\big(F_{\tau^{f,y}_{\Lambda_k},\Lambda_k}\big).
\label{eq:witness-descendant-comparison-normalized}
\end{equation}
\end{definition}

\begin{definition}[Witness-indexed source-compatible descendant realization]
\label{def:source-compatible-witness-realization}
Fix an accepted tick \(k\), a filler tag \(f\), a witness \(y\), and a witness-indexed seam source
\(
\tau^{f,y}_{\Lambda_k}
\).
A representative
\[
J^{f,y}_{\Lambda_k}\in[\Xi^{(1)}_{\Lambda_k}]
\]
is called a \emph{witness-indexed source-compatible descendant realization} if it realizes, on the
same LS witness datum, the Limited Scope output and odd Hamiltonization clauses from item
\textbf{(M1)} and item \textbf{(M4)} of Definition~\ref{main-def:uom-main-theorem-assumptions} and differs
from the source-insertion representative only by an ambiguity direction:
\begin{equation}
J^{f,y}_{\Lambda_k}
-
F_{\tau^{f,y}_{\Lambda_k},\Lambda_k}
\in
\mathcal A^-_{\Lambda_k},
\label{eq:source-compatible-witness-realization}
\end{equation}
where
\(
F_{\tau^{f,y}_{\Lambda_k},\Lambda_k}
\)
is the representable source generator furnished by
\cref{thm:odd-source-representability}.
Equivalently, the witness-indexed transported descendant class and the witness-indexed UOM
source-insertion realization define the same point of the ambiguity quotient
\(
E^-_{\Lambda_k}/\mathcal A^-_{\Lambda_k}
\).
\end{definition}
\begin{theorem}[Same-witness descended-scalar comparison from a source-compatible descendant realization]
\label{thm:witness-descended-scalar-comparison}
Fix an accepted tick \(k\), assume dominant-line realization at \(\Lambda_k\), and let
\(
\tau^{f,y}_{\Lambda_k}
\)
be a witness-indexed seam realization source.
Assume:
\begin{enumerate}[leftmargin=2em]
\item the source \(\tau^{f,y}_{\Lambda_k}\) is represented by the canonical generator
  \(
  F_{\tau^{f,y}_{\Lambda_k},\Lambda_k}
  \)
  from \cref{thm:odd-source-representability};
\item there exists a witness-indexed source-compatible descendant realization
  \(
  J^{f,y}_{\Lambda_k}\in[\Xi^{(1)}_{\Lambda_k}]
  \)
  in the sense of \cref{def:source-compatible-witness-realization};
\item the three vanishing statements of \cref{prop:dominant-quotient-generators} hold at
  \(\Lambda_k\).
\end{enumerate}
Then the raw and normalized comparison identities of
\cref{def:witness-descendant-comparison} hold:
\begin{align}
\big\langle\!\big\langle
\Pi_{\mathrm{phys},\Lambda_k}\mathcal R_{\Lambda_k}\!\big(J^{f,y}_{\Lambda_k}\big),\,
\delta C^{\mathrm{LoG}}_{\Lambda_k}
\big\rangle\!\big\rangle_{C_{\ast,\Lambda_k}}
&=
\big\langle\!\big\langle
\Pi_{\mathrm{phys},\Lambda_k}\mathcal R_{\Lambda_k}\!\big(F_{\tau^{f,y}_{\Lambda_k},\Lambda_k}\big),\,
\delta C^{\mathrm{LoG}}_{\Lambda_k}
\big\rangle\!\big\rangle_{C_{\ast,\Lambda_k}},
\label{eq:witness-descendant-comparison-thm-raw}
\\
\rho^{\mathrm{dom}}_{\Lambda_k}\!\big(J^{f,y}_{\Lambda_k}\big)
&=
\rho^{\mathrm{dom}}_{\Lambda_k}\!\big(F_{\tau^{f,y}_{\Lambda_k},\Lambda_k}\big).
\label{eq:witness-descendant-comparison-thm-normalized}
\end{align}
Moreover,
\begin{equation}
\rho^{\mathrm{dom}}_{\Lambda_k}\!\big(F_{\tau^{f,y}_{\Lambda_k},\Lambda_k}\big)
=
\big\langle\!\big\langle
\Pi_{\mathrm{phys},\Lambda_k}\mathsf T_{\Lambda_k}\!\big[\tau^{f,y}_{\Lambda_k}\big],\,
\widehat{\delta C}^{\mathrm{LoG}}_{\Lambda_k}
\big\rangle\!\big\rangle_{C_{\ast,\Lambda_k}}.
\label{eq:witness-descendant-comparison-thm-seam}
\end{equation}
If quotient descent also holds at \(\Lambda_k\), then the common value equals the descended
dominant scalar on the transported descendant affine class:
\begin{equation}
\overline{\rho}^{\mathrm{dom}}_{\Lambda_k}\big([\![\Xi^{(1)}_{\Lambda_k}]\!]\big)
=
\rho^{\mathrm{dom}}_{\Lambda_k}\!\big(J^{f,y}_{\Lambda_k}\big)
=
\rho^{\mathrm{dom}}_{\Lambda_k}\!\big(F_{\tau^{f,y}_{\Lambda_k},\Lambda_k}\big).
\label{eq:witness-descendant-comparison-thm-descended}
\end{equation}
\end{theorem}

\begin{proof}
Set
\[
A^{f,y}_{\Lambda_k}
:=
J^{f,y}_{\Lambda_k}
-
F_{\tau^{f,y}_{\Lambda_k},\Lambda_k}.
\]
By source compatibility,
\(
A^{f,y}_{\Lambda_k}\in\mathcal A^-_{\Lambda_k}
\).

Because dominant-line realization holds and the three vanishing statements of
\cref{prop:dominant-quotient-generators} are assumed, that proposition gives
\[
\mathcal A^-_{\Lambda_k}\subset\ker\rho^{\mathrm{dom}}_{\Lambda_k}.
\]
Hence
\[
\rho^{\mathrm{dom}}_{\Lambda_k}\!\big(J^{f,y}_{\Lambda_k}\big)
=
\rho^{\mathrm{dom}}_{\Lambda_k}\!\big(
F_{\tau^{f,y}_{\Lambda_k},\Lambda_k}
+
A^{f,y}_{\Lambda_k}
\big)
=
\rho^{\mathrm{dom}}_{\Lambda_k}\!\big(F_{\tau^{f,y}_{\Lambda_k},\Lambda_k}\big),
\]
proving \eqref{eq:witness-descendant-comparison-thm-normalized}.

Since dominant-line realization implies
\(
\delta C^{\mathrm{LoG}}_{\Lambda_k}\neq 0
\)
and
\(
\rho^{\mathrm{dom}}_{\Lambda_k}
=
\kappa_{\Lambda_k}^{-1/2}\widetilde{\rho}^{\mathrm{dom}}_{\Lambda_k},
\)
the raw pairing equality \eqref{eq:witness-descendant-comparison-thm-raw} is just the same identity
multiplied by the positive factor \(\kappa_{\Lambda_k}^{1/2}\).

Equation \eqref{eq:witness-descendant-comparison-thm-seam} is exactly the normalized pairing formula
of \cref{prop:odd-source-pairing} applied to the representable source
\(
\tau^{f,y}_{\Lambda_k}
\).

If quotient descent also holds, then \cref{cor:dominant-quotient-constancy} gives
\[
\rho^{\mathrm{dom}}_{\Lambda_k}\!\big(J^{f,y}_{\Lambda_k}\big)
=
\overline{\rho}^{\mathrm{dom}}_{\Lambda_k}\big([\![\Xi^{(1)}_{\Lambda_k}]\!]\big)
\]
for every representative
\(
J^{f,y}_{\Lambda_k}\in[\Xi^{(1)}_{\Lambda_k}],
\)
which yields \eqref{eq:witness-descendant-comparison-thm-descended} after combining with
\eqref{eq:witness-descendant-comparison-thm-normalized}.
\end{proof}

\begin{theorem}[Fixed-cutoff seam witness criterion for descendant overlap]
\label{thm:fixed-cutoff-seam-witness}
Fix an accepted tick \(k\).
Assume:
\begin{enumerate}[leftmargin=2em]
\item a detector-factorized witness-indexed seam realization at \(\Lambda_k\);
\item a witness \(y\) with
  \[
  \overline{\mathrm{obs}}(y)\neq 0;
  \]
\item the seam coefficient is nonzero,
  \[
  \widetilde c_{\Lambda_k}\neq 0;
  \]
\item for the chosen filler tag \(f\), the source \(\tau^{f,y}_{\Lambda_k}\) is representable on the
  characterization layer in the sense of \cref{def:odd-source-representability};
\item there exists a witness-indexed source-compatible descendant realization in the sense of
  \cref{def:source-compatible-witness-realization}.
\end{enumerate}
Then:
\begin{enumerate}[leftmargin=2em]
\item the compressed welded LoG line is transversely physical at \(\Lambda_k\), hence dominant-line
  realization holds at that cutoff;
\item the chosen 2T representative satisfies
  \[
  \rho^{\mathrm{dom}}_{\Lambda_k}\!\big(J^{f,y}_{\Lambda_k}\big)\neq 0;
  \]
\item therefore the transported descendant class has nonzero overlap with the canonically selected
  compressed dominant odd line at \(\Lambda_k\);
\item if quotient descent also holds at \(\Lambda_k\), then
  \cref{thm:descendant-overlap} applies and the transported descendant class is dominant-active with
  well-defined sign \(\varepsilon_{\Lambda_k}\).
\end{enumerate}
\end{theorem}

\begin{proof}
By \cref{thm:seam-faithfulness-scalar},
\[
\widetilde s_{\Lambda_k}(f,y)
=
\sigma_{\mathrm{fill}}(f)\,
\widetilde c_{\Lambda_k}\,
\overline{\mathrm{obs}}(y)
\neq 0.
\]
The same theorem therefore gives item (1): the compressed line is transversely physical, hence
dominant-line realization holds at \(\Lambda_k\).

By \cref{prop:odd-source-pairing},
\[
\big\langle\!\big\langle
\Pi_{\mathrm{phys},\Lambda_k}\mathcal R_{\Lambda_k}\!\big(F_{\tau^{f,y}_{\Lambda_k},\Lambda_k}\big),\,
\delta C^{\mathrm{LoG}}_{\Lambda_k}
\big\rangle\!\big\rangle_{C_{\ast,\Lambda_k}}
=
\widetilde s_{\Lambda_k}(f,y)
\neq 0.
\]
Applying the comparison identity \eqref{eq:witness-descendant-comparison} now gives
\[
\big\langle\!\big\langle
\Pi_{\mathrm{phys},\Lambda_k}\mathcal R_{\Lambda_k}\!\big(J^{f,y}_{\Lambda_k}\big),\,
\delta C^{\mathrm{LoG}}_{\Lambda_k}
\big\rangle\!\big\rangle_{C_{\ast,\Lambda_k}}
\neq 0.
\]
Because dominant-line realization has already been established, \(\delta C^{\mathrm{LoG}}_{\Lambda_k}\)
is nonzero and \(\widehat{\delta C}^{\mathrm{LoG}}_{\Lambda_k}\) is defined.
By \cref{prop:descendant-overlap-witness}, this is equivalent to
\(
\rho^{\mathrm{dom}}_{\Lambda_k}(J^{f,y}_{\Lambda_k})\neq 0
\),
proving item (2).

Item (3) is exactly the nonzero-overlap condition from
\cref{prop:descendant-overlap-equivalent}.
If quotient descent is also available, apply \cref{thm:descendant-overlap} to obtain item (4).
\end{proof}

\subsection*{Smooth-taper branch continuity and witness transport}
\label{subsec:smooth-taper-branch-continuity}

The previous fixed-cutoff theorem isolates the nonzero witness needed for activity.
To transport that witness along an accepted branch one needs a continuity statement for the welded
response.
At theorem level, the clean analytic choice is the smooth receiver taper.
The hard projector may still be used, but only on branch segments where its rank is constant.

\begin{lemma}[Smooth-taper continuity of the receiver band projector]
\label{lem:smooth-taper-projector}
Let \(I\) be a connected accepted-branch segment.
Suppose the receiver band is implemented by a smooth taper
\[
P^{\mathrm{(rec)}}_{\mathrm{band}}(\Lambda)
=
\sum_{\ell,m} w_\ell(\Lambda)\,|Y_{\ell m}\rangle\langle Y_{\ell m}|,
\qquad
0\le w_\ell(\Lambda)\le 1,
\]
on the receiver one-particle space, where
\[
\sup_{\ell\ge 0}|w_\ell(\Lambda')-w_\ell(\Lambda)|
\longrightarrow 0
\qquad
\text{as }\Lambda'\to\Lambda
\]
for every \(\Lambda\in I\).
Then
\[
\Lambda\longmapsto P^{\mathrm{(rec)}}_{\mathrm{band}}(\Lambda)
\]
is continuous on \(I\) in operator norm.
If instead the hard projector
\(
\sum_{\ell\le \ell_\star(\Lambda)}|Y_{\ell m}\rangle\langle Y_{\ell m}|
\)
is used, the same conclusion holds on every connected subsegment of \(I\) on which
\(\ell_\star(\Lambda)\) is constant.
\end{lemma}

\begin{proof}
In the spherical-harmonic basis \(\{Y_{\ell m}\}\), the difference
\[
P^{\mathrm{(rec)}}_{\mathrm{band}}(\Lambda')
-
P^{\mathrm{(rec)}}_{\mathrm{band}}(\Lambda)
\]
is diagonal with eigenvalues \(w_\ell(\Lambda')-w_\ell(\Lambda)\).
Therefore
\[
\big\|
P^{\mathrm{(rec)}}_{\mathrm{band}}(\Lambda')
-
P^{\mathrm{(rec)}}_{\mathrm{band}}(\Lambda)
\big\|_{\mathrm{op}}
=
\sup_{\ell\ge 0}|w_\ell(\Lambda')-w_\ell(\Lambda)|.
\]
The assumed uniform continuity of the taper coefficients is exactly the statement that this operator
norm tends to zero as \(\Lambda'\to\Lambda\).

For the hard projector, the coefficients take only the values \(0\) and \(1\), so discontinuities can
occur only when \(\ell_\star(\Lambda)\) jumps.
On a subsegment where \(\ell_\star(\Lambda)\) is constant, the projector is constant and therefore
continuous.
\end{proof}

\begin{theorem}[Smooth-taper branch continuity of the welded physical response]
\label{thm:smooth-taper-welded-response}
Let \(I\) be a connected accepted-branch segment.
Write the welded source-to-covariance map in the UOM factorized form
\begin{equation}
\mathsf T_\Lambda
=
\Pi_{L(\Lambda)}M_{\varepsilon,\Lambda}
\circ
\mathsf K_{{\rm comp},\Lambda}
\circ
\mathsf G_{{\rm ret},\Lambda}
\circ
\mathsf W_{{\rm src},\Lambda}.
\label{eq:factorized-welded-map}
\end{equation}
Assume:
\begin{enumerate}[leftmargin=2em]
\item the receiver band is implemented by a smooth taper satisfying
  \cref{lem:smooth-taper-projector};
\item the operator families
  \[
  \Lambda\longmapsto
  \Pi_{L(\Lambda)}M_{\varepsilon,\Lambda},\qquad
  \Lambda\longmapsto
  \mathsf K_{{\rm comp},\Lambda},\qquad
  \Lambda\longmapsto
  \mathsf G_{{\rm ret},\Lambda},\qquad
  \Lambda\longmapsto
  \mathsf W_{{\rm src},\Lambda},\qquad
  \Lambda\longmapsto
  \Pi_{\mathrm{phys},\Lambda}
  \]
  are continuous on \(I\) in operator norm on the corresponding band-limited source and covariance
  tangent spaces;
\item the source family
  \[
  \Lambda\longmapsto \tau_\Lambda
  \]
  is continuous on \(I\) in the welded source norm.
\end{enumerate}
Then the physical covariance response
\begin{equation}
X_\Lambda[\tau_\Lambda]
:=
\Pi_{\mathrm{phys},\Lambda}\mathsf T_\Lambda[\tau_\Lambda]
\label{eq:XLambda-tau-def}
\end{equation}
depends continuously on \(\Lambda\in I\) in the Fisher norm.
If the hard receiver projector is used instead, the same conclusion holds on each connected
subsegment where \(\ell_\star(\Lambda)\) is constant.
\end{theorem}

\begin{proof}
By \cref{lem:smooth-taper-projector}, the first factor
\(
\Pi_{L(\Lambda)}M_{\varepsilon,\Lambda}
\)
is operator-norm continuous along the branch.
By assumption, the remaining three UOM factors in \eqref{eq:factorized-welded-map} and the physical
projector \(\Pi_{\mathrm{phys},\Lambda}\) are also operator-norm continuous.
Hence the composed operator family
\[
\Lambda\longmapsto
\Pi_{\mathrm{phys},\Lambda}\mathsf T_\Lambda
\]
is continuous in operator norm on \(I\).
Applying this continuous family to the continuous source vector \(\tau_\Lambda\) gives continuity of
\(
X_\Lambda[\tau_\Lambda]
\)
in the Fisher norm.

For the hard projector, the same argument applies on every connected subsegment where
\(\ell_\star(\Lambda)\) is constant, because there the first factor is constant.
\end{proof}

\begin{corollary}[Branchwise continuity of the seam coefficient]
\label{cor:branchwise-seam-coeff}
Let \(I\) be a connected accepted-branch segment and assume:
\begin{enumerate}[leftmargin=2em]
\item the hypotheses of \cref{thm:dominant-line-persistence} on \(I\);
\item a detector-factorized branchwise seam realization on \(I\) with continuous unit source
  \[
  \Lambda\longmapsto \tau^{\mathrm{unit}}_{\Lambda};
  \]
\item the hypotheses of \cref{thm:smooth-taper-welded-response} for that unit source.
\end{enumerate}
Define the branchwise normalized seam coefficient by
\begin{equation}
c_\Lambda
:=
\big\langle\!\big\langle
\Pi_{\mathrm{phys},\Lambda}\mathsf T_\Lambda\!\big[\tau^{\mathrm{unit}}_{\Lambda}\big],\,
\widehat{\delta C}^{\mathrm{LoG}}_{\Lambda}
\big\rangle\!\big\rangle_{C_{\ast,\Lambda}}.
\label{eq:branchwise-seam-coeff}
\end{equation}
Then \(c_\Lambda\) is continuous on \(I\).
Consequently, for every fixed filler tag \(f\) and witness \(y\), the normalized seam scalar
\[
s_\Lambda(f,y)
=
\sigma_{\mathrm{fill}}(f)\,
c_\Lambda\,
\overline{\mathrm{obs}}(y)
\]
is continuous on \(I\).
\end{corollary}

\begin{proof}
By \cref{thm:smooth-taper-welded-response}, the response vector
\[
\Lambda\longmapsto
\Pi_{\mathrm{phys},\Lambda}\mathsf T_\Lambda\!\big[\tau^{\mathrm{unit}}_{\Lambda}\big]
\]
is continuous in the Fisher norm.
By \cref{thm:dominant-line-persistence}, the normalized dominant vector
\(
\widehat{\delta C}^{\mathrm{LoG}}_{\Lambda}
\)
is also continuous on \(I\).
Pairing these two continuous vectors in the Fisher metric gives continuity of \(c_\Lambda\).
The final statement is then immediate from
\eqref{eq:normalized-seam-factorization}.
\end{proof}

\begin{theorem}[Branchwise seam witness transport]
\label{thm:branchwise-seam-witness}
Let \(I\) be a connected accepted-branch segment.
Assume:
\begin{enumerate}[leftmargin=2em]
\item the hypotheses of \cref{cor:branchwise-seam-coeff};
\item quotient descent holds for every \(\Lambda\in I\);
\item there exists a fixed LS witness \(y\) with
  \[
  \overline{\mathrm{obs}}(y)\neq 0;
  \]
\item for some fixed filler tag \(f\), there is a branchwise family of representatives
  \[
  J^{f,y}_{\Lambda}\in[\Xi^{(1)}_{\Lambda}]
  \qquad
  (\Lambda\in I)
  \]
  such that, for every \(\Lambda\in I\), the corresponding seam source
  \(\tau^{f,y}_{\Lambda}\) is representable on the characterization layer and the 2T-side
  witness-to-descendant comparison identity
  \eqref{eq:witness-descendant-comparison-normalized} holds.
\end{enumerate}
Then:
\begin{enumerate}[leftmargin=2em]
\item the branchwise dominant-response scalar of the chosen descendant representatives is
  \begin{equation}
  \rho^{\mathrm{dom}}_{\Lambda}\!\big(J^{f,y}_{\Lambda}\big)
  =
  \sigma_{\mathrm{fill}}(f)\,
  c_\Lambda\,
  \overline{\mathrm{obs}}(y);
  \label{eq:branchwise-seam-descendant-scalar}
  \end{equation}
\item the scalar function
  \(
  \Lambda\mapsto \rho^{\mathrm{dom}}_{\Lambda}(J^{f,y}_{\Lambda})
  \)
  is continuous on \(I\);
\item on every connected subsegment \(I'\subset I\) on which \(c_\Lambda\neq 0\), the sign
  \[
  \operatorname{sgn}\!\big(
  \rho^{\mathrm{dom}}_{\Lambda}(J^{f,y}_{\Lambda})
  \big)
  \]
  is constant;
\item because quotient descent holds, the same constant sign is the transported orientation sign of
  the descendant affine class on \(I'\).
\end{enumerate}
\end{theorem}

\begin{proof}
For each \(\Lambda\in I\), representability of \(\tau^{f,y}_{\Lambda}\) and
\cref{prop:odd-source-pairing} give
\[
\rho^{\mathrm{dom}}_{\Lambda}\!\big(F_{\tau^{f,y}_{\Lambda},\Lambda}\big)
=
\big\langle\!\big\langle
\Pi_{\mathrm{phys},\Lambda}\mathsf T_\Lambda\!\big[\tau^{f,y}_{\Lambda}\big],\,
\widehat{\delta C}^{\mathrm{LoG}}_{\Lambda}
\big\rangle\!\big\rangle_{C_{\ast,\Lambda}}
=
s_\Lambda(f,y).
\]
Applying the branchwise comparison identity
\eqref{eq:witness-descendant-comparison-normalized} then yields
\[
\rho^{\mathrm{dom}}_{\Lambda}\!\big(J^{f,y}_{\Lambda}\big)
=
s_\Lambda(f,y).
\]
Now use \cref{cor:branchwise-seam-coeff} to substitute
\(
s_\Lambda(f,y)
=
\sigma_{\mathrm{fill}}(f)c_\Lambda\overline{\mathrm{obs}}(y)
\),
which proves item (1).
Item (2) follows because \(c_\Lambda\) is continuous and
\(\overline{\mathrm{obs}}(y)\) is a fixed nonzero scalar.
On any connected subsegment where \(c_\Lambda\neq 0\), the right-hand side of
\eqref{eq:branchwise-seam-descendant-scalar} is a continuous nonvanishing real-valued function, so
its sign is constant.
This proves item (3).

Finally, quotient descent implies that \(\rho^{\mathrm{dom}}_{\Lambda}\) is constant on
\([\Xi^{(1)}_{\Lambda}]\), so the sign of one representative is the sign of the whole descendant
class.
This gives item (4).
\end{proof}

\begin{corollary}[Seam package discharge of the minimal scalar transport input]
\label{cor:branchwise-seam-minimal-input}
Let \(I\) be a connected accepted-branch segment and assume the hypotheses of
\cref{thm:branchwise-seam-witness}.
Then the branchwise family
\[
J^{f,y}_{\Lambda}\in[\Xi^{(1)}_{\Lambda}]
\qquad
(\Lambda\in I)
\]
satisfies
\[
\Lambda\longmapsto
\rho^{\mathrm{dom}}_{\Lambda}\!\big(J^{f,y}_{\Lambda}\big)
=
\sigma_{\mathrm{fill}}(f)\,
c_\Lambda\,
\overline{\mathrm{obs}}(y)
\]
continuously on \(I\).
On every connected subsegment \(I'\subset I\) such that
\[
c_\Lambda\neq 0
\qquad
\forall \Lambda\in I',
\]
assumption \textup{(3)} of \cref{thm:active-response-transport-minimal} holds in its minimal
scalar form with
\[
J^{\mathrm{desc}}_{\Lambda}:=J^{f,y}_{\Lambda}
\qquad
(\Lambda\in I').
\]
If, in addition, the three generator-wise vanishing statements of
\cref{prop:dominant-quotient-generators} hold for every \(\Lambda\in I'\), then
\cref{thm:active-response-transport-minimal} applies on \(I'\).
\end{corollary}

\begin{proof}
Continuity of the displayed scalar function is exactly item (2) of
\cref{thm:branchwise-seam-witness}, together with the explicit formula of item (1).
On \(I'\), the hypotheses
\(
c_\Lambda\neq 0
\)
and
\(
\overline{\mathrm{obs}}(y)\neq 0
\)
imply
\[
\rho^{\mathrm{dom}}_{\Lambda}\!\big(J^{f,y}_{\Lambda}\big)\neq 0
\qquad
\forall \Lambda\in I'.
\]
Together with the continuity statement already proved above, this is exactly assumption
\textup{(3)} of \cref{thm:active-response-transport-minimal} in its minimal scalar form.
The hypotheses of \cref{thm:branchwise-seam-witness} already include those of
\cref{cor:branchwise-seam-coeff}, and hence the dominant-line persistence input of
\cref{thm:active-response-transport-minimal}.
If the three generator-wise vanishing statements also hold on \(I'\), then every hypothesis of
\cref{thm:active-response-transport-minimal} is in force there.
\end{proof}

\begin{theorem}[Finite-cutoff Active Response Transport Theorem]
\label{thm:active-response-transport-minimal}
Let \(I\) be a connected accepted-branch segment.
For each \(\Lambda\in I\), let
\[
[\Xi^{(1)}_{\Lambda}]
:=
J_{0,\Lambda}+\mathcal A^-_{\Lambda}
\subset E^-_{\Lambda}
\]
be the parity-transported odd descendant affine class modulo the ambiguity space of
\cref{def:transport-ambiguity-subspace}.
Assume:
\begin{enumerate}[leftmargin=2em]
\item the hypotheses of \cref{thm:dominant-line-persistence} hold on \(I\);
\item for every \(\Lambda\in I\), the three vanishing statements of
  \cref{prop:dominant-quotient-generators} hold;
\item there exists a branchwise family of representatives
  \[
  J^{\mathrm{desc}}_{\Lambda}\in[\Xi^{(1)}_{\Lambda}]
  \qquad
  (\Lambda\in I)
  \]
  such that
  \[
  \rho^{\mathrm{dom}}_{\Lambda}(J^{\mathrm{desc}}_{\Lambda})
  \neq 0
  \qquad
  \forall \Lambda\in I,
  \]
  and such that the dominant-response scalar
  \[
  \Lambda\longmapsto
  \rho^{\mathrm{dom}}_{\Lambda}(J^{\mathrm{desc}}_{\Lambda})
  \]
  is continuous on \(I\).
\end{enumerate}
Then, for every \(\Lambda\in I\):
\begin{enumerate}[leftmargin=2em]
\item the compressed dominant covariance tangent is nonzero,
  \[
  \delta C^{\mathrm{LoG}}_{\Lambda}\neq 0,
  \]
  so the dominant line \(L^{\mathrm{dom}}_{\Lambda}\) of
  \cref{def:dominant-response-functional} is canonically defined;
\item the dominant-response coefficient \(\rho^{\mathrm{dom}}_{\Lambda}\) annihilates the odd
  ambiguity space,
  \[
  \rho^{\mathrm{dom}}_{\Lambda}(A)=0
  \qquad
  \forall A\in\mathcal A^-_{\Lambda};
  \]
\item the transported descendant affine class carries a unique nonzero dominant-response value
  \[
  \rho^{\mathrm{dom}}_{\Lambda}(J)
  =
  \overline{\rho}^{\mathrm{dom}}_{\Lambda}\big([\![\Xi^{(1)}_{\Lambda}]\!]\big)
  \neq 0
  \qquad
  \forall J\in[\Xi^{(1)}_{\Lambda}],
  \]
  and hence a well-defined sign
  \[
  \varepsilon_{\Lambda}
  :=
  \operatorname{sgn}\!\Big(
  \overline{\rho}^{\mathrm{dom}}_{\Lambda}\big([\![\Xi^{(1)}_{\Lambda}]\!]\big)
  \Big)
  \in\{\pm1\};
  \]
\item the normalized orientation functional
  \begin{equation}
  \chi_{\Lambda}
  :=
  \nu_{\Lambda}^{-1}\rho^{\mathrm{dom}}_{\Lambda}
  \label{eq:chi-from-rho-dom}
  \end{equation}
  is well defined, annihilates \(\mathcal A^-_{\Lambda}\), and satisfies
  \[
  \chi_{\Lambda}(J)=\varepsilon_{\Lambda}
  \qquad
  \forall J\in[\Xi^{(1)}_{\Lambda}];
  \]
\item along \(I\), the sign \(\varepsilon_{\Lambda}\) is locally constant once an initial
  orientation is fixed.
\end{enumerate}
\end{theorem}

\begin{proof}
Item (1) is exactly the conclusion of \cref{thm:dominant-line-persistence}.

By assumption, the three generator-wise vanishing statements of
\cref{prop:dominant-quotient-generators} hold for every \(\Lambda\in I\).
Therefore \cref{prop:dominant-quotient-generators,thm:dominant-quotient-descent} give item (2)
and the existence of the descended functional
\(
\overline{\rho}^{\mathrm{dom}}_{\Lambda}
\)
at each cutoff.

Applying \cref{prop:descendant-overlap-equivalent,thm:descendant-overlap} pointwise to the
nonvanishing scalar assumption
\(
\rho^{\mathrm{dom}}_{\Lambda}(J^{\mathrm{desc}}_{\Lambda})\neq 0
\)
now yields
item (3): the entire descendant affine class has one common nonzero dominant-response value, and
its sign \(\varepsilon_{\Lambda}\) is well defined.

Since \(\nu_{\Lambda}>0\) by item (3), the normalized functional \(\chi_{\Lambda}\) of
\eqref{eq:chi-from-rho-dom} is well defined.
Item (2) implies that \(\chi_{\Lambda}\) annihilates \(\mathcal A^-_{\Lambda}\), and item (3)
gives
\[
\chi_{\Lambda}(J)
=
\frac{\rho^{\mathrm{dom}}_{\Lambda}(J)}{
\left|
\overline{\rho}^{\mathrm{dom}}_{\Lambda}\big([\![\Xi^{(1)}_{\Lambda}]\!]\big)
\right|
}
=
\varepsilon_{\Lambda}
\qquad
\forall J\in[\Xi^{(1)}_{\Lambda}],
\]
proving item (4).

Finally, assumption (3) provides a continuous nonvanishing real-valued branchwise descendant scalar,
while items (1)--(3) above identify that scalar with the transported orientation sign of the full
descendant affine class.
Since
\(
\Lambda\mapsto \rho^{\mathrm{dom}}_{\Lambda}(J^{\mathrm{desc}}_{\Lambda})
\)
is continuous and nonzero on the connected set \(I\), its sign is constant, hence in particular
locally constant.
Using item (3), this is exactly the transported orientation sign of the full descendant affine
class, proving item (5).
\end{proof}

\begin{proposition}[Ambiguity-triviality on the compressed welded image]
\label{prop:reduced-welded-ambiguity-trivial}
Assume the hypotheses of \cref{thm:active-response-transport-minimal} at an accepted tick \(k\).
Define the reduced welded image of an odd generator \(F\in E^-_{\Lambda_k}\) by
\[
\mathsf W^{\mathrm{red}}_{\Lambda_k}(F)
:=
\rho^{\mathrm{dom}}_{\Lambda_k}(F)\,P_{-,\Lambda_k}^{\mathrm{LoG}},
\]
where \(P_{-,\Lambda_k}^{\mathrm{LoG}}\) is the canonical compressed odd projector of
\cref{cor:compressed-export-selection}.
Then
\[
\mathsf W^{\mathrm{red}}_{\Lambda_k}(F+A)
=
\mathsf W^{\mathrm{red}}_{\Lambda_k}(F)
\qquad
\forall F\in E^-_{\Lambda_k},\ \forall A\in\mathcal A^-_{\Lambda_k}.
\]
Equivalently, \(\mathsf W^{\mathrm{red}}_{\Lambda_k}\) factors through the ambiguity quotient
\(E^-_{\Lambda_k}/\mathcal A^-_{\Lambda_k}\).
In particular, if \(J_1,J_2\in[\Xi^{(1)}_{\Lambda_k}]\), then
\[
\mathsf W^{\mathrm{red}}_{\Lambda_k}(J_1)
=
\mathsf W^{\mathrm{red}}_{\Lambda_k}(J_2).
\]
\end{proposition}

\begin{proof}
Item (2) of \cref{thm:active-response-transport-minimal} gives
\[
\rho^{\mathrm{dom}}_{\Lambda_k}(A)=0
\qquad
\forall A\in\mathcal A^-_{\Lambda_k}.
\]
Therefore, for every \(F\in E^-_{\Lambda_k}\) and \(A\in\mathcal A^-_{\Lambda_k}\),
\[
\mathsf W^{\mathrm{red}}_{\Lambda_k}(F+A)
=
\rho^{\mathrm{dom}}_{\Lambda_k}(F+A)\,P_{-,\Lambda_k}^{\mathrm{LoG}}
=
\big(\rho^{\mathrm{dom}}_{\Lambda_k}(F)+\rho^{\mathrm{dom}}_{\Lambda_k}(A)\big)
P_{-,\Lambda_k}^{\mathrm{LoG}}
=
\rho^{\mathrm{dom}}_{\Lambda_k}(F)\,P_{-,\Lambda_k}^{\mathrm{LoG}}.
\]
Since \cref{cor:compressed-export-selection} fixes \(P_{-,\Lambda_k}^{\mathrm{LoG}}\) canonically on
the compressed, co-centered, band-locked welded channel, the right-hand side depends only on the
ambiguity class of \(F\).
If \(J_2-J_1\in\mathcal A^-_{\Lambda_k}\), apply the same identity with \(F=J_1\) and
\(A=J_2-J_1\).
\end{proof}

\begin{proposition}[Construction of an admissible calibration for the minimal transported interface]
\label{prop:minimal-transport-calibration}
Fix an accepted tick \(k\). Suppose a weak transported-interface package is given:
\[
\mathcal A^-_{\Lambda_k}\subset E^-_{\Lambda_k},
\qquad
[\Xi^{(1)}_{\Lambda_k}]=J_{0,\Lambda_k}+\mathcal A^-_{\Lambda_k},
\]
together with a continuous nonzero linear functional
\[
\chi_{\Lambda_k}:E^-_{\Lambda_k}\to\mathbb R
\]
and a sign
\(
\varepsilon_{\Lambda_k}\in\{\pm1\}
\)
such that
\[
\mathcal A^-_{\Lambda_k}\subset\ker\chi_{\Lambda_k},
\qquad
\chi_{\Lambda_k}(J)=\varepsilon_{\Lambda_k}
\quad
\forall J\in[\Xi^{(1)}_{\Lambda_k}].
\]
Assume in addition that \(\ker\chi_{\Lambda_k}\) admits a continuous symmetric coercive bilinear form
whose induced norm Hilbertizes it.
Choose any vector \(J^\sharp_{\Lambda_k}\in E^-_{\Lambda_k}\) satisfying
\[
\chi_{\Lambda_k}(J^\sharp_{\Lambda_k})=1,
\]
any scale \(\beta_{\Lambda_k}>0\), and any continuous symmetric coercive bilinear form
\[
\langle\cdot,\cdot\rangle^{\mathrm{aux},\chi}_{\Lambda_k}
:
\ker\chi_{\Lambda_k}\times\ker\chi_{\Lambda_k}\to\mathbb R
\]
whose induced norm Hilbertizes \(\ker\chi_{\Lambda_k}\).
Define
\begin{equation}
\mathsf B_{\Lambda_k}(F,G)
:=
\beta_{\Lambda_k}\,\chi_{\Lambda_k}(F)\chi_{\Lambda_k}(G)
+
\left\langle
F-\chi_{\Lambda_k}(F)J^\sharp_{\Lambda_k},
\,
G-\chi_{\Lambda_k}(G)J^\sharp_{\Lambda_k}
\right\rangle^{\mathrm{aux},\chi}_{\Lambda_k}.
\label{eq:B-minimal-transport}
\end{equation}
Then \(\mathsf B_{\Lambda_k}\) is a continuous symmetric positive definite bilinear form on
\(E^-_{\Lambda_k}\) whose induced norm Hilbertizes \(E^-_{\Lambda_k}\), and
\[
\mathsf B_{\Lambda_k}(J^\sharp_{\Lambda_k},K)=0
\qquad
\forall K\in\ker\chi_{\Lambda_k}.
\]
Hence \(\mathsf B_{\Lambda_k}(J^\sharp_{\Lambda_k},K)=0\) in particular for every
\(K\in\mathcal A^-_{\Lambda_k}\), and
\[
\mathsf B_{\Lambda_k}(J^\sharp_{\Lambda_k},J^\sharp_{\Lambda_k})
=
\beta_{\Lambda_k}.
\]
Moreover, for any representative \(J^\Xi_{\Lambda_k}\in[\Xi^{(1)}_{\Lambda_k}]\), the choice
\[
J^\sharp_{\Lambda_k}:=\varepsilon_{\Lambda_k}J^\Xi_{\Lambda_k}
\]
is admissible.
\end{proposition}

\begin{proof}
Because \(\chi_{\Lambda_k}(J^\sharp_{\Lambda_k})=1\), the linear map
\[
P^\chi_{\Lambda_k}(F)
:=
F-\chi_{\Lambda_k}(F)J^\sharp_{\Lambda_k}
\]
takes values in \(\ker\chi_{\Lambda_k}\):
\[
\chi_{\Lambda_k}\!\big(P^\chi_{\Lambda_k}(F)\big)
=
\chi_{\Lambda_k}(F)-\chi_{\Lambda_k}(F)\chi_{\Lambda_k}(J^\sharp_{\Lambda_k})
=
0.
\]
Thus every \(F\in E^-_{\Lambda_k}\) has the decomposition
\[
F=\chi_{\Lambda_k}(F)J^\sharp_{\Lambda_k}+P^\chi_{\Lambda_k}(F),
\qquad
P^\chi_{\Lambda_k}(F)\in\ker\chi_{\Lambda_k}.
\]
This decomposition is unique: if \(F=aJ^\sharp_{\Lambda_k}+K\) with \(K\in\ker\chi_{\Lambda_k}\),
applying \(\chi_{\Lambda_k}\) gives \(a=\chi_{\Lambda_k}(F)\).

Equation \eqref{eq:B-minimal-transport} is therefore well defined, continuous, and symmetric.
If \(\mathsf B_{\Lambda_k}(F,F)=0\), then
\[
\beta_{\Lambda_k}\chi_{\Lambda_k}(F)^2=0
\]
and
\[
\left\langle
P^\chi_{\Lambda_k}(F),P^\chi_{\Lambda_k}(F)
\right\rangle^{\mathrm{aux},\chi}_{\Lambda_k}
=
0.
\]
Since \(\beta_{\Lambda_k}>0\), we get \(\chi_{\Lambda_k}(F)=0\), hence
\(P^\chi_{\Lambda_k}(F)=F\). Coercivity on \(\ker\chi_{\Lambda_k}\) then implies \(F=0\), proving
positive definiteness.

The map
\[
F\longmapsto
\big(
\chi_{\Lambda_k}(F),P^\chi_{\Lambda_k}(F)
\big)
\]
is a continuous linear isomorphism from \(E^-_{\Lambda_k}\) onto
\(
\mathbb R\oplus\ker\chi_{\Lambda_k}
\),
with inverse
\[
(a,K)\longmapsto aJ^\sharp_{\Lambda_k}+K.
\]
Under this isomorphism, \(\mathsf B_{\Lambda_k}\) becomes the product Hilbert form
\[
(a,K),(b,L)\longmapsto
\beta_{\Lambda_k}ab+\langle K,L\rangle^{\mathrm{aux},\chi}_{\Lambda_k},
\]
so its induced norm Hilbertizes \(E^-_{\Lambda_k}\).

Finally, if \(K\in\ker\chi_{\Lambda_k}\), then \(P^\chi_{\Lambda_k}(J^\sharp_{\Lambda_k})=0\) and
\(\chi_{\Lambda_k}(K)=0\), hence
\[
\mathsf B_{\Lambda_k}(J^\sharp_{\Lambda_k},K)=0.
\]
Setting \(K=J^\sharp_{\Lambda_k}\) gives
\(
\mathsf B_{\Lambda_k}(J^\sharp_{\Lambda_k},J^\sharp_{\Lambda_k})=\beta_{\Lambda_k}
\).
If \(J^\Xi_{\Lambda_k}\in[\Xi^{(1)}_{\Lambda_k}]\), then
\(
\chi_{\Lambda_k}(J^\Xi_{\Lambda_k})=\varepsilon_{\Lambda_k}
\),
so
\(
J^\sharp_{\Lambda_k}:=\varepsilon_{\Lambda_k}J^\Xi_{\Lambda_k}
\)
satisfies
\(
\chi_{\Lambda_k}(J^\sharp_{\Lambda_k})=1
\).
\end{proof}

\begin{remark}[Calibration existence versus stiffness matching]
\label{rem:minimal-transport-calibration-scope}
\Cref{prop:minimal-transport-calibration} is a construction theorem, not an unconditional existence
theorem: it extends a previously chosen Hilbertizing form on \(\ker\chi_{\Lambda_k}\) to an
admissible bilinear form on the whole weak transported interface.
It does \emph{not} identify the scale
\(\beta_{\Lambda_k}\) with the split/leak stiffness \(\widetilde\alpha(\Lambda_k)\).
That physical identification belongs to a later strong or physical layer.
\end{remark}

\begin{theorem}[Rank-one descendant-sector upgrade]
\label{thm:active-response-transport-strong}
Assume \cref{thm:active-response-transport-minimal}, and fix an accepted tick \(k\) on the branch.
Suppose in addition that on the transported
descendant sector the projected response map
\[
F\longmapsto
\Pi^{\mathrm{dom}}_{\Lambda_k}\Pi_{\mathrm{phys},\Lambda_k}\mathcal R_{\Lambda_k}(F),
\qquad
\Pi^{\mathrm{dom}}_{\Lambda_k}
:=
\widehat{\delta C}^{\mathrm{LoG}}_{\Lambda_k}\otimes
\widehat{\delta C}^{\mathrm{LoG}}_{\Lambda_k},
\]
has kernel exactly \(\mathcal A^-_{\Lambda_k}\).
Then the stronger conclusions follow:
\begin{enumerate}[leftmargin=2em]
\item the descendant quotient is rank one,
  \[
  E^-_{\Lambda_k}/\mathcal A^-_{\Lambda_k}\cong \mathbb R;
  \]
\item the orientation kernel is exactly the ambiguity space,
  \[
  \ker\chi_{\Lambda_k}
  =
  \ker\rho^{\mathrm{dom}}_{\Lambda_k}
  =
  \mathcal A^-_{\Lambda_k},
  \]
  so the descendant sector satisfies the rank-one kernel-identification package with quotient
  generated by any vector \(J^{\mathrm{act}}_{\Lambda_k}\) satisfying
  \(\chi_{\Lambda_k}(J^{\mathrm{act}}_{\Lambda_k})=1\);
  \item after choosing any such \(J^{\mathrm{act}}_{\Lambda_k}\), any positive scale
  \(\beta_{\Lambda_k}\), and any auxiliary coercive Hilbertizing form on
  \(\ker\chi_{\Lambda_k}\), \cref{prop:minimal-transport-calibration} yields an admissible
  calibration form.
  Recovering the stiffness-matched formula of \cref{prop:char-calibration-extension} requires
  additional physical identifications, namely
  \(\chi_{\Lambda_k}=\rho_{\Lambda_k}\), \(\ker\chi_{\Lambda_k}=\mathcal T^{\mathrm{op}}_{\Lambda_k}\),
  agreement of the chosen normalized generator with the split/leak active generator, and
  \(\beta_{\Lambda_k}=\widetilde\alpha(\Lambda_k)\).
\end{enumerate}
\end{theorem}

\begin{proof}
By \cref{thm:active-response-transport-minimal}, the dominant line is defined at \(\Lambda_k\), the
scalar \(\nu_{\Lambda_k}\) is positive, and
\[
\chi_{\Lambda_k}
=
\nu_{\Lambda_k}^{-1}\rho^{\mathrm{dom}}_{\Lambda_k}.
\]
Hence
\(
\ker\chi_{\Lambda_k}=\ker\rho^{\mathrm{dom}}_{\Lambda_k}
\).

For every \(F\in E^-_{\Lambda_k}\), the definition of \(\rho^{\mathrm{dom}}_{\Lambda_k}\) gives
\[
\Pi^{\mathrm{dom}}_{\Lambda_k}\Pi_{\mathrm{phys},\Lambda_k}\mathcal R_{\Lambda_k}(F)
=
\rho^{\mathrm{dom}}_{\Lambda_k}(F)\,
\widehat{\delta C}^{\mathrm{LoG}}_{\Lambda_k}.
\]
Therefore the kernel of the projected response map is exactly
\(
\ker\rho^{\mathrm{dom}}_{\Lambda_k}
\).
By the additional hypothesis, that kernel is \(\mathcal A^-_{\Lambda_k}\), proving item (2).

By \cref{thm:active-response-transport-minimal}, there exists
\(
J^{\mathrm{desc}}_{\Lambda_k}\in[\Xi^{(1)}_{\Lambda_k}]
\)
with
\(
\rho^{\mathrm{dom}}_{\Lambda_k}(J^{\mathrm{desc}}_{\Lambda_k})\neq 0
\),
hence the projected response map is nonzero.
Its image is automatically contained in the one-dimensional line
\(
L^{\mathrm{dom}}_{\Lambda_k}
\),
so the first isomorphism theorem gives
\[
E^-_{\Lambda_k}/\mathcal A^-_{\Lambda_k}
\cong
\operatorname{im}\!\big(
\Pi^{\mathrm{dom}}_{\Lambda_k}\Pi_{\mathrm{phys},\Lambda_k}\mathcal R_{\Lambda_k}
\big)
\cong
\mathbb R,
\]
which is item (1).
Any vector \(J^{\mathrm{act}}_{\Lambda_k}\) with \(\chi_{\Lambda_k}(J^{\mathrm{act}}_{\Lambda_k})=1\)
lies outside \(\ker\chi_{\Lambda_k}=\mathcal A^-_{\Lambda_k}\), so its quotient class generates this
one-dimensional quotient, completing item (2).

Item (3) is now an application of \cref{prop:minimal-transport-calibration} with the weak-interface
data supplied by \cref{thm:active-response-transport-minimal} and the chosen normalized vector
\(J^{\mathrm{act}}_{\Lambda_k}\).
The final sentence identifies when this weak calibration construction coincides with the
stiffness-matched formula of \cref{prop:char-calibration-extension}.
\end{proof}

\ifdefined\UOMTransportFragmentLoaded\else
\end{document}
\fi
%
  }{%
    \ifdefined\UOMTransportFragmentLoaded\else
\documentclass[11pt]{article}
\usepackage[margin=1in]{geometry}
\usepackage{amsmath,amssymb,amsthm,mathtools}
\usepackage{bm}
\usepackage{microtype}
\usepackage{enumitem}
\usepackage{booktabs}
\usepackage{array}
\usepackage{xcolor}
\usepackage{xr-hyper}
\usepackage{hyperref}
\usepackage[nameinlink,noabbrev]{cleveref}

\providecommand{\Fsloc}{\mathcal F_{\rm s.loc}}
\providecommand{\Qt}{Q_t}
\externaldocument[main-]{UOM/main}
\externaldocument[uomdesc-]{UOM/uom_descendant_note}
\externaldocument{UOM/completion_sector_split_model_note}

\hypersetup{
  colorlinks=true,
  linkcolor=blue!60!black,
  citecolor=blue!60!black,
  urlcolor=blue!60!black
}

\theoremstyle{definition}
\newtheorem{definition}{Definition}[section]
\newtheorem{assumption}{Assumption}[section]
\theoremstyle{plain}
\newtheorem{theorem}[definition]{Theorem}
\newtheorem{lemma}[definition]{Lemma}
\newtheorem{proposition}[definition]{Proposition}
\newtheorem{corollary}[definition]{Corollary}
\theoremstyle{remark}
\newtheorem{remark}[definition]{Remark}
\crefname{definition}{definition}{definitions}
\Crefname{definition}{Definition}{Definitions}
\crefname{assumption}{assumption}{assumptions}
\Crefname{assumption}{Assumption}{Assumptions}
\crefname{theorem}{theorem}{theorems}
\Crefname{theorem}{Theorem}{Theorems}
\crefname{lemma}{lemma}{lemmas}
\Crefname{lemma}{Lemma}{Lemmas}
\crefname{proposition}{proposition}{propositions}
\Crefname{proposition}{Proposition}{Propositions}
\crefname{corollary}{corollary}{corollaries}
\Crefname{corollary}{Corollary}{Corollaries}
\crefname{remark}{remark}{remarks}
\Crefname{remark}{Remark}{Remarks}

\newcommand{\cH}{\mathcal H}
\newcommand{\cA}{\mathcal A}
\newcommand{\cD}{\mathcal D}
\newcommand{\cB}{\mathcal B}
\newcommand{\cC}{\mathcal C}
\newcommand{\cL}{\mathcal L}
\newcommand{\cK}{\mathcal K}
\newcommand{\cR}{\mathcal R}
\newcommand{\cG}{\mathcal G}
\newcommand{\Ztwo}{\mathbb Z_2}
\newcommand{\Sys}{\mathrm{Sys}}
\newcommand{\Time}{\mathrm{Time}}
\newcommand{\RG}{\mathrm{RG}}
\newcommand{\Tr}{\mathrm{Tr}}
\newcommand{\diag}{\mathrm{diag}}
\newcommand{\ad}{\mathrm{ad}}
\newcommand{\id}{\mathrm{id}}
\newcommand{\norm}[1]{\left\lVert #1\right\rVert}
\newcommand{\abs}[1]{\left\lvert #1\right\rvert}
\newcommand{\ip}[2]{\left\langle #1,#2\right\rangle}

\externaldocument{UOM/compressed_log_spectral_selection_note}
\externaldocument{UOM/odd_selector_note}
\externaldocument{UOM/canonical_odd_interface_note}
\externaldocument[splitlift-]{UOM/split_controlled_completion_lift_note}
\title{Finite-Cutoff Active Response Transport Theorem}
\author{}
\date{}
\begin{document}
\maketitle
\fi

\section{Finite-Cutoff Active Response Transport Theorem}
\label{sec:active-response-transport-program}

This section assembles a finite-cutoff transport theorem for the odd descendant / Hamiltonization
lane of the 2T Main assumption package Definition~\ref{main-def:uom-main-theorem-assumptions}.
Accordingly, the phrase ``accepted tick'' remains literal UOM cutoff language throughout this note:
the transport theorem is pointwise on the accepted branch at fixed \(\Lambda_k\).

\begin{remark}[Physical odd sector and branchwise ambient-space conventions]
\label{rem:transport-conventions}
Throughout this note, the superscript \({}^-\) refers to physical \(\Gamma\)-odd parity on the
characterization layer, not to cohomological oddness.
Every ambiguity direction appearing below is therefore understood after passage to the physically
odd characterization sector \(E^-_{\Lambda_k}\).

All branchwise continuity statements are made after the standard accepted-branch identifications of
the relevant finite-cutoff source spaces, covariance-tangent spaces, and physical subspaces with
fixed ambient finite-cutoff spaces.
Operator-norm and Fisher-norm continuity are always understood with respect to those identifications.
\end{remark}

\begin{definition}[Odd ambiguity subspace on the characterization layer]
\label{def:transport-ambiguity-subspace}
For each accepted tick \(k\), let
\[
\mathcal A^-_{\Lambda_k}\subset E^-_{\Lambda_k}
\]
denote the closed linear span of the odd characterization-layer projections of the descendant
ambiguity directions
\[
\big(Q_tX+Q_uY\big)^-,\qquad \big(\partial_\mu K^\mu\big)^-,\qquad \Delta_{\rm ct}^{-},
\]
where \(K^\mu\) is boundary-admissible on the characterization layer and
\(\Delta_{\rm ct}^{-}\) is admissible in the sense of \cref{def:ambiguity-ct-char}.
Equivalently, \(\mathcal A^-_{\Lambda_k}\) is the odd ambiguity space already placed inside
\(\mathcal T^{\mathrm{op}}_{\Lambda_k}\) by \cref{cor:ambiguity-char}.
\end{definition}

\begin{definition}[Dominant compressed response line and scalar coefficient]
\label{def:dominant-response-functional}
Fix an accepted tick \(k\). Let
\[
u^{\mathrm{LoG}}_{\Lambda_k}
\]
and
\[
P_{-,\Lambda_k}^{\mathrm{LoG}}
\]
be the canonically selected compressed odd line and projector furnished by
\cref{cor:compressed-export-selection}. Define the physical covariance tangent induced by that line by
\[
\delta C^{\mathrm{LoG}}_{\Lambda_k}
:=
\Pi_{\mathrm{phys},\Lambda_k}\,\mathsf T_{\Lambda_k}[u^{\mathrm{LoG}}_{\Lambda_k}],
\]
and, whenever this tangent is nonzero, write
\[
L^{\mathrm{dom}}_{\Lambda_k}
:=
\mathbb R\,\widehat{\delta C}^{\mathrm{LoG}}_{\Lambda_k}
\]
for its Fisher-normalized line.
The raw dominant-response coefficient is
\begin{equation}
\widetilde{\rho}^{\mathrm{dom}}_{\Lambda_k}(F)
:=
\big\langle\!\big\langle
\Pi_{\mathrm{phys},\Lambda_k}\mathcal R_{\Lambda_k}(F),\,
\delta C^{\mathrm{LoG}}_{\Lambda_k}
\big\rangle\!\big\rangle_{C_{\ast,\Lambda_k}},
\qquad
F\in E^-_{\Lambda_k}.
\label{eq:rho-dom-raw-def}
\end{equation}
Whenever \(\delta C^{\mathrm{LoG}}_{\Lambda_k}\neq 0\), equivalently whenever dominant-line
realization holds at \(\Lambda_k\), define the normalized dominant-response coefficient by
\begin{equation}
\rho^{\mathrm{dom}}_{\Lambda_k}(F)
:=
\kappa_{\Lambda_k}^{-1/2}\,
\widetilde{\rho}^{\mathrm{dom}}_{\Lambda_k}(F),
\qquad
F\in E^-_{\Lambda_k}.
\label{eq:rho-dom-def}
\end{equation}
\end{definition}

\begin{proposition}[Well-posedness and continuity of the dominant-response coefficients]
\label{prop:rho-dom-continuity}
Fix an accepted tick \(k\).
Then the raw dominant-response coefficient
\[
\widetilde{\rho}^{\mathrm{dom}}_{\Lambda_k}:E^-_{\Lambda_k}\to\mathbb R
\]
is continuous and linear.
If dominant-line realization holds at \(\Lambda_k\), then the normalized dominant-response
coefficient
\[
\rho^{\mathrm{dom}}_{\Lambda_k}:E^-_{\Lambda_k}\to\mathbb R
\]
is also continuous and linear.
\end{proposition}

\begin{proof}
By \cref{prop:intrinsic-cov-response}, the intrinsic response map
\(
\mathcal R_{\Lambda_k}:E^-_{\Lambda_k}\to\mathcal T^{\mathrm{cov}}_{\Lambda_k}
\)
is continuous and linear.
Pairing its physical projection against the fixed covariance tangent
\(
\delta C^{\mathrm{LoG}}_{\Lambda_k}
\)
in the Fisher metric therefore gives a continuous linear functional, which is exactly
\(\widetilde{\rho}^{\mathrm{dom}}_{\Lambda_k}\).
If dominant-line realization holds, then \(\kappa_{\Lambda_k}>0\) by
\cref{prop:dominant-line-equivalent}, so \(\rho^{\mathrm{dom}}_{\Lambda_k}\) is obtained from
\(\widetilde{\rho}^{\mathrm{dom}}_{\Lambda_k}\) by multiplication by the fixed scalar
\(\kappa_{\Lambda_k}^{-1/2}\).
\end{proof}

\subsection*{Dominant-line realization}
\label{subsec:dominant-line-realization}

The compressed spectral-selection note already supplies the one-dimensional source line
\(\mathcal H^{\mathrm{LoG}}_{\Lambda_k}\).
The present transport note needs only the nonvanishing and branchwise continuity of its
\emph{physical} covariance push-forward.

\begin{definition}[Compressed-line transversality]
\label{def:compressed-line-transversality}
For an accepted tick \(k\), say that the welded response is \emph{transversely physical on the
compressed line} if
\begin{equation}
\Pi_{\mathrm{phys},\Lambda_k}\,\mathsf T_{\Lambda_k}\!\restriction_{\mathcal H^{\mathrm{LoG}}_{\Lambda_k}}
\neq 0.
\label{eq:compressed-transversality}
\end{equation}
Equivalently, the unique generator \(u^{\mathrm{LoG}}_{\Lambda_k}\) of the compressed odd line has
nonzero physical covariance response.
\end{definition}

\begin{definition}[Dominant-response norm]
\label{def:dominant-response-norm}
Whenever \(\delta C^{\mathrm{LoG}}_{\Lambda_k}\) is defined, set
\begin{equation}
\kappa_{\Lambda_k}
:=
\big\langle\!\big\langle
\delta C^{\mathrm{LoG}}_{\Lambda_k},\,
\delta C^{\mathrm{LoG}}_{\Lambda_k}
\big\rangle\!\big\rangle_{C_{\ast,\Lambda_k}}.
\label{eq:kappa-dom-def}
\end{equation}
This is the Fisher squared norm of the physical covariance tangent induced by the compressed welded
LoG source line.
\end{definition}

\begin{proposition}[Equivalent formulations of dominant-line realization]
\label{prop:dominant-line-equivalent}
For a fixed accepted tick \(k\), the following are equivalent.
\begin{enumerate}[leftmargin=2em]
\item the compressed welded LoG line is transversely physical in the sense of
  \cref{def:compressed-line-transversality};
\item the distinguished physical covariance tangent is nonzero,
  \[
  \delta C^{\mathrm{LoG}}_{\Lambda_k}\neq 0;
  \]
\item the dominant-response norm is strictly positive,
  \[
  \kappa_{\Lambda_k}>0;
  \]
\item there exists a physical covariance tangent
  \(
  B_{\Lambda_k}\in \mathcal T^{\mathrm{phys}}_{\Lambda_k}
  \)
  such that
  \[
  \big\langle\!\big\langle
  \delta C^{\mathrm{LoG}}_{\Lambda_k},\,B_{\Lambda_k}
  \big\rangle\!\big\rangle_{C_{\ast,\Lambda_k}}
  \neq 0;
  \]
\item in any Gaussian completion realization of
  \cref{cor:compressed-export-selection-completion}, the completion-sector covariance tangent
  representing \(\delta C^{\mathrm{LoG}}_{\Lambda_k}\) is nonzero.
\end{enumerate}
\end{proposition}

\begin{proof}
The equivalence of (1) and (2) is immediate from
\eqref{eq:compressed-transversality} and the definition
\[
\delta C^{\mathrm{LoG}}_{\Lambda_k}
=
\Pi_{\mathrm{phys},\Lambda_k}\mathsf T_{\Lambda_k}[u^{\mathrm{LoG}}_{\Lambda_k}],
\]
because \(\mathcal H^{\mathrm{LoG}}_{\Lambda_k}\) is the one-dimensional line generated by
\(u^{\mathrm{LoG}}_{\Lambda_k}\).

The equivalence of (2) and (3) follows from positive definiteness of the Fisher metric on the
covariance tangent space: a vector has positive Fisher norm if and only if it is nonzero.

Statement (3) implies (4) by taking
\(
B_{\Lambda_k}:=\delta C^{\mathrm{LoG}}_{\Lambda_k}
\),
which belongs to the physical tangent space by construction.
Conversely, if (4) holds then \(\delta C^{\mathrm{LoG}}_{\Lambda_k}\) cannot vanish, so (2) holds.

Finally, \cref{prop:intrinsic-cov-response,cor:compressed-export-selection-completion} identify the
abstract covariance tangent with any Gaussian completion representative of the same tangent object.
Therefore the abstract tangent is nonzero exactly when its completion realization is nonzero, giving
the equivalence of (2) and (5).
\end{proof}

\begin{theorem}[Branchwise realization and persistence of the dominant line]
\label{thm:dominant-line-persistence}
Let \(I\) be a connected accepted-branch segment in cutoff space.
Assume that:
\begin{enumerate}[leftmargin=2em]
\item the compressed projector \(P_{-,\Lambda}^{\mathrm{LoG}}\) varies continuously on \(I\);
\item the physical welded response operator
  \[
  \Lambda\longmapsto \Pi_{\mathrm{phys},\Lambda}\mathsf T_{\Lambda}
  \]
  is continuous on \(I\) in the operator topology induced by the Fisher norm on covariance
  tangents;
\item there exists a constant \(c_I>0\) such that
  \[
  \kappa_{\Lambda}\ge c_I
  \qquad
  \forall \Lambda\in I.
  \]
\end{enumerate}
Then:
\begin{enumerate}[leftmargin=2em]
\item \(\delta C^{\mathrm{LoG}}_{\Lambda}\neq 0\) for every \(\Lambda\in I\);
\item the normalized dominant vector
  \[
  \widehat{\delta C}^{\mathrm{LoG}}_{\Lambda}
  :=
  \kappa_{\Lambda}^{-1/2}\,\delta C^{\mathrm{LoG}}_{\Lambda}
  \]
  is continuous on \(I\);
\item the dominant physical line
  \[
  L^{\mathrm{dom}}_{\Lambda}
  =
  \mathbb R\,\widehat{\delta C}^{\mathrm{LoG}}_{\Lambda}
  \]
  is canonically transported along \(I\);
\item the Fisher-orthogonal rank-one projector
  \[
  \Pi^{\mathrm{dom}}_{\Lambda}
  :=
  \widehat{\delta C}^{\mathrm{LoG}}_{\Lambda}
  \otimes
  \widehat{\delta C}^{\mathrm{LoG}}_{\Lambda}
  \]
  depends continuously on \(\Lambda\in I\).
\end{enumerate}
\end{theorem}

\begin{proof}
By \cref{cor:compressed-export-selection}, the compressed projector
\(
P_{-,\Lambda}^{\mathrm{LoG}}
\)
and hence the normalized source vector \(u^{\mathrm{LoG}}_{\Lambda}\) vary continuously on the branch.
By the assumed continuity of \(\Pi_{\mathrm{phys},\Lambda}\mathsf T_{\Lambda}\), the map
\[
\Lambda\longmapsto
\delta C^{\mathrm{LoG}}_{\Lambda}
=
\Pi_{\mathrm{phys},\Lambda}\mathsf T_{\Lambda}[u^{\mathrm{LoG}}_{\Lambda}]
\]
is continuous in the Fisher norm.

The lower bound \(\kappa_{\Lambda}\ge c_I>0\) implies
\(
\delta C^{\mathrm{LoG}}_{\Lambda}\neq 0
\)
for every \(\Lambda\in I\), proving item (1).
Since the normalization map \(v\mapsto v/\|v\|_{C_\ast}\) is continuous on the complement of the
zero section, continuity of \(\delta C^{\mathrm{LoG}}_{\Lambda}\) and the positive lower bound on its
norm imply continuity of \(\widehat{\delta C}^{\mathrm{LoG}}_{\Lambda}\), proving item (2).
Items (3) and (4) are immediate consequences of item (2).
\end{proof}

\begin{corollary}[Local persistence from one-tick transversality]
\label{cor:dominant-line-local-persistence}
Assume the continuity hypotheses of \cref{thm:dominant-line-persistence}.
If the compressed line is transversely physical at one accepted tick \(\Lambda_{k_0}\), then there
exists a neighborhood \(I_{k_0}\) of \(\Lambda_{k_0}\) on the same accepted branch such that
\[
\delta C^{\mathrm{LoG}}_{\Lambda}\neq 0
\qquad
\forall \Lambda\in I_{k_0}.
\]
Hence the dominant physical line is locally well defined and continuously transported near
\(\Lambda_{k_0}\).
\end{corollary}

\begin{proof}
By \cref{prop:dominant-line-equivalent}, transversality at \(\Lambda_{k_0}\) is equivalent to
\(
\kappa_{\Lambda_{k_0}}>0
\).
Continuity of \(\kappa_{\Lambda}\) then gives a neighborhood \(I_{k_0}\) and a constant
\(c_{k_0}>0\) such that
\(
\kappa_{\Lambda}\ge c_{k_0}
\)
for all \(\Lambda\in I_{k_0}\).
Apply \cref{thm:dominant-line-persistence}.
\end{proof}

\begin{remark}[Practical witness criteria]
\label{rem:dominant-line-witnesses}
For applications, any one of the following suffices to discharge dominant-line realization at a fixed
accepted tick:
\begin{enumerate}[leftmargin=2em]
\item a direct lower bound \(\kappa_{\Lambda_k}>0\);
\item a physical witness tangent \(B_{\Lambda_k}\in\mathcal T^{\mathrm{phys}}_{\Lambda_k}\) with
  nonzero Fisher pairing against \(\delta C^{\mathrm{LoG}}_{\Lambda_k}\);
\item a Gaussian completion calculation showing that the completion covariance of the realized active
  odd mode changes nontrivially under the compressed welded LoG source.
\end{enumerate}
Each witness proves the same theorem-level fact by \cref{prop:dominant-line-equivalent}.
\end{remark}

\subsection*{Descent to the ambiguity quotient}
\label{subsec:ambiguity-quotient-descent}

Once the dominant line has been realized, the next task is purely structural:
show that the dominant-response coefficient is insensitive to ambiguity moves.
This is the exact point at which the transported descendant class becomes an affine class in the
sense required by Main rather than a collection of representatives.

\begin{definition}[Dominant-response quotient and descended descendant point]
\label{def:dominant-response-quotient}
For each accepted tick \(k\), define the dominant-response quotient by
\[
\mathfrak D^-_{\Lambda_k}
:=
E^-_{\Lambda_k}/\mathcal A^-_{\Lambda_k},
\]
with quotient map
\[
q_{\Lambda_k}:E^-_{\Lambda_k}\to \mathfrak D^-_{\Lambda_k}.
\]
If
\[
[\Xi^{(1)}_{\Lambda_k}]
=
J_{0,\Lambda_k}+\mathcal A^-_{\Lambda_k},
\]
then its image in the quotient is a single point,
\[
[\![\Xi^{(1)}_{\Lambda_k}]\!]
:=
q_{\Lambda_k}\!\big(J_{0,\Lambda_k}\big)
\in \mathfrak D^-_{\Lambda_k},
\]
independent of the chosen representative \(J_{0,\Lambda_k}\).
\end{definition}

\begin{proposition}[Universal descent criterion for the dominant-response coefficient]
\label{prop:dominant-quotient-universal}
Assume dominant-line realization at an accepted tick \(k\), so that
\(\rho^{\mathrm{dom}}_{\Lambda_k}\) is defined by \eqref{eq:rho-dom-def}.
Then the following are equivalent.
\begin{enumerate}[leftmargin=2em]
\item there exists a unique continuous linear functional
  \[
  \overline{\rho}^{\mathrm{dom}}_{\Lambda_k}:\mathfrak D^-_{\Lambda_k}\to\mathbb R
  \]
  such that
  \[
  \rho^{\mathrm{dom}}_{\Lambda_k}
  =
  \overline{\rho}^{\mathrm{dom}}_{\Lambda_k}\circ q_{\Lambda_k};
  \]
\item the ambiguity subspace is contained in the kernel:
  \[
  \mathcal A^-_{\Lambda_k}\subset \ker\rho^{\mathrm{dom}}_{\Lambda_k};
  \]
\item the dominant-response coefficient is constant on each ambiguity coset:
  \[
  \rho^{\mathrm{dom}}_{\Lambda_k}(F+A)
  =
  \rho^{\mathrm{dom}}_{\Lambda_k}(F)
  \qquad
  \forall F\in E^-_{\Lambda_k},\ \forall A\in\mathcal A^-_{\Lambda_k}.
  \]
\end{enumerate}
\end{proposition}

\begin{proof}
The equivalence of (2) and (3) is immediate from linearity:
\[
\rho^{\mathrm{dom}}_{\Lambda_k}(F+A)-\rho^{\mathrm{dom}}_{\Lambda_k}(F)
=
\rho^{\mathrm{dom}}_{\Lambda_k}(A).
\]

If (2) holds, define
\[
\overline{\rho}^{\mathrm{dom}}_{\Lambda_k}\big(q_{\Lambda_k}(F)\big)
:=
\rho^{\mathrm{dom}}_{\Lambda_k}(F).
\]
This is well defined precisely because changing representatives by an element of
\(\mathcal A^-_{\Lambda_k}\) does not change the value of \(\rho^{\mathrm{dom}}_{\Lambda_k}\).
Linearity and continuity are inherited from \(\rho^{\mathrm{dom}}_{\Lambda_k}\), and uniqueness
follows because \(q_{\Lambda_k}\) is surjective.

Conversely, if (1) holds and \(A\in\mathcal A^-_{\Lambda_k}\), then \(q_{\Lambda_k}(A)=0\), hence
\[
\rho^{\mathrm{dom}}_{\Lambda_k}(A)
=
\overline{\rho}^{\mathrm{dom}}_{\Lambda_k}\!\big(q_{\Lambda_k}(A)\big)
=
\overline{\rho}^{\mathrm{dom}}_{\Lambda_k}(0)
=
0.
\]
Thus (1) implies (2).
\end{proof}

\begin{corollary}[Constancy on the transported descendant affine class]
\label{cor:dominant-quotient-constancy}
Assume \cref{prop:dominant-quotient-universal}. Then
\[
\rho^{\mathrm{dom}}_{\Lambda_k}(J)
=
\overline{\rho}^{\mathrm{dom}}_{\Lambda_k}\big([\![\Xi^{(1)}_{\Lambda_k}]\!]\big)
\qquad
\forall J\in[\Xi^{(1)}_{\Lambda_k}].
\]
In particular, once the descended quotient functional exists, the descendant affine class carries a
single unambiguous dominant-response value.
\end{corollary}

\begin{proof}
Every \(J\in[\Xi^{(1)}_{\Lambda_k}]\) has the form
\(
J=J_{0,\Lambda_k}+A
\)
with \(A\in\mathcal A^-_{\Lambda_k}\), so
\[
q_{\Lambda_k}(J)=q_{\Lambda_k}(J_{0,\Lambda_k})=[\![\Xi^{(1)}_{\Lambda_k}]\!].
\]
Apply the factorization identity of \cref{prop:dominant-quotient-universal}.
\end{proof}

\begin{remark}[False stronger form and normalized special case]
\label{rem:dominant-quotient-false-stronger-form}
The quotient statement above is intrinsically a statement about the descended dominant scalar
\[
\overline{\rho}^{\mathrm{dom}}_{\Lambda_k}\big([\![\Xi^{(1)}_{\Lambda_k}]\!]\big),
\]
not about a bare \(\chi\)-compatibility condition.
At the weak transported-interface layer one has
\[
\chi_{\Lambda_k}
=
\nu_{\Lambda_k}^{-1}\rho^{\mathrm{dom}}_{\Lambda_k},
\]
so a representative-search condition of the form
\[
\rho^{\mathrm{dom}}_{\Lambda_k}(J)=s
\]
is equivalent to
\[
\chi_{\Lambda_k}(J)=\nu_{\Lambda_k}^{-1}s,
\]
not in general to \(\chi_{\Lambda_k}(J)=s\).
Only after an explicit normalization, for example \(\nu_{\Lambda_k}=1\), may the same scalar
comparison be rewritten in descended \(\bar\chi_{\Lambda_k}\) form without changing the value.
The finite-cutoff Wolfram audit contains exact countermodels to the stronger unscaled
\(\chi\)-level claim.
\end{remark}

\begin{proposition}[Generator-by-generator reduction]
\label{prop:dominant-quotient-generators}
Assume dominant-line realization at an accepted tick \(k\).
Suppose the dominant-response coefficient vanishes on the three generating ambiguity families:
\begin{enumerate}[leftmargin=2em]
\item for every odd-projected BRST/gauge-exact direction \(\big(Q_tX+Q_uY\big)^-\),
  \[
  \rho^{\mathrm{dom}}_{\Lambda_k}\!\big(\big(Q_tX+Q_uY\big)^-\big)=0;
  \]
\item for every odd-projected boundary-admissible divergence \(\big(\partial_\mu K^\mu\big)^-\),
  \[
  \rho^{\mathrm{dom}}_{\Lambda_k}\!\big(\big(\partial_\mu K^\mu\big)^-\big)=0;
  \]
\item for every admissible odd counterterm \(\Delta^-_{\rm ct}\),
  \[
  \rho^{\mathrm{dom}}_{\Lambda_k}(\Delta^-_{\rm ct})=0.
  \]
\end{enumerate}
Then
\[
\mathcal A^-_{\Lambda_k}\subset\ker\rho^{\mathrm{dom}}_{\Lambda_k},
\]
hence \(\rho^{\mathrm{dom}}_{\Lambda_k}\) descends uniquely to the quotient
\(\mathfrak D^-_{\Lambda_k}\).
\end{proposition}

\begin{proof}
By \cref{def:transport-ambiguity-subspace}, \(\mathcal A^-_{\Lambda_k}\) is the closed linear span of
the three families listed in the statement. Since \(\rho^{\mathrm{dom}}_{\Lambda_k}\) is linear, it
vanishes on the algebraic span of those generators. Since it is also continuous, it vanishes on the
closure of that span. Therefore every element of \(\mathcal A^-_{\Lambda_k}\) lies in the kernel.
Now apply \cref{prop:dominant-quotient-universal}.
\end{proof}

\begin{theorem}[Descent to the ambiguity quotient]
\label{thm:dominant-quotient-descent}
Fix an accepted tick \(k\) and assume dominant-line realization.
If the three vanishing statements of \cref{prop:dominant-quotient-generators} hold, then there
exists a unique descended dominant
functional
\[
\overline{\rho}^{\mathrm{dom}}_{\Lambda_k}:\mathfrak D^-_{\Lambda_k}\to\mathbb R
\]
such that
\[
\rho^{\mathrm{dom}}_{\Lambda_k}
=
\overline{\rho}^{\mathrm{dom}}_{\Lambda_k}\circ q_{\Lambda_k},
\]
and the transported descendant affine class has a single well-defined quotient value
\[
\overline{\rho}^{\mathrm{dom}}_{\Lambda_k}\big([\![\Xi^{(1)}_{\Lambda_k}]\!]\big).
\]
\end{theorem}

\begin{proof}
Apply \cref{prop:dominant-quotient-generators,cor:dominant-quotient-constancy}.
\end{proof}

\subsection*{Nonzero overlap on the descendant class}
\label{subsec:descendant-overlap}

After dominant-line realization and quotient descent, the descendant class defines a single point in
the ambiguity quotient, and the dominant-response functional descends to that quotient. This
subsection studies whether that descended descendant point is \emph{active}, i.e.\ not annihilated
by the descended dominant functional.

\begin{definition}[Dominant covariance projector]
\label{def:dominant-cov-projector}
Assume dominant-line realization at an accepted tick \(k\).
Define the Fisher-orthogonal rank-one projector onto the dominant physical line by
\begin{equation}
\Pi^{\mathrm{dom}}_{\Lambda_k}(X)
:=
\big\langle\!\big\langle
X,\,
\widehat{\delta C}^{\mathrm{LoG}}_{\Lambda_k}
\big\rangle\!\big\rangle_{C_{\ast,\Lambda_k}}\,
\widehat{\delta C}^{\mathrm{LoG}}_{\Lambda_k},
\qquad
X\in\mathcal T^{\mathrm{cov}}_{\Lambda_k}.
\label{eq:Pi-dom-def}
\end{equation}
\end{definition}

\begin{proposition}[Representative-level witness criteria]
\label{prop:descendant-overlap-witness}
Assume dominant-line realization at an accepted tick \(k\).
For any odd generator \(F\in E^-_{\Lambda_k}\), the following are equivalent.
\begin{enumerate}[leftmargin=2em]
\item the dominant-response coefficient is nonzero:
  \[
  \rho^{\mathrm{dom}}_{\Lambda_k}(F)\neq 0;
  \]
\item the dominant projection of the physical covariance response is nonzero:
  \[
  \Pi^{\mathrm{dom}}_{\Lambda_k}\Pi_{\mathrm{phys},\Lambda_k}\mathcal R_{\Lambda_k}(F)\neq 0;
  \]
\item the physical covariance response has nonzero Fisher pairing with the dominant tangent:
  \[
  \big\langle\!\big\langle
  \Pi_{\mathrm{phys},\Lambda_k}\mathcal R_{\Lambda_k}(F),\,
  \delta C^{\mathrm{LoG}}_{\Lambda_k}
  \big\rangle\!\big\rangle_{C_{\ast,\Lambda_k}}
  \neq 0;
  \]
\item in any Gaussian completion realization of
  \cref{cor:compressed-export-selection-completion}, the completion representative of
  \(\Pi_{\mathrm{phys},\Lambda_k}\mathcal R_{\Lambda_k}(F)\) has nonzero component on the realized
  active odd line.
\end{enumerate}
\end{proposition}

\begin{proof}
By \cref{def:dominant-cov-projector},
\[
\Pi^{\mathrm{dom}}_{\Lambda_k}\Pi_{\mathrm{phys},\Lambda_k}\mathcal R_{\Lambda_k}(F)
=
\rho^{\mathrm{dom}}_{\Lambda_k}(F)\,
\widehat{\delta C}^{\mathrm{LoG}}_{\Lambda_k}.
\]
Since \(\widehat{\delta C}^{\mathrm{LoG}}_{\Lambda_k}\neq 0\), this proves the equivalence of (1)
and (2).

Statement (3) is equivalent to (1) because
\[
\big\langle\!\big\langle
\Pi_{\mathrm{phys},\Lambda_k}\mathcal R_{\Lambda_k}(F),\,
\delta C^{\mathrm{LoG}}_{\Lambda_k}
\big\rangle\!\big\rangle_{C_{\ast,\Lambda_k}}
=
\|\delta C^{\mathrm{LoG}}_{\Lambda_k}\|_{C_{\ast,\Lambda_k}}\,
\rho^{\mathrm{dom}}_{\Lambda_k}(F),
\]
and the prefactor is strictly positive by dominant-line realization.

Finally, by \cref{prop:intrinsic-cov-response,cor:compressed-export-selection-completion}, the
abstract covariance tangent and its Gaussian completion representative describe the same physical
response object in different coordinates. Therefore the abstract dominant projection is nonzero if
and only if the completion representative has nonzero component on the realized active line,
proving the equivalence of (2) and (4).
\end{proof}

\begin{definition}[Descendant overlap margin]
\label{def:descendant-overlap-margin}
Assume dominant-line realization and quotient descent at an accepted tick \(k\).
Define the descendant overlap margin by
\begin{equation}
\nu_{\Lambda_k}
:=
\left|
\overline{\rho}^{\mathrm{dom}}_{\Lambda_k}\big([\![\Xi^{(1)}_{\Lambda_k}]\!]\big)
\right|.
\label{eq:descendant-overlap-margin}
\end{equation}
We say that the transported descendant class is \emph{dominant-active} if
\(
\nu_{\Lambda_k}>0
\).
\end{definition}

\begin{proposition}[Equivalent formulations of nonzero overlap on the descendant class]
\label{prop:descendant-overlap-equivalent}
Assume dominant-line realization and quotient descent at an accepted tick \(k\).
Then the following are equivalent.
\begin{enumerate}[leftmargin=2em]
\item the transported descendant sector has nonzero overlap with the canonically selected compressed
  dominant odd line;
\item there exists one representative \(J^{\mathrm{desc}}_{\Lambda_k}\in[\Xi^{(1)}_{\Lambda_k}]\)
  such that
  \[
  \Pi^{\mathrm{dom}}_{\Lambda_k}\Pi_{\mathrm{phys},\Lambda_k}
  \mathcal R_{\Lambda_k}(J^{\mathrm{desc}}_{\Lambda_k})
  \neq 0;
  \]
\item every representative \(J\in[\Xi^{(1)}_{\Lambda_k}]\) has the same nonzero dominant-response
  value:
  \[
  \rho^{\mathrm{dom}}_{\Lambda_k}(J)
  =
  \overline{\rho}^{\mathrm{dom}}_{\Lambda_k}\big([\![\Xi^{(1)}_{\Lambda_k}]\!]\big)
  \neq 0;
  \]
\item the descended descendant point is not annihilated by the descended dominant functional:
  \[
  [\![\Xi^{(1)}_{\Lambda_k}]\!]\notin \ker\overline{\rho}^{\mathrm{dom}}_{\Lambda_k};
  \]
\item the descendant overlap margin is strictly positive:
  \[
  \nu_{\Lambda_k}>0.
  \]
\end{enumerate}
\end{proposition}

\begin{proof}
The equivalence of (1) and (2) follows immediately from
\cref{prop:descendant-overlap-witness}.

If (2) holds, then
\(
\rho^{\mathrm{dom}}_{\Lambda_k}(J^{\mathrm{desc}}_{\Lambda_k})\neq 0
\)
by \cref{prop:descendant-overlap-witness}. By
\cref{cor:dominant-quotient-constancy}, every other representative of
\([\Xi^{(1)}_{\Lambda_k}]\) has the same value, proving (3).
Conversely, (3) obviously implies (1).

The equivalence of (3) and (4) is immediate from the identity
\[
\rho^{\mathrm{dom}}_{\Lambda_k}(J)
=
\overline{\rho}^{\mathrm{dom}}_{\Lambda_k}\big([\![\Xi^{(1)}_{\Lambda_k}]\!]\big)
\qquad
\forall J\in[\Xi^{(1)}_{\Lambda_k}]
\]
of \cref{cor:dominant-quotient-constancy}.

Finally, (4) and (5) are equivalent by the definition of \(\nu_{\Lambda_k}\) in
\cref{def:descendant-overlap-margin}.
\end{proof}

\begin{theorem}[Nonzero overlap on the descendant class]
\label{thm:descendant-overlap}
Fix an accepted tick \(k\), and assume dominant-line realization together with quotient descent.
If any one of the equivalent conditions in \cref{prop:descendant-overlap-equivalent} holds, then the
following conclusions obtain.
In particular, it is enough to construct a single representative
\[
J^{\mathrm{desc}}_{\Lambda_k}\in[\Xi^{(1)}_{\Lambda_k}]
\]
whose dominant projected covariance response is nonzero.

When this holds, the quotient descendant point is dominant-active, the scalar value
\[
\overline{\rho}^{\mathrm{dom}}_{\Lambda_k}\big([\![\Xi^{(1)}_{\Lambda_k}]\!]\big)
\]
is nonzero, and the sign
\begin{equation}
\varepsilon_{\Lambda_k}
:=
\operatorname{sgn}\!\Big(
\overline{\rho}^{\mathrm{dom}}_{\Lambda_k}\big([\![\Xi^{(1)}_{\Lambda_k}]\!]\big)
\Big)
\in\{\pm1\}
\label{eq:epsilon-from-overlap}
\end{equation}
is well defined.
Thus the transported descendant sector has nonzero overlap with the canonically selected compressed
dominant odd line.
\end{theorem}

\begin{proof}
The equivalence package is exactly \cref{prop:descendant-overlap-equivalent}.
Once one representative has nonzero dominant projected response, the same proposition gives the
nonzero scalar value on the descended descendant point. The sign in \eqref{eq:epsilon-from-overlap}
is therefore well defined.
\end{proof}

\subsection*{Conditional mod-null reentry on the selected scalar line}
\label{subsec:conditional-mod-null-reentry}

\begin{definition}[Parity-stable mod-null subspace transverse to the selected scalar line]
\label{def:mod-null-transverse-datum}
Fix an accepted tick \(k\), assume dominant-line realization and quotient descent, and let
\[
\overline L^{\mathrm{Main}}_{\Lambda_k}
=
\mathbb R\,\overline\ell^{\mathrm{Main}}_{\Lambda_k}
\subset
\mathfrak D^-_{\Lambda_k}
\]
be a one-dimensional selected scalar line in the ambiguity quotient.
A \emph{parity-stable mod-null datum transverse to}
\(
\overline L^{\mathrm{Main}}_{\Lambda_k}
\)
consists of a linear subspace
\[
N^{\mathrm{mn}}_{\Lambda_k}\subset \mathfrak D^-_{\Lambda_k}
\]
such that:
\begin{enumerate}[leftmargin=2em,label=(\alph*)]
\item \(N^{\mathrm{mn}}_{\Lambda_k}\) is stable under the descended physical parity on
  \(\mathfrak D^-_{\Lambda_k}\);
\item the selected scalar line is transverse to \(N^{\mathrm{mn}}_{\Lambda_k}\):
  \[
  \overline L^{\mathrm{Main}}_{\Lambda_k}\cap N^{\mathrm{mn}}_{\Lambda_k}=\{0\}.
  \]
\end{enumerate}
Define
\[
\mathcal V^{1,\mathrm{mn}}_{\Lambda_k}
:=
\mathfrak D^-_{\Lambda_k}/N^{\mathrm{mn}}_{\Lambda_k}
\]
with quotient map
\[
\pi^{\mathrm{mn}}_{\Lambda_k}:
\mathfrak D^-_{\Lambda_k}\to \mathcal V^{1,\mathrm{mn}}_{\Lambda_k}.
\]
Because \(\mathfrak D^-_{\Lambda_k}\) is itself a quotient of the physical \(\Gamma\)-odd sector,
the descended parity acts on \(\mathfrak D^-_{\Lambda_k}\) by \(-\mathrm{id}\).
Thus the parity-stability clause yields a descended involution
\[
\Gamma^{\mathrm{mn}}_{\Lambda_k}
\bigl(
\pi^{\mathrm{mn}}_{\Lambda_k}(x)
\bigr)
:=
\pi^{\mathrm{mn}}_{\Lambda_k}(-x),
\qquad
x\in \mathfrak D^-_{\Lambda_k}.
\]
\end{definition}

\begin{definition}[Dominant-null parity-stable mod-null datum]
\label{def:dominant-null-mod-null-datum}
Fix an accepted tick \(k\), assume dominant-line realization and quotient descent, and let
\[
\overline L^{\mathrm{Main}}_{\Lambda_k}
=
\mathbb R\,\overline\ell^{\mathrm{Main}}_{\Lambda_k}
\subset
\mathfrak D^-_{\Lambda_k}
\]
be a one-dimensional selected scalar line in the ambiguity quotient.
A parity-stable mod-null datum
\[
N^{\mathrm{mn}}_{\Lambda_k}\subset \mathfrak D^-_{\Lambda_k}
\]
is called \emph{dominant-null} if it is contained in the kernel of the descended dominant
functional:
\[
N^{\mathrm{mn}}_{\Lambda_k}\subset \ker \overline{\rho}^{\mathrm{dom}}_{\Lambda_k}.
\]
It is understood that
\(
N^{\mathrm{mn}}_{\Lambda_k}
\)
is stable under the descended physical parity on
\(
\mathfrak D^-_{\Lambda_k}
\),
and that the associated quotient
\(
\mathfrak D^-_{\Lambda_k}/N^{\mathrm{mn}}_{\Lambda_k}
\)
carries the descended involution from \cref{def:mod-null-transverse-datum}.
The canonical dominant-kernel choice is isolated separately in
\cref{def:canonical-dominant-null-datum}.
The transversality lemma below uses only this kernel-subordination property.
\end{definition}

\begin{definition}[Canonical dominant-null datum]
\label{def:canonical-dominant-null-datum}
Fix an accepted tick \(k\) and assume dominant-line realization together with quotient descent.
Define the \emph{canonical dominant-null datum} by
\[
N^{\mathrm{dom-null}}_{\Lambda_k}
:=
\ker \overline{\rho}^{\mathrm{dom}}_{\Lambda_k}
\subset
\mathfrak D^-_{\Lambda_k}.
\]
This datum depends only on the descended dominant functional on the ambiguity quotient.
No additional reduced-welded or completion-side choice is built into the definition.
\end{definition}

\begin{lemma}[Dominant-kernel criterion for selected-line transversality]
\label{lem:dominant-kernel-transversality}
Fix an accepted tick \(k\), assume dominant-line realization together with quotient descent, and let
\[
\overline L^{\mathrm{Main}}_{\Lambda_k}
=
\mathbb R\,\overline\ell^{\mathrm{Main}}_{\Lambda_k}
\subset
\mathfrak D^-_{\Lambda_k}
\]
be a one-dimensional selected scalar line with
\[
\overline{\rho}^{\mathrm{dom}}_{\Lambda_k}
\!\bigl(
\overline\ell^{\mathrm{Main}}_{\Lambda_k}
\bigr)
\neq 0.
\]
If
\[
N^{\mathrm{mn}}_{\Lambda_k}\subset \mathfrak D^-_{\Lambda_k}
\]
is a dominant-null parity-stable mod-null datum in the sense of
\cref{def:dominant-null-mod-null-datum},
then
\[
\overline L^{\mathrm{Main}}_{\Lambda_k}\cap N^{\mathrm{mn}}_{\Lambda_k}=\{0\}.
\]
So every dominant-null parity-stable mod-null datum is automatically transverse to the selected
scalar line.
\end{lemma}

\begin{proof}
Let
\(
x\in \overline L^{\mathrm{Main}}_{\Lambda_k}\cap N^{\mathrm{mn}}_{\Lambda_k}
\).
Because
\(
\overline L^{\mathrm{Main}}_{\Lambda_k}
=
\mathbb R\,\overline\ell^{\mathrm{Main}}_{\Lambda_k}
\),
there exists \(a\in\mathbb R\) with
\(
x=a\,\overline\ell^{\mathrm{Main}}_{\Lambda_k}
\).
Since
\(
N^{\mathrm{mn}}_{\Lambda_k}\subset
\ker \overline{\rho}^{\mathrm{dom}}_{\Lambda_k}
\),
we have
\[
0
=
\overline{\rho}^{\mathrm{dom}}_{\Lambda_k}(x)
=
a\,
\overline{\rho}^{\mathrm{dom}}_{\Lambda_k}
\!\bigl(
\overline\ell^{\mathrm{Main}}_{\Lambda_k}
\bigr).
\]
The coefficient
\(
\overline{\rho}^{\mathrm{dom}}_{\Lambda_k}
(\overline\ell^{\mathrm{Main}}_{\Lambda_k})
\)
is nonzero by hypothesis, so \(a=0\).
Therefore \(x=0\), proving
\(
\overline L^{\mathrm{Main}}_{\Lambda_k}\cap N^{\mathrm{mn}}_{\Lambda_k}=\{0\}
\).
\end{proof}

\begin{theorem}[The dominant kernel is the canonical dominant-null datum]
\label{thm:canonical-dominant-null-datum}
Fix an accepted tick \(k\), assume dominant-line realization together with quotient descent, and let
\[
\overline L^{\mathrm{Main}}_{\Lambda_k}
=
\mathbb R\,\overline\ell^{\mathrm{Main}}_{\Lambda_k}
\subset
\mathfrak D^-_{\Lambda_k}
\]
be a one-dimensional selected scalar line with
\[
\overline{\rho}^{\mathrm{dom}}_{\Lambda_k}
\!\bigl(
\overline\ell^{\mathrm{Main}}_{\Lambda_k}
\bigr)
\neq 0.
\]
Define
\[
N^{\mathrm{dom-null}}_{\Lambda_k}
:=
\ker \overline{\rho}^{\mathrm{dom}}_{\Lambda_k}
\subset
\mathfrak D^-_{\Lambda_k}.
\]
Then:
\begin{enumerate}[leftmargin=2em]
\item \(N^{\mathrm{dom-null}}_{\Lambda_k}\) is stable under the descended physical parity on
  \(\mathfrak D^-_{\Lambda_k}\);
\item \(N^{\mathrm{dom-null}}_{\Lambda_k}\) is a dominant-null parity-stable mod-null datum in the
  sense of \cref{def:dominant-null-mod-null-datum};
\item the selected scalar line is automatically transverse to the canonical dominant-null datum:
  \[
  \overline L^{\mathrm{Main}}_{\Lambda_k}
  \cap
  N^{\mathrm{dom-null}}_{\Lambda_k}
  =
  \{0\}.
  \]
\end{enumerate}
Hence the canonical dominant kernel is admissible as the mod-null datum for the conditional
reentry theorem.
\end{theorem}

\begin{proof}
Let \(x\in N^{\mathrm{dom-null}}_{\Lambda_k}\).
Because the descended physical parity acts on \(\mathfrak D^-_{\Lambda_k}\) by \(-\mathrm{id}\),
we have
\[
\overline{\rho}^{\mathrm{dom}}_{\Lambda_k}(-x)
=
-\overline{\rho}^{\mathrm{dom}}_{\Lambda_k}(x)
=
0.
\]
So \(-x\in N^{\mathrm{dom-null}}_{\Lambda_k}\), proving parity stability.
Item \textnormal{(2)} is then exactly \cref{def:dominant-null-mod-null-datum} with
\(
N^{\mathrm{mn}}_{\Lambda_k}=N^{\mathrm{dom-null}}_{\Lambda_k}
\).
Item \textnormal{(3)} is \cref{lem:dominant-kernel-transversality} applied to this canonical
choice.
\end{proof}

\begin{lemma}[Transversality reduces mod-null reentry to quotient construction]
\label{lem:mod-null-transversality-reduction}
Fix an accepted tick \(k\), assume dominant-line realization together with quotient descent, and let
\[
\overline L^{\mathrm{Main}}_{\Lambda_k}
=
\mathbb R\,\overline\ell^{\mathrm{Main}}_{\Lambda_k}
\subset
\mathfrak D^-_{\Lambda_k}
\]
be a one-dimensional selected scalar line with
\[
\overline{\rho}^{\mathrm{dom}}_{\Lambda_k}
\!\bigl(
\overline\ell^{\mathrm{Main}}_{\Lambda_k}
\bigr)
\neq 0.
\]
Assume moreover a parity-stable mod-null datum transverse to
\(
\overline L^{\mathrm{Main}}_{\Lambda_k}
\)
in the sense of \cref{def:mod-null-transverse-datum}.
Define
\[
\ell^{\mathrm{sel,mn}}_{\Lambda_k}
:=
\pi^{\mathrm{mn}}_{\Lambda_k}
\!\bigl(
\overline\ell^{\mathrm{Main}}_{\Lambda_k}
\bigr)
\]
and
\[
L^{\mathrm{sel,mn}}_{\Lambda_k}
:=
\mathbb R\,\ell^{\mathrm{sel,mn}}_{\Lambda_k}
\subset
\mathcal V^{1,\mathrm{mn}}_{\Lambda_k}.
\]
Then:
\begin{enumerate}[leftmargin=2em]
\item the restriction of
  \(
  \pi^{\mathrm{mn}}_{\Lambda_k}
  \)
  to
  \(
  \overline L^{\mathrm{Main}}_{\Lambda_k}
  \)
  is injective;
\item \(L^{\mathrm{sel,mn}}_{\Lambda_k}\) is a well-defined one-dimensional line in the mod-null
  quotient, generated by the nonzero class
  \(
  \ell^{\mathrm{sel,mn}}_{\Lambda_k}
  \);
\item the mod-null selected realization map
  \[
  \mathfrak R^{\mathrm{mn}}_{\Lambda_k}:
  L^{\mathrm{sel,mn}}_{\Lambda_k}
  \longrightarrow
  \mathbb R\,u^{\mathrm{LoG}}_{\Lambda_k},
  \qquad
  \mathfrak R^{\mathrm{mn}}_{\Lambda_k}
  \!\bigl(
  a\,\ell^{\mathrm{sel,mn}}_{\Lambda_k}
  \bigr)
  :=
  a\,
  \overline{\rho}^{\mathrm{dom}}_{\Lambda_k}
  \!\bigl(
  \overline\ell^{\mathrm{Main}}_{\Lambda_k}
  \bigr)\,
  u^{\mathrm{LoG}}_{\Lambda_k}
  \]
  is well defined and faithful;
\item the descended involution acts by sign:
  \[
  \Gamma^{\mathrm{mn}}_{\Lambda_k}
  \!\bigl(
  \ell^{\mathrm{sel,mn}}_{\Lambda_k}
  \bigr)
  =
  -
  \ell^{\mathrm{sel,mn}}_{\Lambda_k}.
  \]
\end{enumerate}
In particular, once the transversality condition
\(
\overline L^{\mathrm{Main}}_{\Lambda_k}\cap N^{\mathrm{mn}}_{\Lambda_k}=\{0\}
\)
is known, the remaining mod-null reentry step reduces to quotient construction plus parity descent.
\end{lemma}

\begin{proof}
If
\(
x\in \overline L^{\mathrm{Main}}_{\Lambda_k}
\)
and
\(
\pi^{\mathrm{mn}}_{\Lambda_k}(x)=0
\),
then
\(
x\in N^{\mathrm{mn}}_{\Lambda_k}
\)
by definition of the quotient kernel.
Hence
\(
x\in
\overline L^{\mathrm{Main}}_{\Lambda_k}\cap N^{\mathrm{mn}}_{\Lambda_k}
=
\{0\}
\),
which proves injectivity.
Since
\(
\overline\ell^{\mathrm{Main}}_{\Lambda_k}\neq 0
\),
its image
\(
\ell^{\mathrm{sel,mn}}_{\Lambda_k}
\)
is therefore nonzero, and the image of a one-dimensional line under an injective linear map is
again one-dimensional.
This proves items \textnormal{(1)} and \textnormal{(2)}.

Injectivity on
\(
\overline L^{\mathrm{Main}}_{\Lambda_k}
\)
implies that every element of
\(
L^{\mathrm{sel,mn}}_{\Lambda_k}
\)
has a unique expression
\(
a\,\ell^{\mathrm{sel,mn}}_{\Lambda_k}
\),
so the displayed formula defines
\(
\mathfrak R^{\mathrm{mn}}_{\Lambda_k}
\)
unambiguously.
Its coefficient
\(
\overline{\rho}^{\mathrm{dom}}_{\Lambda_k}
(\overline\ell^{\mathrm{Main}}_{\Lambda_k})
\)
is nonzero by hypothesis, and
\(
u^{\mathrm{LoG}}_{\Lambda_k}\neq 0
\)
by \cref{cor:compressed-export-selection}.
Therefore
\(
\mathfrak R^{\mathrm{mn}}_{\Lambda_k}
(a\,\ell^{\mathrm{sel,mn}}_{\Lambda_k})
=
0
\)
implies
\(
a=0
\),
so the map is faithful.
This proves item \textnormal{(3)}.

By construction of the descended involution,
\[
\Gamma^{\mathrm{mn}}_{\Lambda_k}
\!\bigl(
\ell^{\mathrm{sel,mn}}_{\Lambda_k}
\bigr)
=
\Gamma^{\mathrm{mn}}_{\Lambda_k}
\!\Bigl(
\pi^{\mathrm{mn}}_{\Lambda_k}
\!\bigl(
\overline\ell^{\mathrm{Main}}_{\Lambda_k}
\bigr)
\Bigr)
=
\pi^{\mathrm{mn}}_{\Lambda_k}
\!\bigl(
-\overline\ell^{\mathrm{Main}}_{\Lambda_k}
\bigr)
=
-\ell^{\mathrm{sel,mn}}_{\Lambda_k},
\]
which proves item \textnormal{(4)}.
\end{proof}

\begin{theorem}[Conditional mod-null reentry on the selected scalar line]
\label{thm:conditional-mod-null-reentry}
Fix an accepted tick \(k\) and the next accepted receiver step \(k+1\).
Assume the hypotheses and notation of
\cref{splitlift-thm:splitlift-scalar-descendant-lift}.
In particular, let
\[
\overline L^{\mathrm{Main}}_{\Lambda_{k+1}}
=
\mathbb R\,\overline\ell^{\mathrm{Main}}_{\Lambda_{k+1}}
\subset
\mathfrak D^-_{\Lambda_{k+1}}
\]
be the nonzero selected scalar descendant line supplied there.
Assume moreover a parity-stable mod-null datum transverse to
\(
\overline L^{\mathrm{Main}}_{\Lambda_{k+1}}
\)
in the sense of \cref{def:mod-null-transverse-datum}.
Define
\[
\ell^{\mathrm{sel,mn}}_{\Lambda_{k+1}}
:=
\pi^{\mathrm{mn}}_{\Lambda_{k+1}}
\!\bigl(
\overline\ell^{\mathrm{Main}}_{\Lambda_{k+1}}
\bigr)
\]
and
\[
L^{\mathrm{sel,mn}}_{\Lambda_{k+1}}
:=
\mathbb R\,\ell^{\mathrm{sel,mn}}_{\Lambda_{k+1}}
\subset
\mathcal V^{1,\mathrm{mn}}_{\Lambda_{k+1}}.
\]
Then:
\begin{enumerate}[leftmargin=2em]
\item \(L^{\mathrm{sel,mn}}_{\Lambda_{k+1}}\) is a well-defined one-dimensional selected descendant
  line in the mod-null quotient;
\item the induced mod-null selected realization map
  \[
  \mathfrak R^{\mathrm{mn}}_{\Lambda_{k+1}}:
  L^{\mathrm{sel,mn}}_{\Lambda_{k+1}}
  \longrightarrow
  \mathbb R\,u^{\mathrm{LoG}}_{\Lambda_{k+1}},
  \qquad
  \mathfrak R^{\mathrm{mn}}_{\Lambda_{k+1}}
  \!\bigl(
  a\,\ell^{\mathrm{sel,mn}}_{\Lambda_{k+1}}
  \bigr)
  :=
  a\,
  \overline{\rho}^{\mathrm{dom}}_{\Lambda_{k+1}}
  \!\bigl(
  \overline\ell^{\mathrm{Main}}_{\Lambda_{k+1}}
  \bigr)\,
  u^{\mathrm{LoG}}_{\Lambda_{k+1}}
  \]
  is well defined and faithful;
\item the descended involution acts by sign:
  \[
  \Gamma^{\mathrm{mn}}_{\Lambda_{k+1}}
  \!\bigl(
  \ell^{\mathrm{sel,mn}}_{\Lambda_{k+1}}
  \bigr)
  =
  -
  \ell^{\mathrm{sel,mn}}_{\Lambda_{k+1}}.
  \]
\end{enumerate}
\end{theorem}

\begin{proof}
By \cref{splitlift-thm:splitlift-scalar-descendant-lift},
\[
\overline\ell^{\mathrm{Main}}_{\Lambda_{k+1}}\neq 0
\qquad\text{and}\qquad
\overline{\rho}^{\mathrm{dom}}_{\Lambda_{k+1}}
\!\bigl(
\overline\ell^{\mathrm{Main}}_{\Lambda_{k+1}}
\bigr)
\neq 0.
\]
So every hypothesis of
\cref{lem:mod-null-transversality-reduction}
is in force at \(\Lambda_{k+1}\).
Applying that lemma gives items \textnormal{(1)}--\textnormal{(3)} immediately.
\end{proof}

\begin{corollary}[Canonical dominant-kernel discharge of the mod-null reentry step]
\label{cor:dominant-kernel-mod-null-reentry}
Fix an accepted tick \(k\) and the next accepted receiver step \(k+1\).
Assume the hypotheses and notation of
\cref{splitlift-thm:splitlift-scalar-descendant-lift}.
Define the canonical dominant-null datum
\[
N^{\mathrm{dom-null}}_{\Lambda_{k+1}}
:=
\ker \overline{\rho}^{\mathrm{dom}}_{\Lambda_{k+1}}
\subset
\mathfrak D^-_{\Lambda_{k+1}}.
\]
Then \(N^{\mathrm{dom-null}}_{\Lambda_{k+1}}\) is admissible as the mod-null datum of
\cref{thm:conditional-mod-null-reentry}, the transversality hypothesis of that theorem is
automatic, and hence the conclusions of that theorem hold with this canonical datum.
\end{corollary}

\begin{proof}
By \cref{splitlift-thm:splitlift-scalar-descendant-lift},
\[
\overline{\rho}^{\mathrm{dom}}_{\Lambda_{k+1}}
\!\bigl(
\overline\ell^{\mathrm{Main}}_{\Lambda_{k+1}}
\bigr)
\neq 0.
\]
So \cref{thm:canonical-dominant-null-datum} applies at \(\Lambda_{k+1}\).
It shows that
\(
N^{\mathrm{dom-null}}_{\Lambda_{k+1}}
\)
is a dominant-null parity-stable mod-null datum and that
\[
\overline L^{\mathrm{Main}}_{\Lambda_{k+1}}
\cap
N^{\mathrm{dom-null}}_{\Lambda_{k+1}}
=
\{0\}.
\]
This is exactly the admissibility and transversality clause required by
\cref{thm:conditional-mod-null-reentry}, so that theorem now applies directly with the canonical
datum.
\end{proof}

\begin{remark}[Conditional character of the reentry step]
\label{rem:conditional-mod-null-reentry}
\cref{thm:conditional-mod-null-reentry} is a quotient-level reentry statement.
After
\cref{lem:dominant-kernel-transversality,thm:canonical-dominant-null-datum,cor:dominant-kernel-mod-null-reentry},
the canonical dominant-null datum is fixed as
\(
\ker\overline{\rho}^{\mathrm{dom}}_{\Lambda_{k+1}}
\)
on \(\mathfrak D^-_{\Lambda_{k+1}}\).
What remains is passage to the mod-null quotient together with parity descent on the odd line.
Accordingly, the theorem should not be read as a statement on the raw descendant ambiguity quotient
or on any stronger physical quotient.
\end{remark}

\subsection*{Branchwise orientation transport}
\label{subsec:branchwise-orientation}

This subsection studies transport of the descended descendant sign along a connected accepted
branch. The sign can be transported once the descended descendant value varies continuously with the
cutoff. The compressed spectral-selection theorem already gives continuity of the dominant
projector, and \cref{thm:dominant-line-persistence} gives continuity of the dominant physical line.
We combine those inputs with continuity of the projected descendant response.

\begin{definition}[Branchwise descendant overlap function]
\label{def:branchwise-overlap-function}
Let \(I\) be a connected accepted-branch segment.
Assume dominant-line realization and quotient descent for every \(\Lambda\in I\).
Suppose there is a branchwise choice of representatives
\[
J^{\mathrm{desc}}_{\Lambda}\in[\Xi^{(1)}_{\Lambda}]
\qquad
(\Lambda\in I)
\]
such that the projected covariance response
\[
\Lambda\longmapsto
\Pi_{\mathrm{phys},\Lambda}\mathcal R_{\Lambda}(J^{\mathrm{desc}}_{\Lambda})
\]
is continuous in the Fisher topology.
Define the associated branchwise descendant overlap function by
\begin{equation}
m_I(\Lambda)
:=
\rho^{\mathrm{dom}}_{\Lambda}(J^{\mathrm{desc}}_{\Lambda})
=
\overline{\rho}^{\mathrm{dom}}_{\Lambda}\big([\![\Xi^{(1)}_{\Lambda}]\!]\big).
\label{eq:branchwise-overlap-function}
\end{equation}
\end{definition}

\begin{proposition}[Continuity of the descended overlap scalar]
\label{prop:branchwise-overlap-continuity}
Let \(I\) be a connected accepted-branch segment.
Assume:
\begin{enumerate}[leftmargin=2em]
\item the hypotheses of \cref{thm:dominant-line-persistence} on \(I\);
\item quotient descent holds for every \(\Lambda\in I\);
\item the branchwise choice of representatives in
  \cref{def:branchwise-overlap-function} exists.
\end{enumerate}
Then the scalar function \(m_I:I\to\mathbb R\) is well defined, independent of the chosen
branchwise representatives, and continuous.
Moreover,
\begin{equation}
|m_I(\Lambda)|
=
\nu_{\Lambda}
\qquad
\forall \Lambda\in I.
\label{eq:branchwise-overlap-margin}
\end{equation}
\end{proposition}

\begin{proof}
Independence of the chosen branchwise representatives is exactly the coset constancy proved in
\cref{cor:dominant-quotient-constancy}: if
\(
J'^{\mathrm{desc}}_{\Lambda}-J^{\mathrm{desc}}_{\Lambda}\in\mathcal A^-_{\Lambda}
\),
then
\[
\rho^{\mathrm{dom}}_{\Lambda}(J'^{\mathrm{desc}}_{\Lambda})
=
\rho^{\mathrm{dom}}_{\Lambda}(J^{\mathrm{desc}}_{\Lambda})
=
\overline{\rho}^{\mathrm{dom}}_{\Lambda}\big([\![\Xi^{(1)}_{\Lambda}]\!]\big).
\]
Thus \(m_I\) is well defined.

By \cref{thm:dominant-line-persistence}, the normalized dominant vector
\(
\widehat{\delta C}^{\mathrm{LoG}}_{\Lambda}
\)
and the dominant projector
\(
\Pi^{\mathrm{dom}}_{\Lambda}
\)
depend continuously on \(\Lambda\in I\).
By assumption, the response vector
\[
X_{\Lambda}
:=
\Pi_{\mathrm{phys},\Lambda}\mathcal R_{\Lambda}(J^{\mathrm{desc}}_{\Lambda})
\]
is also continuous on \(I\).
Using \cref{def:dominant-cov-projector}, we have
\[
\Pi^{\mathrm{dom}}_{\Lambda}X_{\Lambda}
=
m_I(\Lambda)\,\widehat{\delta C}^{\mathrm{LoG}}_{\Lambda}.
\]
The left-hand side is continuous because both factors are continuous, while
\(
\widehat{\delta C}^{\mathrm{LoG}}_{\Lambda}\neq 0
\)
for all \(\Lambda\in I\) by \cref{thm:dominant-line-persistence}.
Therefore the scalar coefficient \(m_I(\Lambda)\) is continuous.

Finally, \eqref{eq:branchwise-overlap-margin} is just the definition of \(\nu_{\Lambda}\) from
\cref{def:descendant-overlap-margin} together with \eqref{eq:branchwise-overlap-function}.
\end{proof}

\begin{theorem}[Branchwise orientation transport on the descendant class]
\label{thm:branchwise-orientation-transport}
Let \(I\) be a connected accepted-branch segment and assume the hypotheses of
\cref{prop:branchwise-overlap-continuity}.
Then:
\begin{enumerate}[leftmargin=2em]
\item on every connected subsegment \(I'\subset I\) such that
  \[
  \nu_{\Lambda}>0
  \qquad
  \forall \Lambda\in I',
  \]
  the sign
  \begin{equation}
  \varepsilon_{\Lambda}
  :=
  \operatorname{sgn} m_I(\Lambda)
  =
  \operatorname{sgn}\!\Big(
  \overline{\rho}^{\mathrm{dom}}_{\Lambda}\big([\![\Xi^{(1)}_{\Lambda}]\!]\big)
  \Big)
  \in\{\pm1\}
  \label{eq:branchwise-epsilon}
  \end{equation}
  is constant on \(I'\);
\item if \(\nu_{\Lambda_0}>0\) at one point \(\Lambda_0\in I\), then there exists a neighborhood
  \(I_0\subset I\) of \(\Lambda_0\) on which \(\varepsilon_{\Lambda}\) is constant and equal to
  \(\varepsilon_{\Lambda_0}\);
\item under the present hypotheses, a change of sign can occur only at a cutoff where the
  descendant overlap closes
  \(
  \nu_{\Lambda}=0
  \).
\end{enumerate}
In particular, once dominant-line realization, quotient descent, and nonzero overlap are
established on an accepted branch segment and the projected descendant response varies continuously
there, the branchwise sign statement of \cref{thm:active-response-transport-minimal} follows.
\end{theorem}

\begin{proof}
By \cref{prop:branchwise-overlap-continuity}, \(m_I(\Lambda)\) is a continuous real-valued function
on \(I\).
On any connected subsegment \(I'\) where \(\nu_{\Lambda}=|m_I(\Lambda)|>0\), the function \(m_I\)
never vanishes.
A continuous nonvanishing real-valued function on a connected set has constant sign, proving item
(1).

Item (2) is the local version of the same argument: if \(m_I(\Lambda_0)\neq 0\), continuity gives a
neighborhood \(I_0\) on which \(m_I\) keeps the sign it has at \(\Lambda_0\).

For item (3), any sign change would force \(m_I\) to pass through zero by the intermediate value
theorem, hence \(\nu_{\Lambda}=|m_I(\Lambda)|=0\) at some intermediate cutoff.
\end{proof}

\begin{corollary}[Local orientation persistence from one-tick activity]
\label{cor:branchwise-orientation-local}
Let \(I\) be a connected accepted-branch segment and assume the hypotheses of
\cref{prop:branchwise-overlap-continuity}.
If the transported descendant class is dominant-active at one cutoff \(\Lambda_0\in I\), i.e.
\[
\nu_{\Lambda_0}>0,
\]
then there exists a neighborhood \(I_0\subset I\) of \(\Lambda_0\) such that
\[
\varepsilon_{\Lambda}=\varepsilon_{\Lambda_0}
\qquad
\forall \Lambda\in I_0.
\]
\end{corollary}

\begin{proof}
This is exactly item (2) of \cref{thm:branchwise-orientation-transport}.
\end{proof}

\subsection*{Witness-indexed seam realization and scalar bridge}
\label{subsec:witness-seam-bridge}

The bridge used below factors through a common LS witness datum rather than a direct identification
of a UOM seam source with a descendant representative.
On the UOM side one realizes a filler tag together with a detector witness by an admissible odd
welded source, measures the resulting dominant physical pairing, and then compares that scalar with
the odd representative attached to the same witness.

\begin{definition}[Accepted record payload of a welded odd source]
\label{def:accepted-record-payload}
Fix an accepted tick \(k\).
For an admissible welded odd trace-channel source \(\tau\), let
\[
\rho_{\Lambda_k}[\tau]
\]
denote the receiver state at cutoff \(\Lambda_k\) produced by that source in the welded UOM channel.
Define its accepted record payload by
\begin{equation}
\mathcal P^{\mathrm{acc}}_{\Lambda_k}[\tau]
:=
\Pi_{\Delta(\Lambda_k)}\!\big(\rho_{\Lambda_k}[\tau]\big).
\label{eq:accepted-record-payload}
\end{equation}
This is the diagonal/record part retained by the finite-cutoff acceptance map.
\end{definition}

\begin{definition}[Witness-indexed seam realization at an accepted tick]
\label{def:witness-seam-realization}
Fix an accepted tick \(k\).
Let \(y\) be a witness on the LS difference object, and write \(\mathrm{flip}(y)\) for the induced
LS involution.
A \emph{witness-indexed seam realization} at \(\Lambda_k\) is a choice of admissible welded odd
trace-channel sources
\[
\tau^{f,y}_{\Lambda_k},
\qquad
f\in\{\mathrm{diag},\mathrm{bridge}\},
\]
such that:
\begin{enumerate}[leftmargin=2em]
\item \textbf{filler toggle changes the odd seam sign:}
  \[
  \tau^{\mathrm{bridge},y}_{\Lambda_k}
  =
  -\,\tau^{\mathrm{diag},y}_{\Lambda_k};
  \]
\item \textbf{witness flip changes the odd seam sign:}
  \[
  \tau^{f,\mathrm{flip}(y)}_{\Lambda_k}
  =
  -\,\tau^{f,y}_{\Lambda_k}
  \qquad
  \forall f\in\{\mathrm{diag},\mathrm{bridge}\};
  \]
\item \textbf{accepted record payload is unchanged under filler toggle:}
  \[
  \mathcal P^{\mathrm{acc}}_{\Lambda_k}\!\big[\tau^{\mathrm{bridge},y}_{\Lambda_k}\big]
  =
  \mathcal P^{\mathrm{acc}}_{\Lambda_k}\!\big[\tau^{\mathrm{diag},y}_{\Lambda_k}\big].
  \]
\end{enumerate}
\end{definition}

\begin{remark}[Accepted-record invariance versus seam-scalar use]
\label{rem:accepted-record-vs-seam}
The accepted-record invariance clause in \cref{def:witness-seam-realization} encodes the physical
indistinguishability of filler toggle at the accepted-record layer.
The seam-scalar theorems below use only the induced odd sign relations in the source channel.
\end{remark}

\begin{definition}[Raw and normalized seam pairings]
\label{def:raw-seam-pairing}
Fix an accepted tick \(k\) and a witness-indexed seam realization as in
\cref{def:witness-seam-realization}.
For \(f\in\{\mathrm{diag},\mathrm{bridge}\}\) and witness \(y\), define the \emph{raw seam
pairing} by
\begin{equation}
\widetilde s_{\Lambda_k}(f,y)
:=
\big\langle\!\big\langle
\Pi_{\mathrm{phys},\Lambda_k}\mathsf T_{\Lambda_k}\!\big[\tau^{f,y}_{\Lambda_k}\big],\,
\delta C^{\mathrm{LoG}}_{\Lambda_k}
\big\rangle\!\big\rangle_{C_{\ast,\Lambda_k}}.
\label{eq:raw-seam-pairing}
\end{equation}
Whenever dominant-line realization holds so that
\(
\widehat{\delta C}^{\mathrm{LoG}}_{\Lambda_k}
\)
is defined, also set
\begin{equation}
s_{\Lambda_k}(f,y)
:=
\big\langle\!\big\langle
\Pi_{\mathrm{phys},\Lambda_k}\mathsf T_{\Lambda_k}\!\big[\tau^{f,y}_{\Lambda_k}\big],\,
\widehat{\delta C}^{\mathrm{LoG}}_{\Lambda_k}
\big\rangle\!\big\rangle_{C_{\ast,\Lambda_k}}.
\label{eq:normalized-seam-pairing}
\end{equation}
\end{definition}

\begin{proposition}[Filler and witness parity of the seam scalar]
\label{prop:seam-scalar-signs}
Fix an accepted tick \(k\) and assume a witness-indexed seam realization at \(\Lambda_k\).
Then the raw seam pairing satisfies
\begin{equation}
\widetilde s_{\Lambda_k}(\mathrm{bridge},y)
=
-\,\widetilde s_{\Lambda_k}(\mathrm{diag},y),
\qquad
\widetilde s_{\Lambda_k}(f,\mathrm{flip}(y))
=
-\,\widetilde s_{\Lambda_k}(f,y).
\label{eq:raw-seam-signs}
\end{equation}
If dominant-line realization also holds, then the same identities hold for the normalized seam
pairing:
\begin{equation}
s_{\Lambda_k}(\mathrm{bridge},y)
=
-\,s_{\Lambda_k}(\mathrm{diag},y),
\qquad
s_{\Lambda_k}(f,\mathrm{flip}(y))
=
-\,s_{\Lambda_k}(f,y).
\label{eq:normalized-seam-signs}
\end{equation}
\end{proposition}

\begin{proof}
The map
\(
\tau\mapsto \Pi_{\mathrm{phys},\Lambda_k}\mathsf T_{\Lambda_k}[\tau]
\)
is linear because \(\mathsf T_{\Lambda_k}\) and \(\Pi_{\mathrm{phys},\Lambda_k}\) are linear.
The Fisher pairing is bilinear.
Applying the sign relations of \cref{def:witness-seam-realization} therefore gives
\[
\Pi_{\mathrm{phys},\Lambda_k}\mathsf T_{\Lambda_k}\!\big[\tau^{\mathrm{bridge},y}_{\Lambda_k}\big]
=
-\,\Pi_{\mathrm{phys},\Lambda_k}\mathsf T_{\Lambda_k}\!\big[\tau^{\mathrm{diag},y}_{\Lambda_k}\big]
\]
and
\[
\Pi_{\mathrm{phys},\Lambda_k}\mathsf T_{\Lambda_k}\!\big[\tau^{f,\mathrm{flip}(y)}_{\Lambda_k}\big]
=
-\,\Pi_{\mathrm{phys},\Lambda_k}\mathsf T_{\Lambda_k}\!\big[\tau^{f,y}_{\Lambda_k}\big].
\]
Pairing with \(\delta C^{\mathrm{LoG}}_{\Lambda_k}\) proves \eqref{eq:raw-seam-signs}.
If dominant-line realization holds, then
\(
\widehat{\delta C}^{\mathrm{LoG}}_{\Lambda_k}
\)
is a fixed nonzero normalization of the same vector, so the same argument proves
\eqref{eq:normalized-seam-signs}.
\end{proof}

\begin{definition}[Detector-factorized seam coupling]
\label{def:detector-factorized-seam}
Fix an accepted tick \(k\).
A witness-indexed seam realization is called \emph{detector-factorized} if there exists a single
admissible welded odd trace-channel source
\[
\tau^{\mathrm{unit}}_{\Lambda_k}
\]
and the LS detector factor
\[
\overline{\mathrm{obs}}
:
\text{LS difference object}\to\mathbb R
\]
such that, for every filler tag \(f\) and witness \(y\),
\begin{equation}
\tau^{f,y}_{\Lambda_k}
=
\sigma_{\mathrm{fill}}(f)\,
\overline{\mathrm{obs}}(y)\,
\tau^{\mathrm{unit}}_{\Lambda_k},
\qquad
\sigma_{\mathrm{fill}}(\mathrm{diag})=+1,\quad
\sigma_{\mathrm{fill}}(\mathrm{bridge})=-1.
\label{eq:detector-factorized-seam}
\end{equation}
\end{definition}

\begin{theorem}[Seam faithfulness scalar]
\label{thm:seam-faithfulness-scalar}
Fix an accepted tick \(k\) and assume a detector-factorized witness-indexed seam realization at
\(\Lambda_k\).
Define the fixed-cutoff seam coefficient by
\begin{equation}
\widetilde c_{\Lambda_k}
:=
\big\langle\!\big\langle
\Pi_{\mathrm{phys},\Lambda_k}\mathsf T_{\Lambda_k}\!\big[\tau^{\mathrm{unit}}_{\Lambda_k}\big],\,
\delta C^{\mathrm{LoG}}_{\Lambda_k}
\big\rangle\!\big\rangle_{C_{\ast,\Lambda_k}}.
\label{eq:seam-faithfulness-coeff}
\end{equation}
Then, for every filler tag \(f\) and witness \(y\),
\begin{equation}
\widetilde s_{\Lambda_k}(f,y)
=
\sigma_{\mathrm{fill}}(f)\,
\widetilde c_{\Lambda_k}\,
\overline{\mathrm{obs}}(y).
\label{eq:raw-seam-factorization}
\end{equation}
In particular, if
\[
\widetilde c_{\Lambda_k}\neq 0
\qquad\text{and}\qquad
\overline{\mathrm{obs}}(y)\neq 0,
\]
then the raw seam pairing is nonzero:
\[
\widetilde s_{\Lambda_k}(f,y)\neq 0.
\]
Consequently the compressed welded LoG line is transversely physical at \(\Lambda_k\), hence
dominant-line realization holds at that cutoff by \cref{prop:dominant-line-equivalent}.
When dominant-line realization holds, the normalized seam coefficient
\begin{equation}
c_{\Lambda_k}
:=
\kappa_{\Lambda_k}^{-1/2}\,\widetilde c_{\Lambda_k}
\label{eq:normalized-seam-coeff}
\end{equation}
is defined and
\begin{equation}
s_{\Lambda_k}(f,y)
=
\sigma_{\mathrm{fill}}(f)\,
c_{\Lambda_k}\,
\overline{\mathrm{obs}}(y).
\label{eq:normalized-seam-factorization}
\end{equation}
\end{theorem}

\begin{proof}
Substitute the detector-factorized form \eqref{eq:detector-factorized-seam} into the definition
\eqref{eq:raw-seam-pairing}.
Linearity of \(\mathsf T_{\Lambda_k}\) and bilinearity of the Fisher pairing give
\[
\widetilde s_{\Lambda_k}(f,y)
=
\sigma_{\mathrm{fill}}(f)\,\overline{\mathrm{obs}}(y)\,
\big\langle\!\big\langle
\Pi_{\mathrm{phys},\Lambda_k}\mathsf T_{\Lambda_k}\!\big[\tau^{\mathrm{unit}}_{\Lambda_k}\big],\,
\delta C^{\mathrm{LoG}}_{\Lambda_k}
\big\rangle\!\big\rangle_{C_{\ast,\Lambda_k}},
\]
which is exactly \eqref{eq:raw-seam-factorization}.

If \(\widetilde c_{\Lambda_k}\neq 0\) and \(\overline{\mathrm{obs}}(y)\neq 0\), then
\(
\widetilde s_{\Lambda_k}(f,y)\neq 0
\).
Set
\[
B_{\Lambda_k}
:=
\Pi_{\mathrm{phys},\Lambda_k}\mathsf T_{\Lambda_k}\!\big[\tau^{f,y}_{\Lambda_k}\big].
\]
Then
\[
\big\langle\!\big\langle
\delta C^{\mathrm{LoG}}_{\Lambda_k},\,B_{\Lambda_k}
\big\rangle\!\big\rangle_{C_{\ast,\Lambda_k}}
=
\widetilde s_{\Lambda_k}(f,y)
\neq 0.
\]
By item (4) of \cref{prop:dominant-line-equivalent}, this proves transversality of the compressed
line and hence dominant-line realization at \(\Lambda_k\).

Once dominant-line realization holds, divide \eqref{eq:raw-seam-factorization} by
\(
\|\delta C^{\mathrm{LoG}}_{\Lambda_k}\|_{C_{\ast,\Lambda_k}}
=
\kappa_{\Lambda_k}^{1/2}
\)
to obtain \eqref{eq:normalized-seam-factorization}.
\end{proof}

\begin{remark}[Downstream role of the seam-coefficient hypothesis]
The fixed-cutoff nonvanishing hypothesis
\(
\widetilde c_{\Lambda_k}\neq 0
\)
is exactly the input that, once dominant-line realization defines
\(
c_{\Lambda_k}
\),
is imported downstream as the nondegeneracy clause
\(
c_{\Lambda_k}\neq 0
\)
inside the delayed selected-line package.
This note does not derive that clause from split-threshold analysis alone.
\end{remark}

\begin{definition}[Odd source representability on the characterization layer]
\label{def:odd-source-representability}
Fix an accepted tick \(k\).
Let \(\tau\) be an admissible welded odd trace-channel source.
We say that \(\tau\) is \emph{representable on the characterization layer} if there exists an odd
sufficiently local generator
\[
F_{\tau,\Lambda_k}\in E^-_{\Lambda_k}
\]
such that its intrinsic physical covariance response agrees with the UOM welded response:
\begin{equation}
\Pi_{\mathrm{phys},\Lambda_k}\mathcal R_{\Lambda_k}\!\big(F_{\tau,\Lambda_k}\big)
=
\Pi_{\mathrm{phys},\Lambda_k}\mathsf T_{\Lambda_k}[\tau].
\label{eq:odd-source-representability}
\end{equation}
This is the strong source-insertion realization statement.
\end{definition}

\begin{theorem}[UOM source-insertion representability theorem]
\label{thm:odd-source-representability}
Fix an accepted tick \(k\).
Let \(\tau\) be an admissible welded odd trace-channel source at cutoff \(\Lambda_k\), in the
finite-cutoff UOM sense: its receiver-head profile is band-limited at \(\Lambda_k\) and compactly
supported inside the accepted slab.
Then there exists a canonically associated odd sufficiently local generator
\[
F_{\tau,\Lambda_k}\in E^-_{\Lambda_k}
\]
such that
\begin{equation}
\Pi_{\mathrm{phys},\Lambda_k}\mathcal R_{\Lambda_k}\!\big(F_{\tau,\Lambda_k}\big)
=
\Pi_{\mathrm{phys},\Lambda_k}\mathsf T_{\Lambda_k}[\tau].
\label{eq:odd-source-representability-thm}
\end{equation}
In particular, every witness-indexed seam source used in the present note is representable on the
characterization layer.
\end{theorem}

\begin{proof}
By the finite-cutoff UOM Peierls package, admissible receiver-head sources at cutoff \(\Lambda_k\)
are band-limited boundary test data and the sufficiently local class
\(\mathcal F_{\rm s.loc}^{\Lambda_k}\) contains all compact smearings of local boundary densities.
Let
\[
F_{\tau,\Lambda_k}
\]
be the linear source functional obtained by pairing the odd trace-channel source profile \(\tau\)
with the corresponding odd trace-channel boundary field on the receiver head.
Because this is the linear case of the sufficiently local form and \(\tau\) is physically odd,
\[
F_{\tau,\Lambda_k}\in E^-_{\Lambda_k}.
\]

Now apply the UOM source insertion map from the finite-cutoff retarded-kernel package.
The source \(\tau\) is sent to the bulk forcing term in the linearized equation; the unique
retarded solution is obtained by the retarded Green operator; and boundary restriction gives the
retarded covariance variation. By the UOM identification of the induced retarded kernel with the SK
mixed second derivative, this covariance variation is exactly the welded source-to-covariance
response
\[
\mathsf T_{\Lambda_k}[\tau].
\]

On the other hand, by the definition of the intrinsic covariance-response map, the matrix elements
of
\(
\mathcal R_{\Lambda_k}(F_{\tau,\Lambda_k})
\)
are computed by the Peierls action of \(F_{\tau,\Lambda_k}\) on the sufficiently local quadratic
probes representing boundary covariance entries.
That Peierls variation is precisely the same retarded second-moment variation generated by the
source insertion \(\tau\).
Therefore
\[
\mathcal R_{\Lambda_k}(F_{\tau,\Lambda_k})
=
\mathsf T_{\Lambda_k}[\tau]
\]
as covariance tangents.
Applying the physical projector \(\Pi_{\mathrm{phys},\Lambda_k}\) gives
\eqref{eq:odd-source-representability-thm}.

The witness-indexed seam sources of \cref{def:witness-seam-realization} are, by construction, live
admissible welded odd trace-channel sources at the accepted tick, so the same argument applies to
each of them.
\end{proof}

\begin{proposition}[Dominant pairing for a representable odd source]
\label{prop:odd-source-pairing}
Fix an accepted tick \(k\) and let \(\tau\) be an admissible welded odd trace-channel source that is
representable in the sense of \cref{thm:odd-source-representability}.
Then
\begin{equation}
\big\langle\!\big\langle
\Pi_{\mathrm{phys},\Lambda_k}\mathcal R_{\Lambda_k}\!\big(F_{\tau,\Lambda_k}\big),\,
\delta C^{\mathrm{LoG}}_{\Lambda_k}
\big\rangle\!\big\rangle_{C_{\ast,\Lambda_k}}
=
\big\langle\!\big\langle
\Pi_{\mathrm{phys},\Lambda_k}\mathsf T_{\Lambda_k}[\tau],\,
\delta C^{\mathrm{LoG}}_{\Lambda_k}
\big\rangle\!\big\rangle_{C_{\ast,\Lambda_k}}.
\label{eq:representable-raw-pairing}
\end{equation}
If dominant-line realization also holds, then
\begin{equation}
\rho^{\mathrm{dom}}_{\Lambda_k}\!\big(F_{\tau,\Lambda_k}\big)
=
\big\langle\!\big\langle
\Pi_{\mathrm{phys},\Lambda_k}\mathsf T_{\Lambda_k}[\tau],\,
\widehat{\delta C}^{\mathrm{LoG}}_{\Lambda_k}
\big\rangle\!\big\rangle_{C_{\ast,\Lambda_k}}.
\label{eq:representable-normalized-pairing}
\end{equation}
\end{proposition}

\begin{proof}
Equation \eqref{eq:odd-source-representability} identifies the two physical covariance tangents.
Pairing both sides with \(\delta C^{\mathrm{LoG}}_{\Lambda_k}\) gives
\eqref{eq:representable-raw-pairing}.
If dominant-line realization holds, pair with the normalized vector
\(
\widehat{\delta C}^{\mathrm{LoG}}_{\Lambda_k}
\)
instead and use the definition \eqref{eq:rho-dom-def} of
\(
\rho^{\mathrm{dom}}_{\Lambda_k}
\)
to obtain \eqref{eq:representable-normalized-pairing}.
\end{proof}

\begin{corollary}[Unit welded odd source representability and scalar identification]
\label{cor:unit-source-representability}
Fix an accepted tick \(k\).
Assume a detector-factorized witness-indexed seam realization at \(\Lambda_k\) in the sense of
\cref{def:detector-factorized-seam}, and assume dominant-line realization at \(\Lambda_k\).
Then the unit welded odd source
\(
\tau^{\mathrm{unit}}_{\Lambda_k}
\)
is representable on the characterization layer, with canonical representable generator
\[
F_{\tau^{\mathrm{unit}}_{\Lambda_k},\Lambda_k}\in E^-_{\Lambda_k},
\]
and
\begin{equation}
\rho^{\mathrm{dom}}_{\Lambda_k}
\!\bigl(
F_{\tau^{\mathrm{unit}}_{\Lambda_k},\Lambda_k}
\bigr)
=
\big\langle\!\big\langle
\Pi_{\mathrm{phys},\Lambda_k}\mathsf T_{\Lambda_k}\!\big[\tau^{\mathrm{unit}}_{\Lambda_k}\big],\,
\widehat{\delta C}^{\mathrm{LoG}}_{\Lambda_k}
\big\rangle\!\big\rangle_{C_{\ast,\Lambda_k}}
=
c_{\Lambda_k}.
\label{eq:unit-source-representability-scalar}
\end{equation}
\end{corollary}

\begin{proof}
Representability of
\(
\tau^{\mathrm{unit}}_{\Lambda_k}
\)
is the specialization of \cref{thm:odd-source-representability} to the admissible welded odd trace
source singled out by \cref{def:detector-factorized-seam}.
Because dominant-line realization holds, \cref{prop:odd-source-pairing} applies to that
representable source and gives the first equality in \eqref{eq:unit-source-representability-scalar}.
The second equality is exactly the definition of the normalized seam coefficient
\(
c_{\Lambda_k}
\)
from \cref{thm:seam-faithfulness-scalar}.
\end{proof}

\begin{definition}[Scalar comparison shadow of the witness-indexed bridge]
\label{def:witness-descendant-comparison}
Fix an accepted tick \(k\).
Let \(y\) be an LS witness and let
\(
\tau^{f,y}_{\Lambda_k}
\)
be a witness-indexed seam realization source for a filler tag \(f\).
A \emph{scalar comparison shadow} for that data is a choice of representative
\[
J^{f,y}_{\Lambda_k}\in[\Xi^{(1)}_{\Lambda_k}]
\]
such that the dominant raw pairing of the 2T-side representative agrees with the dominant raw
pairing of the corresponding representable UOM source:
\begin{equation}
\big\langle\!\big\langle
\Pi_{\mathrm{phys},\Lambda_k}\mathcal R_{\Lambda_k}\!\big(J^{f,y}_{\Lambda_k}\big),\,
\delta C^{\mathrm{LoG}}_{\Lambda_k}
\big\rangle\!\big\rangle_{C_{\ast,\Lambda_k}}
=
\big\langle\!\big\langle
\Pi_{\mathrm{phys},\Lambda_k}\mathcal R_{\Lambda_k}\!\big(F_{\tau^{f,y}_{\Lambda_k},\Lambda_k}\big),\,
\delta C^{\mathrm{LoG}}_{\Lambda_k}
\big\rangle\!\big\rangle_{C_{\ast,\Lambda_k}}.
\label{eq:witness-descendant-comparison}
\end{equation}
When dominant-line realization holds, this is equivalently the equality of normalized dominant
pairings
\begin{equation}
\rho^{\mathrm{dom}}_{\Lambda_k}\!\big(J^{f,y}_{\Lambda_k}\big)
=
\rho^{\mathrm{dom}}_{\Lambda_k}\!\big(F_{\tau^{f,y}_{\Lambda_k},\Lambda_k}\big).
\label{eq:witness-descendant-comparison-normalized}
\end{equation}
\end{definition}

\begin{definition}[Witness-indexed source-compatible descendant realization]
\label{def:source-compatible-witness-realization}
Fix an accepted tick \(k\), a filler tag \(f\), a witness \(y\), and a witness-indexed seam source
\(
\tau^{f,y}_{\Lambda_k}
\).
A representative
\[
J^{f,y}_{\Lambda_k}\in[\Xi^{(1)}_{\Lambda_k}]
\]
is called a \emph{witness-indexed source-compatible descendant realization} if it realizes, on the
same LS witness datum, the Limited Scope output and odd Hamiltonization clauses from item
\textbf{(M1)} and item \textbf{(M4)} of Definition~\ref{main-def:uom-main-theorem-assumptions} and differs
from the source-insertion representative only by an ambiguity direction:
\begin{equation}
J^{f,y}_{\Lambda_k}
-
F_{\tau^{f,y}_{\Lambda_k},\Lambda_k}
\in
\mathcal A^-_{\Lambda_k},
\label{eq:source-compatible-witness-realization}
\end{equation}
where
\(
F_{\tau^{f,y}_{\Lambda_k},\Lambda_k}
\)
is the representable source generator furnished by
\cref{thm:odd-source-representability}.
Equivalently, the witness-indexed transported descendant class and the witness-indexed UOM
source-insertion realization define the same point of the ambiguity quotient
\(
E^-_{\Lambda_k}/\mathcal A^-_{\Lambda_k}
\).
\end{definition}
\begin{theorem}[Same-witness descended-scalar comparison from a source-compatible descendant realization]
\label{thm:witness-descended-scalar-comparison}
Fix an accepted tick \(k\), assume dominant-line realization at \(\Lambda_k\), and let
\(
\tau^{f,y}_{\Lambda_k}
\)
be a witness-indexed seam realization source.
Assume:
\begin{enumerate}[leftmargin=2em]
\item the source \(\tau^{f,y}_{\Lambda_k}\) is represented by the canonical generator
  \(
  F_{\tau^{f,y}_{\Lambda_k},\Lambda_k}
  \)
  from \cref{thm:odd-source-representability};
\item there exists a witness-indexed source-compatible descendant realization
  \(
  J^{f,y}_{\Lambda_k}\in[\Xi^{(1)}_{\Lambda_k}]
  \)
  in the sense of \cref{def:source-compatible-witness-realization};
\item the three vanishing statements of \cref{prop:dominant-quotient-generators} hold at
  \(\Lambda_k\).
\end{enumerate}
Then the raw and normalized comparison identities of
\cref{def:witness-descendant-comparison} hold:
\begin{align}
\big\langle\!\big\langle
\Pi_{\mathrm{phys},\Lambda_k}\mathcal R_{\Lambda_k}\!\big(J^{f,y}_{\Lambda_k}\big),\,
\delta C^{\mathrm{LoG}}_{\Lambda_k}
\big\rangle\!\big\rangle_{C_{\ast,\Lambda_k}}
&=
\big\langle\!\big\langle
\Pi_{\mathrm{phys},\Lambda_k}\mathcal R_{\Lambda_k}\!\big(F_{\tau^{f,y}_{\Lambda_k},\Lambda_k}\big),\,
\delta C^{\mathrm{LoG}}_{\Lambda_k}
\big\rangle\!\big\rangle_{C_{\ast,\Lambda_k}},
\label{eq:witness-descendant-comparison-thm-raw}
\\
\rho^{\mathrm{dom}}_{\Lambda_k}\!\big(J^{f,y}_{\Lambda_k}\big)
&=
\rho^{\mathrm{dom}}_{\Lambda_k}\!\big(F_{\tau^{f,y}_{\Lambda_k},\Lambda_k}\big).
\label{eq:witness-descendant-comparison-thm-normalized}
\end{align}
Moreover,
\begin{equation}
\rho^{\mathrm{dom}}_{\Lambda_k}\!\big(F_{\tau^{f,y}_{\Lambda_k},\Lambda_k}\big)
=
\big\langle\!\big\langle
\Pi_{\mathrm{phys},\Lambda_k}\mathsf T_{\Lambda_k}\!\big[\tau^{f,y}_{\Lambda_k}\big],\,
\widehat{\delta C}^{\mathrm{LoG}}_{\Lambda_k}
\big\rangle\!\big\rangle_{C_{\ast,\Lambda_k}}.
\label{eq:witness-descendant-comparison-thm-seam}
\end{equation}
If quotient descent also holds at \(\Lambda_k\), then the common value equals the descended
dominant scalar on the transported descendant affine class:
\begin{equation}
\overline{\rho}^{\mathrm{dom}}_{\Lambda_k}\big([\![\Xi^{(1)}_{\Lambda_k}]\!]\big)
=
\rho^{\mathrm{dom}}_{\Lambda_k}\!\big(J^{f,y}_{\Lambda_k}\big)
=
\rho^{\mathrm{dom}}_{\Lambda_k}\!\big(F_{\tau^{f,y}_{\Lambda_k},\Lambda_k}\big).
\label{eq:witness-descendant-comparison-thm-descended}
\end{equation}
\end{theorem}

\begin{proof}
Set
\[
A^{f,y}_{\Lambda_k}
:=
J^{f,y}_{\Lambda_k}
-
F_{\tau^{f,y}_{\Lambda_k},\Lambda_k}.
\]
By source compatibility,
\(
A^{f,y}_{\Lambda_k}\in\mathcal A^-_{\Lambda_k}
\).

Because dominant-line realization holds and the three vanishing statements of
\cref{prop:dominant-quotient-generators} are assumed, that proposition gives
\[
\mathcal A^-_{\Lambda_k}\subset\ker\rho^{\mathrm{dom}}_{\Lambda_k}.
\]
Hence
\[
\rho^{\mathrm{dom}}_{\Lambda_k}\!\big(J^{f,y}_{\Lambda_k}\big)
=
\rho^{\mathrm{dom}}_{\Lambda_k}\!\big(
F_{\tau^{f,y}_{\Lambda_k},\Lambda_k}
+
A^{f,y}_{\Lambda_k}
\big)
=
\rho^{\mathrm{dom}}_{\Lambda_k}\!\big(F_{\tau^{f,y}_{\Lambda_k},\Lambda_k}\big),
\]
proving \eqref{eq:witness-descendant-comparison-thm-normalized}.

Since dominant-line realization implies
\(
\delta C^{\mathrm{LoG}}_{\Lambda_k}\neq 0
\)
and
\(
\rho^{\mathrm{dom}}_{\Lambda_k}
=
\kappa_{\Lambda_k}^{-1/2}\widetilde{\rho}^{\mathrm{dom}}_{\Lambda_k},
\)
the raw pairing equality \eqref{eq:witness-descendant-comparison-thm-raw} is just the same identity
multiplied by the positive factor \(\kappa_{\Lambda_k}^{1/2}\).

Equation \eqref{eq:witness-descendant-comparison-thm-seam} is exactly the normalized pairing formula
of \cref{prop:odd-source-pairing} applied to the representable source
\(
\tau^{f,y}_{\Lambda_k}
\).

If quotient descent also holds, then \cref{cor:dominant-quotient-constancy} gives
\[
\rho^{\mathrm{dom}}_{\Lambda_k}\!\big(J^{f,y}_{\Lambda_k}\big)
=
\overline{\rho}^{\mathrm{dom}}_{\Lambda_k}\big([\![\Xi^{(1)}_{\Lambda_k}]\!]\big)
\]
for every representative
\(
J^{f,y}_{\Lambda_k}\in[\Xi^{(1)}_{\Lambda_k}],
\)
which yields \eqref{eq:witness-descendant-comparison-thm-descended} after combining with
\eqref{eq:witness-descendant-comparison-thm-normalized}.
\end{proof}

\begin{theorem}[Fixed-cutoff seam witness criterion for descendant overlap]
\label{thm:fixed-cutoff-seam-witness}
Fix an accepted tick \(k\).
Assume:
\begin{enumerate}[leftmargin=2em]
\item a detector-factorized witness-indexed seam realization at \(\Lambda_k\);
\item a witness \(y\) with
  \[
  \overline{\mathrm{obs}}(y)\neq 0;
  \]
\item the seam coefficient is nonzero,
  \[
  \widetilde c_{\Lambda_k}\neq 0;
  \]
\item for the chosen filler tag \(f\), the source \(\tau^{f,y}_{\Lambda_k}\) is representable on the
  characterization layer in the sense of \cref{def:odd-source-representability};
\item there exists a witness-indexed source-compatible descendant realization in the sense of
  \cref{def:source-compatible-witness-realization}.
\end{enumerate}
Then:
\begin{enumerate}[leftmargin=2em]
\item the compressed welded LoG line is transversely physical at \(\Lambda_k\), hence dominant-line
  realization holds at that cutoff;
\item the chosen 2T representative satisfies
  \[
  \rho^{\mathrm{dom}}_{\Lambda_k}\!\big(J^{f,y}_{\Lambda_k}\big)\neq 0;
  \]
\item therefore the transported descendant class has nonzero overlap with the canonically selected
  compressed dominant odd line at \(\Lambda_k\);
\item if quotient descent also holds at \(\Lambda_k\), then
  \cref{thm:descendant-overlap} applies and the transported descendant class is dominant-active with
  well-defined sign \(\varepsilon_{\Lambda_k}\).
\end{enumerate}
\end{theorem}

\begin{proof}
By \cref{thm:seam-faithfulness-scalar},
\[
\widetilde s_{\Lambda_k}(f,y)
=
\sigma_{\mathrm{fill}}(f)\,
\widetilde c_{\Lambda_k}\,
\overline{\mathrm{obs}}(y)
\neq 0.
\]
The same theorem therefore gives item (1): the compressed line is transversely physical, hence
dominant-line realization holds at \(\Lambda_k\).

By \cref{prop:odd-source-pairing},
\[
\big\langle\!\big\langle
\Pi_{\mathrm{phys},\Lambda_k}\mathcal R_{\Lambda_k}\!\big(F_{\tau^{f,y}_{\Lambda_k},\Lambda_k}\big),\,
\delta C^{\mathrm{LoG}}_{\Lambda_k}
\big\rangle\!\big\rangle_{C_{\ast,\Lambda_k}}
=
\widetilde s_{\Lambda_k}(f,y)
\neq 0.
\]
Applying the comparison identity \eqref{eq:witness-descendant-comparison} now gives
\[
\big\langle\!\big\langle
\Pi_{\mathrm{phys},\Lambda_k}\mathcal R_{\Lambda_k}\!\big(J^{f,y}_{\Lambda_k}\big),\,
\delta C^{\mathrm{LoG}}_{\Lambda_k}
\big\rangle\!\big\rangle_{C_{\ast,\Lambda_k}}
\neq 0.
\]
Because dominant-line realization has already been established, \(\delta C^{\mathrm{LoG}}_{\Lambda_k}\)
is nonzero and \(\widehat{\delta C}^{\mathrm{LoG}}_{\Lambda_k}\) is defined.
By \cref{prop:descendant-overlap-witness}, this is equivalent to
\(
\rho^{\mathrm{dom}}_{\Lambda_k}(J^{f,y}_{\Lambda_k})\neq 0
\),
proving item (2).

Item (3) is exactly the nonzero-overlap condition from
\cref{prop:descendant-overlap-equivalent}.
If quotient descent is also available, apply \cref{thm:descendant-overlap} to obtain item (4).
\end{proof}

\subsection*{Smooth-taper branch continuity and witness transport}
\label{subsec:smooth-taper-branch-continuity}

The previous fixed-cutoff theorem isolates the nonzero witness needed for activity.
To transport that witness along an accepted branch one needs a continuity statement for the welded
response.
At theorem level, the clean analytic choice is the smooth receiver taper.
The hard projector may still be used, but only on branch segments where its rank is constant.

\begin{lemma}[Smooth-taper continuity of the receiver band projector]
\label{lem:smooth-taper-projector}
Let \(I\) be a connected accepted-branch segment.
Suppose the receiver band is implemented by a smooth taper
\[
P^{\mathrm{(rec)}}_{\mathrm{band}}(\Lambda)
=
\sum_{\ell,m} w_\ell(\Lambda)\,|Y_{\ell m}\rangle\langle Y_{\ell m}|,
\qquad
0\le w_\ell(\Lambda)\le 1,
\]
on the receiver one-particle space, where
\[
\sup_{\ell\ge 0}|w_\ell(\Lambda')-w_\ell(\Lambda)|
\longrightarrow 0
\qquad
\text{as }\Lambda'\to\Lambda
\]
for every \(\Lambda\in I\).
Then
\[
\Lambda\longmapsto P^{\mathrm{(rec)}}_{\mathrm{band}}(\Lambda)
\]
is continuous on \(I\) in operator norm.
If instead the hard projector
\(
\sum_{\ell\le \ell_\star(\Lambda)}|Y_{\ell m}\rangle\langle Y_{\ell m}|
\)
is used, the same conclusion holds on every connected subsegment of \(I\) on which
\(\ell_\star(\Lambda)\) is constant.
\end{lemma}

\begin{proof}
In the spherical-harmonic basis \(\{Y_{\ell m}\}\), the difference
\[
P^{\mathrm{(rec)}}_{\mathrm{band}}(\Lambda')
-
P^{\mathrm{(rec)}}_{\mathrm{band}}(\Lambda)
\]
is diagonal with eigenvalues \(w_\ell(\Lambda')-w_\ell(\Lambda)\).
Therefore
\[
\big\|
P^{\mathrm{(rec)}}_{\mathrm{band}}(\Lambda')
-
P^{\mathrm{(rec)}}_{\mathrm{band}}(\Lambda)
\big\|_{\mathrm{op}}
=
\sup_{\ell\ge 0}|w_\ell(\Lambda')-w_\ell(\Lambda)|.
\]
The assumed uniform continuity of the taper coefficients is exactly the statement that this operator
norm tends to zero as \(\Lambda'\to\Lambda\).

For the hard projector, the coefficients take only the values \(0\) and \(1\), so discontinuities can
occur only when \(\ell_\star(\Lambda)\) jumps.
On a subsegment where \(\ell_\star(\Lambda)\) is constant, the projector is constant and therefore
continuous.
\end{proof}

\begin{theorem}[Smooth-taper branch continuity of the welded physical response]
\label{thm:smooth-taper-welded-response}
Let \(I\) be a connected accepted-branch segment.
Write the welded source-to-covariance map in the UOM factorized form
\begin{equation}
\mathsf T_\Lambda
=
\Pi_{L(\Lambda)}M_{\varepsilon,\Lambda}
\circ
\mathsf K_{{\rm comp},\Lambda}
\circ
\mathsf G_{{\rm ret},\Lambda}
\circ
\mathsf W_{{\rm src},\Lambda}.
\label{eq:factorized-welded-map}
\end{equation}
Assume:
\begin{enumerate}[leftmargin=2em]
\item the receiver band is implemented by a smooth taper satisfying
  \cref{lem:smooth-taper-projector};
\item the operator families
  \[
  \Lambda\longmapsto
  \Pi_{L(\Lambda)}M_{\varepsilon,\Lambda},\qquad
  \Lambda\longmapsto
  \mathsf K_{{\rm comp},\Lambda},\qquad
  \Lambda\longmapsto
  \mathsf G_{{\rm ret},\Lambda},\qquad
  \Lambda\longmapsto
  \mathsf W_{{\rm src},\Lambda},\qquad
  \Lambda\longmapsto
  \Pi_{\mathrm{phys},\Lambda}
  \]
  are continuous on \(I\) in operator norm on the corresponding band-limited source and covariance
  tangent spaces;
\item the source family
  \[
  \Lambda\longmapsto \tau_\Lambda
  \]
  is continuous on \(I\) in the welded source norm.
\end{enumerate}
Then the physical covariance response
\begin{equation}
X_\Lambda[\tau_\Lambda]
:=
\Pi_{\mathrm{phys},\Lambda}\mathsf T_\Lambda[\tau_\Lambda]
\label{eq:XLambda-tau-def}
\end{equation}
depends continuously on \(\Lambda\in I\) in the Fisher norm.
If the hard receiver projector is used instead, the same conclusion holds on each connected
subsegment where \(\ell_\star(\Lambda)\) is constant.
\end{theorem}

\begin{proof}
By \cref{lem:smooth-taper-projector}, the first factor
\(
\Pi_{L(\Lambda)}M_{\varepsilon,\Lambda}
\)
is operator-norm continuous along the branch.
By assumption, the remaining three UOM factors in \eqref{eq:factorized-welded-map} and the physical
projector \(\Pi_{\mathrm{phys},\Lambda}\) are also operator-norm continuous.
Hence the composed operator family
\[
\Lambda\longmapsto
\Pi_{\mathrm{phys},\Lambda}\mathsf T_\Lambda
\]
is continuous in operator norm on \(I\).
Applying this continuous family to the continuous source vector \(\tau_\Lambda\) gives continuity of
\(
X_\Lambda[\tau_\Lambda]
\)
in the Fisher norm.

For the hard projector, the same argument applies on every connected subsegment where
\(\ell_\star(\Lambda)\) is constant, because there the first factor is constant.
\end{proof}

\begin{corollary}[Branchwise continuity of the seam coefficient]
\label{cor:branchwise-seam-coeff}
Let \(I\) be a connected accepted-branch segment and assume:
\begin{enumerate}[leftmargin=2em]
\item the hypotheses of \cref{thm:dominant-line-persistence} on \(I\);
\item a detector-factorized branchwise seam realization on \(I\) with continuous unit source
  \[
  \Lambda\longmapsto \tau^{\mathrm{unit}}_{\Lambda};
  \]
\item the hypotheses of \cref{thm:smooth-taper-welded-response} for that unit source.
\end{enumerate}
Define the branchwise normalized seam coefficient by
\begin{equation}
c_\Lambda
:=
\big\langle\!\big\langle
\Pi_{\mathrm{phys},\Lambda}\mathsf T_\Lambda\!\big[\tau^{\mathrm{unit}}_{\Lambda}\big],\,
\widehat{\delta C}^{\mathrm{LoG}}_{\Lambda}
\big\rangle\!\big\rangle_{C_{\ast,\Lambda}}.
\label{eq:branchwise-seam-coeff}
\end{equation}
Then \(c_\Lambda\) is continuous on \(I\).
Consequently, for every fixed filler tag \(f\) and witness \(y\), the normalized seam scalar
\[
s_\Lambda(f,y)
=
\sigma_{\mathrm{fill}}(f)\,
c_\Lambda\,
\overline{\mathrm{obs}}(y)
\]
is continuous on \(I\).
\end{corollary}

\begin{proof}
By \cref{thm:smooth-taper-welded-response}, the response vector
\[
\Lambda\longmapsto
\Pi_{\mathrm{phys},\Lambda}\mathsf T_\Lambda\!\big[\tau^{\mathrm{unit}}_{\Lambda}\big]
\]
is continuous in the Fisher norm.
By \cref{thm:dominant-line-persistence}, the normalized dominant vector
\(
\widehat{\delta C}^{\mathrm{LoG}}_{\Lambda}
\)
is also continuous on \(I\).
Pairing these two continuous vectors in the Fisher metric gives continuity of \(c_\Lambda\).
The final statement is then immediate from
\eqref{eq:normalized-seam-factorization}.
\end{proof}

\begin{theorem}[Branchwise seam witness transport]
\label{thm:branchwise-seam-witness}
Let \(I\) be a connected accepted-branch segment.
Assume:
\begin{enumerate}[leftmargin=2em]
\item the hypotheses of \cref{cor:branchwise-seam-coeff};
\item quotient descent holds for every \(\Lambda\in I\);
\item there exists a fixed LS witness \(y\) with
  \[
  \overline{\mathrm{obs}}(y)\neq 0;
  \]
\item for some fixed filler tag \(f\), there is a branchwise family of representatives
  \[
  J^{f,y}_{\Lambda}\in[\Xi^{(1)}_{\Lambda}]
  \qquad
  (\Lambda\in I)
  \]
  such that, for every \(\Lambda\in I\), the corresponding seam source
  \(\tau^{f,y}_{\Lambda}\) is representable on the characterization layer and the 2T-side
  witness-to-descendant comparison identity
  \eqref{eq:witness-descendant-comparison-normalized} holds.
\end{enumerate}
Then:
\begin{enumerate}[leftmargin=2em]
\item the branchwise dominant-response scalar of the chosen descendant representatives is
  \begin{equation}
  \rho^{\mathrm{dom}}_{\Lambda}\!\big(J^{f,y}_{\Lambda}\big)
  =
  \sigma_{\mathrm{fill}}(f)\,
  c_\Lambda\,
  \overline{\mathrm{obs}}(y);
  \label{eq:branchwise-seam-descendant-scalar}
  \end{equation}
\item the scalar function
  \(
  \Lambda\mapsto \rho^{\mathrm{dom}}_{\Lambda}(J^{f,y}_{\Lambda})
  \)
  is continuous on \(I\);
\item on every connected subsegment \(I'\subset I\) on which \(c_\Lambda\neq 0\), the sign
  \[
  \operatorname{sgn}\!\big(
  \rho^{\mathrm{dom}}_{\Lambda}(J^{f,y}_{\Lambda})
  \big)
  \]
  is constant;
\item because quotient descent holds, the same constant sign is the transported orientation sign of
  the descendant affine class on \(I'\).
\end{enumerate}
\end{theorem}

\begin{proof}
For each \(\Lambda\in I\), representability of \(\tau^{f,y}_{\Lambda}\) and
\cref{prop:odd-source-pairing} give
\[
\rho^{\mathrm{dom}}_{\Lambda}\!\big(F_{\tau^{f,y}_{\Lambda},\Lambda}\big)
=
\big\langle\!\big\langle
\Pi_{\mathrm{phys},\Lambda}\mathsf T_\Lambda\!\big[\tau^{f,y}_{\Lambda}\big],\,
\widehat{\delta C}^{\mathrm{LoG}}_{\Lambda}
\big\rangle\!\big\rangle_{C_{\ast,\Lambda}}
=
s_\Lambda(f,y).
\]
Applying the branchwise comparison identity
\eqref{eq:witness-descendant-comparison-normalized} then yields
\[
\rho^{\mathrm{dom}}_{\Lambda}\!\big(J^{f,y}_{\Lambda}\big)
=
s_\Lambda(f,y).
\]
Now use \cref{cor:branchwise-seam-coeff} to substitute
\(
s_\Lambda(f,y)
=
\sigma_{\mathrm{fill}}(f)c_\Lambda\overline{\mathrm{obs}}(y)
\),
which proves item (1).
Item (2) follows because \(c_\Lambda\) is continuous and
\(\overline{\mathrm{obs}}(y)\) is a fixed nonzero scalar.
On any connected subsegment where \(c_\Lambda\neq 0\), the right-hand side of
\eqref{eq:branchwise-seam-descendant-scalar} is a continuous nonvanishing real-valued function, so
its sign is constant.
This proves item (3).

Finally, quotient descent implies that \(\rho^{\mathrm{dom}}_{\Lambda}\) is constant on
\([\Xi^{(1)}_{\Lambda}]\), so the sign of one representative is the sign of the whole descendant
class.
This gives item (4).
\end{proof}

\begin{corollary}[Seam package discharge of the minimal scalar transport input]
\label{cor:branchwise-seam-minimal-input}
Let \(I\) be a connected accepted-branch segment and assume the hypotheses of
\cref{thm:branchwise-seam-witness}.
Then the branchwise family
\[
J^{f,y}_{\Lambda}\in[\Xi^{(1)}_{\Lambda}]
\qquad
(\Lambda\in I)
\]
satisfies
\[
\Lambda\longmapsto
\rho^{\mathrm{dom}}_{\Lambda}\!\big(J^{f,y}_{\Lambda}\big)
=
\sigma_{\mathrm{fill}}(f)\,
c_\Lambda\,
\overline{\mathrm{obs}}(y)
\]
continuously on \(I\).
On every connected subsegment \(I'\subset I\) such that
\[
c_\Lambda\neq 0
\qquad
\forall \Lambda\in I',
\]
assumption \textup{(3)} of \cref{thm:active-response-transport-minimal} holds in its minimal
scalar form with
\[
J^{\mathrm{desc}}_{\Lambda}:=J^{f,y}_{\Lambda}
\qquad
(\Lambda\in I').
\]
If, in addition, the three generator-wise vanishing statements of
\cref{prop:dominant-quotient-generators} hold for every \(\Lambda\in I'\), then
\cref{thm:active-response-transport-minimal} applies on \(I'\).
\end{corollary}

\begin{proof}
Continuity of the displayed scalar function is exactly item (2) of
\cref{thm:branchwise-seam-witness}, together with the explicit formula of item (1).
On \(I'\), the hypotheses
\(
c_\Lambda\neq 0
\)
and
\(
\overline{\mathrm{obs}}(y)\neq 0
\)
imply
\[
\rho^{\mathrm{dom}}_{\Lambda}\!\big(J^{f,y}_{\Lambda}\big)\neq 0
\qquad
\forall \Lambda\in I'.
\]
Together with the continuity statement already proved above, this is exactly assumption
\textup{(3)} of \cref{thm:active-response-transport-minimal} in its minimal scalar form.
The hypotheses of \cref{thm:branchwise-seam-witness} already include those of
\cref{cor:branchwise-seam-coeff}, and hence the dominant-line persistence input of
\cref{thm:active-response-transport-minimal}.
If the three generator-wise vanishing statements also hold on \(I'\), then every hypothesis of
\cref{thm:active-response-transport-minimal} is in force there.
\end{proof}

\begin{theorem}[Finite-cutoff Active Response Transport Theorem]
\label{thm:active-response-transport-minimal}
Let \(I\) be a connected accepted-branch segment.
For each \(\Lambda\in I\), let
\[
[\Xi^{(1)}_{\Lambda}]
:=
J_{0,\Lambda}+\mathcal A^-_{\Lambda}
\subset E^-_{\Lambda}
\]
be the parity-transported odd descendant affine class modulo the ambiguity space of
\cref{def:transport-ambiguity-subspace}.
Assume:
\begin{enumerate}[leftmargin=2em]
\item the hypotheses of \cref{thm:dominant-line-persistence} hold on \(I\);
\item for every \(\Lambda\in I\), the three vanishing statements of
  \cref{prop:dominant-quotient-generators} hold;
\item there exists a branchwise family of representatives
  \[
  J^{\mathrm{desc}}_{\Lambda}\in[\Xi^{(1)}_{\Lambda}]
  \qquad
  (\Lambda\in I)
  \]
  such that
  \[
  \rho^{\mathrm{dom}}_{\Lambda}(J^{\mathrm{desc}}_{\Lambda})
  \neq 0
  \qquad
  \forall \Lambda\in I,
  \]
  and such that the dominant-response scalar
  \[
  \Lambda\longmapsto
  \rho^{\mathrm{dom}}_{\Lambda}(J^{\mathrm{desc}}_{\Lambda})
  \]
  is continuous on \(I\).
\end{enumerate}
Then, for every \(\Lambda\in I\):
\begin{enumerate}[leftmargin=2em]
\item the compressed dominant covariance tangent is nonzero,
  \[
  \delta C^{\mathrm{LoG}}_{\Lambda}\neq 0,
  \]
  so the dominant line \(L^{\mathrm{dom}}_{\Lambda}\) of
  \cref{def:dominant-response-functional} is canonically defined;
\item the dominant-response coefficient \(\rho^{\mathrm{dom}}_{\Lambda}\) annihilates the odd
  ambiguity space,
  \[
  \rho^{\mathrm{dom}}_{\Lambda}(A)=0
  \qquad
  \forall A\in\mathcal A^-_{\Lambda};
  \]
\item the transported descendant affine class carries a unique nonzero dominant-response value
  \[
  \rho^{\mathrm{dom}}_{\Lambda}(J)
  =
  \overline{\rho}^{\mathrm{dom}}_{\Lambda}\big([\![\Xi^{(1)}_{\Lambda}]\!]\big)
  \neq 0
  \qquad
  \forall J\in[\Xi^{(1)}_{\Lambda}],
  \]
  and hence a well-defined sign
  \[
  \varepsilon_{\Lambda}
  :=
  \operatorname{sgn}\!\Big(
  \overline{\rho}^{\mathrm{dom}}_{\Lambda}\big([\![\Xi^{(1)}_{\Lambda}]\!]\big)
  \Big)
  \in\{\pm1\};
  \]
\item the normalized orientation functional
  \begin{equation}
  \chi_{\Lambda}
  :=
  \nu_{\Lambda}^{-1}\rho^{\mathrm{dom}}_{\Lambda}
  \label{eq:chi-from-rho-dom}
  \end{equation}
  is well defined, annihilates \(\mathcal A^-_{\Lambda}\), and satisfies
  \[
  \chi_{\Lambda}(J)=\varepsilon_{\Lambda}
  \qquad
  \forall J\in[\Xi^{(1)}_{\Lambda}];
  \]
\item along \(I\), the sign \(\varepsilon_{\Lambda}\) is locally constant once an initial
  orientation is fixed.
\end{enumerate}
\end{theorem}

\begin{proof}
Item (1) is exactly the conclusion of \cref{thm:dominant-line-persistence}.

By assumption, the three generator-wise vanishing statements of
\cref{prop:dominant-quotient-generators} hold for every \(\Lambda\in I\).
Therefore \cref{prop:dominant-quotient-generators,thm:dominant-quotient-descent} give item (2)
and the existence of the descended functional
\(
\overline{\rho}^{\mathrm{dom}}_{\Lambda}
\)
at each cutoff.

Applying \cref{prop:descendant-overlap-equivalent,thm:descendant-overlap} pointwise to the
nonvanishing scalar assumption
\(
\rho^{\mathrm{dom}}_{\Lambda}(J^{\mathrm{desc}}_{\Lambda})\neq 0
\)
now yields
item (3): the entire descendant affine class has one common nonzero dominant-response value, and
its sign \(\varepsilon_{\Lambda}\) is well defined.

Since \(\nu_{\Lambda}>0\) by item (3), the normalized functional \(\chi_{\Lambda}\) of
\eqref{eq:chi-from-rho-dom} is well defined.
Item (2) implies that \(\chi_{\Lambda}\) annihilates \(\mathcal A^-_{\Lambda}\), and item (3)
gives
\[
\chi_{\Lambda}(J)
=
\frac{\rho^{\mathrm{dom}}_{\Lambda}(J)}{
\left|
\overline{\rho}^{\mathrm{dom}}_{\Lambda}\big([\![\Xi^{(1)}_{\Lambda}]\!]\big)
\right|
}
=
\varepsilon_{\Lambda}
\qquad
\forall J\in[\Xi^{(1)}_{\Lambda}],
\]
proving item (4).

Finally, assumption (3) provides a continuous nonvanishing real-valued branchwise descendant scalar,
while items (1)--(3) above identify that scalar with the transported orientation sign of the full
descendant affine class.
Since
\(
\Lambda\mapsto \rho^{\mathrm{dom}}_{\Lambda}(J^{\mathrm{desc}}_{\Lambda})
\)
is continuous and nonzero on the connected set \(I\), its sign is constant, hence in particular
locally constant.
Using item (3), this is exactly the transported orientation sign of the full descendant affine
class, proving item (5).
\end{proof}

\begin{proposition}[Ambiguity-triviality on the compressed welded image]
\label{prop:reduced-welded-ambiguity-trivial}
Assume the hypotheses of \cref{thm:active-response-transport-minimal} at an accepted tick \(k\).
Define the reduced welded image of an odd generator \(F\in E^-_{\Lambda_k}\) by
\[
\mathsf W^{\mathrm{red}}_{\Lambda_k}(F)
:=
\rho^{\mathrm{dom}}_{\Lambda_k}(F)\,P_{-,\Lambda_k}^{\mathrm{LoG}},
\]
where \(P_{-,\Lambda_k}^{\mathrm{LoG}}\) is the canonical compressed odd projector of
\cref{cor:compressed-export-selection}.
Then
\[
\mathsf W^{\mathrm{red}}_{\Lambda_k}(F+A)
=
\mathsf W^{\mathrm{red}}_{\Lambda_k}(F)
\qquad
\forall F\in E^-_{\Lambda_k},\ \forall A\in\mathcal A^-_{\Lambda_k}.
\]
Equivalently, \(\mathsf W^{\mathrm{red}}_{\Lambda_k}\) factors through the ambiguity quotient
\(E^-_{\Lambda_k}/\mathcal A^-_{\Lambda_k}\).
In particular, if \(J_1,J_2\in[\Xi^{(1)}_{\Lambda_k}]\), then
\[
\mathsf W^{\mathrm{red}}_{\Lambda_k}(J_1)
=
\mathsf W^{\mathrm{red}}_{\Lambda_k}(J_2).
\]
\end{proposition}

\begin{proof}
Item (2) of \cref{thm:active-response-transport-minimal} gives
\[
\rho^{\mathrm{dom}}_{\Lambda_k}(A)=0
\qquad
\forall A\in\mathcal A^-_{\Lambda_k}.
\]
Therefore, for every \(F\in E^-_{\Lambda_k}\) and \(A\in\mathcal A^-_{\Lambda_k}\),
\[
\mathsf W^{\mathrm{red}}_{\Lambda_k}(F+A)
=
\rho^{\mathrm{dom}}_{\Lambda_k}(F+A)\,P_{-,\Lambda_k}^{\mathrm{LoG}}
=
\big(\rho^{\mathrm{dom}}_{\Lambda_k}(F)+\rho^{\mathrm{dom}}_{\Lambda_k}(A)\big)
P_{-,\Lambda_k}^{\mathrm{LoG}}
=
\rho^{\mathrm{dom}}_{\Lambda_k}(F)\,P_{-,\Lambda_k}^{\mathrm{LoG}}.
\]
Since \cref{cor:compressed-export-selection} fixes \(P_{-,\Lambda_k}^{\mathrm{LoG}}\) canonically on
the compressed, co-centered, band-locked welded channel, the right-hand side depends only on the
ambiguity class of \(F\).
If \(J_2-J_1\in\mathcal A^-_{\Lambda_k}\), apply the same identity with \(F=J_1\) and
\(A=J_2-J_1\).
\end{proof}

\begin{proposition}[Construction of an admissible calibration for the minimal transported interface]
\label{prop:minimal-transport-calibration}
Fix an accepted tick \(k\). Suppose a weak transported-interface package is given:
\[
\mathcal A^-_{\Lambda_k}\subset E^-_{\Lambda_k},
\qquad
[\Xi^{(1)}_{\Lambda_k}]=J_{0,\Lambda_k}+\mathcal A^-_{\Lambda_k},
\]
together with a continuous nonzero linear functional
\[
\chi_{\Lambda_k}:E^-_{\Lambda_k}\to\mathbb R
\]
and a sign
\(
\varepsilon_{\Lambda_k}\in\{\pm1\}
\)
such that
\[
\mathcal A^-_{\Lambda_k}\subset\ker\chi_{\Lambda_k},
\qquad
\chi_{\Lambda_k}(J)=\varepsilon_{\Lambda_k}
\quad
\forall J\in[\Xi^{(1)}_{\Lambda_k}].
\]
Assume in addition that \(\ker\chi_{\Lambda_k}\) admits a continuous symmetric coercive bilinear form
whose induced norm Hilbertizes it.
Choose any vector \(J^\sharp_{\Lambda_k}\in E^-_{\Lambda_k}\) satisfying
\[
\chi_{\Lambda_k}(J^\sharp_{\Lambda_k})=1,
\]
any scale \(\beta_{\Lambda_k}>0\), and any continuous symmetric coercive bilinear form
\[
\langle\cdot,\cdot\rangle^{\mathrm{aux},\chi}_{\Lambda_k}
:
\ker\chi_{\Lambda_k}\times\ker\chi_{\Lambda_k}\to\mathbb R
\]
whose induced norm Hilbertizes \(\ker\chi_{\Lambda_k}\).
Define
\begin{equation}
\mathsf B_{\Lambda_k}(F,G)
:=
\beta_{\Lambda_k}\,\chi_{\Lambda_k}(F)\chi_{\Lambda_k}(G)
+
\left\langle
F-\chi_{\Lambda_k}(F)J^\sharp_{\Lambda_k},
\,
G-\chi_{\Lambda_k}(G)J^\sharp_{\Lambda_k}
\right\rangle^{\mathrm{aux},\chi}_{\Lambda_k}.
\label{eq:B-minimal-transport}
\end{equation}
Then \(\mathsf B_{\Lambda_k}\) is a continuous symmetric positive definite bilinear form on
\(E^-_{\Lambda_k}\) whose induced norm Hilbertizes \(E^-_{\Lambda_k}\), and
\[
\mathsf B_{\Lambda_k}(J^\sharp_{\Lambda_k},K)=0
\qquad
\forall K\in\ker\chi_{\Lambda_k}.
\]
Hence \(\mathsf B_{\Lambda_k}(J^\sharp_{\Lambda_k},K)=0\) in particular for every
\(K\in\mathcal A^-_{\Lambda_k}\), and
\[
\mathsf B_{\Lambda_k}(J^\sharp_{\Lambda_k},J^\sharp_{\Lambda_k})
=
\beta_{\Lambda_k}.
\]
Moreover, for any representative \(J^\Xi_{\Lambda_k}\in[\Xi^{(1)}_{\Lambda_k}]\), the choice
\[
J^\sharp_{\Lambda_k}:=\varepsilon_{\Lambda_k}J^\Xi_{\Lambda_k}
\]
is admissible.
\end{proposition}

\begin{proof}
Because \(\chi_{\Lambda_k}(J^\sharp_{\Lambda_k})=1\), the linear map
\[
P^\chi_{\Lambda_k}(F)
:=
F-\chi_{\Lambda_k}(F)J^\sharp_{\Lambda_k}
\]
takes values in \(\ker\chi_{\Lambda_k}\):
\[
\chi_{\Lambda_k}\!\big(P^\chi_{\Lambda_k}(F)\big)
=
\chi_{\Lambda_k}(F)-\chi_{\Lambda_k}(F)\chi_{\Lambda_k}(J^\sharp_{\Lambda_k})
=
0.
\]
Thus every \(F\in E^-_{\Lambda_k}\) has the decomposition
\[
F=\chi_{\Lambda_k}(F)J^\sharp_{\Lambda_k}+P^\chi_{\Lambda_k}(F),
\qquad
P^\chi_{\Lambda_k}(F)\in\ker\chi_{\Lambda_k}.
\]
This decomposition is unique: if \(F=aJ^\sharp_{\Lambda_k}+K\) with \(K\in\ker\chi_{\Lambda_k}\),
applying \(\chi_{\Lambda_k}\) gives \(a=\chi_{\Lambda_k}(F)\).

Equation \eqref{eq:B-minimal-transport} is therefore well defined, continuous, and symmetric.
If \(\mathsf B_{\Lambda_k}(F,F)=0\), then
\[
\beta_{\Lambda_k}\chi_{\Lambda_k}(F)^2=0
\]
and
\[
\left\langle
P^\chi_{\Lambda_k}(F),P^\chi_{\Lambda_k}(F)
\right\rangle^{\mathrm{aux},\chi}_{\Lambda_k}
=
0.
\]
Since \(\beta_{\Lambda_k}>0\), we get \(\chi_{\Lambda_k}(F)=0\), hence
\(P^\chi_{\Lambda_k}(F)=F\). Coercivity on \(\ker\chi_{\Lambda_k}\) then implies \(F=0\), proving
positive definiteness.

The map
\[
F\longmapsto
\big(
\chi_{\Lambda_k}(F),P^\chi_{\Lambda_k}(F)
\big)
\]
is a continuous linear isomorphism from \(E^-_{\Lambda_k}\) onto
\(
\mathbb R\oplus\ker\chi_{\Lambda_k}
\),
with inverse
\[
(a,K)\longmapsto aJ^\sharp_{\Lambda_k}+K.
\]
Under this isomorphism, \(\mathsf B_{\Lambda_k}\) becomes the product Hilbert form
\[
(a,K),(b,L)\longmapsto
\beta_{\Lambda_k}ab+\langle K,L\rangle^{\mathrm{aux},\chi}_{\Lambda_k},
\]
so its induced norm Hilbertizes \(E^-_{\Lambda_k}\).

Finally, if \(K\in\ker\chi_{\Lambda_k}\), then \(P^\chi_{\Lambda_k}(J^\sharp_{\Lambda_k})=0\) and
\(\chi_{\Lambda_k}(K)=0\), hence
\[
\mathsf B_{\Lambda_k}(J^\sharp_{\Lambda_k},K)=0.
\]
Setting \(K=J^\sharp_{\Lambda_k}\) gives
\(
\mathsf B_{\Lambda_k}(J^\sharp_{\Lambda_k},J^\sharp_{\Lambda_k})=\beta_{\Lambda_k}
\).
If \(J^\Xi_{\Lambda_k}\in[\Xi^{(1)}_{\Lambda_k}]\), then
\(
\chi_{\Lambda_k}(J^\Xi_{\Lambda_k})=\varepsilon_{\Lambda_k}
\),
so
\(
J^\sharp_{\Lambda_k}:=\varepsilon_{\Lambda_k}J^\Xi_{\Lambda_k}
\)
satisfies
\(
\chi_{\Lambda_k}(J^\sharp_{\Lambda_k})=1
\).
\end{proof}

\begin{remark}[Calibration existence versus stiffness matching]
\label{rem:minimal-transport-calibration-scope}
\Cref{prop:minimal-transport-calibration} is a construction theorem, not an unconditional existence
theorem: it extends a previously chosen Hilbertizing form on \(\ker\chi_{\Lambda_k}\) to an
admissible bilinear form on the whole weak transported interface.
It does \emph{not} identify the scale
\(\beta_{\Lambda_k}\) with the split/leak stiffness \(\widetilde\alpha(\Lambda_k)\).
That physical identification belongs to a later strong or physical layer.
\end{remark}

\begin{theorem}[Rank-one descendant-sector upgrade]
\label{thm:active-response-transport-strong}
Assume \cref{thm:active-response-transport-minimal}, and fix an accepted tick \(k\) on the branch.
Suppose in addition that on the transported
descendant sector the projected response map
\[
F\longmapsto
\Pi^{\mathrm{dom}}_{\Lambda_k}\Pi_{\mathrm{phys},\Lambda_k}\mathcal R_{\Lambda_k}(F),
\qquad
\Pi^{\mathrm{dom}}_{\Lambda_k}
:=
\widehat{\delta C}^{\mathrm{LoG}}_{\Lambda_k}\otimes
\widehat{\delta C}^{\mathrm{LoG}}_{\Lambda_k},
\]
has kernel exactly \(\mathcal A^-_{\Lambda_k}\).
Then the stronger conclusions follow:
\begin{enumerate}[leftmargin=2em]
\item the descendant quotient is rank one,
  \[
  E^-_{\Lambda_k}/\mathcal A^-_{\Lambda_k}\cong \mathbb R;
  \]
\item the orientation kernel is exactly the ambiguity space,
  \[
  \ker\chi_{\Lambda_k}
  =
  \ker\rho^{\mathrm{dom}}_{\Lambda_k}
  =
  \mathcal A^-_{\Lambda_k},
  \]
  so the descendant sector satisfies the rank-one kernel-identification package with quotient
  generated by any vector \(J^{\mathrm{act}}_{\Lambda_k}\) satisfying
  \(\chi_{\Lambda_k}(J^{\mathrm{act}}_{\Lambda_k})=1\);
  \item after choosing any such \(J^{\mathrm{act}}_{\Lambda_k}\), any positive scale
  \(\beta_{\Lambda_k}\), and any auxiliary coercive Hilbertizing form on
  \(\ker\chi_{\Lambda_k}\), \cref{prop:minimal-transport-calibration} yields an admissible
  calibration form.
  Recovering the stiffness-matched formula of \cref{prop:char-calibration-extension} requires
  additional physical identifications, namely
  \(\chi_{\Lambda_k}=\rho_{\Lambda_k}\), \(\ker\chi_{\Lambda_k}=\mathcal T^{\mathrm{op}}_{\Lambda_k}\),
  agreement of the chosen normalized generator with the split/leak active generator, and
  \(\beta_{\Lambda_k}=\widetilde\alpha(\Lambda_k)\).
\end{enumerate}
\end{theorem}

\begin{proof}
By \cref{thm:active-response-transport-minimal}, the dominant line is defined at \(\Lambda_k\), the
scalar \(\nu_{\Lambda_k}\) is positive, and
\[
\chi_{\Lambda_k}
=
\nu_{\Lambda_k}^{-1}\rho^{\mathrm{dom}}_{\Lambda_k}.
\]
Hence
\(
\ker\chi_{\Lambda_k}=\ker\rho^{\mathrm{dom}}_{\Lambda_k}
\).

For every \(F\in E^-_{\Lambda_k}\), the definition of \(\rho^{\mathrm{dom}}_{\Lambda_k}\) gives
\[
\Pi^{\mathrm{dom}}_{\Lambda_k}\Pi_{\mathrm{phys},\Lambda_k}\mathcal R_{\Lambda_k}(F)
=
\rho^{\mathrm{dom}}_{\Lambda_k}(F)\,
\widehat{\delta C}^{\mathrm{LoG}}_{\Lambda_k}.
\]
Therefore the kernel of the projected response map is exactly
\(
\ker\rho^{\mathrm{dom}}_{\Lambda_k}
\).
By the additional hypothesis, that kernel is \(\mathcal A^-_{\Lambda_k}\), proving item (2).

By \cref{thm:active-response-transport-minimal}, there exists
\(
J^{\mathrm{desc}}_{\Lambda_k}\in[\Xi^{(1)}_{\Lambda_k}]
\)
with
\(
\rho^{\mathrm{dom}}_{\Lambda_k}(J^{\mathrm{desc}}_{\Lambda_k})\neq 0
\),
hence the projected response map is nonzero.
Its image is automatically contained in the one-dimensional line
\(
L^{\mathrm{dom}}_{\Lambda_k}
\),
so the first isomorphism theorem gives
\[
E^-_{\Lambda_k}/\mathcal A^-_{\Lambda_k}
\cong
\operatorname{im}\!\big(
\Pi^{\mathrm{dom}}_{\Lambda_k}\Pi_{\mathrm{phys},\Lambda_k}\mathcal R_{\Lambda_k}
\big)
\cong
\mathbb R,
\]
which is item (1).
Any vector \(J^{\mathrm{act}}_{\Lambda_k}\) with \(\chi_{\Lambda_k}(J^{\mathrm{act}}_{\Lambda_k})=1\)
lies outside \(\ker\chi_{\Lambda_k}=\mathcal A^-_{\Lambda_k}\), so its quotient class generates this
one-dimensional quotient, completing item (2).

Item (3) is now an application of \cref{prop:minimal-transport-calibration} with the weak-interface
data supplied by \cref{thm:active-response-transport-minimal} and the chosen normalized vector
\(J^{\mathrm{act}}_{\Lambda_k}\).
The final sentence identifies when this weak calibration construction coincides with the
stiffness-matched formula of \cref{prop:char-calibration-extension}.
\end{proof}

\ifdefined\UOMTransportFragmentLoaded\else
\end{document}
\fi
%
  }%
}

\appendix
\section{Appendix: BY kernel discretization and scaling with $\Delta I$}
\label{app:by-kernel}

The discretized BY kernel constructs a welded time spectrum $\phi(u)$ and a welded spatial spectrum on $S^3$
harmonics; it then chooses an effective amplitude $A_{\rm eff}$ so that
$A_{\rm eff}^2\propto \Delta I/G(s_t,s_x;\Lambda)$ where $G$ is a (no-amplitude) throughput. The integrated BY penalty is
quadratic in the signal and thus linear in $\Delta I$.
This is the phenomenological origin of the manuscript factorization $\Theta_k=K(\bar\Lambda_k)\Delta I_k$.

\section{Appendix: Minimal bibliography}
\label{app:biblio}

\begin{thebibliography}{99}
\bibitem{Stinespring1955}
W.\ F.\ Stinespring,
``Positive Functions on $C^\ast$-Algebras,''
\emph{Proc.\ Amer.\ Math.\ Soc.}\ \textbf{6} (1955).
DOI: 10.1090/S0002-9939-1955-0069403-4.

\bibitem{Lindblad1976}
G.\ Lindblad,
``On the generators of quantum dynamical semigroups,''
\emph{Commun.\ Math.\ Phys.}\ \textbf{48} (1976).
DOI: 10.1007/BF01608499.

\bibitem{GKS1976}
V.\ Gorini, A.\ Kossakowski, and E.\ C.\ G.\ Sudarshan,
``Completely positive dynamical semigroups of $N$-level systems,''
\emph{J.\ Math.\ Phys.}\ \textbf{17} (1976).
DOI: 10.1063/1.522979.

\bibitem{Spohn1978}
H.\ Spohn,
``Entropy production for quantum dynamical semigroups,''
\emph{J.\ Math.\ Phys.}\ \textbf{19} (1978).
DOI: 10.1063/1.523789.

\bibitem{BrownYork1993}
J.\ D.\ Brown and J.\ W.\ York,
``Quasilocal energy and conserved charges derived from the gravitational action,''
\emph{Phys.\ Rev.\ D}\ \textbf{47} (1993).
DOI: 10.1103/PhysRevD.47.1407.

\bibitem{Bhatia1997}
R.\ Bhatia,
\emph{Matrix Analysis},
Springer (1997).
DOI: 10.1007/978-1-4612-0653-8.

\bibitem{PGWP2006}
D.\ P\'erez-Garc\'{\i}a, M.\ M.\ Wolf, D.\ Petz, and M.\ B.\ Ruskai,
``Contractivity of positive and trace-preserving maps under $L_p$ norms,''
\emph{J.\ Math.\ Phys.}\ \textbf{47} (2006).
DOI: 10.1063/1.2218675.
\end{thebibliography}

\end{document}
